\documentclass[11pt]{article}
\usepackage[T1]{fontenc}
\usepackage{amsmath,amssymb}
\usepackage{mathtools}
\usepackage[margin=1in]{geometry}
\usepackage{enumitem}
\usepackage{amsthm}
\usepackage[hidelinks]{hyperref}
\hypersetup{linktocpage}
\newtheoremstyle{ntstyle}{\medskipamount}{\medskipamount}{\upshape\statementbreaks}{0pt}{\bfseries}{:}{ }{\thmname{#1}\thmnumber{ #2}\thmnote{ (#3)}}
\theoremstyle{ntstyle}
\newtheorem{theorem}{Theorem}[subsection]
\newtheorem{definition}{Definition}[subsection]
\newcommand{\statementitems}{\mbox{}}
\makeatletter
\renewcommand*\l@subsection{\@dottedtocline{2}{1.5em}{3.3em}}
\raggedbottom
\clubpenalty\@M \@clubpenalty\clubpenalty \widowpenalty\@M \displaywidowpenalty\@M
\newcommand*\statementbreaks{{\@@par}\if@nobreak\else\addpenalty\@beginparpenalty\fi
  \@beginparpenalty\@M
  \dimen@\pagegoal \advance\dimen@-\pagetotal
  \ifdim\dimen@<10\baselineskip
    \def\nt@inside{\@M}\@endparpenalty\@M
  \else
    \def\nt@inside{500 }\ifnum\@endparpenalty<\@highpenalty \@endparpenalty\@highpenalty \fi
  \fi
  \interlinepenalty\nt@inside \postdisplaypenalty\nt@inside \@itempenalty\nt@inside
  \let\nt@par\@par \@setpar{\nt@par\ifvmode\addpenalty\nt@inside\fi}\let\nt@esphack\@esphack \def\@esphack{\nt@esphack\ifvmode\nobreak\fi}}
\let\nt@closedisplay\]
\def\nt@guardedclose{\@ifnextchar\par{\nt@closeafterpar}{\@ifnextchar\end{\postdisplaypenalty\@M\nt@closedisplay}{\nt@closedisplay}}}
\def\nt@closeafterpar\par{\@ifnextchar\par{\nt@closeafterpar}{\@ifnextchar\end{\postdisplaypenalty\@M\nt@closedisplay\par}{\nt@closedisplay\par}}}
\renewenvironment{proof}[1][\proofname]{\par
  \pushQED{\qed}\normalfont\topsep\medskipamount\edef\nt@afterproof{\@endparpenalty=\the\@endparpenalty\relax}\@beginparpenalty\@M\@endparpenalty\@M\let\]\nt@guardedclose\trivlist\item[]\textit{#1\@addpunct{.}} \ignorespaces}{\ifvmode\addpenalty\@M\fi\popQED\nt@afterproof\endtrivlist\@endpefalse}
\makeatother
\newenvironment{altproof}[1][\proofname]{\begin{proof}[#1]\renewcommand{\qedsymbol}{}}{\end{proof}}
\emergencystretch=1em
\allowdisplaybreaks\newcommand{\leg}[2]{\left(\frac{#1}{#2}\right)}
\newcommand{\pleg}[3]{\left(\frac{#1}{#2}\right)_{#3}}
\newcommand{\Primes}{\mathbb{P}}
\newcommand{\ZmodUnits}[1]{\mathbb{Z}_{#1}^{\times}}
\DeclareMathOperator{\lcm}{lcm}
\DeclareMathOperator{\ind}{ind}
\DeclareMathOperator{\tr}{tr}
\DeclareMathOperator{\id}{id}
\DeclareMathOperator{\area}{area}
\DeclareMathOperator{\disc}{disc}
\DeclareMathOperator{\sgn}{sgn}
\setlength{\parindent}{0pt}
\setlength{\parskip}{0.5em}
\setlist{topsep=0.2em,itemsep=0.2em}

\title{Elementary Number Theory Notes}
\author{}
\date{}

\begin{document}
\maketitle
\newpage
\begingroup
\setlength{\parskip}{0pt}
\tableofcontents
\endgroup
\newpage

\section*{Notation and Conventions}
\addcontentsline{toc}{section}{Notation and Conventions}

\medskip
\noindent Unless stated otherwise, the following conventions apply throughout.

\noindent \textbf{Number systems.} The symbols $\mathbb{Z},\mathbb{Q},\mathbb{R},\mathbb{C}$ denote the integers, rationals, reals, and complex numbers, respectively. We write $\mathbb{Z}^+=\{1,2,3,\ldots\}$, $\mathbb{Z}_{\ge0}=\{0,1,2,\ldots\}$, and $\mathbb{Z}^-=\{-1,-2,-3,\ldots\}$; to avoid ambiguity, the symbol $\mathbb{N}$ is not used. For a prime $p$, a nonzero integer $m$, and $a\in\mathbb{Z}_{\ge0}$, we write $p^{a}\parallel m$ to mean $p^{a}\mid m$ but $p^{a+1}\nmid m$, that is, $p^{a}$ is the exact power of $p$ dividing $m$. When a product is indexed by the condition $p^a\parallel m$, only prime powers with $a\ge1$ are included. Without further qualification, ``prime'' means a positive rational prime; primes in other rings are named explicitly. We write $\Primes=\{2,3,5,7,\ldots\}$ for the set of primes.

\noindent \textbf{Assumed background.} We freely use, without separate proof, a few facts that are logically prior to the theory developed here: the arithmetic and order properties of $\mathbb{Z}$, $\mathbb{Q}$, and $\mathbb{R}$ as ordered rings and fields, the completeness of $\mathbb{R}$, and the Archimedean property, so that $\lfloor x\rfloor$ exists for every $x\in\mathbb{R}$; the well-ordering principle and (strong) induction on $\mathbb{Z}_{\ge0}$; the pigeonhole principle, both in its finite form and in the form that if infinitely many objects are placed into finitely many classes, then some class contains infinitely many of them; elementary finite counting, including the facts that a map between two finite sets with the same number of elements is injective if and only if it is surjective, and that $\binom{n}{k}$ is the number of $k$-element subsets of an $n$-element set; the binomial theorem and inclusion--exclusion; elementary complex arithmetic, including complex conjugation, the absolute value and the triangle inequality, Euler's formula $e^{i\theta}=\cos\theta+i\sin\theta$ for real $\theta$ (so that $|e^{i\theta}|=1$), and the fact that $e^{2\pi i x}=1$ if and only if $x\in\mathbb{Z}$; basic linear algebra over a field (matrices, triangularity, invertibility, trace, and determinant); the formal derivative of a polynomial and its use in detecting repeated roots; and Newton's identities. In the analytic sections we use standard real analysis (limits, continuity, differentiation, and integration of functions of one real variable, the mean value theorem, the AM--GM inequality, rearrangement of absolutely convergent series, and interchange of the order of summation in an absolutely convergent double series), together with the elementary theory of the functions $\exp$, $\ln$, $\sin$, $\cos$, $\tan$, and $\cot$: their derivatives, monotonicity on the usual intervals, and power series, the angle-addition formulas, radian measure with $\pi$ the least positive zero of $\sin$, and the positivity of $\sin$ on $(0,\pi)$. Where areas occur (in the proof of Theorem~\ref{thm:5.7.3} and in the sections on lattice points and areas) we use the elementary measure theory and geometry of the plane: area is defined on the regions considered there, including convex sets, and is nonnegative, monotone, finitely and countably additive, and invariant under translations; an invertible linear map $T$ scales area by $|\det T|$; a circular sector of radius $1$ and angle $\theta\in(0,2\pi]$ has area $\theta/2$, so a disk of radius $r$ has area $\pi r^2$; a simple polygon bounds a well-defined interior region (the Jordan curve theorem for polygons), its interior angle at a vertex of its convex hull is less than $\pi$, and cutting it along a polygonal path that joins two of its boundary points through its interior divides it into two simple polygons. Asymptotic notation and the notions of average order and density are recalled where they are first used.

\noindent \textbf{Algebraic vocabulary.} A few standard terms are used in passing before the subsections in which they are formally introduced: group, abelian group, subgroup, cyclic group, generator, and the order of an element (Definitions~\ref{def:6.6.1}, \ref{def:6.6.2}, and~\ref{def:6.6.3}); homomorphism, isomorphism, automorphism, kernel, coset, and commutative ring with identity (Definition~\ref{def:6.6.4}); and field, finite field, and characteristic (Definition~\ref{def:6.10.1}). Earlier arguments use the elementary laws of these structures, notably inverses in groups and fields; the later subsections state the definitions explicitly and develop the structural results needed in what follows.

\noindent \textbf{Numbering.} Results are numbered in the form section.subsection.item; definitions and theorems have separate counters.

\newpage

\section{Divisibility}

\subsection{Basic Divisibility Properties}

\begin{definition}[divisibility]\label{def:1.1.1} Let $d, n \in
		\mathbb{Z}$. We say that $d$ divides $n$, written $d \mid n$, if there
	exists $k \in \mathbb{Z}$ such that $n = dk$; in this case we also
	say ``$n$ is divisible by $d$'' or ``$d$ is a
	divisor of $n$''. Otherwise, we write $d \nmid n$ to mean that $d \mid n$ does not hold. Every integer divides $0$ (take $k
		= 0$), so $d \mid 0$ for all $d \in \mathbb{Z}$; in particular $0
		\mid 0$. In contrast, $0 \mid n$ if and only if $n = 0$. The trivial
	divisibility relations $1 \mid n$ and $n \mid n$ hold for every $n \in \mathbb{Z}$.

\end{definition}
\begin{theorem}\label{thm:1.1.1} If $a \mid b$ and $b
		\neq 0$, then $|a| \le |b|$.

\end{theorem}

\begin{proof} Write $b = ak$ for some $k\in\mathbb{Z}$. Since $b \neq 0$, both $a$ and $k$
	are nonzero, so $|k| \ge 1$ and $|b| = |a|\,|k| \ge |a|$.

\end{proof}
\begin{theorem}\label{thm:1.1.2}\statementitems

	\begin{enumerate}[label={\alph*)}]
		\item (cancellation) If $c \neq 0$ and $ac \mid bc$, then $a \mid
			      b$ $(a, b, c \in \mathbb{Z})$.
		\item (transitivity) If $a \mid b$ and $b \mid c$, then $a \mid c$
		      $(a, b, c \in \mathbb{Z})$.
		\item (linear combinations) If $a \mid b_i$ for all $1 \le i \le
			      m$, then $a \mid \sum_{i=1}^{m} w_i b_i$ for any weights $w_i \in
			      \mathbb{Z}$ $(a, b_i \in \mathbb{Z})$.
		\item (mutual divisibility) If $a \mid b$ and $b \mid a$, then $a = \pm
			      b$ $(a, b \in \mathbb{Z})$; in particular, if $a, b \in
			      \mathbb{Z}^+$, then $a = b$.
		\item The divisibility relation $a \mid b$ is a partial order on $\mathbb{Z}^+$.
	\end{enumerate}

\end{theorem}

\begin{proof}

	\begin{enumerate}[label={\alph*)}]
		\item Suppose $c \neq 0$ and $ac \mid bc$. Then $bc = k(ac)$ for
		      some $k \in \mathbb{Z}$, that is, $(b - ka)c = 0$. A product of two integers is zero only when one of its
		      factors is zero. Since $c \neq 0$, we get $b - ka = 0$, that is, $b = ka$, and hence $a \mid b$.
		\item There exist $k_1, k_2 \in \mathbb{Z}$ such that $b = k_1 a$
		      and $c = k_2 b$, so $c = (k_1 k_2) a$; since $k_1 k_2 \in
			      \mathbb{Z}$, we have $a \mid c$.
		\item For each $i$ there exists $q_i \in \mathbb{Z}$ such that
		      $b_i = a q_i$. Then $\sum_{i=1}^{m} w_i b_i = \sum_{i=1}^{m} w_i a
			      q_i = \left(\sum_{i=1}^{m} w_i q_i\right) a$, and since
		      $\sum_{i=1}^{m} w_i q_i \in \mathbb{Z}$, we have $a \mid
			      \sum_{i=1}^{m} w_i b_i$.
		\item If $a = 0$, then $a \mid b$ implies $b = 0 = \pm a$;
		      the case $b = 0$ is symmetric. If $a, b \neq 0$, Theorem~\ref{thm:1.1.1} gives
		      $|a| \le |b|$ and $|b| \le |a|$, so $|a| = |b|$ and $a = \pm b$.
		      If $a, b \in \mathbb{Z}^+$, then both are positive, so $a = b$.
		\item Reflexivity ($a \mid a$), transitivity (b), and antisymmetry
		      (d) together show that the divisibility relation $a \mid b$ is a partial order on
		      $\mathbb{Z}^+$.
	\end{enumerate}

\end{proof}
\subsection{The Division Algorithm}

\begin{definition}[floor]\label{def:1.2.1} For $x\in\mathbb{R}$, the floor $\lfloor x\rfloor$ is the greatest integer
	not exceeding $x$.

\end{definition}
\begin{definition}[ceiling]\label{def:1.2.2} For $x\in\mathbb{R}$,
	the ceiling $\lceil x\rceil$ is the least integer greater than or equal to $x$.

\end{definition}
\begin{definition}[fractional part]\label{def:1.2.3} For
	$x\in\mathbb{R}$, the fractional part of $x$ is
	$\{x\}=x-\lfloor x\rfloor\in[0,1)$.

\end{definition}
\begin{definition}[distance to the nearest integer]\label{def:1.2.4}
	For $x\in\mathbb{R}$, put
	$\lVert x\rVert=\min_{k\in\mathbb{Z}}|x-k|$.

\end{definition}
\begin{theorem}[division algorithm]\label{thm:1.2.1} Given integers
	$a, b$ with $b > 0$, there exist unique integers $q, r \in \mathbb{Z}$
	with $0 \le r < b$ such that $a = qb + r$ (the division form).

\end{theorem}

\begin{proof} We first prove existence for $a\ge0$ by induction. For $a = 0$, we have $a = 0
		\cdot b + 0$, so $0$ can be written in division form. Now
	assume the claim holds for $a = k$, with $k \in \mathbb{Z}_{\ge0}$ fixed, that is, $k = qb + r$ in division form. When $a = k+1$,
	if $r + 1 < b$, then $k+1 = qb + (r+1)$ is in division form; if $r+1
		= b$, then $k+1 = (q+1)b + 0$ is in division form. By induction,
	every $a \ge 0$ can be written in division form. When $a < 0$, we
	have $-a > 0$, so there exist $q, r \in \mathbb{Z}$ with $0 \le r <
		b$ such that $-a = qb + r$. If $r = 0$, then $a = (-q)b + 0$ is in
	division form; if $0 < r < b$, then $a = (-q-1)b + (b - r)$ is in
	division form. Therefore, for all $a, b \in \mathbb{Z}$ with $b > 0$,
	the integer $a$ can be written in division form. For uniqueness, suppose $a = qb
		+ r = q'b + r'$ with $0 \le r, r' < b$; then $b \mid (r-r')$. If $r-r'=0$, then $r=r'$ immediately. Otherwise, Theorem~\ref{thm:1.1.1} gives $b\le |r-r'|$, contradicting $|r-r'|<b$, which holds because $0\le r,r'<b$ gives $-b<r-r'<b$. Hence $r=r'$ in all cases; then $(q-q')b=r'-r=0$ with $b\neq0$, so $q=q'$.

\end{proof}
\newpage

\section{Primes, the Euclidean Algorithm, and Diophantine Equations}

\subsection{Primes and Their Infinitude}

\begin{definition}[prime]\label{def:2.1.1} An integer $p \ge 2$
	is a prime if its only positive divisors are $1$ and $p$ itself.

\end{definition}
\begin{definition}[composite]\label{def:2.1.2} An integer $n \ge 2$ is a
	composite number if $n$ is not prime; equivalently, $n$ can be written as $n = ab$ with $a, b \in \mathbb{Z}^+$ and $1 < a, b < n$.

\end{definition}
\begin{theorem}\label{thm:2.1.1} Every integer greater than $1$ has
	a prime factor.

\end{theorem}

\begin{proof} Suppose, for contradiction, that some
	integer $n > 1$ has no prime factor, and let $S$ be the set of all
	such integers. By the well-ordering principle, $S$ has a smallest
	element $m$. If $m$ is prime, then $m$ has a prime factor, namely
	itself. If $m$ is composite, then $m$ has a factor
	$k$ with $1 < k < m$, so $k \notin S$; hence $k$ has a prime factor,
	and since $k \mid m$, that prime also divides $m$. Either way, $m$ has
	a prime factor, a contradiction. Therefore every integer greater than
	$1$ has a prime factor.

\end{proof}
\begin{theorem}\label{thm:2.1.2} There are infinitely many primes.

\end{theorem}

\begin{proof} Assume there are only finitely many primes, say $p_1,
		p_2, \ldots, p_k$. Let $N = \prod_{i=1}^{k} p_i + 1$. The integer $N > 1$
	has a prime factor (Theorem~\ref{thm:2.1.1}), so $p_i \mid N$ for some $1 \le
		i \le k$. Since $p_i \mid \prod_{j=1}^{k} p_j$, we get $p_i \mid N -
		\prod_{j=1}^{k} p_j$, that is, $p_i \mid 1$, which is impossible since
	$p_i\ge2$. Thus there are infinitely many primes.

\end{proof}
\begin{theorem}\label{thm:2.1.3}
	There are infinitely many primes of the form $4k+3$, and infinitely many
	primes of the form $6k+5$.

\end{theorem}

\begin{proof} By Theorem~\ref{thm:2.1.1}, repeated division by prime
	factors expresses every integer greater than $1$ as a product of primes:
	the remaining positive integer decreases at each step and so must
	eventually reach $1$.

	Suppose $p_1,\ldots,p_k$ were all the primes
	of the form $4m+3$, and set $N=4p_1\cdots p_k-1$. Then $N>1$; it is odd and
	$N=4(p_1\cdots p_k-1)+3$ is itself of the form $4m+3$. Every prime factor
	of $N$ is odd, hence of the form $4m+1$ or $4m+3$. A product of factors
	of the form $4m+1$ is again of that form, since
	$(4a+1)(4b+1)=4(4ab+a+b)+1$; so if every prime factor of $N$ were of the
	form $4m+1$, then $N$ would be of the form $4m+1$, a contradiction. Hence
	some prime $q$ of the form $4m+3$ divides $N$. But $q\neq p_i$ for every $i$: since $p_i\mid N+1$, the relation $p_i\mid N$ would force $p_i\mid1$. This contradicts the assumption that the list contains all such primes.

	For the second class, suppose $p_1,\ldots,p_k$ were all the primes of the
	form $6m+5$, and set $N=6p_1\cdots p_k-1=6(p_1\cdots p_k-1)+5$. Neither
	$2$ nor $3$ divides $N$: such a prime would divide both $N$ and
	$N+1=6p_1\cdots p_k$, hence would divide $1$. So every prime factor of
	$N$ exceeds $3$ and is therefore of the form $6m+1$ or $6m+5$. A product
	of factors of the form $6m+1$ is again of that form, since
	$(6a+1)(6b+1)=6(6ab+a+b)+1$; so if every prime factor of $N$ were of the
	form $6m+1$, then $N$ would be of the form $6m+1$, a contradiction. Hence
	some prime $q$ of the form $6m+5$ divides $N$, and as before $q\neq p_i$
	for all $i$.\qedhere

\end{proof}
\subsection{GCD, LCM, and the Euclidean Algorithm}

\begin{definition}[gcd and lcm]\label{def:2.2.1} For integers $a,
		b$ not both zero, the greatest common divisor $\gcd(a, b)$ is the greatest positive integer dividing both $a$ and $b$; for nonzero integers $a, b$,
	the least common multiple $\lcm(a, b)$ is the least positive integer that is
	divisible by both $a$ and $b$. By convention, $\gcd(0, 0)=0$ and $\lcm(0, n)=\lcm(n, 0)=0$
	for every integer $n$; in particular, $\lcm(0, 0)=0$.

\end{definition}
\begin{definition}[coprime]\label{def:2.2.2} Integers $a$ and $b$ are
	coprime if $\gcd(a, b) = 1$.

\end{definition}
\begin{definition}[gcd of several integers and pairwise coprime integers]\label{def:2.2.3} Let
	$n\in\mathbb{Z}^+$ and let $a_1, \ldots, a_n$ be integers. If they are not all zero,
	their greatest common divisor $\gcd(a_1, \ldots, a_n)$ is the greatest positive
	integer that divides every $a_i$; by convention, $\gcd(0, \ldots, 0)=0$. The
	integers $a_1, \ldots, a_n$ are coprime if $\gcd(a_1, \ldots, a_n)=1$, and pairwise
	coprime if $\gcd(a_i, a_j)=1$ for all $i\neq j$; for $n=2$, both notions reduce
	to Definition~\ref{def:2.2.2}. For $n\ge2$, pairwise coprime integers are coprime,
	because every positive common divisor of all the $a_i$ divides both $a_1$ and
	$a_2$ and hence equals $1$. The converse fails: $\gcd(6, 10, 15)=1$, although
	no two of $6$, $10$, $15$ are coprime.

\end{definition}
\begin{theorem}\label{thm:2.2.1}
	For integers $a, b$ not both zero, $\gcd(a, b) = \gcd(a, a+b)$;
	equivalently, $\gcd(a, b) = \gcd(a+b, b)$ and $\gcd(a, b) = \gcd(a-b, b)$.

\end{theorem}

\begin{proof} Every common divisor of $a$ and $b$ divides
	$a+b$ (Theorem~\ref{thm:1.1.2}(c)), hence is a common divisor of $a$ and $a+b$;
	conversely, every common divisor of $a$ and $a+b$ divides $(a+b) - a =
		b$ (Theorem~\ref{thm:1.1.2}(c)), hence is a common divisor of $a$ and $b$. Thus
	the pairs $\{a, b\}$ and $\{a, a+b\}$ have exactly the same common
	divisors, and in particular the same greatest common divisor, so $\gcd(a, b) =
		\gcd(a, a+b)$. The identity $\gcd(a, b) = \gcd(a+b, b)$ follows from the
	symmetry $\gcd(x, y) = \gcd(y, x)$, and applying the former identity with $a$ replaced by
	$a-b$ gives $\gcd(a-b, b) = \gcd(a, b)$.

\end{proof}
\begin{theorem}[Euclidean algorithm]\label{thm:2.2.2} Let $a, b \in
		\mathbb{Z}^+$ and start with the pair $(a,b)$. By the division
	algorithm, compute the unique integers $q, r$ such that $a = qb + r$
	with $0 \le r < b$. If $r = 0$, output $b$; otherwise repeat this
	step with $(b, r)$ in place of $(a,b)$. The final output is $\gcd(a, b)$.

\end{theorem}

\begin{proof} The value of $b$ strictly decreases at each
	step, so the algorithm must terminate. At each step, $\gcd(a, b) = \gcd(b, r)$: since $r = a - qb$, applying
	Theorem~\ref{thm:2.2.1} repeatedly to subtract $b$ from the first argument gives
	$\gcd(a, b) = \gcd(a-b, b) = \cdots = \gcd(a-qb, b) = \gcd(r, b) =
		\gcd(b, r)$. Therefore each step preserves the gcd. In the final step, the current
	pair is $(a', b')$ with remainder $0$, that is, $a' = q b' + 0$ for some
	quotient $q$, so $\gcd(a, b) = \gcd(a', b') = \gcd(b', 0) = b'$,
	which is precisely the final output.

\end{proof}
\begin{theorem}\label{thm:2.2.3} The least common multiple of
	$a,b\in\mathbb{Z}^+$ is a divisor of every common multiple
	$m\in\mathbb{Z}$ of $a$ and $b$.

\end{theorem}

\begin{proof} Since $a,b>0$, their product $ab$ is a positive common multiple.
	Thus the well-ordering principle gives a least positive common multiple
	$L=\lcm(a,b)$. Let $m\in \mathbb{Z}^+$ be any common multiple of $a$
	and $b$, so $a \mid m$ and $b \mid m$. As $L>0$, the division
	algorithm gives $m=qL+r$ with $q,r\in\mathbb{Z}$, $0\le r<L$.
	Since $a\mid m$ and $a\mid L$, we have $a\mid (m-qL)=r$; similarly
	$b\mid r$. Hence $r$ is also a common multiple of $a$ and $b$. If
	$1\le r<L$, then $r$ would be a positive common multiple smaller than
	$L$, contradicting the definition of $L$. Therefore $r=0$, so $L\mid
		m$. For $m=0$, clearly $L\mid 0$. If $m\in\mathbb{Z}^-$ is a common
	multiple, then $-m\in\mathbb{Z}^+$ is also a common multiple, so
	$L\mid (-m)$, hence $L\mid m$. Thus the least common multiple divides every common multiple.

\end{proof}
\begin{theorem}[linear combination criterion]\label{thm:2.2.4} For
	$a, b, d \in \mathbb{Z}$ with $a,b$ not both zero, $\gcd(a, b) \mid d$ if and only
	if there exist $x, y \in \mathbb{Z}$ such that $ax+by=d$.

\end{theorem}

\begin{proof} First we prove that if $\gcd(a, b) \mid d$,
	then there exist $x, y \in \mathbb{Z}$ such that $ax+by=d$. We may
	assume $a, b > 0$: replacing $a, b$ by $|a|, |b|$ changes neither
	$\gcd(a, b)$ nor the set of values $ax+by$ attainable with $x, y \in
		\mathbb{Z}$, and the cases where $a$ or $b$ is $0$ are immediate. Let
	$r_0 = \max\{a, b\}$ and $r_1 = \min\{a, b\}$, and for each $i \ge 1$ with $r_i>0$ let
	$r_{i+1} = r_{i-1} \bmod r_i$ (here $u \bmod v$ denotes the least nonnegative remainder of $u$ on division by $v$; the indices increase along the algorithm, while the back-substitution below proceeds from $r_n$ back to $r_0$ and $r_1$). Let $n$ be the smallest positive
	integer such that $r_{n+1} = 0$; by the Euclidean algorithm, $r_n =
		\gcd(a, b)$. We show that for all $0 \le j \le n$, there
	exist $x_j, y_j \in \mathbb{Z}$ such that $r_{n-j} x_j + r_{n-j+1}
		y_j = \gcd(a, b)$. We prove this by induction on $j$. When $j=0$, we have $r_n = \gcd(a, b)$
	and $r_{n+1}=0$, so $x_0=1$ and $y_0=0$ satisfy the equation. Now assume
	that for $j=k$, where $0\le k<n$, there exist $x_k, y_k \in \mathbb{Z}$
	such that $r_{n-k} x_k + r_{n-k+1} y_k = \gcd(a, b)$. Since $r_{n-k-1}
		= q_{n-k} r_{n-k} + r_{n-k+1}$ for some $q_{n-k} \in \mathbb{Z}$ (the
	quotient at this step, indexed to match the dividend $r_{n-k-1}$ and
	the divisor $r_{n-k}$), we have $r_{n-k+1} = r_{n-k-1} - q_{n-k} r_{n-k}$. Therefore
	\[
		\begin{aligned} r_{n-k} x_k + (r_{n-k-1} - q_{n-k} r_{n-k}) y_k & =
		                \gcd(a, b)                                          \\ r_{n-k-1} y_k + r_{n-k}(x_k - q_{n-k} y_k) & = \gcd(a, b).
		\end{aligned}
	\]

	Setting $x_{k+1} = y_k$ and $y_{k+1} = x_k - q_{n-k} y_k$ gives integers
	$x_{k+1}, y_{k+1} \in \mathbb{Z}$ with $r_{n-(k+1)}
		x_{k+1} + r_{n-(k+1)+1} y_{k+1} = \gcd(a, b)$, which is the claim for
	$j=k+1$. By induction, the integers $x'=x_n$ and $y'=y_n$ satisfy
	$r_0x'+r_1y'=\gcd(a, b)$. If $a\ge b$, set $(X,Y)=(x',y')$; if $b>a$, set $(X,Y)=(y',x')$. Then $aX+bY=\gcd(a, b)$. Since there exists $t\in\mathbb{Z}$ such that $d=t\gcd(a, b)$, taking $x=tX$ and $y=tY$
	gives $ax+by=d$. Conversely, if there exist $x, y \in \mathbb{Z}$
	such that $ax+by=d$, then since $\gcd(a, b) \mid a$ and $\gcd(a, b)
		\mid b$, we have $\gcd(a, b) \mid ax+by$, that is, $\gcd(a, b) \mid d$.

\end{proof}
\begin{theorem}[B{\'e}zout's identity]\label{thm:2.2.5} For integers
	$a, b$ not both zero, there exist $x, y \in \mathbb{Z}$ such that
	$ax + by = \gcd(a, b)$. More precisely,
	\[
		\{ax+by \mid x,y\in\mathbb{Z}\}=\gcd(a, b)\,\mathbb{Z};
	\]
	in particular, $\gcd(a, b)$ is the least positive integer of the form
	$ax+by$ with $x,y\in\mathbb{Z}$.

\end{theorem}

\begin{proof} This is the special case $d = \gcd(a, b)$ of
	the linear combination criterion (Theorem~\ref{thm:2.2.4}): since $\gcd(a, b)
		\mid \gcd(a, b)$, that criterion provides $x, y \in \mathbb{Z}$ with
	$ax + by = \gcd(a, b)$.
	More generally, the same criterion shows that an integer $d$ can be written as
	$ax+by$ with $x,y\in\mathbb{Z}$ if and only if $\gcd(a, b)\mid d$, so the set
	of such integers is exactly the set of multiples of $\gcd(a, b)$. Since
	$\gcd(a, b)$ is positive (Definition~\ref{def:2.2.1}), its least positive
	multiple is $\gcd(a, b)$ itself.

\end{proof}
\begin{theorem}\label{thm:2.2.6} For integers $a, b$ not both zero, $\gcd(a, b)$ is divisible by every common divisor of $a$ and $b$.

\end{theorem}

\begin{proof} By B{\'e}zout's identity (Theorem~\ref{thm:2.2.5}), $ax+by = \gcd(a, b)$ for some $x, y \in
		\mathbb{Z}$. Let $d$ be a common
	divisor of $a$ and $b$. Since $d \mid a$ and $d \mid b$, we have $d
		\mid (ax+by) = \gcd(a, b)$.

\end{proof}
\begin{theorem}\label{thm:2.2.7} Let $a,b$ be distinct positive integers and
	$d=\gcd(a, b)$. Then $d\le\min\{a, b\}$, $d\le\tfrac{1}{2}\max\{a, b\}$,
	$d\le|a-b|$, and $d\le\tfrac{a+b}{3}$.

\end{theorem}

\begin{proof} Since $d=\gcd(a, b)$ divides both $a$ and
	$b$, each of $a$ and $b$ is a positive multiple of $d$ and hence at least
		$d$; therefore $d\le\min\{a, b\}$. For the other three bounds, assume
	without loss of generality that $a>b$ (which is possible since $a\neq b$), and
	write $a=d\alpha$ and $b=d\beta$, where $\alpha=a/d$ and
	$\beta=b/d$ are positive integers. From $a>b$ we get
	$\alpha>\beta\ge1$, so $\alpha\ge2$; hence
	$\max\{a, b\}=a=d\alpha\ge2d$, giving $d\le\tfrac{1}{2}\max\{a, b\}$. Next,
	$d\mid a$ and $d\mid b$ yield $d\mid(a-b)$, and since
	$a-b>0$ this positive multiple of $d$ is at least $d$, so $d\le
		a-b=|a-b|$. Finally, $a+b=d(\alpha+\beta)$ with $\alpha\ge2$ and
	$\beta\ge1$ gives $\alpha+\beta\ge3$, hence $a+b\ge3d$, that is, $d\le\tfrac{a+b}{3}$.

\end{proof}
\begin{theorem}[Euclid's lemma]\label{thm:2.2.8} If $a$ and $b$ are coprime and
	$a \mid bc$, then $a \mid c$.

\end{theorem}

\begin{proof} Since $\gcd(a, b) = 1$, B{\'e}zout's identity (Theorem~\ref{thm:2.2.5})
	gives $ax + by = 1$ for some $x, y \in \mathbb{Z}$.
	Multiplying this identity by $c$ yields $acx+bcy=c$; since $a \mid acx$ and $a
		\mid bcy$ (the latter because $a \mid bc$), we have $a \mid c$.

\end{proof}
\begin{theorem}\label{thm:2.2.9} If $\gcd(a, b) = \gcd(a, c) = 1$,
	then $\gcd(a, bc)=1$.

\end{theorem}

\begin{proof} By B{\'e}zout's identity (Theorem~\ref{thm:2.2.5}), write
	$a x_1 + b x_2 = 1$ and $a y_1 + c y_2 = 1$ with $x_1, x_2, y_1, y_2 \in \mathbb{Z}$.
	Multiplying these two identities and collecting terms gives
	\[
		a\bigl(a x_1 y_1 + c x_1 y_2 + b x_2 y_1\bigr) + bc\,(x_2 y_2) = 1.
	\]
	Hence every common divisor of $a$ and $bc$ divides $1$, so $\gcd(a, bc) = 1$.

\end{proof}
\begin{theorem}\label{thm:2.2.10} If $b \mid a$, $c \mid a$, and
	$\gcd(b, c)=1$, then $bc \mid a$.

\end{theorem}

\begin{proof} Write $a = kb$ with $k \in
		\mathbb{Z}$, so $c \mid kb$. As $\gcd(b, c)=1$,
	Euclid's lemma (Theorem~\ref{thm:2.2.8}) gives $c \mid k$, and then $bc \mid kb =
		a$.

\end{proof}
\begin{theorem}\label{thm:2.2.11} For
	all $a, b \in \mathbb{Z}^+$, we have $\gcd(a, b)\lcm(a, b) = ab$.

\end{theorem}

\begin{proof} Put $g = \gcd(a, b)$ and $m = ab/g$, a
	positive integer because $g \mid a$. From $m = a\,(b/g) = b\,(a/g)$ we
	see that $m$ is a common multiple of $a$ and $b$. Let $n$ be an arbitrary common multiple, say $n = au = bv$. By B{\'e}zout's identity (Theorem~\ref{thm:2.2.5}), write
	$g = ax + by$; then
	\[
		\frac{n}{m} = \frac{ng}{ab} = \frac{n(ax + by)}{ab} = \frac{n}{b}\,x +
		\frac{n}{a}\,y = vx + uy \in \mathbb{Z},
	\]
	so $m \mid n$. In particular $m \le n$ for every positive common
	multiple $n$, so $m = \lcm(a, b)$ and
	$\gcd(a, b)\lcm(a, b) = gm = ab$.

\end{proof}
\begin{theorem}\label{thm:2.2.12} For
	all $a, b, c \in \mathbb{Z}^+$,
	\[
		\gcd(ab, ac) = a\gcd(b, c), \qquad \lcm(ab, ac) =
		a\lcm(b, c).
	\]

\end{theorem}

\begin{proof} Put $g = \gcd(b, c)$. Since $g \mid b$ and $g
		\mid c$, the integer $ag$ divides both $ab$ and $ac$, so $ag \mid
		\gcd(ab, ac)$ by Theorem~\ref{thm:2.2.6}. Conversely, by B{\'e}zout's identity there
	are integers $x, y$ with $bx + cy = g$, hence $ag = (ab)x + (ac)y$;
	since $\gcd(ab, ac)$ divides $ab$ and $ac$, it divides $ag$. Two
	positive integers each dividing the other are equal, so $\gcd(ab, ac) =
		ag$. For the second identity, Theorem~\ref{thm:2.2.11} gives
	\[
		\lcm(ab, ac) = \frac{(ab)(ac)}{\gcd(ab, ac)} =
		\frac{a^{2}bc}{a\gcd(b, c)} = a\cdot\frac{bc}{\gcd(b, c)} =
		a\lcm(b, c).
	\]

\end{proof}
\begin{theorem}\label{thm:2.2.13} Let $a, b \in
		\mathbb{Z}^+$ and $d = \gcd(a, b)$. Then $a = da_1$ and $b = db_1$ for
	some $a_1, b_1 \in \mathbb{Z}^+$ with $\gcd(a_1, b_1) = 1$.

\end{theorem}

\begin{proof} Since $d \mid a$ and $d \mid b$, write $a =
		da_1$ and $b = db_1$ with $a_1, b_1 \in \mathbb{Z}^+$. By Theorem~\ref{thm:2.2.12}, $d = \gcd(a, b) = \gcd(da_1, db_1) = d\gcd(a_1, b_1)$;
	cancelling $d > 0$ gives $\gcd(a_1, b_1) = 1$.

\end{proof}
\subsection{The Fundamental Theorem of Arithmetic}

\begin{theorem}\label{thm:2.3.1} Every integer $n\ge2$ can be written as a
	product of primes.

\end{theorem}

\begin{proof} We use strong induction on $n$. For $n = 2$,
	since $2$ is prime, the statement is true. Now assume
	every integer $k$ with $1 < k < n$ can be written as a product of
	primes. If $n$ is prime, it is itself a product of
	primes; if $n$ is not prime, then $n = ab$ where $a, b \in
		\mathbb{Z}$ and $1 < a, b < n$, so $a$ and $b$ can each be written as
	a product of primes, hence so can $n$. By induction, every integer
	greater than $1$ can be written as a product of primes.

\end{proof}
\begin{theorem}[Euclid's lemma for primes]\label{thm:2.3.2} If $p$ is prime and $p \mid ab$, then $p \mid a$ or $p \mid b$.

\end{theorem}

\begin{proof} If $p \mid a$, the statement holds. If $p
		\nmid a$, then since $\gcd(a, p) \mid p$ and $p$ is prime, $\gcd(a, p)
		\in \{1, p\}$; it cannot equal $p$ (that would give $p \mid a$), so
	$\gcd(a, p)=1$. Since $p \mid ab$, Euclid's lemma (Theorem~\ref{thm:2.2.8}) gives
	$p \mid b$.

\end{proof}
\begin{theorem}[fundamental theorem of arithmetic]\label{thm:2.3.3}
	Every integer greater than $1$ can be written as a product of primes, and
	this factorization is unique up to the order of the factors.

\end{theorem}

\begin{proof} We use strong induction on $n$. For $n = 2$,
	since $2$ is prime, the statement is true. Assume every integer $k$
	with $1 < k < n$ has a unique prime factorization. By Theorem~\ref{thm:2.3.1}, $n$ can be written as a product of primes, say
	$n=p_1p_2\cdots p_r$. Suppose $n$ has another factorization $n=q_1q_2
		\cdots q_s$; then $p_1 \mid q_1 q_2 \cdots q_s$. By induction on $s$
	using Theorem~\ref{thm:2.3.2} (from $p_1 \mid q_1(q_2 \cdots q_s)$ we get $p_1
		\mid q_1$ or $p_1 \mid q_2 \cdots q_s$, and then we repeat the argument with the shorter
	product), $p_1 \mid q_t$ for some $1 \le t \le s$. Since both $p_1$
	and $q_t$ are prime, $p_1 = q_t$. If $n$ is prime, uniqueness is
	immediate ($r = 1$ and every factorization consists of the single factor $n$). If $n$ is
	composite, cancelling the common factor $p_1 = q_t$ from $p_1 p_2
		\cdots p_r = q_1 \cdots q_s$ gives $p_2 \cdots p_r = q_1 \cdots
		q_{t-1} q_{t+1} \cdots q_s$ (which is valid since $\mathbb{Z}$ has no zero
	divisors and $p_1 \neq 0$), and this common value is $\frac{n}{p_1}
		\in \mathbb{Z}$ with $1 < \frac{n}{p_1} < n$. By the induction
	hypothesis the factorization of $\frac{n}{p_1}$ is unique. Thus $r
		- 1 = s - 1$ and, after reordering, $p_i = q_i$ for $i \ge 2$,
	so the factorization of $n$ itself is unique. By induction, every
	integer greater than $1$ has a unique prime factorization.

\end{proof}
\begin{theorem}\label{thm:2.3.4} For every $a \in \mathbb{Z}$ and every
	prime $p$, we have $\gcd(a, p) = 1$ if and only if $p \nmid a$.

\end{theorem}

\begin{proof} We first prove that if $\gcd(a, p) = 1$,
	then $p \nmid a$. Suppose instead that $p \mid a$; then, since $p
		\mid p$ as well, $p$ is a common divisor of $a$ and $p$, so $p \le
		\gcd(a, p) = 1$. But a prime is greater than $1$, a contradiction, so
	$p \nmid a$. Conversely, suppose $p \nmid a$. Since
	$\gcd(a, p) \mid p$ and $p$ is prime, $\gcd(a, p) = 1$ or $\gcd(a, p)
		= p$. The second option would give $\gcd(a, p) \mid a$, hence $p \mid a$,
	which is impossible. Thus $\gcd(a, p) = 1$.

\end{proof}
\begin{theorem}\label{thm:2.3.5} For $x, y \in \mathbb{Z}^+$,
	there exist $A, B \in \mathbb{Z}^+$ such that
	$\lcm(x, y) = AB$, $A \mid x$, $B \mid y$, and $\gcd(A, B) = 1$.

\end{theorem}

\begin{proof} By the fundamental theorem of arithmetic (Theorem~\ref{thm:2.3.3}),
	write $x = \prod_{p\in\Primes} p^{\alpha_p}$ and $y = \prod_{p\in\Primes}
		p^{\beta_p}$ over all primes $p$ (with only finitely many exponents
	nonzero). We first note that for positive integers $u = \prod_{p\in\Primes}
		p^{\gamma_p}$ and $v = \prod_{p\in\Primes} p^{\delta_p}$, we have $u \mid v$ if and only if
	$\gamma_p \le \delta_p$ for every prime $p$: if $v = uw$ with $w =
		\prod_{p\in\Primes} p^{\epsilon_p}$ (by the fundamental theorem of arithmetic), then $\delta_p = \gamma_p + \epsilon_p
		\ge \gamma_p$; conversely, $w = \prod_{p\in\Primes} p^{\delta_p - \gamma_p}$ is a
	positive integer with $uw = v$. Hence a positive integer is a common
	multiple of $x$ and $y$ if and only if its exponent at each $p$ is at least
		$\max\{\alpha_p, \beta_p\}$, and the least such integer is
	$\lcm(x, y) = \prod_{p\in\Primes} p^{\max\{\alpha_p, \beta_p\}}$. Let
	$A = \prod_{p:\ \alpha_p \ge \beta_p} p^{\alpha_p}$ and $B =
		\prod_{p:\ \alpha_p < \beta_p} p^{\beta_p}$; these two products range
	over disjoint sets of primes, so $\gcd(A, B) = 1$, and $AB = \prod_{p\in\Primes}
		p^{\max\{\alpha_p, \beta_p\}} = \lcm(x, y)$. Moreover, $A
		\mid x$, since for each prime $p$ with $\alpha_p \ge \beta_p$ the
	exponent of $p$ in $A$ is $\alpha_p$ (equal to its exponent in $x$),
	and for the remaining primes $A$ has exponent $0 \le \alpha_p$;
	similarly $B \mid y$.

\end{proof}
\begin{theorem}\label{thm:2.3.6} Let $k,r\in\mathbb{Z}^+$ and let
	$a_1,\ldots,a_r$ be pairwise coprime
	positive integers with $a_1a_2\cdots a_r=n^{k}$ for some
	$n\in\mathbb{Z}^+$. Then each $a_i$ is itself a $k$-th power.

\end{theorem}

\begin{proof} By the fundamental theorem of arithmetic
	(Theorem~\ref{thm:2.3.3}), write $n=q_1^{e_1}\cdots q_s^{e_s}$ with distinct
	primes $q_1,\ldots,q_s$, so that $n^{k}=q_1^{ke_1}\cdots q_s^{ke_s}$;
	thus every prime in the factorization of $n^{k}$ occurs with an
	exponent that is a multiple of $k$. Fix an index $i$ and a prime $p$
	with $p\mid a_i$. Because the factors are pairwise coprime, $p$
	divides no $a_j$ with $j\neq i$, so all occurrences of $p$ in the
	product $a_1a_2\cdots a_r=n^{k}$ come from $a_i$ alone; hence the
	exponent of $p$ in $a_i$ equals its exponent in $n^{k}$, a multiple
	of $k$, say $k\,m_p$ with $m_p\ge1$. As this holds for every prime
	dividing $a_i$, unique factorization gives $a_i=\prod_{p\mid
			a_i}p^{\,k\,m_p}=\Big(\prod_{p\mid a_i}p^{\,m_p}\Big)^{k}$, a $k$-th
	power (where the empty product gives $a_i=1=1^{k}$ when $a_i$ has no prime
	factor).

\end{proof}
\begin{theorem}\label{thm:2.3.7} If an integer $d$ is the square of a rational
	number, then $d$ is the square of an integer. In particular, if
	$d\in\mathbb{Z}^+$ is not the square of an integer, then $\sqrt{d}$ is
	irrational.

\end{theorem}

\begin{proof} Suppose $d=x^2$ with $x\in\mathbb{Q}$. If $x=0$, then
	$d=0^2$. Otherwise, write $|x|=a/b$ with $a,b\in\mathbb{Z}^+$ and
	$\gcd(a,b)=1$, which is possible by Theorem~\ref{thm:2.2.13}. Then
	$db^2=a^2$. If $b>1$, then $b$ has a prime factor $p$
	(Theorem~\ref{thm:2.1.1}); since $p\mid db^2=a^2$, Euclid's lemma for primes
	(Theorem~\ref{thm:2.3.2}) gives $p\mid a$, contradicting $\gcd(a,b)=1$.
	Hence $b=1$ and $d=a^2$. For the second assertion, if $\sqrt{d}$ were
	rational, then $d$ would be the square of a rational number, and hence of
	an integer, contrary to the hypothesis.

\end{proof}
\subsection{Fermat Numbers}
\begin{definition}[Fermat numbers]\label{def:2.4.1} The Fermat numbers are defined by $\mathcal{F}_n =
			2^{2^{n}} + 1$ for $n\in\mathbb{Z}_{\ge0}$.

\end{definition}
\begin{theorem}\label{thm:2.4.1}
	For all $n\ge0$, we have $\prod_{i=0}^{n-1}\mathcal{F}_i=\mathcal{F}_n-2$;
	consequently $\gcd(\mathcal{F}_m, \mathcal{F}_n)=1$ for $m\neq n$.

\end{theorem}

\begin{proof} We argue by induction on $n$. For $n=0$, the empty
	product equals $1 = 3 - 2 = \mathcal{F}_0 - 2$. For the inductive step,
	assuming that the formula holds for some $n\ge0$, we obtain
	$\prod_{i=0}^{n}\mathcal{F}_i = (\mathcal{F}_n-2)\mathcal{F}_n =
			2^{2^{n+1}} - 1 = \mathcal{F}_{n+1}-2$. To prove coprimality, suppose $m<n$ and
	let $d=\gcd(\mathcal{F}_m,\mathcal{F}_n)$. Then $d\mid \prod_{i=0}^{n-1}\mathcal{F}_i = \mathcal{F}_n - 2$, so
	$d\mid 2$; but $\mathcal{F}_n$ is odd, hence $d=1$.

\end{proof}
\begin{altproof}[Alternative proof of Theorem~\ref{thm:2.1.2}]
	For every $n\ge0$, choose a prime divisor $q_n$ of $\mathcal{F}_n$.
	If $m\neq n$, Theorem~\ref{thm:2.4.1} gives
	$\gcd(\mathcal{F}_m, \mathcal{F}_n)=1$, so $q_m\neq q_n$. Hence the terms of the
	infinite sequence $\mathcal{F}_0,\mathcal{F}_1,\mathcal{F}_2,\ldots$
	have infinitely many distinct prime divisors, and therefore there are
	infinitely many primes.

\end{altproof}
\begin{theorem}\label{thm:2.4.2} For $n\ge2$, the last
	decimal digit of $\mathcal{F}_n$ is $7$; equivalently, $10\mid
		\mathcal{F}_n-7$.

\end{theorem}

\begin{proof} We argue using divisibility alone. First, by
	induction on $m\ge1$ we show $10\mid 16^{m}-6$. For $m=1$ this is
	$16-6=10$. Assuming $16^{m}-6=10t$ for some integer $t$, we get
	$16^{m+1}-6=16\cdot16^{m}-6=16(6+10t)-6=90+160t=10(9+16t)$, so $10\mid
		16^{m+1}-6$, which completes the induction. Now let $n\ge2$. Since $4\mid
		2^{n}$, write $2^{n}=4m$ with $m=2^{\,n-2}\ge1$; then
	$2^{2^{n}}=2^{4m}=16^{m}$, and the preceding claim gives $10\mid 2^{2^{n}}-6$.
	Hence $10\mid(2^{2^{n}}+1)-7=\mathcal{F}_n-7$, and since $0\le 7<10$ the
	last decimal digit of $\mathcal{F}_n$ is $7$.

\end{proof}
\subsection{The Euclidean Algorithm and the Fibonacci Numbers}

\begin{definition}[Fibonacci numbers]\label{def:2.5.1} The Fibonacci numbers are defined by $F_0 = 0$,
	$F_1 = 1$, and $F_{n+1} = F_n + F_{n-1}$ for $n \ge 1$.

\end{definition}
\begin{theorem}\label{thm:2.5.1} For all integers
	$m \ge 1$ and $n \ge 0$, we have $F_{m+n} = F_m F_{n+1} + F_{m-1} F_n$.

\end{theorem}

\begin{proof} Fix $m \ge 1$ and proceed by induction on $n$. For $n =
		0$, we have $F_m F_1 + F_{m-1} F_0 = F_m\cdot 1 + F_{m-1}\cdot 0 = F_m$. For
	$n = 1$, we have $F_m F_2 + F_{m-1} F_1 = F_m + F_{m-1} = F_{m+1}$. Assuming
	the identity for $n-1$ and $n$ (with $n \ge 1$), we obtain
	\[
		\begin{aligned}
			F_{m+n+1} &= F_{m+n} + F_{m+n-1} \\
			&= (F_m F_{n+1} + F_{m-1} F_n) + (F_m F_n + F_{m-1} F_{n-1}) \\
			&= F_m(F_{n+1} + F_n) + F_{m-1}(F_n + F_{n-1}) \\
			&= F_m F_{n+2} + F_{m-1} F_{n+1}.
		\end{aligned}
	\]

\end{proof}
\begin{theorem}\label{thm:2.5.2}
	$\gcd(F_n, F_{n+1}) = 1$ for all $n \ge 0$.

\end{theorem}

\begin{proof} For $n = 0$ this is $\gcd(F_0, F_1) = \gcd(0, 1) = 1$. For $n \ge 1$, we have $\gcd(F_{n+1}, F_n) =
		\gcd(F_n + F_{n-1}, F_n) = \gcd(F_{n-1}, F_n)$ by Theorem~\ref{thm:2.2.1}. Repeating this reduction gives
	$\gcd(F_{n+1}, F_n) = \gcd(F_0, F_1) = 1$.

\end{proof}
\begin{theorem}\label{thm:2.5.3} For integers $m
		> n \ge 1$, we have $\gcd(F_m, F_n) = \gcd(F_{m-n}, F_n)$.

\end{theorem}

\begin{proof} By Theorem~\ref{thm:2.5.1}, with the first index equal to $m-n\ge1$,
	\[
		F_m=F_{m-n}F_{n+1}+F_{m-n-1}F_n.
	\]
	Subtracting an integer multiple of $F_n$ does not change the set of common divisors (Theorem~\ref{thm:1.1.2}(c)), so
	\[
		\gcd(F_m,F_n)=\gcd(F_{m-n}F_{n+1},F_n).
	\]
	By Theorem~\ref{thm:2.5.2}, $\gcd(F_n,F_{n+1})=1$.
	Every positive divisor $d$ of $F_n$ is consequently coprime to
	$F_{n+1}$; Euclid's lemma (Theorem~\ref{thm:2.2.8}) therefore gives
	\[
		d\mid F_{m-n}F_{n+1}\quad\Longleftrightarrow\quad d\mid F_{m-n}.
	\]
	Thus $(F_{m-n}F_{n+1},F_n)$ and $(F_{m-n},F_n)$ have the same
	positive common divisors and hence the same gcd, proving the claim.

\end{proof}
\begin{theorem}\label{thm:2.5.4} For
	all integers $m, n \ge 0$ not both zero, $\gcd(F_m, F_n) = F_{\gcd(m, n)}$.

\end{theorem}

\begin{proof} If $n = 0$ then $\gcd(F_m, F_0) = \gcd(F_m, 0) = F_m = F_{\gcd(m, 0)}$ (as $\gcd(m, 0) = m$); the case $m
		= 0$ is symmetric. So assume $m, n \ge 1$ with $m \ge n$, and run the Euclidean
	algorithm on the indices: set $r_0 = m$, $r_1 = n$, and $r_{k-1} =
		q_k r_k + r_{k+1}$ with $0 \le r_{k+1} < r_k$, stopping at the last
	nonzero remainder $d = \gcd(m, n)$. At a non-terminal step (where
	$r_{k+1} > 0$), before each of the $q_k$ successive subtractions of $r_k$,
	the first index is strictly greater than $r_k$ (the smallest
	such first index is $r_{k-1} - (q_k-1)r_k = r_k + r_{k+1} >
		r_k$), so Theorem~\ref{thm:2.5.3} (which requires a strictly larger first
	index) applies at each subtraction and gives $\gcd(F_{r_{k-1}}, F_{r_k}) = \gcd(F_{r_{k-1} - r_k}, F_{r_k}) = \cdots =
		\gcd(F_{r_{k+1}}, F_{r_k})$. At the terminal step, $r_{k+1} = 0$, so
	$r_{k-1} = q_k d$ with $r_k = d$; here Theorem~\ref{thm:2.5.3} is applied only
	$q_k - 1$ times, giving $\gcd(F_{q_k d}, F_d) =
		\gcd(F_{(q_k-1)d}, F_d) = \cdots = \gcd(F_d, F_d)$ (where every
	intermediate first index is a multiple of $d$ that exceeds $d$),
	and then $\gcd(F_d, F_d) =
		F_d$ directly, so Theorem~\ref{thm:2.5.3} is never invoked with equal
	indices. Thus $\gcd(F_m, F_n)$ is preserved throughout and equals
	$F_d = F_{\gcd(m, n)}$.

\end{proof}
\begin{theorem}\label{thm:2.5.5} For $m
		\ge 3$ and any $n \ge 1$, we have $F_m \mid F_n$ if and only if $m \mid n$.

\end{theorem}

\begin{proof} By Theorem~\ref{thm:2.5.4}, $F_m \mid F_n$ if and only if
		$\gcd(F_m, F_n) = F_m$, that is, if and only if $F_{\gcd(m, n)} = F_m$. The initial values $F_1 = F_2 = 1$ and the recurrence show by induction that
	$F_k>0$ for every $k\ge1$. Thus
	$F_{k+1} = F_k + F_{k-1} > F_k$ for every $k \ge 2$, so
	the sequence is strictly increasing from $F_2$ onward. Since $F_3 = 2$,
	we have $F_k = 1$ only for $k \in \{1, 2\}$ and $F_k \ge 2$ for $k \ge 3$. Since $m \ge 3$, we have $F_m \ge 2$, so
	the equality $F_{\gcd(m, n)} = F_m$ forces $\gcd(m, n) \ge 3$, and
	strict monotonicity then forces $\gcd(m, n) = m$, that is, $m \mid n$.
	Conversely, $m \mid n$ gives $\gcd(m, n) = m$, whence $\gcd(F_m, F_n)
		= F_m$, that is, $F_m \mid F_n$.

\end{proof}
\subsection{Linear Diophantine Equations and Nonnegative Integer Representations}

\begin{theorem}\label{thm:2.6.1}
	Let $a,b,c\in\mathbb{Z}$ with $a,b$ not both zero, and let $d=\gcd(a, b)$. Then
	$ax+by=c$ has an integer solution if and only if $d\mid c$; and if $(x_0,y_0)$
	is one solution, then all solutions are given by
	$x=x_0+\tfrac{b}{d}t,\ y=y_0-\tfrac{a}{d}t$ for $t\in\mathbb{Z}$.

\end{theorem}

\begin{proof} The solvability criterion is precisely Theorem~\ref{thm:2.2.4}. For the parametrization, if $a=0$ then $d=|b|$ and the solutions are exactly $y=c/b$ with $x$ arbitrary, matching the displayed formula; the case $b=0$ is symmetric. Assume now $ab\neq0$. If $ax+by=ax_0+by_0$ then $a(x-x_0)=b(y_0-y)$; dividing by $d$ gives $\tfrac{a}{d}(x-x_0)=\tfrac{b}{d}(y_0-y)$ with
	$\gcd(\tfrac{a}{d}, \tfrac{b}{d})=1$. By Euclid's lemma (Theorem~\ref{thm:2.2.8}), $\tfrac{a}{d}\mid
		y_0-y$, so $y_0-y=\tfrac{a}{d} t$ and then $x-x_0=\tfrac{b}{d} t$;
	conversely, every pair of this form is a solution.

\end{proof}
\begin{theorem}[Frobenius coin theorem, two-variable case]\label{thm:2.6.2} For coprime $a,b\ge2$, the largest integer not
	representable as $ax+by$ with $x,y\in\mathbb{Z}_{\ge0}$ is $ab-a-b$, and
	exactly $\tfrac{(a-1)(b-1)}{2}$ nonnegative integers are non-representable.

\end{theorem}

\begin{proof}
	For a nonnegative integer $n$, the $b$ integers $n-ax$ with $x\in\{0,\ldots,b-1\}$ have pairwise
	distinct remainders on division by $b$: if $b$ divided $a(x-x')$ with
	$x\neq x'$, Euclid's lemma (Theorem~\ref{thm:2.2.8}) would force
	$b\mid x-x'$, which is impossible when $0<|x-x'|<b$. By the pigeonhole
	principle, every remainder occurs exactly once; in particular, each $n$
	has a unique $x_n\in\{0,\ldots,b-1\}$ for which $b$ divides $n-ax_n$.
	Then $n$ is representable if and only if $n\ge ax_n$: if $n\ge ax_n$
	take $y=(n-ax_n)/b\ge0$; conversely, from $n=ax+by$ with $x,y\ge0$ we
	get $b\mid n-ax$, so $b\mid a(x-x_n)$; Euclid's lemma (Theorem~\ref{thm:2.2.8}) gives
	$b\mid x-x_n$, and then $x\ge x_n$ (otherwise $x\le x_n-b<0$), hence
	$n\ge ax\ge ax_n$. So $n$ is
	non-representable if and only if $0\le n<ax_n$.
	For each fixed $x\in\{0,\ldots,b-1\}$, the non-representable integers $n$ with $x_n=x$ are therefore precisely
	\[
		\{ax-jb\mid 1\le j\le\lfloor ax/b\rfloor\}\subseteq\mathbb{Z}_{\ge0}.
	\]
	This class contains $\lfloor ax/b\rfloor$ integers. When it is nonempty, its largest element is $ax-b$.
	Since $a(b-1)\ge b$, the class with $x=b-1$ is nonempty, and its largest element is greater than that of any other nonempty class. Thus the largest non-representable integer is $a(b-1)-b=ab-a-b$.
	Summing over these disjoint classes, with the class $x=0$ empty, gives a total of
	$\sum_{x=1}^{b-1}\lfloor ax/b\rfloor$ non-representable nonnegative integers. Since $\gcd(a, b)=1$,
	we have $ax/b\notin\mathbb{Z}$ for $0<x<b$, so $\{ax/b\}+\{a(b-x)/b\}=1$ and
	hence $\lfloor ax/b\rfloor+\lfloor
		a(b-x)/b\rfloor=\tfrac{ax}{b}+\tfrac{a(b-x)}{b}-1=a-1$. Summing
	this identity over $x=1,\ldots,b-1$ and using the fact that the map $x\mapsto b-x$ permutes these indices gives $2\sum_{x=1}^{b-1}\lfloor ax/b\rfloor
	=(a-1)(b-1)$, as required.

\end{proof}
\begin{theorem}[stars and bars]\label{thm:2.6.3} For positive integers $n,m$, the number of solutions of
	$x_1+\cdots+x_n=m$ in positive integers is $\binom{m-1}{n-1}$ (with the convention $\binom{a}{b}=0$ when $a<b$, so the count is $0$ if $m<n$), and
	the number of solutions in nonnegative integers is
	$\binom{n+m-1}{n-1}=\binom{n+m-1}{m}$.

\end{theorem}

\begin{proof} For positive solutions, arrange $m$ units in a row
	and choose $n-1$ of the $m-1$ gaps between consecutive units in which
	to place dividers. The $n$ groups have positive sizes $x_1,\ldots,x_n$ summing
	to $m$. This gives a bijection between choices of dividers and positive
	integer solutions, so there are $\binom{m-1}{n-1}$ such solutions.
	For nonnegative solutions, set $y_i=x_i+1$. This gives a bijection
	with the positive integer solutions of $y_1+\cdots+y_n=m+n$, of which
	there are $\binom{m+n-1}{n-1}=\binom{n+m-1}{m}$.

\end{proof}
\newpage

\section{Modular Arithmetic}

\subsection{Congruences and Basic Properties}

\begin{definition}[congruence]\label{def:3.1.1} Let $a, b \in \mathbb{Z}$ and let $n$ be a
	positive integer. If $n \mid (a - b)$, we say that $a$ and $b$ are congruent modulo $n$, written $a \equiv b \pmod{n}$.

	For each fixed $n$, congruence modulo $n$ is an equivalence relation on
	$\mathbb{Z}$: $n\mid(a-a)$ gives reflexivity; $n\mid(a-b)$ implies
	$n\mid(b-a)$, giving symmetry; and $n\mid(a-b)$ together with
	$n\mid(b-c)$ implies $n\mid(a-c)$, giving transitivity.

\end{definition}
\begin{theorem}\label{thm:3.1.1} If $a \equiv x \pmod{n}$ and $b
		\equiv y \pmod{n}$, then the following hold:

	\begin{enumerate}[label={\alph*)}]
		\item $a \pm b \equiv x \pm y \pmod{n}$;
		\item $ab \equiv xy \pmod{n}$;
		\item $a^{c} \equiv x^{c} \pmod{n}$, where $c \in \mathbb{Z}^+$.
	\end{enumerate}

\end{theorem}

\begin{proof}

	\begin{enumerate}[label={\alph*)}]
		\item By the definition of congruence, $n \mid (a - x)$ and $n \mid (b
			      - y)$, so $n \mid \big((a \pm b) - (x \pm y)\big)$; therefore $a \pm
			      b \equiv x \pm y \pmod{n}$.
		\item Since $n \mid (a-x)$, we have $n \mid (a-x)b$; since $n \mid (b - y)$, we have $n
			      \mid x(b-y)$. Then $n \mid \big[(a-x)b + x(b-y)\big] = ab - xy$;
		      therefore $ab \equiv xy \pmod{n}$.
		\item We argue by induction on $c$. When $c = 1$, the congruence $a
			      \equiv x \pmod{n}$ holds by assumption. Now assume $a^{k}
			      \equiv x^{k} \pmod{n}$; then $a^{k+1} \equiv a \cdot a^{k} \equiv x
			      \cdot x^{k} \equiv x^{k+1} \pmod{n}$. By induction, the statement
		      holds for all positive integers $c$.
	\end{enumerate}

\end{proof}
\begin{theorem}\label{thm:3.1.2} Let $m \in \mathbb{Z}^+$ and $k,
		a, a' \in \mathbb{Z}$. If $\gcd(k, m) = d$, then $ka
		\equiv ka' \pmod{m}$ if and only if $a \equiv a' \pmod{m/d}$.

\end{theorem}

\begin{proof} Since $m>0$, we have $d>0$. Write $k=dk_1$ and
	$m=dm_1$, where $k_1\in\mathbb{Z}$ and $m_1\in\mathbb{Z}^+$.
	If $g=\gcd(k_1,m_1)>1$, then $dg$ would be a common divisor of $k$
	and $m$ greater than $d$, a contradiction. Thus $\gcd(k_1,m_1)=1$.
	Cancelling the positive factor $d$ (Theorem~\ref{thm:1.1.2}(a)) shows that
	$m\mid k(a-a')$ if and only if $m_1\mid k_1(a-a')$.
	By Euclid's lemma (Theorem~\ref{thm:2.2.8}), the latter is equivalent to
	$m_1\mid(a-a')$, that is, $a\equiv a'\pmod{m/d}$.
	This also covers $k=0$: then $d=m$ and $m_1=1$, and both congruences
	in the statement hold for all $a,a'$.

\end{proof}
\begin{theorem}\label{thm:3.1.3} Let $a, b \in
		\mathbb{Z}$ and $m, n \in \mathbb{Z}^+$. Then $a \equiv b \pmod{m}$
	and $a \equiv b \pmod{n}$ hold simultaneously if and only if $a \equiv b
		\pmod{\lcm(m, n)}$.

\end{theorem}

\begin{proof} The two congruences on the left say $m \mid (a
		- b)$ and $n \mid (a - b)$, that is, $a - b$ is a common multiple of $m$
	and $n$. By Theorem~\ref{thm:2.2.3}, the least common multiple divides every common
	multiple, so this forces $\lcm(m, n) \mid (a - b)$;
	conversely, every multiple of $\lcm(m, n)$ is a multiple of
	both $m$ and $n$. The two conditions are therefore equivalent.

\end{proof}
\subsection{Inverses, Linear Congruences, and Residue Systems}

\begin{definition}[inverse]\label{def:3.2.1} Let $n \in
		\mathbb{Z}^+$ and $a \in \mathbb{Z}$. If
	$au \equiv 1 \pmod{n}$ for some $u \in \mathbb{Z}$, we call such a $u$ an inverse of $a$ modulo $n$, denoted $a^{-1} \pmod{n}$ (or simply $a^{-1}$ when
	$n$ is clear from context).
	The case $n=1$ is included: since $1\mid au-1$ for all integers $a,u$,
	every integer is an inverse of every integer modulo $1$.

\end{definition}
\begin{theorem}\label{thm:3.2.1} For $a \in \mathbb{Z}$ and $n\in\mathbb{Z}^+$, an inverse of
	$a$ modulo $n$ exists if and only if $a$ and $n$ are coprime.

\end{theorem}

\begin{proof} Suppose first that some $u \in \mathbb{Z}$ satisfies
	$au \equiv 1 \pmod{n}$. Then $n \mid (au - 1)$, that is, $au - 1 = kn$
	for some $k \in \mathbb{Z}$, and hence $au - kn = 1$. Let $\gcd(a, n) =
		d$; since $d \mid a$ and $d \mid n$, we get $d \mid (au - kn) = 1$, so
	$d = 1$. Conversely, if $a$
	and $n$ are coprime, B{\'e}zout's identity (Theorem~\ref{thm:2.2.5}) supplies $x, y \in
		\mathbb{Z}$ with $ax + ny = 1$, hence $ax \equiv 1 \pmod{n}$,
	and $x$ is an inverse of $a$.

\end{proof}
\begin{theorem}\label{thm:3.2.2} For $a \in \mathbb{Z}$ and $n\in\mathbb{Z}^+$, the inverse of
	$a$ modulo $n$, when it exists, is unique modulo $n$.

\end{theorem}

\begin{proof} If $x_1, x_2$ are both inverses of $a$,
	then $x_1 \equiv x_1 a x_2 \equiv x_2 \pmod{n}$, so the inverse is
	unique. Thus the notation $a^{-1} \pmod{n}$ in Definition~\ref{def:3.2.1} is
	unambiguous modulo $n$.

\end{proof}
\begin{theorem}[linear congruence theorem]\label{thm:3.2.3} Let $m \in \mathbb{Z}^+$ and $a, b \in \mathbb{Z}$. The congruence $ax \equiv
		b \pmod{m}$ has a solution if and only if $\gcd(a, m) \mid b$. In that case, setting $d = \gcd(a, m)$, let $u$ be an integer satisfying $u\equiv(\tfrac{a}{d})^{-1}\pmod{m/d}$. Then the congruence has exactly $d$
	distinct solutions modulo $m$, namely
	$x\equiv\tfrac{b}{d}u+\tfrac{m}{d}t\pmod{m}$ for
	$t=0,\ldots,d-1$.
	Moreover, the solution is unique modulo $m$ if $\gcd(a, m) = 1$.

\end{theorem}

\begin{proof} If $ax \equiv b \pmod{m}$ has a solution
	$x_0$, then $ax_0 - b = km$ for some $k \in \mathbb{Z}$; letting $d =
		\gcd(a, m)$, we have $d \mid a$ and $d \mid m$, so $d \mid b$. Conversely,
	suppose $d = \gcd(a, m) \mid b$ and write $a = d a'$, $m = d m'$, $b =
		d b'$ with $a',b'\in\mathbb{Z}$ and $m'\in\mathbb{Z}^+$. B{\'e}zout's identity (Theorem~\ref{thm:2.2.5}) gives $r, s \in \mathbb{Z}$ with $a r + m s = d$;
	dividing by $d$ yields $a' r + m' s = 1$, whence $\gcd(a', m') = 1$.
	Dividing $ax - b = my$ through by $d$ shows that
	$ax \equiv b \pmod{m}$ holds if and only if $a' x \equiv b'
		\pmod{m'}$. Since $\gcd(a', m') = 1$, Theorem~\ref{thm:3.2.1} provides an integer
	$u$ with $a' u \equiv 1 \pmod{m'}$, and the reduced congruence is
	equivalent to $x \equiv b'u \pmod{m'}$. Thus $x_1=b'u$ is a solution,
	and all integer solutions are $x=x_1+m'v$ with $v\in\mathbb{Z}$.
	Write $v=qd+t$ with $0\le t<d$. Since $dm'=m$, this gives
	$x\equiv x_1+m't\pmod m$, so every solution belongs to one of the
	$d$ displayed classes. Conversely, every integer in any of these
	classes is congruent to $x_1$ modulo $m'$ and is therefore a solution.
	For $0\le t,t'<d$, the difference $m'(t-t')$ is divisible by
	$m=dm'$ if and only if $d\mid t-t'$, which in this range holds
	if and only if $t=t'$. Hence the congruence has exactly $d$ distinct
	solutions modulo $m$, namely
	$x\equiv\tfrac{b}{d}u+\tfrac{m}{d}t\pmod m$ for $t=0,\ldots,d-1$.
	For $d=1$ this is the single solution $x\equiv a^{-1}b\pmod m$.

\end{proof}
\begin{definition}[equivalence class]\label{def:3.2.2} For $a\in\mathbb{Z}$ and $n\in\mathbb{Z}^+$, the
	equivalence class (residue class) of $a$ modulo $n$ is
	$a+n\mathbb{Z}=\{a+kn\mid k\in\mathbb{Z}\}$, the set of all integers
	congruent to $a$ modulo $n$. We write $\mathbb{Z}_n=\mathbb{Z}/n\mathbb{Z}$
	for the set of residue classes modulo $n$.

	The sum and product of the classes of $a$ and $b$ are the classes of $a+b$ and $ab$, respectively.
	Theorem~\ref{thm:3.1.1} shows that these operations do not depend on the
	chosen representatives; the ring laws are inherited from $\mathbb{Z}$.

	The group of units $\ZmodUnits{n}$ consists of the classes admitting
	a multiplicative inverse; equivalently, by Theorem~\ref{thm:3.2.1},
	it consists of the classes of integers coprime to $n$. For a prime $p$, we write
	$\mathbb{F}_p=\mathbb{Z}/p\mathbb{Z}$.

\end{definition}
\begin{definition}[complete residue system]\label{def:3.2.3} For $n\in\mathbb{Z}^+$, a
	complete residue system modulo $n$ is a set of integers containing
	exactly one representative from each equivalence class modulo $n$.

\end{definition}
\begin{theorem}\label{thm:3.2.4} For every $n\in\mathbb{Z}^+$, there are exactly $n$ residue
	classes modulo $n$. Both $\{0,1,\ldots,n-1\}$ and
	$\{1,2,\ldots,n\}$ are complete residue systems modulo $n$, and every
	complete residue system modulo $n$ has exactly $n$ elements.

\end{theorem}

\begin{proof} By the division algorithm (Theorem~\ref{thm:1.2.1}),
	every integer is congruent modulo $n$ to exactly one of
	$0,1,\ldots,n-1$, so these integers represent exactly the $n$ residue
	classes. Replacing the representative $0$ by the congruent integer $n$
	gives the second listed complete residue system. Since every complete
	residue system contains exactly one representative from each of the $n$
	classes, it has exactly $n$ elements.

\end{proof}
\begin{theorem}\label{thm:3.2.5} If $\{c_1, c_2, \ldots, c_n\}$ is a
	complete residue system modulo $n$, then for all $k, h \in
		\mathbb{Z}$ with $\gcd(k, n) = 1$, the set $\{k c_1 + h, k c_2 + h,
		\ldots, k c_n + h\}$ is also a complete residue system modulo $n$.

\end{theorem}

\begin{proof} Since $\gcd(k, n) = 1$, Theorem~\ref{thm:3.2.1} gives
	an inverse $k^{-1}$ of $k$ modulo $n$. If $k c_i + h \equiv k c_j + h
		\pmod{n}$ for some $i \neq j$, then $k c_i \equiv k c_j \pmod{n}$, and
	multiplying by $k^{-1}$ gives $c_i \equiv c_j \pmod{n}$, contradicting
	the fact that $c_1, c_2, \ldots, c_n$ represent distinct equivalence classes.
	Hence the $n$ integers $k c_1 + h, \ldots, k c_n + h$ lie in pairwise
	distinct classes modulo $n$; as there are exactly $n$ such classes
	(Theorem~\ref{thm:3.2.4}), each class contains exactly one of them, so
	$\{kc_1 + h, kc_2 + h, \ldots, kc_n + h\}$ is a complete residue
	system.

\end{proof}
\begin{definition}[reduced residue system]\label{def:3.2.4} Let $n\in\mathbb{Z}^+$ and let $C$
	be a complete residue system modulo $n$. The subset consisting of those elements of $C$
	coprime to $n$ is a
	reduced residue system modulo $n$.

\end{definition}
\begin{theorem}\label{thm:3.2.6} For $n\in\mathbb{Z}^+$, all reduced residue systems modulo $n$ have the same number
	of elements, independent of the complete residue system containing
	them.

\end{theorem}

\begin{proof} Whether an integer is coprime to $n$ depends only on its
	residue class: if $x\equiv y\pmod n$, write $y=x+tn$ with $t\in\mathbb{Z}$.
	Every common divisor of $x$ and $n$ divides $y=x+tn$, and every common
	divisor of $y$ and $n$ divides $x=y-tn$. Thus $\gcd(x,n)=\gcd(y,n)$. Hence coprimality to $n$
	is a well-defined property of the $n$ equivalence classes modulo $n$.
	A complete residue system contains exactly one representative of each
	of these $n$ classes (Theorem~\ref{thm:3.2.4}), so the number of its elements coprime to $n$
	equals the number of equivalence classes coprime to $n$, a
	quantity that depends only on $n$. Since a reduced residue system is
	obtained by selecting precisely these elements from a complete residue system, every reduced residue system modulo
	$n$ has the same number of elements.

\end{proof}
\begin{definition}[Euler's totient $\varphi$]\label{def:3.2.5} For $n\in\mathbb{Z}^+$, Euler's totient
	$\varphi(n)$ is the number of elements in a reduced residue
	system modulo $n$; Theorem~\ref{thm:3.2.6} shows that this number is independent of
	the chosen system.

\end{definition}
\begin{theorem}\label{thm:3.2.7} For every $n\in\mathbb{Z}^+$,
	\[
		\varphi(n)=\big|\{a\in\mathbb{Z} \mid 1\le a\le n,\ \gcd(a,n)=1\}\big|.
	\]

\end{theorem}

\begin{proof} The set on the right is the subset of the
	complete residue system $\{1,2,\ldots,n\}$ (Theorem~\ref{thm:3.2.4}) consisting
	of the elements coprime to $n$, and hence is a reduced residue system by
	Definition~\ref{def:3.2.4}. Its cardinality is therefore $\varphi(n)$ by
	Definition~\ref{def:3.2.5}. In particular, $\varphi(1)=1$, since the
	displayed set is $\{1\}$ when $n=1$.

\end{proof}
\begin{theorem}\label{thm:3.2.8} Let $n\in\mathbb{Z}^+$. If $\{r_1, r_2, \ldots,
		r_{\varphi(n)}\}$ is a
	reduced residue system modulo $n$ and $a \in \mathbb{Z}$ with
	$\gcd(a, n)=1$, then $\{a r_1, a r_2, \ldots, a r_{\varphi(n)}\}$ is also a
	reduced residue system modulo $n$.

\end{theorem}

\begin{proof} Write $k=\varphi(n)$. By Definition~\ref{def:3.2.4}, the
	given reduced residue system consists of the elements coprime to $n$
	in some complete residue system. Label the remaining elements of that complete residue system
	$r_{k+1},\ldots,r_n$, so that $\gcd(r_i,n)=1$ exactly when $1\le i\le k$. For every $1 \le i \le n$, the common
	divisors of $ar_i$ and $n$ are exactly the common divisors of $r_i$
	and $n$: if $e\mid ar_i$ and $e\mid n$, then $\gcd(a, e)=1$ because
	$e\mid n$ and $\gcd(a, n)=1$, so B{\'e}zout's identity (Theorem~\ref{thm:2.2.5})
	gives integers $u,v$ with
	$ua+ve=1$, and multiplying by $r_i$ yields $r_i=u(ar_i)+v(er_i)$, so
	$e\mid r_i$; the converse is immediate. Hence $\gcd(ar_i, n) =
		\gcd(r_i, n)$ for all $1 \le i \le n$, so this value equals $1$
	exactly for $1 \le i \le k$. Since $\gcd(a, n)=1$, the set $\{a
		r_1, a r_2, \ldots, a r_n\}$ is also a complete residue system
	(Theorem~\ref{thm:3.2.5} with $h = 0$), and $\{a r_1, a r_2, \ldots, a r_k\}$ is
	precisely
	its subset of integers coprime to $n$, which is a reduced residue
	system by Definition~\ref{def:3.2.4}.

\end{proof}
\subsection{The Fermat--Euler and Wilson Theorems}

\begin{theorem}[Fermat--Euler theorem]\label{thm:3.3.1} If $a \in \mathbb{Z}$
	and $n \in \mathbb{Z}^+$ satisfy $\gcd(a, n)=1$, then $a^{\varphi(n)} \equiv
		1 \pmod{n}$.

\end{theorem}

\begin{proof} Let $\{r_1,\ldots,r_{\varphi(n)}\}$ be a reduced residue
	system modulo $n$. By Theorem~\ref{thm:3.2.8}, the products $ar_i$ also
	form a reduced residue system. Thus there is a permutation $\pi$ of
	$\{1,\ldots,\varphi(n)\}$ such that $ar_i\equiv r_{\pi(i)}\pmod n$
	for every $i$. Put $t=\prod_{i=1}^{\varphi(n)}r_i$. Multiplying these
	congruences gives
	\[
		a^{\varphi(n)}t
		=\prod_{i=1}^{\varphi(n)}ar_i
		\equiv\prod_{i=1}^{\varphi(n)}r_{\pi(i)}
		=t\pmod n.
	\]
	Since $\gcd(r_i,n)=1$ for every $i$, induction on the number of factors
	using Theorem~\ref{thm:2.2.9} gives $\gcd(t,n)=1$. Hence $t^{-1}$ exists
	modulo $n$, and $a^{\varphi(n)}\equiv a^{\varphi(n)}tt^{-1}
	\equiv tt^{-1}\equiv1\pmod n$.

\end{proof}
\begin{samepage}
\begin{theorem}[Fermat's little theorem]\label{thm:3.3.2} When $n$
	is a prime and $\gcd(a, n)=1$, we have $a^{n-1} \equiv 1 \pmod{n}$.

\end{theorem}

\begin{proof} When $n$ is prime, every integer $r$ with $1\le r\le n-1$
	is coprime to $n$: its gcd with $n$ is a positive divisor of $n$
	smaller than $n$, hence $1$. The remaining integer $n$ is not coprime
	to $n$, so Theorem~\ref{thm:3.2.7} gives $\varphi(n)=n-1$.
	Since $\gcd(a,n)=1$, the Fermat--Euler theorem (Theorem~\ref{thm:3.3.1})
	now gives $a^{n-1}\equiv1\pmod n$.

\end{proof}
\end{samepage}
\begin{definition}[Fermat pseudoprime]\label{def:3.3.1} Let $a\in\mathbb{Z}$. A composite
	integer $n$ is a Fermat pseudoprime to base $a$, or briefly a pseudoprime to
	base $a$, if $a^{n-1}\equiv1\pmod{n}$. In that case $\gcd(a, n)=1$: the
	congruence shows that $a^{n-2}$ is an inverse of $a$ modulo $n$, so
	Theorem~\ref{thm:3.2.1} applies.

\end{definition}
\begin{theorem}[pseudoprimes to base $2$]\label{thm:3.3.3} The converse of Fermat's
	little theorem is false: for an integer $n>1$, a congruence
	$a^{n-1}\equiv1\pmod{n}$ with $\gcd(a, n)=1$ does not imply that $n$ is
	prime, even when $a=2$. More precisely, for every prime $p>3$, the
	integer
	\[
		n_p=\frac{4^{p}-1}{3}
	\]
	is a pseudoprime to base $2$ (Definition~\ref{def:3.3.1}). In particular,
	$n_5=341=11\cdot31$ satisfies $2^{340}\equiv1\pmod{341}$, and there are
	infinitely many pseudoprimes to base $2$.

\end{theorem}

\begin{proof} Fix a prime $p>3$; then $p$ is odd and $p\ge5$. Since
	$2\equiv-1\pmod{3}$, Theorem~\ref{thm:3.1.1} gives
	$2^{p}+1\equiv(-1)^{p}+1=0\pmod{3}$, so $u=\tfrac{2^{p}+1}{3}$ is an
	integer and
	\[
		n_p=\frac{(2^{p}-1)(2^{p}+1)}{3}=(2^{p}-1)\,u.
	\]
	Both factors exceed $1$, because $2^{p}-1\ge31$ and $u\ge11$. Hence $n_p$
	is composite (Definition~\ref{def:2.1.2}).

	Next, $2^{2p}=4^{p}=3n_p+1$, so $2^{2p}\equiv1\pmod{n_p}$. We claim that
	$2p\mid n_p-1$. Since $4\equiv1\pmod{3}$, the number
	$w=\tfrac{4^{p-1}-1}{3}$ is an integer, and
	\[
		n_p-1=\frac{4^{p}-4}{3}=4w.
	\]
	Since $p$ is odd, $p\nmid2$, so $\gcd(2, p)=1$ (Theorem~\ref{thm:2.3.4}), and
	Fermat's little theorem (Theorem~\ref{thm:3.3.2}) gives
	$2^{p-1}\equiv1\pmod{p}$. Therefore $3w=4^{p-1}-1=(2^{p-1})^{2}-1$ is
	divisible by $p$. Since $p>3$, we have $p\nmid3$, so $p\mid w$ by
	Theorem~\ref{thm:2.3.2}, and hence $p\mid4w$. Also $2\mid4w$. As
	$\gcd(2, p)=1$, Theorem~\ref{thm:2.2.10} gives $2p\mid4w=n_p-1$.

	Write $n_p-1=2pt$ with $t\in\mathbb{Z}^+$. Then
	$2^{n_p-1}=(2^{2p})^{t}\equiv1\pmod{n_p}$ by Theorem~\ref{thm:3.1.1}(c), so
	$n_p$ is a pseudoprime to base $2$. For $p=5$, we get
	$n_5=\tfrac{1023}{3}=341=11\cdot31$ and $n_5-1=340$. Finally, $n_p<n_q$
	for primes $p<q$, and there are infinitely many primes
	(Theorem~\ref{thm:2.1.2}); thus the integers $n_p$, as $p$ ranges over the primes greater than $3$, form infinitely many distinct pseudoprimes to base $2$.

	The restriction $p>3$ cannot be dropped: $n_3=21$ is composite, but
	$2^{6}=64\equiv1\pmod{21}$ gives $2^{20}=(2^{6})^{3}\cdot4\equiv4\pmod{21}$.

\end{proof}
\begin{theorem}[Wilson's theorem]\label{thm:3.3.4} An integer
	$p>1$ is prime if and only if $(p-1)!\equiv-1\pmod{p}$.

\end{theorem}

\begin{proof} For the forward implication, we first prove that, for a prime
	$p$, if $x^2 \equiv 1 \pmod{p}$, then $x
		\equiv \pm 1 \pmod{p}$. If $x^2 \equiv 1 \pmod{p}$, then $p \mid (x^2 - 1)$,
	and hence $p \mid (x-1)(x+1)$; by Theorem~\ref{thm:2.3.2}, either $p \mid (x-1)$ or $p
		\mid (x+1)$, hence $x \equiv \pm 1 \pmod{p}$. When $p = 2$, we have $1! = 1
		\equiv -1 \pmod{2}$, so the statement is true. When $p\ge3$, every integer in
	$\{1,\ldots,p-1\}$ is coprime to $p$, so we may take its inverse to be the
	unique representative in this same set. For $x\in\{2,\ldots,p-2\}$,
	the implication proved above gives $x^2\not\equiv1\pmod p$, so
	$x^{-1}\not\equiv x\pmod p$. Moreover, $x^{-1}$ also lies in
	$\{2,\ldots,p-2\}$: it cannot equal $1$, since that would give
	$x\equiv1\pmod p$, and it cannot equal $p-1$, since that would give
	$x\equiv(p-1)^{-1}\equiv p-1\pmod p$; here we use $(p-1)^2\equiv1\pmod p$.
	Taking inverses twice returns the original representative, because
	$x$ is the inverse of $x^{-1}$. Hence inversion partitions the $p-3$ elements
	of $\{2,\ldots,p-2\}$ into $\tfrac{p-3}{2}$ disjoint pairs.
	Choose one member $x_i$ from each pair. Each pair has product
	$x_ix_i^{-1}\equiv1\pmod p$, so
	\[
		(p-1)!
		=1\cdot(p-1)\prod_{i=1}^{(p-3)/2}x_ix_i^{-1}
		\equiv-1\pmod p.
	\]
	When $p=3$, the product over pairs is empty and equals $1$.

	Conversely, suppose $(p-1)!\equiv-1\pmod{p}$ but $p$ is
		not prime. As $p > 1$, the integer $p$ is composite, so it has a proper
		divisor $d$ with $1 < d < p$. Since $d \mid p$ and $d \le p-1$, we have $d
	\mid (p-1)!$, so $(p-1)! \equiv 0 \pmod{d}$. However, $(p-1)! \equiv -1
	\pmod{p}$ together with $d \mid p$ implies $(p-1)! \equiv -1 \pmod{d}$.
		Therefore, $0 \equiv -1 \pmod{d}$, which means $d \mid 1$. This forces $d = 1$,
		contradicting $1 < d$. Hence $p$ must be prime.

\end{proof}
\subsection{The Chinese Remainder Theorem, Euler's Totient, and Linear Congruences}

\begin{theorem}[Chinese remainder theorem]\label{thm:3.4.1} Let
	$k\in\mathbb{Z}^+$, let $n_1, n_2, \ldots, n_k \in \mathbb{Z}^+$ be
	pairwise coprime, and let $a_1, a_2, \ldots, a_k \in \mathbb{Z}$. Then
	the system $\{x \equiv a_1 \pmod{n_1},\ \ldots,\ x \equiv a_k \pmod{n_k}\}$ has
	a unique solution modulo $n_1 n_2 \cdots n_k$.

\end{theorem}

\begin{proof} Let $n = \prod_{i=1}^{k} n_i$ and $m_i=n/n_i$.
	The factors $n_j$ with $j\neq i$ are each coprime to $n_i$, so
	Theorem~\ref{thm:2.2.9} gives $\gcd(n_i,m_i)=1$; when $k=1$,
	we have $m_1=1$ and this is immediate. Choose an integer $u_i$ with
	$m_i u_i\equiv1\pmod{n_i}$, and set $c_i=m_i u_i$ and
	$x=\sum_{i=1}^{k}a_i c_i$. For $j\neq i$, we have $n_i\mid m_j$, so
	$c_j\equiv0\pmod{n_i}$, whereas $c_i\equiv1\pmod{n_i}$.
	Thus $x\equiv a_i\pmod{n_i}$ for every $i$, proving existence. For
	uniqueness, if $x_1, x_2$ both satisfy the system, then $x_1
		\equiv x_2 \pmod{n_i}$ for all $i$, so $n_i \mid (x_2 - x_1)$ for all
	$i$. Since the $n_i$ are pairwise coprime, each partial product
	$n_1\cdots n_j$ is coprime to $n_{j+1}$ (Theorem~\ref{thm:2.2.9}), so repeated application of Theorem~\ref{thm:2.2.10} shows inductively that $n_1\cdots n_k\mid(x_2-x_1)$, that is,
	$x_2 \equiv x_1 \pmod{n}$.

\end{proof}
\begin{theorem}[multiplicativity of $\varphi$]\label{thm:3.4.2} If $a, b \in \mathbb{Z}^+$
	and $\gcd(a, b) = 1$, then $\varphi(ab) = \varphi(a)\varphi(b)$.

\end{theorem}

\begin{proof} For each positive integer $q$, put
	$R_q=\{i\in\mathbb{Z}\mid 0\le i<q,\ \gcd(i,q)=1\}$, a reduced
	residue system modulo $q$. For every integer $i$,
	we have $\gcd(i,ab)=1$ if and only if $\gcd(i,a)=\gcd(i,b)=1$:
	one direction follows because $a,b\mid ab$, and the other follows from
	Theorem~\ref{thm:2.2.9}. Also, $i$ and $i\bmod a$ have the same
	common divisors with $a$, and similarly for $b$.
	Hence the map
	\[
		f:R_{ab}\longrightarrow R_a\times R_b,\qquad
		f(i)=(i\bmod a,\ i\bmod b)
	\]
	is well defined. By the Chinese remainder theorem
	(Theorem~\ref{thm:3.4.1}), each pair in $R_a\times R_b$ determines
	exactly one integer $i$ with $0\le i<ab$ satisfying the two
	congruences. The coprimality criterion above shows that this $i$ belongs to
	$R_{ab}$, so $f$ is a bijection. Taking cardinalities gives
	$\varphi(ab)=\varphi(a)\varphi(b)$. This includes the cases $a=1$ and $b=1$,
	since $R_1=\{0\}$.

\end{proof}
\begin{theorem}\label{thm:3.4.3} If $n = p_1^{a_1} p_2^{a_2} \cdots
		p_k^{a_k}$, where $p_1, p_2, \ldots, p_k$ are distinct primes and
	$a_1, a_2, \ldots, a_k \in \mathbb{Z}^+$, then $\varphi(n) = n \prod_{i=1}^{k} \left( 1 -
		\frac{1}{p_i} \right)$.

\end{theorem}

\begin{proof} By multiplicativity of $\varphi$ (Theorem~\ref{thm:3.4.2}), $\varphi(n) = \prod_{i=1}^{k} \varphi(p_i^{a_i})$. For a prime $p$ and $e
		\ge 1$,
	we have $\varphi(p^{e}) = |\{ i \mid 1 \le i \le p^{e},\ p \nmid i \}| = p^{e} -
		p^{e-1} = p^{e} \left( 1 - \frac{1}{p} \right)$. So $\varphi(n) = \prod_{i=1}^{k} p_i^{a_i} \left( 1 -
		\frac{1}{p_i} \right) = n \prod_{i=1}^{k} \left( 1 - \frac{1}{p_i}
		\right)$. For $n=1$, the same formula holds with an empty product,
	since $\varphi(1)=1$.

\end{proof}
\begin{altproof}[Alternative proof of Theorem~\ref{thm:3.4.3}]
	Among the integers $1,2,\ldots,n$, an integer is not coprime to $n$
	if and only if it is divisible by at least one of the distinct prime
	divisors $p_1,\ldots,p_k$ of $n$. For each $1\le i\le k$, let
	$A_i=\{m\in\{1,2,\ldots,n\} \mid p_i\mid m\}$. If
	$\varnothing\neq I\subseteq\{1,2,\ldots,k\}$, then the primes
	$\{p_i \mid i\in I\}$ are distinct and their product divides $n$. An integer
	is divisible by every $p_i$ ($i\in I$) if and only if it is divisible by
	their product (repeated application of Theorem~\ref{thm:2.2.10}), and the
	multiples of a divisor $q$ of $n$ in $\{1,2,\ldots,n\}$ are
	$q,2q,\ldots,(n/q)q$, of which there are $n/q$; hence
	\[
		\left|\bigcap_{i\in I}A_i\right|
		=\frac{n}{\prod_{i\in I}p_i}.
	\]
	By the inclusion--exclusion principle, with the empty product interpreted as $1$,
	\[
		\begin{aligned}
			\varphi(n)
			 & =n-\left|\bigcup_{i=1}^{k}A_i\right|                \\
			 & =n+\sum_{\varnothing\neq I\subseteq\{1,2,\ldots,k\}}
			      (-1)^{|I|}\frac{n}{\prod_{i\in I}p_i} \\
			 & =n\sum_{I\subseteq\{1,2,\ldots,k\}}
			     (-1)^{|I|}\prod_{i\in I}\frac{1}{p_i}             \\
			 & =n\prod_{i=1}^{k}\left(1-\frac{1}{p_i}\right).
		\end{aligned}
	\]

\end{altproof}
\begin{theorem}[generalized Chinese remainder theorem]\label{thm:3.4.4} For $a, b \in \mathbb{Z}$ and $m, n
		\in \mathbb{Z}^+$, the system $\{x \equiv a \pmod{m},\ x \equiv b \pmod{n}\}$
	has a solution if and only if $\gcd(m, n) \mid (a - b)$. Furthermore,
	if a solution exists, it is unique modulo $\lcm(m, n)$.

\end{theorem}

\begin{proof} If a solution $x$ exists, then $x = k_1 m + a =
		k_2 n + b$, so $a - b = k_2 n - k_1 m$, and since $\gcd(m, n)$ divides both $m$ and
		$n$, we get $\gcd(m, n) \mid (a-b)$. Conversely, if $\gcd(m, n) \mid (a-b)$, B{\'e}zout's identity gives $u, v$ with $mu + nv = \gcd(m, n)$. Write $a - b
		= k\gcd(m, n)$, and put $u' = ku$ and $v' = kv$; then $mu' + nv' = a - b$, so $x
		= -mu' + a = nv' + b$ satisfies both congruences. For uniqueness, if
	$x_1, x_2$ are solutions, then $m \mid (x_2-x_1)$ and $n \mid (x_2-x_1)$,
	so $\lcm(m, n) \mid (x_2-x_1)$.

\end{proof}
\begin{theorem}\label{thm:3.4.5} Let $a \in \mathbb{Z}$, $m \in
		\mathbb{Z}^+$, and let $n\in\mathbb{Z}^+$ with $n \ge \log_2 m$. Then
	\[
		a^{n} \equiv a^{(n \bmod \varphi(m)) + \varphi(m)} \pmod{m}.
	\]
\end{theorem}

\begin{proof} If $m=1$, the congruence is automatic. Assume
	$m>1$. Put $r = n \bmod \varphi(m)$, so
	$0 \le r < \varphi(m)$ and $n \equiv r \pmod{\varphi(m)}$; we must show
	$a^{n} \equiv a^{r + \varphi(m)} \pmod{m}$. By the Chinese remainder
	theorem (Theorem~\ref{thm:3.4.1}), it suffices to prove this modulo every prime
	power $p^{k} \parallel m$. Write $m = p^{k}m'$ with $p \nmid m'$.
	Multiplicativity of $\varphi$ (Theorem~\ref{thm:3.4.2}) gives
	$\varphi(m) = \varphi(p^{k})\,\varphi(m')$, so
	$\varphi(p^{k}) \mid \varphi(m)$ and
	$\varphi(m) \ge \varphi(p^{k})$. First, $n \ge \log_2 m \ge \log_2 p^{k} =
		k \log_2 p \ge k$, so $n \ge k$; and $r + \varphi(m) \ge \varphi(m)
		\ge \varphi(p^{k}) = p^{k-1}(p-1) \ge 2^{k-1} \ge k$, so
	$r + \varphi(m) \ge k$.

	If $p \nmid a$, then $\gcd(a, p^{k}) = 1$, and
	since $\varphi(p^{k}) \mid \varphi(m)$, the Fermat--Euler theorem gives
	$a^{\varphi(m)} \equiv 1 \pmod{p^{k}}$. From $n \equiv r
		\pmod{\varphi(m)}$ and $\varphi(p^{k}) \mid \varphi(m)$ we get
	$n \equiv r \pmod{\varphi(p^{k})}$, so $a^{n} \equiv a^{r} \equiv
		a^{r} \cdot a^{\varphi(m)} = a^{r + \varphi(m)} \pmod{p^{k}}$.

	If $p \mid a$, then $p^{k} \mid a^{k} \mid a^{n}$
	(as $n \ge k$), so $a^{n} \equiv 0 \pmod{p^{k}}$; likewise
	$r + \varphi(m) \ge k$ gives $a^{r + \varphi(m)} \equiv 0
		\pmod{p^{k}}$. Hence $a^{n} \equiv a^{r + \varphi(m)} \pmod{p^{k}}$.

	In both cases the congruence holds modulo $p^{k}$; by the Chinese
	remainder theorem it holds modulo $m$.

	Unlike the Fermat--Euler theorem (Theorem~\ref{thm:3.3.1}), this result
	does not require the coprimality hypothesis $\gcd(a,m)=1$. It replaces the exponent
	$n$ by an integer $E=(n\bmod\varphi(m))+\varphi(m)$ satisfying
	$\varphi(m)\le E<2\varphi(m)$; this new exponent need not be smaller than $n$.

\end{proof}
\subsection{Binomial Congruences and Wolstenholme's Theorem}

\begin{theorem}\label{thm:3.5.1} For $n\in\mathbb{Z}_{\ge0}$,
	the number $\binom{m}{n} = \frac{m(m-1) \cdots (m-n+1)}{n!}$ is an integer when
	$m\in\mathbb{Z}_{\ge0}$, and $\binom{-m}{n}
		= (-1)^n \frac{m(m+1) \cdots (m+n-1)}{n!}$ is an integer when
	$m\in\mathbb{Z}^+$; in both formulas the product is empty when $n=0$.

\end{theorem}

\begin{proof} For $n=0$, both displayed products are empty and both quotients equal $1$. Assume $n\ge1$. Since $\binom{-m}{n} = (-1)^n
		\binom{m+n-1}{n}$, the statement is equivalent to the assertion that the rising factorial $(m)_n := m(m+1)
		\cdots (m+n-1)$ (here this subscript denotes the rising factorial, not the falling factorial) is a multiple of $n!$ for all positive integers
	$m$. This also settles $\binom{m}{n}$: for $m \ge n$, the falling product
	$m(m-1)\cdots(m-n+1)$ equals the rising factorial $(m-n+1)_n$, and is therefore a
	multiple of $n!$; for $m < n$, it contains the factor $0$, so
	$\binom{m}{n} = 0 \in \mathbb{Z}$. We prove the rising-factorial claim by induction. When $n = 1$, we have $(m)_1 =
		m$, so $1! \mid (m)_1$. Now assume $(N-1)! \mid
		(m)_{N-1}$ for all $m$. To prove $N! \mid (m)_N$ for all $m$, proceed by induction
	on $m$. For $m = 1$, we have $(1)_N = N!$. For $M \ge 2$, assume $N! \mid (M-1)_N$; then $(M)_N - (M-1)_N =
			(M+N-1)(M)_{N-1} - (M-1)(M)_{N-1} = N \cdot
		(M)_{N-1}$, and $(N-1)! \mid (M)_{N-1}$ gives $N! \mid N
			(M)_{N-1}$. Combined with the hypothesis $N! \mid (M-1)_N$, this gives $N! \mid (M)_N$. By
	induction, $n! \mid (m)_n$ for all $m, n \in \mathbb{Z}^+$.

\end{proof}
\begin{theorem}\label{thm:3.5.2} If $p$ is a prime, then $\binom{p}{1},
		\binom{p}{2}, \ldots, \binom{p}{p-1}$ are all divisible by $p$.

\end{theorem}

\begin{proof} For $1 \le n \le p-1$, Theorem~\ref{thm:3.5.1} shows that
	$\binom{p}{n} =
		\frac{p(p-1)\cdots(p-n+1)}{n!}$
	is an integer. Hence $n! \mid
		p(p-1)\cdots(p-n+1)$. Since $p$ is prime and $1 \le n < p$, none of
	$1, 2, \ldots, n$ is divisible by $p$, so $\gcd(n!, p) = 1$; Euclid's
	lemma (Theorem~\ref{thm:2.2.8}) then gives $n! \mid
		(p-1)(p-2)\cdots(p-n+1)$. Hence $\binom{p}{n} = p \cdot
		\frac{(p-1)\cdots(p-n+1)}{n!}$ with
	$\frac{(p-1)\cdots(p-n+1)}{n!} \in \mathbb{Z}$, so $p \mid
		\binom{p}{n}$.

\end{proof}
\begin{theorem}[Vandermonde's identity]\label{thm:3.5.3} For
	nonnegative integers $m, n, r$,
	\[
		\binom{m+n}{r} = \sum_{k=0}^{r} \binom{m}{k}\binom{n}{r-k}
	\]
	(with the convention $\binom{a}{b} = 0$ for $b > a$).

\end{theorem}

\begin{proof} Compare coefficients of $x^{r}$ in
	$(1+x)^{m}(1+x)^{n} = (1+x)^{m+n}$. On the right, the
	coefficient of $x^{r}$ is $\binom{m+n}{r}$. On the left, write
	$(1+x)^{m} = \sum_{i}\binom{m}{i}x^{i}$ and $(1+x)^{n} =
		\sum_{j}\binom{n}{j}x^{j}$; collecting terms with $i + j = r$
	gives $\sum_{k=0}^{r}\binom{m}{k}\binom{n}{r-k}$. Equating the two coefficients yields the identity.

\end{proof}
\begin{definition}[localization at a prime]\label{def:3.5.1} For a
	prime $p$, the localization of $\mathbb{Z}$ at $p$ is
	\[
		\mathbb{Z}_{(p)}=\left\{\frac{a}{b}\in\mathbb{Q} \mid a,b\in\mathbb{Z},\
		p\nmid b\right\}.
	\]
	For $x,y\in\mathbb{Z}_{(p)}$ and $r\in\mathbb{Z}^+$, write
	$x\equiv y\pmod{p^r}$ when $x-y\in p^r\mathbb{Z}_{(p)}$. An element $u$ of a
	commutative ring $R$ with identity is a unit if $uv=1$ for some $v\in R$;
	this is the sense in which the word is used throughout, and for
	$R=\mathbb{Z}_n$ it recovers the group $\ZmodUnits{n}$ of
	Definition~\ref{def:3.2.2}.

\end{definition}
\begin{theorem}[congruence of fractions]\label{thm:3.5.4}
	Let $p$ be a prime and $r\in\mathbb{Z}^+$. For $x=a/b$ and $y=c/d$ with
	$a,b,c,d\in\mathbb{Z}$ and $p\nmid bd$,
	\[
		\frac{a}{b}\equiv\frac{c}{d}\pmod{p^r}
		\quad\Longleftrightarrow\quad p^r\mid(ad-bc).
	\]
	Consequently, for integers this congruence agrees with ordinary congruence
	modulo $p^r$, and $a/b$ is congruent to $a$ times an inverse of $b$ modulo
	$p^r$. In particular, a fraction in $\mathbb{Z}_{(p)}$ is congruent to $0$
	modulo $p^r$ exactly when its numerator in lowest terms is divisible by $p^r$.

\end{theorem}

\begin{proof} Indeed, if $ad-bc=p^r h$ with $h\in\mathbb{Z}$, then
	$x-y=p^r h/(bd)\in p^r\mathbb{Z}_{(p)}$. Conversely, if
	$(ad-bc)/(bd)=p^r e/f$ with $p\nmid f$, then
	$f(ad-bc)=p^r ebd$; since $\gcd(f,p^r)=1$, Euclid's lemma gives
	$p^r\mid(ad-bc)$. Thus, for integers, this definition agrees with ordinary congruence,
	and $a/b$ reduces to $a$ times an inverse of $b$ modulo $p^r$.
	In particular, a fraction in $\mathbb{Z}_{(p)}$ is congruent to $0$
	modulo $p^r$ exactly when its numerator in lowest terms is divisible by $p^r$.

\end{proof}

\begin{theorem}\label{thm:3.5.5}
	If $b\in\mathbb{Z}$ and $p\nmid b$, then $b$ is a unit in
		$\mathbb{Z}_{(p)}$. Consequently, a denominator coprime to $p$ may be
	inverted when reducing modulo $p^r$.

\end{theorem}

\begin{proof} Both $b=b/1$ and $b^{-1}=1/b$ belong to
	$\mathbb{Z}_{(p)}$, and their product is $1$. Multiplication by $b^{-1}$
	preserves congruences modulo $p^r$, because if
	$x-y=p^r z$ with $z\in\mathbb{Z}_{(p)}$, then
	$b^{-1}x-b^{-1}y=p^r(b^{-1}z)$ and $b^{-1}z\in\mathbb{Z}_{(p)}$.

\end{proof}
\begin{theorem}\label{thm:3.5.6} For every prime $p\ge5$,
	\[
		\sum_{k=1}^{p-1}\frac{1}{k^{2}}\equiv0\pmod{p},
	\]
	where the fractions are interpreted in $\mathbb{Z}_{(p)}$ as in
	Definition~\ref{def:3.5.1}.

\end{theorem}

\begin{proof} Since $\gcd(i, p)=1$ for all $i\in\{1,\ldots,p-1\}$,
	the inverse $i^{-1}$ modulo $p$ exists and is unique. For $i\neq j$ in that range,
	$i^{-1}\not\equiv j^{-1}\pmod{p}$, since $i^{-1}\equiv j^{-1}$ would give
	$i\equiv ijj^{-1}\equiv iji^{-1}\equiv j\pmod{p}$. Inversion therefore
	permutes $\{1,2,\ldots,p-1\}$ modulo $p$, so, by Theorem~\ref{thm:3.5.4},
	\[
		\sum_{k=1}^{p-1}\frac{1}{k^{2}}\equiv\sum_{k=1}^{p-1}(k^{-1})^{2}
		\equiv\sum_{k=1}^{p-1}k^{2}=\frac{(p-1)p(2p-1)}{6}\pmod{p}.
	\]
	The closed form follows from the classical identity
	$\sum_{k=1}^{m}k^{2}=m(m+1)(2m+1)/6$, verified by induction on $m$ since
	$m(m+1)(2m+1)-(m-1)m(2m-1)=6m^{2}$ (the companion identity
	$\sum_{k=1}^{m}k=m(m+1)/2$, used below, follows in the same way from
	$m(m+1)-(m-1)m=2m$); and $(p-1)p(2p-1)/6$ is an integer
	divisible by $p$ because $\gcd(6, p)=1$ for $p\ge5$.

\end{proof}
\begin{theorem}[Wolstenholme's theorem]\label{thm:3.5.7} For all
	primes $p \ge 5$, we have $1 + 2^{-1} + \cdots + (p-1)^{-1} \equiv 0 \pmod{p^2}$.

\end{theorem}

\begin{proof} Work in $\mathbb{Z}_{(p)}$ and put
	$H=\sum_{k=1}^{p-1}1/k$ and $T=\sum_{k=1}^{p-1}1/[k(p-k)]$.
	All denominators are coprime to $p$. Since $k\mapsto p-k$ permutes
	$\{1,\ldots,p-1\}$, we have the exact identity
	\[
		2H=\sum_{k=1}^{p-1}\left(\frac{1}{k}+\frac{1}{p-k}\right)
		=p\sum_{k=1}^{p-1}\frac{1}{k(p-k)}=pT.
	\]
	Modulo $p$, we have $k(p-k)\equiv-k^2$, so
	\[
		T\equiv-\sum_{k=1}^{p-1}\frac{1}{k^2}\equiv0\pmod p
	\]
	by Theorem~\ref{thm:3.5.6}. Hence $2H=pT\equiv0\pmod{p^2}$.
	Since $p\ge5$, the integer $2$ is a unit in $\mathbb{Z}_{(p)}$
	(Theorem~\ref{thm:3.5.5}); cancelling it gives $H\equiv0\pmod{p^2}$.

\end{proof}
\begin{theorem}[Wolstenholme's theorem, binomial form]\label{thm:3.5.8} For
	every prime $p \ge 5$, we have $\binom{2p-1}{p-1} \equiv 1 \pmod{p^{3}}$.

\end{theorem}

\begin{proof} For $1\le k\le p-1$, set $c_k=\binom pk/p$,
	an integer by Theorem~\ref{thm:3.5.2}. Vandermonde's identity
	(Theorem~\ref{thm:3.5.3}) and $\binom pk=\binom p{p-k}$ give
	\[
		\binom{2p}{p}=\sum_{k=0}^{p}\binom pk^2
		=2+p^2\sum_{k=1}^{p-1}c_k^2.
	\]
	The factorial formula gives the exact identity
	\[
		k\binom pk=\frac{p!}{(k-1)!\,(p-k)!}=p\binom{p-1}{k-1},
	\]
	so $kc_k=\binom{p-1}{k-1}$. In $\mathbb{Z}_{(p)}$,
	\[
		\binom{p-1}{k-1}
		=\prod_{i=1}^{k-1}\frac{p-i}{i}
		\equiv(-1)^{k-1}\pmod p,
	\]
	including $k=1$ by the empty-product convention. Thus
	$c_k^2\equiv1/k^2\pmod p$, and Theorem~\ref{thm:3.5.6} yields
	\[
		\sum_{k=1}^{p-1}c_k^2
		\equiv\sum_{k=1}^{p-1}\frac{1}{k^2}\equiv0\pmod p.
	\]
	The left-hand sum is an integer, so it is divisible by $p$.
	Consequently $\binom{2p}{p}\equiv2\pmod{p^3}$.
	Finally, $\binom{2p}{p}=2\binom{2p-1}{p-1}$, and $2$ is invertible
	modulo $p^3$, giving $\binom{2p-1}{p-1}\equiv1\pmod{p^3}$.

\end{proof}
\begin{theorem}\label{thm:3.5.9} Let $p$ be an odd prime. Then:

	\begin{enumerate}[label={\alph*)}]
		\item (Fermat quotient) $\dfrac{2^{p-1} - 1}{p} \equiv 1 +
			      \dfrac{1}{3} + \dfrac{1}{5} + \cdots + \dfrac{1}{p-2} \pmod{p}$ (where each
		      fraction $1/k$ denotes the inverse of $k$ modulo $p$).
		\item (Wilson modulo $p^{2}$) $(p-1)! \equiv (-1)^{\frac{p-1}{2}}\,
			      2^{2p-2} \big[\big(\tfrac{p-1}{2}\big)!\big]^{2} \pmod{p^{2}}$.
	\end{enumerate}

\end{theorem}

\begin{proof} Put $q_2=(2^{p-1}-1)/p$, an integer by Fermat's little
	theorem (Theorem~\ref{thm:3.3.2}). For the first part, the binomial theorem gives
	\[
		2^p=(1+1)^p=\sum_{k=0}^{p}\binom{p}{k}=2+\sum_{k=1}^{p-1}\binom{p}{k},
	\]
	so $2^p-2=\sum_{k=1}^{p-1}\binom{p}{k}$. For $1\le k\le p-1$,
	\[
		\binom{p}{k}=\frac{p!}{k!(p-k)!}=\frac{p}{k}\frac{(p-1)!}{(k-1)!(p-k)!}=\frac{p}{k}\binom{p-1}{k-1}.
	\]
	Consequently $\frac{2^{p-1} - 1}{p} =
		\frac{2^{p} - 2}{2p} =
		\frac{1}{2}\sum_{k=1}^{p-1}\frac{1}{k}\binom{p-1}{k-1}$. Since $\binom{p-1}{k-1} = \prod_{i=1}^{k-1}\tfrac{p-i}{i} \equiv \prod_{i=1}^{k-1}\tfrac{-i}{i} = (-1)^{k-1} \pmod{p}$, the Fermat quotient is
	congruent to
		$\frac{1}{2}\sum_{k=1}^{p-1}\frac{(-1)^{k-1}}{k} =
		\frac{1}{2}(S_{\text{odd}} - S_{\text{even}}) \pmod{p}$, where
	$S_{\text{odd}} = \sum_{k\ \text{odd}} \frac{1}{k}$ and $S_{\text{even}} =
		\sum_{k\ \text{even}} \frac{1}{k}$, with both sums taken over $1 \le k \le p-1$. Since $k
		\mapsto k^{-1}$ permutes the nonzero residues modulo $p$, we have
	$\sum_{k=1}^{p-1}\frac{1}{k} \equiv \sum_{k=1}^{p-1} k =
		\frac{p(p-1)}{2} \equiv 0 \pmod{p}$, where the last congruence holds because $p$ is odd,
	so that $\tfrac{p-1}{2}$ is an integer. Hence
	$S_{\text{odd}} + S_{\text{even}} = \sum_{k=1}^{p-1}\frac{1}{k} \equiv 0
		\pmod{p}$, and therefore $S_{\text{odd}} - S_{\text{even}} \equiv 2
		S_{\text{odd}} \pmod{p}$, so $\frac{2^{p-1}-1}{p} \equiv S_{\text{odd}} = 1 +
		\frac{1}{3} + \cdots + \frac{1}{p-2} \pmod{p}$.

	We next establish the required congruence modulo $p^2$. Let
	$h=\tfrac{p-1}{2}$.
	Then $(p-1)! = h!\cdot\prod_{j=1}^{h}(p - j)$. In the expansion of
	$\prod_{j=1}^{h}(j - p)$ modulo $p^{2}$, only the constant term and
	the terms linear in $p$ survive; the coefficient of $-p$ is
	$\sum_{j}\prod_{i\neq j}i = \sum_{j}\frac{h!}{j} = h! H_h$, where
	$H_h = \sum_{j=1}^{h}\frac{1}{j}$. These are equalities of rational
	numbers. Since $p \nmid j$ for $1\le j\le h$, each fraction $1/j$ can
	be interpreted modulo $p^{2}$ as the multiplicative inverse of $j$. Hence
	$\prod_{j=1}^{h}(j - p) \equiv h!\big(1 - p H_h\big) \pmod{p^{2}}$,
	so $\prod_{j=1}^{h}(p-j) =
		(-1)^{h}\prod_{j=1}^h (j-p) \equiv (-1)^{h} h!\,(1 - p H_h) \pmod{p^{2}}$, giving
	$(p-1)! \equiv (-1)^{h}(h!)^{2}(1 - p H_h) \pmod{p^{2}}$. Now $H_h =
		\sum_{j=1}^{h}\frac{1}{j} = 2 S_{\text{even}}$, and by the
	first part of the proof, $S_{\text{odd}} \equiv -S_{\text{even}} \pmod{p}$, so $q_2 \equiv
		S_{\text{odd}} \equiv -\tfrac{1}{2} H_h \pmod{p}$, that is, $1 - p H_h \equiv
		1 + 2p\,q_2 \equiv (1 + p q_2)^{2} = (2^{p-1})^{2} = 2^{2p-2}
		\pmod{p^{2}}$. Therefore $(p-1)! \equiv (-1)^{h} 2^{2p-2} (h!)^{2} =
		(-1)^{\frac{p-1}{2}} 2^{2p-2}\big[\big(\tfrac{p-1}{2}\big)!\big]^{2}
		\pmod{p^{2}}$.

\end{proof}
\begin{definition}[formal power-series ring]\label{def:3.5.2} The
	formal power-series ring $\mathbb{Z}[[x]]$ consists of all coefficient
	sequences $(a_0,a_1,\ldots)$ with $a_i\in\mathbb{Z}$, written
	$\sum_{i\ge0}a_ix^i$. Addition is coefficientwise, and multiplication is
	given by the finite convolution
	\[
		\left(\sum_{i\ge0}a_ix^i\right)\left(\sum_{j\ge0}b_jx^j\right)
		=\sum_{n\ge0}\left(\sum_{i=0}^{n}a_i b_{n-i}\right)x^n.
	\]

	Equality is coefficientwise; no convergence of numerical series is required.
	The same definitions with coefficients in $\mathbb{Q}$ or $\mathbb{C}$ give the rings
	$\mathbb{Q}[[x]]$ and $\mathbb{C}[[x]]$; their basic properties are collected in
	Theorem~\ref{thm:3.5.10}.

\end{definition}
\begin{theorem}[formal power series]\label{thm:3.5.10}
	Let $K$ be one of $\mathbb{Z}$, $\mathbb{Q}$, and $\mathbb{C}$.

	\begin{enumerate}[label={\alph*)}]
		\item $K[[x]]$ is a commutative ring with identity $1$.
		\item A series $\sum_{i\ge0}a_ix^i\in K[[x]]$ whose constant term $a_0$ is
		      invertible in $K$ has exactly one inverse in $K[[x]]$.
		\item For $c\in K$, the substitution $f(x)\mapsto f(cx)$, that is,
		      $\sum_{i\ge0}a_ix^i\mapsto\sum_{i\ge0}a_ic^ix^i$, preserves sums and products.
	\end{enumerate}

\end{theorem}

\begin{proof} These operations make $\mathbb{Z}[[x]]$ a commutative ring with identity
	$1$: commutativity and distributivity are checked coefficientwise, and
	multiplication is associative because both bracketings of a product of
	three series have coefficient $\sum_{i+j+k=n}a_ib_jc_k$ at $x^n$.
	The same computations apply with coefficients in $\mathbb{Q}$ or $\mathbb{C}$.
	In each of these rings, a series
	$\sum_{i\ge0}a_ix^i$ whose constant term $a_0$ is invertible in the
	coefficient ring has exactly one inverse $\sum_{i\ge0}b_ix^i$: the
	conditions that $\sum_{i=0}^{n}a_ib_{n-i}$ equal $1$ for $n=0$ and $0$ for
	$n\ge1$ determine $b_0=a_0^{-1}$ and, recursively,
	$b_n=-a_0^{-1}\sum_{i=1}^{n}a_ib_{n-i}$ for $n\ge1$. For a constant $c$ in the
	coefficient ring, the substitution $\sum_{i\ge0}a_ix^i\mapsto\sum_{i\ge0}a_ic^ix^i$,
	written $f(x)\mapsto f(cx)$, preserves sums and products, because
	$c^ic^{n-i}=c^n$.

\end{proof}
\begin{theorem}\label{thm:3.5.11} In the formal power-series ring
	$\mathbb{Z}[[x]]$ of Definition~\ref{def:3.5.2}, if $p$ is a prime, all coefficients of $(1-x)^{-p}$ except the coefficients of the monomials $1, x^{p}, x^{2p}, \ldots$ are divisible by $p$, while the coefficients of $1, x^{p}, x^{2p},
		\ldots$ are congruent to $1 \pmod{p}$.

\end{theorem}

\begin{proof} In $\mathbb{Z}[[x]]$ (Theorem~\ref{thm:3.5.10}) we have
	$(1-x)^{-1} = \sum_{j\ge0} x^{j}$, since
	$(1-x)\sum_{j\ge0}x^{j} = 1$ coefficientwise. For $n\ge1$, the coefficient of $x^{n}$ in the $p$-fold product
	$(1-x)^{-p}$ equals the number of
	$p$-tuples $(j_1,\ldots,j_p)$ of nonnegative integers satisfying
	$j_1+\cdots+j_p=n$. By stars and bars (Theorem~\ref{thm:2.6.3}), this
	number is $\binom{n+p-1}{p-1} = \binom{p+n-1}{n}$. The constant term is $1$; hence
	$(1-x)^{-p} = 1 + \sum_{n=1}^{\infty}
		\binom{p+n-1}{n} x^{n}$ has integer coefficients. Since
	$(1-x^{p})^{-1} = 1 + x^{p} + x^{2p} + \cdots$, we only need to prove that
	$(1-x^{p})^{-1} - (1-x)^{-p}$ has every coefficient divisible by $p$.
	Now $(1-x^{p})^{-1} - (1-x)^{-p} =
		(1-x^{p})^{-1}(1-x)^{-p}\big[(1-x)^{p} - 1 + x^{p}\big]$. When $p =
		2$, we have $(1-x)^{2} - 1 + x^{2} = -2x + 2x^{2}$. When $p \ge 3$,
	we have $(1-x)^{p} - 1 + x^{p} = \sum_{r=1}^{p-1} (-1)^{r} \binom{p}{r}
		x^{r}$. In either case, all coefficients of the bracketed factor are divisible by $p$
	(Theorem~\ref{thm:3.5.2}); since $(1-x^{p})^{-1}(1-x)^{-p}$ has integer coefficients, each
	coefficient of the product is a finite $\mathbb{Z}$-linear combination of the coefficients of the bracketed factor,
	hence is still divisible by $p$.

\end{proof}
\begin{theorem}\label{thm:3.5.12} Let $p$
	be a prime. If $a \in \mathbb{Z}^+$ and $m \equiv 1 \pmod{p^{a}}$, then $m^{p}
		\equiv 1 \pmod{p^{a+1}}$.

\end{theorem}

\begin{proof} Write $m=1+kp^a$ with $k\in\mathbb{Z}$. Then
	\[
		m^p=1+kp^{a+1}+\sum_{i=2}^{p}\binom pi k^i p^{ai}.
	\]
	The linear term is divisible by $p^{a+1}$. For every $i\ge2$,
	we have $ai\ge2a\ge a+1$ because $a\ge1$, so each remaining term is also
	divisible by $p^{a+1}$. Hence $m^p\equiv1\pmod{p^{a+1}}$.\qedhere

\end{proof}
\begin{altproof}[Alternative proof of Theorem~\ref{thm:3.3.1}]
	Let $a\in\mathbb{Z}$ and $m\in\mathbb{Z}^+$ with $\gcd(a,m)=1$.
	If $m=1$, the congruence is automatic. Assume $m>1$ and, by the
	fundamental theorem of arithmetic (Theorem~\ref{thm:2.3.3}), write
	$m=\prod_{i=1}^{r}p_i^{e_i}$ with distinct primes $p_i$ and $e_i\ge1$.
	By the Chinese remainder theorem (Theorem~\ref{thm:3.4.1}), it suffices
	to prove $a^{\varphi(m)}\equiv1\pmod{p_i^{e_i}}$ for every $i$.

	For any prime $p$, Theorem~\ref{thm:3.5.2} and the binomial theorem give
	$(x+y)^p\equiv x^p+y^p\pmod p$ for all integers $x,y$.
	Induction on the number of summands extends this to any finite sum.
	In particular, for $1\le b\le p-1$, taking $b$ summands equal to $1$
	gives $b^p\equiv b\pmod p$; for $b=0$ this is immediate.
	Both sides depend only on the residue class of $b$, so
	$b^p\equiv b\pmod p$ for every integer $b$.

	Now fix $i$ and take $p=p_i$. Since $p_i\mid m$ and $\gcd(a,m)=1$,
	we have $\gcd(a,p_i)=1$. Cancelling $a$ from $a^{p_i}\equiv a\pmod{p_i}$
	gives $a^{p_i-1}\equiv1\pmod{p_i}$.
	Successive applications of Theorem~\ref{thm:3.5.12}, raising the
	current integer to the $p_i$-th power each time, give, after $e_i-1$ steps,
	\[
		a^{p_i^{e_i-1}(p_i-1)}\equiv1\pmod{p_i^{e_i}}.
	\]
	By Theorem~\ref{thm:3.4.3}, the exponent is $\varphi(p_i^{e_i})$,
	which divides $\varphi(m)$. Raising this congruence to the positive
	integer power $\varphi(m)/\varphi(p_i^{e_i})$ yields
	$a^{\varphi(m)}\equiv1\pmod{p_i^{e_i}}$ for every $i$.
	The Chinese remainder theorem now gives the result modulo $m$.

\end{altproof}
\subsection{Polynomial Congruences and Lagrange's Theorem}

\begin{definition}[polynomial congruences]\label{def:3.6.1}
	For a prime $p$ and polynomials $f,g\in\mathbb{Z}[x]$, a polynomial identity
	$f(x)\equiv g(x)\pmod{p}$ is understood coefficientwise: every coefficient
	of $f-g$ is divisible by $p$, or equivalently, their reductions are equal
	in $\mathbb{F}_p[x]$. This is stronger than agreement of their values
	modulo $p$: for example, $x^2-x$ takes an even value at every integer $x$, but its
	coefficient of $x^2$ is not divisible by $2$. In an equation
	$f(x)\equiv0\pmod p$ to be solved for $x$, the congruence instead concerns
	the value of $f$ at $x$. The roots of this congruence are counted as distinct residue classes;
	we use integer representatives when writing them as $x_1,\ldots,x_k$.

\end{definition}

\begin{theorem}[factor theorem modulo $p$]\label{thm:3.6.1} Let $p$ be a prime and $f(x)$ a
	polynomial with integer coefficients. If $f(x) \equiv 0 \pmod{p}$ has
	$k$ distinct roots $x_1, \ldots, x_k$ modulo $p$, then there exists a
	polynomial $g(x)$ with integer coefficients such that $f(x) \equiv (x
		- x_1) \cdots (x - x_k)\, g(x) \pmod{p}$.

\end{theorem}

\begin{proof} We argue by induction on $k$. When $k=0$,
	take $g=f$. For the second base case, when $k=1$, dividing the integer
	polynomial $f$ by the monic linear
	polynomial $x - x_1$ using the usual algorithm (at each step subtracting an integer multiple of a monomial in $x - x_1$ to cancel the current leading term, which keeps every coefficient an integer) yields
	$f(x) = (x - x_1) q(x) + r$ with $q(x)$ an integer-coefficient
	polynomial and $r \in \mathbb{Z}$; then $0 \equiv
		f(x_1) \equiv r \pmod{p}$, so $f(x) \equiv (x - x_1) q(x) \pmod{p}$.
	Now assume the claim holds for some $k = t\ge1$. If $f$ has $t+1$
	distinct roots, then $f(x) \equiv (x - x_1) \cdots (x - x_t)\, h(x)
		\pmod{p}$; substituting $x = x_{t+1}$ gives $0 \equiv (x_{t+1} - x_1)
		\cdots (x_{t+1} - x_t)\, h(x_{t+1}) \pmod{p}$, and since $x_{t+1} -
		x_i \not\equiv 0$, repeated application of Theorem~\ref{thm:2.3.2} gives $h(x_{t+1}) \equiv 0 \pmod{p}$, so
	$h(x) \equiv (x - x_{t+1}) g(x) \pmod{p}$, giving the claim.

\end{proof}
\begin{theorem}[Lagrange's theorem on polynomial congruences]\label{thm:3.6.2} Let $F$ be a
	field and let $f\in F[x]$ be a polynomial of degree $n$ with nonzero
	leading coefficient. Then $f$ has at most $n$ roots in $F$. In particular,
	for a prime $p$, a polynomial $f(x)=\sum_{i=0}^{n} a_i x^{i}$ with integer
	coefficients and $p \nmid a_n$ has at most $n$ roots modulo $p$.

\end{theorem}

\begin{proof} We argue by induction on $n$. If $n=0$, then $f$ is a nonzero constant
	and has no root. Let $n\ge1$. If $f$ has no root, there is nothing to prove,
	so let $r\in F$ be a root. Dividing $f$ by the monic polynomial $x-r$ using the
	usual algorithm (at each step subtracting a multiple of $x-r$ to cancel
	the current leading term, which keeps every coefficient in $F$) gives
	$f(x)=(x-r)q(x)+c$ with $q\in F[x]$ of degree $n-1$ and $c\in F$; evaluating
	at $r$ gives $c=0$, so $f=(x-r)q$. If $s\neq r$ is another root, then
	$0=f(s)=(s-r)q(s)$ with $s-r$ nonzero, hence invertible, so $q(s)=0$. By the
	induction hypothesis $q$ has at most $n-1$ roots, so $f$ has at most $n$.

	For the special case modulo $p$, the ring $\mathbb{F}_p=\mathbb{Z}/p\mathbb{Z}$ is a field:
	every nonzero class has an integer representative not divisible by $p$,
	hence coprime to $p$, so the class is invertible by
	Theorem~\ref{thm:3.2.1}. The roots of $f$ modulo $p$ are exactly the
	roots in $\mathbb{F}_p$ of its reduction $\bar f$, and $\bar f$ has degree
	$n$ because $p\nmid a_n$.

\end{proof}
\begin{altproof}[Alternative proof of Theorem~\ref{thm:3.6.2}, modulo $p$] Assume $f(x)$ has $n+1$ distinct roots
	$x_1, \ldots, x_{n+1}$ modulo $p$. By Theorem~\ref{thm:3.6.1}, $f(x) \equiv (x
		- x_1) \cdots (x - x_n)\, g(x) \pmod{p}$ for some integer-coefficient
	polynomial $g$. Let $b x^m$ be the highest-degree term of $g$ whose
	coefficient is not divisible by $p$ (if no such term existed, then the
	right-hand side would be $0 \pmod{p}$, which is impossible because $p\nmid a_n$).
	Then, modulo $p$, the highest-degree term of the product on the right is
	$b x^{n+m}$, since the monic factor $(x-x_1)\cdots(x-x_n)$ has degree
	$n$. But $f$ has degree $n$ modulo $p$ because $p\nmid a_n$. Hence
	$n+m=n$, so $m=0$: modulo $p$, the polynomial $g$ is a nonzero constant, necessarily
	congruent to $a_n$ by comparing leading coefficients. Therefore
	$f(x) \equiv a_n \prod_{i=1}^{n}(x - x_i) \pmod{p}$. Substituting
	$x = x_{n+1}$ gives $0 \equiv a_n \prod_{i=1}^{n}(x_{n+1} - x_i)
		\pmod{p}$. Since $p \nmid a_n$ and each $x_{n+1} - x_i \not\equiv 0$,
	the product is not congruent to $0$, a contradiction. Therefore $f(x)$ has
	at most $n$ roots modulo $p$.

\end{altproof}
\newpage

\section{Digit Expansions, Floor Functions, and \texorpdfstring{$p$}{p}-adic Valuations}

\subsection{Base Representations and Lucas's Theorem}

\begin{definition}[base-$b$ representation]\label{def:4.1.1} Let $b\ge2$ be an integer.
	A base-$b$ representation of $n\in\mathbb{Z}^+$ is an expression
	$n=\sum_{i=0}^{k}d_i b^i$, where $k\in\mathbb{Z}_{\ge0}$,
	$d_i\in\{0,1,\ldots,b-1\}$ for every $i$, and $d_k\neq0$.
	For $n=0$, the representation is $(0)_b$. We write
	$n=(d_kd_{k-1}\cdots d_0)_b$ and call $d_i$ the $i$-th digit of $n$ in base $b$.

\end{definition}
\begin{theorem}\label{thm:4.1.1} Every integer $n \ge 0$ has a
	unique base-$b$ representation.

\end{theorem}

\begin{proof} We prove the claim by strong induction on $n$. By the
	division algorithm, $n = qb + r$ where $0 \le r < b$. For
	$0 \le n < b$, we have $q = 0$ and $r = n$, so $n = (d_0)_b$ with $d_0 = n$
	(in particular $n = (0)_b$ when $n = 0$) is the unique
	representation: a representation with at least two digits and no leading
	zero would have value at least $b$. Now assume that $n \ge b$ and that every integer
	less than $n$ has a unique base-$b$ representation. Then $q\ge1$ and $q<n$ (because $n=qb+r$ with $b\ge2$ and $0\le r<b$), so $q$ has a unique base-$b$
	representation $q = \sum_{i=0}^{k} d_i b^{i}$ with leading digit $d_k
		\neq 0$. Then $n = qb + r = \sum_{i=0}^{k} d_i b^{i+1} + r$; setting
	$d_0' = r$ and $d_i' = d_{i-1}$ for $1\le i\le k+1$ gives a base-$b$
	representation of $n$ with leading digit $d_{k+1}' = d_k \neq 0$ (no leading
	zero). Uniqueness follows by the same inductive argument: in any representation
	$n = \sum_{i=0}^{k}d_i b^i$ with $0 \le d_i < b$, reducing modulo $b$
	forces $d_0 \equiv n \pmod{b}$ with $0 \le d_0 < b$, so $d_0$ is the
	division-algorithm remainder $r$; then $(n - d_0)/b = q$ has
	the representation $\sum_{i \ge 1} d_i b^{i-1}$, which is unique by the inductive
	hypothesis, so every higher digit is determined as well.

\end{proof}
\begin{theorem}[Lucas's theorem]\label{thm:4.1.2} Let $p$ be a
	prime and $a, b \in \mathbb{Z}_{\ge0}$. Write both in base $p$ with $k+1$ digits
	(padding the shorter one with leading zeros): $a = (a_k
		\cdots a_0)_p$, $b = (b_k \cdots b_0)_p$ with $0 \le a_i, b_i < p$.
	Then $\binom{a}{b} \equiv \prod_{i=0}^{k} \binom{a_i}{b_i} \pmod{p}$.
	If $b_i>a_i$ for some $i$, by convention $\binom{a_i}{b_i}=0$ (and then $\binom{a}{b}\equiv 0 \pmod{p}$).

\end{theorem}

\begin{proof} All polynomial congruences in this proof are coefficientwise.
	By Theorem~\ref{thm:3.5.2}, $(1+x)^p\equiv1+x^p\pmod{p}$.
	Starting at $i=0$, raising to the $p$-th power and applying the same identity
	with $x$ replaced by $x^{p^i}$ gives, by induction,
	$(1+x)^{p^i}\equiv1+x^{p^i}\pmod{p}$ for every $i\ge0$. Hence
	\[
		\begin{aligned}
			(1+x)^a
			&=\prod_{i=0}^{k}\big((1+x)^{p^i}\big)^{a_i}\\
			&\equiv\prod_{i=0}^{k}(1+x^{p^i})^{a_i}
			=\prod_{i=0}^{k}\left(\sum_{j=0}^{a_i}\binom{a_i}{j}x^{jp^i}\right)
			\pmod{p}.
		\end{aligned}
	\]
	Let $b = \sum_{i=0}^{k} b_i p^{i}$. On
	expanding the right-hand side, we find that its coefficient of $x^{b}$ is
	$\sum \prod_{i=0}^{k}\binom{a_i}{j_i}$ over all tuples $(j_0,
		\ldots, j_k)$ with $\sum_{i} j_i p^{i} = b$, where each $j_i$ ranges
	over $0 \le j_i \le a_i \le p-1$ (terms with $j_i > a_i$ vanish). For
	such tuples, $\sum_{i=0}^{k}j_i p^i = \sum_{i=0}^{k}b_i p^i$ with all digits in
	$[0, p-1]$ forces $j_i = b_i$ for every $i$ by uniqueness of the
	base-$p$ representation (Theorem~\ref{thm:4.1.1}), so at most one tuple
	contributes (none if $b_i > a_i$ for some $i$, matching the convention
	$\binom{a_i}{b_i}=0$ above). Hence the left-hand side gives $\binom{a}{b}$ and the right-hand side
	gives $\prod_{i=0}^{k}\binom{a_i}{b_i}$. Therefore
	$\binom{a}{b} \equiv \prod_{i=0}^{k}\binom{a_i}{b_i} \pmod{p}$.

\end{proof}
\begin{theorem}\label{thm:4.1.3}
	Let $b\ge2$ be an integer and $x\in[0,1)$. Then there are digits
	$d_1,d_2,\ldots\in\{0,1,\ldots,b-1\}$ with $x=\sum_{i\ge1}d_ib^{-i}$. The
	expansion is unique except when $x=m/b^{k}$ for some integers $k\ge1$ and
	$0<m<b^{k}$, in which case $x$
	has exactly two expansions, one ending in all $0$'s and one ending in all
	$(b-1)$'s.

\end{theorem}

\begin{proof} For existence, set
	$d_i=\lfloor b^i x\rfloor-b\lfloor b^{i-1}x\rfloor$ for $i\ge1$.
	Splitting $b^{i-1}x$ into its integer and fractional parts gives
	$d_i=\lfloor b\{b^{i-1}x\}\rfloor\in\{0,1,\ldots,b-1\}$.
	Since $\lfloor x\rfloor=0$, the partial sums telescope, and for every $k\ge1$,
	\[
		\sum_{i=1}^{k}d_i b^{-i}=b^{-k}\lfloor b^k x\rfloor,
		\qquad
		0\le x-\sum_{i=1}^{k}d_i b^{-i}=b^{-k}\{b^k x\}<b^{-k}.
	\]
	Letting $k\to\infty$ establishes existence.

	For uniqueness, all digit series converge absolutely by comparison with
	$\sum_{i\ge1}(b-1)b^{-i}$. Suppose two distinct digit sequences $(d_i)$ and
	$(e_i)$ give the same value. Let $k$ be the first index at which they differ and,
	after interchanging them if necessary, suppose $d_k>e_k$. Then
	\[
		\begin{aligned}
			0&=\sum_{i\ge k}(d_i-e_i)b^{-i}\\
			 &\ge(d_k-e_k)b^{-k}-\sum_{i>k}(b-1)b^{-i}\\
			 &=(d_k-e_k-1)b^{-k}\ge0.
		\end{aligned}
	\]
	Equality throughout forces $d_k-e_k=1$, as well as $d_i=0$ and $e_i=b-1$ for every
	$i>k$. Indeed, each tail term attains its lower bound; otherwise the first
	inequality would be strict. Thus one expansion terminates, with its last nonzero digit in position $k$, and the common value is $m/b^k$ with $0<m<b^k$.

	Conversely, if $x=m/b^k$ with $0<m<b^k$, the constructed expansion has
	$d_i=0$ for all $i>k$. Let $\ell$ be the position of its last nonzero digit.
	Keeping the earlier digits, replacing $d_\ell$ by $d_\ell-1$, and replacing
	all later zeros by $b-1$ gives a second expansion, since
	$\sum_{i>\ell}(b-1)b^{-i}=b^{-\ell}$.
	There is only one terminating digit sequence: two such sequences could be
	compared after multiplication by a common power of $b$, using uniqueness
	of finite base-$b$ representations with leading zeros omitted
	(Theorem~\ref{thm:4.1.1}). When any other expansion is compared with this terminating
	one, the equality case above forces the first difference to occur at
	$\ell$ and forces precisely the second expansion just constructed.
	Hence there are exactly two expansions in this case, and there is no ambiguity in any other case. In particular, $x=0$ has only the all-zero expansion.

\end{proof}
\subsection{Digit Sums and Divisibility}

\begin{definition}[digit sum]\label{def:4.2.1} For an integer $b\ge2$ and $n\in
		\mathbb{Z}_{\ge0}$, the digit sum of $n$ in base $b$ is $S_b(n) =
		\sum_{i=0}^{k} d_i$, where $n = (d_k d_{k-1} \cdots d_0)_b$ is the
	base-$b$ representation of $n$. Leading zeros contribute nothing, so the value does not depend on how many of them are written,
	and $S_b(0) = 0$. For a prime $p$, the digit sum $S_p(n)$ is thus computed from the
	unique base-$p$ representation $n=\sum_{i\ge0}d_ip^i$ of
	Theorem~\ref{thm:4.1.1}, where $0\le d_i\le p-1$ and all but finitely many
	$d_i$ vanish.

\end{definition}
\begin{theorem}\label{thm:4.2.1} For any polynomial $P$ with integer
	coefficients and any $a, b \in \mathbb{Z}$, we have $(a - b) \mid (P(a) - P(b))$.

\end{theorem}

\begin{proof} Let $P(x) = \sum_{i=0}^{n} c_i x^{i}$ with
	$c_i \in \mathbb{Z}$. Then $P(a) - P(b) = \sum_{i=1}^{n} c_i (a^{i} -
		b^{i})$. Since $(a - b) \mid (a^{i} -
		b^{i})=(a-b)(a^{i-1}+a^{i-2}b+\cdots+b^{i-1})$ for all $i \ge 1$, we
	have $(a - b) \mid (P(a) - P(b))$.

\end{proof}
\begin{theorem}\label{thm:4.2.2} For $b \ge 2$ and $n \in
		\mathbb{Z}_{\ge0}$, we have $(b - 1) \mid (n - S_b(n))$.

\end{theorem}

\begin{proof} Let $n = (d_k \cdots d_0)_b$ and $P(x) =
		\sum_{i=0}^{k} d_i x^{i}$. Then $P(b) = n$ and $P(1) = S_b(n)$. By
	Theorem~\ref{thm:4.2.1}, $(b - 1) \mid (P(b) - P(1)) = n - S_b(n)$.

\end{proof}
\begin{theorem}\label{thm:4.2.3}
	Let $b\ge2$ and let $x_1,x_2,\ldots,x_r$ be nonnegative integers. Then
	\[
		S_b\!\left(\sum_{j=1}^{r}x_j\right)\le\sum_{j=1}^{r}S_b(x_j).
	\]
	In particular, $S_b(x+y)\le S_b(x)+S_b(y)$ for all nonnegative integers
	$x,y$.

\end{theorem}

\begin{proof} It suffices to prove the two-summand case,
	since the general case then follows by induction on $r$. Write
	$x=\sum_{i\ge0}x_i b^i$ and $y=\sum_{i\ge0}y_i b^i$,
	where $0\le x_i,y_i\le b-1$ and all
	but finitely many digits vanish. Set $c_0=0$ and, recursively, divide
	$x_i+y_i+c_i$ by $b$: let $c_{i+1}$ be the quotient and let
	$d_i\in\{0,\ldots,b-1\}$ be the remainder. These are the outgoing carry
	and the digit in the $i$-th position, respectively, so
	\[
		x_i+y_i+c_i=d_i+b c_{i+1}
	\]
	for every $i\ge0$. By induction, $0\le c_i\le1$ (since
	$0\le x_i+y_i+c_i\le2b-1$). Choose $M$ such that $x_i=y_i=0$ for all
	$i\ge M$. Then $d_M=c_M$ and $c_{M+1}=0$, since $c_M<b$.
	Thus $c_i=d_i=0$ for all $i\ge M+1$, so all sums below are finite.
	Multiplying the carry relation by $b^{i}$ and summing
	over $i\ge0$ gives
	$x+y+\sum_{i\ge1}c_i b^{i}=\sum_{i\ge0}d_i b^{i}+\sum_{i\ge1}c_i b^{i}$
	(the carry sums coincide after the shift $i\mapsto i+1$), so
	$x+y=\sum_{i\ge0}d_i b^{i}$. By uniqueness of the base-$b$
	expansion (Theorem~\ref{thm:4.1.1}), the $d_i$ are the base-$b$ digits
	of $x+y$, that is, the algorithm computes $x+y$. Summing the carry
	relation over all $i\ge0$ gives
	\[
		S_b(x)+S_b(y)+\sum_{i\ge1}c_i
		=S_b(x+y)+b\sum_{i\ge1}c_i.
	\]
	Rearranging gives
	$S_b(x+y)=S_b(x)+S_b(y)-(b-1)\sum_{i\ge1}c_i\le S_b(x)+S_b(y)$.

\end{proof}
\begin{theorem}\label{thm:4.2.4} Let $b\ge2$ and $n\ge1$ be integers. Then:

	\begin{enumerate}[label={\alph*)}]
		\item for every integer $k$ with $1\le k\le b^{n}$,
		      we have $S_b\big(k(b^{n}-1)\big)=n(b-1)$;
		\item every positive multiple $m$ of $b^{n}-1$ satisfies
		      $S_b(m)\ge n(b-1)$.
	\end{enumerate}

\end{theorem}

\begin{proof}

	\begin{enumerate}[label={\alph*)}]
		\item Write
		      $k(b^{n}-1)=(k-1)b^{n}+(b^{n}-k)$. Both $k-1$ and $b^{n}-k$ lie in
		      $\{0,1,\ldots,b^{n}-1\}$, so the two summands occupy disjoint blocks
		      of base-$b$ digits, where the lower block is padded on the left with zeros
		      so that it contains exactly $n$ digits. Hence
		      \[
			      S_b\big(k(b^{n}-1)\big)=S_b(k-1)+S_b(b^{n}-k).
		      \]
		      Put $j=k-1$. Since $b^{n}-k=(b^{n}-1)-j$, subtracting the $n$-digit
		      base-$b$ representation of $j$ from the $n$-digit number
		      $\big((b-1)(b-1)\cdots(b-1)\big)_b$ replaces each digit $d$ by $b-1-d$.
		      Consequently $S_b(j)+S_b((b^{n}-1)-j)=n(b-1)$.
		\item Group the base-$b$ digits of $m$ into blocks of $n$ digits
		      from the right:
		      \[
			      m=\sum_{i=0}^{r}B_i b^{ni},\qquad 0\le B_i\le b^{n}-1,
		      \]
		      and put $m'=\sum_{i=0}^{r}B_i$. Because
		      $b^{n}\equiv1\pmod{b^{n}-1}$, we have
		      $m'\equiv m\equiv0\pmod{b^{n}-1}$. Moreover, $m'>0$, and Theorem~\ref{thm:4.2.3}
		      gives
		      \[
			      S_b(m')\le\sum_{i=0}^{r}S_b(B_i)=S_b(m).
		      \]
		      If $m\ge b^{n}$, then $B_i>0$ for some $i\ge1$, so
		      $B_i b^{ni}>B_i$ and therefore $m'<m$. Repeating the operation thus
		      produces a strictly decreasing sequence of positive multiples of
		      $b^{n}-1$ without increasing the digit sum. It must terminate at a
		      positive multiple below $b^{n}$, which can only be $b^{n}-1$.
		      Therefore $S_b(m)\ge S_b(b^{n}-1)=n(b-1)$.
	\end{enumerate}

\end{proof}
\begin{theorem}\label{thm:4.2.5}
	Let $n\ge1$ be an integer and write $A=10x+y$ with $x\in\mathbb{Z}$ and $y\in\{0,1,\ldots,9\}$. Then
	\[
		10n+1\mid A\iff 10n+1\mid x-ny,\qquad
		10n-1\mid A\iff 10n-1\mid x+ny.
	\]
	The case $n=1$ gives the tests for $11$ and for $9$; the case $n=2$ gives the
	tests for $21$ and for $19$ and, since $7\mid21$, also yields the classical test:
	$7\mid A$ if and only if $7\mid x-2y$; the case $n=3$ shows that
	$29\mid A$ if and only if $29\mid x+3y$.

\end{theorem}

\begin{proof} Put $m=10n+1$. Then
	\[
		10(x-ny)=10x-10ny=(A-y)-y(10n)=A-y(10n+1)=A-ym,
	\]
	so $10(x-ny)\equiv A\pmod{m}$. Since $\gcd(10,m)=1$, we get $m\mid A$ if and only
	if $m\mid 10(x-ny)$, that is, if and only if $m\mid x-ny$.

	Now put $m=10n-1$. Then
	\[
		10(x+ny)=(A-y)+10ny=A+y(10n-1)=A+ym,
	\]
	so $10(x+ny)\equiv A\pmod{m}$, and again $\gcd(10,m)=1$ gives the claim.
	The cases $n=1$ and $n=3$ are direct substitutions. For the divisibility test for $7$ when
	$n=2$, note that $7\mid 10n+1=21$, and $10(x-2y)=A-21y$ with $\gcd(10,7)=1$
	shows that $7\mid A$ if and only if $7\mid x-2y$.

\end{proof}
\begin{theorem}[Kempner's theorem]\label{thm:4.2.6} Fix a digit
	$d\in\{0,1,\ldots,9\}$ and let $\mathcal{A}_d$ be the set of positive integers
	whose decimal expansion contains no digit equal to $d$. Then
	$\sum_{a\in\mathcal{A}_d}\frac1a$ converges; in fact, the sum is less than $90$.

\end{theorem}

\begin{proof} Count the elements of $\mathcal{A}_d$ with exactly
	$j$ digits. The leading digit may be any of the at most $9$ admissible nonzero
	digits, and each of the remaining $j-1$ digits may be any of the at most $9$
	admissible digits; hence there are at most $9\cdot 9^{\,j-1}$ such integers.
	Each of them is at least $10^{\,j-1}$, so their reciprocals contribute at most
	\[
		\frac{9\cdot 9^{\,j-1}}{10^{\,j-1}}=9\left(\frac{9}{10}\right)^{j-1}.
	\]
	Summing the geometric series over $j\ge 1$ gives
	$\sum_{a\in\mathcal{A}_d}\frac1a\le 9\cdot\frac{1}{1-9/10}=90$. The
	inequality is strict: in the one-digit case, if $d\neq0$ there are only $8$
	admissible leading digits, and if $d=0$ the term $a=2$ contributes $\tfrac12$,
	strictly less than the bound $1$ used for it; hence the total is less than
	$90$.

\end{proof}
\begin{theorem}\label{thm:4.2.7} Fix a nonempty finite string $B$ of $L\ge1$
	decimal digits. Then the set of positive integers whose decimal expansion does
	not contain $B$ as a substring has natural density $0$: the proportion of all positive integers not exceeding $x$
	that belong to this set tends to $0$ as $x\to\infty$.

\end{theorem}

\begin{proof} For $j\ge1$, let $N_j$ be the number of positive integers with
	exactly $j$ decimal digits whose (ordinary, unpadded) decimal expansion does
	not contain $B$. Write $j=qL+r$ with $q\ge0$ and $0\le r<L$, and cut the
	$j$-digit string into a prefix of $r$ digits followed by $q$ consecutive
	blocks of length $L$. If the expansion avoids $B$, then in particular none of
	these $q$ blocks equals $B$, so
	\[
		N_j\le 10^{r}\bigl(10^{L}-1\bigr)^{q}
		=10^{\,j}\bigl(1-10^{-L}\bigr)^{q}.
	\]
	This gives only an upper bound: the count allows leading zeros and does not
	exclude occurrences of $B$ crossing block boundaries. Every admissible
	expansion is still counted. Since $q=(j-r)/L\ge j/L-1$, this gives $N_j\le c\,(10\rho)^{\,j}$, where
	$\rho=\bigl(1-10^{-L}\bigr)^{1/L}<1$ and $c=\bigl(1-10^{-L}\bigr)^{-1}$.
	Now let $x\ge10$ and choose $J$ with $10^{\,J-1}\le x<10^{\,J}$, so that
	every positive integer not exceeding $x$ has at most $J$ digits. Since
	$10\rho\ge9>1$, summing the geometric series gives
	\[
		\sum_{j=1}^{J}N_j\le c\sum_{j=1}^{J}(10\rho)^{\,j}
		\le\frac{10c\rho}{10\rho-1}\,10^{\,J}\rho^{\,J}.
	\]
	Since the number of positive integers not exceeding $x$ is exactly
	$\lfloor x\rfloor\ge10^{\,J-1}$, the proportion of those integers
	that avoid $B$ is at most a constant multiple of $\rho^{\,J}$,
	which tends to $0$ as $x\to\infty$ because $\rho<1$. Hence the density
	is $0$.

\end{proof}
\subsection{The Floor and Nearest-Integer Functions}

\begin{theorem}\label{thm:4.3.1} For
	$x, y\in\mathbb{R}$, $n\in\mathbb{Z}$, and $m\in\mathbb{Z}^+$:

	\begin{enumerate}[label={\alph*)}]
		\item $\lfloor n+x\rfloor = n + \lfloor x\rfloor$;
		\item $\big\lfloor \tfrac{\lfloor x\rfloor}{m}\big\rfloor =
			      \big\lfloor \tfrac{x}{m}\big\rfloor$;
		\item $\lfloor x\rfloor + \lfloor y\rfloor \le \lfloor x+y\rfloor
			      \le \lfloor x\rfloor + \lfloor y\rfloor + 1$;
		\item $-\lfloor -x\rfloor = \lceil x\rceil$, and $\lfloor x\rfloor
			      + \lfloor -x\rfloor$ equals $0$ if $x\in\mathbb{Z}$ and $-1$ otherwise;
		\item if $x\le y$, then $\lfloor x\rfloor \le \lfloor y\rfloor$;
		\item if $x\ge0$ and $y\ge0$, then $\lfloor x\rfloor\,\lfloor y\rfloor
			      \le \lfloor xy\rfloor$;
		\item $\lVert x\rVert = \min\{\{x\}, 1-\{x\}\}$ and $\lVert
			      x+y\rVert \le \lVert x\rVert + \lVert y\rVert$.
	\end{enumerate}

\end{theorem}

\begin{proof}
	\begin{enumerate}[label={\alph*)}]
		\item Let $q=\lfloor x\rfloor$. Then $q\le x<q+1$, so
		      $n+q\le n+x<n+q+1$ and $\lfloor n+x\rfloor=n+q$.
		\item By the division algorithm, write $\lfloor x\rfloor=mq+s$
		      with $q\in\mathbb{Z}$ and $0\le s<m$. Then
		      $x=mq+s+\{x\}$ with $0\le s+\{x\}<m$.
		      Thus both $\lfloor\lfloor x\rfloor/m\rfloor$ and
		      $\lfloor x/m\rfloor$ equal $q$. The same argument applies to negative $x$.
		\item By part (a),
		      \[
			      \lfloor x+y\rfloor=\lfloor x\rfloor+\lfloor y\rfloor
			      +\lfloor\{x\}+\{y\}\rfloor.
		      \]
		      Since $0\le\{x\}+\{y\}<2$, the last term is $0$ or $1$,
		      giving the two inequalities.
		\item The inequalities $\lfloor-x\rfloor\le-x<\lfloor-x\rfloor+1$
		      are equivalent to $-\lfloor-x\rfloor-1<x\le-\lfloor-x\rfloor$.
		      Hence $\lceil x\rceil=-\lfloor-x\rfloor$.
		      If $x\in\mathbb{Z}$, then $\lfloor x\rfloor+\lfloor-x\rfloor=0$;
		      otherwise $\lceil x\rceil=\lfloor x\rfloor+1$, so the sum is $-1$.
		\item By Definition~\ref{def:1.2.1}, $\lfloor x\rfloor\le x\le y$, so
		      $\lfloor x\rfloor$ is an integer not exceeding $y$. Since
		      $\lfloor y\rfloor$ is the greatest such integer,
		      $\lfloor x\rfloor\le\lfloor y\rfloor$.
		\item Since $x,y\ge0$, we have $0\le\lfloor x\rfloor\le x$ and
		      $0\le\lfloor y\rfloor\le y$, so
		      $\lfloor x\rfloor\lfloor y\rfloor\le xy$. The left-hand side is an
		      integer not exceeding $xy$, hence it is at most $\lfloor xy\rfloor$.
		      The hypothesis cannot be dropped: for $x=y=-\tfrac12$, we have
		      $\lfloor x\rfloor\lfloor y\rfloor=1$, whereas
		      $\lfloor xy\rfloor=\lfloor\tfrac14\rfloor=0$.
		\item Put $q=\lfloor x\rfloor$. Every integer at most $q$ is at least
		      as far from $x$ as $q$, and every integer at least $q+1$ is at
		      least as far from $x$ as $q+1$. Thus the minimum distance is
		      attained at $q$ or $q+1$, whose distances from $x$ are $\{x\}$ and
		      $1-\{x\}$. Hence $\lVert x\rVert=\min\{\{x\},1-\{x\}\}$;
		      in particular, this gives $0$ when $x$ is an integer.
		      Choose integers $j,k$ attaining the minima for $x,y$, respectively. Then
		      \[
			      \lVert x+y\rVert\le |(x+y)-(j+k)|
			      \le |x-j|+|y-k|=\lVert x\rVert+\lVert y\rVert.
		      \]
	\end{enumerate}
\end{proof}
\begin{theorem}\label{thm:4.3.2} For integers $b\ge2$ and
	$n\ge1$,
	\[
		S_b(n)\le(b-1)(\lfloor\log_b n\rfloor+1).
	\]

\end{theorem}

\begin{proof} Write $n=(d_kd_{k-1}\cdots d_0)_b$ with
	$d_k\neq0$. Then
	\[
		b^k\le n\le\sum_{i=0}^{k}(b-1)b^i=b^{k+1}-1<b^{k+1}.
	\]
	Since $b>1$, taking logarithms to base $b$ preserves the inequalities,
	so $k\le\log_b n<k+1$. By Definition~\ref{def:1.2.1},
	$k=\lfloor\log_b n\rfloor$. As every digit satisfies
	$0\le d_i\le b-1$,
	\[
		S_b(n)=\sum_{i=0}^{k}d_i\le(b-1)(k+1)
		=(b-1)(\lfloor\log_b n\rfloor+1).
	\]

\end{proof}
\begin{theorem}[Hermite's identity]\label{thm:4.3.3} For
	$x\in\mathbb{R}$ and $m\in\mathbb{Z}^+$, we have $\lfloor mx\rfloor =
		\sum_{i=0}^{m-1}\Big\lfloor x + \tfrac{i}{m}\Big\rfloor$.

\end{theorem}

\begin{proof} Let $g(x) = \sum_{i=0}^{m-1}\lfloor
		x+\tfrac{i}{m}\rfloor - \lfloor mx\rfloor$. Then
	$g\big(x+\tfrac{1}{m}\big) = \sum_{i=1}^{m}\lfloor x+\tfrac{i}{m}\rfloor
		- \lfloor mx+1\rfloor = g(x) + \lfloor x+1\rfloor - \lfloor x\rfloor
		- 1 = g(x)$, so $g$ has period $\tfrac{1}{m}$. For $x\in[0,\tfrac{1}{m})$,
	we have $\lfloor x+\tfrac{i}{m}\rfloor = 0$ ($0\le i\le m-1$) and
	$\lfloor mx\rfloor = 0$, so $g(x) = 0$. Hence the identity holds for
	all $x$.

\end{proof}
\begin{theorem}\label{thm:4.3.4}
	Let $m,n\ge 1$ be integers, let $x\in\mathbb{Z}$, and put $d=\gcd(m,n)$. Then
	\[
		\sum_{i=0}^{n-1}\Bigl\lfloor\frac{x+mi}{n}\Bigr\rfloor
		=\sum_{j=0}^{m-1}\Bigl\lfloor\frac{x+nj}{m}\Bigr\rfloor
		=x-(x\bmod d)+\frac{mn-m-n+d}{2}.
	\]
	For $m=1$, this is Hermite's identity (Theorem~\ref{thm:4.3.3})
	at the argument $x/n$.

\end{theorem}

\begin{proof} Write $r_i=(x+mi)\bmod n$, so that
	$\lfloor (x+mi)/n\rfloor=(x+mi-r_i)/n$ and
	\[
		\sum_{i=0}^{n-1}\Bigl\lfloor\frac{x+mi}{n}\Bigr\rfloor
		=\frac{1}{n}\Bigl(nx+m\frac{n(n-1)}{2}-\sum_{i=0}^{n-1}r_i\Bigr).
	\]
	As $i$ runs over a complete residue system modulo $n$, the residue $mi\bmod n$
	runs over the multiples of $d$ modulo $n$, each attained exactly $d$ times:
	by the linear congruence theorem (Theorem~\ref{thm:3.2.3}), $mi\equiv r\pmod{n}$
	is solvable if and only if $d\mid r$, and then it has exactly $d$ solutions
	modulo $n$.
	Hence the sequence $r_0,\ldots,r_{n-1}$ contains each of the $n/d$ residues
	$x_0,x_0+d,\ldots,x_0+n-d$ exactly $d$ times, where $x_0:=x\bmod d$. Therefore
	\[
		\sum_{i=0}^{n-1}r_i=d\sum_{s=0}^{n/d-1}(x_0+ds)
		=d\Bigl(\frac{n}{d}x_0+d\cdot\frac{(n/d)(n/d-1)}{2}\Bigr)
		=nx_0+\frac{n(n-d)}{2}.
	\]
	Substituting back gives
	\[
		\sum_{i=0}^{n-1}\Bigl\lfloor\frac{x+mi}{n}\Bigr\rfloor
		=x-x_0+\frac{m(n-1)-(n-d)}{2}
		=x-x_0+\frac{mn-m-n+d}{2}.
	\]
	The final expression is symmetric under interchange of $m$ and $n$ (note that
	$d=\gcd(m,n)$ and $x_0=x\bmod d$ are symmetric), so the two sums
	agree. When $m=1$, we have $d=1$, and the formula becomes
	$\sum_{i=0}^{n-1}\lfloor x/n+i/n\rfloor=x$; this is
	Theorem~\ref{thm:4.3.3} with the argument $x/n$ and $n$ summands.

\end{proof}
\subsection{Base-\texorpdfstring{$m$}{m} Expansions of Fractions: Digits, Remainders, and Periodicity}

\begin{definition}[base-$m$ expansion of a fraction]\label{def:4.4.1} Fix an integer $m\ge2$. For $n\in\mathbb{Z}^+$ and $a\in\mathbb{Z}$ with $0\le a<n$, define $r_0=a$ and, for $k\ge1$,
	\[
		d_k=\left\lfloor\frac{m r_{k-1}}{n}\right\rfloor,\qquad r_k=m r_{k-1}-d_kn.
	\]
	The base-$m$ expansion of $a/n$ is $0.d_1d_2d_3\ldots$, with digits
	interpreted in base $m$. Throughout this subsection, we use the digit sequence
	chosen by this recurrence. Once a remainder is $0$, every subsequent digit
	and remainder is $0$, so a terminating expansion is continued with zeros.
	For a general rational number, first separate its sign, then split its
	absolute value into a nonnegative integer part and a fractional part in
	$[0,1)$. Apply the recurrence to that fractional part; all later statements
	about periodicity concern its digit sequence.

\end{definition}
\begin{theorem}\label{thm:4.4.1}
	For the digits and remainders in Definition~\ref{def:4.4.1},
	$0\le d_k\le m-1$ and $0\le r_k<n$ for every integer $k\ge1$.
	Moreover, for every integer $K\ge0$,
	\[
		\frac{a}{n}=\sum_{k=1}^{K}\frac{d_k}{m^k}
		+\frac{r_K}{nm^K},\qquad
		r_K=n\left\{\frac{m^Ka}{n}\right\},
	\]
	where the sum is empty when $K=0$. Consequently,
	\[
		\frac{a}{n}=\sum_{k\ge1}\frac{d_k}{m^k},
	\]
	and a positive integer $P$ is a period from the first digit, meaning that
	$d_{P+t}=d_t$ for every $t\ge1$, if and only if $r_P=r_0$.

\end{theorem}

\begin{proof} If $0\le r_{k-1}<n$, then
	$0\le mr_{k-1}<mn$. Division with remainder by $n$ therefore gives
	$0\le d_k\le m-1$ and $0\le r_k<n$; the induction starts from
	$0\le r_0=a<n$. The recurrence also gives
	\[
		\frac{r_{k-1}}{nm^{k-1}}=\frac{d_k}{m^k}
		+\frac{r_k}{nm^k}.
	\]
	Summing from $k=1$ to $K$ gives the finite identity by telescoping. Since
	$0\le r_K<n$ and $m\ge2$, its remainder term tends to $0$ as
	$K\to\infty$, proving the infinite expansion. Multiplying the finite
	identity by $m^K$ gives
	\[
		\frac{m^Ka}{n}=\sum_{k=1}^{K}d_km^{K-k}+\frac{r_K}{n}.
	\]
	The sum is an integer and $0\le r_K/n<1$, which proves the formula for
	$r_K$.

	If $r_P=r_0$, the recurrence, in which each remainder uniquely determines
	the next digit and remainder, gives $r_{P+t}=r_t$ for $t\ge0$ and
	$d_{P+t}=d_t$ for $t\ge1$. Conversely, suppose $P$ is a period from the
	first digit and put $D=d_1m^{P-1}+\cdots+d_P$. Summing the repeated blocks
	in the infinite expansion gives
	\[
		\frac{a}{n}=\frac{D}{m^P-1},
	\]
	so $m^Pa/n=D+a/n$. The finite identity with $K=P$ also gives
	$m^Pa/n=D+r_P/n$. Hence $r_P=a=r_0$.

\end{proof}
\begin{theorem}\label{thm:4.4.2} Let $m,n,a$ be integers with $m\ge2$,
	$n\ge1$, and $0\le a<n$. The sequence $(d_k)_{k\ge1}$ in
	Definition~\ref{def:4.4.1} is eventually periodic: there exist integers
	$L\ge0$ and $P\ge1$ with $d_{k+P}=d_k$ for every integer $k>L$.

\end{theorem}

\begin{proof} Since $r_k\in\{0,1,\ldots,n-1\}$ for every $k$, two of $r_0,r_1,\ldots,r_n$ coincide, say $r_i=r_j$ with $0\le i<j\le n$. The recurrence shows that each remainder uniquely determines the next digit and the next remainder. Therefore $r_{i+t}=r_{j+t}$ and $d_{i+t+1}=d_{j+t+1}$ for every $t\ge0$. Taking $L=i$ and $P=j-i$ proves the assertion.

\end{proof}
\begin{theorem}\label{thm:4.4.3} Let $m\ge2$ be an integer, and let
	$d_1,d_2,\ldots\in\{0,1,\ldots,m-1\}$ be digits for which there are
	integers $L\ge0$ and $P\ge1$ with $d_{k+P}=d_k$ for every integer $k>L$.
	Then $x=\sum_{k\ge1}d_km^{-k}$ is rational; explicitly,
	\[
		x=\frac{1}{m^{L}}\left(A+\frac{B}{m^{P}-1}\right),\qquad
		A=\sum_{k=1}^{L}d_km^{L-k},\qquad
		B=\sum_{k=1}^{P}d_{L+k}m^{P-k},
	\]
	where $A=0$ when $L=0$. Consequently, a real number $x\in[0,1)$ is
	rational if and only if $x=\sum_{k\ge1}d_km^{-k}$ for some eventually
	periodic sequence of base-$m$ digits $d_k\in\{0,1,\ldots,m-1\}$.

\end{theorem}

\begin{proof} The series converges, because its terms are nonnegative and
	bounded by $(m-1)m^{-k}$. Multiplying by $m^{L}$ gives
	\[
		m^{L}x=\sum_{k=1}^{L}d_km^{L-k}+\sum_{j\ge1}e_jm^{-j}=A+y,
		\qquad e_j=d_{L+j},\quad y=\sum_{j\ge1}e_jm^{-j}.
	\]
	The sequence $(e_j)$ satisfies $e_{j+P}=e_j$ for every $j\ge1$. A
	convergent series of nonnegative terms may be summed in consecutive
	blocks of $P$ terms, and the block with indices $tP+1,\ldots,tP+P$ has
	sum $m^{-tP}\sum_{i=1}^{P}e_im^{-i}=m^{-tP}B/m^{P}$. Therefore
	\[
		y=\frac{B}{m^{P}}\sum_{t\ge0}m^{-tP}
		=\frac{B}{m^{P}}\cdot\frac{1}{1-m^{-P}}
		=\frac{B}{m^{P}-1},
	\]
	and $x=m^{-L}\bigl(A+B/(m^{P}-1)\bigr)$ is rational.

	For the final assertion, the first part shows that an eventually
	periodic digit sequence gives a rational number. Conversely, if
	$x\in[0,1)$ is rational, write $x=a/n$ with $n\in\mathbb{Z}^+$ and
	$0\le a<n$. The digits of Definition~\ref{def:4.4.1} satisfy
	$x=\sum_{k\ge1}d_km^{-k}$ by Theorem~\ref{thm:4.4.1}, and they are
	eventually periodic by Theorem~\ref{thm:4.4.2}.

\end{proof}
\begin{definition}[purely periodic]\label{def:4.4.2} A base-$m$ expansion is
	purely periodic if there is a positive integer $P$ such that
	$d_{k+P}=d_k$ for every integer $k\ge1$. Here $P$ need not be the least
	positive period.

\end{definition}
\begin{theorem}\label{thm:4.4.4} Let $m,n,a$ be integers with $m\ge2$,
	$0<a<n$, and $\gcd(a,n)=1$. The base-$m$ expansion of $a/n$ is purely
	periodic if and only if $\gcd(m,n)=1$.

\end{theorem}

\begin{proof} By Theorem~\ref{thm:4.4.1}, $r_k$ is the least
	nonnegative residue of $m^ka$ modulo $n$, and a positive integer $P$ is a
	period from the first digit if and only if $r_P=r_0$. If $\gcd(m,n)=1$,
	multiplication by $m$ is a bijection modulo $n$. Hence the finite orbit of
	$a$ under this bijection returns to $a$: among the residues
	$a, ma, m^{2}a, \ldots$ modulo $n$, two must coincide, say
	$m^{i}a \equiv m^{j}a \pmod{n}$ with $i<j$, and cancelling the unit
	$m^{i}$ gives $m^{j-i}a \equiv a \pmod{n}$, that is, $r_{j-i}=r_0$;
	so the expansion is purely
	periodic. Conversely, if the expansion has period $P$, then $r_P=r_0=a$,
	so $m^Pa\equiv a\pmod{n}$. Since $\gcd(a,n)=1$, Euclid's lemma (Theorem~\ref{thm:2.2.8}) gives
	$n\mid m^P-1$. A prime dividing both $m$ and $n$ would then divide both
	$m^P$ and $m^P-1$, hence would divide $1$, a contradiction. Therefore
	$\gcd(m,n)=1$.

\end{proof}
\subsection{The \texorpdfstring{$p$}{p}-adic Valuation}

\begin{definition}[$p$-adic valuation]\label{def:4.5.1}
	For a prime $p$ and a nonzero integer $n$, the $p$-adic valuation
	$v_p(n)$ is the exponent of $p$ in the prime factorization of $|n|$;
	equivalently, $v_p(n)$ is the largest integer $\alpha \ge 0$ with
	$p^{\alpha} \mid n$. By convention, $v_p(0)=+\infty$, taken to be greater than every integer.
	This convention allows the valuation inequality to apply when a sum is zero; the power and quotient
	formulas retain the hypothesis that their inputs are nonzero.
	If $p_1,p_2,\ldots,p_r$ are the distinct prime divisors of a nonzero
	integer $n$, then the fundamental theorem of arithmetic (Theorem~\ref{thm:2.3.3}) gives
	$|n|=\prod_{i=1}^{r}p_i^{v_{p_i}(n)}$.

\end{definition}
\begin{theorem}\label{thm:4.5.1} Fix a prime $p$, and let $a,b$ be
	nonzero integers. Then:

	\begin{enumerate}[label={\alph*)}]
		\item For every $\alpha\in\mathbb{Z}_{\ge0}$, we have $v_p(a)=\alpha$ if and only if
		      $a=kp^{\alpha}$ for some $k\in\mathbb{Z}$ with $\gcd(k,p)=1$.
		\item $v_p(ab) = v_p(a) + v_p(b)$.
		\item $v_p(a + b) \ge \min\{v_p(a), v_p(b)\}$, with equality when
		      $v_p(a) \neq v_p(b)$.
		\item $v_p(a^{k}) = k\, v_p(a)$ for every integer $k \ge 0$.
	\end{enumerate}

\end{theorem}

\begin{proof}

	\begin{enumerate}[label={\alph*)}]
		\item ($\Rightarrow$) If $v_p(a) = \alpha$, then $p^{\alpha} \mid
			      a$ but $p^{\alpha+1} \nmid a$, so $a = k p^{\alpha}$ with $p \nmid
			      k$. ($\Leftarrow$) If $a = k p^{\alpha}$ with $\gcd(k, p) = 1$, then
		      $p^{\alpha} \mid a$ and $p^{\alpha+1} \nmid a$, so $v_p(a) = \alpha$.
		\item Write $a = k_1 p^{v_p(a)}$, $b = k_2 p^{v_p(b)}$ with
		      $\gcd(k_1, p) = \gcd(k_2, p) = 1$. Then $ab = k_1 k_2 p^{v_p(a) +
					      v_p(b)}$ with $\gcd(k_1 k_2, p) = 1$, so $v_p(ab) = v_p(a) + v_p(b)$.
		\item If $a+b=0$, then $v_p(a+b)=+\infty\ge\min\{v_p(a),v_p(b)\}$.
		      Here $b=-a$, so $v_p(a)=v_p(b)$ and the equality clause does not apply.
		      Otherwise, interchange $a,b$ if necessary so that
		      $\alpha=v_p(a)\ge v_p(b)=\beta$. By part (a), write
		      $a=k_1p^{\alpha}$ and $b=k_2p^{\beta}$ with $p\nmid k_1k_2$.
		      Then $a+b=p^{\beta}(k_1p^{\alpha-\beta}+k_2)$, so
		      $v_p(a+b)\ge\beta$. If $\alpha>\beta$, the expression in parentheses
		      is congruent to $k_2\not\equiv0\pmod p$, so $v_p(a+b)=\beta$.
		\item For $k=0$, we have $v_p(a^0)=v_p(1)=0$.
		      The induction step follows from
		      $v_p(a^{k+1})=v_p(a^k)+v_p(a)$ by part (b).\qedhere
	\end{enumerate}

\end{proof}
\begin{theorem}\label{thm:4.5.2} For $x, y \in \mathbb{Z}$, we have $x \mid y$ if and only if $v_p(x) \le v_p(y)$
	for all primes $p$.

\end{theorem}

\begin{proof} ($\Rightarrow$) If $x=0$ then $y=0$ and
	$v_p(x)=v_p(y)=+\infty$ for every $p$; if $y=0$ then
	$v_p(y)=+\infty\ge v_p(x)$ for every $p$. So assume $x,y\neq0$; then
	$y=cx$ with $c\neq0$, and Theorem~\ref{thm:4.5.1}(b) gives
	$v_p(y)=v_p(c)+v_p(x)\ge v_p(x)$ for all $p$.
	($\Leftarrow$) Suppose $v_p(x) \le v_p(y)$ for all primes $p$. If $x = 0$, then $v_p(x) = +\infty$ forces $v_p(y) = +\infty$ for all $p$, so $y = 0$ and $x \mid y$; if $x \neq 0$ but $y = 0$, then $x \mid 0 = y$; and if $x,y$ are both nonzero, put $\lambda_p:=v_p(y)-v_p(x)\ge0$.
	Only finitely many $\lambda_p$ are nonzero, by unique factorization, so $|y| = \big(\prod_{p\in\Primes} p^{\lambda_p}\big)|x|$ and hence $y = \pm\big(\prod_{p\in\Primes} p^{\lambda_p}\big) x$, giving $x \mid y$ for either choice of sign.

\end{proof}
\begin{theorem}\label{thm:4.5.3} For $x, y \in \mathbb{Z}$, we have $v_p(x) = v_p(y)$ for all primes $p$ if and only if $|x| = |y|$ (equivalently, $x = \pm y$); in particular, for $x, y \in \mathbb{Z}^+$, we have $x = y$ if and only if $v_p(x) = v_p(y)$ for all primes $p$.

\end{theorem}

\begin{proof} Note that $v_p(x) = v_p(y)$ for all $p$ if and only if $v_p(x) \le v_p(y)$ and $v_p(y) \le v_p(x)$ for all $p$, which holds if and only if $x \mid y$ and $y \mid x$ (Theorem~\ref{thm:4.5.2}), which in turn holds if and only if $x = \pm y$ (Theorem~\ref{thm:1.1.2}(d)), that is, $|x| = |y|$; when $x, y > 0$ this reads $x = y$.

\end{proof}
\begin{definition}[$p$-adic valuation on $\mathbb{Q}^{\times}$]\label{def:4.5.2} Let $q \in \mathbb{Q}^{\times}$, the multiplicative group
	of nonzero rationals. Choose
	$a,b\in\mathbb{Z}^+$ with $q=\pm a/b$ and set
	$v_p(q)=v_p(a)-v_p(b)$.

\end{definition}
\begin{theorem}\label{thm:4.5.4} The value prescribed in Definition~\ref{def:4.5.2} is
	independent of the representation $q=\pm a/b$; hence
	$v_p:\mathbb{Q}^{\times}\to\mathbb{Z}$ is well defined. Moreover,
	$v_p(q)=v_p(-q)$ for every $q\in\mathbb{Q}^{\times}$.

\end{theorem}

\begin{proof} Suppose $q=\pm a/b=\pm c/d$, with
	$a,b,c,d \in \mathbb{Z}^+$ and the same sign in both fractions, as in
	Definition~\ref{def:4.5.2}. Cross-multiplication gives $ad = bc$. Hence $v_p(ad) =
		v_p(bc)$, that is, $v_p(a) + v_p(d) = v_p(b) + v_p(c)$, so
	$v_p(a) - v_p(b) = v_p(c) - v_p(d)$. Replacing $q$ by $-q$ changes only
	the sign preceding $a/b$, which does not enter
	the definition of the valuation; therefore $v_p(q)=v_p(-q)$.

\end{proof}
\begin{theorem}\label{thm:4.5.5} For $q, r \in \mathbb{Q}^{\times}$ and
	$p \in \Primes$:

	\begin{enumerate}[label={\alph*)}]
		\item $v_p(qr) = v_p(q) + v_p(r)$;
		\item if $q + r \neq 0$, then $v_p(q + r) \ge \min\{v_p(q), v_p(r)\}$, with equality if
		      $v_p(q) \neq v_p(r)$ (with the convention $v_p(0) = +\infty$, the inequality holds unconditionally);
		\item $v_p(q^{k}) = k\, v_p(q)$ for all $k \in \mathbb{Z}$.
	\end{enumerate}

\end{theorem}

\begin{proof} Write $q=a/b$ and $r=c/d$, where $a,c$ are nonzero
	integers and $b,d$ are positive integers. The invariance under a change of sign established in
	Theorem~\ref{thm:4.5.4} allows us to use these signed numerators.
	\begin{enumerate}[label={\alph*)}]
		\item The integer product formula gives
		      \[
			      v_p(qr)=v_p(ac)-v_p(bd)
			      =v_p(a)+v_p(c)-v_p(b)-v_p(d)=v_p(q)+v_p(r).
		      \]
		\item If $q+r=0$, the extended inequality reads
		      $+\infty\ge\min\{v_p(q),v_p(r)\}$.
		      In this case $r=-q$, so the two valuations are equal.
		      Now assume $q+r\neq0$. Since $ad+bc\neq0$, the integer sum
		      inequality gives
		      \[
			      \begin{aligned}
			      v_p(q+r)&=v_p(ad+bc)-v_p(bd)\\
			      &\ge\min\{v_p(ad),v_p(bc)\}-v_p(bd)\\
			      &=\min\{v_p(q),v_p(r)\}.
			      \end{aligned}
		      \]
		      Here $v_p(ad)=v_p(q)+v_p(bd)$ and
		      $v_p(bc)=v_p(r)+v_p(bd)$. Thus $v_p(q)\neq v_p(r)$ implies
		      $v_p(ad)\neq v_p(bc)$, so equality follows from
		      Theorem~\ref{thm:4.5.1}(c).
		\item For $k=0$ the formula is $v_p(1)=0$; for $k>0$ it follows by
		      induction using part (a). Also, by part (a),
		      $0=v_p(1)=v_p(qq^{-1})=v_p(q)+v_p(q^{-1})$, so
		      $v_p(q^{-1})=-v_p(q)$. If $k<0$, apply the positive-exponent
		      case to $(q^{-1})^{-k}$ to obtain $v_p(q^k)=k\,v_p(q)$.
	\end{enumerate}
\end{proof}
\begin{theorem}\label{thm:4.5.6}
	For every integer $n\ge2$, the sum $H_n=1+\tfrac12+\cdots+\tfrac1n$ is not an integer.

\end{theorem}

\begin{proof} Let $k\ge1$ be the unique integer with
	$2^{k}\le n<2^{k+1}$. For $1\le m\le n$, we have $v_2(1/m)=-v_2(m)$. If
	$m\neq2^{k}$ and $v_2(m)\ge k$, then $m$ is a multiple of $2^{k}$ other than
	$2^{k}$ itself, so $m\ge2\cdot2^{k}=2^{k+1}>n$, which is impossible. Hence
	$v_2(m)<k$, that is, $v_2(1/m)>-k$, for every $m\le n$ with $m\neq2^{k}$,
	whereas $v_2(1/2^{k})=-k$. Since the $2$-adic valuation satisfies
	$v_2(x+y)\ge\min\{v_2(x),v_2(y)\}$ with equality when the two valuations differ
	(Theorem~\ref{thm:4.5.5}(b)), iterating over the $n-1$ terms other than
	$1/2^{k}$ shows that their sum $S$ has
	$v_2(S)\ge\min_{m\neq2^{k}}v_2(1/m)>-k$; since $v_2(S)\neq v_2(1/2^{k})$,
	the equality case of Theorem~\ref{thm:4.5.5}(b) then gives
	$v_2(H_n)=\min\{v_2(S),-k\}=-k<0$. An integer has nonnegative
	$2$-adic valuation, so $H_n\notin\mathbb{Z}$.

\end{proof}
\subsection{Prime-Power Formulas for gcd and lcm: A Valuation-Theoretic Review}

\begin{theorem}\label{thm:4.6.1} For all $a,b\in\mathbb{Z}^+$ and
	every prime $p$,
	\[
		v_p(\gcd(a, b))=\min\{v_p(a),v_p(b)\},\qquad
		v_p(\lcm(a, b))=\max\{v_p(a),v_p(b)\}.
	\]
	Consequently,
	\[
		\gcd(a, b)=\prod_{p\in\Primes}p^{\min\{v_p(a), v_p(b)\}},\qquad
		\lcm(a, b)=\prod_{p\in\Primes}p^{\max\{v_p(a), v_p(b)\}},
	\]
	where all but finitely many factors equal $1$.

\end{theorem}

\begin{proof} Put $g=\gcd(a, b)$. Since $g\mid a$ and
	$g\mid b$, Theorem~\ref{thm:4.5.2} gives
	$v_p(g)\le\min\{v_p(a), v_p(b)\}$ for every prime $p$. Conversely, the
	positive integer
	$g_0=\prod_{p\in\Primes}p^{\min\{v_p(a), v_p(b)\}}$ divides both $a$
	and $b$, so Theorem~\ref{thm:2.2.6} gives $g_0\mid g$. Hence
	$\min\{v_p(a), v_p(b)\}=v_p(g_0)\le v_p(g)$ for every $p$, proving the
	gcd valuation formula. By Theorem~\ref{thm:4.5.3}, that formula also proves the
	stated product formula for $g$.

	Put $L=\lcm(a, b)$. Since $a\mid L$ and $b\mid L$,
	we have $v_p(L)\ge\max\{v_p(a), v_p(b)\}$. Conversely,
	$L_0=\prod_{p\in\Primes}p^{\max\{v_p(a), v_p(b)\}}$ is a common
	multiple of $a$ and $b$, so Theorem~\ref{thm:2.2.3} gives $L\mid L_0$. Thus
	$v_p(L)\le v_p(L_0)=\max\{v_p(a), v_p(b)\}$, proving both the lcm
	valuation formula and, by Theorem~\ref{thm:4.5.3}, the lcm product formula.

\end{proof}
\begin{theorem}\label{thm:4.6.2} For all $a, b, c\in\mathbb{Z}^+$,
	\[
		\gcd(a^{c}, b^{c})=(\gcd(a, b))^{c},\qquad
		\lcm(a^{c}, b^{c})=(\lcm(a, b))^{c}.
	\]

\end{theorem}

\begin{proof} Since $c>0$,
	\[
		\min\{c\,v_p(a),c\,v_p(b)\}
		=c\min\{v_p(a), v_p(b)\}.
	\]
	Theorem~\ref{thm:4.5.1}(d), Theorem~\ref{thm:4.6.1}, and
	Theorem~\ref{thm:4.5.3} therefore give the gcd identity;
	replacing $\min$ by $\max$ gives the lcm identity.

\end{proof}
\begin{theorem}\label{thm:4.6.3} For all $a, b\in\mathbb{Z}$ and $c\in\mathbb{Z}^+$, we have
	$a\mid b$ if and only if $a^{c}\mid b^{c}$.

\end{theorem}

\begin{proof} If $a=0$, then $a^{c}=0$, and each of the conditions $a\mid b$
	and $a^{c}\mid b^{c}$ holds exactly when $b=0$
	(Definition~\ref{def:1.1.1}); if $a\neq0$ and $b=0$, both conditions hold.
	Assume now that $a,b\neq0$. By Theorem~\ref{thm:4.5.1}(d), every prime $p$
	satisfies $v_p(a^{c})=c\,v_p(a)$ and $v_p(b^{c})=c\,v_p(b)$, and since
	$c>0$, the inequality $c\,v_p(a)\le c\,v_p(b)$ is equivalent to
	$v_p(a)\le v_p(b)$. Applying Theorem~\ref{thm:4.5.2} to the pair $(a,b)$
	and to the pair $(a^{c},b^{c})$ now gives the claim.

\end{proof}
\begin{theorem}\label{thm:4.6.4} If $a,b\in\mathbb{Z}^+$ satisfy $\gcd(a, b) = 1$ and $c\in\mathbb{Z}$, then
	$\gcd(ab, c) = \gcd(a, c)\,\gcd(b, c)$, and the two factors on the
	right are coprime.

\end{theorem}

\begin{proof} Assume first that $c>0$, and fix a prime $p$. Since $\gcd(a, b) = 1$, at
	least one of $v_p(a), v_p(b)$ is $0$ (Theorem~\ref{thm:4.6.1}), so
	$v_p(ab) = v_p(a) + v_p(b) = \max\{v_p(a), v_p(b)\}$, and by Theorem~\ref{thm:4.6.1},
	\[
		v_p(\gcd(ab, c)) = \min\{\max\{v_p(a), v_p(b)\},\, v_p(c)\}.
	\]
	If $v_p(a) = 0$, this equals $\min\{v_p(b), v_p(c)\} = v_p(\gcd(b, c))$,
	while $v_p(\gcd(a, c)) = \min\{0, v_p(c)\} = 0$; the case $v_p(b) = 0$
	is symmetric. In either case
	$v_p(\gcd(ab, c)) = v_p(\gcd(a, c)) + v_p(\gcd(b, c))$ and
	$\min\{v_p(\gcd(a, c)), v_p(\gcd(b, c))\} = 0$. As these hold for every
	prime $p$, Theorem~\ref{thm:4.5.3} gives $\gcd(ab, c) = \gcd(a, c)\,\gcd(b, c)$
	and $\gcd(\gcd(a, c), \gcd(b, c)) = 1$. If $c<0$, replacing $c$ by
	$|c|$ leaves all three gcds unchanged, so the conclusion follows from
	the positive case. If $c=0$, then
	\[
		\gcd(ab, 0) = ab = \gcd(a, 0)\,\gcd(b, 0).
	\]
	Here the two factors on the right are $a$ and $b$, which are coprime by hypothesis.

\end{proof}
\begin{altproof}[Alternative proof of Theorem~\ref{thm:2.2.11}]
	By Theorems~\ref{thm:4.6.1} and~\ref{thm:4.5.1}(b), together with
	$\min\{x,y\}+\max\{x,y\}=x+y$,
	we have $v_p\big(\gcd(a,b)\lcm(a,b)\big)=v_p(a)+v_p(b)=v_p(ab)$ for every prime
	$p$; Theorem~\ref{thm:4.5.3} then gives $\gcd(a,b)\lcm(a,b)=ab$.

\end{altproof}
\begin{altproof}[Alternative proof of Theorem~\ref{thm:2.2.12}]
	By Theorem~\ref{thm:4.6.1}, $\min\{v_p(a)+v_p(b),\,v_p(a)+v_p(c)\}
		=v_p(a)+\min\{v_p(b),v_p(c)\}$, and likewise with $\max$ in place of
	$\min$; Theorems~\ref{thm:4.5.1}(b) and~\ref{thm:4.5.3} then give both
	identities.

\end{altproof}
\begin{altproof}[Alternative proof of Theorem~\ref{thm:2.2.13}]
	Write $d=\gcd(a,b)$, $a=da_1$, $b=db_1$. Theorem~\ref{thm:4.6.1} gives
	\[
		\min\{v_p(a_1),v_p(b_1)\}=\min\{v_p(a)-v_p(d),\,v_p(b)-v_p(d)\}=0
	\]
	for every prime $p$, whence $\gcd(a_1,b_1)=1$ by Theorem~\ref{thm:4.5.3}.

\end{altproof}
\subsection{Legendre's Formula and Kummer's Theorem}

\begin{theorem}[Legendre's formula]\label{thm:4.7.1} For a prime
	$p$ and a nonnegative integer $n$, we have $v_p(n!) = \sum_{i=1}^{\infty}
		\left\lfloor \frac{n}{p^{i}} \right\rfloor$.

\end{theorem}

\begin{proof} If $n=0$, then
	$v_p(0!)=v_p(1)=0$, while every term
	$\left\lfloor 0/p^i\right\rfloor$ on the right is $0$. Assume
	$n\ge1$. Repeated application of Theorem~\ref{thm:4.5.1}(b), by
	induction on the number of factors, gives
	\[
		v_p(n!)=\sum_{j=1}^{n}v_p(j).
	\]
	Consider the set
	\[
		T=\{(i,j)\in\mathbb{Z}^+\times\{1,2,\ldots,n\} \mid p^i\mid j\}.
	\]
	It is finite because $(i,j)\in T$ implies $p^i\le j\le n$. For a fixed
	$j$, Definition~\ref{def:4.5.1} shows that the number of positive integers $i$
	with $(i,j)\in T$ is $v_p(j)$. Thus, counting $T$ by its second
	coordinate gives $|T|=\sum_{j=1}^{n}v_p(j)$. For a fixed $i$, the
	possible second coordinates are exactly
	$p^i,2p^i,\ldots,\lfloor n/p^i\rfloor p^i$, so counting by the first
	coordinate gives
	$|T|=\sum_{i=1}^{\infty}\left\lfloor n/p^i\right\rfloor$; this sum is
	finite because its terms vanish when $p^i>n$. Equating the two counts
	proves the formula.

\end{proof}
\begin{altproof}[Alternative proof of Theorem~\ref{thm:4.7.1}]
	Put $m=\lfloor n/p\rfloor$. The multiples of $p$ in
	$\{1,2,\ldots,n\}$ are $p,2p,\ldots,mp$, while every other integer in
	this set has $p$-adic valuation $0$. Hence Theorem~\ref{thm:4.5.1}(b) gives
	\[
		\begin{aligned}
			v_p(n!)
			 & =\sum_{j=1}^{m}v_p(pj)
			=\sum_{j=1}^{m}(1+v_p(j))                \\
			 & =\left\lfloor\frac{n}{p}\right\rfloor
			+v_p\!\left(\left\lfloor\frac{n}{p}\right\rfloor!\right).
		\end{aligned}
	\]
	This recursion proves the formula by strong induction on $n$. The case
	$n=0$ is immediate. If $n\ge1$, then $m=\lfloor n/p\rfloor<n$, so the
	induction hypothesis and the identity
	$\left\lfloor\lfloor n/p\rfloor/p^i\right\rfloor
		=\left\lfloor n/p^{i+1}\right\rfloor$ give
	\[
		\begin{aligned}
			v_p(n!)
			 & =\left\lfloor\frac{n}{p}\right\rfloor
			+\sum_{i=1}^{\infty}
				\left\lfloor\frac{\lfloor n/p\rfloor}{p^i}\right\rfloor \\
			 & =\left\lfloor\frac{n}{p}\right\rfloor
			+\sum_{i=1}^{\infty}\left\lfloor\frac{n}{p^{i+1}}\right\rfloor
			=\sum_{i=1}^{\infty}\left\lfloor\frac{n}{p^i}\right\rfloor.
		\end{aligned}
	\]

\end{altproof}

\begin{theorem}[Legendre's digit-sum formula]\label{thm:4.7.2} For
	all nonnegative integers $n$ and primes $p$, we have $v_p(n!) = \frac{n -
			S_p(n)}{p - 1}$.

\end{theorem}

\begin{proof} By Theorem~\ref{thm:4.1.1}, write
	$n=\sum_{i=0}^{m}d_ip^i$ with $0\le d_i\le p-1$, allowing $m=0$ and
	$d_0=0$ when $n=0$. Then $n<p^{m+1}$. By Theorem~\ref{thm:4.7.1},
	$v_p(n!)=\sum_{j\ge1}\left\lfloor n/p^j\right\rfloor$; for $j>m$
	we have $n<p^{m+1}\le p^j$, so these terms vanish. For
	$1\le j\le m$,
	\[
		\frac{n}{p^j}=\sum_{i=j}^{m}d_ip^{i-j}
		+\frac{1}{p^j}\sum_{i=0}^{j-1}d_ip^i,
	\]
	and
	$0\le\sum_{i=0}^{j-1}d_ip^i\le(p-1)\sum_{i=0}^{j-1}p^i=p^j-1<p^j$.
	Thus
	$\left\lfloor n/p^j\right\rfloor=\sum_{i=j}^{m}d_ip^{i-j}$, and hence
	\[
		\begin{aligned}
			v_p(n!)
			 & =\sum_{j=1}^{m}\sum_{i=j}^{m}d_ip^{i-j}
			=\sum_{i=1}^{m}d_i\sum_{j=1}^{i}p^{i-j}    \\
			 & =\sum_{i=1}^{m}d_i\frac{p^i-1}{p-1}
			=\frac{n-S_p(n)}{p-1}.
		\end{aligned}
	\]
	For the last equality, the $i=0$ term $d_0(p^0-1)$ is $0$, so the
	numerator is
	$\sum_{i=0}^{m}d_i(p^i-1)=n-S_p(n)$. Empty sums cover the case $m=0$.

\end{proof}
\begin{theorem}[Kummer's theorem]\label{thm:4.7.3} Let
	$n,k\in\mathbb{Z}_{\ge0}$ with $k\le n$, and let $p$ be a prime. Then
	\[
		v_p\!\left(\binom{n}{k}\right)
		=\frac{S_p(k)+S_p(n-k)-S_p(n)}{p-1},
	\]
	and this quantity equals the number of carries that occur when $k$
	and $n-k$ are added in base $p$.

\end{theorem}

\begin{proof} The assumption $k\le n$ gives $n-k\ge0$.
	By the factorial formula for $\binom{n}{k}$, Theorem~\ref{thm:4.5.5}(a), and
	Theorem~\ref{thm:4.7.2},
	\[
		\begin{aligned}
			v_p\!\left(\binom{n}{k}\right)
			 & =v_p(n!)-v_p(k!)-v_p((n-k)!)               \\
			 & =\frac{n-S_p(n)}{p-1}-\frac{k-S_p(k)}{p-1}
			-\frac{(n-k)-S_p(n-k)}{p-1}                   \\
			 & =\frac{S_p(k)+S_p(n-k)-S_p(n)}{p-1}.
		\end{aligned}
	\]
	It remains to identify the last quantity with the number of carries.
	By Theorem~\ref{thm:4.1.1}, write
	\[
		k=\sum_{i\ge0}a_i p^i,\qquad
		n-k=\sum_{i\ge0}b_i p^i,\qquad
		n=\sum_{i\ge0}d_i p^i,
	\]
	where $0\le a_i,b_i,d_i\le p-1$ and all but finitely many digits are
	zero. Define the carry digits $c_i$ and output digits $e_i$ recursively
	by setting $c_0=0$ and, for each $i\ge0$, using the division algorithm to write uniquely
	\[
		a_i+b_i+c_i=e_i+p c_{i+1},\qquad 0\le e_i\le p-1.
	\]
	Every $c_i$ belongs to $\{0,1\}$. Indeed, this holds for $c_0$; if
	$c_i\in\{0,1\}$, then
	$0\le a_i+b_i+c_i\le(p-1)+(p-1)+1=2p-1<2p$, so
	$c_{i+1}=\left\lfloor(a_i+b_i+c_i)/p\right\rfloor\in\{0,1\}$.
	Choose $M$ such that $a_i=b_i=0$ for every $i\ge M$. Then
	$c_{M+1}=\lfloor c_M/p\rfloor=0$, and the same recurrence shows that
	$c_i=e_i=0$ for all sufficiently large $i$. Multiplying the recurrence
	by $p^i$ and summing over $i\ge0$ gives
	\[
		\sum_{i\ge0}e_ip^i
		=k+(n-k)+\sum_{i\ge0}c_ip^i-
		\sum_{i\ge0}c_{i+1}p^{i+1}=n,
	\]
	because $c_0=0$ makes the last two sums equal. Thus the $e_i$ form a
	base-$p$ representation of $n$, and uniqueness in Theorem~\ref{thm:4.1.1} gives
	$e_i=d_i$ for every $i$. Summing the unweighted recurrences now gives
	\[
		S_p(k)+S_p(n-k)+\sum_{i\ge1}c_i
		=S_p(n)+p\sum_{i\ge1}c_i;
	\]
	on the left, we used $c_0=0$, while on the right, we reindexed
	$\sum_{i\ge0}c_{i+1}=\sum_{i\ge1}c_i$. Rearranging gives
	\[
		S_p(k)+S_p(n-k)-S_p(n)
		=(p-1)\sum_{i\ge1}c_i.
	\]
	There is a carry into position $i$ exactly when $c_i=1$. Since each
	$c_i\in\{0,1\}$, the number of carries is
	$\#\{i\ge1 \mid c_i=1\}=\sum_{i\ge1}c_i$. Substitution into the first
	formula proves the result.

	This also gives a valuation-based proof of integrality without assuming it in
	advance: the calculation applies to the positive rational number
	$Q=n!/(k!(n-k)!)$ and shows that $v_p(Q)\ge0$ for every prime $p$.
	Write $Q=A/B$ in lowest terms with $A,B\in\mathbb{Z}^+$.
	If $B>1$, a prime divisor $p$ of $B$ would satisfy $p\nmid A$ and
	$v_p(Q)=-v_p(B)<0$, a contradiction. Hence $B=1$, so $Q$ is an integer.

\end{proof}
\begin{theorem}\label{thm:4.7.4} Let $p$ be a prime. Then:

	\begin{enumerate}[label={\alph*)}]
		\item $n\equiv S_p(n)\pmod{p-1}$ for every
		      $n\in\mathbb{Z}_{\ge0}$.
		\item $v_p(n!)<\dfrac{n}{p-1}$ for every $n\in\mathbb{Z}^+$.
		\item If $n,k\in\mathbb{Z}_{\ge0}$ and $k\le n$, then
		      $p\mid\binom{n}{k}$ if and only if at least one carry occurs when $k$
		      and $n-k$ are added in base $p$; equivalently,
		      $p\nmid\binom{n}{k}$ if and only if no carry occurs.
	\end{enumerate}

\end{theorem}

\begin{proof}

	\begin{enumerate}[label={\alph*)}]
		\item Write $n=\sum_{i\ge0}d_ip^i$ in base $p$. Since
		      $p\equiv1\pmod{p-1}$,
		      we have $n\equiv\sum_{i\ge0}d_i=S_p(n)\pmod{p-1}$.
		\item If $n\ge1$, then $S_p(n)\ge1$. Theorem~\ref{thm:4.7.2} therefore gives
		      \[
			      v_p(n!)=\frac{n-S_p(n)}{p-1}
			      \le\frac{n-1}{p-1}<\frac{n}{p-1}.
		      \]
		\item By Definition~\ref{def:4.5.1},
		      $p\mid\binom{n}{k}$ if and only if
		      $v_p\!\left(\binom{n}{k}\right)\ge1$. Theorem~\ref{thm:4.7.3} identifies this
		      valuation with the number of carries, proving both assertions.
	\end{enumerate}

\end{proof}
\subsection{The Lifting-the-Exponent Lemma}

\begin{theorem}[lifting the exponent (LTE), coprime case]\label{thm:4.8.1} Let $p$ be a
	prime, and let $a, b \in \mathbb{Z}$ with $\gcd(a, p) = 1$, $a \neq b$, and $p \mid (a - b)$.
	If $n \in \mathbb{Z}^+$ and $\gcd(n, p) = 1$, then $v_p(a^{n} - b^{n}) = v_p(a - b)$.

\end{theorem}

\begin{proof} Let $v_p(a - b) = \lambda$; since $p \mid (a
		- b)$, we have $\lambda \ge 1$. Write $a - b =
		k p^{\lambda}$ with $\gcd(k, p)=1$, so $a = b + k p^{\lambda}$. The
	conditions $\gcd(a,p)=1$ and $p\mid(a-b)$ give $p\nmid b$: otherwise
	$p\mid(a-b)+b=a$. Expanding gives
	$a^{n} - b^{n} = \sum_{i=1}^{n} \binom{n}{i} k^{i} p^{i\lambda}
		b^{n-i} = p^{\lambda} s$, where $s = \sum_{i=1}^{n} \binom{n}{i} k^{i}
		p^{\lambda(i-1)} b^{n-i} = n k b^{n-1} + s'$ with $s' =
		\sum_{i=2}^{n} \binom{n}{i} k^{i} p^{\lambda(i-1)} b^{n-i}$. Now
	$\gcd(n, p) = \gcd(k, p) = 1$ and $p \nmid b$, so $p \nmid n k b^{n-1}$;
	and every term of $s'$ contains the factor $p^{\lambda(i-1)}$ with $i
		\ge 2$ and $\lambda \ge 1$, so $p \mid s'$. Hence $p \nmid s$; in
	particular $s \neq 0$, so that $a^{n} - b^{n} = p^{\lambda} s \neq 0$ and $v_p(s) = 0$; hence $v_p(a^{n} - b^{n}) =
		\lambda = v_p(a - b)$.

\end{proof}
\begin{theorem}[lifting the exponent, prime case]\label{thm:4.8.2} Let $p$ be an odd
	prime, and let $a, b \in \mathbb{Z}$ with $\gcd(a, p) = 1$, $a \neq b$, and $p \mid (a - b)$.
	Then $v_p(a^{p} - b^{p}) = v_p(a - b) + 1$.

\end{theorem}

\begin{proof} Put $\lambda=v_p(a-b)\ge1$ and write
	$a-b=kp^\lambda$ with $p\nmid k$. Since $p\nmid a$ and $p\mid(a-b)$,
	we also have $p\nmid b$. The binomial expansion gives
	\[
		a^p-b^p=p b^{p-1}(kp^\lambda)
		+\sum_{i=2}^{p-1}\binom pi b^{p-i}(kp^\lambda)^i
		+(kp^\lambda)^p.
	\]
	The first term has valuation exactly $\lambda+1$.
	For $2\le i\le p-1$, Theorem~\ref{thm:3.5.2} gives
	$p\mid\binom pi$, so each term in the sum is divisible by
	$p^{1+2\lambda}$, and $1+2\lambda\ge\lambda+2$.
	The last term is divisible by $p^{p\lambda}$, with
	$p\lambda\ge\lambda+2$ because $p\ge3$ and $\lambda\ge1$.
	Consequently, for some $C\in\mathbb{Z}$,
	\[
		a^p-b^p=p^{\lambda+1}\bigl(kb^{p-1}+pC\bigr).
	\]
	The factor in parentheses is not divisible by $p$, so it is nonzero
	and has valuation $0$. Hence $v_p(a^p-b^p)=\lambda+1$.

\end{proof}
\begin{theorem}[lifting the exponent, odd prime, general $n$]\label{thm:4.8.3} Let $p$ be an odd prime, and let $a, b \in \mathbb{Z}$ with $\gcd(a, p)=1$, $a
		\neq b$, and $p \mid (a-b)$. Then $v_p(a^{n} - b^{n}) = v_p(a-b) +
		v_p(n)$ for all $n \in \mathbb{Z}^+$.

\end{theorem}

\begin{proof} Write $n=p^{t}m$ with $p\nmid m$, so that
	$t=v_p(n)$, and put $A_j=a^{mp^{j}}$, $B_j=b^{mp^{j}}$ for $0\le j\le
		t$; thus $A_0=a^{m}$, $B_0=b^{m}$, while $A_t=a^{n}$ and $B_t=b^{n}$.
	Since $\gcd(a,p)=1$, we have $p\nmid a$. The relation $p\mid(a-b)$ then
	gives $p\nmid b$. Hence neither $A_j$ nor $B_j$ is divisible by $p$. Moreover,
	$A_j\neq B_j$ for every $j$: otherwise $|a|^{mp^{j}}=|b|^{mp^{j}}$
	would force $|a|=|b|$, that is, $b=\pm a$; the case $b=a$ is excluded
	by hypothesis, while $b=-a$ would give $p\mid a-b=2a$ and hence $p\mid a$ since $p$
	is odd, contrary to $p\nmid a$. We prove by induction on
	$j$ that
	\[v_p(A_j-B_j)=v_p(a-b)+j\qquad(0\le j\le t).\]
	For $j=0$ this is Theorem~\ref{thm:4.8.1} applied with the exponent $m$, which
	is coprime to $p$. Assume the identity for some $j<t$. Then
	$v_p(A_j-B_j)=v_p(a-b)+j\ge1$, so $p\mid A_j-B_j$, and $A_j,B_j$
	satisfy every hypothesis of Theorem~\ref{thm:4.8.2}; since $A_{j+1}=A_j^{\,p}$
	and $B_{j+1}=B_j^{\,p}$, that theorem gives
	$v_p(A_{j+1}-B_{j+1})=v_p(A_j-B_j)+1=v_p(a-b)+j+1$. Taking $j=t$ and
	recalling that $A_t=a^{n}$ and $B_t=b^{n}$, we obtain
	\[v_p(a^{n}-b^{n})=v_p(a-b)+t=v_p(a-b)+v_p(n);\]
	for $t=0$, no induction step is needed and the assertion is Theorem~\ref{thm:4.8.1}
	itself.

\end{proof}
\begin{theorem}[lifting the exponent, sum version]\label{thm:4.8.4} Let $p$ be an
	odd prime, and let $a, b \in \mathbb{Z}$ with $\gcd(a, p)=1$, $a \neq -b$, and $p \mid (a+b)$. Then
	$v_p(a^{n} + b^{n}) = v_p(a+b) + v_p(n)$ for all positive odd $n$,
	whereas $v_p(a^{n} + b^{n}) = 0$ for all positive even $n$.

\end{theorem}

\begin{proof} Put $c = a$ and $d = -b$; then $\gcd(c, p) =
		1$, $c \neq d$, and $p \mid (c-d)$. If $n$ is odd,
	Theorem~\ref{thm:4.8.3} gives
	$v_p(a^{n} + b^{n}) = v_p(c^{n} - d^{n}) = v_p(c - d) + v_p(n) =
		v_p(a+b) + v_p(n)$. If instead $n$ is even, then $b \equiv -a
		\pmod{p}$ gives $a^{n} + b^{n} \equiv a^{n}\big(1 + (-1)^{n}\big) =
		2 a^{n} \pmod{p}$; as $p$ is odd and $p \nmid a$, the sum is not
	divisible by $p$, so $v_p(a^{n} + b^{n}) = 0$.

\end{proof}
\begin{theorem}[lifting the exponent, $p=2$, $n=2$]\label{thm:4.8.5}
	Let $a,b$ be odd integers with $a\neq\pm b$. Then
	$v_2(a^2-b^2)=v_2(a-b)+v_2(a+b)$; when $4\mid(a-b)$, this is
	$v_2(a-b)+1$; otherwise it is $v_2(a+b)+1$.

\end{theorem}

\begin{proof} Since $a\neq\pm b$, both $a-b$ and $a+b$
	are nonzero. From the identity $a^{2}-b^{2}=(a-b)(a+b)$ and the
	additivity of $v_2$ (Theorem~\ref{thm:4.5.1}(b)),
	\[
		v_2(a^2-b^2)=v_2(a-b)+v_2(a+b).
	\]
	Since $a$ and $b$ are odd, $a-b$ and $a+b$ are both even, so
	$v_2(a-b)\ge1$ and $v_2(a+b)\ge1$; and the two valuations cannot both be at least
	$2$, since $(a-b)+(a+b)=2a$ and $v_2(2a)=1$ as $a$ is odd. Hence
	$\min\{v_2(a-b),v_2(a+b)\}=1$. If $4\mid(a-b)$, then
	$v_2(a-b)\ge2$, so $v_2(a+b)=1$ and
	\[
		v_2(a^2-b^2)=v_2(a-b)+1;
	\]
	otherwise $v_2(a-b)=1$ and
	\[
		v_2(a^2-b^2)=v_2(a+b)+1.
	\]

\end{proof}
\begin{theorem}[lifting the exponent, $p=2$]\label{thm:4.8.6} Let $a, b$
	be odd integers with $a \neq \pm b$, and let $n \in \mathbb{Z}^+$. Then
	\[
		v_2(a^{n} - b^{n}) =
		\begin{cases} v_2(a - b),           & \text{if } 2 \nmid n, \\ v_2(a-b) +
		              v_2(a+b) + v_2(n) - 1, & \text{if } 2 \mid n.
		\end{cases}
	\]

\end{theorem}

\begin{proof} The hypothesis $a\neq\pm b$ guarantees
	$a^{n}-b^{n}\neq0$ for every $n\ge1$. For odd $n$ this holds because $a\neq b$
	and $x\mapsto x^{n}$ is injective on $\mathbb{Z}$; for even $n$,
	the equality $a^{n}=b^{n}$ would force $|a|=|b|$, that is, $a=\pm b$.
	Consequently every factor of $a^{n}-b^{n}$ occurring below is nonzero,
	so all the $2$-adic valuations are finite. Write $n=2^{s}m$
	with $m$ odd, so that $s=v_2(n)$.

	If $s=0$, then $n=m$ is odd, so $\gcd(n,2)=1$ and Theorem~\ref{thm:4.8.1},
	applied with $p=2$, gives $v_2(a^{n}-b^{n})=v_2(a-b)$.

	Now suppose $s\ge1$. Repeatedly applying the identity $x^{2}-y^{2}=(x-y)(x+y)$ with
	$x=a^{2^{j}m}$, $y=b^{2^{j}m}$ for $j=s-1,s-2,\ldots,0$ yields
	\[
		a^{n}-b^{n}=(a^{m}-b^{m})\prod_{j=0}^{s-1}\big(a^{2^{j}m}+b^{2^{j}m}\big),
	\]
	so by additivity of $v_2$ (Theorem~\ref{thm:4.5.1}(b)) it suffices to evaluate
	the valuation of each factor. First, $m$ is odd, so Theorem~\ref{thm:4.8.1} with $p=2$ gives
	$v_2(a^{m}-b^{m})=v_2(a-b)$. Next,
	\[
		a^m+b^m=(a+b)\sum_{i=0}^{m-1}(-1)^i a^{m-1-i}b^i.
	\]
	The second factor is a sum of $m$ odd terms and is therefore odd, since $m$
	is odd. This also covers $m=1$, when the sum is $1$.
	Thus $v_2(a^m+b^m)=v_2(a+b)$. Finally, for
	$1\le j\le s-1$ the integers $a^{2^{j}m}$ and $b^{2^{j}m}$ are squares
	of odd integers, hence each is congruent to $1\pmod{4}$, so
	$a^{2^{j}m}+b^{2^{j}m}\equiv2\pmod{4}$ and
	$v_2\big(a^{2^{j}m}+b^{2^{j}m}\big)=1$; there are $s-1$ such factors,
	where this range is empty when $s=1$. Adding the contributions gives
	\[
		v_2(a^{n}-b^{n})=v_2(a-b)+v_2(a+b)+(s-1)=v_2(a-b)+v_2(a+b)+v_2(n)-1.
	\]

\end{proof}
\begin{theorem}\label{thm:4.8.7}
	Let $p$ be an odd prime and let $a,b\in\mathbb{Z}$ satisfy
	$p\nmid ab$, $a\neq b$, and $p\mid(a-b)$. Put
	$\lambda=v_p(a-b)$. For $n,r\in\mathbb{Z}^+$ with
	$r\ge\lambda$,
	\[
		p^r\mid a^n-b^n \quad\Longleftrightarrow\quad
		p^{r-\lambda}\mid n.
	\]

\end{theorem}

\begin{proof} Since $p$ is prime and $p\nmid ab$, we have $p\nmid a$ and
	$p\nmid b$, hence $\gcd(a,p)=\gcd(b,p)=1$. By the lifting-the-exponent lemma
	(Theorem~\ref{thm:4.8.3}),
	$v_p(a^n-b^n)=\lambda+v_p(n)$. Thus $p^r\mid a^n-b^n$ if and only
	if $\lambda+v_p(n)\ge r$, equivalently
	$v_p(n)\ge r-\lambda$, which is exactly
	$p^{r-\lambda}\mid n$.

\end{proof}
\subsection{Applications of Valuations to Elementary Exponential Diophantine Equations}

\begin{definition}[exponential Diophantine equation]\label{def:4.9.1}
	An exponential Diophantine equation is an equation in which one or more unknowns occur as exponents, and whose integer solutions are sought.

\end{definition}
\begin{theorem}\label{thm:4.9.1}
	The only integer solutions of $2^x-y^2=1$ with $x\ge0$ are $(x,y)=(0,0)$
	and $(1,\pm1)$.

\end{theorem}

\begin{proof} If $x\ge2$ then $2^{x}\equiv0\pmod{4}$, so
	$y^{2}\equiv-1\equiv3\pmod{4}$, which is impossible since squares are
	congruent to $0$ or $1\pmod{4}$. Checking $x=0,1$ directly gives exactly the stated
	solutions.

\end{proof}
\begin{theorem}\label{thm:4.9.2}
	The only solutions of $3^x-2^y=1$ in positive integers are $(x,y)=(1,1)$
	and $(x,y)=(2,3)$.

\end{theorem}

\begin{proof} If $y=1$, then $3^{x}=3$, so $x=1$. If
	$y\ge2$, reducing modulo $4$ gives $3^{x}\equiv1\pmod{4}$ (since
	$2^{y}\equiv0\pmod{4}$), and as $3^{x}\equiv(-1)^{x}\pmod{4}$, the exponent $x$ is even, say
	$x=2k$. Then $2^{y}=3^{2k}-1=(3^{k}-1)(3^{k}+1)$; the two factors are
	positive integers differing by $2$ whose product is a power of $2$,
	hence both are powers of $2$. If the smaller were $2^r$ with
	$r\ge2$, both powers would be divisible by $4$, so their difference
	would also be divisible by $4$, whereas it equals $2$. Thus the smaller factor
	must be $2$, so
	$3^{k}-1=2$, which gives $k=1$ and $x=2$; then $3^{k}+1=4=2^{2}$ gives
	$y=3$.

\end{proof}
\begin{altproof}[Alternative proof of Theorem~\ref{thm:4.9.2}] Write the equation as $2^{y}+1=3^{x}$. Reduction modulo $3$
	shows that $y$ is odd. The sum version of the lifting-the-exponent
	lemma (Theorem~\ref{thm:4.8.4}), with $p=3$, $a=2$, $b=1$, and
	$n=y$, therefore gives
	\[
		x=v_3(2^y+1)=v_3(2+1)+v_3(y)=1+v_3(y).
	\]
	If $y=1$, the equation gives $x=1$. Now suppose $y>1$, so $y\ge3$.
	Modulo $8$, the left-hand side is $1$; hence $3^{x}\equiv1\pmod{8}$ and
	$x$ is even. Therefore $v_3(y)=x-1$ is a positive odd integer, so
	$3\mid y$. Write $y=3u$; since $y$ is odd, so is $u$. Factoring the
	sum of cubes gives
	\[
		3^x=2^{3u}+1=(2^u+1)(2^{2u}-2^u+1).
	\]
	Both positive factors are powers of $3$. Moreover,
	\[
		\gcd(2^u+1,2^{2u}-2^u+1)=\gcd(2^u+1,3)=3.
	\]
	Indeed, the first equality follows by reducing the second argument modulo
	$2^u+1$, and the second holds because $u$ is odd, so
	$2^u\equiv-1\pmod{3}$. Thus at least one of the two factors equals
	$3$. If $2^u+1=3$, then $2^u=2$; if
	$2^{2u}-2^u+1=3$, putting $t=2^u$ gives
	$(t-2)(t+1)=0$, again forcing $2^u=2$. Hence $u=1$, $y=3$, and the
	original equation gives $x=2$. The pairs $(1,1)$ and $(2,3)$ indeed
	satisfy the equation.

\end{altproof}
\begin{theorem}\label{thm:4.9.3}
	The only solution of $3^x+1=2^y$ in positive integers is $(x,y)=(1,2)$.

\end{theorem}

\begin{proof} If $x$ is even, then
	$3^x+1\equiv2\pmod{8}$, so the equality would force $y=1$, which
	is impossible because $3^x+1>2$. Hence $x$ is odd. Apply Theorem~\ref{thm:4.8.1} with $p=2$, $a=3$, $b=-1$, and $n=x$. Since
	$3^x-(-1)^x=3^x+1$, it gives
	\[
		y=v_2(3^x+1)=v_2(3-(-1))=v_2(4)=2.
	\]
	Thus $3^x=3$, so $x=1$; direct substitution verifies the solution.

\end{proof}
\newpage

\section{Number-Theoretic Functions}

\subsection{Divisor Functions and Basic Divisor-Sum Identities}

\begin{definition}[arithmetic function and multiplicative function]\label{def:5.1.1} An
	arithmetic function is a function $f:\mathbb{Z}^+\to\mathbb{C}$. An arithmetic
	function $f$ that is not identically zero is called multiplicative
	if $f(mn)=f(m)f(n)$ for all $m,n\in\mathbb{Z}^+$
	with $\gcd(m, n)=1$.

\end{definition}
\begin{theorem}\label{thm:5.1.1}
	Every multiplicative function $f$ satisfies $f(1)=1$.

\end{theorem}

\begin{proof} Since $f$ is not identically zero, choose $m$ with
	$f(m)\neq0$. Then $f(m)=f(m\cdot1)=f(m)f(1)$, so cancellation gives
	$f(1)=1$.

\end{proof}
\begin{definition}[$\tau$ and $\sigma$]\label{def:5.1.2} For $n \in \mathbb{Z}^+$, put $\tau(n) =
		\sum_{d \mid n} 1$ and $\sigma(n) = \sum_{d \mid n} d$; here and throughout, the index of a divisor sum $\sum_{d \mid n}$ ranges over the positive divisors of $n$ only.

\end{definition}
\begin{theorem}\label{thm:5.1.2}\statementitems

	\begin{enumerate}[label={\alph*)}]
		\item $\tau(n) = 2$ if and only if $n$ is a prime; $\sigma(n) = n
			      + 1$ if and only if $n$ is a prime.
		\item The arithmetic functions $\tau(n)$ and $\sigma(n)$ are multiplicative.
		\item If $n=\prod_{i=1}^{k}p_i^{a_i}$, where $p_1,p_2,\ldots,p_k$
		      are distinct primes and $a_1,a_2,\ldots,a_k\in\mathbb{Z}^+$, then
		      $\tau(n)=\prod_{i=1}^{k}(a_i+1)$ and
		      $\sigma(n)=\prod_{i=1}^{k}\dfrac{p_i^{a_i+1}-1}{p_i-1}$.
	\end{enumerate}

\end{theorem}

\begin{proof}

	\begin{enumerate}[label={\alph*)}]
		\item If $n$ is prime, then $1$ and $n$ are its only positive divisors,
		      so $\tau(n)=2$ and $\sigma(n)=1+n$. Conversely, if $\tau(n)=2$, then
		      $n\neq1$ and $n$ has no positive divisors other than $1$ and itself,
		      so $n$ is prime; if
		      $\sigma(n)=n+1$, then $n\neq1$, and a divisor $d \neq 1, n$
		      would give $\sigma(n)\ge 1+d+n > n+1$, so $n$ is prime.
		\item For $\gcd(m, n)=1$, every divisor $d\mid mn$ factors uniquely
		      as $d=d_1d_2$ with $d_1\mid m$, $d_2\mid n$. By the fundamental
		      theorem of arithmetic (Theorem~\ref{thm:2.3.3}), each prime divides at
		      most one of $m, n$, so the full $p$-part of $d$ must divide $d_1$ if $p\mid m$
		      and $d_2$ if $p\mid n$, which gives existence and uniqueness of $(d_1,d_2)$. Hence
		      $\tau(mn)=\tau(m)\tau(n)$ and $\sigma(mn)=\big(\sum_{d_1\mid
				      m}d_1\big)\big(\sum_{d_2\mid n}d_2\big)=\sigma(m)\sigma(n)$.
		\item Every positive divisor of
		      $n=\prod_{i=1}^{k}p_i^{a_i}$ has a unique representation
		      $\prod_{i=1}^{k}p_i^{b_i}$ with $0\le b_i\le a_i$ (again by Theorem~\ref{thm:2.3.3}). Therefore
		      \[
			      \tau(n)=\prod_{i=1}^{k}(a_i+1)
		      \]
		      and
		      \[
			      \sigma(n)=\prod_{i=1}^{k}(1+p_i+\cdots+p_i^{a_i})
			      =\prod_{i=1}^{k}\frac{p_i^{a_i+1}-1}{p_i-1}.
		      \]
	\end{enumerate}

\end{proof}
\begin{theorem}\label{thm:5.1.3} If $f$ is a multiplicative
	function, then $F(n) = \sum_{d \mid n} f(d)$ is also multiplicative.

\end{theorem}

\begin{proof} Let $m,n\in\mathbb{Z}^+$ with $\gcd(m, n)=1$. By unique prime factorization, every $d\mid mn$ has a unique factorization $d=d_1d_2$, where $d_1$ contains the prime-power factors belonging to $m$ and $d_2$ those belonging to $n$; thus $d_1\mid m$, $d_2\mid n$,
	and this correspondence is a bijection. Hence $F(mn) = \sum_{d_1 \mid
			m} \sum_{d_2 \mid n} f(d_1 d_2) = \big(\sum_{d_1 \mid m}
		f(d_1)\big)\big(\sum_{d_2 \mid n} f(d_2)\big) = F(m) F(n)$, using
	$f(d_1 d_2) = f(d_1) f(d_2)$ since $\gcd(d_1, d_2) = 1$. Moreover, $F(1) = f(1) = 1$, so $F$ is not identically zero and is multiplicative.

\end{proof}
\begin{theorem}\label{thm:5.1.4} Let $n\in\mathbb{Z}^+$ and let
	$d\in\mathbb{Z}^+$ with $d\mid n$. Among the fractions
	$1/n,2/n,\ldots,n/n$, exactly $\varphi(d)$ reduce to a fraction with
	denominator $d$.

\end{theorem}

\begin{proof} For $1\le k\le n$, put $g=\gcd(k, n)$.
	Dividing numerator and denominator by $g$ gives
	\[
		\frac{k}{n}=\frac{k/g}{n/g},\qquad
		\gcd(k/g, n/g)=1,
	\]
	so the reduced denominator is $n/g$. It equals $d$ if and only if
	$g=n/d$. In that case $k=(n/d)r$ for a unique integer $r$ with
	$1\le r\le d$, and
	\[
		\gcd(k, n)=\frac{n}{d}
		\iff \gcd\!\left(\frac{n}{d}r,\frac{n}{d}d\right)=\frac{n}{d}
		\iff \gcd(r, d)=1.
	\]
	Thus $k\mapsto r=k/(n/d)$ induces a bijection from the fractions whose
	reduced denominator is $d$ to
	$\{r\in\mathbb{Z} \mid 1\le r\le d,\ \gcd(r, d)=1\}$, a set of cardinality
	$\varphi(d)$.

\end{proof}
\begin{theorem}\label{thm:5.1.5} For every $n\in\mathbb{Z}^+$,
	\[
		\sum_{d\mid n}\varphi(d)=n.
	\]

\end{theorem}

\begin{proof} The $n$ fractions
	$1/n,2/n,\ldots,n/n$ are partitioned by reduced
	denominator. Every reduced denominator divides $n$, and by Theorem~\ref{thm:5.1.4}
	the class with reduced denominator $d$ has exactly
	$\varphi(d)$ members. Summing the cardinalities of these disjoint
	classes over all $d\mid n$ gives
	$n=\sum_{d\mid n}\varphi(d)$.

\end{proof}
\begin{altproof}[Alternative proof of Theorem~\ref{thm:5.1.5}]
	Since $\varphi$ is multiplicative, Theorem~\ref{thm:5.1.3} shows that
	$F(n)=\sum_{d\mid n}\varphi(d)$ is multiplicative. If
	$n=\prod_{i=1}^{k}p_i^{a_i}$, where the $p_i$ are distinct primes and
	$a_i\in\mathbb{Z}^+$, then
	\[
		\begin{aligned}
			F(n)
			 & =\prod_{i=1}^{k}F(p_i^{a_i})
			=\prod_{i=1}^{k}\sum_{j=0}^{a_i}\varphi(p_i^j)                      \\
			 & =\prod_{i=1}^{k}\left(1+\sum_{j=1}^{a_i}(p_i^j-p_i^{j-1})\right)
			=\prod_{i=1}^{k}p_i^{a_i}=n.
		\end{aligned}
	\]

\end{altproof}
\begin{definition}[divisor power sums $\sigma_k$]\label{def:5.1.3} For $k\ge0$ and $n\in\mathbb{Z}^+$,
	put $\sigma_k(n) = \sum_{d\mid n} d^{k}$; thus $\sigma_0 = \tau$ and
	$\sigma_1 = \sigma$.

\end{definition}
\begin{theorem}\label{thm:5.1.6} For every integer $k\ge0$, the
	function $\sigma_k$ is multiplicative.

\end{theorem}

\begin{proof} The function $d\mapsto d^k$ is
	multiplicative, so Theorem~\ref{thm:5.1.3} applied to its divisor sum gives the
	claim.

\end{proof}
\begin{theorem}\label{thm:5.1.7} If
	$n=\prod_{i=1}^{r}p_i^{a_i}$, where $p_1,p_2,\ldots,p_r$ are
	distinct primes and $a_1,a_2,\ldots,a_r\in\mathbb{Z}^+$, then
	$\sigma_k(n)=\prod_{i=1}^{r}\tfrac{p_i^{k(a_i+1)}-1}{p_i^k-1}$ for
	$k\ge1$.

\end{theorem}

\begin{proof} Since $\sigma_k$ is multiplicative
	(Theorem~\ref{thm:5.1.6}),
	$\sigma_k(n)=\prod_{i=1}^{r}\sigma_k(p_i^{a_i})$. For a prime
	$p$ and $a\in\mathbb{Z}^+$, the divisors of $p^{a}$ are exactly
	$1,p,\ldots,p^{a}$, so
	\[
		\sigma_k(p^{a})=\sum_{j=0}^{a}(p^{j})^{k}=\sum_{j=0}^{a}(p^{k})^{j}
		=\frac{(p^{k})^{a+1}-1}{p^{k}-1}=\frac{p^{k(a+1)}-1}{p^{k}-1},
	\]
	where the geometric-series identity
	$\sum_{j=0}^{a}x^{j}=\frac{x^{a+1}-1}{x-1}$ is applied with $x=p^{k}$;
	this is valid because $k\ge1$ and $p\ge2$ give $x=p^{k}\ge2$, so
	$x\neq1$. Multiplying these factors over $i=1,\ldots,r$ yields
	$\sigma_k(n)=\prod_{i=1}^{r}\tfrac{p_i^{k(a_i+1)}-1}{p_i^k-1}$.

\end{proof}
\begin{theorem}\label{thm:5.1.8} For $n\in\mathbb{Z}^+$, the number $\tau(n)$
	is odd if and only if $n$ is a perfect square.

\end{theorem}

\begin{proof} Write
	$n=\prod_{i=1}^{r}p_i^{a_i}$, where $p_1,p_2,\ldots,p_r$ are distinct
	primes and $a_1,a_2,\ldots,a_r\in\mathbb{Z}^+$. Then
	$\tau(n)=\prod_{i=1}^{r}(a_i+1)$ is odd if and only if every $a_i$ is
	even, which holds if and only if $n$ is a perfect square.

\end{proof}
\begin{theorem}\label{thm:5.1.9} For
	every $n\ge1$, we have $\tau(n)\le 2\sqrt{n}$.

\end{theorem}

\begin{proof} By Definition~\ref{def:5.1.2}, $\tau(n)=\sum_{d\mid n}1$
	counts the positive divisors of $n$. The map
	$d\mapsto n/d$ is an involution on those divisors, since $n/(n/d)=d$,
	and groups them into
	pairs $\{d,\,n/d\}$, with a singleton when $d=n/d$ (that is, $d=\sqrt{n}$).
	In each pair, at least one member is
	at most $\sqrt{n}$: if both $d>\sqrt{n}$ and $n/d>\sqrt{n}$, then $n=d\cdot(n/d)>\sqrt{n}\cdot\sqrt{n}=n$,
	which is impossible. So every divisor of $n$ is either at most $\sqrt{n}$ or
	equal to $n/d$ for some divisor $d\le\sqrt{n}$, and $\tau(n)$ is at most twice the number of divisors not exceeding $\sqrt{n}$. Those divisors lie among $1,2,\ldots,\lfloor\sqrt{n}\rfloor$, so there are at most $\lfloor\sqrt{n}\rfloor$ of them.
	Hence $\tau(n)\le 2\lfloor\sqrt{n}\rfloor\le 2\sqrt{n}$. If $n$
	is a perfect square, the middle divisor $\sqrt n$ is counted only once,
	so the inequality still holds, although the estimate is not sharp.

\end{proof}
\begin{theorem}\label{thm:5.1.10} Let $f$ and $g$ be multiplicative
	functions. Then $f(n)=\prod_{p\mid n}f\bigl(p^{v_p(n)}\bigr)$ for every
	$n\in\mathbb{Z}^+$; consequently $f$ is determined by its values at $1$ and
	at the prime powers, and $f=g$ whenever $f(p^{a})=g(p^{a})$ for every prime $p$ and every $a\ge1$.

\end{theorem}

\begin{proof} By Theorem~\ref{thm:5.1.1}, $f(1)=1$, which is the
	empty product and settles the case $n=1$. For $n>1$ write
	$n=\prod_{i=1}^{k}p_i^{a_i}$ with distinct primes $p_i$
	(Theorem~\ref{thm:2.3.3}) and induct on $k$: the factors $p_1^{a_1}$ and
	$\prod_{i\ge2}p_i^{a_i}$ are coprime, so multiplicativity gives
	$f(n)=f(p_1^{a_1})\,f\big(\prod_{i\ge2}p_i^{a_i}\big)$, and the induction
	hypothesis gives the expansion of the second factor. The remaining assertions follow at
	once, since the displayed product involves only values at prime powers.

\end{proof}
\subsection{Perfect and Amicable Numbers}

\begin{definition}[perfect number]\label{def:5.2.1}
	A positive integer $n$ is perfect if $\sigma(n) = 2n$.

\end{definition}
\begin{theorem}[Euclid--Euler theorem]\label{thm:5.2.1} An even positive integer
	$n$ is perfect if and only if
	$n = 2^{p-1}(2^{p}-1)$ for some prime $p$ such that $2^{p}-1$ is a Mersenne prime.

\end{theorem}

\begin{proof} ($\Leftarrow$) If $q = 2^p-1$ is prime,
	then $\sigma(2^{p-1}q) = \sigma(2^{p-1})\sigma(q) = (2^p-1)(q+1) =
		q\cdot 2^p = 2\cdot 2^{p-1}q$, so $n$ is perfect.

	($\Rightarrow$) Write $n=2^km$ with $k\ge1$ and $m$ a positive odd integer. Then $2n = \sigma(n) =
		(2^{k+1}-1)\sigma(m)$, so $2^{k+1}\mid (2^{k+1}-1)\sigma(m)$; as
	$\gcd(2^{k+1}-1, 2^{k+1})=1$, we get $2^{k+1}\mid\sigma(m)$, say
	$\sigma(m)=2^{k+1}t$ with $t\in\mathbb{Z}^+$. It follows that
	$m=(2^{k+1}-1)t$ and $\sigma(m)=m+t$.

	If $t>1$,
	then $1, t, m$ are three distinct divisors of $m$ (since $t\mid m$
	via $m=(2^{k+1}-1)t$, and $1<t<m$), so $\sigma(m)\ge 1+t+m$,
	contradicting $\sigma(m)=m+t$. Thus $t=1$ and $\sigma(m) = m+1$, so $m$
	is prime and $m=2^{k+1}-1$.

	Put $p=k+1\ge2$. If $p$ is composite, write $p =
		ab$ with $1 < a < p$; then $2^{p}-1 = (2^{a})^{b} - 1 =
			(2^{a}-1)\big((2^{a})^{b-1} + (2^{a})^{b-2} + \cdots + 1\big)$, so
	$2^{a}-1$ is a proper nontrivial divisor of $2^{p}-1$, with $1 < 2^{a}-1 < 2^{p}-1$, contradicting the primality of $2^{p}-1 = m$.
	Hence $p$ is prime.

\end{proof}
\begin{theorem}\label{thm:5.2.2} Let $d,n\in\mathbb{Z}^+$ with $d\mid n$ and
	$d<n$. Then
	$\tfrac{\sigma(d)}{d}<\tfrac{\sigma(n)}{n}$. In particular, if $n$ is
	a perfect number (so $\sigma(n)/n=2$ by Definition~\ref{def:5.2.1}), then
	$\sigma(d)/d<2$ for every proper divisor $d$ of $n$, while
	$\sigma(M)/M>2$ for every proper multiple $M\in\mathbb{Z}^+$ of $n$.

\end{theorem}

\begin{proof} First rewrite each ratio as a sum of
	reciprocals of divisors. Since $e\mapsto n/e$ is a bijection of the
	set of divisors of $n$ onto itself,
	$\tfrac{\sigma(n)}{n}=\tfrac{1}{n}\sum_{e\mid n}e=\sum_{e\mid
			n}\tfrac{n/e}{n}=\sum_{e\mid n}\tfrac{1}{e}$, and the same argument
	for $d$ gives $\tfrac{\sigma(d)}{d}=\sum_{e\mid d}\tfrac{1}{e}$.
	Every divisor of $d$ divides $n$ (since $d\mid n$), so
	the divisors of $d$ form a subset of those of $n$, and
	$\sum_{e\mid d}\tfrac{1}{e}$ is a subsum of
	$\sum_{e\mid n}\tfrac{1}{e}$. The inclusion is strict: $n$
	divides itself, contributing $\tfrac{1}{n}$ to $\sum_{e\mid
			n}\tfrac{1}{e}$, yet $n\nmid d$ since $0<d<n$, so $\tfrac{1}{n}$ is absent
	from $\sum_{e\mid d}\tfrac{1}{e}$. As all terms $1/e$ are strictly
	positive, omitting at least the term $1/n$ strictly decreases the
	sum: $\sum_{e\mid d}\tfrac{1}{e}<\sum_{e\mid n}\tfrac{1}{e}$, that is, $\tfrac{\sigma(d)}{d}<\tfrac{\sigma(n)}{n}$. Now suppose $n$ is
	perfect, so $\sigma(n)/n=2$. For a proper divisor $d$ of $n$ the
	inequality gives $\sigma(d)/d<2$; for a proper multiple $M$ of $n$ we
	have $n\mid M$ and $n<M$, so applying the inequality with the pair
	$(n,M)$ in place of $(d,n)$ yields $2=\sigma(n)/n<\sigma(M)/M$.

\end{proof}
\begin{theorem}\label{thm:5.2.3} If $n$ is an even perfect number with $3\mid
		n$, then $n=6$.

\end{theorem}

\begin{proof} By the Euclid--Euler theorem (Theorem~\ref{thm:5.2.1}), $n=2^{p-1}(2^{p}-1)$ with $2^{p}-1$ prime. As $3\nmid
		2^{p-1}$, the hypothesis $3\mid n$ forces $3\mid 2^{p}-1$, that is, $2^{p}\equiv1\pmod{3}$; since $2\equiv-1\pmod{3}$, this requires $p$ to be even,
	and the only even prime is $p=2$, giving $n=2(2^{2}-1)=6$.

\end{proof}
\begin{theorem}\label{thm:5.2.4} Every odd perfect number $n$ satisfies
	$n\equiv1\pmod{4}$.

\end{theorem}

\begin{proof} Write $n=\prod_{i=1}^{r}p_i^{a_i}$, where
	$p_1,p_2,\ldots,p_r$ are distinct odd primes and
	$a_1,a_2,\ldots,a_r\in\mathbb{Z}^+$. Then
	$\prod_{i=1}^{r}\sigma(p_i^{a_i})=\sigma(n)=2n$. As $n$ is odd,
	$2n\equiv2\pmod{4}$, so $v_2(\sigma(n))=1$: exactly one factor $\sigma(p_i^{a_i})$ is even, and that factor is congruent to $2\pmod{4}$, while the remaining factors are odd.
	Now $\sigma(p^{a})=1+p+\cdots+p^{a}$ has $a+1$ odd terms, so it is
	odd if and only if $a$ is even. Thus exactly one exponent, say $a_1$, is odd and
	the others are even. For the exceptional prime $p=p_1$, if $p\equiv3\pmod{4}$, then
	$1,p,p^2,\ldots\equiv1,3,1,3,\ldots\pmod{4}$ and the $a+1$ terms (an even number of them)
	group into pairs $(1,3)$, which gives $\sigma(p^{a})\equiv0\pmod{4}$ and
	contradicts the earlier congruence $\sigma(p^{a})\equiv2\pmod{4}$; hence $p\equiv1\pmod{4}$, and then
	$\sigma(p^{a})\equiv a+1\pmod{4}$, so $\sigma(p^{a})\equiv2\pmod{4}$ forces
	$a\equiv1\pmod{4}$. Consequently $p_1^{a_1}\equiv1\pmod{4}$ (as
	$p_1\equiv1$), and $\prod_{i=2}^{r}p_i^{a_i}$ is the square of an odd integer,
	also congruent to $1\pmod{4}$; therefore $n\equiv1\pmod{4}$.

\end{proof}
\begin{definition}[amicable pair]\label{def:5.2.2} For
	$a,b\in\mathbb{Z}^+$ with $a\neq b$, the pair $(a,b)$ is amicable if
	$\sigma(a) = \sigma(b) = a+b$ (equivalently,
	$s(a)=b$ and $s(b)=a$, where $s(n)=\sigma(n)-n$ is the sum of the
	proper positive divisors of $n$).

\end{definition}
\begin{theorem}[Th{\=a}bit's rule]\label{thm:5.2.5} Let $n$ be an integer with $n\ge2$, and put
	\[
		p=3\cdot2^{n-1}-1,\qquad q=3\cdot2^{n}-1,\qquad r=9\cdot2^{2n-1}-1.
	\]
	If $p,q,r$ are all prime, then $(2^{n}pq,\ 2^{n}r)$ is an amicable pair.

\end{theorem}

\begin{proof} The three primes are odd and distinct
	($p<q<r$ for $n\ge2$), so $2^{n},p,q$ are pairwise coprime, as are
	$2^{n}$ and $r$. Since $\sigma$ is multiplicative (Theorem~\ref{thm:5.1.2}(b))
	with $\sigma(2^{n})=2^{n+1}-1$ and $\sigma(\ell)=\ell+1$ for a prime $\ell$,
	\[
		\sigma(2^{n}pq)=(2^{n+1}-1)(p+1)(q+1)
		=(2^{n+1}-1)\cdot3\cdot2^{n-1}\cdot3\cdot2^{n}
		=(2^{n+1}-1)\,9\cdot2^{2n-1},
	\]
	and $\sigma(2^{n}r)=(2^{n+1}-1)(r+1)=(2^{n+1}-1)\,9\cdot2^{2n-1}$ as well.
	For the sum, $pq=9\cdot2^{2n-1}-9\cdot2^{n-1}+1$, so
	\[
		pq+r=\big(9\cdot2^{2n-1}-9\cdot2^{n-1}+1\big)+\big(9\cdot2^{2n-1}-1\big)
		=9\cdot2^{n-1}\big(2^{n+1}-1\big),
	\]
	whence $2^{n}pq+2^{n}r=2^{n}(pq+r)=(2^{n+1}-1)\,9\cdot2^{2n-1}$. Thus
	$\sigma(2^{n}pq)=\sigma(2^{n}r)=2^{n}pq+2^{n}r$. Finally, the two numbers
	are distinct, since $pq-r=2-9\cdot2^{n-1}<0$ for $n\ge2$.

\end{proof}
\subsection{Dirichlet Convolution and M{\"o}bius Inversion}

\begin{definition}[Dirichlet convolution]\label{def:5.3.1} For
	functions $f,g:\mathbb{Z}^+\to\mathbb{C}$, define
	$(f*g)(n) = \sum_{d \mid n} f(d)\, g\big(\tfrac{n}{d}\big) = \sum_{de =
			n} f(d) g(e)$, where the last sum is over ordered pairs of positive
	integers with product $n$.

\end{definition}
\begin{definition}[constant-one function]\label{def:5.3.2}
	The constant-one arithmetic function is $\mathbf{1}(n)=1$.

\end{definition}
\begin{definition}[convolution identity]\label{def:5.3.3}
	The arithmetic function $\varepsilon$ is given by
	$\varepsilon(1)=1$ and $\varepsilon(n)=0$ for $n>1$.

\end{definition}
\begin{definition}[identity function]\label{def:5.3.4}
	The identity arithmetic function is $\id(n)=n$.

\end{definition}
\begin{definition}[logarithm function]\label{def:5.3.5}
	In Dirichlet convolutions, we write $\log$ for the arithmetic function
	$\log(n)=\ln n$; elsewhere $\log_b$ retains its usual meaning.

\end{definition}
\begin{theorem}\label{thm:5.3.1}
	Dirichlet convolution is commutative and associative, with two-sided
	identity $\varepsilon$ (that is, $f*\varepsilon=\varepsilon*f=f$ for
	all $f$).

\end{theorem}

\begin{proof} Commutativity follows from $(f*g)(n) = \sum_{de=n}
		f(d)g(e) = \sum_{ed=n} g(e)f(d) = (g*f)(n)$. For associativity, all
	sums involved are finite, and
	\[
		\begin{aligned}
			((f*g)*h)(n)
			 & =\sum_{mc=n}(f*g)(m)h(c)
			=\sum_{mc=n}\sum_{ab=m}f(a)g(b)h(c) \\
			 & =\sum_{abc=n}f(a)g(b)h(c),
		\end{aligned}
	\]
	the sum over all ordered triples of positive integers with product
	$n$. Expanding $(f*(g*h))(n)$ groups the same triples first
	by $bc$, so the two sides agree. For the identity,
	$(f*\varepsilon)(n) = \sum_{d\mid n}
		f(d)\,\varepsilon(n/d) = f(n)$, since $\varepsilon(n/d)$ vanishes
	unless $d=n$; commutativity gives $\varepsilon*f=f$ as well.

\end{proof}
\begin{definition}[M{\"o}bius function]\label{def:5.3.6} For $n
		\in \mathbb{Z}^+$, define the M{\"o}bius function by
	\[
		\mu(n) =
		\begin{cases} 1, & \text{if } n = 1,      \\ 0, & \text{if } p^{2} \mid n
		              \text{ for some prime } p, \\ (-1)^{r}, & \text{if } n =
		              \prod_{i=1}^{r} p_i,\ p_i \text{ distinct primes}.
		\end{cases}
	\]

\end{definition}
\begin{theorem}\label{thm:5.3.2}\statementitems

	\begin{enumerate}[label={\alph*)}]
		\item $\sum_{d \mid n} \mu(d) =
			      \begin{cases} 1, & \text{if } n = 1, \\ 0, & \text{if } n \ge 2
			      \end{cases}$.
		\item The arithmetic function $\mu(n)$ is multiplicative.
		\item For functions $f,F:\mathbb{Z}^+\to\mathbb{C}$,
		      we have $F(n) = \sum_{d \mid n} f(d)$ for every $n\ge1$ if and only if
		      $f(n) = \sum_{d \mid n} \mu(d) F\big(\tfrac{n}{d}\big)
			      = \sum_{d \mid n}\mu\big(\tfrac{n}{d}\big) F(d)$ for every $n\ge1$
		      (the M{\"o}bius inversion formula). Equivalently, $F=f*\mathbf{1}$ if and only if
		      $f=F*\mu$.
	\end{enumerate}

\end{theorem}

\begin{proof}

	\begin{enumerate}[label={\alph*)}]
		\item If $n = 1$, the sum is $\mu(1) = 1$. Let $n\ge2$, and let
		      $k\ge1$ be the number of distinct prime factors of $n$. By unique
		      prime factorization, the squarefree divisors of $n$ (those divisible by
		      no square $p^2$ with $p$ prime) correspond exactly
		      to subsets of these $k$ primes; a subset of size $j$ contributes
		      $(-1)^j$. All nonsquarefree divisors contribute $0$, so
		      $\sum_{d \mid n} \mu(d) = \sum_{j=0}^{k} \binom{k}{j}(-1)^{j} = (1 -
			      1)^{k} = 0$.
		\item Let $m,n\in\mathbb{Z}^+$ with $\gcd(m,n)=1$. If $m=1$ or
		      $n=1$, the result is immediate. If some $p^2$ divides $m$ or $n$, then
		      $p^2\mid mn$, so both $\mu(mn)$ and $\mu(m)\mu(n)$ are $0$. Otherwise,
		      $m=\prod_{i=1}^{r}p_i$ and $n=\prod_{j=1}^{s}q_j$ are squarefree.
		      Coprimality and unique prime factorization imply that the two sets of
		      primes are disjoint, so $mn$ has exactly $r+s$ distinct prime factors.
		      Thus $\mu(mn)=(-1)^{r+s}=\mu(m)\mu(n)$.
		\item Suppose first that $F(m)=\sum_{c\mid m}f(c)$ for every
		      $m\ge1$. Then
		      \[
			      \begin{aligned}
				      \sum_{d\mid n}\mu(d)F\big(\tfrac{n}{d}\big)
				       & =\sum_{d\mid n}\mu(d)\sum_{c\mid\frac{n}{d}}f(c)       \\
				       & =\sum_{c\mid n}f(c)\sum_{d\mid\frac{n}{c}}\mu(d)=f(n).
			      \end{aligned}
		      \]
		      All these sums are finite. The middle equality regroups the sum over
		      all pairs $(d,c)$ such that $d\mid n$ and $c\mid n/d$, or equivalently
		      $dc\mid n$, by summing over $c$ first; the last equality follows from
		      (a), since the inner sum is $1$ exactly when $c=n$ and is $0$
		      otherwise. The map $d\mapsto n/d$ is an involution on the positive
		      divisors of $n$, so we may equivalently write
		      $f(n)=\sum_{d\mid n}\mu\big(\tfrac{n}{d}\big)F(d)$.

		      Conversely, suppose that
		      $f(m)=\sum_{d\mid m}\mu(d)F(m/d)$ for every $m\ge1$. Then, again
		      regrouping only finite sums,
		      \[
			      \begin{aligned}
				      \sum_{c\mid n}f(c)
				       & =\sum_{c\mid n}\sum_{d\mid c}\mu(d)F\big(\tfrac{c}{d}\big) \\
				       & =\sum_{e\mid n}F(e)\sum_{d\mid\frac{n}{e}}\mu(d)=F(n),
			      \end{aligned}
		      \]
		      where in the middle sum we wrote $c=de$; the last equality again
		      follows from (a). This proves the converse and hence the equivalence.
	\end{enumerate}

\end{proof}
\begin{altproof}[Alternative proof of Theorem~\ref{thm:5.3.2}(c)]
	Part (a) and commutativity give
	$\mu*\mathbf{1}=\mathbf{1}*\mu=\varepsilon$. Therefore, by
	associativity and the identity law from Theorem~\ref{thm:5.3.1}, convolving both
	sides with $\mu$ gives
	\[
		F=f*\mathbf{1}\ \Rightarrow\ F*\mu
		=(f*\mathbf{1})*\mu=f*(\mathbf{1}*\mu)
		=f*\varepsilon=f.
	\]
	Conversely, convolving both sides of $f=F*\mu$ with $\mathbf{1}$ gives
	\[
		f*\mathbf{1}=(F*\mu)*\mathbf{1}
		=F*(\mu*\mathbf{1})=F*\varepsilon=F.
	\]
	Thus $F=f*\mathbf{1}$ if and only if $f=F*\mu$.

\end{altproof}
\begin{theorem}\label{thm:5.3.3} Let $f$ be an
	arithmetic function. Then, as formal power series in $\mathbb{C}[[x]]$
	(and for $|x|<1$ whenever the right-hand side converges absolutely),
	\[
		\sum_{n=1}^{\infty}\frac{f(n)x^{n}}{1-x^{n}}
		=\sum_{N=1}^{\infty}\Bigl(\sum_{d\mid N}f(d)\Bigr)x^{N}.
	\]
	In particular, $\sum_{n\ge1}\varphi(n)x^{n}/(1-x^{n})=x/(1-x)^{2}$ and
	$\sum_{n\ge1}\mu(n)x^{n}/(1-x^{n})=x$.

\end{theorem}

\begin{proof} Formal series with complex coefficients use the same
	coefficientwise addition and finite-convolution multiplication as in
	Definition~\ref{def:3.5.2}. The geometric-series expansion gives
	$\dfrac{x^{n}}{1-x^{n}}=\sum_{m=1}^{\infty}x^{nm}$, so
	\[
		\sum_{n=1}^{\infty}\frac{f(n)x^{n}}{1-x^{n}}
		=\sum_{n=1}^{\infty}\sum_{m=1}^{\infty}f(n)x^{nm}.
	\]
	The terms contributing to the coefficient of $x^{N}$ are precisely those
	indexed by pairs $(n,m)$ with $nm=N$. For each divisor $d\mid N$, the pair
	$(d,N/d)$ contributes $f(d)$, giving the coefficient $\sum_{d\mid N}f(d)$.

	For the analytic identity, the case $x=0$ is immediate. Put $r=|x|\in(0,1)$ and
	$F(N)=\sum_{d\mid N}f(d)$, and suppose $\sum_{N\ge1}|F(N)|r^N<\infty$.
	M{\"o}bius inversion and $|\mu(d)|\le1$ give
	$|f(n)|\le\sum_{d\mid n}|F(d)|$.

	For each positive integer $M$,
	\[
		\begin{aligned}
		\sum_{n=1}^{M}\frac{|f(n)|r^n}{1-r^n}
		&\le\frac{1}{1-r}\sum_{d=1}^{M}|F(d)|
		       \sum_{j=1}^{\lfloor M/d\rfloor}r^{dj}\\
		&\le\frac{1}{1-r}\sum_{d\ge1}|F(d)|\frac{r^d}{1-r^d}\\
		&\le\frac{1}{(1-r)^2}\sum_{d\ge1}|F(d)|r^d<\infty.
		\end{aligned}
	\]
	The first inequality follows by applying the bound on $|f(n)|$ and
	the inequality $1-r^n\ge1-r$, then regrouping a finite sum. The subsequent estimates
	use nonnegative geometric series and $1-r^d\ge1-r$. This bound is independent of $M$, so the double
	series $\sum_{n,m\ge1}f(n)x^{nm}$ converges absolutely. It may therefore
	be regrouped by $N=nm$, proving the analytic identity.

	For the two special cases,
	$\sum_{d\mid N}\varphi(d)=N$ (Theorem~\ref{thm:5.1.5}) together with
	$\sum_{N\ge1}Nx^{N}=x/(1-x)^{2}$ (from
	$(1-x)\sum_{N\ge1}Nx^{N}=\sum_{N\ge1}x^{N}=x/(1-x)$, where the first equality
	holds coefficientwise by telescoping) gives the first identity. For the second identity,
	$\sum_{d\mid N}\mu(d)$ equals $1$ when $N=1$ and $0$ when $N>1$ (Theorem~\ref{thm:5.3.2}(a)); hence only the term $x$ remains.

\end{proof}
\begin{theorem}\label{thm:5.3.4} Let $m,n\in\mathbb{Z}^+$. A positive integer
	$d$ is a cyclic period of a word $w=(w_0,\ldots,w_{n-1})$ if
	$w_{j+d}=w_j$ for every $j$, with indices read modulo $n$.
	For words on an alphabet of $m$ colors, let $A_m(n)$ count those of
	length $n$ whose least positive cyclic period is $n$. Then
	\[
		A_m(n)=\sum_{d\mid n}\mu\big(\tfrac{n}{d}\big)m^d.
	\]
	Consequently, the number $L_m(n)$ of aperiodic necklaces of length $n$ and the number $N_m(n)$ of all necklaces of length $n$, where words are identified exactly when one is a cyclic rotation of the other, are
	\[
		L_m(n)=\frac1n\sum_{d\mid n}\mu(d)m^{n/d},
		\qquad
		N_m(n)=\frac1n\sum_{d\mid n}\varphi(d)m^{n/d}.
	\]

\end{theorem}

\begin{proof} Every word has a least positive cyclic period $d$, since
	a shift by $n$ fixes the word. If a nonnegative shift $s$ also fixes it,
	write $s=qd+r$ with $0\le r<d$. Shifting by $d$, or by its inverse,
	fixes the word, so shifting by $r=s-qd$ does too. Minimality of $d$
	forces $r=0$. Thus the shifts fixing the word are exactly the multiples
	of $d$; in particular, $d\mid n$.

	A word of least cyclic period $d\mid n$ is obtained by repeating its
	first $d$ letters. This block has least cyclic period $d$, since a
	smaller cyclic period of the block would fix the entire word.
	Conversely, repeating any such block produces a word of least cyclic
	period $d$: the least period of the resulting word divides $d$, by the preceding argument,
	and cannot be smaller. This gives a bijection between such words and the $A_m(d)$ blocks.
	Partitioning all words by their least cyclic period gives
	\[
		m^n=\sum_{d\mid n}A_m(d).
	\]
	M{\"o}bius inversion (Theorem~\ref{thm:5.3.2}(c)) gives
	\[
		A_m(n)=\sum_{d\mid n}\mu\big(\tfrac{n}{d}\big)m^d
		=\sum_{d\mid n}\mu(d)m^{n/d}.
	\]
	If a word has least cyclic period $d$, its distinct rotations are the
	$d$ shifts by $0,1,\ldots,d-1$. Thus each class of words identified by
	rotation has exactly $d$ elements. In particular, $L_m(n)=A_m(n)/n$.
	For each $d\mid n$, the $A_m(d)$ words of length $n$ with least cyclic
	period $d$ form $A_m(d)/d$ such classes. Hence
	\[
		N_m(n)=\sum_{d\mid n}\frac{A_m(d)}{d}.
	\]
	Substituting the formula for $A_m(d)$ and writing $d=ek$ gives
	\[
		\begin{aligned}
			N_m(n)
			 & =\sum_{e\mid n}m^e\sum_{k\mid n/e}\frac{\mu(k)}{ek}    \\
			 & =\frac1n\sum_{e\mid n}\varphi\big(\tfrac{n}{e}\big)m^e
			=\frac1n\sum_{d\mid n}\varphi(d)m^{n/d},
		\end{aligned}
	\]
	because M{\"o}bius inversion applied to $\sum_{d\mid r}\varphi(d)=r$ (Theorem~\ref{thm:5.1.5}) gives $\varphi(r)=r\sum_{k\mid r}\mu(k)/k$. This proves all three formulas.

	For a word of least cyclic
	period $d$, the shifts fixing it form a subgroup of $(\mathbb{Z}_n,+)$.
	The $d$ cosets of this subgroup correspond to the distinct rotations
	(Definition~\ref{def:6.6.4}), as the proof above shows directly.

\end{proof}
\begin{theorem}\label{thm:5.3.5} If $f$ and $g$ are multiplicative functions, then
	$f*g$ is multiplicative.

\end{theorem}

\begin{proof} Let $m,n\in\mathbb{Z}^+$ with
	$\gcd(m,n)=1$. By unique prime factorization, every divisor of $mn$ has
	a unique factorization $d_1d_2$ with $d_1\mid m$ and $d_2\mid n$; moreover,
	$\gcd(d_1,d_2)=\gcd(m/d_1,n/d_2)=1$. Hence
	\[
		\begin{aligned}
			(f*g)(mn)
			 & =\sum_{d_1\mid m}\sum_{d_2\mid n}
			f(d_1d_2)g\big(\tfrac{mn}{d_1d_2}\big)                         \\
			 & =\sum_{d_1\mid m}\sum_{d_2\mid n}
			f(d_1)f(d_2)g\big(\tfrac{m}{d_1}\big)g\big(\tfrac{n}{d_2}\big) \\
			 & =(f*g)(m)(f*g)(n).
		\end{aligned}
	\]
	Also $(f*g)(1)=f(1)g(1)=1$, so $f*g$ is not identically zero and is
	multiplicative.

\end{proof}
\begin{theorem}\label{thm:5.3.6} A Dirichlet inverse of an arithmetic function
	$f$ is an arithmetic function $g$ such that $f*g=\varepsilon$.
	A function $f:\mathbb{Z}^+\to\mathbb{C}$ has such an inverse if and only if
	$f(1)\neq0$. The inverse is unique. If $f$ is multiplicative, then its
	Dirichlet inverse is multiplicative. Consequently, the invertible
	functions form an abelian group under Dirichlet convolution, the
	multiplicative functions form a subgroup, and
	$\mu=\mathbf{1}^{-1}$.

\end{theorem}

\begin{proof} If $f*g=\varepsilon$, then evaluating at
	$1$ gives $f(1)g(1)=1$, so $f(1)\neq0$. Conversely, suppose
	$f(1)\neq0$. The equations $(f*g)(n)=\varepsilon(n)$ determine $g$
	recursively and uniquely by
	\[
		g(1)=\frac{1}{f(1)},\qquad
		g(n)=-\frac{1}{f(1)}
		\sum_{\substack{d\mid n\\d>1}}f(d)g\big(\tfrac{n}{d}\big)
		\quad(n>1),
	\]
	because $n/d<n$ when $d>1$. The resulting function satisfies
	$f*g=\varepsilon$, and commutativity gives $g*f=\varepsilon$.

	Now suppose that $f$ is multiplicative. Define $h(1)=1$ and, for each prime $p$, define recursively
	\[
		h(p^a)=-\sum_{j=1}^{a}f(p^j)h(p^{a-j})\quad(a\ge1).
	\]
	For $n>1$, define $h(n)=\prod_{p\mid n}h(p^{v_p(n)})$.
	This is well defined by unique factorization; coprime integers have
	disjoint sets of prime factors, so $h$ is multiplicative and $h(1)=1$.
	Theorem~\ref{thm:5.3.5} shows that $f*h$ is multiplicative, while the recursion gives
	$(f*h)(p^a)=0$ for $a\ge1$ and $(f*h)(1)=1$.
	The function $\varepsilon$ is also multiplicative and has the same values at $1$ and all prime powers. Theorem~\ref{thm:5.1.10} therefore
	gives $f*h=\varepsilon$, and uniqueness gives
	$h=f^{-1}$; hence $f^{-1}$ is multiplicative. The group assertions
	now follow from Theorem~\ref{thm:5.3.1} and Theorem~\ref{thm:5.3.5}. Finally, Theorem~\ref{thm:5.3.2}(a) says $\mu*\mathbf{1}=\varepsilon$, so
	$\mu=\mathbf{1}^{-1}$.

\end{proof}

\subsection{Incidence Algebras and M{\"o}bius Inversion}

\begin{definition}[incidence algebra and the M{\"o}bius function of a poset]\label{def:5.4.1} Let $(P,\le)$ be a locally finite partially ordered set (poset), meaning that every interval $[x,y]=\{z\in P \mid x\le z\le y\}$ is finite. The incidence algebra $I(P)$ over $\mathbb{C}$ consists of the functions $F(x,y)$ defined for $x\le y$, with pointwise addition and convolution
	\[
		(F*G)(x,y)=\sum_{x\le z\le y}F(x,z)G(z,y).
	\]
	Its identity is $\delta(x,y)=[\,x=y\,]$ (Iverson's bracket: $[Q]=1$ if $Q$ holds and $0$ otherwise), and its zeta function is $\zeta_P(x,y)=1$ for $x\le y$. The M{\"o}bius function $\mu_P$ of $P$ is the convolution inverse of $\zeta_P$.

\end{definition}
\begin{theorem}[M{\"o}bius inversion on a locally finite poset]\label{thm:5.4.1} The function $\zeta_P$ is invertible in $I(P)$. Its inverse is determined recursively by
	\[
		\mu_P(x,x)=1,\qquad
		\mu_P(x,y)=-\sum_{x\le z<y}\mu_P(x,z)\quad(x<y).
	\]
	Equivalently, it satisfies
	\[
		\mu_P(x,y)=-\sum_{x<z\le y}\mu_P(z,y)\quad(x<y).
	\]
	Let $f,g:P\to\mathbb{C}$, and assume that every principal lower set $\{x \mid x\le y\}$ is finite. Then
	\[
		g(y)=\sum_{x\le y}f(x)
		\quad\Longleftrightarrow\quad
		f(y)=\sum_{x\le y}\mu_P(x,y)g(x).
	\]
	For a finite poset, extend $\zeta_P(x,y)$ and $\mu_P(x,y)$ by $0$ when $x\nleq y$. Once a linear extension has been chosen, the matrix $(\zeta_P(x,y))$ is upper triangular with diagonal entries $1$, and $(\mu_P(x,y))$ is its inverse matrix.

\end{theorem}

\begin{proof} The convolution is well defined because every interval is finite. It is associative: for $x\le y$, both $((F*G)*H)(x,y)$ and $(F*(G*H))(x,y)$ equal the finite sum
	\[
		\sum_{x\le u\le v\le y}F(x,u)G(u,v)H(v,y).
	\]
	Also
	\[
		(F*\delta)(x,y)=\sum_{x\le z\le y}F(x,z)[\,z=y\,]=F(x,y)
	\]
	and similarly $\delta*F=F$, so $\delta$ is the identity.

	Define $\mu_R(x,x)=1$ and, for $x<y$, define $\mu_R(x,y)$ by the first recursion. Induction on the finite number $|[x,y]|$ is legitimate, since every interval $[x,z]$ with $z<y$ has fewer elements than $[x,y]$. The recursion says exactly
	\[
		(\mu_R*\zeta_P)(x,y)=\sum_{x\le z\le y}\mu_R(x,z)=[\,x=y\,],
	\]
	so $\mu_R*\zeta_P=\delta$. Similarly, the second recursion defines a function $\mu_L$ with $\zeta_P*\mu_L=\delta$. Associativity gives
	\[
		\mu_L=\delta*\mu_L=(\mu_R*\zeta_P)*\mu_L
		=\mu_R*(\zeta_P*\mu_L)=\mu_R,
	\]
	so the common function $\mu_P$ is the two-sided inverse of $\zeta_P$ and satisfies both recursions.

	Suppose $g(y)=\sum_{x\le y}f(x)$. The sums involved are finite, so they may be rearranged:
	\[
		\begin{aligned}
			\sum_{x\le y}\mu_P(x,y)g(x)
			 & =\sum_{x\le y}\mu_P(x,y)\sum_{u\le x}f(u)      \\
			 & =\sum_{u\le y}f(u)\sum_{u\le x\le y}\mu_P(x,y)
			=\sum_{u\le y}f(u)\,\delta(u,y)=f(y),
		\end{aligned}
	\]
	where the penultimate equality follows from $\zeta_P*\mu_P=\delta$. Conversely, if $f(y)=\sum_{x\le y}\mu_P(x,y)g(x)$, then
	\[
		\begin{aligned}
			\sum_{u\le y}f(u)
			 & =\sum_{u\le y}\sum_{x\le u}\mu_P(x,u)g(x)      \\
			 & =\sum_{x\le y}g(x)\sum_{x\le u\le y}\mu_P(x,u)
			=\sum_{x\le y}g(x)\,\delta(x,y)=g(y),
		\end{aligned}
	\]
	using $\mu_P*\zeta_P=\delta$. This proves the equivalence without assuming that $P$ itself is finite.
	For finite $P$, choose a linear extension by successively removing a minimal
	element. Set $Z=(\zeta_P(x,y))$ and $M=(\mu_P(x,y))$ in this order,
	with row index $x$, column index $y$, and entries $0$ when $x\nleq y$. Then $Z$ is upper triangular with
	diagonal entries $1$. Matrix multiplication agrees with convolution, so
	$MZ=ZM=I$. With $f,g$ as column vectors, the lower-sum formulas read
	$g=Z^{\mathsf T}f$ and $f=M^{\mathsf T}g$; the transposes arise because
	$y$ is the second index of $\zeta_P(x,y)$ and $\mu_P(x,y)$.

\end{proof}
\begin{altproof}[Alternative proof of Theorem~\ref{thm:5.3.2}(c)] Put $P=\mathbb{Z}^+$ with $x\le y$ if and only if $x\mid y$. This poset is locally finite because $[x,y]$ is empty unless $x\mid y$, and if $x\mid y$, then every element of $[x,y]$ is a divisor of $y$. The map $z\mapsto z/x$ identifies $[x,y]$ with $[1,y/x]$, so $\mu_P(x,y)=\mu_P(1,y/x)$.

	We compute the latter value from prime factorization. If $n=\prod_{i=1}^{r}p_i^{a_i}$, then the map $d=\prod_i p_i^{b_i}\mapsto(b_1,\ldots,b_r)$ identifies $[1,n]$ with the product of the chains $0\le b_i\le a_i$. For locally finite posets $P_1,P_2$, the function
	\[
		\nu\big((x_1,x_2),(y_1,y_2)\big)=\mu_{P_1}(x_1,y_1)\mu_{P_2}(x_2,y_2)
	\]
	is the inverse of the zeta function on $P_1\times P_2$, because the sum over $(z_1,z_2)$ factors as
	\[
		\left(\sum_{x_1\le z_1\le y_1}\mu_{P_1}(x_1,z_1)\right)
		\left(\sum_{x_2\le z_2\le y_2}\mu_{P_2}(x_2,z_2)\right)
		=\delta(x_1,y_1)\delta(x_2,y_2).
	\]
	Hence the M{\"o}bius function of a product is the product of the M{\"o}bius functions of its factors. On the chain $1\mid p\mid p^2\mid\cdots$, the recursion gives
	\[
		\mu_P(1,1)=1,\qquad \mu_P(1,p)=-1,\qquad
		\mu_P(1,p^a)=-\sum_{j=0}^{a-1}\mu_P(1,p^j)=0\quad(a\ge2),
	\]
	where the last equality follows inductively because the sum is $1-1+0+\cdots+0$. Therefore
	\[
		\mu_P(1,n)=\prod_{i=1}^{r}\mu_P(1,p_i^{a_i})
		=\begin{cases}
			1,     & n=1,                                               \\
			(-1)^r, & n\text{ is a product of }r\text{ distinct primes}, \\
			0,     & n\text{ is divisible by a prime square},
		\end{cases}
	\]
	which is exactly the arithmetic M{\"o}bius function $\mu(n)$ of Definition~\ref{def:5.3.6}. Thus $\mu_P(d,n)=\mu(n/d)$, and the general inversion formula becomes
	\[
		F(n)=\sum_{d\mid n}f(d)
		\quad\Longleftrightarrow\quad
		f(n)=\sum_{d\mid n}\mu\big(\tfrac{n}{d}\big)F(d),
	\]
	which is Theorem~\ref{thm:5.3.2}(c). This also explains, from the poset calculation, why a square factor forces $\mu(n)=0$: a prime-power chain contributes the value $0$ on that interval as soon as its exponent is at least $2$.

\end{altproof}
\subsection{The von Mangoldt Function}

\begin{definition}[von Mangoldt function]\label{def:5.5.1}
	For $n \in \mathbb{Z}^+$, put $\Lambda(n) = \ln p$ if $n = p^{k}$
	for a prime $p$ and an integer $k\ge1$, and $\Lambda(n) = 0$ otherwise.

\end{definition}
\begin{theorem}\label{thm:5.5.1} For every $n\in\mathbb{Z}^+$, we have $\sum_{d \mid n} \Lambda(d) = \ln
		n$, and $\Lambda(n) = -\sum_{d \mid n}\mu(d)\ln d$.

\end{theorem}

\begin{proof} Only prime-power divisors contribute, so
	$\sum_{d \mid n}\Lambda(d)=\sum_{p\in\Primes}v_p(n)\,\ln p
		=\ln\!\left(\prod_{p\in\Primes}p^{v_p(n)}\right)=\ln n$, where the last equality follows from
	the prime factorization $n=\prod_{p\in\Primes} p^{v_p(n)}$ (Definition~\ref{def:4.5.1}). M{\"o}bius inversion (Theorem~\ref{thm:5.3.2}(c)) gives $\Lambda(n) = \sum_{d \mid n}
		\mu(d)\ln\frac{n}{d} = \big(\sum_{d \mid n}\mu(d)\big)\ln n - \sum_{d
			\mid n}\mu(d)\ln d$. Since $\big(\sum_{d \mid n}\mu(d)\big)\ln n = 0$
	both when $n=1$ and when $n \ge 2$, we get $\Lambda(n) = -\sum_{d \mid
			n}\mu(d)\ln d$.

\end{proof}
\begin{theorem}\label{thm:5.5.2}
	$\varphi=\id*\mu$ and $\Lambda=\mu*\log$.

\end{theorem}

\begin{proof} The identity
	$\sum_{d\mid n}\varphi(d)=n$ (Theorem~\ref{thm:5.1.5}) and M{\"o}bius inversion
	(Theorem~\ref{thm:5.3.2}(c)) give $\varphi=\id*\mu$.
	The identity $\sum_{d\mid n}\Lambda(d)=\ln n$ (Theorem~\ref{thm:5.5.1}) and the
	same inversion formula give $\Lambda=\log*\mu=\mu*\log$, using the
	commutativity of Dirichlet convolution (Theorem~\ref{thm:5.3.1}).

\end{proof}
\subsection{Ramanujan's Sums}

\begin{definition}[Ramanujan sum]\label{def:5.6.1} For $q \in \mathbb{Z}^+$ and $n \in \mathbb{Z}$, the Ramanujan sum is
	\[
		c_q(n) = \sum_{\substack{a=1 \\ \gcd(a, q)=1}}^{q} e^{2\pi i a n / q},
	\]
	the sum of the $n$-th powers of the primitive $q$-th roots of unity.

\end{definition}
\begin{theorem}\label{thm:5.6.1} For all $q \in \mathbb{Z}^+$ and $n \in \mathbb{Z}$,
	\[
		c_q(n) = \sum_{d \mid \gcd(q, n)} d\, \mu\!\left(\tfrac{q}{d}\right).
	\]
	In particular, $c_q(n) \in \mathbb{Z}$; it depends on $n$ only through $\gcd(q, n)$, and we have $c_q(1) = \mu(q)$ and $c_q(q) = \varphi(q)$.

\end{theorem}

\begin{proof} For $m\in\mathbb{Z}$, put
	$G_q(m)=\sum_{a=1}^{q}e^{2\pi i a m/q}$ and $z=e^{2\pi i m/q}$.
	If $q\mid m$, then $z=1$ and $G_q(m)=q$. Otherwise $z\neq1$ and
	$z^q=1$, so the finite geometric-sum formula gives
	\[
		G_q(m)=\sum_{a=1}^{q}z^a=\frac{z(z^q-1)}{z-1}=0.
	\]
	Thus $G_q(m)=q\,[\,q\mid m\,]$.

	Group the terms of $G_q(n)$ by $g=\gcd(a,q)$. For each divisor $g\mid q$,
	the relevant integers are exactly $a=gb$ with $1\le b\le q/g$ and
	$\gcd(b,q/g)=1$. Their contribution is $c_{q/g}(n)$, since
	$e^{2\pi i a n/q}=e^{2\pi i b n/(q/g)}$. Hence
	\[
		\sum_{d\mid q}c_d(n)=G_q(n)=q\,[\,q\mid n\,].
	\]
	Fixing $n$ and applying M{\"o}bius inversion (Theorem~\ref{thm:5.3.2}(c))
	to this divisor sum gives
	\[
		c_q(n)=\sum_{d\mid q}\mu(q/d)\,d\,[\,d\mid n\,]
		=\sum_{d\mid\gcd(q,n)}d\,\mu(q/d).
	\]
	This is an integer and depends on $n$ only through $\gcd(q,n)$.
	For $n=1$, only $d=1$ contributes, so $c_q(1)=\mu(q)$.
	For $n=q$, every exponential in the defining sum is $1$, giving
	$c_q(q)=\varphi(q)$. The same argument also gives $c_q(0)=\varphi(q)$;
	neither the geometric-sum calculation nor the inversion requires $n>0$.

\end{proof}
\begin{theorem}\label{thm:5.6.2} For each fixed $n \in \mathbb{Z}$, the map $q \mapsto c_q(n)$ is multiplicative: if $\gcd(q_1, q_2) = 1$, then $c_{q_1 q_2}(n) = c_{q_1}(n)\, c_{q_2}(n)$.

\end{theorem}

\begin{proof} Let $\gcd(q_1, q_2) = 1$. Then $\gcd(q_1 q_2, n) = \gcd(q_1, n)\, \gcd(q_2, n)$ (Theorem~\ref{thm:4.6.4}), and the two factors on the right are coprime. Each divisor $d \mid \gcd(q_1 q_2, n)$ therefore factors uniquely as $d = d_1 d_2$ with $d_i \mid \gcd(q_i, n)$, and $\gcd(q_1/d_1, q_2/d_2) = 1$ yields $\mu(q_1 q_2/d) = \mu(q_1/d_1)\, \mu(q_2/d_2)$. Theorem~\ref{thm:5.6.1} now gives
	\[
		\begin{aligned}
			c_{q_1 q_2}(n) & = \sum_{d_1 \mid \gcd(q_1, n)} \sum_{d_2 \mid \gcd(q_2, n)} d_1 d_2\, \mu(q_1/d_1)\, \mu(q_2/d_2)                                                  \\
			               & = \left(\sum_{d_1 \mid \gcd(q_1, n)} d_1 \mu(q_1/d_1)\right) \left(\sum_{d_2 \mid \gcd(q_2, n)} d_2 \mu(q_2/d_2)\right) = c_{q_1}(n)\, c_{q_2}(n).
		\end{aligned}\qedhere
	\]

\end{proof}
\begin{theorem}[H{\"o}lder's formula]\label{thm:5.6.3} For $q\in\mathbb{Z}^+$ and $n\in\mathbb{Z}$, write $d = \gcd(q, n)$. Then $c_q(n) = \dfrac{\varphi(q)}{\varphi(q/d)}\, \mu\!\left(\tfrac{q}{d}\right)$.

\end{theorem}

\begin{proof} Fix $n$. The left-hand side $c_q(n)$ is multiplicative in $q$ by Theorem~\ref{thm:5.6.2}. For the right-hand side, let $\gcd(q_1, q_2) = 1$ and put $d_i = \gcd(q_i, n)$. Theorem~\ref{thm:4.6.4} gives $d := \gcd(q_1 q_2, n) = d_1 d_2$, and since $d_i \mid q_i$, we have $\gcd(d_1, d_2) = 1$ as well as $\gcd(q_1/d_1, q_2/d_2) = 1$. Multiplicativity of $\varphi$ and $\mu$ on coprime arguments yields $\varphi(q_1 q_2) = \varphi(q_1)\varphi(q_2)$, $\varphi(q_1 q_2/d) = \varphi(q_1/d_1)\varphi(q_2/d_2)$, and $\mu(q_1 q_2/d) = \mu(q_1/d_1)\mu(q_2/d_2)$, so the right-hand side $\tfrac{\varphi(q)}{\varphi(q/d)}\,\mu(q/d)$ at $q = q_1 q_2$ is the product of its values at $q_1$ and $q_2$. Both sides equal
	$1$ at $q = 1$. By Theorem~\ref{thm:5.1.10}, it is enough to
	check the identity at prime powers $q = p^{k}$ ($k \ge 1$). Let $v=v_p(n)$, with $v=+\infty$ when $n=0$, so
	$d=p^{\min\{k,v\}}$. In $c_{p^{k}}(n) = \sum_{i=0}^{\min\{k, v\}}
	p^{i}\, \mu(p^{k-i})$, only the terms with $k - i \in \{0,1\}$ survive. If $v \ge k$,
	the terms $i = k, k-1$ give $p^{k} - p^{k-1} = \varphi(p^{k})$,
	matching $\tfrac{\varphi(p^{k})}{\varphi(1)}\mu(1)$. If $v = k-1$,
	only $i = k-1$ contributes, which gives $-p^{k-1}$ and matches
	$\tfrac{\varphi(p^{k})}{\varphi(p)}\mu(p) = -p^{k-1}$. If $v \le k-2$,
	no term contributes, so $c_{p^{k}}(n) = 0$, matching $\mu(p^{k-v}) = 0$
	as $k - v \ge 2$.

\end{proof}
\subsection{Average Orders of Arithmetic Functions}

\begin{definition}[average order]\label{def:5.7.1} An arithmetic
	function $g$ is an average order of an arithmetic function $f$ if
	$\sum_{n\le x}f(n)\sim\sum_{n\le x}g(n)$ as $x\to\infty$. In the ratio definitions below, $v(x)$ is assumed nonzero for all sufficiently large $x$. Here $u(x)\sim v(x)$ means that $u(x)/v(x)\to1$ as $x\to\infty$; $u(x)=O(v(x))$ means that $|u(x)|\le C\,|v(x)|$ for a constant $C>0$ and all sufficiently large $x$; and $u(x)=o(v(x))$ means that $u(x)/v(x)\to0$ as $x\to\infty$. Unless stated otherwise, implied constants may depend on parameters fixed before the asymptotic variable is introduced, and are uniform in all remaining parameters; a constant independent of all such parameters is called absolute.

\end{definition}
\begin{theorem}\label{thm:5.7.1} Let $\gamma$ be Euler's constant. As
	$x\to\infty$:
	\begin{enumerate}[label={\alph*)}]
		\item $\displaystyle \sum_{n\le x}\frac{1}{n}=\ln x+\gamma+O\!\left(\frac{1}{x}\right)$;
		\item $\displaystyle \sum_{n>x}\frac{1}{n^{2}}=O\!\left(\frac{1}{x}\right)$.
	\end{enumerate}
	In particular, $\sum_{n\le x}\tfrac{1}{n}=\ln x+O(1)$.

\end{theorem}

\begin{proof} (a) For each integer $d\ge1$, comparison with the endpoint
	values of the integrand gives
	\[
		0<b_d:=\frac1d-\int_d^{d+1}\frac{dt}{t}
		<\frac1d-\frac1{d+1}=\frac1{d(d+1)}.
	\]
	Thus the sequence $\gamma_N:=H_N-\ln(N+1)=\sum_{d=1}^N b_d$, where
	$H_N=\sum_{d=1}^N1/d$, is increasing and bounded above by $1$.
	Let its limit be $\gamma$. The telescoping majorant gives
	$0\le\gamma-\gamma_N\le\sum_{d>N}1/[d(d+1)]=1/(N+1)$.
	Since $1/(N+1)<\ln(1+1/N)<1/N$, we obtain
	\[
		0<H_N-\ln N-\gamma
		=\ln(1+1/N)-(\gamma-\gamma_N)<\frac1N.
	\]
	In particular, $H_N-\ln N$ has the same limit $\gamma$, Euler's
	constant, and the error is $O(1/N)$. For real $x\ge1$, set
	$N=\lfloor x\rfloor$. Then $N\ge x/2$ and
	$0\le\ln x-\ln N\le1/N$, so the integer estimate yields
	$\sum_{n\le x}1/n=\ln x+\gamma+O(1/x)$.

	(b) For $n\ge2$, we have $1/n^2\le1/[n(n-1)]=1/(n-1)-1/n$.
	Thus, with $N=\lfloor x\rfloor\ge1$,
	\[
		\sum_{n>x}\frac1{n^2}
		\le\sum_{n=N+1}^{\infty}\left(\frac1{n-1}-\frac1n\right)
		=\frac1N\le\frac2x,
	\]
	which proves (b).

\end{proof}
\begin{theorem}[Dirichlet's divisor-sum estimate]\label{thm:5.7.2} $\sum_{n\le x}\tau(n)=x\ln
		x+(2\gamma-1)x+O(\sqrt{x})$, where $\gamma$ is Euler's constant. In particular, the average order of $\tau(n)$ is $\ln n$.

\end{theorem}

\begin{proof} We have $\sum_{n\le
			x}\tau(n)=\#\{(d,e)\in(\mathbb{Z}^+)^2 \mid de\le x\}$. Counting the
	pairs with $d\le\sqrt{x}$, then those with $e\le\sqrt{x}$, and
	subtracting the pairs with both $d,e\le\sqrt{x}$ (counted twice) gives
	$\sum_{n\le x}\tau(n)=2\sum_{d\le\sqrt{x}}\big\lfloor\tfrac{x}{d}\big\rfloor-\lfloor\sqrt{x}\rfloor^2$. With $\lfloor
		x/d\rfloor=x/d+O(1)$, the harmonic estimate $\sum_{d\le\sqrt{x}}\tfrac{1}{d}=\tfrac{1}{2}\ln x+\gamma+O(1/\sqrt{x})$ (Theorem~\ref{thm:5.7.1}(a) applied with $x$ replaced by $\sqrt{x}$, using $\ln\sqrt{x}=\tfrac12\ln x$), and $\lfloor\sqrt
		{x}\rfloor^2=x+O(\sqrt{x})$, the original sum equals $2x(\tfrac{1}{2}\ln
		x+\gamma)-x+O(\sqrt{x})=x\ln x+(2\gamma-1)x+O(\sqrt{x})$.

	To justify the ``in particular'' claim, compare $\tau$ with $g(n)=\ln n$. Since $\ln t$ is increasing, the left- and right-endpoint estimates give $\int_1^{N}\ln t\,dt\le\sum_{n=1}^{N}\ln n\le\int_1^{N}\ln t\,dt+\ln N$ for every integer $N\ge1$, that is, $\sum_{n=1}^{N}\ln n-\int_1^{N}\ln t\,dt=O(\ln N)$. With $\int_1^{N}\ln t\,dt=N\ln N-N+1$ and $N=\lfloor x\rfloor$, this yields $\sum_{n\le x}\ln n=x\ln x-x+O(\ln x)\sim x\ln x$. Since also $\sum_{n\le x}\tau(n)=x\ln x+(2\gamma-1)x+O(\sqrt{x})\sim x\ln x$, the two partial sums are asymptotically equal, so $\ln n$ is an average order of $\tau(n)$ in the sense of Definition~\ref{def:5.7.1}.

\end{proof}
\begin{theorem}[Basel problem]\label{thm:5.7.3} Define $\zeta(s)=\sum_{n=1}^{\infty}n^{-s}$ for real $s>1$. Then
	$\sum_{n=1}^{\infty}\frac{1}{n^2}=\zeta(2)=\frac{\pi^2}{6}$.

\end{theorem}

\begin{proof} Fix
	$m\in\mathbb{Z}^+$ and set $n=2m+1$. De Moivre's identity
	$(\cos\theta+i\sin\theta)^{n}=\cos n\theta+i\sin n\theta$ follows by
	induction on $n$ from the angle-addition formulas; combining it with
	the binomial
	theorem gives, for real $\theta$ with $\sin\theta\neq0$,
	\[
		\sin(n\theta)=\operatorname{Im}\big((\cos\theta+i\sin\theta)^n\big)=\sum_{j=0}^{m}(-1)^j\binom{n}{2j+1}\cos^{n-2j-1}\theta\,\sin^{2j+1}\theta,
	\]
	so dividing by $\sin^{n}\theta$ gives
	$\frac{\sin(n\theta)}{\sin^{n}\theta}=P(\cot^2\theta)$ with
	$P(t)=\sum_{j=0}^{m}(-1)^j\binom{n}{2j+1}t^{\,m-j}$ a polynomial of
	degree $m$. For $\theta_k=\frac{k\pi}{2m+1}$ ($1\le k\le m$),
	we have $\sin(n\theta_k)=\sin(k\pi)=0$ while $\sin\theta_k\neq0$. All $\theta_k$
	lie in $(0,\pi/2)$, where $\cot^2\theta$ is strictly decreasing, so the
	$m$ values $t_k=\cot^2\theta_k$ are distinct roots of the degree-$m$
	polynomial $P$. By Theorem~\ref{thm:3.6.2}, these are all its roots.
	Both $P(t)$ and $\binom{n}{1}\prod_{k=1}^{m}(t-t_k)$ have degree $m$,
	leading coefficient $\binom{n}{1}$, and the same $m$ distinct roots, so
	their difference has degree less than $m$ but at least $m$ roots; by
	Theorem~\ref{thm:3.6.2} it is the zero polynomial, and they coincide. Comparing coefficients of $t^{m-1}$ (Vieta's
	formulas) gives
	$\sum_{k=1}^{m}\cot^2\theta_k=\binom{n}{3}\big/\binom{n}{1}=\frac{m(2m-1)}{3}$,
	and since $\csc^2\theta=1+\cot^2\theta$,
	we get $\sum_{k=1}^{m}\csc^2\theta_k=\frac{m(2m-1)}{3}+m=\frac{2m(m+1)}{3}$.
	For $0<\theta<\frac{\pi}{2}$ we have $\sin\theta<\theta<\tan\theta$
	(the triangle with vertices $(0,0)$, $(1,0)$, $(\cos\theta,\sin\theta)$
	is contained in the circular sector of angle $\theta$, which is
	contained in the triangle with vertices $(0,0)$, $(1,0)$, $(1,\tan\theta)$, and the
	three areas are $\tfrac12\sin\theta$, $\tfrac{\theta}{2}$,
	$\tfrac12\tan\theta$),
	hence $\cot^2\theta<\frac{1}{\theta^2}<\csc^2\theta$. Applying this at
	each $\theta_k$ and summing gives
	$\frac{m(2m-1)}{3}<\frac{(2m+1)^2}{\pi^2}\sum_{k=1}^{m}\frac{1}{k^2}<\frac{2m(m+1)}{3}$,
	that is, $\frac{\pi^2}{(2m+1)^2}\cdot\frac{m(2m-1)}{3}<\sum_{k=1}^{m}\frac{1}{k^2}<\frac{\pi^2}{(2m+1)^2}\cdot\frac{2m(m+1)}{3}$.
	As $m\to\infty$, both outer terms tend to $\frac{\pi^2}{6}$, so by the
	squeeze theorem $\sum_{n=1}^{\infty}\frac{1}{n^2}=\frac{\pi^2}{6}$.

\end{proof}
\begin{altproof}[Alternative proof of Theorem~\ref{thm:2.1.2}] Fix $s>1$. For a finite set of primes $P$, expanding each absolutely convergent geometric series and multiplying gives
	\[
		\prod_{p\in P}(1-p^{-s})^{-1}=\prod_{p\in P}\sum_{j=0}^{\infty}p^{-js}=\sum_{\substack{n\ge1\\\text{every prime divisor of }n\text{ lies in }P}}\frac1{n^s},
	\]
	because the fundamental theorem of arithmetic (Theorem~\ref{thm:2.3.3}) ensures that every admissible $n$ occurs exactly once. Letting $P$ expand to include all primes is justified by absolute convergence, since $\sum_p p^{-s}\le\sum_{n=2}^{\infty}n^{-s}<\infty$, and yields Euler's product
	\[
		\zeta(s)=\sum_{n=1}^{\infty}\frac1{n^s}=\prod_{p\in\Primes}(1-p^{-s})^{-1}.
	\]
	If there were only finitely many primes $p_1,\ldots,p_r$, this would be a finite product and would tend to the finite value $\prod_{i=1}^{r}(1-p_i^{-1})^{-1}$ as $s\to1^+$. But $x^{-s}$ is decreasing, so
	\[
		\zeta(s)=1+\sum_{n=2}^{\infty}n^{-s}\ge\int_1^{\infty}x^{-s}\,dx=\frac1{s-1}\longrightarrow\infty,
	\]
	a contradiction.

\end{altproof}
\begin{theorem}\label{thm:5.7.4} For an arithmetic function $f$, write
	\[
		\mathcal{D}_f(s)=\sum_{n=1}^{\infty}\frac{f(n)}{n^s}
	\]
	for every real $s$ for which the series converges absolutely. Then:
	\begin{enumerate}[label={\alph*)}]
		\item if $\mathcal{D}_f(s)$ and $\mathcal{D}_g(s)$ converge absolutely, then $\mathcal{D}_{f*g}(s)=\mathcal{D}_f(s)\mathcal{D}_g(s)$;
		\item if $f$ is multiplicative and $\mathcal{D}_f(s)$ converges absolutely, then
		      \[
			      \mathcal{D}_f(s)=\prod_{p\in\Primes}\left(\sum_{a=0}^{\infty}\frac{f(p^a)}{p^{as}}\right);
		      \]
		\item for $s>1$,
		      \[
			      \zeta(s)=\prod_{p\in\Primes}(1-p^{-s})^{-1},
			      \qquad
			      \sum_{n=1}^{\infty}\frac{\mu(n)}{n^s}=\frac1{\zeta(s)};
		      \]
		      in particular,
		      \[
			      \sum_{n=1}^{\infty}\frac{\mu(n)}{n^2}=\frac6{\pi^2},
			      \qquad
			      \sum_{n>x}\frac{\mu(n)}{n^2}=O\!\left(\frac1x\right);
		      \]

		\item for $s>2$,
		      \[
			      \sum_{n=1}^{\infty}\frac{\varphi(n)}{n^s}=\frac{\zeta(s-1)}{\zeta(s)}.
		      \]
	\end{enumerate}

\end{theorem}

\begin{proof} Under the hypotheses of (a),
	\[
		\sum_{n\ge1}\frac{|(f*g)(n)|}{n^s}
		\le\sum_{d,e\ge1}\frac{|f(d)|\,|g(e)|}{(de)^s}
		=\left(\sum_{d\ge1}\frac{|f(d)|}{d^s}\right)
		 \left(\sum_{e\ge1}\frac{|g(e)|}{e^s}\right)<\infty.
	\]
	Thus the convolution series also converges absolutely, and the double
	series may be rearranged. Consequently
	\[
		\begin{aligned}
			\mathcal{D}_f(s)\mathcal{D}_g(s)
			 & =\sum_{d,e\ge1}\frac{f(d)g(e)}{(de)^s}
			=\sum_{n\ge1}\frac1{n^s}\sum_{de=n}f(d)g(e)                \\
			 & =\sum_{n\ge1}\frac{(f*g)(n)}{n^s}=\mathcal{D}_{f*g}(s),
		\end{aligned}
	\]
	which proves (a). For a finite set of primes $P$, expansion of the finite product gives
	\[
		\prod_{p\in P}\left(\sum_{a=0}^{\infty}\frac{f(p^a)}{p^{as}}\right)
		=\sum_{\substack{n\ge1\\\text{every prime divisor of }n\text{ lies in }P}}\frac{f(n)}{n^s},
	\]
	because Theorem~\ref{thm:5.1.10} gives
	$f(n)=\prod_{p^a\parallel n}f(p^a)$. Take $P$ to be the primes at most
	$y\ge2$, and denote this finite product by $\Pi_y(s)$. Any omitted
	integer has a prime divisor greater than $y$ and is therefore itself
	greater than $y$. Hence
	\[
		\left|\mathcal{D}_f(s)-\Pi_y(s)\right|
		\le\sum_{n>y}\frac{|f(n)|}{n^s}\longrightarrow0
		\qquad(y\to\infty).
	\]
	This proves (b) directly from absolute convergence, even if a local
	factor is zero.

	For $f=\mathbf{1}$, each local factor is the geometric series $\sum_{a\ge0}p^{-as}=(1-p^{-s})^{-1}$, which gives the Euler product for $\zeta(s)$. For $f=\mu$, Definition~\ref{def:5.3.6} gives the local factor
	\[
		\sum_{a=0}^{\infty}\frac{\mu(p^a)}{p^{as}}=1-p^{-s};
	\]
	hence its Euler product is $\prod_p(1-p^{-s})=1/\zeta(s)$. Both series converge absolutely for $s>1$ because $|\mu(n)|\le1$.

	For $f=\varphi$, Theorem~\ref{thm:3.4.2} and $\varphi(p^a)=p^a-p^{a-1}$ for $a\ge1$ give, when $s>2$,
	\[
		\begin{aligned}
			\sum_{a=0}^{\infty}\frac{\varphi(p^a)}{p^{as}}
			 & =1+(1-p^{-1})\sum_{a=1}^{\infty}p^{a(1-s)} \\
			 & =1+(1-p^{-1})\frac{p^{1-s}}{1-p^{1-s}}
			=\frac{1-p^{-s}}{1-p^{1-s}}.
		\end{aligned}
	\]
	The Dirichlet series converges absolutely because $0\le\varphi(n)\le n$, so (b) yields
	\[
		\mathcal{D}_{\varphi}(s)=\prod_p\frac{1-p^{-s}}{1-p^{1-s}}
		=\frac{\zeta(s-1)}{\zeta(s)}.
	\]
	Multiplying by $\zeta(s)$ and using (a) gives
	\[
		\mathcal{D}_{\varphi*\mathbf{1}}(s)=\mathcal{D}_{\id}(s),
	\]
	since $\mathcal{D}_{\id}(s)=\sum n^{1-s}=\zeta(s-1)$. To justify coefficient
	comparison, suppose $\sum a_n n^{-s}$ and $\sum b_n n^{-s}$ converge
	absolutely and agree for every sufficiently large real $s$. If their
	coefficients differ, let $N$ be the first index at which they differ. Subtracting
	the series and multiplying by $N^s$ gives
	\[
		0=a_N-b_N+\sum_{n>N}(a_n-b_n)(N/n)^s.
	\]
	Fix a sufficiently large $s_0$ and put
	$C=N^{s_0}\sum_{n>N}|a_n-b_n|n^{-s_0}<\infty$. For $s\ge s_0$,
	\[
		\left|\sum_{n>N}(a_n-b_n)(N/n)^s\right|
		\le C\left(\frac{N}{N+1}\right)^{s-s_0}\longrightarrow0,
	\]
	because $N/n\le N/(N+1)<1$ for every $n>N$. This contradicts
	$a_N-b_N\neq0$, proving coefficient uniqueness. Therefore
	$(\varphi*\mathbf{1})(n)=\id(n)=n$, recovering
	$\sum_{d\mid n}\varphi(d)=n$.

	Finally, take $s=2$ in the M{\"o}bius identity and use $\zeta(2)=\pi^2/6$ (Theorem~\ref{thm:5.7.3}). The tail estimate follows from
	\[
		\left|\sum_{n>x}\frac{\mu(n)}{n^2}\right|\le\sum_{n>x}\frac1{n^2}=O\!\left(\frac1x\right)
	\]
	by Theorem~\ref{thm:5.7.1}(b).

\end{proof}

\begin{theorem}\label{thm:5.7.5}
	$\sum_{n\le x}\varphi(n)=\dfrac{3}{\pi^2}x^2+O(x\ln x)$; in particular, the average order of $\varphi(n)$ is $\dfrac{6}{\pi^{2}}n$.

\end{theorem}

\begin{proof} For real $y\ge1$, put $N=\lfloor y\rfloor$. Since
	$y-1\le N\le y$,
	\[
		\frac{y^2-y}{2}\le\sum_{e\le y}e=\frac{N(N+1)}{2}
		\le\frac{y^2+y}{2}.
	\]
	Hence $\sum_{e\le y}e=y^2/2+O(y)$ with an absolute implied constant,
	uniformly for $y\ge1$; in particular, this estimate applies to every $y=x/d$ with
	$1\le d\le x$.
	M{\"o}bius inversion (Theorem~\ref{thm:5.3.2}(c)) applied to
	$\sum_{d\mid n}\varphi(d)=n$ (Theorem~\ref{thm:5.1.5}) gives
	$\varphi(n)=\sum_{d\mid n}\mu(d)n/d$. Writing $n=de$ and using
	$|\mu(d)|\le1$, we obtain
	\[
		\begin{aligned}
		\sum_{n\le x}\varphi(n)
		&=\sum_{d\le x}\mu(d)\sum_{e\le x/d}e\\
		&=\frac{x^2}{2}\sum_{d\le x}\frac{\mu(d)}{d^2}
		  +O\!\left(x\sum_{d\le x}\frac1d\right).
		\end{aligned}
	\]
	Theorem~\ref{thm:5.7.1}(a) shows that the error is $O(x\ln x)$, and
	Theorem~\ref{thm:5.7.4} gives
	$\sum_{d\le x}\mu(d)/d^2=6/\pi^2+O(1/x)$. Multiplying this last estimate
	by $x^2/2$ contributes $3x^2/\pi^2+O(x)$. Combining the errors gives
	$\sum_{n\le x}\varphi(n)=3x^2/\pi^2+O(x\ln x)$.
	Finally,
	$\sum_{n\le x}(6/\pi^2)n=3x^2/\pi^2+O(x)\sim3x^2/\pi^2$,
	so $(6/\pi^2)n$ is an average order of $\varphi$ in the sense of
	Definition~\ref{def:5.7.1}.

\end{proof}
\begin{theorem}\label{thm:5.7.6} $\sum_{n\le x}\sigma(n)=\dfrac{\pi^2}{12}x^2+O(x\ln x)$;
	in particular, the average order of $\sigma(n)$ is $\tfrac{\pi^2}{6} n$.

\end{theorem}

\begin{proof} Use the uniform estimate
	$\sum_{d\le y}d=y^2/2+O(y)$ for $y\ge1$, established in the preceding
	proof. Writing $n=de$ with $d\mid n$ gives
	\[
		\begin{aligned}
		\sum_{n\le x}\sigma(n)
		&=\sum_{e\le x}\sum_{d\le x/e}d\\
		&=\frac{x^2}{2}\sum_{e\le x}\frac1{e^2}
		  +O\!\left(x\sum_{e\le x}\frac1e\right).
		\end{aligned}
	\]
	Theorem~\ref{thm:5.7.1} gives
	$\sum_{e\le x}1/e=O(\ln x)$ and
	$\sum_{e\le x}1/e^2=\zeta(2)+O(1/x)$. Since
	$\zeta(2)=\pi^2/6$ (Theorem~\ref{thm:5.7.3}), the main term is
	$\pi^2x^2/12+O(x)$, while the remaining error is $O(x\ln x)$. Hence
	$\sum_{n\le x}\sigma(n)=\pi^2x^2/12+O(x\ln x)$.
	Also, $\sum_{n\le x}(\pi^2/6)n=\pi^2x^2/12+O(x)$ has the same
	leading term, so $(\pi^2/6)n$ is an average order of $\sigma$
	(Definition~\ref{def:5.7.1}).

\end{proof}
\begin{theorem}\label{thm:5.7.7} Let $Q(x)=\#\{n\le x \mid n\text{ squarefree}\}=\sum_{n\le
			x}\mu^2(n)$. Then $Q(x)=\dfrac{6}{\pi^2}x+O(\sqrt{x})$; in particular,
	the squarefree integers have natural density $1/\zeta(2)=6/\pi^2$
	(a set $A\subseteq\mathbb{Z}^+$ has natural density $\delta$ if
	$\#\{n\le x \mid n\in A\}/x\to\delta$ as $x\to\infty$).

\end{theorem}

\begin{proof} First we show that $\mu^2(n)=\sum_{d^2\mid n}\mu(d)$.
	Define $h(k)=\mu(\sqrt{k})$ when $k$ is a perfect square and $h(k)=0$
	otherwise. A product of coprime positive integers is a square if
	and only if each factor is a square, so the multiplicativity of $\mu$
	shows that $h$ is multiplicative. Hence
	\[
		\sum_{d^2\mid n}\mu(d)=(\mathbf{1}*h)(n)
	\]
	is multiplicative by Theorem~\ref{thm:5.3.5}; $\mu^2$ is multiplicative as well.
	At a prime power $p^{a}$, the displayed sum equals $\mu(1)+[a\ge2]\,\mu(p)$,
	which is $1$ when $a\le1$ and $0$ when $a\ge2$, exactly
	$\mu^2(p^{a})$. The two multiplicative functions agree at $1$ and at
	every prime power, so Theorem~\ref{thm:5.1.10} shows that the identity
	holds for all positive integers.
	Regrouping the double
	sum by $d$ (each pair $(d,n)$ satisfies $d^2\mid n\le x$, so
	$d\le\sqrt{x}$) gives $Q(x)=\sum_{n\le x}\sum_{d^2\mid
			n}\mu(d)=\sum_{d\le\sqrt{x}}\mu(d)\big\lfloor x/d^2\big\rfloor$. Using
	$\lfloor x/d^2\rfloor=x/d^2+O(1)$, we obtain $Q(x)=x\sum_{d\le\sqrt
			{x}}\tfrac{\mu(d)}{d^2}+O(\sqrt
		{x})=x\sum_{d\ge1}\tfrac{\mu(d)}{d^2}+O\big(x\sum_{d>\sqrt
			{x}}\tfrac{1}{d^2}\big)+O(\sqrt{x})$. By Theorem~\ref{thm:5.7.4}, $\sum_{d\ge1}\mu(d)/d^2=1/\zeta(2)$; and, by Theorem~\ref{thm:5.7.1}(b), the tail estimate
	$\sum_{d>\sqrt{x}}1/d^2=O(1/\sqrt{x})$ shows that the middle error term is
	$O(x\cdot 1/\sqrt{x})=O(\sqrt{x})$, so the resulting expression is $x/\zeta(2)+O(\sqrt{x})$. Finally,
	$\zeta(2)=\pi^2/6$ (the Basel problem, Theorem~\ref{thm:5.7.3}) gives
	$Q(x)=\tfrac{6}{\pi^2}x+O(\sqrt{x})$.

\end{proof}
\newpage

\section{Orders, Primitive Roots, and Multiplicative Structure Modulo \texorpdfstring{$m$}{m}}

\subsection{The Order of an Element; Period Lengths of Base-\texorpdfstring{$m$}{m} Expansions}\label{subsec:6.1}

\begin{definition}[order]\label{def:6.1.1} Let $n\in\mathbb{Z}^+$ and
	$a\in\mathbb{Z}$ with $\gcd(a,n)=1$. The order of $a$ modulo $n$,
	denoted $\delta_n(a)$, is the smallest positive integer $r$ such that
	$a^r\equiv1\pmod n$. Such an $r$ exists by the Fermat--Euler theorem
	(Theorem~\ref{thm:3.3.1}). The order depends only on the residue class of
	$a$, since congruent integers have congruent positive powers.
	It is the order of this class in the multiplicative group $\ZmodUnits{n}$.
	Modular powers of a unit are understood in this group: its zeroth power
	is the identity, and negative powers use the inverse. In particular,
	$\delta_1(a)=1$ for every integer $a$; the sole class modulo $1$ is the identity.

\end{definition}
\begin{theorem}\label{thm:6.1.1} Let $m\in\mathbb{Z}^+$ and let $a,b,n\in\mathbb{Z}$ with
	$\gcd(a,m)=\gcd(b,m)=1$. Powers are taken in the unit group
	$\ZmodUnits{m}$ as in Definition~\ref{def:6.1.1}; in particular,
	negative exponents use inverses. Then:

	\begin{enumerate}[label={\alph*)}]
		\item $a^{n} \equiv a^{\,n \bmod \delta_m(a)} \pmod{m}$.
		\item $a^{0}, a^{1}, \ldots, a^{\delta_m(a) - 1}$ are pairwise
		      distinct modulo $m$.
		\item $a^{n} \equiv 1 \pmod{m}$ if and only if $\delta_m(a) \mid n$.
		\item $\delta_m(a^{n}) = \dfrac{\delta_m(a)}{\gcd(\delta_m(a), n)}$; for $n=0$ the formula gives $1$, since $\delta_m(1)=1$ and $\gcd(\delta_m(a),0)=\delta_m(a)$.
		\item $\dfrac{\lcm(\delta_m(a), \delta_m(b))}{\gcd(\delta_m(a), \delta_m(b))} \mid \delta_m(ab) \mid
			      \lcm(\delta_m(a), \delta_m(b))$.
		\item $\delta_m(ab) = \delta_m(a)\,\delta_m(b)$ if and only if
		      $\gcd(\delta_m(a), \delta_m(b)) = 1$.
		\item There always exists $c \in \mathbb{Z}^+$ with $\gcd(c, m) =
			      1$ such that $\delta_m(c) = \lcm(\delta_m(a), \delta_m(b))$.
		\item For $m_1, m_2 \in \mathbb{Z}^+$ with $\gcd(m_1, m_2) = 1$
		      and $\gcd(a, m_1m_2) = 1$,
		      \[
			      \delta_{m_1m_2}(a)
			      =\lcm(\delta_{m_1}(a),\delta_{m_2}(a)).
		      \]
		\item If $m = \prod_{i=1}^{k} p_i^{\alpha_i}$ is the prime factorization
		      of $m$ and $\gcd(a,m)=1$, then
		      $\delta_m(a)=\lcm_{i=1}^{k}\big(\delta_{p_i^{\alpha_i}}(a)\big)$.
		      For $m=1$, the empty lcm is taken to be $1$.
	\end{enumerate}

	Parts (a)--(d) hold for any element of finite order in a group,
	and (e)--(g) hold for any two commuting elements of finite order, with
	$\delta_m(\cdot)$ interpreted as the order of the element. The statement above
	is the special case of the group $\ZmodUnits{m}$.

\end{theorem}

\begin{proof}

	\begin{enumerate}[label={\alph*)}]
		\item By the division algorithm, write $n=\delta_m(a)l+r$ with
		      $l\in\mathbb{Z}$ and $0\le r<\delta_m(a)$. In the unit group,
		      $(a^{\delta_m(a)})^l=1$ even when $l<0$, so
		      $a^n\equiv a^r=a^{\,n\bmod\delta_m(a)}\pmod m$.
		      If $\delta_m(a)\mid n$, then $r=0$ and the right-hand side is the identity.
		\item If $a^{i} \equiv a^{j} \pmod{m}$ with $0 \le i < j \le
			      \delta_m(a) - 1$, then $a^{j - i} \equiv 1$ with $0 < j - i <
			      \delta_m(a)$, contradicting minimality.
		\item ($\Leftarrow$) If $\delta_m(a)\mid n$, then $a^n\equiv1$.
		      ($\Rightarrow$) If $a^n\equiv1$, write $n=q\delta_m(a)+r$; by (a),
		      $a^r\equiv1$, so $r=0$.
		\item For $t\in\mathbb{Z}^+$, we have $(a^n)^t \equiv 1$ if and only if $\delta_m(a) \mid nt$, that is, if and only if
			      $\frac{\delta_m(a)}{\gcd(\delta_m(a), n)} \mid
			      \frac{n}{\gcd(\delta_m(a), n)} t$; since the two quotients are
		      coprime, Euclid's lemma (Theorem~\ref{thm:2.2.8}) gives $\frac{\delta_m(a)}{\gcd(\delta_m(a), n)}
			      \mid t$, so the least positive such $t$ is $\frac{\delta_m(a)}{\gcd(\delta_m(a), n)}$.
		\item Write $L = \lcm(\delta_m(a), \delta_m(b))$, $G
			      = \gcd(\delta_m(a), \delta_m(b))$. First, $(ab)^{L} \equiv 1$, so
		      $\delta_m(ab) \mid L$. Next, since $(ab)^{\delta_m(ab)} \equiv 1$ and $b^{\delta_m(b)} \equiv 1$, we have
		      $(ab)^{\delta_m(ab)\delta_m(b)} = a^{\delta_m(ab)\delta_m(b)}\,b^{\delta_m(ab)\delta_m(b)} \equiv
			      a^{\delta_m(ab)\delta_m(b)} \equiv 1$ (the $b$-factor is $1$ because $\delta_m(b)$ divides the exponent, and the left-hand side is $1$ because $\delta_m(ab)$ divides it), so $\delta_m(a) \mid
			      \delta_m(ab)\delta_m(b)$. Dividing by $G$ and using coprimality gives
		      $\frac{\delta_m(a)}{G} \mid \delta_m(ab)$, and symmetrically
		      $\frac{\delta_m(b)}{G} \mid \delta_m(ab)$. As
		      $\gcd\big(\tfrac{\delta_m(a)}{G},\tfrac{\delta_m(b)}{G}\big)=1$
		      (dividing two integers by their gcd always leaves coprime quotients),
		      their product satisfies $\frac{L}{G} \mid \delta_m(ab)$.
		\item ($\Leftarrow$) If $G = 1$, then $\frac{L}{G} = L =
			      \delta_m(a)\delta_m(b)$, so (e) forces $\delta_m(ab) =
			      \delta_m(a)\delta_m(b)$. ($\Rightarrow$) If $\delta_m(ab) =
			      \delta_m(a)\delta_m(b)$, then by the second divisibility relation in (e), $\delta_m(a)\delta_m(b)=\delta_m(ab)\mid L=\frac{\delta_m(a)\delta_m(b)}{G}$, forcing $G=1$.
		\item By Theorem~\ref{thm:2.3.5}, write $\lcm(\delta_m(a), \delta_m(b))=AB$ with $A\mid\delta_m(a)$,
		      $B\mid\delta_m(b)$, and $\gcd(A, B)=1$.
		      Let $a' = a^{\delta_m(a)/A}$, $b' = b^{\delta_m(b)/B}$; by (d),
		      $\delta_m(a') = A$, $\delta_m(b') = B$, and by (f), $\delta_m(a'b') = AB$.
		      Choose $c\in\{1,\ldots,m\}$ representing the unit class of $a'b'$.
		      Then $\gcd(c,m)=1$ and $\delta_m(c)=AB$, including in the case $m=1$.
		\item By the Chinese remainder theorem (Theorem~\ref{thm:3.4.1}), $a^{n} \equiv 1 \pmod{m_1 m_2}$ if and only if $a^{n} \equiv 1
			      \pmod{m_1}$ and $a^n\equiv1\pmod{m_2}$, that is, if and only if $\delta_{m_1}(a) \mid n$ and
		      $\delta_{m_2}(a) \mid n$, or equivalently $\lcm(\delta_{m_1}(a), \delta_{m_2}(a)) \mid n$; the least positive such $n$ gives the claim.
		\item By the Chinese remainder theorem and (c), $a^{n} \equiv 1 \pmod{m}$ if and only if
			      $\delta_{p_i^{\alpha_i}}(a) \mid n$ for all $i$, or equivalently
			      $\lcm_i\big(\delta_{p_i^{\alpha_i}}(a)\big) \mid n$; the
		      least positive such $n$ is $\delta_m(a)$.
	\end{enumerate}

\end{proof}
\begin{theorem}\label{thm:6.1.2} Let $G$ be a
	group with identity element $e$, let $x\in G$, let $q$ be a prime, and let
	$t\in\mathbb{Z}^+$. If $x^{q^{t}}=e$ and $x^{q^{t-1}}\neq e$, then $x$ has
	order exactly $q^{t}$. In particular, in the group $\ZmodUnits{m}$, if
	$\gcd(a,m)=1$, $a^{q^{t}}\equiv1\pmod{m}$, and
	$a^{q^{t-1}}\not\equiv1\pmod{m}$, then $\delta_m(a)=q^{t}$.

\end{theorem}

\begin{proof} Since $x^{q^{t}}=e$, the element $x$ has finite
	order; let $d$ be that order, the least positive integer with $x^{d}=e$.
	Division with remainder gives $q^{t}=d\ell+r$ with $0\le r<d$, and then
	$x^{r}=x^{q^{t}}\big(x^{d}\big)^{-\ell}=e$, so minimality of $d$ forces $r=0$;
	that is, $d\mid q^{t}$. Since $q$ is prime, unique factorization gives
	$d=q^{s}$ with $0\le s\le t$. If $s\le t-1$, then $q^{s}\mid q^{t-1}$, so
	$x^{q^{t-1}}=\big(x^{q^{s}}\big)^{q^{t-1-s}}=e$, contradicting the hypothesis.
	Thus $s=t$ and $d=q^{t}$. The claim for $\ZmodUnits{m}$ is
	the case $x=a$, $e=1$, since $\delta_m(a)$ is the order of
	$a$ in that group (Definition~\ref{def:6.1.1}).

\end{proof}
\begin{theorem}\label{thm:6.1.3} Let $m,n\ge2$ be integers with
	$\gcd(m,n)=1$. For every integer $a$ with $\gcd(a,n)=1$, the minimal
	period length of the base-$m$ expansion of $a/n$ equals $\delta_n(m)$.

\end{theorem}

\begin{proof} For negative $a$, first separate the sign, as in
	Definition~\ref{def:4.4.1}, and replace $a$ by $|a|$. Then replace this
	positive numerator by its remainder modulo $n$, discarding the integer
	part. Neither operation changes the period of the digit sequence under
	consideration, and coprimality ensures that the remainder is nonzero.
	Thus we may assume $0<a<n$. By
	Theorem~\ref{thm:4.4.1}, a positive integer $P$ is a period precisely when
	$r_P=r_0$, equivalently $m^Pa\equiv a\pmod{n}$. Since $a$ is a unit
	modulo $n$, this is equivalent to $m^P\equiv1\pmod{n}$, which by Theorem~\ref{thm:6.1.1}(c) is equivalent to $\delta_n(m)\mid P$. The least positive period
	is therefore $\delta_n(m)$.

\end{proof}
\begin{theorem}\label{thm:6.1.4} Let $m\ge2$ and $n\ge1$ be integers.
	Write $m=\prod_{i=1}^{r}p_i^{\alpha_i}$ for the prime factorization of the base,
	and write
	\[
		n=n'\prod_{i=1}^{r}p_i^{\beta_i},\qquad
		n'\in\mathbb{Z}^+,\quad \beta_i\in\mathbb{Z}_{\ge0},\quad \gcd(m,n')=1.
	\]
	Put $L=\max_{1\le i\le r}\left\lceil\beta_i/\alpha_i\right\rceil$ (with $L=0$ when every $\beta_i=0$), and put $P=\delta_{n'}(m)$, with $P=1$ if $n'=1$. Then the base-$m$ expansion of every $a/n$ with $a\in\mathbb{Z}$ and
	$0\le a<n$ satisfies $d_{k+P}=d_k$ for every integer $k>L$.

\end{theorem}

\begin{proof} Since $L\alpha_i\ge\beta_i$ for every $i$, the integer
	\[
		c=\prod_{i=1}^{r}p_i^{L\alpha_i-\beta_i}
	\]
	satisfies $m^L/n=c/n'$. Thus $m^La/n=ca/n'$. Write $ca=Qn'+R$ with $0\le R<n'$. After $L$ long-division steps, the remaining fractional part equals both $r_L/n$ and $R/n'$, hence
	\[
		r_L=R\prod_{i=1}^{r}p_i^{\beta_i}.
	\]
	Let $(\widetilde d_j,\widetilde r_j)$ be the digit--remainder recursion in base $m$ for $R/n'$, and write $B=\prod_{i=1}^{r}p_i^{\beta_i}$, so $n=Bn'$. Induction on $j$ gives
	\[
		r_{L+j}=B\,\widetilde r_j,\qquad d_{L+j}=\widetilde d_j\quad(j\ge1):
	\]
	the case $j=0$ is the identity $r_L=B\widetilde r_0$ proved above, and if
	$r_{L+j-1}=B\widetilde r_{j-1}$, then
	$m r_{L+j-1}=B m\widetilde r_{j-1}=B(n'\widetilde d_j+\widetilde r_j)=n\widetilde d_j+B\widetilde r_j$
	with $0\le B\widetilde r_j<Bn'=n$, so uniqueness of quotient and
	remainder upon division by $n$ forces $d_{L+j}=\widetilde d_j$ and
	$r_{L+j}=B\widetilde r_j$.
	Because $\gcd(m,n')=1$ and $m^P\equiv1\pmod{n'}$, multiplication by $m^P$ fixes every residue modulo $n'$. Hence $\widetilde r_{j+P}=\widetilde r_j$ and $\widetilde d_{j+P}=\widetilde d_j$. Therefore $d_{k+P}=d_k$ for every $k>L$.

\end{proof}
\begin{theorem}[Midy's theorem in base $m$]\label{thm:6.1.5} Let $m\ge2$
	be an integer, let $p$ be a prime with $p\nmid m$, and suppose the
	base-$m$ expansion of $1/p$ has least period length $2k$, where
	$k\in\mathbb{Z}^+$. Split this repeating block into two $k$-digit blocks
	with base-$m$ values $A$ and $B$, retaining any leading zeros.
	Then $A+B=m^k-1$.

\end{theorem}

\begin{proof} By Theorem~\ref{thm:6.1.3}, the minimal period length of the base-$m$ expansion of $1/p$ equals $\delta_p(m)$, so the hypothesis gives $\delta_p(m)=2k$. Hence $m^{2k}\equiv1\pmod{p}$ but $m^k\not\equiv1\pmod{p}$. The factorization $(m^k-1)(m^k+1)\equiv0\pmod{p}$ therefore gives $m^k\equiv-1\pmod{p}$. The complete block has value $Am^k+B=(m^{2k}-1)/p$. Put $q=(m^k+1)/p$. Then
	\[
		Am^k+B=(m^k-1)q=(q-1)m^k+(m^k-q).
	\]
	Because $p\ge2$ and $p\nmid m$, we have $1\le q<m^k$; both $q-1$ and $m^k-q$ therefore lie in $[0,m^k)$. Uniqueness of quotient and remainder upon division by $m^k$ yields $A=q-1$ and $B=m^k-q$, so $A+B=m^k-1$.

\end{proof}
\subsection{Primitive Roots}

\begin{definition}[primitive root]\label{def:6.2.1} For $m \in
		\mathbb{Z}^+$ and an integer $g$ with $\gcd(g, m) = 1$, if $\delta_m(g) =
		\varphi(m)$, we call $g$ a primitive root modulo $m$.

\end{definition}
\begin{theorem}\label{thm:6.2.1} Let $m \in \mathbb{Z}^+$ and
	suppose that $\gcd(g, m) = 1$. Then $g$ is a primitive root modulo $m$ if and only
	if the $\varphi(m)$ terms $g^{0}, g^{1}, \ldots, g^{\varphi(m) - 1}$
	form a reduced residue system modulo $m$.

\end{theorem}

\begin{proof} ($\Rightarrow$) If $\delta_m(g) =
		\varphi(m)$, Theorem~\ref{thm:6.1.1}(b) shows that the powers $g^{0}, \ldots,
		g^{\varphi(m) - 1}$ are pairwise incongruent and coprime to $m$. Since there are
	$\varphi(m)$ of them, they represent every unit class exactly once and
	hence form a reduced residue system (Definition~\ref{def:3.2.4}).
	($\Leftarrow$) Suppose they form a reduced residue system and put $D =
		\delta_m(g)$. If $D < \varphi(m)$, then $g^{D} \equiv g^{0} \pmod{m}$
	with $0 \le D \le \varphi(m) - 1$, contradicting pairwise
	incongruence; hence $D \ge \varphi(m)$. The Fermat--Euler theorem (Theorem~\ref{thm:3.3.1}) gives
	$g^{\varphi(m)} \equiv 1 \pmod{m}$, so $D \mid \varphi(m)$ by Theorem~\ref{thm:6.1.1}(c) and therefore $D \le \varphi(m)$. Hence $\delta_m(g) =
		\varphi(m)$.

\end{proof}
\begin{theorem}\label{thm:6.2.2} If there exists a primitive root
	modulo $n$, then there are exactly $\varphi(\varphi(n))$ pairwise
	incongruent primitive roots modulo $n$.

\end{theorem}

\begin{proof} Let $g$ be a primitive root. By Theorem~\ref{thm:6.2.1}, every residue coprime to $n$ is congruent to exactly one of
	$g^{0}, \ldots, g^{\varphi(n) - 1}$, so every primitive root modulo
	$n$ is congruent to $g^{k}$ for some $0 \le k \le \varphi(n) - 1$. By
	Theorem~\ref{thm:6.1.1}(d), $\delta_n(g^{k}) =
		\frac{\varphi(n)}{\gcd(\varphi(n), k)}$,
	so $g^{k}$ is a primitive root if and only if $\gcd(\varphi(n), k) =
		1$. The number of such $k$ with $0 \le k \le \varphi(n) - 1$ is
	$\varphi(\varphi(n))$.

\end{proof}
\begin{altproof}[Alternative proof of Theorem~\ref{thm:6.2.2}]
	With $g$ a primitive root, $\ZmodUnits{n}$ is cyclic of order
	$\varphi(n)$, generated by $g$. For $d \mid \varphi(n)$, let $N(d)$
	be the number of residue classes modulo $n$ of order $d$. Since
	$\ZmodUnits{n}$ is cyclic, $x^{d}
		\equiv 1$ has exactly $d$ solutions, namely $\{g^{i\varphi(n)/d} \mid
		0\le i<d\}$; by Theorem~\ref{thm:6.1.1}(c) these are exactly the classes whose
	order divides $d$, so $\sum_{e \mid d} N(e) = d$. Comparing with
	$\sum_{e \mid d}\varphi(e) = d$ and applying M{\"o}bius inversion gives
	$N(d) = \varphi(d)$. Taking $d = \varphi(n)$ gives $N(\varphi(n)) =
		\varphi(\varphi(n))$.

\end{altproof}
\begin{theorem}[Lucas's primitive root test]\label{thm:6.2.3} For $n
		\in \mathbb{Z}^+$ with $\gcd(n, g) = 1$, the integer $g$ is a primitive root
	modulo $n$ if and only if $g^{\varphi(n)/p} \not\equiv 1 \pmod{n}$ for
	every prime $p \mid \varphi(n)$.

\end{theorem}

\begin{proof} ($\Rightarrow$) Let $\delta_n(g) =
		\varphi(n)$ and let $p \mid \varphi(n)$ be prime. Then
	$\frac{\varphi(n)}{p} \in \mathbb{Z}^+$ and $\frac{\varphi(n)}{p} <
		\varphi(n) = \delta_n(g)$, so $\delta_n(g) \nmid \frac{\varphi(n)}{p}$
	and Theorem~\ref{thm:6.1.1}(c) gives $g^{\varphi(n)/p} \not\equiv 1$. ($\Leftarrow$) The Fermat--Euler theorem gives $g^{\varphi(n)}\equiv1\pmod{n}$, so Theorem~\ref{thm:6.1.1}(c) gives $\delta_n(g)\mid\varphi(n)$. If $\delta_n(g) < \varphi(n)$, write $\varphi(n) = \delta_n(g)\,t$ with $t > 1$ and choose a prime $p \mid t$; then
	$\delta_n(g) \mid \frac{\varphi(n)}{p}$, so $g^{\varphi(n)/p}
		\equiv 1$ by Theorem~\ref{thm:6.1.1}(c), a contradiction.

\end{proof}
\begin{theorem}\label{thm:6.2.4} Let $e\ge3$, $k$, and $j$ be integers
	with $3\le j\le 2^{e-2}$. Then
	\[
		v_2\Big(\tbinom{2^{e-2}}{j}(2k)^{j}\Big)\ge e.
	\]

\end{theorem}

\begin{proof} When $e=3$, there is no $j$ in the stated range.
	If $k=0$, the expression inside $v_2$ is zero and its valuation is
	$+\infty$, so the claim holds. Assume $k\neq0$.
	Since $(2k)^j=2^jk^j$, we have $v_2\big((2k)^j\big)\ge j$.
	For the binomial coefficient, Lucas's theorem (Theorem~\ref{thm:4.1.2}) already shows that $\binom{2^{e-2}}{j}$ is even for $0 < j < 2^{e-2}$: in base
	$2$, the number $2^{e-2}$ is a $1$ followed by $e-2$ zeros, so
	$\binom{2^{e-2}}{j} \equiv \prod_{i=0}^{e-2}\binom{(2^{e-2})_i}{j_i} \pmod{2}$
	vanishes unless $j \in \{0, 2^{e-2}\}$. Kummer's theorem (Theorem~\ref{thm:4.7.3}) refines this to the exact value
	\[
		v_2\binom{2^{e-2}}{j} = S_2(j) + S_2(2^{e-2}-j) - S_2(2^{e-2}) = (e-2) - v_2(j),
	\]
	where the last equality holds trivially when $j=2^{e-2}$ (both sides are $0$),
	and holds for $3\le j<2^{e-2}$ because adding $j$ and $2^{e-2}-j$ (whose sum is
	the single-bit number $2^{e-2}$) produces exactly $(e-2)-v_2(j)$ carries in
	base $2$: with $t=v_2(j)$, which satisfies $t\le e-3$ in this range,
	positions $0,\ldots,t-1$ of both summands are $0$ and generate no carry;
	at position $t$ both bits are $1$ (since $2^{e-2}-j\equiv-j\equiv2^{t}
	\pmod{2^{t+1}}$), generating a carry; and at each of the $e-3-t$ positions
	$t+1,\ldots,e-3$ the two bits are complementary, so the incoming carry is
	passed on, one further carry per position, for a total of
	$1+(e-3-t)=(e-2)-v_2(j)$ carries (the last carry, into position $e-2$,
	produces the leading $1$ of $2^{e-2}$). Hence
	\[
		v_2\big(\binom{2^{e-2}}{j}(2k)^{j}\big) \ge (e-2) - v_2(j) + j = e +
		\big(j - v_2(j) - 2\big).
	\]

	Finally, $j - v_2(j) \ge 2$ for every $j \ge 3$: if $v_2(j) = 0$ then
	$j - v_2(j) = j \ge 3$; if $v_2(j) \ge 1$ then $j \ge 2^{v_2(j)} \ge
		2\,v_2(j)$, so $j - v_2(j) \ge v_2(j) \ge 1$, and $j - v_2(j) = 1$
	would force $v_2(j) = 1,\ j = 2$, excluded by $j \ge 3$. Therefore
	$v_2\big(\binom{2^{e-2}}{j}(2k)^{j}\big) \ge e$.

\end{proof}
\begin{theorem}\label{thm:6.2.5} Let $p$ be an odd prime and let $g$ be a
	primitive root modulo $p$ with $v_p(g^{p-1}-1)=1$. Then $g$ is a primitive
	root modulo $p^{e}$ for every integer $e\ge1$.

\end{theorem}

\begin{proof} Since $g^{\varphi(p^e)}\equiv1\pmod{p^e}$ by the
	Fermat--Euler theorem, Theorem~\ref{thm:6.1.1}(c) gives
	$\delta_{p^e}(g)\mid\varphi(p^e)=p^{e-1}(p-1)$; and
	$g^{\delta_{p^e}(g)}\equiv1\pmod{p^e}$ implies
	$g^{\delta_{p^e}(g)}\equiv1\pmod{p}$, so
	$(p-1)=\delta_p(g)\mid\delta_{p^e}(g)$. Every divisor of $p^{e-1}(p-1)$ that
	is a multiple of $p-1$ has the form $p^{k}(p-1)$, so
	$\delta_{p^e}(g)=p^{k}(p-1)$ for some $0\le k\le e-1$. The lifting-the-exponent lemma applies to the pair $g^{p-1},1$: Fermat's little
	theorem (Theorem~\ref{thm:3.3.2}) gives $p\mid(g^{p-1}-1)$, we have
	$p\nmid g$, and
	$v_p(g^{p-1}-1)=1$, which in particular gives $g^{p-1}\neq1$. Theorem~\ref{thm:4.8.3} then gives $v_p\big(g^{p^{k}(p-1)}-1\big)=1+k$, so
	$g^{p^{k}(p-1)}\equiv1\pmod{p^e}$ if and only if $k+1\ge e$; the least such
	$k$ is $e-1$, giving $\delta_{p^e}(g)=p^{e-1}(p-1)=\varphi(p^e)$.

\end{proof}
\begin{theorem}[primitive root theorem]\label{thm:6.2.6}
	For $n \in \mathbb{Z}^+$, a primitive root modulo $n$ exists if and
	only if $n = 1, 2, 4,
		p^{e}$, or $2 p^{e}$, where $p$ is an odd prime and $e \in \mathbb{Z}^+$.

\end{theorem}

\begin{proof}
	We begin with Case 1, where $n = 1, 2, 4$. The integer $1$ is a primitive root
	modulo $1$ and modulo $2$, while $3$ is a primitive root modulo $4$.

	In Case 2, let $n = p^{e}$ or $2p^{e}$, where $p$ is odd. We proceed in four steps.

	The first step proves that a primitive root modulo $p$ exists. Let $N(d)$ be the number
	of residue classes modulo $p$ of order $d$. By the polynomial root bound (Theorem~\ref{thm:3.6.2}), $x^d\equiv1\pmod{p}$ has at most $d$ roots; if $N(d) > 0$, an element $a$ of
	order $d$ gives, by Theorem~\ref{thm:6.1.1}(b), $d$ pairwise incongruent roots $a^{0}, \ldots, a^{d-1}$ of $x^{d}
		\equiv 1$, among which, by Theorem~\ref{thm:6.1.1}(d), $a^{i}$ has order
	$\frac{d}{\gcd(d, i)}$, equal to $d$ precisely when $\gcd(i, d) = 1$,
	so exactly $\varphi(d)$ of them have order $d$. Every element of
	order $d$ is a root of $x^{d} \equiv 1$ and hence one of these
	$a^{i}$, so $N(d) \in \{0, \varphi(d)\}$, and thus $N(d) \le
		\varphi(d)$. By Fermat's little theorem and Theorem~\ref{thm:6.1.1}(c) every
	residue class coprime to $p$ has order dividing $p-1$, so $\sum_{d
			\mid p-1} N(d) = p - 1 = \sum_{d \mid p-1}\varphi(d)$ (Theorem~\ref{thm:5.1.5}). Were $N(d) < \varphi(d)$ for even one $d$, the first
	sum would be strictly smaller than the second; hence equality holds
	throughout, and in particular $N(p-1) = \varphi(p-1) \ge 1$.

	In the second step, we choose a primitive root $g$ modulo $p$ with $g^{p-1} \not\equiv 1
		\pmod{p^2}$. Start with any primitive root $g$ modulo $p$ and write
	$g^{p-1} \equiv 1 + kp \pmod{p^2}$. If $p \nmid k$ then $g^{p-1}
		\not\equiv 1 \pmod{p^2}$. If $p \mid k$, replace $g$ by another lift
	modulo $p^2$: among the $p$ integers $g, g+p, \ldots, g+(p-1)p$, all
	congruent to $g$ modulo $p$ (hence still primitive roots modulo $p$)
	but pairwise distinct modulo $p^2$, take $g' = g + p$. Since $g' \equiv g \not\equiv 0 \pmod{p}$, we
	also have $p \nmid g'$. By the binomial theorem
	$g'^{p-1} = (g+p)^{p-1} = \sum_{i=0}^{p-1}
		\binom{p-1}{i} g^{p-1-i} p^{i} \equiv g^{p-1} + (p-1) g^{p-2} p
		\pmod{p^2}$, since every term with $i \ge 2$ contains a factor
	$p^{2}$. Hence $g'^{p-1} \equiv (1 + kp) + (p-1) g^{p-2} p = 1 +
		\big(k + (p-1) g^{p-2}\big) p \pmod{p^2}$. Because $p \mid k$ and
	$(p-1) g^{p-2} \equiv -g^{p-2} \pmod{p}$, we get $k + (p-1) g^{p-2}
		\equiv -g^{p-2} \not\equiv 0 \pmod{p}$ (as $p \nmid g$), so $g'^{p-1}
		\not\equiv 1 \pmod{p^2}$.

	For the third step, let $e \ge 2$ and choose a primitive root $g$
	modulo $p$ with $g^{p-1} \not\equiv 1 \pmod{p^2}$, so $v_p(g^{p-1} -
		1) = 1$; then $g$ is a primitive root modulo $p^{e}$ by
	Theorem~\ref{thm:6.2.5}.

	For the fourth step, take an odd primitive root $g$ modulo $p^e$
	(replace $g$ by $g + p^e$ if it is even; since $g+p^e\equiv g\pmod{p^e}$,
	this does not change the order modulo $p^e$, so $g+p^e$ remains a
	primitive root modulo $p^e$). Since $g$ is odd, it is automatically a
	unit modulo $2p^e$. If $g^t\equiv1\pmod{2p^e}$, then in particular
	$g^t\equiv1\pmod{p^e}$, so $\delta_{p^e}(g)=\varphi(p^e)\mid t$; applying this to
	the least such $t=\delta_{2p^e}(g)$ shows that $\delta_{2p^e}(g)$ is a
	multiple of $\varphi(p^e)$. On the other hand, the Fermat--Euler theorem gives
	$\delta_{2p^e}(g)\mid\varphi(2p^e)=\varphi(p^e)$, so
	$\delta_{2p^e}(g)=\varphi(2p^e)$.

	In Case 3, let $n = 2^{e}$ with $e \ge 3$; we show that no primitive root exists. We prove by induction on $e$ that $a^{2^{e-2}}\equiv1\pmod{2^e}$ for every odd integer $a$.
	For $e = 3$, write $a = 2t+1$, so $a^{2} = 4t(t+1) +
		1$ with $2 \mid t(t+1)$, and $a^{2} \equiv 1 \pmod{8}$ for all odd
	$a$. Now assume the claim for some $e = i \ge 3$, that is, $a^{2^{i-2}} \equiv 1 \pmod{2^{i}}$, and write $a^{2^{i-2}} = 1 + k 2^{i}$.
	Then $a^{2^{i-1}} = \big(a^{2^{i-2}}\big)^{2} = (1 + k 2^{i})^{2} = 1
		+ k 2^{i+1} + k^{2} 2^{2i} \equiv 1 \pmod{2^{i+1}}$, so the statement
	holds for $e = i+1$.
	Hence $a^{2^{e-2}} \equiv 1 \pmod{2^e}$ for every $a$ coprime to
	$2^{e}$, so Theorem~\ref{thm:6.1.1}(c) gives $\delta_{2^e}(a) \mid
		2^{e-2}$ and therefore $\delta_{2^e}(a) \le 2^{e-2} < 2^{e-1} =
		\varphi(2^e)$. As $a$ was an arbitrary unit modulo $2^{e}$, no
	element has order $\varphi(2^e)$, and no primitive root modulo
	$2^{e}$ exists.

	Finally, Case 4 covers every remaining $n$. It remains to show
	that any $n$ not of the form $1, 2, 4, p^{e}, 2p^{e}$ (with $p$ an
	odd prime and $e \ge 1$) and not a power of $2$ has no primitive
	root. Factor $n = 2^{c} p_1^{e_1} \cdots p_r^{e_r}$ with $c \ge 0$
	and the $p_j$ distinct odd primes; since $n$ is not a power of $2$,
	we have $r \ge 1$. We claim that $n$ factors as $n = n_1 n_2$ with $\gcd(n_1, n_2) = 1$ and with both
	$\varphi(n_1)$ and $\varphi(n_2)$ even; the two subcases below are
	exhaustive:
	\begin{itemize}
		\item If $r \ge 2$, take $n_1 = p_1^{e_1}$ and $n_2 = n/n_1$. Then
		      $\varphi(n_1) = p_1^{e_1-1}(p_1-1)$ is even (as $p_1$ is odd), and
		      $n_2$ is divisible by the odd prime $p_2$, so $\varphi(n_2)$ is a
		      multiple of $\varphi(p_2^{e_2}) = p_2^{e_2-1}(p_2-1)$ and hence even.
		\item If $r = 1$, then $n = 2^{c} p_1^{e_1}$; since $n \neq
			      p_1^{e_1}, 2p_1^{e_1}$ we must have $c \ge 2$. Take $n_1 = 2^{c}$,
		      $n_2 = p_1^{e_1}$. Then $\varphi(n_1) = 2^{c-1}$ is even (as $c \ge
			      2$) and $\varphi(n_2) = p_1^{e_1-1}(p_1-1)$ is even.
	\end{itemize}
	In either subcase, $\gcd(n_1, n_2) = 1$ and $2 \mid
		\gcd(\varphi(n_1), \varphi(n_2))$. Since $\varphi$ is multiplicative and $\gcd(n_1, n_2) = 1$,
	we have $\varphi(n) = \varphi(n_1)\varphi(n_2)$, so
	\[
		L := \lcm(\varphi(n_1), \varphi(n_2)) =
		\frac{\varphi(n_1)\varphi(n_2)}{\gcd(\varphi(n_1), \varphi(n_2))} =
		\frac{\varphi(n)}{\gcd(\varphi(n_1), \varphi(n_2))} \le
		\frac{\varphi(n)}{2} < \varphi(n).
	\]
	For every $a$ coprime to $n$, the Fermat--Euler theorem gives $a^{\varphi(n_1)}
		\equiv 1 \pmod{n_1}$ and $a^{\varphi(n_2)} \equiv 1 \pmod{n_2}$; as
	$L$ is a common multiple of $\varphi(n_1)$ and $\varphi(n_2)$, we get
	$a^{L} \equiv 1 \pmod{n_1}$ and $a^{L} \equiv 1 \pmod{n_2}$, whence
	$a^{L} \equiv 1 \pmod{n}$ (since $n = n_1 n_2$ with $\gcd(n_1, n_2) =
		1$). Thus $\delta_n(a) \le L < \varphi(n)$ for every such $a$, so no
	primitive root exists.

	Combining all cases completes the proof.

\end{proof}
\begin{altproof}[Alternative proof of Theorem~\ref{thm:6.2.6}, Case 2, Step 1: root count] First we prove that for all $d \mid p-1$, the congruence $x^{d}
		\equiv 1 \pmod{p}$ has exactly $d$ distinct solutions. Since $d \mid
		p-1$, write $x^{p-1} - 1 = (x^{d} - 1)\,h(x)$ with $\deg h = p - 1 -
		d$. By Fermat's little theorem $x^{p-1} - 1 \equiv 0 \pmod{p}$ has
	$p-1$ distinct roots; by the polynomial root bound
	(Theorem~\ref{thm:3.6.2}), $x^{d} - 1$ has at most $d$ roots
	and $h(x)$ at most $p-1-d$ roots. The two root sets are disjoint: if $r^{d}\equiv1\pmod{p}$, then $h(r)\equiv1+r^{d}+\cdots+r^{p-1-d}\equiv\frac{p-1}{d}\not\equiv0\pmod{p}$, where the last step follows because $1\le\frac{p-1}{d}\le p-1$. Their union consists of all $p-1$ nonzero residues, so both upper bounds are attained, and $x^d-1$ has exactly $d$ roots. By Theorem~\ref{thm:6.1.1}(c) the solutions of $x^{d}\equiv1\pmod{p}$ are exactly the residue classes whose order divides $d$, so $\sum_{c \mid d} N(c) = d$
	for all $d \mid p-1$. Comparing with $\sum_{c \mid d} \varphi(c) = d$ (Theorem~\ref{thm:5.1.5})
	and applying M{\"o}bius inversion gives $N(d) = \varphi(d)$. Thus
	$N(p-1) = \varphi(p-1) > 0$ and a primitive root modulo $p$ exists.

\end{altproof}
\begin{altproof}[Alternative proof of Theorem~\ref{thm:6.2.6}, Case 2, Step 1: prime-power assembly]
	As shown in the alternative proof above, for $d \mid p-1$ the congruence $x^{d}
		\equiv 1 \pmod{p}$ has exactly $d$ solutions in
	$\ZmodUnits{p}$. We claim that if $q^{c} \mid p-1$ with $q$ prime
	and $c \ge 1$, then modulo $p$ there are exactly $q^{c} - q^{c-1} >
		0$ elements of order $q^{c}$; in particular, such an element exists.
	Indeed, the $q^{c}$ solutions of $x^{q^{c}} \equiv 1 \pmod{p}$ all
	have order dividing $q^{c}$; those whose order properly divides
	$q^{c}$ have order dividing $q^{c-1}$, hence satisfy $x^{q^{c-1}}
		\equiv 1 \pmod{p}$, and there are exactly $q^{c-1}$ of these.
	Subtracting these counts shows that exactly $q^{c} - q^{c-1} > 0$ elements have order
	$q^{c}$. Now write $p-1 = \prod_{i=1}^{t} q_i^{c_i}$ with the $q_i$
	distinct primes, and use this claim to choose $a_i$ of order
	$q_i^{c_i}$ for each $i$. The orders $q_1^{c_1}, \ldots, q_t^{c_t}$
	are pairwise coprime, so applying Theorem~\ref{thm:6.1.1}(f) repeatedly gives
	$\delta_p(a_1 a_2 \cdots a_t) = \prod_{i=1}^{t} q_i^{c_i} = p-1$.
	Hence $g := a_1 a_2 \cdots a_t$ is a primitive root modulo $p$.

\end{altproof}
\begin{altproof}[Alternative proof of Theorem~\ref{thm:6.2.6}, Case 3: the power congruence] Write $a = 2k+1$. Then
	\[
		a^{2^{e-2}} = (2k+1)^{2^{e-2}} \equiv 1 + \binom{2^{e-2}}{1}(2k) +
		\binom{2^{e-2}}{2}(2k)^{2} \pmod{2^e}
	\]
	(for each $j \ge 3$, Theorem~\ref{thm:6.2.4} gives
	$v_2\big(\binom{2^{e-2}}{j}(2k)^{j}\big) \ge e$, so the $j$-th term vanishes
		modulo $2^e$). Simplifying gives
		\[
		a^{2^{e-2}} \equiv 1 + 2^{e-1}k + 2^{e-1}(2^{e-2}-1)k^{2} = 1 +
			2^{e-1}\big(k + (2^{e-2}-1)k^{2}\big) \pmod{2^e}.
	\]

	Since $2^{e-2}-1$ is odd for $e \ge 3$, we have $k + (2^{e-2}-1)k^{2} \equiv
		k + k^{2} = k(k+1) \equiv 0 \pmod{2}$, so $a^{2^{e-2}} \equiv 1 \pmod{2^e}$.

	Then Theorem~\ref{thm:6.1.1}(c) gives $\delta_{2^e}(a) \mid 2^{e-2}$, so $\delta_{2^e}(a) \le 2^{e-2}
		< 2^{e-1} = \varphi(2^e)$; since $a$ was an arbitrary unit modulo
	$2^{e}$, no element has order $\varphi(2^e)$, and hence no primitive root
	modulo $2^e$ exists.

\end{altproof}
\begin{altproof}[Alternative proof of Theorem~\ref{thm:6.2.4}]
	The case $k=0$ is immediate, so assume $k\neq0$.
	Put $s = e-2$, so the
	binomial coefficient in question is $\binom{2^{s}}{j}$. From $\binom{2^{s}}{j} =
		\frac{(2^{s})!}{j!\,(2^{s}-j)!}$ we obtain the absorption identity
	$j\binom{2^{s}}{j} = 2^{s}\binom{2^{s}-1}{j-1}$; applying $v_2$ to both sides gives
	\[
		v_2(j) + v_2\binom{2^{s}}{j} = s + v_2\binom{2^{s}-1}{j-1}.
	\]

	The integer $2^{s}-1 = (11\cdots1)_2$ has all
	$s$ of its base-$2$ digits equal to $1$, so for any $0 \le j-1 \le
		2^{s}-1$ every digit satisfies $(j-1)_i \le 1$ and hence, by Lucas's theorem (Theorem~\ref{thm:4.1.2}),
	$\binom{2^{s}-1}{j-1} \equiv \prod_{i=0}^{s-1}\binom{1}{(j-1)_i} \equiv 1
		\pmod{2}$; that is, $\binom{2^{s}-1}{j-1}$ is odd, so
	$v_2\binom{2^{s}-1}{j-1} = 0$. Therefore
	\[
		v_2\binom{2^{s}}{j} = s - v_2(j) = (e-2) - v_2(j),
	\]
	the same value obtained above using Kummer's theorem. Hence
	$v_2\big(\binom{2^{e-2}}{j}(2k)^{j}\big) \ge (e-2) - v_2(j) + j = e +
		\big(j - v_2(j) - 2\big) \ge e$, again using $j - v_2(j) \ge 2$ for $j \ge 3$.

\end{altproof}
\begin{altproof}[Alternative proof of Theorem~\ref{thm:6.2.6}, Case 3: noncyclicity] Consider $f : \ZmodUnits{2^e} \to
		\ZmodUnits{8}$, $f(a) = a \bmod 8$, which is a
	homomorphism; it is surjective because $e\ge3$, so each odd residue
	modulo $8$ is the reduction of the same odd integer modulo $2^e$. The image of a cyclic group under a homomorphism is
	cyclic: a generator maps to a generator of the image. But every element of $\ZmodUnits{8} = \{1, 3, 5, 7\}$
	satisfies $x^{2} \equiv 1 \pmod{8}$, so $\ZmodUnits{8}$ has no
	element of order $4=|\ZmodUnits{8}|$, and hence it is not cyclic.
	More explicitly, the map $(u,v)\mapsto(-1)^u3^v\pmod8$, for
	$u,v\in\{0,1\}$ with addition modulo $2$, is a bijective homomorphism
	from $\mathbb{Z}_2\times\mathbb{Z}_2$ onto $\ZmodUnits{8}$:
	the four images are $1,3,7,5$, and $(-1)^2\equiv3^2\equiv1\pmod8$.
	If $\ZmodUnits{2^e}$ were cyclic, then
	$\ZmodUnits{8}$ (its homomorphic image) would also be cyclic,
	a contradiction. So $\ZmodUnits{2^e}$ is not cyclic and no
	primitive root exists.

\end{altproof}
\begin{altproof}[Alternative proof of Theorem~\ref{thm:6.2.6}, Case 3: LTE] For any $a \in \mathbb{Z}^+$ with
	$2 \nmid a$, we have $\gcd(a, 2^e) = 1$; the case $a = 1$ is trivial, so assume $a > 1$, whence $a \neq \pm 1$. Since $e \ge 3$, the exponent $2^{e-2}$ is even, so Theorem~\ref{thm:4.8.6} gives $v_2\big(a^{2^{e-2}} -
		1\big) = v_2(a-1) + v_2(a+1) + v_2(2^{e-2}) - 1 = v_2(a-1) + v_2(a+1)
		+ e - 3 \ge e$, where the last step follows because $a-1$ and $a+1$ are consecutive
	even integers, so one of them is divisible by $4$ and $v_2(a-1) +
		v_2(a+1) \ge 3$. So
	$a^{2^{e-2}} \equiv 1 \pmod{2^e}$, hence $\delta_{2^e}(a) \mid
		2^{e-2}$ by Theorem~\ref{thm:6.1.1}(c), so $\delta_{2^e}(a) \le 2^{e-2} < 2^{e-1} = \varphi(2^e)$ and no primitive root modulo $2^{e}$ exists.

\end{altproof}
\begin{theorem}\label{thm:6.2.7}
	Suppose $m$ has a primitive root. For $d \in \mathbb{Z}^+$, there
	is an element of order $d$ modulo $m$ if and only if $d \mid
		\varphi(m)$; in that case, the number of pairwise incongruent elements
	of order exactly $d$ is $\varphi(d)$.

\end{theorem}

\begin{proof} Let $g$ be a primitive root. By Theorem~\ref{thm:6.2.1}, every unit is congruent to $g^{k}$ for a unique $0 \le k <
		\varphi(m)$, with $\delta_m(g^{k})
		= \frac{\varphi(m)}{\gcd(\varphi(m), k)}$ by Theorem~\ref{thm:6.1.1}(d).
	($\Rightarrow$) The Fermat--Euler theorem and Theorem~\ref{thm:6.1.1}(c) show that every
	order divides $\varphi(m)$. ($\Leftarrow$) If $d \mid \varphi(m)$, then
	$g^{k}$ has order $d$ if and only if $\gcd(\varphi(m), k) =
		\frac{\varphi(m)}{d}$; writing $k = \frac{\varphi(m)}{d}j$ with
	$0\le j<d$, we see that this holds if and only if $\gcd(j,d)=1$, giving exactly $\varphi(d)$ such
	elements. Taking $d = \varphi(m)$ recovers Theorem~\ref{thm:6.2.2}.

\end{proof}
\begin{definition}[Fibonacci primitive root]\label{def:6.2.2} For a
	prime $p$, an integer $g$ is a Fibonacci primitive root modulo $p$ if $g$ is a
	primitive root modulo $p$ and $g^{2} \equiv g + 1 \pmod{p}$.

\end{definition}
\begin{theorem}\label{thm:6.2.8} Let $p$ be an odd prime, and suppose that
	$g$ is a Fibonacci primitive root modulo $p$. Then $g-1$ is also a primitive
	root modulo $p$. The existence of such a $g$ requires $p\ge5$.

\end{theorem}

\begin{proof} The defining congruence gives directly
	\[
		g(g-1)=g^2-g\equiv1\pmod p.
	\]
	Thus $g-1$ is a unit and represents $g^{-1}$ modulo $p$.
	For every positive integer $k$,
	we have $(g^{-1})^k\equiv1\pmod p$ if and only if $g^k\equiv1\pmod p$,
	so the two units have the same order. Consequently
	\[
		\delta_p(g-1)=\delta_p(g^{-1})=\delta_p(g)=p-1,
	\]
	and $g-1$ is a primitive root. If $p=3$, the only primitive-root
	class is $2$, but $2^2\equiv1\not\equiv0\equiv2+1\pmod3$.
	Since $p$ is odd, the existence hypothesis therefore forces $p\ge5$.

\end{proof}
\subsection{The Index (Discrete Logarithm)}\label{subsec:6.3}

\begin{definition}[index (discrete logarithm)]\label{def:6.3.1} For
	$m \in \mathbb{Z}^+$ with a primitive root $g$ and with $\gcd(a, m) =
		1$, an integer $r$ satisfying $g^{r} \equiv a \pmod{m}$ is called an
	index or discrete logarithm of $a$ to the base $g$ modulo $m$.
	Here powers are taken in $\ZmodUnits{m}$; in particular, a negative
	exponent denotes a power of the inverse class.

\end{definition}
\begin{theorem}\label{thm:6.3.1} If $\gcd(a, m) = \gcd(b, m) = 1$
	and $g$ is a primitive root modulo $m$, there is exactly one integer
	$0 \le r < \varphi(m)$ with $g^{r} \equiv a \pmod{m}$. Denote it by
	$\ind_{m,g}(a)$; the same notation is used for the unit class of $a$. Then:

	\begin{enumerate}[label={\alph*)}]
		\item For every $r\in\mathbb{Z}$, we have $g^{r} \equiv a \pmod{m}$ if and only if $r \equiv
			      \ind_{m,g}(a) \pmod{\varphi(m)}$.
		\item $\ind_{m,g}(ab) \equiv
			      \ind_{m,g}(a) + \ind_{m,g}(b) \pmod{\varphi(m)}$.
		\item For every $n\in\mathbb{Z}$, we have $\ind_{m,g}(a^{n}) \equiv n\,
			      \ind_{m,g}(a) \pmod{\varphi(m)}$.
		\item $\ind_{m,g}(a) \equiv
			      \ind_{m,g'}(a) \cdot \ind_{m,g}(g')
			      \pmod{\varphi(m)}$, where $g'$ is another primitive root.
		\item $\delta_m(a) = \dfrac{\varphi(m)}{\gcd(\varphi(m), \, \ind_{m,g}(a))}$.
	\end{enumerate}

\end{theorem}

\begin{proof} Throughout the proof, we abbreviate
	$\ind_{m,g}$ as $\ind$ when the base $g$
	and modulus $m$ are clear from context. The
	$\varphi(m)$ powers $g^{0},\ldots,g^{\varphi(m)-1}$ are pairwise
	distinct by Theorem~\ref{thm:6.1.1}(b), and all are units. Since there
	are exactly $\varphi(m)$ unit classes, these powers represent every
	unit class exactly once. This proves the existence and uniqueness of
	$\ind_{m,g}(a)$.

	\begin{enumerate}[label={\alph*)}]
		\item Write $r=q\varphi(m)+t$ with $q\in\mathbb{Z}$ and
		      $0\le t<\varphi(m)$. In the unit group,
		      $g^r=(g^{\varphi(m)})^qg^t\equiv g^t\pmod m$,
		      including when $q<0$. By uniqueness in the indicated range,
		      $g^r\equiv a\pmod m$ if and only if $t=\ind_{m,g}(a)$,
		      equivalently $r\equiv\ind_{m,g}(a)\pmod{\varphi(m)}$.
		\item Since $g^{\ind(ab)} \equiv ab \equiv
			      g^{\ind(a) + \ind(b)} \pmod{m}$, part (a) shows that the indices
		      agree modulo $\varphi(m)$.
		\item In the unit group, $a\equiv g^{\ind(a)}$ implies
		      $a^n\equiv g^{n\ind(a)}$ for every $n\in\mathbb{Z}$,
		      using inverses when $n<0$. Part (a) therefore gives
		      $\ind(a^n)\equiv n\ind(a)\pmod{\varphi(m)}$.
		\item Since $g^{\ind_{m,g}(a)} \equiv a \equiv
			      g'^{\ind_{m,g'}(a)} \equiv
			      g^{\ind_{m,g}(g')\, \ind_{m,g'}(a)} \pmod{m}$,
		      part (a) gives the identity.
		\item By Theorem~\ref{thm:6.1.1}(d) with $a \equiv
			      g^{\ind_{m,g}(a)}$ and $\delta_m(g) = \varphi(m)$,
		      we get $\delta_m(a) = \frac{\varphi(m)}{\gcd(\varphi(m), \, \ind_{m,g}(a))}$.
	\end{enumerate}

\end{proof}
\begin{altproof}[Alternative proof of Theorem~\ref{thm:6.2.2}] Fix a modulus $m$ admitting a primitive root, and choose
	a primitive root $g_0$. A unit $a$ is a primitive root if and only if
	$\delta_m(a)=\varphi(m)$, equivalently
	$\gcd(\varphi(m),\ind_{m,g_0}(a))=1$ by Theorem~\ref{thm:6.3.1}(e). Since
	$x\mapsto\ind_{m,g_0}(x)\bmod\varphi(m)$ is a
	bijection $\ZmodUnits{m}\to\mathbb{Z}_{\varphi(m)}$, the
	number of primitive roots equals the number of $r$ with $0 \le r <
		\varphi(m)$ and $\gcd(r, \varphi(m)) = 1$, which is $\varphi(\varphi(m))$.

\end{altproof}
\subsection{The Carmichael Function}

\begin{definition}[Carmichael function $\lambda(m)$]\label{def:6.4.1} For $m \in \mathbb{Z}^+$, the Carmichael function $\lambda(m)$ is
	the smallest positive integer $n$ satisfying $a^n\equiv1\pmod m$
	for every integer $a$ with $\gcd(a,m)=1$.
	The Fermat--Euler theorem supplies such an exponent, namely $\varphi(m)$.

\end{definition}
\begin{theorem}\label{thm:6.4.1} For $m\in\mathbb{Z}^+$, we have $\lambda(m) =
		\lcm_{\gcd(a, m)=1}\big(\delta_m(a)\big) =
		\max_{\gcd(a, m)=1}\{\delta_m(a)\}$.

\end{theorem}

\begin{proof} Orders depend only on residue classes, so the lcm and
	maximum may be taken over the finite nonempty set of representatives
	$1\le a\le m$ with $\gcd(a,m)=1$. Put
	$L=\lcm_{\gcd(a,m)=1}\delta_m(a)$.
	For $n\ge1$, Theorem~\ref{thm:6.1.1}(c) shows that $a^n\equiv1\pmod m$
	for every $a$ coprime to $m$ if and only if every order $\delta_m(a)$
	divides $n$, or equivalently $L\mid n$. Hence $\lambda(m)=L$.

	It remains to show that some unit has order $L$. If $L=1$, take $a=1$.
	Otherwise, write $L = \prod_{j=1}^{s}
		q_j^{c_j}$ with distinct primes $q_j$ and $c_j \ge 1$. For each $j$,
	there is some $a_j$ coprime to $m$ with $q_j^{c_j} \mid
		\delta_m(a_j)$, since otherwise $v_{q_j}(\delta_m(a)) \le c_j - 1$ for
	all $a$ and $q_j^{c_j} \nmid L$. Put $b_j = a_j^{\,\delta_m(a_j) /
				q_j^{c_j}}$; then $\gcd\big(\delta_m(a_j), \delta_m(a_j)/q_j^{c_j}\big)
		= \delta_m(a_j)/q_j^{c_j}$, so Theorem~\ref{thm:6.1.1}(d) gives $\delta_m(b_j) =
		q_j^{c_j}$. The orders $\delta_m(b_1), \ldots, \delta_m(b_s)$ are
	pairwise coprime, so $s-1$ applications of Theorem~\ref{thm:6.1.1}(f) give
	$\delta_m(b_1b_2\cdots b_s) = \prod_{j=1}^{s} q_j^{c_j} = L$. Hence
	$\lambda(m) = L = \max_{\gcd(a, m)=1}\{\delta_m(a)\}$.
	Also, the Fermat--Euler theorem gives $\delta_m(a)\mid\varphi(m)$ for
	every unit, so taking the lcm of these orders yields $\lambda(m)\mid\varphi(m)$.

\end{proof}
\begin{theorem}\label{thm:6.4.2} At $1$ and at prime powers, $\lambda$ takes the values
	\[
		\lambda(1)=1,\qquad \lambda(2)=1,\qquad \lambda(4)=2,
	\]
	\[
		\lambda(2^{e})=2^{e-2}\ (e\ge3),\qquad \lambda(p^{e})=\varphi(p^{e})\ (p\ \text{odd prime},\ e\ge1).
	\]
	More generally, if $m=p_1^{e_1}\cdots p_k^{e_k}>1$ is the prime
	factorization of $m$, then
	\[
		\lambda(m)=\lcm\big(\lambda(p_1^{e_1}),\ldots,\lambda(p_k^{e_k})\big).
	\]

\end{theorem}

\begin{proof} We first handle Case 1, a general modulus $m>1$. By
	Theorem~\ref{thm:6.1.1}(i), $\delta_m(a) =
		\lcm_i(\delta_{p_i^{e_i}}(a))$. We claim that for each
	$i$, the set of orders $\{\delta_{p_i^{e_i}}(a) \mid \gcd(a, m)=1\}$
	equals $\{\delta_{p_i^{e_i}}(b) \mid \gcd(b, p_i^{e_i})=1\}$: the forward
	inclusion is immediate since $\gcd(a, m)=1$ implies
	$\gcd(a, p_i^{e_i})=1$; for the reverse inclusion, given $b$ with
	$\gcd(b, p_i^{e_i})=1$, the Chinese remainder theorem supplies
	$a \in \mathbb{Z}^+$ with $a \equiv b \pmod{p_i^{e_i}}$ and $a
		\equiv 1 \pmod{p_j^{e_j}}$ for all $j \neq i$; then $\gcd(a, m)=1$ and
	$\delta_{p_i^{e_i}}(a) = \delta_{p_i^{e_i}}(b)$, so
	$\delta_{p_i^{e_i}}(b)$ lies in the first set. Hence $\lambda(m) =
		\lcm_i(\lambda(p_i^{e_i}))$. Next consider Case 2, where $m = 2^{e}$ and $e \ge 3$. Every odd $a$ satisfies
	$a^{2^{e-2}} \equiv 1 \pmod{2^e}$ (Case 3 in the proof of
	Theorem~\ref{thm:6.2.6}), so
	$\delta_{2^e}(a) \mid 2^{e-2}$ for every such $a$, whence
	$\lambda(2^e) \mid 2^{e-2}$. It remains to exhibit an element of
	order exactly $2^{e-2}$; we show that $5$ is one. For $e = 3$,
	we have $5^{2^{e-3}} = 5$ and $v_2(5-1) = 2 = e-1$; for $e \ge 4$, the exponent $2^{e-3}$
	is even, so the $p=2$ lifting-the-exponent lemma (Theorem~\ref{thm:4.8.6}) gives $v_2\big(5^{2^{e-3}}
		- 1\big) = v_2(5-1) + v_2(5+1) + v_2(2^{e-3}) - 1 = 2 + 1 + (e-3) - 1
		= e-1$. In both cases, $v_2\big(5^{2^{e-3}} - 1\big) = e-1 < e$, so
	$5^{2^{e-3}} \not\equiv 1 \pmod{2^e}$. Since $\delta_{2^e}(5) \mid
		2^{e-2}$ while $\delta_{2^e}(5) \nmid 2^{e-3}$ (every proper divisor
	of $2^{e-2}$ divides $2^{e-3}$), we conclude that $\delta_{2^e}(5) =
		2^{e-2}$. Therefore $\lambda(2^{e}) = 2^{e-2} =
		\tfrac{1}{2}\varphi(2^{e})$. It remains to handle Case 3, consisting of $m = 1,2,4$ and $m=p^e$ for odd $p$. For $m=1,2,4$ the equality
	$\lambda(m)=\varphi(m)$ is immediate. For $m=p^e$ with $p$ odd,
	Theorem~\ref{thm:6.2.6} provides a primitive root modulo $p^e$, say $g$, so
	$\delta_{p^e}(g)=\varphi(p^e)$. Since we always have
	$\lambda(p^e)=\max_{\gcd(a, p^e)=1}\delta_{p^e}(a)$ (Theorem~\ref{thm:6.4.1}) and
	$\lambda(p^e)\mid\varphi(p^e)$, it follows that
	$\lambda(p^e)=\varphi(p^e)$.

\end{proof}
\begin{theorem}\label{thm:6.4.3} For $m\in\mathbb{Z}^+$, a primitive root modulo $m$ exists
	if and only if $\lambda(m) = \varphi(m)$.

\end{theorem}

\begin{proof} For every $a$ coprime to $m$, the Fermat--Euler theorem and
	Theorem~\ref{thm:6.1.1}(c) give
	$\delta_m(a)\mid\varphi(m)$. Taking the least common multiple over all
	such $a$ shows that $\lambda(m)\mid\varphi(m)$. By
	Theorem~\ref{thm:6.4.1}, some unit $g$ has
	$\delta_m(g)=\lambda(m)$. Thus if $\lambda(m)=\varphi(m)$, this $g$ is
	a primitive root. Conversely, if $g$ is a primitive root, then
	$\varphi(m)=\delta_m(g)\mid\lambda(m)$; together with
	$\lambda(m)\mid\varphi(m)$, this gives $\lambda(m)=\varphi(m)$.

\end{proof}
\begin{theorem}\label{thm:6.4.4} Let $e \ge 2$ be an integer. Then the residues
	$\{\pm5^k \mid 0\le k<2^{e-2}\}$ are pairwise incongruent and form a reduced
	residue system modulo $2^{e}$.

\end{theorem}

\begin{proof} For $e = 2$, we have $\{1, 3\} = \{1, -1\}$. For $e
		\ge 3$, the proof of Theorem~\ref{thm:6.4.2} shows that $5$ has
	order $2^{e-2}$, so $1, 5, \ldots,
		5^{2^{e-2}-1}$ are distinct and congruent to $1 \pmod{4}$, while their negatives are
	congruent to $3 \pmod{4}$; the $\pm 5^{k}$ give $2^{e-1} = \varphi(2^{e})$
	incongruent residues, a reduced residue system.

\end{proof}
\subsection{Carmichael Numbers}

\begin{definition}[Carmichael number]\label{def:6.5.1} A composite
	$n \in \mathbb{Z}^+$ is a Carmichael number if $a^{n-1} \equiv 1
		\pmod{n}$ for all $a$ with $\gcd(a, n) = 1$.

\end{definition}
\begin{theorem}\label{thm:6.5.1} A composite number $n$ is a
	Carmichael number if and only if $\lambda(n) \mid (n-1)$.

\end{theorem}

\begin{proof} By Theorem~\ref{thm:6.1.1}(c), $a^{n-1} \equiv 1
		\pmod{n}$ for all $a$ coprime to $n$ if and only if $\delta_n(a) \mid
		(n-1)$ for all such $a$, that is, if and only if $\lambda(n) =
		\lcm_{\gcd(a, n)=1}(\delta_n(a)) \mid (n-1)$.

\end{proof}
\begin{theorem}[Korselt's criterion]\label{thm:6.5.2} A composite
	number $n$ is a Carmichael number if and only if $n$ is squarefree
	and $(p-1) \mid (n-1)$ for every prime $p \mid n$. All Carmichael
	numbers are odd and have at least $3$ distinct prime factors.

\end{theorem}

\begin{proof} ($\Leftarrow$) If $n = p_1 \cdots p_k$ is
	squarefree with $(p_i-1) \mid (n-1)$, then $\lambda(n) =
		\lcm(p_1-1, \ldots, p_k-1) \mid (n-1)$, so, by Theorem~\ref{thm:6.5.1}, $n$ is a Carmichael number. ($\Rightarrow$) By Theorem~\ref{thm:6.5.1},
	$\lambda(n) \mid (n-1)$. We first show that $n$ is squarefree. If $p^{2} \mid n$ for
	some prime $p$, write $n = p^{a} n'$ with $a \ge 2$, $\gcd(p, n') =
		1$; by Theorem~\ref{thm:6.4.2}, $\lambda(n) =
		\lcm(\lambda(p^{a}), \lambda(n'))$. Since $a \ge 2$,
	the value $\lambda(p^{a})$ equals $p^{a-1}(p-1)$ when $p$ is odd, or $2^{a-2}$
	(with $\lambda(4)=2$ at $a=2$) when $p=2$; in either case $p \mid
		\lambda(p^{a})$. Hence $p \mid \lambda(p^{a}) \mid \lambda(n) \mid
		(n-1)$, but $p \mid n$ gives $p \nmid (n-1)$, a contradiction.
	Next fix a prime $p \mid n$ to prove the prime-factor condition. Take a primitive root $g$ modulo $p$ and, since $n$ is
	squarefree and hence $\gcd(p, n/p) = 1$, use the Chinese remainder
	theorem to choose an integer $a$ with $a \equiv g \pmod{p}$ and $a
		\equiv 1 \pmod{n/p}$. Then $\gcd(a, n) = 1$, so the Fermat property of
	$n$ gives $a^{n-1} \equiv 1 \pmod{n}$, and reduction modulo $p$ yields
	$g^{n-1} \equiv a^{n-1} \equiv 1 \pmod{p}$. As $\delta_p(g) = p-1$,
	Theorem~\ref{thm:6.1.1}(c) gives $(p-1) \mid (n-1)$. The integer $n$ must be odd. Indeed, if $n$ were even, then, since $n$ is squarefree
	and composite, $n$ would not be a power of $2$ (the only even squarefree
	power of $2$ is $2$ itself, which is prime), so $n$ would have an odd prime
	factor $p$; then $(p-1) \mid (n-1)$ with $n-1$ odd forces $p-1$
	to be odd, which is impossible. Finally, $n$ has at least three prime factors. Since $n$ is squarefree, if it had
	exactly two prime factors we could write $n=p_1p_2$ with $p_1<p_2$.
	The prime-factor condition for $p_2$ would give
	$(p_2-1)\mid(n-1)$. But
	$n-1=p_1(p_2-1)+(p_1-1)$, so subtracting the first term shows
	$(p_2-1)\mid(p_1-1)$, which is impossible because
	$0<p_1-1<p_2-1$.

\end{proof}
\subsection{The Group-Theoretic Viewpoint}\label{subsec:6.6}

\begin{definition}[group and abelian group]\label{def:6.6.1}
	A group is a set $G$ with an associative binary operation, an identity
	element, and an inverse for each element. A group is abelian if its
	operation is commutative.

\end{definition}
\begin{definition}[subgroup]\label{def:6.6.2}
	A subgroup is a subset of a group that contains the identity and is
	closed under the group operation and under inversion.

\end{definition}
\begin{definition}[cyclic group, generator, and element order]\label{def:6.6.3}
	For an element $g$ of a group, the cyclic subgroup generated by $g$ is
	$\langle g\rangle=\{g^k\mid k\in\mathbb{Z}\}$. The order of $g$ is the
	size of $\langle g\rangle$; when this size is finite, it is equivalently
	the least positive integer $r$ with $g^r=1$. A group is cyclic if it is
	generated by one element, and such an element is called a generator.
	For $n\in\mathbb{Z}^+$, the symbol $C_n$ denotes a cyclic group of order $n$;
	it is determined up to isomorphism by Theorem~\ref{thm:6.6.1}.
	A direct product of groups has componentwise multiplication; the empty
	direct product is the one-element group.
	In particular, $\mathbb{Z}/n\mathbb{Z}$ is a cyclic group under addition,
	and $\ZmodUnits{n}$ is a group under multiplication modulo $n$. For an
	integer $a$ coprime to $m$, the residue class of $a$ generates
	$\langle a\rangle\subseteq\ZmodUnits{m}$, and its order is precisely
	$\delta_m(a)$ of Definition~\ref{def:6.1.1}, since both are the least
	positive $r$ with $a^r\equiv1\pmod m$.

\end{definition}
\begin{theorem}[uniqueness of $C_n$]\label{thm:6.6.1}
	Let $n\in\mathbb{Z}^+$. Any two cyclic groups of order $n$ are isomorphic.

\end{theorem}

\begin{proof} If $g$ and $h$ generate cyclic
	groups of order $n$, then $g^k=g^l$ and $h^k=h^l$ each hold exactly when
	$n\mid k-l$ (Theorem~\ref{thm:6.1.1}(c), which holds for any element of finite
	order), so $g^k\mapsto h^k$ is a well-defined bijective homomorphism.

\end{proof}
\begin{definition}[homomorphism, kernel, coset, and commutative ring]\label{def:6.6.4}
	A homomorphism between groups is a map preserving the operation; it is
	an isomorphism if it is also a bijection, an endomorphism if source and
	target coincide, and an automorphism if it is both an isomorphism and an endomorphism. The kernel of a
	homomorphism is the preimage of the identity. For a subgroup $H\le G$
	and $g\in G$, the left coset $gH$ is $\{gh\mid h\in H\}$.

	A commutative ring with identity is a set whose addition makes it an
	abelian group, and whose multiplication is associative, commutative,
	distributive over addition, and has an identity $1$.
	This is the sense in which the word ring is used throughout, with
	one exception: the proof of Theorem~\ref{thm:8.1.8} uses quaternions
	and $2 \times 2$ complex matrices, whose multiplication is not
	commutative.
	The same terminology as for groups is used for rings, with a
	homomorphism required
	to preserve addition, multiplication, and $1$. For a ring homomorphism, the kernel is the preimage
	of the additive identity $0$. A ring has no zero divisors if a product
	of nonzero elements is always nonzero. A subring is a subset that
	contains $0$ and $1$ and is closed under addition, negation, and
	multiplication.

\end{definition}


\begin{theorem}[Lagrange's theorem for finite groups]\label{thm:6.6.2} If $G$ is a
	finite group and $H \le G$ is a subgroup, then $|H|$ divides $|G|$, and
	the number of left cosets $[G:H] = |G|/|H|$ is the index of $H$ in
	$G$. In particular, the order of any element $g \in G$ (the size of
	the cyclic subgroup $\langle g\rangle$) divides $|G|$.

\end{theorem}

\begin{proof} Declare $a \sim b$ if and only if $a^{-1}b \in H$;
	this is an equivalence relation whose classes are the left cosets
	$gH = \{gh \mid h \in H\}$, so the cosets partition $G$. For each $g$,
	the map $h \mapsto gh$ is a bijection $H \to gH$ (its inverse is left
	multiplication by $g^{-1}$), so every coset has exactly $|H|$
	elements. Summing over the $[G:H]$ cosets gives $|G| = [G:H]\,|H|$,
	hence $|H| \mid |G|$. Taking $H = \langle g\rangle$, whose order is
	the order of $g$, shows that the order of $g$ divides $|G|$.

\end{proof}
\begin{theorem}\label{thm:6.6.3} $\ZmodUnits{2} \cong C_1$,
	$\ZmodUnits{4} \cong C_2$, $\ZmodUnits{2^{e}} \cong
		C_2 \times C_{2^{e-2}}$ ($e \ge 3$), and, for an odd prime $p$ and
	an integer $e\ge1$,
	we have $\ZmodUnits{p^{e}} \cong C_{p^{e-1}(p-1)}$ and $\ZmodUnits{2
			p^{e}} \cong C_{\varphi(2p^{e})}$.

\end{theorem}

\begin{proof} For $m = 2$, the group $\ZmodUnits{2}
		= \{1\}$ is trivial, and $\ZmodUnits{4} = \{1, 3\}$ is cyclic of
	order $2$, generated by $3$. For an odd prime $p$ and $e \ge 1$, the
	group $\ZmodUnits{p^{e}}$ is cyclic of order $\varphi(p^{e}) =
		p^{e-1}(p-1)$ because $p^{e}$ has a primitive root (Theorem~\ref{thm:6.2.6}). For
	$m = 2p^{e}$ we have $\gcd(2, p^{e}) = 1$, so the Chinese remainder
	theorem gives a ring isomorphism $\mathbb{Z}_{2p^{e}} \cong \mathbb{Z}_2
		\times \mathbb{Z}_{p^{e}}$ and hence a group isomorphism
	$\ZmodUnits{2p^{e}} \cong \ZmodUnits{2} \times
		\ZmodUnits{p^{e}} \cong \ZmodUnits{p^{e}}$, which is
	cyclic of order $\varphi(p^{e}) = \varphi(2)\varphi(p^{e}) =
		\varphi(2p^{e})$. For $e \ge 3$, the proof of
	Theorem~\ref{thm:6.4.2} shows that $5$ has order $2^{e-2}$:
	$\delta_{2^e}(5) \mid 2^{e-2}$, while $v_2\big(5^{2^{e-3}} - 1\big) =
		e-1 < e$ gives $5^{2^{e-3}} \not\equiv 1 \pmod{2^e}$. Moreover, $1, 5,
		\ldots, 5^{2^{e-2}-1}$ are all congruent to $1 \pmod{4}$, while $-1 \equiv 3
		\pmod{4}$, so $-1 \notin \langle 5 \rangle$; hence $\{\pm 5^{k}\}$
	gives $2^{e-1} = \varphi(2^{e})$ distinct units. The subgroups
	$\langle-1\rangle$ and $\langle5\rangle$ intersect trivially (a common
	element is a power of $5$ congruent to $\pm1$, and
	$-1\notin\langle5\rangle$), so the product map
	$\langle-1\rangle\times\langle5\rangle\to\ZmodUnits{2^{e}}$,
	$(u,v)\mapsto uv$ (a homomorphism since the group is abelian), is
	injective, and the count $2\cdot2^{e-2}=2^{e-1}=\varphi(2^{e})$ makes
	it bijective; therefore
	$\ZmodUnits{2^{e}} \cong C_2 \times C_{2^{e-2}}$.

\end{proof}
\begin{theorem}\label{thm:6.6.4}
	For the prime factorization $m = p_1^{e_1} \cdots p_k^{e_k}$ of a positive integer $m$ (allowing the empty product when $m=1$), we have $\ZmodUnits{m} \cong
		\ZmodUnits{p_1^{e_1}} \times \cdots \times
		\ZmodUnits{p_k^{e_k}}$; combined with Theorem~\ref{thm:6.6.3}, this
	determines the structure of $\ZmodUnits{m}$ for every $m$.

\end{theorem}

\begin{proof} By the Chinese remainder theorem, the ring
	isomorphism $\mathbb{Z}_m \cong \prod_{i=1}^{k}\mathbb{Z}_{p_i^{e_i}}$
	restricts to a group isomorphism on units, since a residue is coprime
	to $m$ if and only if it is coprime to each $p_i^{e_i}$.

\end{proof}
\begin{theorem}\label{thm:6.6.5}\statementitems

	\begin{enumerate}[label={\alph*)}]
		\item Let $f:G\to H$ be a homomorphism of groups with kernel $K$. If
		      $h=f(g)$ with $g\in G$, then $\{x\in G\mid f(x)=h\}=gK$.
		      Consequently, if $G$ is finite, then $|G|=|K|\cdot|f(G)|$.
		\item Let $G$ be a cyclic group of finite order $N$ with generator $g$,
		      let $m\in\mathbb{Z}^+$, and put $d=\gcd(m,N)$. Then the elements
		      $x\in G$ with $x^m=1$ form the subgroup $\langle g^{N/d}\rangle$,
		      which has exactly $d$ elements. In particular, for every divisor $n$
		      of $N$, the subgroup $\{x\in G\mid x^n=1\}$ is the unique subgroup of
		      $G$ of order $n$.
	\end{enumerate}

\end{theorem}

\begin{proof}

	\begin{enumerate}[label={\alph*)}]
		\item Since $f$ is a homomorphism, $f(x)=f(g)$ holds if and only if
		      $f(g^{-1}x)=1$, that is, $g^{-1}x\in K$, or equivalently $x\in gK$.
		      Thus the sets $\{x\in G\mid f(x)=h\}$ with $h\in f(G)$ are left
		      cosets of $K$; they partition $G$, and each has $|K|$ elements,
		      because $k\mapsto gk$ is a bijection $K\to gK$. Counting the elements
		      of $G$ gives $|G|=|K|\cdot|f(G)|$.
		\item The element $g$ has order $N$, so every $x\in G$ equals $g^j$ for
		      exactly one $j$ with $0\le j<N$ (Theorem~\ref{thm:6.1.1}(b), which holds
		      for any element of finite order). By Theorem~\ref{thm:6.1.1}(c),
		      $x^m=g^{jm}=1$ if and only if $N\mid jm$, that is,
		      $\frac{N}{d}\mid j\cdot\frac{m}{d}$. Since $\gcd(\frac{N}{d},\frac{m}{d})=1$
		      (Theorem~\ref{thm:2.2.13}), Euclid's lemma (Theorem~\ref{thm:2.2.8})
		      shows that this holds exactly when $\frac{N}{d}\mid j$. The integers $j$
		      with $0\le j<N$ and $\frac{N}{d}\mid j$ are $i\frac{N}{d}$ with
		      $0\le i<d$, and the corresponding powers $g^{iN/d}$ are exactly the
		      elements of $\langle g^{N/d}\rangle$, since every exponent of $g$ may be
		      reduced modulo $N$ (Theorem~\ref{thm:6.1.1}(a)). Hence the solution set
		      is this subgroup and has $d$ elements. If $n\mid N$, then $d=n$; and if
		      $H$ is any subgroup of order $n$, Lagrange's theorem
		      (Theorem~\ref{thm:6.6.2}) and Theorem~\ref{thm:6.1.1}(c) give $x^n=1$ for
		      every $x\in H$, so $H$ is contained in the solution set, which also has
		      $n$ elements; therefore $H$ equals it.
	\end{enumerate}

\end{proof}
\begin{altproof}[Alternative proof of Theorem~\ref{thm:6.2.6}, Case 4] We reprove that
	$\ZmodUnits{m}$ is cyclic (equivalently, a primitive root
	modulo $m$ exists) if and only if $m = 1, 2, 4, p^{e}$, or
	$2 p^{e}$ ($p$ an odd prime, $e \ge 1$) by assembling the general
	modulus from the prime-power cases already established in Theorem~\ref{thm:6.2.6} (Cases 1--3), which provides an alternative to its Case 4. By
	Theorem~\ref{thm:6.6.4}, $\ZmodUnits{m}
		\cong \prod_{i=1}^{k}\ZmodUnits{p_i^{e_i}}$. A finite product of cyclic groups $C_{n_1} \times \cdots \times
		C_{n_k}$ is cyclic if and only if the orders $n_i$ are pairwise
	coprime. Indeed, the order of an element $(x_1, \ldots, x_k)$ is the
	$\lcm$ of the orders of its components, because
	$(x_1,\ldots,x_k)^r$ is the identity exactly when $x_i^r=1$ for every
	$i$, that is (Theorem~\ref{thm:6.1.1}(c), applied in each factor), exactly
	when every component order divides $r$; so the largest
	order attained by any element is $\lcm(n_1, \ldots, n_k)$ (take each $x_i$ to be a generator of $C_{n_i}$); the product is
	cyclic exactly when some element has order equal to the group order
	$\prod_{i=1}^{k}n_i$, that is, when $\lcm(n_1, \ldots, n_k) =
		\prod_{i=1}^{k}n_i$, which, by comparison of $p$-adic valuations, holds if and
	only if the $n_i$ are pairwise coprime. By the prime-power cases of Theorem~\ref{thm:6.2.6}, each
	$\ZmodUnits{p_i^{e_i}}$ is cyclic when $p_i$ is odd or
	$p_i^{e_i} \in \{2, 4\}$ (Cases 1--2, where a primitive root exists),
	while $\ZmodUnits{2^{e}}$ ($e \ge 3$) is noncyclic
	(Case 3). Thus for $\ZmodUnits{m}$ to be cyclic we need,
	first, every factor $\ZmodUnits{p_i^{e_i}}$ to be cyclic,
	which excludes any factor with $p_i^{e_i} = 2^{e}$ and $e \ge 3$; hence
	the power of $2$ dividing $m$ is $1$, $2$, or $4$. Second, the factor
	orders must be pairwise coprime; but $\varphi(p^{e})$ is even for every
	odd prime $p$ and $\varphi(4) = 2$ is even, so at most one factor may
	have even order, which forbids two distinct odd prime factors and
	forbids an odd prime factor when $4 \mid m$. Combining the two
	conditions, we see that $\ZmodUnits{m}$ is cyclic exactly when $m = 1, 2,
		4, p^{e}$, or $2 p^{e}$ ($p$ an odd prime).

\end{altproof}
\subsection{Group-Theoretic Results and a Generalization of Wilson's Theorem}

\begin{theorem}[Cauchy's theorem]\label{thm:6.7.1} If $q$ is a
	prime and $q \mid |G|$ for a finite group $G$, then there exists $g
		\in G$ whose order equals $q$.

\end{theorem}

\begin{proof} Consider the set $S = \{(a_1, \ldots, a_q)
		\in G^{q} \mid a_1 a_2 \cdots a_q = 1\}$. Choosing $a_1, \ldots,
		a_{q-1}$ freely and setting $a_q = (a_1 \cdots a_{q-1})^{-1}$ gives
	$|S| = |G|^{q-1}$. The cyclic group $\mathbb{Z}/q\mathbb{Z}$ acts on $S$ by
	cyclic rotation $(a_1, \ldots, a_q) \mapsto (a_2, \ldots, a_q, a_1)$;
	this preserves $S$ because $a_1 \cdots a_q = 1$ implies $a_2 \cdots
		a_q a_1 = a_1^{-1}(a_1 \cdots a_q)a_1 = 1$. Every orbit of this action
	has size $1$ or $q$, and the size-$1$ orbits are exactly the constant
	tuples $(a, \ldots, a)$ with $a^{q} = 1$. Indeed, no nonzero shift fixes
	a nonconstant tuple. If the shift
	by some $k$ with $k \not\equiv 0 \pmod{q}$ fixes a tuple, then, since $q$
	is prime, $\gcd(k, q) = 1$, so the multiples $k, 2k, \ldots, qk$ run
	through all residues modulo $q$. Hence every shift fixes that tuple,
	which forces $a_1 = a_2 = \cdots = a_q$. Conversely, a constant
	tuple is fixed by every shift. Thus a nonconstant tuple is moved by
	every nonzero shift, so its $q$ shifts are pairwise distinct and its
	orbit has exactly $q$ elements, while a constant tuple $(a, \ldots, a)$
	lies in $S$ precisely when $a^{q} = 1$ and forms an orbit of size $1$.
	(No orbit--stabilizer or Lagrange argument is needed here.) If $F$ is the
	number of fixed points, then $|S| \equiv F \pmod{q}$; as $q \mid |G|$,
	we have $q \mid |G|^{q-1} = |S|$, so $q \mid F$. Since $(1, \ldots,
		1)$ is a fixed point, $F \ge 1$, hence $F \ge q \ge 2$, so there
	exists $a \neq 1$ with $a^{q} = 1$. As $q$ is prime, $a$ has order
	exactly $q$.

\end{proof}
\begin{theorem}\label{thm:6.7.2} A finite abelian group $G$ in which
	every element has order $1$ or $2$ has $|G| = 2^{k}$ for some
	$k\in\mathbb{Z}_{\ge0}$ and is isomorphic to $(C_2)^{k}$.

\end{theorem}

\begin{proof} If an odd prime $q$ divides $|G|$, then by
	Theorem~\ref{thm:6.7.1} there is an element of order $q > 2$, a contradiction;
	hence $|G| = 2^{k}$. We prove $G \cong (C_2)^{k}$ by induction on
	$k$. When $k = 0$, we have $G = \{1\}$. For the inductive step, let
	$k\ge1$. Then $|G|>1$, so choose $b_1\neq
		1$; then $b_1^2=1$ and $H = \{1, b_1\} \cong C_2$. Put
	$Q=G/H=\{xH\mid x\in G\}$ and define $(xH)(yH)=xyH$.
	This multiplication is well defined: for $h_1,h_2\in H$, commutativity gives
	$(xh_1)(yh_2)H=xy(h_1h_2)H=xyH$.
	Associativity and commutativity follow from those of $G$; the identity is
	$H$, and the inverse of $xH$ is $x^{-1}H$. Thus $Q$ is an abelian group.
	By Theorem~\ref{thm:6.6.2}, it has $|G|/|H|=2^{k-1}$ elements.
	Since $(xH)^2=x^2H=H$, every element has order $1$ or $2$.
	The induction hypothesis gives $Q\cong(C_2)^{k-1}$. Choose lifts $b_2,\ldots,b_k$ of generators $b_2H,\ldots,b_kH$ of $Q$. The map $(C_2)^k\to G$ given by $(\varepsilon_1,\ldots,\varepsilon_k)\mapsto b_1^{\varepsilon_1}\cdots b_k^{\varepsilon_k}$ is surjective: the lifted generators produce every coset of $H$, and multiplication by $b_1$ interchanges the two elements of each coset. Both groups have $2^k$ elements, so the surjection is bijective; consequently, the displayed representation is unique. The map is a homomorphism (as $G$ is
	abelian), so $G \cong (C_2)^{k}$.

\end{proof}
\begin{theorem}[Gauss's generalization of Wilson's theorem]\label{thm:6.7.3} Let $T$ be a reduced residue system modulo $m$; then
	$\prod_{a \in T} a \equiv -1 \pmod{m}$ when $m = 1, 2, 4, p^{e}$, or $2
		p^{e}$ ($p$ an odd prime, $e \ge 1$), and $\prod_{a \in T} a \equiv 1 \pmod{m}$ otherwise.

\end{theorem}

\begin{proof} First, consider Case 1, in which a primitive root exists.
	Let $g$ be a primitive root, so the classes represented by $T$ are $g^{0}, \ldots,
		g^{\varphi(m)-1}$, and $\prod_{a \in T} a \equiv
		g^{\varphi(m)(\varphi(m)-1)/2} \pmod{m}$. For $m \neq 1, 2$,
	the totient $\varphi(m)$ is even and $\varphi(m)-1$ is odd, so
	$\frac{\varphi(m)(\varphi(m)-1)}{2} \equiv \frac{\varphi(m)}{2}
		\pmod{\varphi(m)}$, giving $\prod_{a \in T}a\equiv g^{\varphi(m)/2}\pmod{m}$. For $m>2$, we have $-1\not\equiv1\pmod{m}$ and $(-1)^2\equiv1$, so $-1$ has order $2$; a cyclic group has exactly one element of order $2$, namely $g^{\varphi(m)/2}$: the order of $g^{i}$ is $\varphi(m)/\gcd(i,\varphi(m))$ (Theorem~\ref{thm:6.1.1}(d)), which equals $2$ exactly when $\gcd(i,\varphi(m))=\varphi(m)/2$, and for $0\le i<\varphi(m)$ this forces $i=\varphi(m)/2$. Hence the product is $-1$ modulo $m$. For $m = 1, 2$,
	we have $\prod_{a \in T} a \equiv 1 \equiv -1 \pmod{m}$. Now consider Case 2, in which no
	primitive root exists. Let $G = \ZmodUnits{m}$ and $T_0 = \{x \in
		G \mid x^{2} = 1\}$. Elements $b \notin T_0$ pair with $b^{-1} \neq b$
	(product $1$), so $\prod_{x \in G} x = \prod_{x \in T_0} x$. Because $G$
	is abelian, $T_0$ is a subgroup: it contains $1$, each of its elements is
	its own inverse, and $(xy)^2=x^2y^2=1$ for $x,y\in T_0$. Every element
	of $T_0$ has order $1$ or $2$, so by
	Theorem~\ref{thm:6.7.2}, $T_0 \cong (C_2)^{k}$. Here $m>2$, so $-1\in T_0$ is not the identity and $k\ge1$. Fixing $v \in T_0 \setminus
		\{1\}$ and pairing $\{x, xv\}$ partitions $T_0$ into
		$|T_0|/2 = 2^{k}/2 = 2^{k-1}$ pairs, each contributing
	$x\cdot xv = v$ since $x^{2} = 1$; hence $\prod_{x \in T_0} x =
		v^{2^{k-1}}$, which is $1$ when $k \ge 2$ and $v$ when $k = 1$. By
	Theorem~\ref{thm:6.2.6}, ``no primitive root exists'' means $m \notin \{1, 2,
		4, p^{e}, 2p^{e}\}$; writing $m = 2^{c} t$ with $t$ odd, either $t = 1$
	and then $c \ge 3$, or $t$ has two distinct prime factors and
	separating them gives two coprime factors greater than $1$ with even totients
	(the factor $2^{c}$ may be attached to either of them), or $t = p^{e}$ and then $c \ge 2$
	(as $m \neq p^{e}, 2p^{e}$), so $m = 2^{c} \cdot p^{e}$ with both
	$\varphi(2^{c})$ and $\varphi(p^{e})$ even. Thus $m$ falls into one of
	two cases, and we show that $k \ge 2$ in each. First consider Case 2a, where $m = 2^{e}$ and $e \ge 3$. We show directly that $x^{2}
		\equiv 1 \pmod{2^{e}}$ has exactly the four solutions $x \equiv \pm
			1,\ \pm 1 + 2^{e-1} \pmod{2^{e}}$, so $|T_0| = 4$ and $k = 2$.
	Any solution is odd. If $x\equiv1\pmod4$, write $x+1=2u$
	with $u$ odd. Since $u$ is coprime to $2^{e-1}$,
	\[
		2^e\mid(x-1)(x+1)
		\iff 2^{e-1}\mid u(x-1)
		\iff 2^{e-1}\mid x-1.
	\]
	This gives $x\equiv1,\ 1+2^{e-1}\pmod{2^e}$.
	If $x\equiv-1\pmod4$, write $x-1=2u$ with $u$ odd; the same
	cancellation gives $2^{e-1}\mid x+1$, hence
	$x\equiv-1,\ -1+2^{e-1}\pmod{2^e}$.
	These divisibility arguments include $x=1$ and $x=-1$ directly,
	without applying a nonzero-factor valuation formula to a zero factor.
	Since $e \ge 3$ gives $2^{e-1} \ge 4$, the four residues $1, -1,
			1+2^{e-1}, -1+2^{e-1}$ are pairwise distinct modulo $2^{e}$, so $|T_0| =
		4$, that is, $k = 2 \ge 2$. In Case 2b, there are two coprime factors $m_1, m_2 > 1$ with even
	totients. Here $m = m_1 m_2$ with
	$\gcd(m_1, m_2) = 1$, $2 \mid \varphi(m_1)$, $2 \mid \varphi(m_2)$;
	by the Chinese remainder theorem the elements $a \equiv -1 \pmod{m_1}, a \equiv 1 \pmod{m_2}$
	and $b \equiv 1 \pmod{m_1}, b \equiv -1 \pmod{m_2}$ give $\{1, a, b,
		ab\} \cong C_2 \times C_2 \le T_0$, so $k \ge 2$. In both cases $k
		\ge 2$, hence $\prod_{a \in T} a \equiv 1 \pmod{m}$.

\end{proof}
\subsection{Power Sums, Bernoulli Numbers, and the von Staudt--Clausen Theorem}

\begin{theorem}\label{thm:6.8.1} Let $p$ be a prime and $m \ge 1$ an integer. Then $\sum_{a=1}^{p-1} a^{m} \equiv -1 \pmod{p}$ if $(p-1) \mid m$, and $\sum_{a=1}^{p-1} a^{m} \equiv 0 \pmod{p}$ otherwise.

\end{theorem}

\begin{proof} Let $g$ be a primitive root modulo $p$ (Theorem~\ref{thm:6.2.6}). Every nonzero residue modulo $p$ is represented exactly once by $g^j$ with $0\le j\le p-2$. Hence
	\[
		\sum_{a=1}^{p-1}a^m\equiv\sum_{j=0}^{p-2}g^{jm}\pmod{p}.
	\]
	If $(p-1)\mid m$, every summand is $1$, so the sum is $p-1\equiv-1\pmod{p}$. If $(p-1)\nmid m$, then $g^m\not\equiv1\pmod{p}$ by Theorem~\ref{thm:6.1.1}(c), while
	\[
		(g^m-1)\sum_{j=0}^{p-2}g^{jm}=g^{m(p-1)}-1\equiv0\pmod{p}.
	\]
	Since $g^m-1$ is invertible modulo $p$, the sum is $0\pmod{p}$.

\end{proof}
\begin{altproof}[Alternative proof of Theorem~\ref{thm:6.8.1}] If $(p-1) \mid m$, then for each $a \in \{1,\ldots,p-1\}$ Fermat's little theorem (Theorem~\ref{thm:3.3.2}) gives $a^{m} = (a^{p-1})^{m/(p-1)} \equiv 1 \pmod{p}$, so the sum is congruent to $p-1 \equiv -1 \pmod{p}$. If $(p-1) \nmid m$, put $d = \gcd(m, p-1)$, a proper divisor of $p-1$. If $a^{m} \equiv 1 \pmod{p}$ held for every $a \in \{1,\ldots,p-1\}$, then choosing integers $u,v$ with $um + v(p-1) = d$ (B{\'e}zout's identity, Theorem~\ref{thm:2.2.5}) would give $a^{d} = (a^{m})^{u} (a^{p-1})^{v} \equiv 1 \pmod{p}$ for all such $a$ (using Theorem~\ref{thm:3.3.2}), so $x^{d} - 1$ would have $p-1 > d$ roots modulo $p$, contradicting the polynomial root bound (Theorem~\ref{thm:3.6.2}). Hence some $g \in \{1,\ldots,p-1\}$ has $g^{m} \not\equiv 1 \pmod{p}$. As $a \mapsto ga \bmod p$ permutes $\{1,\ldots,p-1\}$, we have $\sum_{a} a^{m} \equiv \sum_{a} (ga)^{m} = g^{m} \sum_{a} a^{m} \pmod{p}$, so $(g^{m}-1)\sum_{a} a^{m} \equiv 0 \pmod{p}$; since $g^{m}-1 \not\equiv 0$, it follows that $\sum_{a} a^{m} \equiv 0 \pmod{p}$.

\end{altproof}
\begin{definition}[Bernoulli numbers]\label{def:6.8.1} We work in the formal
	power-series ring $\mathbb{Q}[[t]]$ of Definition~\ref{def:3.5.2}. For
	$a\in\mathbb{Q}$, the formal exponential series is
	$e^{at}=\sum_{k\ge0}a^kt^k/k!$, with $a^0=1$ for every $a$, so that $e^{0t}=1$.
	By the binomial theorem, the coefficient of $t^n$ in $e^{at}e^{bt}$ is
	$\sum_{i=0}^{n}\frac{a^ib^{n-i}}{i!\,(n-i)!}=\frac{(a+b)^n}{n!}$, so
	$e^{at}e^{bt}=e^{(a+b)t}$ for all $a,b\in\mathbb{Q}$; in particular
	$e^{t}e^{-t}=1$. The series $D(t)=\sum_{k\ge0}t^k/(k+1)!$ satisfies
	$tD(t)=e^{t}-1$ and has constant term $1$, so it is invertible in
	$\mathbb{Q}[[t]]$ (Theorem~\ref{thm:3.5.10}(b)). The Bernoulli numbers $B_0, B_1, B_2, \ldots \in \mathbb{Q}$
	are the coefficients of its inverse:
	\[
		\frac{t}{e^{t} - 1}:=D(t)^{-1} = \sum_{m=0}^{\infty} B_m \frac{t^{m}}{m!}.
	\]

\end{definition}
\begin{theorem}\label{thm:6.8.2} $B_0 = 1$ and $\sum_{j=0}^{m} \binom{m+1}{j} B_j = 0$ for every $m \ge 1$; consequently $B_1 = -\tfrac{1}{2}$, and $B_m = 0$ for every odd $m \ge 3$.

\end{theorem}

\begin{proof} From $(e^{t}-1)\sum_{j \ge 0} B_j t^{j}/j! = tD(t)D(t)^{-1}=t$ and $e^{t}-1 = \sum_{k \ge 1} t^{k}/k!$, comparing coefficients of $t^{m+1}$ gives $\sum_{j=0}^{m} \frac{B_j}{j!\,(m+1-j)!} = [\,m=0\,]$; multiplying by $(m+1)!$ yields $\sum_{j=0}^{m} \binom{m+1}{j} B_j = (m+1)!\,[\,m=0\,]$, that is, $B_0 = 1$ for $m = 0$ and the stated recurrence for $m \ge 1$. Taking $m = 1$ gives $B_0 + 2 B_1 = 0$, so $B_1 = -\tfrac{1}{2}$. Finally, substituting $-t$ for $t$ in $tD(t)=e^{t}-1$, which preserves sums and products (Theorem~\ref{thm:3.5.10}(c)), gives $-tD(-t)=e^{-t}-1$, and hence
	\[
		tD(-t)=1-e^{-t}=e^{-t}(e^{t}-1)=e^{-t}\,tD(t).
	\]
	Multiplication by $t$ shifts every coefficient up by one place, so the factor $t$ may be cancelled, giving $D(-t)=e^{-t}D(t)$. The same substitution carries the identity $D(t)D(t)^{-1}=1$ to
	\[
		D(-t)\sum_{m\ge0}(-1)^mB_m\frac{t^m}{m!}=1,
	\]
	so the inverse of $D(-t)$ is $\sum_{m\ge0}(-1)^mB_mt^m/m!$. Inverting $D(-t)=e^{-t}D(t)$ and using $e^{t}e^{-t}=1$ and $e^{t}=1+tD(t)$, we obtain
	\[
		\sum_{m\ge0}(-1)^mB_m\frac{t^m}{m!}=e^{t}D(t)^{-1}=\big(1+tD(t)\big)D(t)^{-1}=\sum_{m\ge0}B_m\frac{t^m}{m!}+t.
	\]
	For odd $m\ge3$, comparing the coefficients of $t^m$ gives $-B_m=B_m$, so $B_m = 0$.

\end{proof}
\begin{theorem}[Faulhaber's formula]\label{thm:6.8.3} For all integers $m \ge 1$ and $N \ge 1$,
	\[
		\sum_{i=0}^{N-1} i^{m} = \frac{1}{m+1} \sum_{j=0}^{m} \binom{m+1}{j} B_j\, N^{\,m+1-j}.
	\]

\end{theorem}

\begin{proof} We work in $\mathbb{Q}[[t]]$ with the notation of Definition~\ref{def:6.8.1}. Since $e^{rt}e^{t}=e^{(r+1)t}$, the sum telescopes:
	\[
		(e^{t}-1)\sum_{r=0}^{N-1}e^{rt}
		=\sum_{r=0}^{N-1}\big(e^{(r+1)t}-e^{rt}\big)
		=e^{Nt}-1.
	\]
	Put $D_N(t)=\sum_{a\ge0}N^{a+1}t^a/(a+1)!$, so that $tD_N(t)=e^{Nt}-1$. Since $e^{t}-1=tD(t)$, cancelling the factor $t$ gives $D(t)\sum_{r=0}^{N-1}e^{rt}=D_N(t)$, and multiplying by $D(t)^{-1}$ gives
	\[
		\sum_{r=0}^{N-1}e^{rt}
		=\left(\sum_{a\ge0}\frac{N^{a+1}t^a}{(a+1)!}\right)
		 \left(\sum_{j\ge0}\frac{B_jt^j}{j!}\right).
	\]
	The left-hand side has
	constant term $N$ and coefficient $\frac{1}{m!}\sum_{r=0}^{N-1}r^m$
	at $t^m$ for $m\ge1$; the term $r=0$ is $e^{0t}=1$ and contributes nothing to these coefficients.
	Comparing coefficients of $t^m$ for $m\ge1$ and multiplying by $m!$ gives
	\[
		\sum_{r=0}^{N-1}r^m
		=\sum_{j=0}^m\frac{1}{m+1-j}\binom mj B_jN^{m+1-j}
		=\frac{1}{m+1}\sum_{j=0}^m\binom{m+1}{j}B_jN^{m+1-j},
	\]
	since $\binom mj/(m+1-j)=m!/[j!(m+1-j)!]
	=\binom{m+1}{j}/(m+1)$. This proves the formula.

\end{proof}
\begin{theorem}[von Staudt--Clausen theorem]\label{thm:6.8.4} For every even integer $m \ge 2$,
	\[
		B_m + \sum_{\substack{p \ \mathrm{prime} \\ (p-1) \mid m}} \frac{1}{p} \in \mathbb{Z}.
	\]

\end{theorem}

\begin{proof} Work in $\mathbb{Z}_{(p)}$ of Definition~\ref{def:3.5.1},
	whose elements we call $p$-integers. In terms of the rational valuation
	(Definition~\ref{def:4.5.2}), $x$ is a $p$-integer exactly when
	$v_p(x)\ge0$, and $x\equiv y\pmod p$ exactly when $v_p(x-y)\ge1$.
	Here we retain $v_p(0)=+\infty$, so these criteria include zero.
	Thus this terminology uses the same ring and congruence as before. A rational number that is a $p$-integer for every prime $p$ lies in $\mathbb{Z}$: writing such a number as $a/b$ in lowest terms with $b \ge 1$, we see that if a prime $p$ divided $b$ then $v_p(b) \ge 1$ while $v_p(a) = 0$ (as $\gcd(a, b) = 1$), so $v_p(a/b) = -v_p(b) < 0$, contradicting $p$-integrality; hence $b = 1$ and $a/b \in \mathbb{Z}$. We prove the following claim: for every prime $p$ and even $m \ge 2$, the product $p B_m$ is a $p$-integer and $p B_m \equiv -[\,(p-1) \mid m\,] \pmod{p}$.

	We argue by strong induction on the even index $m$. Since $B_0 = 1$, $B_1 = -\tfrac{1}{2}$, and $B_j = 0$ for odd $j \ge 3$ (Theorem~\ref{thm:6.8.2}), the induction hypothesis (namely the claim for even indices below $m$) gives $v_p(B_j) \ge -1$ for all $j < m$. In Faulhaber's formula (Theorem~\ref{thm:6.8.3}) with $N = p$, the $j = m$ term is $\frac{1}{m+1}\binom{m+1}{m} B_m\, p = B_m\, p$, so
	\[
		p B_m = \sum_{a=1}^{p-1} a^{m} - \sum_{j=0}^{m-1} T_j, \qquad T_j = \frac{1}{m+1-j} \binom{m}{j} B_j\, p^{\,m+1-j}.
	\]
	We claim that $v_p(T_j)\ge1$ for $0\le j<m$.
	If $B_j=0$, then $T_j=0$ and the claim is immediate.
	For the remaining indices, $B_j$ and $T_j$ are nonzero, so their
	valuations are finite. Put $s=m+1-j\ge2$; since
	$\binom{m}{j}\in\mathbb{Z}$ and $v_p(B_j)\ge-1$,
	\[
		v_p(T_j) \ge v_p(B_j) + (m+1-j) - v_p(m+1-j) \ge -1 + \big(s - v_p(s)\big).
	\]
	For any prime $p$ and integer $s \ge 2$, we have $v_p(s) \le s-2$ except when $(p,s) = (2,2)$, where $v_2(2) = 1$.
	Indeed, $v_p(s) \ge s-1$ forces $p^{s-1} \le s$.
	When $s = 2$, this gives $p \le 2$, so $p = 2$.
	When $s \ge 3$, we have $p^{s-1} \ge 2^{s-1} > s$.
	The last inequality follows by induction: $2^{2} = 4 > 3$, and
	$2^{s} = 2 \cdot 2^{s-1} > 2s \ge s+1$ for $s \ge 3$.
	Thus $p^{s-1} \le s$ holds only when $p = s = 2$.
	Except in the case $(p,s)=(2,2)$, we have $s - v_p(s) \ge 2$, so $v_p(T_j) \ge 1$. In the exceptional case $p=2$, $s=2$, the index $j=m-1$ is odd.
	For $m\ge4$ this is a zero Bernoulli term, already handled above;
	for $m=2$,
	we have $T_1=\frac12\binom21 B_1\,2^2=-2$, whose $2$-adic valuation is $1$. Thus $v_p(T_j) \ge 1$ in all cases, so $\sum_{j} T_j$ is a $p$-integer congruent to $0 \pmod{p}$. Since $\sum_{a=1}^{p-1} a^{m} \in \mathbb{Z}$, it follows that $p B_m$ is a $p$-integer and $p B_m \equiv \sum_{a=1}^{p-1} a^{m} \pmod{p}$.
	By Theorem~\ref{thm:6.8.1},
	$\sum_{a=1}^{p-1} a^{m} \equiv -[\,(p-1) \mid m\,] \pmod{p}$. This proves the induction claim.

	The sum in the statement is finite, since $(p-1)\mid m$ implies $p\le m+1$. Now fix a prime $q$. If $(q-1) \nmid m$, then the claim gives $v_q(B_m) \ge 0$, and every term $1/p$ in the sum has $p \neq q$, hence is a $q$-integer; so $B_m + \sum 1/p$ is a $q$-integer. If $(q-1) \mid m$, then $q B_m \equiv -1 \pmod{q}$, so $q B_m = -1 + q\alpha$ with $\alpha$ a $q$-integer, whence $B_m + \tfrac{1}{q} = \alpha$ is a $q$-integer while the other terms $1/p$ ($p \neq q$) are $q$-integers; again $B_m + \sum 1/p$ is a $q$-integer. As this holds for every prime $q$, we conclude that $B_m + \sum_{(p-1) \mid m} 1/p \in \mathbb{Z}$.

\end{proof}
\subsection{Symmetric-Function Congruences and Fermat Quotients}

\begin{definition}[the set $t(m)$]\label{def:6.9.1}
	For an integer $m\ge2$, let $t(m)$ denote the set of integers in
	$\{1,2,\ldots,m-1\}$ that are coprime to $m$. Thus $t(m)$ is the
	reduced residue system of least positive representatives and has
	$\varphi(m)$ elements.

\end{definition}
\begin{theorem}[Bauer's congruential identity]\label{thm:6.9.1} Let
	$f_m(x) = \prod_{t \in t(m)}(x - t)$.

	\begin{enumerate}[label={\alph*)}]
		\item If $p$ is an odd prime with $p\mid m$, say $p^{a} \parallel m$, then $f_m(x)
			      \equiv (x^{p-1} - 1)^{\varphi(m)/(p-1)} \pmod{p^{a}}$; in particular
		      $f_{p^{a}}(x) \equiv (x^{p-1} - 1)^{p^{a-1}} \pmod{p^{a}}$.
		\item If $2 \mid m$ and $m > 2$, with $2^{a} \parallel m$, then $f_m(x)
			      \equiv (x^{2} - 1)^{\varphi(m)/2} \pmod{2^{a}}$; in particular, for
		      $a > 1$, we have $f_{2^{a}}(x) \equiv (x^{2} - 1)^{2^{a-2}} \pmod{2^{a}}$.
	\end{enumerate}

\end{theorem}

\begin{proof} We first prove the prime-power cases by
	induction on $a$, then reduce the general case to these using the Chinese remainder theorem.

	First, consider an odd prime $p$ and the modulus $p^a$. For the base case $a =
		1$, the nonzero residues modulo $p$ are exactly the roots of $x^{p-1}
		- 1$, and both $f_p(x)$ and $x^{p-1} - 1$ are monic of degree $p-1$
	with the same roots, so $f_p(x) \equiv x^{p-1} - 1 =
			(x^{p-1}-1)^{p^{0}} \pmod{p}$. For the inductive step from $a$ to $a+1$, the
	residues in $\{1,\ldots,p^{a+1}-1\}$ coprime to $p$ are exactly $r+jp^a$ with $r\in t(p^a)$ and $0\le j<p$: each is nonzero modulo $p$ because $r$ is, and congruence modulo $p^{a+1}$ first forces the same $r$ modulo $p^a$ and then the same $j$ modulo $p$. Hence $f_{p^{a+1}}(x) =
		\prod_{r \in t(p^{a})} \prod_{j=0}^{p-1}(x - r - j p^{a})$. Writing
	$y = x - r$, we obtain $\prod_{j=0}^{p-1}(y - j p^{a}) \equiv y^{p} -
		p^{a}\big(\textstyle\sum_{j=0}^{p-1} j\big) y^{p-1} = y^{p} -
		p^{a}\cdot\frac{p(p-1)}{2}y^{p-1} \equiv y^{p} \pmod{p^{a+1}}$, since
	$p^{a}\cdot\frac{p(p-1)}{2} = p^{a+1}\cdot\frac{p-1}{2}
		\equiv 0 \pmod{p^{a+1}}$, where $\frac{p-1}{2} \in \mathbb{Z}$ because
	$p$ is an odd prime. Here every term involving at least two of the quantities $jp^a$ is
	an elementary symmetric term containing a factor $p^{ak}$ with
	$k\ge2$; since $ak\ge2a\ge a+1$, all such terms also vanish modulo
	$p^{a+1}$. Hence $f_{p^{a+1}}(x) \equiv
		\prod_{r\in t(p^a)}(x-r)^p=f_{p^a}(x)^p \pmod{p^{a+1}}$. Writing
	$f_{p^{a}}(x) = (x^{p-1}-1)^{p^{a-1}} + p^{a} h(x)$ from the
	induction hypothesis, we obtain $f_{p^{a}}(x)^{p} \equiv (x^{p-1}-1)^{p^{a}}
		\pmod{p^{a+1}}$ (the cross terms are divisible by $p^{a+1}$). Thus
	$f_{p^{a+1}}(x) \equiv (x^{p-1}-1)^{p^{a}} \pmod{p^{a+1}}$.

	Next, take $p=2$ and the modulus $2^a$ ($a \ge 2$). For the base case $a = 2$, we have $f_4(x) = (x-1)(x-3) = x^{2} - 4x + 3
		\equiv x^{2} - 1 \pmod{4}$. For the inductive step from $a$ to $a+1$ ($a \ge 2$),
	the set $t(2^{a})$ is closed under $t \mapsto 2^{a} - t$, which partitions it
	into pairs $\{r, 2^{a}-r\}$. For such a pair the four factors modulo
	$2^{a+1}$ are $x - r,\ x - (r + 2^{a}),\ x - (2^{a} - r),\ x -
			(2^{a+1} - r)$; grouping them gives $(x-r)(x-(2^{a+1}-r)) \equiv x^{2} - r^{2}$
	and $(x-(r+2^{a}))(x-(2^{a}-r)) = (x-2^{a})^{2} - r^{2} \equiv x^{2}
		- r^{2} \pmod{2^{a+1}}$, so the four factors contribute $(x^{2}-r^{2})^{2}$.
	Hence $f_{2^{a+1}}(x) \equiv
		\big(\prod_{\text{pairs}}(x^{2}-r^{2})\big)^{2} \pmod{2^{a+1}}$.
	Put $G(x)=\prod_{\text{pairs}}(x^{2}-r^{2})$. We have $G(x)\equiv f_{2^{a}}(x)\pmod{2^{a}}$. Moreover $G(x)\equiv f_{2^{a}}(x)\pmod{2}$: modulo $2$, every odd $r$ is $1$, so both polynomials are congruent to $(x^2-1)^{2^{a-2}}$. Hence $G-f_{2^a}$ has all coefficients divisible by $2^a$, while $G+f_{2^a}$ has all coefficients even. Therefore $G^2-f_{2^a}^2=(G-f_{2^a})(G+f_{2^a})$ has all coefficients divisible by $2^{a+1}$, and thus $f_{2^{a+1}}(x) \equiv f_{2^{a}}(x)^{2}
		\pmod{2^{a+1}}$; with $f_{2^{a}}(x) = (x^{2}-1)^{2^{a-2}} + 2^{a}
		h(x)$ this is $(x^{2}-1)^{2^{a-1}} \pmod{2^{a+1}}$.

	Finally, consider a general $m$. Fix a prime power $p^{a} \parallel m$.
	By the Chinese remainder theorem the reduction map $t(m) \to t(p^{a})$, $t \mapsto t \bmod
		p^{a}$, is surjective and exactly $\varphi(m)/\varphi(p^{a})$-to-one.
	Since $x - t \equiv x - (t \bmod p^{a}) \pmod{p^{a}}$, we have $f_m(x) \equiv
		\prod_{s \in t(p^{a})}(x - s)^{\varphi(m)/\varphi(p^{a})} =
		f_{p^{a}}(x)^{\varphi(m)/\varphi(p^{a})} \pmod{p^{a}}$. For odd $p$
	this equals $(x^{p-1}-1)^{p^{a-1}\cdot \varphi(m)/\varphi(p^{a})} =
		(x^{p-1}-1)^{\varphi(m)/(p-1)}$; for $p = 2$ and $m > 2$ with $a \ge 2$, it equals
	$(x^{2}-1)^{2^{a-2}\cdot\varphi(m)/\varphi(2^{a})} =
		(x^{2}-1)^{\varphi(m)/2}$. In the remaining subcase
	$2 \parallel m$ (that is, $a = 1$), we have $t(2) = \{1\}$, so
	$f_{2}(x) = x - 1$ and $\varphi(2) = 1$; the reduction then gives
	$f_m(x) \equiv f_{2}(x)^{\varphi(m)/\varphi(2)} = (x-1)^{\varphi(m)}
		\pmod{2}$, and since $x + 1 \equiv x - 1 \pmod{2}$ we get
	$(x^{2}-1)^{\varphi(m)/2} = \big((x-1)(x+1)\big)^{\varphi(m)/2} \equiv
		(x-1)^{\varphi(m)} \pmod{2}$, so $f_m(x) \equiv (x^{2}-1)^{\varphi(m)/2}
		\pmod{2}$ as claimed.

\end{proof}
\begin{theorem}\label{thm:6.9.2} Let $p$ be an odd prime and let $a\ge1$ be an integer. Let $A_{k}$
	denote the $k$-th elementary symmetric sum of the elements of
	$t(p^{a})$ (the sum of all products of $k$ distinct elements). If $\nu\ge0$ is an integer with $2\nu+1\le\varphi(p^{a})$ and $2\nu$ is not a multiple of $p-1$, then $A_{2\nu+1} \equiv 0 \pmod{p^{2a}}$.

\end{theorem}

\begin{proof} Let $m = p^{a}$ and $n = \varphi(p^{a}) =
		p^{a-1}(p-1)$, which is even. Since $t(m)$ is closed under $t \mapsto
		m - t$ (with $m - t \neq t$), group $f_m(x) = \prod_{t \in
			t(m)}(x-t)$ into $n/2$ pairs: for each pair, $(x-t)(x-(m-t)) = x^{2}
		- m x + t\cdot(m-t) = u + b_t$, where $u = x^{2} - m x$ and $b_t = t\cdot(m-t)
		\equiv -t^{2} \pmod{m}$. Thus $f_m(x) = \prod_{\text{pairs}}(u + b_t)
		= \sum_{k} C_k\, u^{n/2 - k}$ with $C_k = e_k(\{b_t\})$ the
	elementary symmetric sums of the $n/2$ values $b_t$.

	In the expansion $u^{j} = (x^{2} - p^{a} x)^{j} = \sum_{i}
		\binom{j}{i}(-p^{a})^{i} x^{2j - i}$, an odd power of $x$ requires $i$ to be odd. Thus the possibilities are $i=1,3,5,\ldots$; modulo $p^{2a}$ only $i=1$ survives, because every remaining odd-indexed term has $i\ge3$ and hence contains the factor $p^{ai}$, which is divisible by $p^{3a}$ and therefore by $p^{2a}$. The surviving term contributes $-j\,p^{a} x^{2j-1}$. If we write
	$[x^{k}] g(x)$ for the coefficient of $x^{k}$ in a polynomial $g$, then only
	$j = \mu+1$ contributes to the coefficient of $x^{2\mu+1}$, where $\mu = n/2 - \nu - 1$, and $[x^{2\mu+1}] f_m(x) \equiv -(\mu+1)\,p^{a}\, C_{\nu}
		\pmod{p^{2a}}$. Since $f_m(x) = \sum_k (-1)^{k} A_k x^{n-k}$ and $n -
		(2\nu+1) = 2\mu+1$, we have $[x^{2\mu+1}] f_m(x) = -A_{2\nu+1}$, so
	$A_{2\nu+1} \equiv (\mu+1)\,p^{a}\, C_{\nu} \pmod{p^{2a}}$.

	It remains to prove $C_{\nu} \equiv 0 \pmod{p^{a}}$. By Bauer's identity
	(Theorem~\ref{thm:6.9.1}), $f_m(x) \equiv (x^{p-1}-1)^{p^{a-1}} \pmod{p^{a}}$;
	we compare this with $f_m(x) = \prod_{\text{pairs}}(u+b_t) = \sum_k C_k
		u^{n/2-k}$, where $u = x^2 - p^a x \equiv x^2 \pmod{p^a}$ (so
	$u^{n/2-k} \equiv x^{n-2k} \pmod{p^a}$). Every nonzero term of
	$(x^{p-1}-1)^{p^{a-1}}$ has degree equal to a multiple of $p-1$. Since
	$(p-1)\mid n$, the coefficient of $x^{n-2\nu}$ is forced to be $0
		\pmod{p^a}$ whenever $(p-1) \nmid 2\nu$, which holds by hypothesis.
	Hence $C_{\nu} \equiv 0 \pmod{p^{a}}$ directly, and therefore
	$A_{2\nu+1} \equiv (\mu+1)\,p^{a}\,C_{\nu} \equiv 0 \pmod{p^{2a}}$.

\end{proof}
\begin{theorem}[Lerch's formula]\label{thm:6.9.3} For an odd prime
	$p$, let $q_p(a) = \tfrac{a^{p-1}-1}{p}$ (the Fermat quotient,
	$\gcd(a, p)=1$) and $W_p = \tfrac{(p-1)!+1}{p}$ (the Wilson quotient).
	Then $\sum_{a=1}^{p-1} q_p(a) \equiv W_p \pmod{p}$, equivalently
	$\sum_{a=1}^{p-1} a^{p-1} \equiv (p-1)! + p \pmod{p^{2}}$.

\end{theorem}

\begin{proof} For $1\le a\le p-1$, write
	$a^{p-1}=1+p\,q_p(a)$. Then
	$\prod_{a=1}^{p-1}a^{p-1}=\big((p-1)!\big)^{p-1}$, while
	$\prod_{a=1}^{p-1}(1+p\,q_p(a))\equiv
		1+p\sum_{a=1}^{p-1}q_p(a)\pmod{p^2}$. By Wilson's theorem (Theorem~\ref{thm:3.3.4})
	$(p-1)!=-1+pW_p$, so
	$\big((p-1)!\big)^{p-1}=(-1+pW_p)^{p-1}\equiv(-1)^{p-1}+(p-1)(-1)^{p-2}pW_p\equiv
		1+pW_p\pmod{p^2}$ (using $p-1\equiv-1$). Comparing these expressions gives $1+p\sum
		q_p(a)\equiv 1+pW_p\pmod{p^2}$, so $\sum_{a=1}^{p-1}q_p(a)\equiv
		W_p\pmod{p}$. The equivalent form follows since $\sum_{a=1}^{p-1}
		a^{p-1}=(p-1)+p\sum q_p(a)\equiv(p-1)+pW_p=(p-1)!+p\pmod{p^2}$.

\end{proof}
\begin{definition}[Wieferich prime]\label{def:6.9.2} Let
	$b\in\mathbb{Z}$ and let $p$ be an odd prime with $p\nmid b$. We call
	$p$ a Wieferich prime to base $b$ if
	$b^{p-1}\equiv1\pmod{p^2}$. When no base is specified, a Wieferich prime refers to one to base $2$.

\end{definition}
\begin{theorem}\label{thm:6.9.4} Let $p$ be an odd prime and, for $p\nmid a$, define $q_p(a)=(a^{p-1}-1)/p$. Then:
	\begin{enumerate}[label={\alph*)}]
		\item $q_p(ab)\equiv q_p(a)+q_p(b)\pmod{p}$ whenever $p\nmid ab$;
		\item if $g$ is a primitive root modulo $p$, then $g$ is a primitive root modulo $p^e$ for every $e\ge2$ if and only if $q_p(g)\not\equiv0\pmod{p}$;
		\item if $g$ is a primitive root modulo $p$, then among
		      $g,g+p,\ldots,g+(p-1)p$, exactly one is not a primitive root modulo
		      $p^2$;
		\item if $g$ is a primitive root modulo $p^2$, $r\in\mathbb{Z}$, and $a\equiv g^r\pmod{p^2}$, then
		      \[
			      q_p(a)\equiv r q_p(g)\pmod{p},
			      \qquad
			      r\equiv q_p(a)q_p(g)^{-1}\pmod{p};
		      \]
		      thus $q_p$ is an additive, modulo-$p$ analogue of the index (discrete logarithm) of Subsection~\ref{subsec:6.3};
		\item the map $Q_p:\ZmodUnits{p^2}\to\mathbb{F}_p$, $Q_p(a)=q_p(a)\bmod p$, is a surjective group homomorphism to the additive group, and its kernel is the $(p-1)$-element subgroup of solutions of $x^{p-1}\equiv1\pmod{p^2}$;
		\item if $2$ is a primitive root modulo $p$, then $2$ lifts to a primitive root modulo every $p^e$ if and only if $p$ is not a Wieferich prime to base $2$.
	\end{enumerate}

\end{theorem}

\begin{proof} Write $a^{p-1}=1+pq_p(a)$ and $b^{p-1}=1+pq_p(b)$. Then
	\[
		(ab)^{p-1}=(1+pq_p(a))(1+pq_p(b))
		\equiv1+p\big(q_p(a)+q_p(b)\big)\pmod{p^2},
	\]
	which proves (a) after subtracting $1$ and dividing by $p$.

	Suppose that $g$ is primitive modulo $p$. Its order modulo $p^2$ is a multiple of $p-1$ and divides $p(p-1)$, so it is either $p-1$ or $p(p-1)$. The first case is equivalent to $g^{p-1}\equiv1\pmod{p^2}$, which is equivalent to $q_p(g)\equiv0\pmod{p}$. Thus $g$ is primitive modulo $p^2$ if and only if $q_p(g)\not\equiv0\pmod{p}$. In the latter case, $v_p(g^{p-1}-1)=1$, and Theorem~\ref{thm:6.2.5} proves that $g$ is primitive modulo $p^e$ for every $e\ge2$. Conversely, if $g$ is primitive modulo $p^e$ for some $e\ge2$, then, since $g$ generates $\ZmodUnits{p^e}$ and the reduction homomorphism $\ZmodUnits{p^e}\to\ZmodUnits{p^2}$ is surjective (an integer coprime to $p$ represents a unit modulo both moduli), the image of $g$ generates $\ZmodUnits{p^2}$ and so has order $\varphi(p^2)$; hence $g$ is primitive modulo $p^2$ and $q_p(g)\not\equiv0\pmod{p}$. This proves (b) and makes explicit the role of the Fermat quotient in the lifting criterion in Theorem~\ref{thm:6.2.5}.

	For $0\le k<p$, the binomial theorem gives
	\[
		(g+kp)^{p-1}\equiv g^{p-1}+(p-1)g^{p-2}kp\pmod{p^2}.
	\]
	Subtracting $1$ from both sides and dividing by $p$ gives
	\[
		q_p(g+kp)\equiv q_p(g)-kg^{-1}\pmod{p}.
	\]
	Exactly one residue $k\pmod{p}$ makes the right-hand side zero; part (b) with $e=2$ proves (c).

	For (d), let $r\in\mathbb{Z}$ and choose $r_0\in\{0,\ldots,p(p-1)-1\}$
	with $r_0\equiv r\pmod{p(p-1)}$. Negative powers are interpreted in
	$\ZmodUnits{p^2}$, so $g^r\equiv g^{r_0}\pmod{p^2}$.
	Part (a), applied to nonnegative powers (with $q_p(1)=0$), gives
	$q_p(g^{r_0})\equiv r_0q_p(g)\pmod p$.
	Since $a\equiv g^{r_0}\pmod{p^2}$, raising to the power $p-1$,
	subtracting $1$, and dividing by $p$ gives
	$q_p(a)\equiv q_p(g^{r_0})\pmod p$. Also $r_0\equiv r\pmod p$.
	Part (b) shows that $q_p(g)$ is invertible modulo $p$, which proves both congruences in (d). Note that $a\equiv g^r\pmod{p}$ alone does not suffice here: $q_p(a)\bmod p$ depends on the residue of $a$ modulo $p^2$, so the index must also be taken modulo $p^2$.

	For (e), we first check that $Q_p$ is well defined on residue classes. If
	$a'=a+p^2t$, then the binomial theorem gives
	\[
		(a')^{p-1}-a^{p-1}\equiv(p-1)a^{p-2}p^2t\pmod{p^3}.
	\]
	After division by $p$, this says $q_p(a')\equiv q_p(a)\pmod{p}$.
	Part (a) now says exactly that $Q_p$ is a homomorphism. The binomial theorem gives
	\[
		(1+pk)^{p-1}\equiv1+(p-1)pk\pmod{p^2},
		\qquad q_p(1+pk)\equiv-k\pmod{p},
	\]
	so every element of $\mathbb{F}_p$ occurs and $Q_p$ is surjective. By Theorem~\ref{thm:6.6.5}(a), its kernel has $\varphi(p^2)/p=p-1$ elements, and by definition it is precisely the solution set of $x^{p-1}\equiv1\pmod{p^2}$, proving (e). By Definition~\ref{def:6.9.2}, $p$ is a Wieferich prime to base $2$ if and only if $q_p(2)\equiv0\pmod{p}$; hence (f) is (b) with $g=2$. The first two such primes are $1093$ and $3511$. No primitive-root hypothesis is part of Definition~\ref{def:6.9.2}; the primitive-root interpretation in (f) applies when that additional hypothesis holds.

\end{proof}
\begin{theorem}[Lerch's formula, primitive-root form]\label{thm:6.9.5} Let $p$ be an odd prime, let $g$ be a primitive root modulo $p^2$, and, for each $a\in\{1,\ldots,p-1\}$, choose $r_a$ with $a\equiv g^{r_a}\pmod{p^2}$. Then
	\[
		W_p\equiv q_p(g)\sum_{a=1}^{p-1}r_a\pmod{p},
	\]
	where $W_p=((p-1)!+1)/p$. Equivalently, Lerch's formula (Theorem~\ref{thm:6.9.3}) is the sum of the logarithmic identities $q_p(a)\equiv r_aq_p(g)\pmod{p}$.

\end{theorem}

\begin{proof} Theorem~\ref{thm:6.9.4}(d) gives $q_p(a)\equiv r_aq_p(g)\pmod{p}$ for every $a$. Summing over $1\le a\le p-1$ and applying Lerch's formula (Theorem~\ref{thm:6.9.3}) yields
	\[
		W_p\equiv\sum_{a=1}^{p-1}q_p(a)
		\equiv q_p(g)\sum_{a=1}^{p-1}r_a\pmod{p}.
	\]
	Changing $r_a$ by a multiple of $p(p-1)$ does not change its residue modulo $p$, so the right-hand side is independent of the chosen integer representatives.

\end{proof}
\subsection{A Polynomial Application of Primitive Roots: The Chevalley--Warning Theorem}

\begin{definition}[finite field and characteristic]\label{def:6.10.1}
	A field is a commutative ring with identity $1\ne0$ in which every
	nonzero element is invertible. A finite field has finitely many elements. The
	characteristic of a field is the least positive integer $r$ with
	$r\cdot1=0$, or $0$ if no such integer exists. For a finite field,
	such an $r$ exists because its additive group is finite.
	For a prime $p$, the ring $\mathbb{F}_p=\mathbb{Z}/p\mathbb{Z}$ is the field of
	residue classes modulo $p$. The nonzero elements of a field form its
	multiplicative group.

\end{definition}
\begin{definition}[polynomial evaluation]\label{def:6.10.2}
	For $f\in\mathbb{F}_p[x_1,\ldots,x_n]$ and
	$u=(u_1,\ldots,u_n)\in\mathbb{F}_p^n$, we write $f(u)$ for the value obtained by
	substituting $u_i$ for $x_i$. We use $0^0=1$ in monomial evaluation, so
	the constant monomial $x_1^0\cdots x_n^0$ has value $1$ at every point.

\end{definition}
\begin{definition}[common zero]\label{def:6.10.3}
	Let $f_1,\ldots,f_k\in\mathbb{F}_p[x_1,\ldots,x_n]$. A point
	$u\in\mathbb{F}_p^n$ is a common zero of these polynomials if
	$f_i(u)=0$ for every $i$. It is nonzero if $u\neq(0,\ldots,0)$.

\end{definition}
\begin{definition}[total degree]\label{def:6.10.4}
	The total degree of $x_1^{e_1}\cdots x_n^{e_n}$ is
	$e_1+\cdots+e_n$. For a nonzero polynomial $f$, the total degree $\deg f$ is the largest
	total degree of a monomial with nonzero coefficient; for the zero polynomial we write $\deg0=-\infty$.

\end{definition}
\begin{theorem}\label{thm:6.10.1} The multiplicative group of any finite
	field is cyclic.

\end{theorem}

\begin{proof} Let $F$ be a finite field and $G=F^{\times}$, and let
	$e$ be the least common multiple of the orders of the elements of $G$. For
	each prime $\ell\mid e$, choose $g_\ell\in G$ whose order $r_\ell$ is
	divisible by $\ell^{v_\ell(e)}$ (such an element exists, since $e$ is the least
	common multiple of the element orders) and set
	$h_\ell=g_\ell^{r_\ell/\ell^{v_\ell(e)}}$; then $h_\ell$ has exact order
	$\ell^{v_\ell(e)}$ by Theorem~\ref{thm:6.1.1}(d). In an abelian group, the
	product of two elements of coprime orders $r,s$ has order $rs$
	(Theorem~\ref{thm:6.1.1}(f)); applying this repeatedly to the pairwise coprime orders
	$\ell^{v_\ell(e)}$ shows that the product $h=\prod_{\ell\mid e}h_\ell$ has order $e$.
	Every $x\in G$ satisfies $x^{e}=1$, an equation of degree $e$ with at most
	$e$ roots in $F$ (Theorem~\ref{thm:3.6.2}); hence $|G|\le e$, while
	$e\mid|G|$ by Lagrange's theorem (Theorem~\ref{thm:6.6.2}) applied to every element order. Thus $e=|G|$ and $h$ generates $G$. When $e=1$, the product defining $h$ is the identity and the same argument gives $|G|=1$.

\end{proof}
\begin{theorem}[Chevalley--Warning theorem]\label{thm:6.10.2} Let $f_1,
		\ldots, f_k \in \mathbb{F}_p[x_1, \ldots, x_n]$ satisfy
	$\sum_{\substack{1\le i\le k\\ f_i\neq0}} \deg f_i < n$ (with the
	empty sum equal to $0$), and let $Z = \{x \in
		\mathbb{F}_p^{\,n} \mid f_1(x) = \cdots = f_k(x) = 0\}$. Then $|Z|
		\equiv 0 \pmod{p}$. In particular, if no $f_i$ has a constant term, then $(0,\ldots,0)\in Z$ and a nonzero common zero exists.

\end{theorem}

\begin{proof} First, for an integer $e \ge 0$, the sum $\sum_{x
			\in \mathbb{F}_p} x^{e}$ equals $-1$ if $(p-1) \mid e$ and $e > 0$, and
	is congruent to $0$ otherwise: if $e > 0$ and $(p-1) \mid e$, every nonzero
	$x$ has $x^{e} = 1$, so the sum is $p - 1 \equiv -1$; if $e > 0$ and
	$(p-1) \nmid e$, a primitive root $g$ gives $\sum_{x \neq 0} x^{e} =
		\frac{g^{(p-1)e} - 1}{g^{e} - 1} = 0$; and $e = 0$ gives $p \equiv
		0$. Set $P(x) = \prod_{i=1}^{k}(1 - f_i(x)^{p-1})$: by Fermat's
	little theorem $P(x) = 1$ for $x \in Z$ and $0$ otherwise, so $|Z|
		\equiv \sum_{x} P(x) \pmod{p}$. Since $\deg P \le (p-1)\sum_{i:f_i\neq0} \deg f_i
		< (p-1)n$, every monomial $c\, x_1^{e_1}\cdots x_n^{e_n}$ of $P$ has
	some $e_j<p-1$, so
	\[
		\sum_{x\in\mathbb{F}_p^n}x_1^{e_1}\cdots x_n^{e_n}
		=\prod_{j=1}^{n}\left(\sum_{x_j\in\mathbb{F}_p}x_j^{e_j}\right)
	\]
	has a zero factor. Hence $\sum_x P(x)
		\equiv 0$, that is, $|Z| \equiv 0 \pmod{p}$. If all constant terms vanish, then $(0,\ldots,0)\in Z$, so $|Z|\ge1$. Since the positive integer $|Z|$ is divisible by $p$, in fact $|Z|\ge p\ge2$; therefore $Z$ contains a point different from $(0,\ldots,0)$.\qedhere

\end{proof}
\begin{theorem}\label{thm:6.10.3} For every odd prime $p$ there exist $x, y$ with $1 + x^{2} +
		y^{2} \equiv 0 \pmod{p}$.

\end{theorem}

\begin{proof} The form $x_0^{2} + x_1^{2} + x_2^{2}$ has
	degree $2 < 3$ and no constant term, so by Theorem~\ref{thm:6.10.2} there is a
	nonzero $(x_0, x_1, x_2)$ with $x_0^{2} + x_1^{2} + x_2^{2} \equiv 0
		\pmod{p}$. Some coordinate, say $x_0$, is nonzero; multiplying by
	$x_0^{-2}$ gives $1 + x^{2} + y^{2} \equiv 0 \pmod{p}$.

\end{proof}
\begin{altproof}[Alternative proof of Theorem~\ref{thm:6.10.3}] The $\tfrac{p+1}{2}$ values $\{x^{2} \mid 0
		\le x \le \tfrac{p-1}{2}\}$ are pairwise distinct modulo $p$: if
	$x^{2} \equiv y^{2}$ with $0 \le x, y \le \tfrac{p-1}{2}$, then $p
		\mid (x-y)(x+y)$, and since $|x-y| < p$ and $0 \le x + y \le p-1 < p$,
	either $x = y$ or $x + y = 0$, and the latter also gives $x = y = 0$.
	By the same argument, the $\tfrac{p+1}{2}$ values $\{-1 - y^{2} \mid 0
		\le y \le \tfrac{p-1}{2}\}$ are pairwise distinct. Two subsets of
	size $\tfrac{p+1}{2}$ in $\mathbb{Z}_p$ (of size $p$) intersect,
	since $2 \cdot \tfrac{p+1}{2} = p + 1 > p$; an element in the
	intersection gives $x^{2} \equiv -1 - y^{2}$, that is, $1 + x^{2} + y^{2}
		\equiv 0 \pmod{p}$.

\end{altproof}
\subsection{Orders and Primitive Prime Factors: Zsigmondy's Theorem and Cyclotomic Polynomials}

\begin{definition}[cyclotomic polynomial]\label{def:6.11.1}
	For an integer $n\ge1$, put $\Phi_n(x) = \prod_{\substack{1 \le k \le n \\ \gcd(k, n) = 1}}\big(x
		- \zeta_n^{\,k}\big)$, where $\zeta_n = e^{2\pi i / n}$.

\end{definition}
\begin{theorem}\label{thm:6.11.1} For every $n\in\mathbb{Z}^+$, we have $x^{n} - 1 = \prod_{d \mid n} \Phi_d(x)$ in $\mathbb{C}[x]$.

\end{theorem}

\begin{proof} The $n$ numbers $\zeta_n^k$ with $1\le k\le n$ are roots of
	$x^n-1$, and they are distinct, because $\zeta_n^k=\zeta_n^l$ means
	$e^{2\pi i(k-l)/n}=1$, that is, $n\mid k-l$. By Theorem~\ref{thm:3.6.2},
	applied over $\mathbb{C}$, there are no other roots. If $r_1,\ldots,r_j$
	are distinct roots of a polynomial $f$ and $f=(x-r_1)\cdots(x-r_{j-1})q$,
	then $q(r_j)=0$ because each $r_j-r_i$ is nonzero; hence the factor step
	in the proof of Theorem~\ref{thm:3.6.2}, applied $n$ times, together with a
	comparison of leading coefficients, gives
	$x^n-1=\prod_{k=1}^{n}(x-\zeta_n^k)$. Every $n$-th root of unity
	$\zeta$ has a well-defined multiplicative order $d \mid n$
	(Theorem~\ref{thm:6.1.1}(c), applied in the group $\mathbb{C}^{\times}$),
	and the roots of $\Phi_d$ are exactly the $n$-th roots of unity of
	order precisely $d$: each $\zeta_d^k$ with $\gcd(k,d)=1$ has order
	$d/\gcd(d,k)=d$ by Theorem~\ref{thm:6.1.1}(d), and conversely a root of
	order $d$ is a root of $x^d-1$, hence equals $\zeta_d^k$ for some
	$1\le k\le d$ by the first part of the proof, and its order $d$
	forces $\gcd(k,d)=1$. Grouping the $n$ linear factors of $x^{n} - 1$
	by the order of their roots gives $x^{n} - 1 = \prod_{d \mid n} \Phi_d(x)$.

\end{proof}
\begin{theorem}\label{thm:6.11.2} For every $n\in\mathbb{Z}^+$, the polynomial $\Phi_n(x)$ lies in $\mathbb{Z}[x]$, is
	monic, and has degree $\varphi(n)$.

\end{theorem}

\begin{proof} We argue by induction on $n$. For $n=1$, we have $\Phi_1 = x - 1$.
	Under the induction hypothesis for indices less than $n$, the product $D = \prod_{d
			\mid n,\, d < n}\Phi_d$ is monic with integer coefficients, and
	$x^{n} - 1 = D\cdot\Phi_n$ (Theorem~\ref{thm:6.11.1}). Dividing the monic
	integer polynomial $x^{n} - 1$ by the monic integer polynomial $D$
	using the usual algorithm keeps every coefficient integral at each step
	(since the leading coefficient of $D$ is $1$, no denominators are ever
	introduced), so the quotient $\Phi_n = \frac{x^{n} - 1}{D}$ lies in
	$\mathbb{Z}[x]$ and is monic; and
	$\sum_{d \mid n}\varphi(d) = n$ (Theorem~\ref{thm:5.1.5}) gives $\deg \Phi_n =
		\varphi(n)$.

\end{proof}
\begin{theorem}\label{thm:6.11.3} For every $n\in\mathbb{Z}^+$, we have $\Phi_n(x) = \prod_{d \mid n}(x^{d} -
		1)^{\mu(n/d)}$. Moreover, write
	$\Phi_n(x)=\sum_{j=0}^{\varphi(n)}c_jx^j$ and, for integers $a,b$,
	define the homogenized polynomial by
	\[
		\Phi_n(a,b):=\sum_{j=0}^{\varphi(n)}c_j a^j
		b^{\varphi(n)-j}\in\mathbb{Z}.
	\]
	Then $a^{n} - b^{n} = \prod_{d \mid n} \Phi_d(a, b)$. When $b\neq0$,
	the definition agrees with $b^{\varphi(n)}\Phi_n(a/b)$.

\end{theorem}

\begin{proof} We first record the product form of M{\"o}bius inversion in the multiplicative group $\mathbb{Q}(x)^{\times}$ of nonzero rational functions, that is, of quotients $f/g$ of nonzero polynomials $f,g\in\mathbb{Q}[x]$. If $U_n=\prod_{d\mid n}V_d$, then
	\[
		\prod_{d\mid n}U_d^{\mu(n/d)}=\prod_{e\mid n}V_e^{\sum_{e\mid d\mid n}\mu(n/d)}=V_n,
	\]
	because the inner sum is $1$ for $e=n$ and $0$ otherwise (Theorem~\ref{thm:5.3.2}(a)). Applying this to $U_d=x^d-1$ and $V_d=\Phi_d(x)$ proves the first formula in the theorem and explains the meaning of negative exponents.

	For $b\neq0$, substitution in $x^n-1=\prod_{d\mid n}\Phi_d(x)$ and
	multiplication by $b^n=\prod_{d\mid n}b^{\varphi(d)}$ give
	$a^n-b^n=\prod_{d\mid n}\Phi_d(a,b)$. For $b=0$, only the term
	$j=\varphi(d)$ of $\Phi_d(a,0)$ survives, and its coefficient is $1$ because
	$\Phi_d$ is monic of degree $\varphi(d)$ (Theorem~\ref{thm:6.11.2}); hence
	$\prod_{d\mid n}\Phi_d(a,0)=a^{\sum_{d\mid n}\varphi(d)}=a^n$ by
	Theorem~\ref{thm:5.1.5}, which is the left-hand side.

\end{proof}
\begin{theorem}\label{thm:6.11.4} If the roots of $\Phi_q$ are $\zeta_1,\ldots,\zeta_{\varphi(q)}$, then for every integer $n\ge1$,
	\[
		c_q(n)=\sum_{\substack{1\le k\le q\\\gcd(k,q)=1}}\zeta_q^{kn}=\sum_{j=1}^{\varphi(q)}\zeta_j^n.
	\]
	Consequently, Newton's identities express every coefficient of $\Phi_q$ as a polynomial with rational coefficients in $c_q(1),\ldots,c_q(\varphi(q))$, and
	\[
		[x^{\varphi(q)-1}]\Phi_q(x)=-c_q(1)=-\mu(q).
	\]

\end{theorem}

\begin{proof} By Definition~\ref{def:6.11.1}, the roots of $\Phi_q$ are precisely $\zeta_q^k$ with $1\le k\le q$ and $\gcd(k,q)=1$, so the first identity is a direct reindexing. Newton's identities relate the elementary symmetric functions of the roots, which are the signed coefficients of the monic polynomial $\Phi_q$, to the power sums of the roots. The coefficient of $x^{\varphi(q)-1}$ is minus the first power sum, hence is $-c_q(1)$. Theorem~\ref{thm:5.6.1} gives $c_q(1)=\mu(q)$.

\end{proof}
\begin{theorem}\label{thm:6.11.5} For all $q,n\in\mathbb{Z}^+$,
	\[
		\sum_{k=1}^{q}\zeta_q^{kn}=\sum_{d\mid q}c_d(n).
	\]

\end{theorem}

\begin{proof} Every $q$-th root of unity has a unique order $d\mid q$, so grouping the $q$-th roots of unity by their exact orders and summing their $n$-th powers gives the stated identity.

	The same grouping by exact order underlies three identities already proved:
	multiplying the linear factors instead of summing the powers gives
	$x^{q}-1=\prod_{d\mid q}\Phi_d(x)$ (Theorem~\ref{thm:6.11.1}); multiplicative
	M{\"o}bius inversion of that product gives
	$\Phi_q(x)=\prod_{d\mid q}(x^{d}-1)^{\mu(q/d)}$
	(Theorem~\ref{thm:6.11.3}); and ordinary M{\"o}bius inversion of the displayed
	sum gives $c_q(n)=\sum_{d\mid\gcd(q,n)}d\,\mu(q/d)$
	(Theorem~\ref{thm:5.6.1}).

\end{proof}

\begin{theorem}\label{thm:6.11.6} Let $p$ be prime and $d\mid p-1$. The roots of $\Phi_d(x)$ in $\mathbb{F}_p$ are exactly the elements of order $d$ in $\mathbb{F}_p^{\times}$; there are $\varphi(d)$ of them. In particular, the roots of $\Phi_{p-1}(x)$ are exactly the primitive roots modulo $p$.

\end{theorem}

\begin{proof} For every $m\in\mathbb{Z}^+$, the identity $x^{m}-1=\prod_{e\mid m}\Phi_e(x)$ of Theorem~\ref{thm:6.11.1} has all its factors in $\mathbb{Z}[x]$ (Theorem~\ref{thm:6.11.2}), so it also holds in $\mathbb{F}_p[x]$ after reduction of coefficients. For $e\mid p-1$, let $N(e)$ be the number of roots of $\Phi_e$ in $\mathbb{F}_p$. Since $\Phi_e$ is monic of degree $\varphi(e)$, Theorem~\ref{thm:3.6.2} gives $N(e)\le\varphi(e)$. By Fermat's little theorem, each of the $p-1$ nonzero elements of $\mathbb{F}_p$ is a root of $x^{p-1}-1=\prod_{e\mid p-1}\Phi_e(x)$, hence, as $\mathbb{F}_p$ is a field, a root of at least one $\Phi_e$. Therefore
	\[
		p-1\le\sum_{e\mid p-1}N(e)\le\sum_{e\mid p-1}\varphi(e)=p-1
	\]
	by Theorem~\ref{thm:5.1.5}. Equality throughout shows that $N(e)=\varphi(e)$ for every $e\mid p-1$ and that every nonzero element of $\mathbb{F}_p$ is a root of exactly one $\Phi_e$.

	Now let $\alpha\in\mathbb{F}_p^{\times}$ have order $k$, so $k\mid p-1$. Since $\alpha^k=1$ and $x^k-1=\prod_{c\mid k}\Phi_c(x)$, the element $\alpha$ is a root of $\Phi_c$ for some $c\mid k$. As $\Phi_c$ divides $x^c-1$, also $\alpha^c=1$, so $k\mid c$ by Theorem~\ref{thm:6.1.1}(c); hence $c=k$. Thus every element of order $k$ is a root of $\Phi_k$, and of no other $\Phi_e$. Consequently, the roots of $\Phi_d$ are exactly the elements of order $d$, and there are $N(d)=\varphi(d)$ of them.

\end{proof}
\begin{definition}[primitive prime factor]\label{def:6.11.2} Let $n\ge1$ be an integer and let $a,b$ be coprime positive integers. A prime $p$ is a primitive prime factor of $a^n-b^n$ if
	\[
		p\mid a^n-b^n\qquad\text{and}\qquad p\nmid a^k-b^k\quad(1\le k<n).
	\]

\end{definition}
\begin{theorem}\label{thm:6.11.7} Under the conditions of Definition~\ref{def:6.11.2}, a prime $p$ is a primitive prime factor of $a^n-b^n$ if and only if $p\nmid ab$ and $ab^{-1}$ has order exactly $n$ modulo $p$. If $n>1$, then $p$ is odd and $n\mid p-1$.

\end{theorem}

\begin{proof} If $p\mid a^n-b^n$ and $p\mid b$, then $p\mid a^n$ and hence $p\mid a$, contradicting $\gcd(a,b)=1$; similarly $p\nmid a$. Thus $b^{-1}$ exists modulo $p$, and
	\[
		p\mid a^k-b^k\iff (ab^{-1})^k\equiv1\pmod{p}.
	\]
	Therefore the least positive exponent $k$ for which $p\mid a^k-b^k$ is exactly $\delta_p(ab^{-1})$, proving the equivalence. If $n>1$, the prime $p$ cannot be $2$, since $\ZmodUnits{2}$ has order $1$. For odd $p$, Lagrange's theorem
	(Theorem~\ref{thm:6.6.2}) gives
	$n=\delta_p(ab^{-1})\mid|\mathbb{F}_p^{\times}|=p-1$.

\end{proof}
\begin{theorem}\label{thm:6.11.8} Let $a,b$ be coprime positive integers and let $n\in\mathbb{Z}^+$.

	\begin{enumerate}[label={\alph*)}]
		\item Let $p$ be an odd prime with $p \nmid ab$, and let $d$ be the least
		      positive integer with $p \mid a^{d} - b^{d}$ (the order of
		      $ab^{-1}$ modulo $p$, so $d\mid p - 1$ and $p\nmid d$). Then
		      $p \mid \Phi_n(a, b)$ if and only if $n = d\,p^{k}$ for some
		      $k \ge 0$, and
		      \[
			      v_p\big(\Phi_n(a, b)\big) =
			      \begin{cases} v_p(a^{d} - b^{d}) & n = d, \\ 1 & n = d\,p^{k},\ k \ge
			                    1,                          \\ 0 & \text{otherwise.}
			      \end{cases}
		      \]
		\item Let $p = 2$ and let $a$ and $b$ both be odd. Then
		      \[
			      v_2\big(\Phi_n(a, b)\big) =
			      \begin{cases} v_2(a - b) & n = 1, \\ v_2(a+b) & n = 2, \\ 1 & n =
				                                                              2^{k},\ k \ge 2, \\ 0 & \text{otherwise.}
			      \end{cases}
		      \]
	\end{enumerate}

\end{theorem}

\begin{proof} First suppose $a=b$. Since $\gcd(a,b)=1$, for positive $a,b$ this forces $a=b=1$. Then $\Phi_1(1,1)=0$, $\Phi_{p^k}(1,1)=p$ for every prime power $p^k$, and $\Phi_n(1,1)=1$ when $n>1$ is not a prime power: dividing $x^n-1=\prod_{e\mid n}\Phi_e(x)$ by $x-1=\Phi_1(x)$ and letting $x\to1$ gives $n=\prod_{e\mid n,\,e>1}\Phi_e(1,1)$, and, arguing by induction on $n$, we find that the proper divisors $e>1$ of $n$ contribute a factor $q$ for each prime-power divisor $q^{j}$ (and $1$ for each non-prime-power divisor); these contributions multiply to $n/p$ when $n=p^{k}$ (the divisors $p,p^{2},\ldots,p^{k-1}$) and to $n$ when $n$ has at least two distinct prime factors, leaving $\Phi_n(1,1)=p$ in the first case and $\Phi_n(1,1)=1$ in the second. These values give the stated formulas directly (with $v_p(0)=+\infty$). Hence assume $a\neq b$ whenever LTE is used. By Theorem~\ref{thm:6.11.7} (or Theorem~\ref{thm:6.1.1}(c) applied to $ab^{-1}$), we have $p \mid a^{n} - b^{n}$ if and only if $d \mid
		n$. If $d \nmid n$, then $p \nmid a^{n} - b^{n} = \prod_{e \mid
			n}\Phi_e(a, b)$, so $v_p(\Phi_n) = 0$. If $d \mid n$, write $n = dm$,
	$s = v_p(m)$, and $A = a^{d}, B = b^{d}$ (so $p \mid A - B$, $p \nmid
		AB$); the odd-prime lifting-the-exponent lemma (Theorem~\ref{thm:4.8.3}) gives $v_p(a^{n} - b^{n}) =
		v_p(A^{m} - B^{m}) = t + s$ where $t = v_p(a^{d} - b^{d})$. We now prove the index restriction that will be used below. If $p\mid\Phi_e(a,b)$, then $\Phi_e(a,b)\mid a^e-b^e$ implies $p\mid a^e-b^e$, and therefore $d\mid e$; write $e=df$. If a prime $q\neq p$ divided $f$, then the cyclotomic factorizations imply
	\[
		\Phi_e(a,b)\mid\frac{a^e-b^e}{a^{e/q}-b^{e/q}}.
	\]
	Indeed, $\Phi_e$ occurs in the factorization indexed by the divisors of $e$ but not in that indexed by the divisors of $e/q$. Applying LTE to numerator and denominator, with $A=a^d$ and $B=b^d$, gives
	\[
		v_p\left(\frac{a^e-b^e}{a^{e/q}-b^{e/q}}\right)=\big(t+v_p(f)\big)-\big(t+v_p(f/q)\big)=v_p(q)=0,
	\]
	contradicting $p\mid\Phi_e(a,b)$. Hence $f$ has no prime divisor other than $p$, so $e=dp^j$ for some $j\ge0$. Those indices dividing $n=dm$ are exactly $dp^j$ with $0\le j\le s$. Applying this to the index $d$ itself (where $s = 0$)
	gives $v_p(\Phi_d) = t$; proceeding inductively through the indices $d p^{j}$ ($j=1,2,\ldots$), we successively subtract the previously determined valuations from the total $v_p(a^{d
				p^{j}} - b^{d p^{j}}) = t + j$ to obtain $v_p(\Phi_{d
			p^{j}}) = 1$ for $j \ge 1$. For our $n = dm$, the sum $\sum_{e
			\mid n} v_p(\Phi_e) = t + s$ is exhausted by the indices $\{d
		p^{j}\}_{j=0}^{s}$, so every other $e$ (including $n$ itself when $n
		\neq d p^{k}$) has $v_p(\Phi_e) = 0$.

	Now let $p=2$, and write $n=2^km$ with $m$ odd.
	For $n=1$, we have $\Phi_1(a,b)=a-b$ directly. For $n=2$, the identity
	$a^{2}-b^{2}=(a-b)(a+b)=\Phi_1(a,b)\Phi_2(a,b)$ gives $\Phi_2(a,b)=a+b$.
	If $n>1$ is odd, then by Theorem~\ref{thm:4.8.1} (lifting the exponent, coprime case, valid at $p=2$
	since $\gcd(n, 2)=1$), we have $v_2(a^{n}-b^{n})=v_2(a-b)$; since
	$a^{n}-b^{n}=\prod_{e\mid n}\Phi_e(a,b)$ and $\Phi_1(a,b)=a-b$
	already accounts for this entire valuation, every other factor
	$\Phi_e(a,b)$ ($e\mid n$, $e>1$) is odd, so $v_2(\Phi_n(a,b))=0$. For
	$n=2^k$ with $k\ge2$, we argue by induction on $k$. For $k=2$,
	we have $\Phi_4(a,b)=a^{2}+b^{2}$ (from
	$a^{4}-b^{4}=(a-b)(a+b)(a^{2}+b^{2})$), and since $a,b$ are odd,
	$a^{2}\equiv b^{2}\equiv1\pmod{8}$, so $a^{2}+b^{2}\equiv2\pmod{8}$,
	giving $v_2(\Phi_4(a,b))=1$. Now assume
	$v_2(\Phi_{2^{i}}(a,b))=1$ for $2\le i<k$. From
	$a^{2^{k}}-b^{2^{k}}=\prod_{j=0}^{k}\Phi_{2^{j}}(a,b)$ and Theorem~\ref{thm:4.8.6}, $v_2(a^{2^{k}}-b^{2^{k}})=v_2(a-b)+v_2(a+b)+k-1$; subtracting
	$v_2(\Phi_1)+v_2(\Phi_2)=v_2(a-b)+v_2(a+b)$ and the inductive
	contributions $\sum_{i=2}^{k-1}v_2(\Phi_{2^{i}}(a,b))=k-2$ leaves
	$v_2(\Phi_{2^{k}}(a,b))=(k-1)-(k-2)=1$. For $n=2^{k}m$ with $m>1$
	odd, $k\ge0$: if $k=0$ this is the odd case above; if $k\ge1$,
	comparing $v_2(a^{n}-b^{n})=v_2(a-b)+v_2(a+b)+k-1$ (Theorem~\ref{thm:4.8.6})
	with the sum
	$v_2(\Phi_1(a,b))+v_2(\Phi_2(a,b))+\sum_{j=2}^{k}v_2(\Phi_{2^{j}}(a,b))=v_2(a-b)+v_2(a+b)+(k-1)$
	over only the power-of-$2$ divisors of $n$ shows that these two totals
	already agree; since $a^{n}-b^{n}=\prod_{e\mid n}\Phi_e(a,b)$ is the
	product over all divisors of $n$, every remaining divisor $e\mid n$
	with odd part exceeding $1$ contributes $v_2(\Phi_e(a,b))=0$; in
	particular, $v_2(\Phi_n(a,b))=0$.

\end{proof}
\begin{theorem}[Zsigmondy's theorem]\label{thm:6.11.9} Let $a > b \ge 1$ be
	coprime and $n \ge 1$. Then $a^{n} - b^{n}$ has a primitive prime
	factor except in exactly the following cases: (1) $n = 1$ and $a - b = 1$; (2) $n = 2$
	and $a + b$ is a power of $2$; (3) $n = 6$ and $(a, b) = (2, 1)$.

\end{theorem}

\begin{proof} If $n = 1$, every prime factor of $a - b$
	is primitive, so there is no primitive prime factor exactly when
	$a - b = 1$. If $n = 2$, every odd prime factor of $a + b$ is
	primitive, because
	\[
		\gcd(a - b, a + b) \mid 2\gcd(a, b) = 2.
	\]
	Thus there is no primitive prime factor exactly when $a + b$ is a
	power of $2$ (in that case $a$ and $b$ are odd, so $2 \mid a - b$).
	Hence assume $n \ge 3$ and suppose that $a^{n} - b^{n}$ has no
	primitive prime factor.

	We first reduce this assumption to a condition involving a single cyclotomic value. For
	$x > y \ge 1$ and every root of unity $\zeta$,
	\[
		|x - \zeta y| \ge x - y,
	\]
	with strict inequality when $\zeta \neq 1$. Indeed, $|\zeta|=1$ gives
	$|1-\zeta|^2=2-2\operatorname{Re}\zeta$, so
	\[
		|x-\zeta y|^2=x^2-2xy\operatorname{Re}\zeta+y^2=(x-y)^2+xy\,|1-\zeta|^2,
	\]
	and $xy>0$. For $n\ge3$, primitive $n$-th roots occur in nonreal conjugate pairs, so $\Phi_n(a,b)$ is positive and equals the product of the magnitudes $|a-\zeta b|$. Since no such root is $1$, the product formula gives
	$\Phi_n(a, b) > 1$. Let $q$ be a prime factor of $\Phi_n(a, b)$. By
	assumption, $q$ is not a primitive prime factor. Note first that $q \nmid ab$: indeed
	$q$ divides $\Phi_n(a, b)$, which divides $a^{n} - b^{n}$, and a prime
	dividing both $ab$ and $a^{n} - b^{n}$ would divide both $a$ and $b$,
	contradicting $\gcd(a, b) = 1$. For the same reason, $a$ and $b$ are both
	odd when $q = 2$. Non-primitivity means that $q \mid a^{k}-b^{k}$ for
	some $1 \le k < n$, so the least positive integer $d$ with
	$q \mid a^{d}-b^{d}$ satisfies $d \le k < n$; in particular, $d \neq n$.
	If $q$ is odd, Theorem~\ref{thm:6.11.8}(a) gives $n = d\,q^{j}$ with
	$j \ge 0$, and $j = 0$ would mean $n = d$; hence
	\[
		n = d\,q^{j},\qquad j \ge 1,\qquad d \mid q - 1,
	\]
	and $v_q(\Phi_n(a, b)) = 1$; in particular, $q$ is the largest prime
	factor of $n$. If $q = 2$, Theorem~\ref{thm:6.11.8}(b) (applicable because
	$a$ and $b$ are both odd) gives $n = 2^{j}$ with $j \ge 1$ (the
	cases $n = 1, 2$ of that formula do not arise, as $n \ge 3$) and again
	$v_2(\Phi_n(a, b)) = 1$, so the same conclusion holds. Consequently
	all prime factors of $\Phi_n(a, b)$ equal the largest prime factor
	$p$ of $n$, and
	\[
		\Phi_n(a, b) = p.
	\]
	Writing
	\[
		n = k\,p^{\alpha},
	\]
	we have $\alpha \ge 1$ and $k \mid p - 1$ (with $k = 1$ when
	$p = 2$).

	Suppose first that $a - b \ge 2$. The strict root estimate above gives
	\[
		p = \Phi_n(a, b) > (a - b)^{\varphi(n)} \ge 2^{\varphi(n)}
		\ge 1 + \varphi(n) \ge 1 + \varphi(p) = p,
	\]
	a contradiction. Here
	$\varphi(n) = p^{\alpha-1}(p - 1)\varphi(k) \ge p - 1$.

	It remains to consider $a - b = 1$, so $a = b + 1$. First let
	$\alpha \ge 2$. To obtain the cyclotomic identities below, let $z$ be
	a complex root of unity of order $d$. Then $z^p$ has order
	$d/\gcd(d,p)$, so $z^p$ is a primitive $m$-th root exactly when $d=mp$ in the case $p\mid m$, and exactly when $d=m$ or $d=mp$ in the case $p\nmid m$.
	The polynomial $\Phi_m(x^p)$ has distinct roots, since the roots of
	$\Phi_m$ are distinct and nonzero. Comparing these roots and the
	leading coefficients gives the corresponding identities in one variable.
	Homogenizing as in Theorem~\ref{thm:6.11.3} then gives
	\[
		\Phi_{mp}(x, y) = \Phi_m(x^p, y^p)\quad (p \mid m),\qquad
		\Phi_{mp}(x, y) = \frac{\Phi_m(x^p, y^p)}{\Phi_m(x, y)}\quad (p \nmid m).
	\]
	Since $\alpha \ge 2$, the first identity applies with
	$m = kp^{\alpha-1}$, and together with the root estimate it gives
	\[
		\begin{aligned}
			p = \Phi_{kp^\alpha}(a, b)
			 & = \Phi_{kp^{\alpha-1}}(a^p, b^p)                \\
			 & \ge a^p - b^p = (b + 1)^p - b^p \ge pb + 1 > p,
		\end{aligned}
	\]
	a contradiction.

	Thus $\alpha = 1$. Since $p \nmid k$, the second identity above gives
	$\Phi_{kp}(x, y) = \Phi_k(x^p, y^p)/\Phi_k(x, y)$. For every
	primitive $k$-th root $\zeta$,
	\[
		\frac{|a^p - \zeta b^p|}{|a - \zeta b|}
		\ge \frac{a^p - b^p}{a + b},
	\]
	and $(a^p - b^p)/(a + b) \ge 1$. Taking the product over the
	$\varphi(k)$ primitive $k$-th roots gives
	\[
		p = \Phi_{kp}(a, b)
		\ge \left(\frac{a^p - b^p}{a + b}\right)^{\varphi(k)}
		\ge \frac{(b + 1)^p - b^p}{2b + 1}.
	\]
	If $b = 1$, then
	\[
		p \ge \frac{2^p - 1}{3},
	\]
	which forces $p \le 3$. If $b \ge 2$, the preceding inequality and
	$(b + 1)^p - b^p > p b^{p-1}$ give
	\[
		2b + 1 > b^{p-1};
	\]
	this also forces $p \le 3$, since $b^4 > 2b + 1$ for $b \ge 2$.
	Therefore $p \in \{2, 3\}$.

	If $p = 2$, then $k \mid p - 1$ gives $k = 1$, hence $n = 2$,
	contrary to $n \ge 3$. If $p = 3$, then $k \mid 2$, so $k = 1$ or
	$k = 2$. For $k = 1$, the equality $3 = \Phi_3(a, b)$ is impossible
	because
	\[
		\Phi_3(b + 1, b) = 3b^2 + 3b + 1 \equiv 1 \pmod{3}.
	\]
	For $k = 2$, we have $n = 6$, and
	\[
		3 = \Phi_6(a, b) = a^2 - ab + b^2 = b^2 + b + 1,
	\]
	so $b = 1$ and $a = 2$. Conversely,
	$2^6 - 1 = 63 = 3^2 \cdot 7$, and both $3$ and $7$ divide earlier
	terms ($3 \mid 2^2 - 1$ and $7 \mid 2^3 - 1$), so
	$(a, b, n) = (2, 1, 6)$ is indeed an exception. This completes the
	proof.

\end{proof}
\begin{theorem}\label{thm:6.11.10} For every $n \ge 1$, there are infinitely many
	primes $p \equiv 1 \pmod{n}$.

\end{theorem}

\begin{proof} For $n = 1$, this is the infinitude of
	primes, so assume $n \ge 2$. Any primitive prime factor $p$ of $a^{n} -
		1$ satisfies $n \mid p - 1$, that is, $p \equiv 1 \pmod{n}$ (Theorem~\ref{thm:6.11.7}, whose last assertion applies because $n \ge 2$). Suppose, for contradiction, that $q_1, \ldots, q_r$ are all
	the primes congruent to $1 \pmod{n}$. Choose $a\equiv1\pmod{q_1\cdots q_r}$ with $a>2$. If $n=2$, choose $a$ even; then $a+1>2$ is odd and cannot be a power of $2$, so the exceptional pair for exponent $2$ is avoided. If $n=6$, the condition $a>2$ avoids the exceptional base $a=2$; for every other $n$ the listed Zsigmondy exceptions do not apply. By Theorem~\ref{thm:6.11.9}, $a^{n} - 1$ has a primitive
	prime factor $p$, which satisfies $p \equiv 1 \pmod{n}$ and
	$\delta_p(a) = n$. But $a \equiv 1 \pmod{q_i}$ gives
	$\delta_{q_i}(a) = 1 \neq n$, so $p \notin \{q_1, \ldots,
		q_r\}$. Thus $p$ is a new prime congruent to $1 \pmod{n}$, a contradiction. Hence there
	are infinitely many such primes.

\end{proof}
\newpage

\section{Quadratic Residues}

\subsection{Quadratic Residues and Euler's Criterion}

\begin{definition}[quadratic residue]\label{def:7.1.1} Let $m \ge 2$ be an integer and let
	$a$ be an integer with $\gcd(a, m) = 1$. If $x^{2} \equiv a \pmod{m}$
	has a solution, $a$ is called a quadratic residue modulo $m$; if it
	has none, $a$ is called a quadratic non-residue modulo $m$. Integers
	$a$ with $\gcd(a, m) > 1$ are classified as neither residues nor non-residues: the terms
	quadratic residue and quadratic non-residue are used only for $a$
	coprime to the modulus. In this subsection, the modulus is an odd prime
	$p$, so the condition $\gcd(a, p) = 1$ reads $p \nmid a$, and the
	classification of the $p-1$ nonzero classes into residues and
	non-residues is exhaustive.

\end{definition}
\begin{theorem}\label{thm:7.1.1} Let $p$ be an odd prime with $p \nmid a$. The congruence
	$ax^{2} + bx + c \equiv 0 \pmod{p}$ is solvable if and only if the
	discriminant $d = b^{2} - 4ac$ is a quadratic residue modulo $p$ or
	$d \equiv 0 \pmod{p}$.

\end{theorem}

\begin{proof} Multiplying by $4a$, which is coprime to $p$,
	preserves the solution set and gives $4a^{2}x^{2} + 4abx + 4ac \equiv
		0$, that is, $(2ax + b)^{2} \equiv b^{2} - 4ac = d \pmod{p}$. Set $y =
		2ax + b$. Since $\gcd(2a, p) = 1$, the affine map $x \mapsto 2ax + b$
	is a bijection modulo $p$, so the original congruence is solvable if
	and only if $y^{2} \equiv d \pmod{p}$ is solvable, that is, if and only if
	$d$ is a quadratic residue or $d \equiv 0$.

\end{proof}
\begin{theorem}[Euler's criterion]\label{thm:7.1.2} Let $p$ be an
	odd prime, and suppose $\gcd(a, p) = 1$. Then $a^{\frac{p-1}{2}} \equiv 1
		\pmod{p}$ if $a$ is a quadratic residue, and $a^{\frac{p-1}{2}} \equiv -1 \pmod{p}$ if
	$a$ is a non-residue.

\end{theorem}

\begin{proof} Since $p$ is odd, $\tfrac{p-1}{2}$ is an
	integer, so the assertion makes sense. We argue in three steps.

	First, the power is $\pm 1$. By Fermat's little theorem (Theorem~\ref{thm:3.3.2}), $a^{p-1}
		\equiv 1 \pmod{p}$, so $y = a^{(p-1)/2}$ satisfies $y^{2} \equiv 1
		\pmod{p}$; the only solutions of $y^{2} \equiv 1 \pmod{p}$ are $y
		\equiv \pm 1$ (Theorem~\ref{thm:2.3.2} applied to $p \mid (y-1)(y+1)$). Hence $a^{(p-1)/2} \equiv 1$ or $a^{(p-1)/2} \equiv
		-1$, and not both, because $1 \equiv -1$ would force $p \mid 2$.

	Second, the quadratic residues are exactly the even powers of a
	primitive root. Let $g$ be a primitive root modulo $p$ (Theorem~\ref{thm:6.2.6}), so
	$\delta_p(g) = p-1$ and every $a$ with $p \nmid a$ satisfies $a
		\equiv g^{k} \pmod{p}$ for exactly one $k$ with $0 \le k \le p-2$. If
	$k = 2j$ is even, then $a \equiv (g^{j})^{2}$, so $a$ is a quadratic
	residue. Conversely, let $a \equiv x^{2} \pmod{p}$. Then $p \nmid x$,
	so $x \equiv g^{i}$ for some $i$, and $g^{2i} \equiv g^{k} \pmod{p}$
	forces $2i \equiv k \pmod{p-1}$ because $g$ has order $p-1$. Writing
	$k = 2i + (p-1)t$ and using the fact that $p-1$ is even, we see that $k$ is
	even.

	Finally, $a^{(p-1)/2} \equiv g^{k(p-1)/2}
		\pmod{p}$, and since $\delta_p(g) = p-1$, this is congruent to $1$ if and
	only if $(p-1) \mid k\,\tfrac{p-1}{2}$, that is, if and only if $k$ is
	even. By the second step, this happens exactly when $a$ is a quadratic residue,
	and by the first step, the power equals $-1$ in the remaining case, that is,
	when $a$ is a non-residue.

\end{proof}
\begin{altproof}[Alternative proof of Theorem~\ref{thm:7.1.2}] By Fermat's
	little theorem, all $p-1$ nonzero residue classes are roots of
	\[
		x^{p-1} - 1 = \big(x^{\frac{p-1}{2}} - 1\big)\big(x^{\frac{p-1}{2}} +
		1\big) \pmod{p}.
	\]
	No class is a root of both factors, since that would give $1 \equiv
		-1 \pmod{p}$, which is impossible for $p > 2$. By the polynomial root bound
	(Theorem~\ref{thm:3.6.2}), each factor has at most $\tfrac{p-1}{2}$ roots, and
	together they account for all $p-1$ nonzero classes, so each factor
	has exactly $\tfrac{p-1}{2}$ roots. Next, the set $Q$ of quadratic
	residues has exactly $\tfrac{p-1}{2}$ elements: the squaring map $x
		\mapsto x^{2}$ on the $p-1$ nonzero classes has image $Q$ and is
	two-to-one, because $x^{2} \equiv y^{2} \pmod{p}$ means $p \mid (x -
		y)(x + y)$, that is, $y \equiv \pm x$, and $x \not\equiv -x$ for $p$
	odd and $p \nmid x$. Finally, every $a \in Q$, say $a \equiv x^{2}$,
	satisfies $a^{(p-1)/2} \equiv x^{p-1} \equiv 1$, so $Q$ is contained in the
	root set of the first factor; both sets have $\tfrac{p-1}{2}$
	elements, so they coincide. Consequently, the non-residues are
	exactly the roots of the second factor, that is, they satisfy
	$a^{(p-1)/2} \equiv -1$.

\end{altproof}
\begin{altproof}[Alternative proof of Theorem~\ref{thm:7.1.2}] Let $S = \{1,
		2, \ldots, p-1\}$, whose elements represent the nonzero classes
	modulo $p$. For $x \in S$, let $x' \in S$ be the unique element with
	$x\,x' \equiv a \pmod{p}$; it exists and is unique because $x$ is
	invertible modulo $p$, namely $x' \equiv a x^{-1}$, and $x'
		\not\equiv 0$ because $p \nmid a$. The map $x \mapsto x'$ is an
	involution of $S$, since $(x')' \equiv a(ax^{-1})^{-1} \equiv x$. Its
	fixed points are precisely the $x \in S$ with $x^{2} \equiv a
		\pmod{p}$, and we determine them exactly.

	First suppose $a$ is a non-residue. Then $x^{2} \equiv a$ has no
	solution, so the involution has no fixed point, and consequently
	every orbit $\{x, x'\}$ has exactly two elements. These orbits
	partition $S$, so there are $\tfrac{p-1}{2}$ of them, and the product
	of the two elements of each orbit is congruent to $a$. Multiplying these orbit products gives $(p-1)! \equiv a^{(p-1)/2} \pmod{p}$, and
	Wilson's theorem (Theorem~\ref{thm:3.3.4}) gives $a^{(p-1)/2} \equiv -1$.

	Now suppose $a$ is a quadratic residue and fix $r$ with $r^{2} \equiv
		a \pmod{p}$; then $p \nmid r$, so we may take $r \in S$. The fixed
	points are exactly $r$ and $p - r$. Indeed, $x^{2} \equiv a \equiv
		r^{2}$ means $p \mid (x - r)(x + r)$, so since $p$ is prime, $x \equiv
		r$ or $x \equiv -r \equiv p - r \pmod{p}$; both $r$ and $p - r$ lie in
	$S$, and they are distinct, because $r \equiv -r$ would give $p \mid
		2r$, which is impossible for $p$ odd with $p \nmid r$. Hence the involution has
	exactly two fixed points. The complement $S \setminus \{r, p-r\}$ has
	$p - 3$ elements and is stable under the involution (an involution
	maps non-fixed points to non-fixed points), and it contains no fixed
	point, so it is partitioned into $\tfrac{p-3}{2}$ two-element orbits,
	each with product congruent to $a$; note that $\tfrac{p-3}{2}$ is an integer
	because $p$ is odd. Therefore
	\[
		(p-1)! \equiv r(p-r)\,a^{\frac{p-3}{2}} \equiv -r^{2}a^{\frac{p-3}{2}}
		\equiv -a\,a^{\frac{p-3}{2}} = -a^{\frac{p-1}{2}} \pmod{p},
	\]
	and Wilson's theorem gives $-a^{(p-1)/2} \equiv -1$, that is,
	$a^{(p-1)/2} \equiv 1$.

\end{altproof}
\begin{theorem}\label{thm:7.1.3} Let $p$ be an odd prime and let $g$ be a
	primitive root modulo $p$. For $0\le k\le p-2$, the class $g^{k}$ is a
	quadratic residue modulo $p$ if and only if $k$ is even.

\end{theorem}

\begin{proof} If $k=2j$ is even, then $g^{k}\equiv(g^{j})^{2}$ is a
	quadratic residue. Conversely, let $g^{k}\equiv x^{2}\pmod{p}$. Then
	$p\nmid x$, so $x\equiv g^{i}$ for some $i$, and $g^{2i}\equiv g^{k}$
	forces $2i\equiv k\pmod{p-1}$ because $g$ has order $p-1$
	(Theorem~\ref{thm:6.1.1}(a)--(b)). Writing $k=2i+(p-1)t$ and using the fact that $p-1$ is
	even, we conclude that $k$ is even.

\end{proof}
\subsection{The Legendre Symbol and the Supplementary Laws}
\begin{definition}[Legendre symbol]\label{def:7.2.1} For an odd
	prime $p$ and an integer $a$, the Legendre symbol $\leg{a}{p}$ is $1$ if $a$ is a quadratic
	residue modulo $p$, $-1$ if $a$ is a non-residue, and $0$ if $p \mid a$.

\end{definition}
\begin{theorem}\label{thm:7.2.1} Let $p$ be an odd prime, and suppose $p \nmid
		a$. Then
	\[
		(p-1)! \equiv \begin{cases} -a^{\frac{p-1}{2}} \pmod{p}, & \text{if
	$a$ is a quadratic residue modulo } p,\\[2pt] +a^{\frac{p-1}{2}}
		\pmod{p}, & \text{if $a$ is a quadratic non-residue modulo } p.
		\end{cases}
	\]
	Equivalently, with the Legendre symbol of Definition~\ref{def:7.2.1}, the
	two cases can be written as the single congruence
	\[
		(p-1)! \equiv -\leg{a}{p}a^{\frac{p-1}{2}} \pmod{p}.
	\]
	In particular, the case $a = 1$ gives $(p-1)! \equiv -1 \pmod{p}$; combining the two cases with Wilson's theorem gives $a^{\frac{p-1}{2}} \equiv \leg{a}{p} \pmod{p}$.

\end{theorem}

\begin{proof} Put $S = \{1, 2, \ldots, p-1\}$ and, for $x
		\in S$, let $x' \in S$ be the unique element with $x\,x' \equiv a
		\pmod{p}$, that is, $x' \equiv ax^{-1}$; this is well defined, and $x'
		\neq 0$ because $p \nmid a$. The map $x \mapsto x'$ is an involution
	of $S$ whose fixed points are the solutions of $x^{2} \equiv a
		\pmod{p}$.

	If $a$ is a non-residue, there are no fixed points, so $S$ is
	partitioned into $\tfrac{p-1}{2}$ two-element orbits $\{x, x'\}$,
	each with product congruent to $a$, and multiplying these products gives $(p-1)! \equiv
		a^{(p-1)/2} \pmod{p}$.

	If $a$ is a quadratic residue, fix $r \in S$ with $r^{2} \equiv a$.
	As $x^{2} \equiv r^{2}$ forces $p \mid (x-r)(x+r)$, the fixed points
	are exactly $r$ and $p - r$, and they are distinct because $p \mid 2r$
	is impossible for odd $p$ with $p \nmid r$. The other $p-3$ elements
	form $\tfrac{p-3}{2}$ two-element orbits, each with product congruent to $a$, so
	\[
		(p-1)! \equiv r(p-r)a^{\frac{p-3}{2}} \equiv -r^{2}a^{\frac{p-3}{2}}
		\equiv -a^{\frac{p-1}{2}} \pmod{p}.
	\]
	The two displayed cases combine into the single congruence because
	$\leg{a}{p} = 1$ in the quadratic residue case and $\leg{a}{p} = -1$ in the
	non-residue case. For $a = 1$, which is a quadratic residue since $1 = 1^{2}$,
	the congruence just proved gives $(p-1)! \equiv -1^{(p-1)/2} = -1 \pmod{p}$,
	which is Wilson's theorem (Theorem~\ref{thm:3.3.4}).

\end{proof}
\begin{altproof}[Alternative proof of Theorem~\ref{thm:7.1.2}] Combining Wilson's theorem with the two cases of
	Theorem~\ref{thm:7.2.1}, we obtain
	$-a^{(p-1)/2} \equiv (p-1)! \equiv -1 \pmod{p}$
	for a quadratic residue $a$, and
	$a^{(p-1)/2} \equiv (p-1)! \equiv -1 \pmod{p}$
	for a non-residue.
	Hence $a^{(p-1)/2} \equiv 1$ for residues and $a^{(p-1)/2} \equiv -1$
	for non-residues, which is Euler's criterion (Theorem~\ref{thm:7.1.2}).

\end{altproof}
\begin{theorem}\label{thm:7.2.2} Among the $p - 1$ nonzero
	residues modulo an odd prime $p$, exactly $\tfrac{p-1}{2}$ are
	quadratic residues and $\tfrac{p-1}{2}$ are non-residues; the
	quadratic residues are exactly the even powers $g^{0}, g^{2}, \ldots,
		g^{p-3}$ of a primitive root $g$, equivalently the values $1^{2},
			2^{2}, \ldots, \big(\tfrac{p-1}{2}\big)^{2}$.

\end{theorem}

\begin{proof} Fix a primitive root $g$ modulo $p$
	(Theorem~\ref{thm:6.2.6}).
	Every nonzero residue class is $g^{k}$ for exactly one $k$ with $0 \le
		k \le p-2$, and by Theorem~\ref{thm:7.1.3} the class
	$g^{k}$ is a quadratic residue if and only if $k$ is even. The range
	$0 \le k \le p-2$ contains $p-1$ integers, of which exactly
	$\tfrac{p-1}{2}$ are even, namely $k = 0, 2, \ldots, p-3$; this gives
	both the count $\tfrac{p-1}{2}$ of quadratic residues, hence
	$\tfrac{p-1}{2}$ non-residues for the remaining classes, and the
	description by even powers of $g$.

	For the last description, consider the squares $1^{2}, 2^{2}, \ldots,
		\big(\tfrac{p-1}{2}\big)^{2}$. They are pairwise incongruent modulo
	$p$: if $x^{2} \equiv y^{2}$ with $1 \le x, y \le \tfrac{p-1}{2}$,
	then $p \mid (x-y)(x+y)$, while $|x - y| < p$ and $2 \le x + y \le p -
		1 < p$, so $x = y$. Thus they represent $\tfrac{p-1}{2}$ distinct
	quadratic residues, and since there are only $\tfrac{p-1}{2}$
	quadratic residues in all, they represent every one of them.

\end{proof}
\begin{altproof}[Alternative proof of Theorem~\ref{thm:7.2.2}] The map $x
		\mapsto x^{2}$ from the $p-1$ nonzero classes onto the set of
	quadratic residues is two-to-one: $x^{2} \equiv y^{2} \pmod{p}$ holds if and
	only if $p \mid (x-y)(x+y)$, that is, if and only if $y \equiv \pm x$,
	and the two preimages $\pm x$ are distinct because $x \equiv -x$ would
	give $p \mid 2x$, which is impossible for odd $p$ with $p \nmid x$. Hence the
	image has exactly $\tfrac{p-1}{2}$ elements, and the remaining
	$\tfrac{p-1}{2}$ nonzero classes are the non-residues.

\end{altproof}
\begin{theorem}[multiplication rules]\label{thm:7.2.3} Modulo an
	odd prime,
	\[
		\begin{aligned}
			\text{quadratic residue} \cdot \text{quadratic residue}         & = \text{quadratic residue},     \\
			\text{quadratic residue} \cdot \text{quadratic non-residue}     & = \text{quadratic non-residue}, \\
			\text{quadratic non-residue} \cdot \text{quadratic non-residue} & = \text{quadratic residue}.
		\end{aligned}
	\]

\end{theorem}

\begin{proof} By Euler's criterion, $(ab)^{(p-1)/2} =
		a^{(p-1)/2} b^{(p-1)/2}$, so the $\pm 1$ values multiply according to the usual sign rule.

\end{proof}
\begin{altproof}[Alternative proof of Theorem~\ref{thm:7.2.3}] Let
	$Q$ be the set of quadratic residues and $N$ the set of non-residues
	modulo $p$, so that $\ZmodUnits{p} = Q \sqcup N$ and, by
	Theorem~\ref{thm:7.2.2}, $|Q| = |N| = \tfrac{p-1}{2}$.

	First, $Q \cdot Q \subseteq Q$ follows directly from the definition: if $a
		\equiv x^{2} \pmod{p}$ and $b \equiv y^{2}$, then $ab \equiv (xy)^{2}$.

	Next, let $a \in Q$. Multiplication by $a$ is a bijection of
	$\ZmodUnits{p}$ (its inverse is multiplication by $a^{-1}$,
	which also lies in $Q$, since $a^{-1} \equiv a(a^{-1})^{2} \pmod{p}$ and the right-hand side is a
		product of the residue $a$ with a square). By the previous paragraph,
	$aQ \subseteq Q$, and $|aQ| = |Q|$, so $aQ = Q$. Since multiplication
		by $a$ permutes $\ZmodUnits{p} = Q \sqcup N$ and fixes the
		subset $Q$ setwise, it must map $N$ onto $N$: for $b \in N$ we get $ab
	\in N$. This is the second rule.

		Finally, let $b \in N$. Multiplication by $b$ is again a bijection of
	$\ZmodUnits{p}$, and by the second rule $bQ \subseteq N$;
		since $|bQ| = |Q| = \tfrac{p-1}{2} = |N|$, we get $bQ = N$. As
		multiplication by $b$ permutes $Q \sqcup N$ and sends $Q$ onto $N$, it
		must send $N$ onto $Q$: for $c \in N$ we get $bc \in Q$. This is the
	third rule.

\end{altproof}
\begin{theorem}\label{thm:7.2.4} For an odd
	prime $p$, the congruence $x^{2} \equiv -1 \pmod{p}$ is solvable if and
	only if $p \equiv 1 \pmod{4}$.

\end{theorem}

\begin{proof} By Euler's criterion, $-1$ is a residue if
	and only if $(-1)^{(p-1)/2} = 1$, that is, if and only if
	$\tfrac{p-1}{2}$ is even, that is, if and only if $p \equiv 1 \pmod{4}$.

\end{proof}
\begin{theorem}[supplementary laws for the Legendre symbol]\label{thm:7.2.5} For an odd prime $p$
	and integers $a, b$:

	\begin{enumerate}[label={\alph*)}]
		\item $\leg{a}{p} \equiv a^{(p-1)/2} \pmod{p}$.

		\item If $a \equiv b \pmod{p}$, then $\leg{a}{p} =
			      \leg{b}{p}$.

		\item $\leg{ab}{p} =
			      \leg{a}{p}\leg{b}{p}$.

		\item $\leg{1}{p} = 1$, and $\leg{a^{2}}{p}
			      = 1$ when $p \nmid a$.

		\item (first supplementary law) $\leg{-1}{p} = (-1)^{(p-1)/2}$, that is,
		      $1$ if $p \equiv 1 \pmod{4}$ and $-1$ if $p \equiv 3 \pmod{4}$.

		\item (second supplementary law) $\leg{2}{p} = (-1)^{(p^{2}-1)/8}$, that
		      is, $1$ if $p \equiv \pm 1 \pmod{8}$ and $-1$ if $p \equiv \pm 3
			      \pmod{8}$.
	\end{enumerate}

\end{theorem}

\begin{proof} Property (a) is Euler's criterion (Theorem~\ref{thm:7.1.2}) when $p \nmid a$, and when $p \mid a$ both sides are congruent to $0
		\pmod{p}$. Property (b) holds because whether $a$ is a quadratic residue depends
	only on the residue class of $a$ modulo $p$.

	For (c), we use the multiplication rules (Theorem~\ref{thm:7.2.3}). If $p \mid ab$ then $p \mid a$ or $p \mid b$, so
	both sides are $0$. Otherwise, all three symbols are $\pm 1$, and the
	three rules read: a residue times a residue is a residue, so $1 \cdot
		1 = 1$; a residue times a non-residue is a non-residue, so $1 \cdot
		(-1) = -1$ and $(-1) \cdot 1 = -1$; a non-residue times a non-residue
	is a residue, so $(-1) \cdot (-1) = 1$. In each of the four cases, the
	value of $\leg{ab}{p}$ is the product
	$\leg{a}{p}\leg{b}{p}$.

	Property (d) follows from $1 = 1^{2}$ and, for $p \nmid a$, from the fact
		that $a^{2}$ is itself a square. Assertion (e) follows by applying (a) with $a =
		-1$: both $\leg{-1}{p}$ and $(-1)^{(p-1)/2}$ are $\pm 1$ and they are
	congruent modulo $p > 2$, hence equal; and $(-1)^{(p-1)/2} = 1$ if and
	only if $\tfrac{p-1}{2}$ is even, that is, $p \equiv 1 \pmod{4}$.

	For (f), consider the $\tfrac{p-1}{2}$ even numbers $2k$ with $1 \le k
		\le \tfrac{p-1}{2}$, so that $2 \le 2k \le p-1$. Define
	\[
		(\varepsilon_k, r_k) = \begin{cases} (+1,\ 2k), & \text{if } 2k <
		              \tfrac{p}{2},                \\ (-1,\ p - 2k), & \text{if } 2k > \tfrac{p}{2};
		\end{cases}
	\]
	the case $2k = \tfrac{p}{2}$ cannot occur because $p$ is odd. In both
	cases $r_k \in \{1, \ldots, \tfrac{p-1}{2}\}$ and $2k \equiv
		\varepsilon_k r_k \pmod{p}$. The $r_k$ are pairwise distinct: if $r_k
		= r_l$ then $2k \equiv \pm 2l \pmod{p}$, hence $k \equiv \pm l
		\pmod{p}$; but $1 \le k, l \le \tfrac{p-1}{2}$ gives $|k - l| < p$ and
	$2 \le k + l \le p-1 < p$, so $k = l$. Since these are $\tfrac{p-1}{2}$ distinct
	elements of $\{1, \ldots, \tfrac{p-1}{2}\}$, they exhaust that set.
	Multiplying the congruences $2k \equiv \varepsilon_k r_k$ over $k = 1,
		\ldots, \tfrac{p-1}{2}$ therefore gives
	\[
		2^{\frac{p-1}{2}}\Big(\tfrac{p-1}{2}\Big)! \equiv
		(-1)^{\mu}\Big(\tfrac{p-1}{2}\Big)! \pmod{p}, \qquad \mu = \#\{k \mid
		\varepsilon_k = -1\},
	\]
	and cancelling the invertible factor $\big(\tfrac{p-1}{2}\big)!$
	yields $2^{(p-1)/2} \equiv (-1)^{\mu} \pmod{p}$, so $\leg{2}{p} =
		(-1)^{\mu}$ by (a). Now $\varepsilon_k = -1$ exactly when $k >
		\tfrac{p}{4}$, so $\mu = \tfrac{p-1}{2} - \big\lfloor \tfrac{p}{4}
		\big\rfloor$. Writing $p = 8m + r$ with $r \in \{1, 3, 5, 7\}$, we
	determine the parity in each class:

	\begin{itemize}
		\item $r = 1$: $\mu = 4m - 2m = 2m$ (even), and $\tfrac{p^{2}-1}{8} =
			      8m^{2} + 2m$ (even);
		\item $r = 3$: $\mu = (4m+1) - 2m = 2m+1$ (odd), and
		      $\tfrac{p^{2}-1}{8} = 8m^{2} + 6m + 1$ (odd);
		\item $r = 5$: $\mu = (4m+2) - (2m+1) = 2m+1$ (odd), and
		      $\tfrac{p^{2}-1}{8} = 8m^{2} + 10m + 3$ (odd);
		\item $r = 7$: $\mu = (4m+3) - (2m+1) = 2m+2$ (even), and
		      $\tfrac{p^{2}-1}{8} = 8m^{2} + 14m + 6$ (even).
	\end{itemize}

	In every case, $\mu \equiv \tfrac{p^{2}-1}{8} \pmod{2}$, so
	$\leg{2}{p} = (-1)^{\mu} = (-1)^{(p^{2}-1)/8}$, and this is $+1$
	exactly for $r = 1, 7$, that is, for $p \equiv \pm 1 \pmod{8}$.

\end{proof}
\subsection{Gauss's Lemma and Half-Systems}
\begin{theorem}[Gauss's lemma on half-systems]\label{thm:7.3.1} Let $p$ be an odd
	prime, and suppose $\gcd(a, p) = 1$. For $k = 1, \ldots, \tfrac{p-1}{2}$,
	reduce $ak$ to its least positive residue modulo $p$, and let $\mu$
	be the number of these residues exceeding $\tfrac{p}{2}$ (the letter
	$\mu$ here denotes this count, not the M{\"o}bius function of
	Definition~\ref{def:5.3.6}). Then
	$\leg{a}{p} = (-1)^{\mu}$.

\end{theorem}

\begin{proof} Write $ak \equiv \varepsilon_k r_k \pmod{p}$
	with $r_k \in \{1, \ldots, \tfrac{p-1}{2}\}$ and $\varepsilon_k = \pm
		1$; here $\varepsilon_k = -1$ exactly when the least residue exceeds
	$\tfrac{p}{2}$, so $\mu = \#\{k \mid \varepsilon_k = -1\}$. The $r_k$
	are distinct (if $r_k = r_l$ then $ak \equiv \pm al$, so $k \equiv
		\pm l$, forcing $k = l$), so $\{r_k\} = \{1, \ldots,
		\tfrac{p-1}{2}\}$. Taking the product of $ak \equiv \varepsilon_k
		r_k$ over all $k$ gives $a^{(p-1)/2}\big(\tfrac{p-1}{2}\big)! \equiv
		(-1)^{\mu}\big(\tfrac{p-1}{2}\big)!$, so $a^{(p-1)/2} \equiv
		(-1)^{\mu}$; by Euler's criterion $\leg{a}{p} =
		(-1)^{\mu}$.

\end{proof}
\begin{theorem}[floor-sum form]\label{thm:7.3.2} For an odd prime $p$
	and an odd integer $a$ with $\gcd(a, p) = 1$, we have $\leg{a}{p} =
		(-1)^{\sum_{k=1}^{(p-1)/2} \lfloor ak/p \rfloor}$.

\end{theorem}

\begin{proof} We introduce fresh notation. For $k = 1,
		\ldots, \tfrac{p-1}{2}$, divide $ak$ by $p$ with remainder:
	\[
		ak = p\Big\lfloor \tfrac{ak}{p} \Big\rfloor + s_k, \qquad 0 \le s_k <
		p,
	\]
	and $s_k \neq 0$ because $p \nmid ak$. Since $p$ is odd, $s_k \neq
		\tfrac{p}{2}$, so exactly one of $s_k < \tfrac{p}{2}$, $s_k >
		\tfrac{p}{2}$ holds; define
	\[
		(\varepsilon_k, r_k) = \begin{cases} (+1,\ s_k), & \text{if } s_k <
		              \tfrac{p}{2},                  \\ (-1,\ p - s_k), & \text{if } s_k > \tfrac{p}{2}.
		\end{cases}
	\]
	Then $r_k \in \{1, \ldots, \tfrac{p-1}{2}\}$, $ak \equiv \varepsilon_k
		r_k \pmod{p}$, and $\mu := \#\{k \mid \varepsilon_k = -1\} = \#\{k
		\mid s_k > \tfrac{p}{2}\}$ is exactly the count $\mu$ in Gauss's
	lemma, since $s_k$ is the least positive residue of $ak$. The $r_k$
	are pairwise distinct: $r_k = r_l$ gives $ak \equiv \pm al \pmod{p}$,
	hence $k \equiv \pm l \pmod{p}$ because $\gcd(a, p) = 1$; here $|k -
		l| < p$ and $2 \le k + l \le p - 1 < p$, so $k = l$. Consequently,
	$\{r_1, \ldots, r_{(p-1)/2}\} = \{1, \ldots, \tfrac{p-1}{2}\}$ and in
	particular $\sum_k r_k = \sum_{k=1}^{(p-1)/2} k$.

	Now compare the two expressions for $s_k$: $s_k = r_k$ when
	$\varepsilon_k = 1$ and $s_k = p - r_k$ when $\varepsilon_k = -1$.
	Hence
	\[
		\sum_k s_k = \sum_k r_k + \sum_{\varepsilon_k = -1}(p - 2r_k) \equiv
		\sum_{k=1}^{(p-1)/2} k + \mu \pmod{2},
	\]
	using the fact that $p$ is odd and $2r_k$ is even. Summing the division
	identities over $k$ gives
	\[
		a\sum_{k=1}^{(p-1)/2} k = p\sum_{k=1}^{(p-1)/2} \Big\lfloor
		\tfrac{ak}{p} \Big\rfloor + \sum_k s_k,
	\]
	and reducing modulo $2$ with $a$ and $p$ odd yields
	\[
		\sum_{k=1}^{(p-1)/2} k \equiv \sum_{k=1}^{(p-1)/2} \Big\lfloor
		\tfrac{ak}{p} \Big\rfloor + \sum_{k=1}^{(p-1)/2} k + \mu \pmod{2},
	\]
	so $\mu \equiv \sum_{k=1}^{(p-1)/2} \lfloor ak/p \rfloor \pmod{2}$.
	Gauss's lemma (Theorem~\ref{thm:7.3.1}) gives
	\[
		\leg{a}{p} = (-1)^{\mu} = (-1)^{\sum_k \lfloor ak/p \rfloor}.
	\]

\end{proof}
\begin{definition}[half-system]\label{def:7.3.1} Let $n > 2$ be an
	integer. A subset $I \subset \ZmodUnits{n}$ is a half-system modulo
	$n$ if $I \cup (-I) = \ZmodUnits{n}$ and $I \cap (-I) =
		\varnothing$, that is, $I$ contains exactly one element of $\{a, -a\}$ for every $a
		\in \ZmodUnits{n}$. When we write a half-system using integer representatives, we identify those representatives with their residue classes.

\end{definition}
\begin{theorem}\label{thm:7.3.3} Every
	integer $n>2$ has a half-system. For an odd prime $p$, the set
	$\{1,\ldots,\tfrac{p-1}{2}\}$ is a half-system modulo $p$; it is the
	standard half-system used in Gauss's lemma (Theorem~\ref{thm:7.3.1}) and its
	floor-sum form (Theorem~\ref{thm:7.3.2}).

\end{theorem}

\begin{proof} If a unit $a$ satisfied
	$a\equiv-a\pmod{n}$, then multiplication by its inverse would give
	$2\equiv0\pmod{n}$, contrary to $n>2$. Thus $a\mapsto-a$ is a
	fixed-point-free involution of $\ZmodUnits{n}$; choosing one
	element from each of its two-element orbits gives a half-system. For an
	odd prime $p$, every nonzero class has exactly one representative among
	$\pm1,\ldots,\pm\tfrac{p-1}{2}$, so the displayed set is a
	half-system.

\end{proof}
\begin{theorem}[Gauss's lemma, general form]\label{thm:7.3.4}
	Let $p$ be an odd prime, let $I$ be any half-system modulo $p$
	(Definition~\ref{def:7.3.1}), and let $n$ be an integer with $\gcd(n, p) = 1$. Writing $J = \{i \in I \mid ni \in -I\}$, we
	have $\leg{n}{p} = (-1)^{|J|}$. Taking $I = \{1,
		\ldots, \tfrac{p-1}{2}\}$ recovers Theorem~\ref{thm:7.3.1}.

\end{theorem}

\begin{proof} Note first that $|I| = \tfrac{p-1}{2}$,
	since $I$ and $-I$ are disjoint with union $\ZmodUnits{p}$
	and $|{-I}| = |I|$. Fix $i \in I$. Since $\gcd(n, p) = 1$, the class
	$ni$ lies in $\ZmodUnits{p} = I \cup (-I)$, and exactly one of
	the alternatives $ni \in I$, $ni \in -I$ holds because $I \cap (-I) =
		\varnothing$. Accordingly, define $\eta_i = +1$ and $s_i = ni$
	in the first case, and $\eta_i = -1$ and $s_i = -ni$ in the
	second, with these products taken in $\mathbb{F}_p$; in both cases $s_i \in I$, $\eta_i = \pm 1$, and $ni \equiv
		\eta_i s_i \pmod{p}$, and $\eta_i = -1$ holds exactly when $i \in J$,
	so $\#\{i \mid \eta_i = -1\} = |J|$.

	The $s_i$ are pairwise distinct: if $s_i = s_j$ for $i, j \in I$, then
	$ni \equiv \pm nj \pmod{p}$, so $i \equiv \pm j \pmod{p}$ since
	$\gcd(n, p) = 1$; the case $i \equiv -j$ would put $j \in I \cap (-I)
		= \varnothing$ (because $i \in I$ forces $-i \in -I$), so $i = j$.
	Hence $\{s_i \mid i \in I\}$ is a subset of $I$ with $|I|$ elements,
	that is, $\{s_i \mid i \in I\} = I$, and therefore $\prod_{i \in I}
		s_i = \prod_{i \in I} i$.

	Multiplying the congruences $ni \equiv \eta_i s_i$ over all $i \in I$
	gives
	\[
		n^{\frac{p-1}{2}} \prod_{i \in I} i \equiv \Big(\prod_{i \in I}
		\eta_i\Big) \prod_{i \in I} s_i = (-1)^{|J|} \prod_{i \in I} i
		\pmod{p}.
	\]
	Cancelling $\prod_{i \in I} i$, which is invertible modulo $p$, gives
	$n^{(p-1)/2} \equiv (-1)^{|J|} \pmod{p}$; Euler's criterion
	(Theorem~\ref{thm:7.1.2}) then yields $\leg{n}{p} = (-1)^{|J|}$, since both sides
	are $\pm 1$ and $p > 2$.

\end{proof}
\subsection{Wilson Pairing and Square Roots of \texorpdfstring{$-1$}{-1}}
\begin{theorem}\label{thm:7.4.1} For an odd prime
	$p$,
	\[
		(-1)^{\frac{p-1}{2}}\big[\big(\tfrac{p-1}{2}\big)!\big]^{2} \equiv -1
		\pmod{p}.
	\]

\end{theorem}

\begin{proof} Pair $k$ with $p - k$ for $k = 1, \ldots,
		\tfrac{p-1}{2}$: these pairs partition $\{1, \ldots, p-1\}$, and $p - k
		\equiv -k \pmod{p}$, so Wilson's theorem (Theorem~\ref{thm:3.3.4}) gives
	\[
		-1 \equiv (p-1)! \equiv \prod_{k=1}^{(p-1)/2} k(p-k) \equiv
		(-1)^{\frac{p-1}{2}}\big[\big(\tfrac{p-1}{2}\big)!\big]^{2} \pmod{p}.
	\]

\end{proof}
\begin{theorem}\label{thm:7.4.2} When $p$ is a prime with $p \equiv 1 \pmod{4}$, the integer $x =
		\big(\tfrac{p-1}{2}\big)!$ satisfies $x^{2} \equiv -1 \pmod{p}$.

\end{theorem}

\begin{proof} For $p \equiv 1 \pmod{4}$, the sign
	$(-1)^{(p-1)/2}$ in Theorem~\ref{thm:7.4.1} equals $+1$, so that theorem reads
	$\big[\big(\tfrac{p-1}{2}\big)!\big]^{2} \equiv -1 \pmod{p}$.

\end{proof}
\begin{theorem}\label{thm:7.4.3}
	For a prime $p \equiv 3 \pmod{4}$, we have
	$\big[\big(\tfrac{p-1}{2}\big)!\big]^{2} \equiv 1 \pmod{p}$, hence
	$\big(\tfrac{p-1}{2}\big)! \equiv \pm 1 \pmod{p}$; more precisely,
	$\big(\tfrac{p-1}{2}\big)! \equiv (-1)^{v} \pmod{p}$, where $v$ is the
	number of integers in $\{1, \ldots, \tfrac{p-1}{2}\}$ that are
	quadratic non-residues modulo $p$.

\end{theorem}

\begin{proof} Write $P = \big(\tfrac{p-1}{2}\big)!$. By
	Theorem~\ref{thm:7.4.1}, $(-1)^{(p-1)/2}P^{2} \equiv -1 \pmod{p}$. For $p \equiv
		3 \pmod{4}$, the sign equals $(-1)^{(p-1)/2} = -1$, so $P^{2} \equiv 1$ and
	therefore $P \equiv \pm 1 \pmod{p}$ (Theorem~\ref{thm:2.3.2} applied to $p \mid
		(P-1)(P+1)$). It remains to
	identify the sign. By Euler's criterion, $k^{(p-1)/2} \equiv
		\leg{k}{p} \pmod{p}$ for each $k = 1, \ldots, \tfrac{p-1}{2}$;
	multiplying these congruences together gives
	\[
		P^{\frac{p-1}{2}} \equiv \prod_{k=1}^{(p-1)/2}\leg{k}{p} = (-1)^{v}
		\pmod{p},
	\]
	since exactly $v$ of the factors equal $-1$ and the rest equal $+1$.
	Since $p \equiv 3 \pmod{4}$, the exponent $\tfrac{p-1}{2}$ is odd, and
	$P \equiv \pm 1$, so $P^{(p-1)/2} \equiv P$ (raising either $1$ or $-1$ to an odd power leaves it unchanged). Hence $P \equiv (-1)^{v} \pmod{p}$.

\end{proof}
\subsection{The Law of Quadratic Reciprocity and Its Applications}

\begin{theorem}[quadratic reciprocity]\label{thm:7.5.1} For
	distinct odd primes $p, q$,
	we have $\leg{p}{q}\leg{q}{p} =
		(-1)^{\frac{p-1}{2}\cdot\frac{q-1}{2}}$; equivalently,
	$\leg{p}{q} = \leg{q}{p}$ unless $p
		\equiv q \equiv 3 \pmod{4}$, in which case they have opposite signs.

\end{theorem}

\begin{proof} By the
	floor-sum form of Gauss's lemma (Theorem~\ref{thm:7.3.2}),
	\[
		\leg{q}{p} = (-1)^{S_1}, \quad S_1 = \sum_{x=1}^{(p-1)/2} \Big\lfloor
		\tfrac{qx}{p} \Big\rfloor, \qquad \leg{p}{q} = (-1)^{S_2}, \quad S_2 =
		\sum_{y=1}^{(q-1)/2} \Big\lfloor \tfrac{py}{q} \Big\rfloor,
	\]
	where Theorem~\ref{thm:7.3.2} applies because $q$ is odd and coprime to $p$, and likewise
	$p$ is odd and coprime to $q$. It therefore suffices to prove $S_1 + S_2 =
		\tfrac{p-1}{2}\cdot\tfrac{q-1}{2}$, and we do so by counting the
	lattice points in the rectangle
	\[
		E = \Big\{(x, y) \in \mathbb{Z}^{2} \ \Big|\ 1 \le x \le
		\tfrac{p-1}{2},\ 1 \le y \le \tfrac{q-1}{2}\Big\}, \qquad |E| =
		\tfrac{p-1}{2}\cdot\tfrac{q-1}{2},
	\]
	in two ways.

	No point of $E$ lies on the line $py = qx$: if $py = qx$ then $p \mid
		qx$, and $\gcd(p, q) = 1$ forces $p \mid x$, which is impossible for $1 \le x
		\le \tfrac{p-1}{2}$. So each point of $E$ lies strictly below or
	strictly above that line, and it suffices to show that $S_1$ counts
	the points below and $S_2$ the points above.

	Fix $x$ with $1 \le x \le \tfrac{p-1}{2}$. A point $(x, y)$ lies below
	the line, that is, $y < qx/p$, exactly when $1 \le y \le \lfloor qx/p
		\rfloor$, because $qx/p$ is not an integer. All the corresponding points $(x,y)$ lie in
	the rectangle: $qx/p \le \tfrac{q(p-1)}{2p} < \tfrac{q}{2}$, so
	$\lfloor qx/p \rfloor \le \tfrac{q-1}{2}$. Hence the column with
	abscissa $x$ contains exactly $\lfloor qx/p \rfloor$ points below the
	line, and summing over $x$ shows that $S_1$ is the total number of
	points of $E$ below the line. Interchanging the roles of $p$ and $q$,
	that is, fixing $y$ and counting the $x$ with $x < py/q$, shows in the
	same way that $S_2$ is the number of points of $E$ above the line.
	Therefore $S_1 + S_2 = |E| = \tfrac{p-1}{2}\cdot\tfrac{q-1}{2}$, which
	proves the law.

\end{proof}
\begin{altproof}[Alternative proof of Theorem~\ref{thm:7.5.1}] Write $G = \ZmodUnits{pq}$; since $\gcd(p, q)
		= 1$, the Chinese remainder theorem (Theorem~\ref{thm:3.4.1}) provides a group
	isomorphism $f : \ZmodUnits{p} \times \ZmodUnits{q}
		\to G$ under which $(-1, -1)$ corresponds to $-1$, so that $f(-a, -b)
		= -f(a, b)$, and $|G| = (p-1)(q-1)$.

	As $pq > 2$, the map $x \mapsto
		-x$ is a fixed-point-free involution of $G$, so the half-systems modulo
	$pq$ in the sense of Definition~\ref{def:7.3.1} are exactly the subsets containing
	exactly one element of each pair $\{x, -x\}$, and each of them has
	$\tfrac{(p-1)(q-1)}{2}$ elements.

	If $H$ and $H'$
	are half-systems, then each $x \in H'$ equals $\pm y$ for exactly one
	$y \in H$, whence $\prod_{x \in H'} x = (-1)^{m} \prod_{y \in H} y$ in
	$G$, where $m = \#\{y \in H \mid -y \in H'\}$; this is a single
	identity in $G$, so under $f^{-1}$ the same scalar $(-1)^{m}$ relates
	the two products in each of the two coordinates.

	We apply this identity to two explicit
	half-systems: $H_1 = \{\, k \in \mathbb{Z} \mid 1 \le k \le
		\tfrac{pq-1}{2},\ p \nmid k,\ q \nmid k \,\}$, which is a half-system
	because $k \mapsto pq - k$ is a bijection from $\{1, \ldots,
		\tfrac{pq-1}{2}\}$ onto $\{\tfrac{pq+1}{2}, \ldots, pq-1\}$
	preserving coprimality with $pq$, so that each pair $\{x, -x\}$ of $G$
	has exactly one representative in the lower half; and $H_2 =
		f\big(\ZmodUnits{p} \times \{1, \ldots,
		\tfrac{q-1}{2}\}\big)$, which is a half-system because for $1 \le b
		\le \tfrac{q-1}{2}$ the residue $-b \bmod q$ lies in
	$\{\tfrac{q+1}{2}, \ldots, q-1\}$, so $H_2 \cap (-H_2) = \varnothing$,
	while $|H_2| = \tfrac{(p-1)(q-1)}{2}$ forces $H_2 \cup (-H_2) = G$.

	Now we compute the two products in both coordinates. Modulo $p$, the
	first coordinate of $H_2$ runs over all of $\ZmodUnits{p}$
	once for each of the $\tfrac{q-1}{2}$ admissible values of $b$, so
	Wilson's theorem (Theorem~\ref{thm:3.3.4}) gives $\prod_{x \in H_2} x \equiv
		\big((p-1)!\big)^{(q-1)/2} \equiv (-1)^{(q-1)/2} \pmod{p}$. Modulo $q$
	the second coordinate takes each value $b \in \{1, \ldots,
		\tfrac{q-1}{2}\}$ exactly $p-1$ times, so $\prod_{x \in H_2} x \equiv
		\big(\big(\tfrac{q-1}{2}\big)!\big)^{p-1} \pmod{q}$. Since
	$(q-1)! = \big(\tfrac{q-1}{2}\big)! \prod_{b=1}^{(q-1)/2}(q - b)
		\equiv (-1)^{(q-1)/2}\big(\big(\tfrac{q-1}{2}\big)!\big)^{2}
		\pmod{q}$, Wilson's theorem turns this into
	$\big(\big(\tfrac{q-1}{2}\big)!\big)^{2} \equiv (-1)^{(q+1)/2}
		\pmod{q}$ and hence $\prod_{x \in H_2} x \equiv
		\big(\big(\tfrac{q-1}{2}\big)!\big)^{2 \cdot \frac{p-1}{2}} \equiv
		(-1)^{\frac{q+1}{2}\cdot\frac{p-1}{2}} \pmod{q}$.

	For $H_1$, use
	$\tfrac{pq-1}{2} = \tfrac{q-1}{2}\,p + \tfrac{p-1}{2}$.
	Split the integers in $[1, \tfrac{pq-1}{2}]$ into
	$\tfrac{q-1}{2}$ complete blocks of $p$ consecutive integers and a
	final block of $\tfrac{p-1}{2}$ integers. In each complete block,
	the integers coprime to $p$ contribute the residues $1, \ldots, p-1$,
	whose product is $(p-1)! \equiv -1 \pmod{p}$. The final block
	contributes the residues $1, \ldots, \tfrac{p-1}{2}$.
	Hence the product of the integers in the interval that are coprime
	to $p$ is congruent to $(-1)^{(q-1)/2}\big(\tfrac{p-1}{2}\big)! \pmod{p}$. From these integers we must
	still delete the multiples of $q$, namely $q, 2q, \ldots,
		\tfrac{p-1}{2}q$, since $\big\lfloor \tfrac{pq-1}{2q} \big\rfloor =
		\tfrac{p-1}{2}$. Their product is congruent to
		$q^{(p-1)/2}\big(\tfrac{p-1}{2}\big)! \pmod{p}$, so dividing gives
	$\prod_{k \in H_1} k \equiv (-1)^{(q-1)/2}q^{-(p-1)/2} \equiv
		(-1)^{(q-1)/2}\leg{q}{p} \pmod{p}$, where the last step follows from Euler's criterion
	(Theorem~\ref{thm:7.1.2}) together with $\leg{q}{p} = \pm 1$, whose inverse is
	itself. The same computation with $p$ and $q$ interchanged, based on
	$\tfrac{pq-1}{2} = \tfrac{p-1}{2}\,q + \tfrac{q-1}{2}$, gives
	$\prod_{k \in H_1} k \equiv (-1)^{(p-1)/2}\leg{p}{q} \pmod{q}$.
	Writing $\prod_{k \in H_1} k = (-1)^{m}\prod_{x \in H_2} x$ and
	reading off the two coordinates, we obtain
	$(-1)^{(q-1)/2}\leg{q}{p} \equiv (-1)^{m}(-1)^{(q-1)/2} \pmod{p}$ and
	$(-1)^{(p-1)/2}\leg{p}{q} \equiv
		(-1)^{m}(-1)^{\frac{q+1}{2}\frac{p-1}{2}} \pmod{q}$; in each of these congruences
	both sides are $\pm 1$ and the moduli are primes greater than $2$, so they
	are genuine equalities of signs, giving $(-1)^{m} = \leg{q}{p}$ from
	the first and $(-1)^{m} =
		(-1)^{(p-1)/2}\leg{p}{q}(-1)^{\frac{q+1}{2}\frac{p-1}{2}}$ from the
	second, whence
	\[
		\leg{q}{p}\leg{p}{q} = (-1)^{\frac{p-1}{2} +
			\frac{q+1}{2}\frac{p-1}{2}} = (-1)^{\frac{p-1}{2}\cdot\frac{q-1}{2}},
	\]
	where the last equality holds because $\frac{p-1}{2} +
		\frac{q+1}{2}\frac{p-1}{2} = \frac{p-1}{2}\cdot\frac{q-1}{2} + (p-1)$
	with $p-1$ even. This proves the law.

\end{altproof}
\begin{altproof}[Alternative proof of Theorem~\ref{thm:7.5.1}]
	Let
	\[
		E = \Big\{(x, y)\in\mathbb{Z}^{2} \Big| 1 \le x \le \tfrac{p-1}{2},\ 1 \le y \le
		\tfrac{q-1}{2}\Big\},
	\]
	so $|E| = \tfrac{p-1}{2}\cdot\tfrac{q-1}{2}$, and split $E$ according
	to the value of $v(x, y) = qx - py$:
	\[
		\begin{aligned} E_1 & = \{v < -\tfrac{p}{2}\}, & E_2 & =
		                \{-\tfrac{p}{2} < v < 0\},               \\ E_3 & = \{0 < v < \tfrac{q}{2}\}, & E_4
		                    & = \{v > \tfrac{q}{2}\}.\end{aligned}
	\]
	These four sets are pairwise disjoint and cover $E$: indeed, $v \neq 0$
	(else $p \mid qx$, so $p \mid x$, which is impossible), and $v \neq
		-\tfrac{p}{2}$, $v \neq \tfrac{q}{2}$ because $p, q$ are odd while $v$
	is an integer; a negative $v$ satisfies exactly one of $v < -p/2$, $-p/2
		< v < 0$, and a positive $v$ exactly one of $0 < v < q/2$, $v > q/2$.
	Write $t_i = |E_i|$, so $t_1 + t_2 + t_3 + t_4 = |E|$.

	We claim that $\leg{q}{p} = (-1)^{t_2}$. Fix $x$ with $1 \le x \le
		\tfrac{p-1}{2}$ and let $s$ be the least positive residue of $qx$
	modulo $p$, so $s \neq 0$ and $s \neq \tfrac{p}{2}$. A point $(x, y)
		\in E$ lies in $E_2$ exactly when $qx < py < qx + \tfrac{p}{2}$, that
	is, when $\tfrac{qx}{p} < y < \tfrac{qx}{p} + \tfrac12$. An open
	interval of length $\tfrac12$ contains at most one integer, and the interval above
	contains one precisely when the next integer above $qx/p$, namely $y_0
		= \lfloor qx/p \rfloor + 1$, satisfies $y_0 < qx/p + \tfrac12$, that
	is, when the fractional part $\{qx/p\} = s/p$ exceeds $\tfrac12$, that
	is, when $s > \tfrac{p}{2}$. In that case $y_0$ automatically lies in
	the admissible range: $y_0 > qx/p > 0$ gives $y_0 \ge 1$, and $y_0 <
		qx/p + \tfrac12 < \tfrac{q}{2} + \tfrac12$ gives $y_0 \le
		\tfrac{q-1}{2}$, using $qx/p \le \tfrac{q(p-1)}{2p} < \tfrac{q}{2}$.
	Hence each $x$ contributes exactly one point to $E_2$ when $s >
		\tfrac{p}{2}$ and none otherwise, so $t_2$ equals the number $\mu$ in
	Gauss's lemma for $a = q$, and Theorem~\ref{thm:7.3.1} gives $\leg{q}{p} =
		(-1)^{t_2}$. Exchanging the roles of $p$ and $q$ (that is, replacing
	$v$ by $-v$) gives $\leg{p}{q} = (-1)^{t_3}$ in the same way.

	Finally, the map $\iota(x, y) = \big(\tfrac{p+1}{2} - x,\
		\tfrac{q+1}{2} - y\big)$ sends $E$ to $E$ and satisfies $\iota \circ
		\iota = \mathrm{id}$, and
	\[
		v(\iota(x, y)) = q\tfrac{p+1}{2} - qx - p\tfrac{q+1}{2} + py =
		\tfrac{q-p}{2} - v(x, y).
	\]
	If $v < -\tfrac{p}{2}$ then $v(\iota(x,y)) > \tfrac{q-p}{2} +
		\tfrac{p}{2} = \tfrac{q}{2}$, so $\iota(E_1) \subseteq E_4$, and
	symmetrically $\iota(E_4) \subseteq E_1$; since $\iota$ is an involution,
	it restricts to a bijection from $E_1$ onto $E_4$, and hence $t_1 = t_4$.
	Consequently
	\[
		\leg{p}{q}\leg{q}{p} = (-1)^{t_2 + t_3} = (-1)^{t_1 + t_2 + t_3 +
			t_4} = (-1)^{|E|} = (-1)^{\frac{p-1}{2}\frac{q-1}{2}},
	\]
	since $t_1 + t_4 = 2t_1$ is even.

\end{altproof}
\begin{theorem}\label{thm:7.5.2} For an odd prime $p$:

	\begin{enumerate}[label={\alph*)}]
		\item $\leg{-1}{p} = 1$ if and only if $p \equiv 1 \pmod{4}$.
		\item $\leg{2}{p} = 1$ if and only if $p \equiv \pm 1 \pmod{8}$.
		\item $\leg{-2}{p} = 1$ if and only if $p \equiv 1, 3 \pmod{8}$.
		\item $\leg{3}{p} = 1$ if and only if $p \equiv \pm 1
			      \pmod{12}$ (for $p \neq 3$).
		\item $\leg{-3}{p} = 1$ if and only if $p \equiv 1 \pmod{3}$
		      (for $p \neq 3$).
		\item $\leg{5}{p} = 1$ if and only if $p \equiv \pm 1 \pmod{5}$ (for $p \neq
			      5$).
		\item $\leg{-2}{p} = \leg{3}{p} = 1$
			      if and only if $p \equiv 1, 11 \pmod{24}$ (for $p \neq 3$).
	\end{enumerate}

	For $p = 3$ in (d), (e), (g) and for $p = 5$ in (f), each symbol whose
	numerator is divisible by $p$ is $0$.

\end{theorem}

\begin{proof} Parts (a) and (b) are the two supplementary laws,
	that is, properties (e) and (f) of Theorem~\ref{thm:7.2.5}.

	(c) By multiplicativity, $\leg{-2}{p} =
		\leg{-1}{p}\leg{2}{p}$, which equals $+1$ exactly when the two
	factors agree. They are both $+1$ when $p \equiv 1 \pmod{4}$ and $p
		\equiv \pm 1 \pmod{8}$, which, since $7 \equiv 3 \pmod{4}$, leaves $p
		\equiv 1 \pmod{8}$; they are both $-1$ when $p \equiv 3 \pmod{4}$ and
	$p \equiv \pm 3 \pmod{8}$, which, since $5 \equiv 1 \pmod{4}$, leaves $p
		\equiv 3 \pmod{8}$. Hence $\leg{-2}{p} = 1$ if and only if $p \equiv 1, 3
		\pmod{8}$.

	(d) Let $p \neq 3$. Quadratic reciprocity (Theorem~\ref{thm:7.5.1}) gives
	$\leg{3}{p} = (-1)^{\frac{p-1}{2}\frac{3-1}{2}}\leg{p}{3} =
		(-1)^{\frac{p-1}{2}}\leg{p}{3}$. Here $\leg{p}{3} = 1$ if and only if $p \equiv 1
		\pmod{3}$, and $\leg{p}{3} = -1$ if and only if $p \equiv 2 \pmod{3}$, because $1$
	is the only nonzero square modulo $3$; and $(-1)^{\frac{p-1}{2}} = 1$
		if and only if $p \equiv 1 \pmod{4}$, while $(-1)^{\frac{p-1}{2}} = -1$ if and only if $p \equiv 3 \pmod{4}$. The
	product is $+1$ exactly when the two factors have the same sign: either $p
		\equiv 1 \pmod{3}$ and $p \equiv 1 \pmod{4}$, which by the Chinese
	remainder theorem means $p \equiv 1 \pmod{12}$; or $p \equiv 2
		\pmod{3}$ and $p \equiv 3 \pmod{4}$, which means $p \equiv 11 \equiv
		-1 \pmod{12}$. The two remaining classes coprime to $12$, namely $p
		\equiv 5, 7 \pmod{12}$, give factors of opposite sign and hence
	$\leg{3}{p} = -1$. Therefore $\leg{3}{p} = 1$ if and only if $p \equiv \pm 1
		\pmod{12}$.

	(e) For $p \neq 3$, by (d) and multiplicativity,
	\[
		\leg{-3}{p} = \leg{-1}{p}\leg{3}{p} =
		(-1)^{\frac{p-1}{2}}(-1)^{\frac{p-1}{2}}\leg{p}{3} = \leg{p}{3},
	\]
	which is $1$ if and only if $p \equiv 1 \pmod{3}$.

	(f) Let $p \neq 5$. Since $5 \equiv 1 \pmod{4}$, the exponent
	$\frac{p-1}{2}\frac{5-1}{2} = p-1$ is even, so reciprocity gives
	$\leg{5}{p} = \leg{p}{5}$. The nonzero squares modulo $5$ are $1$ and
	$4$, that is, $\pm 1$, so $\leg{p}{5} = 1$ if and only if $p \equiv \pm 1
		\pmod{5}$.

	(g) Let $p \neq 3$. By (c) and (d), the two conditions $\leg{-2}{p} =
		1$ and $\leg{3}{p} = 1$ say that $p \equiv 1, 3 \pmod{8}$ and $p
		\equiv \pm 1 \pmod{12}$. Since $p$ is coprime to $24$, we may combine these conditions
	by the Chinese remainder theorem modulo $24$: the condition $p \equiv
		1, 3 \pmod{8}$ with $p$ odd and $3 \nmid p$ leaves $p \equiv 1, 11,
		17, 19 \pmod{24}$, while $p \equiv \pm 1 \pmod{12}$ leaves $p \equiv
		1, 11, 13, 23 \pmod{24}$. The classes satisfying both are $p \equiv 1$
	(from $p \equiv 1 \pmod{8}$ and $p \equiv 1 \pmod{12}$) and $p \equiv
		11 \pmod{24}$ (from $p \equiv 3 \pmod{8}$ and $p \equiv -1
		\pmod{12}$); all other combinations are inconsistent. Hence
	$\leg{-2}{p} = \leg{3}{p} = 1$ if and only if $p \equiv 1, 11 \pmod{24}$.

\end{proof}
\begin{theorem}\label{thm:7.5.3} There are infinitely many primes $p \equiv 1 \pmod{4}$.

\end{theorem}

\begin{proof} Suppose $p_1, \ldots, p_k$ are all the primes congruent to $1$ modulo $4$,
	and put $N = (2 p_1 \cdots p_k)^2 + 1$. Then $N$ is odd and greater than $1$, so
	it has an odd prime factor $q$. From $(2p_1\cdots p_k)^2 \equiv -1
		\pmod{q}$ we get $\leg{-1}{q} = 1$, hence $q \equiv 1
		\pmod{4}$. But $N \equiv 1 \pmod{p_i}$, so $q \neq p_i$. Thus $q$ is a new prime
	congruent to $1$ modulo $4$, a contradiction.

\end{proof}

\begin{theorem}\label{thm:7.5.4} There are infinitely many primes $p \equiv 1 \pmod{3}$.

\end{theorem}

\begin{proof} Suppose $p_1, \ldots, p_k$ are all such
	primes, and set $n = 3 p_1 \cdots p_k$ and $N = n^2 + n + 1$. Since $n$
	is odd, $N$ is odd and $4N = (2n+1)^2 + 3$. Let $q$ be a prime divisor of $N$ (necessarily odd, since $N$ is odd); then
	$(2n+1)^2 \equiv -3 \pmod{q}$. If $q = 3$ then $3 \mid (2n+1)^2$,
	which is impossible since $3 \mid n$ gives $2n+1 \equiv 1 \pmod{3}$. Hence $q
		\neq 3$ and $\leg{-3}{q} = 1$, so $q \equiv 1 \pmod{3}$. Finally, $N \equiv 1 \pmod{3}$ and $N \equiv 1 \pmod{p_i}$, so $q
		\notin \{3, p_1, \ldots, p_k\}$. Thus $q$ is a new prime congruent to $1 \pmod{3}$,
	a contradiction.

\end{proof}

\subsection{Jacobi and Kronecker Symbols; Taking Square Roots}

\begin{definition}[Jacobi symbol]\label{def:7.6.1} For a positive
	odd integer $m = p_1^{\alpha_1} \cdots p_k^{\alpha_k}$ and an integer
	$a$, put $\leg{a}{m} := \prod_{i=1}^{k}\leg{a}{p_i}^{\alpha_i}$, with $\leg{a}{1} = 1$; the
	factors are Legendre symbols, and when $m$ is an odd prime, the two symbols coincide.

\end{definition}
\begin{theorem}\label{thm:7.6.1}
	For positive odd $m, n$ and integers $a, b$:

	\begin{enumerate}[label={\alph*)}]
		\item if $a \equiv b \pmod{m}$, then $\leg{a}{m} =
			      \leg{b}{m}$;
		\item $\leg{a}{m} = 0$ if and only if $\gcd(a, m) > 1$; otherwise, its value is $\pm 1$;
		\item $\leg{ab}{m} = \leg{a}{m}\leg{b}{m}$ and $\leg{a}{mn} = \leg{a}{m}\leg{a}{n}$;
		\item (reciprocity, flipped form, valid for all positive odd $m, n$)
		      $\leg{m}{n} = (-1)^{\frac{m-1}{2}\frac{n-1}{2}}\leg{n}{m}$; when
		      $\gcd(m, n) = 1$ this is equivalent to the product form
		      $\leg{m}{n}\leg{n}{m} = (-1)^{\frac{m-1}{2}\frac{n-1}{2}}$, whereas
		      when $\gcd(m, n) > 1$ both symbols are $0$ by property (b), so the flipped reciprocity form still holds (since both sides are $0$), while the product form fails because its right-hand side is $\pm 1$;
		\item (supplementary laws) $\leg{-1}{m} =
			      (-1)^{\frac{m-1}{2}}$ and $\leg{2}{m} = (-1)^{\frac{m^2-1}{8}}$.
	\end{enumerate}

\end{theorem}

\begin{proof} Properties (a)--(c) are immediate from the
	definition and the properties of the Legendre symbol. For (d) and (e), use the parity identities for
	odd $x, y$, namely $\frac{x-1}{2} + \frac{y-1}{2}
		\equiv \frac{xy-1}{2} \pmod{2}$ and $\frac{x^2-1}{8} + \frac{y^2-1}{8}
		\equiv \frac{(xy)^2-1}{8} \pmod{2}$ (the first because $(x-1)(y-1)
		\equiv 0 \pmod{4}$, the second because $8$ divides both $x^2-1$ and $y^2-1$), which
	extend by induction to any number of odd factors. Writing $m=\prod_{i=1}^{r}p_i$ and
	$n=\prod_{j=1}^{s}q_j$ as products of odd primes, where repetitions
	are allowed so that multiplicities are retained, we obtain
	$\leg{-1}{m}=\prod_{i=1}^{r} (-1)^{\frac{p_i-1}{2}} =
		(-1)^{\frac{m-1}{2}}$ and $\leg{2}{m} =
		(-1)^{\frac{m^2-1}{8}}$; and, when $\gcd(m, n) = 1$, we obtain $\leg{m}{n}\leg{n}{m}=\prod_{i=1}^{r}\prod_{j=1}^{s}(-1)^{\frac{p_i-1}{2}\frac{q_j-1}{2}} =
		(-1)^{\frac{m-1}{2}\frac{n-1}{2}}$ using quadratic reciprocity (Theorem~\ref{thm:7.5.1}) applied to each pair of distinct primes $p_i \neq q_j$; this is
	the product form of (d). The flipped form of (d) follows in
	full generality: if $\gcd(m, n) = 1$, multiply the product form by
	$\leg{n}{m}$ and use $\leg{n}{m}^{2} = 1$, which holds because
	$\leg{n}{m} = \pm 1$ by property (b); if $\gcd(m, n) > 1$, then property (b)
	gives $\leg{m}{n} = \leg{n}{m} = 0$ and both sides of the flipped form
	are $0$.

\end{proof}
\begin{theorem}\label{thm:7.6.2} Let $n > 1$ be odd and let $a$ be an integer with
	$\gcd(a, n) = 1$; in accordance with Definition~\ref{def:7.1.1}, call $a$ a
	quadratic residue modulo $n$ if $x^{2} \equiv a \pmod{n}$ is solvable
	and a quadratic non-residue otherwise. (The coprimality assumption
	entails no loss of generality here: by property (b) of Theorem~\ref{thm:7.6.1} a symbol
	$\leg{a}{n}$ equal to $\pm 1$ already forces $\gcd(a, n) = 1$, and for
	$\gcd(a, n) > 1$ the symbol is $0$ and the terms residue and
	non-residue are not used.) Then:

	\begin{enumerate}[label={\alph*)}]
		\item if $\leg{a}{n} = -1$, then $a$ is not a quadratic residue modulo
		      $n$;
		\item in general, the converse fails: for $n = 15$ and $a = 2$ one has
		      $\leg{2}{15} = 1$, yet $2$ is not a quadratic residue modulo $15$.
	\end{enumerate}

\end{theorem}

\begin{proof} (a) Suppose $x^{2} \equiv a \pmod{n}$ for
	some $x$, and write $n = p_1^{\alpha_1}\cdots p_k^{\alpha_k}$ with odd
	primes $p_i$. For each $i$, we have $x^{2} \equiv a \pmod{p_i}$, so
	$\leg{a}{p_i} \in \{0, 1\}$, and hence, by Definition~\ref{def:7.6.1},
	\[
		\leg{a}{n} = \prod_{i=1}^{k}\leg{a}{p_i}^{\alpha_i} \in \{0, 1\}.
	\]
	Thus $\leg{a}{n} = -1$ is incompatible with the solvability of $x^{2}
		\equiv a \pmod{n}$, which is assertion (a) in contrapositive form.

	(b) Take $n = 15$ and $a = 2$. Then $\leg{2}{15} =
		\leg{2}{3}\leg{2}{5} = (-1)(-1) = 1$, since $2$
	is a non-residue modulo $3$ and modulo $5$. On the other hand, $x^{2}
		\equiv 2 \pmod{15}$ would force $x^{2} \equiv 2 \pmod{3}$, whereas the
	squares modulo $3$ are $0$ and $1$; so no solution exists and $2$ is
	not a quadratic residue modulo $15$. The mechanism is general: an even
	number of factors $-1$ in Definition~\ref{def:7.6.1} cancel and thereby mask
	the non-residues.

\end{proof}
\begin{theorem}\label{thm:7.6.3} Let $n \ge 1$ be odd and let $a$ be an
	integer. Define finite sequences $(n_i)_{i \ge 0}$, $(a_i)_{i \ge 0}$,
	$(e_i)$, $(u_i)$ of integers as follows:
	\[
		n_0 = n, \qquad a_0 = a \bmod n_0,
	\]
	and, as long as $n_i > 1$ and $a_i \neq 0$, write the unique
	factorization
	\[
		a_i = 2^{e_i}u_i, \qquad e_i \ge 0, \quad u_i \text{ odd}, \quad u_i
		\ge 1,
	\]
	and set
	\[
		n_{i+1} = u_i, \qquad a_{i+1} = n_i \bmod u_i,
	\]
	where $x \bmod y$ denotes the least nonnegative remainder. The
	recursion stops at the first index $t$ with $n_t = 1$ or $a_t = 0$.
	Such a $t$ exists because $0 \le a_{i+1} < n_{i+1} = u_i \le a_i <
		n_i$ shows that the odd positive integers $n_0 > n_1 > \cdots$
	strictly decrease; moreover $t \le \log_2 n$. With these data,
	\[
		\leg{a}{n} = \begin{cases} 0, & \text{if } a_t = 0 \text{ and } n_t >
		              1,                                                        \\[4pt] \displaystyle (-1)^{\sum_{i=0}^{t-1}\left(e_i\,\frac{n_i^{2}
				                                                                                - 1}{8} + \frac{u_i - 1}{2}\cdot\frac{n_i - 1}{2}\right)}, & \text{if
		                                                                                                                                               } n_t = 1.\end{cases}
	\]
	The formula uses no factorization of $a$ or $n$ beyond removing
	powers of $2$. Each step consists of one division with remainder,
	so the procedure is a variant of the Euclidean algorithm with an
	accompanying sign calculation. In
	particular, when $n = p$ is an odd prime, the Jacobi symbol coincides
	with the Legendre symbol (Definition~\ref{def:7.6.1}), so the formula
	evaluates $\leg{a}{p}$ without factoring and without modular
	exponentiation.

\end{theorem}

\begin{proof} All steps rest on Theorem~\ref{thm:7.6.1}, and each
	line below is an identity of Jacobi symbols with an odd positive lower
	argument, so the recursion never leaves the domain of definition.

	First, $\leg{a}{n} = \leg{a_0}{n_0}$ by property (a) (invariance modulo
	the lower argument). Next, fix $i$ with $n_i > 1$ and $a_i \neq 0$.
	Property (c) (complete multiplicativity in the upper argument) and
	property (e) give
	\[
		\leg{a_i}{n_i} = \leg{2}{n_i}^{e_i}\leg{u_i}{n_i} =
		(-1)^{e_i\frac{n_i^{2}-1}{8}}\leg{u_i}{n_i},
	\]
	and reciprocity (property (d)), applied to the odd positive integers
	$u_i$ and $n_i$, gives
	\[
		\leg{u_i}{n_i} = (-1)^{\frac{u_i-1}{2}\frac{n_i-1}{2}}\leg{n_i}{u_i} =
		(-1)^{\frac{u_i-1}{2}\frac{n_i-1}{2}}\leg{a_{i+1}}{n_{i+1}},
	\]
	where the last equality again follows from property (a), since $n_{i+1} = u_i$ and
	$a_{i+1} = n_i \bmod u_i$. Here property (d) is used in its flipped form,
	which is valid for all positive odd $u_i, n_i$ with no coprimality
	hypothesis; in particular, the case $\gcd(u_i, n_i) > 1$, where both
	symbols vanish, needs no separate treatment. If $u_i = 1$, then
	$\leg{u_i}{n_i} = 1$, the sign $(-1)^{\frac{u_i-1}{2}\frac{n_i-1}{2}}$
	is $1$, $n_{i+1} = 1$, and $a_{i+1} = 0$, so the identity again holds and
	the recursion terminates at the next index.

	Successively substituting the identities for $i = 0, 1, \ldots, t-1$ gives
	\[
		\leg{a}{n} =
		(-1)^{\sum_{i=0}^{t-1}\left(e_i\frac{n_i^{2}-1}{8} +
			\frac{u_i-1}{2}\frac{n_i-1}{2}\right)}\leg{a_t}{n_t}.
	\]
	It remains to evaluate the final symbol. If $n_t = 1$ then
	$\leg{a_t}{1} = 1$ by Definition~\ref{def:7.6.1} and the second formula follows.
	If $a_t = 0$ and $n_t > 1$ then $\gcd(a_t, n_t) = n_t > 1$, so
	$\leg{a_t}{n_t} = 0$ by property (b), and the whole product is $0$; this
	is the first formula. Since the sequence terminates, exactly one of
	these two cases occurs.

	Finally, we prove the bound $t\le\log_2 n$. Strict decrease alone gives
	only a linear estimate; a step reduces $n_i$ by at least half except in one
	case, for which the following step compensates. If $e_i\ge1$, then
	$n_{i+1}=u_i\le a_i/2<n_i/2$. If $e_i=0$ and $a_i\le n_i/2$, then
	$n_{i+1}=u_i=a_i\le n_i/2$. In the remaining case $e_i=0$ and
	$a_i>n_i/2$, both $a_i$ and $n_i$ are odd. Since $n_i\ge3$,
	we have $n_{i+1}=a_i>n_i/2\ge3/2$, so $n_{i+1}>1$.
	Also $n_i<2a_i$, so $a_{i+1}=n_i\bmod u_i=n_i-a_i$ is positive and even.
	The step at index $i+1$ therefore exists, with $e_{i+1}\ge1$ and
	$n_{i+2}\le a_{i+1}/2=(n_i-a_i)/2<n_i/4$.
	Call a step of the third kind deficient. Each deficient step is
	followed by a step of the first kind, so no two deficient steps are adjacent.
	Pair each deficient index $i$ with $i+1$, and let every remaining index
	form a singleton block. These blocks partition $0,1,\ldots,t-1$.
	Each singleton block $\{i\}$ satisfies $n_{i+1}\le n_i/2$, and each
	paired block $\{i,i+1\}$ satisfies $n_{i+2}<n_i/4$.
	If there are $k$ paired blocks and $m$ singleton blocks, then $t=2k+m$.
	Multiplying the estimates along the blocks gives
	$n_0/n_t\ge4^k2^m=2^t$. Since $n_t\ge1$, we get $t\le\log_2 n$.
	This also covers $t=0$, when the partition is empty, including $n=1$.

\end{proof}
\begin{definition}[Kronecker symbol]\label{def:7.6.2} Let $a \in \mathbb{Z}$.
	For $n \neq 0$, write $n = u\,p_1 p_2 \cdots p_k$ with $u = \pm 1$ and
	with $p_1, \ldots, p_k$ primes (not necessarily distinct: each prime
	occurs as many times as its multiplicity in $|n|$), where $k = 0$ and the
	empty product is $1$ when $n = \pm 1$. Define
	\[
		\leg{a}{n} = \leg{a}{u}\prod_{i=1}^{k}\leg{a}{p_i},
	\]
	where $\leg{a}{1} = 1$; $\leg{a}{-1} = 1$ if $a \ge 0$ and $-1$ if $a
		< 0$; $\leg{a}{2} = 0$ if $a$ is even, $+1$ if $a \equiv \pm 1
		\pmod{8}$, and $-1$ if $a \equiv \pm 3 \pmod{8}$; and $\leg{a}{p}$ is
	the Legendre symbol for odd primes $p$. For $n = 0$, define
	$\leg{a}{0} = 1$ if $a = \pm 1$ and $\leg{a}{0} = 0$ otherwise.

\end{definition}
\begin{theorem}\label{thm:7.6.4} For every $a\in\mathbb{Z}$, the Kronecker symbol is well
	defined. For all nonzero $m,n\in\mathbb{Z}$, it is completely
	multiplicative in the lower argument,
	\[
		\leg{a}{mn}=\leg{a}{m}\leg{a}{n}.
	\]
	For every odd $n>0$, it coincides with the Jacobi symbol of Definition~\ref{def:7.6.1}.

\end{theorem}

\begin{proof} The value at the lower argument $0$ is prescribed
	explicitly. For a nonzero lower argument, unique prime factorization
	determines its sign and the multiset of prime factors of its absolute
	value, so the defining product is independent of every choice. The
	prime-factor multiset of $|mn|$ is the union, with multiplicity, of
	those for $|m|$ and $|n|$. If $a\ge0$, the rule on the unit
	$u\in\{\pm1\}$ is the trivial character; if $a<0$, it is the sign
	character $u\mapsto u$. In either case, it is multiplicative. This proves complete
	multiplicativity. If $n>0$ is odd, its unit is $+1$ and it has no factor
	$2$, so the product is exactly the Jacobi-symbol product of Definition~\ref{def:7.6.1}.

\end{proof}
\begin{theorem}\label{thm:7.6.5} Let $p$ be a prime with $p \equiv 3
		\pmod{4}$ and let $a$ be an integer with $\leg{a}{p} = 1$. Then $x
		\equiv a^{(p+1)/4} \pmod{p}$ satisfies $x^{2} \equiv a \pmod{p}$.

\end{theorem}

\begin{proof} By Euler's criterion, $x^2 \equiv a^{(p+1)/2} = a \cdot a^{(p-1)/2}
		\equiv a\leg{a}{p} = a \pmod{p}$.

\end{proof}
\begin{theorem}\label{thm:7.6.6} Let $p$ be a prime with $p \equiv 5
		\pmod{8}$ and let $a$ be an integer with $\leg{a}{p} = 1$. Then
	$a^{(p-1)/4} \equiv \pm 1 \pmod{p}$, and a square root of $a$ is $x \equiv
		a^{(p+3)/8}$ if $a^{(p-1)/4} \equiv 1$, and $x \equiv
		2a\,(4a)^{(p-5)/8}$ if $a^{(p-1)/4} \equiv -1$.

\end{theorem}

\begin{proof} Since $\big(a^{(p-1)/4}\big)^2 =
		a^{(p-1)/2} \equiv 1$, we have $a^{(p-1)/4} \equiv \pm 1$. If $a^{(p-1)/4}\equiv1$, then $x^2 \equiv
		a^{(p+3)/4} = a \cdot a^{(p-1)/4} \equiv a$. If $a^{(p-1)/4}\equiv-1$, then $x^2 \equiv
		4a^2(4a)^{(p-5)/4} = a \cdot 2^{(p-1)/2} \cdot a^{(p-1)/4} \equiv
		a\leg{2}{p}(-1) = a$, using $\leg{2}{p} = -1$ for
	$p \equiv 5 \pmod{8}$.

\end{proof}
\begin{theorem}[Tonelli--Shanks algorithm]\label{thm:7.6.7} Let $p$ be an odd prime,
	write $p - 1 = 2^{s}q$ with $q$ odd, let $z$ be a quadratic
	non-residue modulo $p$, and let $a$ satisfy $\leg{a}{p} = 1$. Define a
	finite sequence of quadruples $(M_j, c_j, t_j, R_j) \in
		\mathbb{Z}_{\ge0} \times (\mathbb{F}_p^{\times})^{3}$ by
	\[
		(M_0, c_0, t_0, R_0) = \big(s,\ z^{q},\ a^{q},\ a^{(q+1)/2}\big),
	\]
	and, for $j \ge 0$ with $t_j \neq 1$,
	\[
		i_j = \min\{\,i \ge 1 \mid t_j^{2^{i}} = 1\,\},\qquad
		b_j = c_j^{\,2^{M_j - i_j - 1}},
	\]
	\[
		(M_{j+1}, c_{j+1}, t_{j+1}, R_{j+1}) = \big(i_j,\ b_j^{2},\ t_j
		b_j^{2},\ R_j b_j\big),
	\]
	where the recursion stops at the first index $J$ with $t_J = 1$. Then
	$1 \le i_j < M_j$ for every update index $0\le j<J$, so $s = M_0 > M_1 > \cdots > M_J \ge
		0$ and $J \le s$; the sequence is therefore well defined and finite,
	and
	\[
		R_J = a^{(q+1)/2}\prod_{j=0}^{J-1} b_j \quad\text{satisfies}\quad
		R_J^{2} = a \ \text{ in } \mathbb{F}_p.
	\]

\end{theorem}

\begin{proof} Work in $\mathbb{F}_p^{\times}$, which is cyclic of
	order $p-1=2^s q$; its elements of $2$-power order form a cyclic
	$2$-group of order $2^{s}$, in which $-1$ is the unique element of
	order $2$. We maintain three invariants: (a) $c_j$ has order $2^{M_j}$;
	(b) $t_j^{2^{M_j-1}} = 1$; (c) $R_j^{2} = a\,t_j$. At $j = 0$,
	since $z$ is a non-residue, $z^{(p-1)/2} = -1$, so $c_0 = z^{q}$ has
	order exactly $2^{s} = 2^{M_0}$; $t_0^{2^{s-1}} = a^{(p-1)/2} = 1$ by
	Euler's criterion; and $R_0^{2} = a^{q+1} = a\,a^{q} = a\,t_0$.

	Assume that the three invariants hold at index $j$ and that $t_j \neq 1$. Note that $M_j
		\ge 2$, since $M_j = 1$ would make (b) read $t_j = 1$; hence by (b) the
	set $\{i \ge 1 \mid t_j^{2^{i}} = 1\}$ contains $M_j - 1 \ge 1$, so $i_j$
	exists and $i_j \le M_j - 1 < M_j$; and $i_j \ge 1$ because $t_j \neq
		1$. By minimality, $t_j^{2^{i_j-1}}$ is a square root of $1$ different
	from $1$, hence $t_j^{2^{i_j-1}} = -1$. The exponent $M_j - i_j - 1$ is
	nonnegative, so $b_j = c_j^{2^{M_j-i_j-1}}$ is defined and, by (a), has order
	$2^{i_j+1}$; consequently $b_j^{2^{i_j}} = c_j^{2^{M_j-1}} = -1$ and
	$b_j^{2}$ has order $2^{i_j}$. The invariants persist: (a) holds at $j+1$ because
	$c_{j+1} = b_j^{2}$ has order $2^{i_j} = 2^{M_{j+1}}$; (b) holds
	because $t_{j+1}^{2^{M_{j+1}-1}} = (t_jb_j^{2})^{2^{i_j-1}} =
		t_j^{2^{i_j-1}}b_j^{2^{i_j}} = (-1)(-1) = 1$; and (c) holds because
	$R_{j+1}^{2} = R_j^{2}b_j^{2} = a\,t_j\,b_j^{2} = a\,t_{j+1}$.

	As long as the recursion runs, $M_{j+1} = i_j$ with $1 \le i_j < M_j$,
	so $s = M_0 > M_1 > \cdots$ with all values in $\{1, \ldots, s\}$;
	moreover $M_j = 1$ forces $t_j = 1$ by (b). The recursion therefore
	stops at some $J \le s$ with $t_J = 1$; then (c) gives $R_J^{2} = a\,t_J = a$. Unwinding $R_{j+1} = R_jb_j$
	from $R_0 = a^{(q+1)/2}$ gives the displayed product formula.

\end{proof}
\begin{theorem}[Atkin's square-root algorithm]\label{thm:7.6.8} Let $p \equiv 5 \pmod{8}$ be prime and
	let $\leg{a}{p} = 1$. Put $b \equiv (2a)^{(p-5)/8} \pmod{p}$, $c \equiv 2a\,b^2 \pmod{p}$, and $x \equiv a\,b\,(c-1) \pmod{p}$.
	Then $x^2 \equiv a \pmod{p}$, and $c$ is a square root of $-1$ modulo $p$.

\end{theorem}

\begin{proof} Since $p \equiv 5 \pmod{8}$, the exponent
	$(p-5)/8$ is a nonnegative integer, and $2$ is a quadratic
	non-residue modulo $p$ (the second supplementary law, Theorem~\ref{thm:7.2.5}(f)), so Euler's
	criterion gives $2^{(p-1)/2} \equiv -1 \pmod{p}$. From $b \equiv
		(2a)^{(p-5)/8}$ we obtain $c \equiv 2a\,b^2 \equiv (2a)^{1 + (p-5)/4}
		= (2a)^{(p-1)/4} \pmod{p}$, hence $c^2 \equiv (2a)^{(p-1)/2} =
		2^{(p-1)/2}\,a^{(p-1)/2} \equiv (-1)\cdot 1 = -1 \pmod{p}$, using
	$\leg{a}{p} = 1$ and Euler's criterion. Now $2a\,b^2
		\equiv c$ gives $a\,b^2 \equiv c\cdot 2^{-1} \pmod{p}$, so $x^2 \equiv
		a^2 b^2 (c-1)^2 \equiv a\,(a b^2)\,(c-1)^2 \equiv \tfrac{a
			c}{2}\,(c^2 - 2c + 1) \equiv \tfrac{a c}{2}\,(-2c) \equiv -a\,c^2
		\equiv a \pmod{p}$.

\end{proof}
\subsection{Higher Power Residues, Roots of Unity, and the Power-Residue Symbol}

\begin{definition}[$n$-th power residue]\label{def:7.7.1}
	Let $p$ be an odd prime, suppose $\gcd(a, p) = 1$, and let $n \ge 1$. If $x^{n} \equiv a
		\pmod{p}$ has a solution, $a$ is an $n$-th power residue modulo $p$;
	otherwise, it is an $n$-th power non-residue. For $n = 1$, every $a$ is a
	first power residue (take $x = a$). For $n = 2, 3, 4$, these are
	the quadratic, cubic, and quartic (biquadratic) residues.

\end{definition}
\begin{definition}[power-residue set]\label{def:7.7.2} For an odd
	prime $p$ and an integer $k\ge1$, define
	\[
		\mathcal{R}^{(k)}=\{x^k \mid x\in\ZmodUnits{p}\}
		\subseteq\ZmodUnits{p}.
	\]
	To emphasize the modulus, we write $\mathcal{R}_p^{(k)}$.

\end{definition}
\begin{theorem}\label{thm:7.7.1} Let $p$ be an odd
	prime with primitive root $g$, let $n\ge1$ be an integer, suppose
	$\gcd(a, p) = 1$, and put $d = \gcd(n, p-1)$. Then:

	\begin{enumerate}[label={\alph*)}]
		\item (criterion) $a$ is an $n$-th power residue if and only if $d \mid
			      \ind_{p,g}(a)$;
		\item (Euler criterion) $a$ is an $n$-th power residue if and only if
			      $a^{\frac{p-1}{d}} \equiv 1 \pmod{p}$;
		\item (count) the number of $n$-th power residues in a reduced residue system is exactly
		      $\frac{p-1}{d}$, and when the congruence $x^{n} \equiv a \pmod{p}$ is solvable, there are
		      exactly $d$ roots $g^{\,y_0 + i\frac{p-1}{d}}$ ($i = 0, \ldots, d-1$),
		      where $y_0\in\mathbb{Z}$ is any solution of
		      $ny_0\equiv\ind_{p,g}(a)\pmod{p-1}$;
		\item (structure) the power-residue set $\mathcal{R}^{(n)}$ of
		      Definition~\ref{def:7.7.2} is the subgroup $\langle g^{d}\rangle$ of
		      $\ZmodUnits{p}$; in
		      particular, it depends only on $d = \gcd(n, p-1)$ and not on $n$ itself.
	\end{enumerate}

\end{theorem}

\begin{proof} Setting $x \equiv g^{y}$ and $a \equiv
		g^{c}$ with $c = \ind_{p,g}(a)$, we see that the congruence $x^{n}
		\equiv a$ is equivalent to the linear congruence $ny \equiv c
		\pmod{p-1}$, which is solvable if and only if $d \mid c$ (Theorem~\ref{thm:3.2.3}), with exactly
	$d$ solutions $y$ modulo $p-1$ when solvable. (When $d \mid c$, write
		$n = d n'$ and $p - 1 = d M$ with $\gcd(n', M) = 1$; then $ny \equiv
		c \pmod{p-1}$ is equivalent to $n' y \equiv c/d \pmod{M}$, which by
	$\gcd(n', M) = 1$ has a unique solution $y \equiv y_0 \pmod{M}$, and
	the $d$ residues $y_0, y_0 + M, \ldots, y_0 + (d-1)M$ are
	exactly its solutions modulo $p-1$.) As $a$ runs over a reduced
	residue system, $c$ runs over the integers in $\{0, \ldots, p-2\}$, of which
	$\frac{p-1}{d}$ are divisible by $d$. Next, $a^{\frac{p-1}{d}}
		\equiv g^{c\frac{p-1}{d}} \equiv 1$ if and only if $(p-1) \mid c\tfrac{p-1}{d}$,
		that is, if and only if $d \mid c$. Finally, the $n$-th powers are
	$\{g^{ny} \mid y \in \mathbb{Z}\}$. Every $ny$ is divisible by $d$.
	Conversely, B{\'e}zout's identity gives $u,v\in\mathbb{Z}$ with
	$d=un+v(p-1)$, so every multiple $jd$ satisfies
	$jd\equiv n(ju)\pmod{p-1}$. Thus $\{ny\bmod(p-1) \mid y\in\mathbb{Z}\}$
	is exactly the set of multiples of $d$ modulo $p-1$, and the $n$-th
	power residues form the subgroup $\langle g^{d}\rangle$, which depends
	only on $d$.

\end{proof}
\begin{definition}[root of unity]\label{def:7.7.3}
	For $k,m\in\mathbb{Z}^+$, a solution of $x^k\equiv1\pmod{m}$ is a $k$-th root of unity modulo
	$m$.

\end{definition}
\begin{definition}[primitive root of unity]\label{def:7.7.4}
	For $k,m\in\mathbb{Z}^+$, a $k$-th root of unity modulo $m$ is called primitive if its multiplicative order is
	$k$.

\end{definition}
\begin{theorem}\label{thm:7.7.2} For
	an odd prime $p$ and $k \ge 2$, a primitive $k$-th root of unity
	exists modulo $p$ if and only if $k \mid p - 1$. In particular, a primitive cube
	root of unity exists modulo $p$ if and only if $3 \mid p - 1$.

\end{theorem}

\begin{proof} A primitive $k$-th root of unity is exactly
	an element of order $k$ in $\ZmodUnits{p}$. Since
	$\ZmodUnits{p}$ is cyclic of order $p-1$ (Theorem~\ref{thm:6.2.6}), it
	contains
	an element of order $k$ if and only if $k \mid p-1$ (Theorem~\ref{thm:6.2.7}). Taking $k
		= 3$ gives the cube-root statement.


\end{proof}
\begin{definition}[power-residue symbol]\label{def:7.7.5} Let $p$
	be an odd prime, let $n\in\mathbb{Z}^+$ satisfy $n\mid p-1$, and let
	$a\in\mathbb{Z}$ satisfy $p\nmid a$. Define $\pleg{a}{p}{n}$
	to be the residue class $a^{(p-1)/n}\pmod{p}$.

\end{definition}
\begin{theorem}\label{thm:7.7.3} Let $p$ be an odd prime, let $n\in\mathbb{Z}^+$
	satisfy $n\mid p-1$, and let $a,b\in\mathbb{Z}$ satisfy $p\nmid ab$.
	Then:
	\begin{enumerate}[label={\alph*)}]
		\item the value $\pleg{a}{p}{n}$ is an $n$-th root of unity and lies
		      in the unique subgroup of order $n$ in $\ZmodUnits{p}$, namely
		      the subgroup of all $n$-th roots of unity;
		\item $\pleg{a}{p}{n}=1$ if and only if $a$ is an $n$-th power
		      residue modulo $p$;
		\item $\pleg{ab}{p}{n}=\pleg{a}{p}{n}\pleg{b}{p}{n}$;
		\item for $n=2$, the power-residue symbol is the residue class of the integer-valued Legendre symbol.
	\end{enumerate}

\end{theorem}

\begin{proof} Fermat's little theorem gives
	$\big(a^{(p-1)/n}\big)^n=a^{p-1}\equiv1\pmod{p}$. Since
	$\ZmodUnits{p}$ is cyclic of order $p-1$, Theorem~\ref{thm:6.6.5}(b) shows that it has a unique
	subgroup of order $n$, consisting of all roots of $x^n=1$, which proves
	(a). Part (b) is the Euler criterion for $n$-th powers (Theorem~\ref{thm:7.7.1}). Part (c) follows from
	$a^{(p-1)/n}b^{(p-1)/n}=(ab)^{(p-1)/n}$. For $n=2$, Euler's criterion
	(Theorem~\ref{thm:7.1.2}) identifies $a^{(p-1)/2}$ with $\leg{a}{p}$, proving
	(d).

\end{proof}
\begin{theorem}\label{thm:7.7.4} For an odd
	prime $p$, a primitive root $g$ modulo $p$, and an integer $a$ with $p
		\nmid a$, we have $\leg{a}{p} = (-1)^{\ind_{p,g}(a)}$.

\end{theorem}

\begin{proof} This is the case $n = 2$ of Theorem~\ref{thm:7.7.1}:
	$a$ is a quadratic residue if and only if $2 \mid \ind_{p,g}(a)$,
	that is, $(-1)^{\ind_{p,g}(a)} = 1$; the non-residue case
	gives $-1$.

\end{proof}
\begin{theorem}\label{thm:7.7.5} For an odd prime
	$p$, with $d = \gcd(3, p-1)$:

	\begin{enumerate}[label={\alph*)}]
		\item if $3 \nmid p - 1$, that is, if $p \equiv 2 \pmod{3}$ or $p = 3$
		      (in both cases $d = 1$), then $x \mapsto x^{3}$ is a bijection of
		      $\ZmodUnits{p}$, so every reduced residue is a cubic residue
		      with a unique cube root;
		\item if $p \equiv 1 \pmod{3}$ ($d = 3$), there are exactly
		      $\frac{p-1}{3}$ cubic residues; for every $a\in\mathbb{Z}$ with
		      $p\nmid a$, the criterion is $a^{(p-1)/3} \equiv 1\pmod{p}$,
		      and, when this criterion is satisfied, the three cube roots are $x_0, \omega x_0,
		      \omega^{2} x_0$, where $x_0$ is any one cube root of $a$ and $\omega$ is a primitive cube root of unity.
	\end{enumerate}

\end{theorem}

\begin{proof} Both cases follow from Theorem~\ref{thm:7.7.1} with $n = 3$:
	the count is $\frac{p-1}{d}$; for each $a$ with $p\nmid a$, the criterion
	is $a^{(p-1)/d} \equiv 1 \pmod{p}$, and there are $d$ roots when it is
	satisfied. For $d = 1$, the map is a bijection; for
	$d = 3$, the three roots form the coset $x_0 \cdot \{1, \omega,
		\omega^{2}\}$, where $\omega$ has order $3$ (such an element exists since $3 \mid
		p-1$) and is a root of $x^{2} + x + 1 \equiv 0 \pmod{p}$.

\end{proof}
\begin{theorem}\label{thm:7.7.6} For an odd prime $p$,
	we have $\mathcal{R}^{(4)}\subseteq\mathcal{R}^{(2)}$. More precisely:

	\begin{enumerate}[label={\alph*)}]
		\item if $p \equiv 3 \pmod{4}$, the quartic residues
		      coincide with the quadratic residues;
		\item if $p \equiv 1 \pmod{4}$, the quartic residues form a
		      subgroup of index $2$ in the quadratic residues; there are exactly
		      $\frac{p-1}{4}$ of them; for $p\nmid a$, the criterion is
		      $a^{(p-1)/4} \equiv 1 \pmod{p}$;
		\item $-1$ is a quartic residue modulo $p$ if and only if
		      $p\equiv1\pmod8$.
	\end{enumerate}

\end{theorem}

\begin{proof} If $a\equiv x^4=(x^2)^2\pmod p$, then $a$ is a
	quadratic residue, which proves the inclusion. Let $g$ be a primitive
	root modulo $p$. By Theorem~\ref{thm:7.7.1},
	\[
		\mathcal{R}^{(4)}=\langle g^{\gcd(4,p-1)}\rangle,
		\qquad
		\mathcal{R}^{(2)}=\langle g^2\rangle.
	\]
	If $p\equiv3\pmod4$, then $\gcd(4,p-1)=2$, so the two sets are
	equal. If $p\equiv1\pmod4$, then $\gcd(4,p-1)=4$, so
	$\langle g^4\rangle$ has index $2$ in $\langle g^2\rangle$; the count
	and criterion are again those of Theorem~\ref{thm:7.7.1}.

	For (c), if $p\equiv3\pmod4$, then by (a) the quartic residues are the
	quadratic residues, while $-1$ is not a quadratic residue by the first
	supplementary law (Theorem~\ref{thm:7.2.5}(e)). If $p\equiv1\pmod4$,
	the criterion in (b) gives
	\[
		-1\in\mathcal{R}^{(4)}
		\iff (-1)^{(p-1)/4}=1
		\iff \frac{p-1}{4}\text{ is even}
		\iff p\equiv1\pmod8.
	\]

\end{proof}
\begin{theorem}\label{thm:7.7.7} Let $p$ be an odd
	prime with primitive root $g$, let $k \ge 2$, and write $d = \gcd(k, p-1)$.

	\begin{enumerate}[label={\alph*)}]
		\item (roots of unity) $x^{k} \equiv 1 \pmod{p}$ has exactly $d$ solutions;
		\item (subgroup structure) the power-residue set $\mathcal{R}^{(k)}$ of
		      Definition~\ref{def:7.7.2} is a subgroup of $\ZmodUnits{p}$ of index
		      $d$ and order $\frac{p-1}{d}$, depending only on $d$;
		\item (unified Euler criterion) for every $a\in\mathbb{Z}$ with
		      $p\nmid a$, the integer $a$ is a $k$-th power residue if and only if
		      $a^{(p-1)/d} \equiv 1 \pmod{p}$;
		\item (intersection) $\mathcal{R}^{(j)} \cap \mathcal{R}^{(\ell)} =
			      \mathcal{R}^{(\lcm(j, \ell))}$ for all $j, \ell \ge 1$; in
		      particular, for every $a\in\mathbb{Z}$ with $p\nmid a$, the integer $a$ is a
		      sixth-power residue if and only if it is simultaneously a quadratic and a
		      cubic residue.
	\end{enumerate}

\end{theorem}

\begin{proof} The first item is the count in Theorem~\ref{thm:7.7.1}
	applied to $a = 1$, the second is its structure statement, and the third is
	its Euler criterion. For the intersection, that
	structure statement gives $\mathcal{R}^{(j)} = \langle g^{\gcd(j,
				p-1)}\rangle$, the set of $g^{c}$ with $\gcd(j, p-1) \mid c$, so
	$\mathcal{R}^{(j)} \cap \mathcal{R}^{(\ell)}$ is the set of $g^{c}$ with
	$\lcm(\gcd(j, p-1), \gcd(\ell, p-1)) \mid c$. For every prime $q$,
	Theorem~\ref{thm:4.6.1} gives the $q$-adic valuations
	$\max\{\min\{v_q(j),v_q(p-1)\},\min\{v_q(\ell),v_q(p-1)\}\}$ and
	$\min\{\max\{v_q(j),v_q(\ell)\},v_q(p-1)\}$ for $\lcm(\gcd(j, p-1), \gcd(\ell, p-1))$
	and $\gcd(\lcm(j, \ell), p-1)$, respectively. These agree, because for integers
	$x,y,z$ both $\max\{\min\{x,z\},\min\{y,z\}\}$ and $\min\{\max\{x,y\},z\}$ equal
	$\max\{x,y\}$ when $z\ge\max\{x,y\}$ and equal $z$ otherwise. Hence
	$\lcm(\gcd(j, p-1), \gcd(\ell, p-1)) =
		\gcd(\lcm(j, \ell), p-1)$ by Theorem~\ref{thm:4.5.3}, so this set is
	$\mathcal{R}^{(\lcm(j, \ell))}$. Taking $j = 2$ and $\ell = 3$
	gives $\gcd(6, p-1)$, so sixth-power residues are exactly the residues
	that are both quadratic and cubic.

\end{proof}
\subsection{Power Congruences Modulo Prime Powers and Composite Numbers}

\begin{theorem}\label{thm:7.8.1} Let $m\in\mathbb{Z}^+$ admit a primitive root
	$g$, and let $k\ge1$. By Theorem~\ref{thm:6.2.6}, the existence of such a primitive root is equivalent to
	$m\in\{1,2,4,p^e,2p^e\}$, where $p$ is an odd prime and $e\ge1$.
	Write $d = \gcd(k, \varphi(m))$ and
	$d' = \frac{\varphi(m)}{d}$. Then
	for $\gcd(a,m)=1$, the congruence $x^k\equiv a\pmod{m}$ is solvable
	if and only if $a^{d'} \equiv 1\pmod{m}$, or equivalently $d \mid \ind_{m,g}(a)$; and when
	the congruence is solvable, there
	are exactly $d$ roots $g^{\,y_0 + id'}$ ($0 \le i < d$, for a fixed
	particular solution $y_0$ of $ky\equiv\ind_{m,g}(a)\pmod{\varphi(m)}$).

\end{theorem}

\begin{proof} Since $g$ is primitive, Definition~\ref{def:6.2.1}
	gives $\delta_m(g)=\varphi(m)$, and Theorem~\ref{thm:6.2.1} says that the powers
	of $g$ represent all units modulo $m$. Put
	$c=\ind_{m,g}(a)$. Writing $x\equiv g^y$ reduces
	$x^k\equiv a\pmod{m}$ to the linear congruence
	$ky\equiv c\pmod{\varphi(m)}$. By Theorem~\ref{thm:3.2.3}, this is solvable if
	and only if $d\mid c$, and then its $d$ solutions are
	$y\equiv y_0+id'\pmod{\varphi(m)}$ for $0\le i<d$. Finally,
	$a^{d'}\equiv g^{cd'}\equiv1\pmod{m}$ if and only if
	$\varphi(m)\mid cd'$, equivalently $d\mid c$.

\end{proof}
\begin{theorem}\label{thm:7.8.2}
	Let
	$m=2^\alpha\prod_{i=1}^r p_i^{e_i}$, where $\alpha\ge0$, the $p_i$
	are distinct odd primes, and $e_i\in\mathbb{Z}^+$. Let $k\ge1$ and
	$\gcd(a,m)=1$. Then:

	\begin{enumerate}[label={\alph*)}]
		\item (localization principle) $x^{k}\equiv a\pmod{m}$ is solvable
		      if and only if it is solvable modulo $2^{\alpha}$ and modulo each
		      $p_i^{e_i}$, and the number of solutions modulo $m$ is the product
		      of the numbers of solutions modulo these prime-power factors;
		\item (explicit count) put
		      \[
			      N_i=p_i^{e_i-1}(p_i-1),\qquad d_i=\gcd(k,N_i).
		      \]
		      For $\alpha=2$, put
		      $\delta_{\mathrm{sgn}}=\gcd(k,2)$. For $\alpha\ge3$, write
		      uniquely
		      \[
			      a\equiv(-1)^b5^c\pmod{2^\alpha},\qquad
			      b\in\{0,1\},\quad c\in\mathbb{Z},\quad
			      0\le c<2^{\alpha-2},
		      \]
		      and put
		      \[
			      \delta_{\mathrm{sgn}}=\gcd(k,2),\qquad
			      \delta_{\mathrm{cyc}}=\gcd(k,2^{\alpha-2}).
		      \]
		      Define
		      \[
			      \varepsilon_2(k,a,2^\alpha)=
			      \begin{cases}
				      1, & \alpha\le1,\\
				      \delta_{\mathrm{sgn}},
				      & \alpha=2\ \text{and }a^{2/\delta_{\mathrm{sgn}}}
				        \equiv1\pmod4,\\
				      \delta_{\mathrm{sgn}}\delta_{\mathrm{cyc}},
				      & \alpha\ge3,\
				        \delta_{\mathrm{sgn}}\mid b\
				        \text{and }\delta_{\mathrm{cyc}}\mid c,\\
				      0, & \text{otherwise}.
			      \end{cases}
		      \]
		      If $a^{N_i/d_i}\equiv1\pmod{p_i^{e_i}}$ for every $i$, then
		      \[
			      \#\{x\bmod m \mid x^k\equiv a\pmod{m}\}
			      =\varepsilon_2(k,a,2^\alpha)\prod_{i=1}^r d_i;
		      \]
		      if any of the odd-prime-power conditions fails, the number is $0$.
	\end{enumerate}

\end{theorem}

\begin{proof} By the Chinese remainder theorem (Theorem~\ref{thm:3.4.1}),
	reduction is a ring isomorphism
	\[
		\mathbb{Z}_m\cong\mathbb{Z}_{2^\alpha}\times
		\prod_{i=1}^r\mathbb{Z}_{p_i^{e_i}}.
	\]
	It therefore identifies the solution set with the product of its local
	solution sets; this is the localization principle. For an odd prime
	power $p_i^{e_i}$, Theorem~\ref{thm:7.8.1}
	applies with group order $N_i$: the local congruence is solvable exactly
	when $a^{N_i/d_i}\equiv1\pmod{p_i^{e_i}}$, and then it has $d_i$
	solutions.

	For $\alpha\le1$ the power-of-$2$ unit group is trivial. For
	$\alpha=2$ it is cyclic of order $2$, so Theorem~\ref{thm:7.8.1} gives the
	second branch of $\varepsilon_2$. Finally, for $\alpha\ge3$, Theorem~\ref{thm:6.6.3} gives
	$\ZmodUnits{2^\alpha}\cong C_2\times C_{2^{\alpha-2}}$, and
	Theorem~\ref{thm:6.4.4} gives the unique expressions $(-1)^u5^v$. Thus $x^k\equiv a\pmod{2^\alpha}$
	is equivalent to
	\[
		ku\equiv b\pmod{2},\qquad
		kv\equiv c\pmod{2^{\alpha-2}}.
	\]
	By Theorem~\ref{thm:3.2.3}, these congruences are simultaneously
	solvable exactly when $\delta_{\mathrm{sgn}}\mid b$ and
	$\delta_{\mathrm{cyc}}\mid c$, and then they have, respectively,
	$\delta_{\mathrm{sgn}}$ and $\delta_{\mathrm{cyc}}$ solutions.
	Multiplying all local counts proves the formula.

\end{proof}
\begin{theorem}\label{thm:7.8.3}
	Let $n\ge2$, let $p$ be an odd prime with $p\nmid a$, and let
	$k\ge1$. If $p\nmid n$, then $x^{n}\equiv a\pmod{p^{k}}$ is solvable
	if and only if $x^{n}\equiv a\pmod{p}$ is solvable; moreover, every
	solution modulo $p$ lifts uniquely to a solution modulo each $p^k$.

\end{theorem}

\begin{proof} Assume that $p\nmid n$, and fix a solution $x_1$
	modulo $p$. For $1\le j<k$, suppose that
	$x_j^n\equiv a\pmod{p^j}$. We seek a lift $x_{j+1}=x_j+t p^j$. The
	binomial theorem gives
	\[
		x_{j+1}^n-a\equiv
		x_j^n-a+t p^j n x_j^{n-1}\pmod{p^{j+1}}.
	\]
	Writing $x_j^n-a=p^j s$, we obtain the lift condition
	\[
		s+t n x_j^{n-1}\equiv0\pmod{p}.
	\]
	Because $p\nmid a$, every solution satisfies $p\nmid x_j$; together with
	$p\nmid n$, this makes $n x_j^{n-1}$ invertible modulo $p$, so there
	is exactly one $t\pmod{p}$. Induction gives existence and uniqueness
	at every level, while reduction modulo $p$ proves the converse.

	The hypothesis $p\nmid n$ cannot be dropped: if $p\mid n$, solvability modulo
	$p$ need not imply solvability modulo $p^k$. Take $a=1+p$ and exponent $n=p$.
	Modulo $p$, the value $x=1$ is a solution of $x^p\equiv a$. Since $x^p\equiv x\pmod{p}$,
	every solution modulo $p$ satisfies $x\equiv1\pmod{p}$. Hence every lift has the form
	$x=1+pt$, and the binomial theorem gives $x^p\equiv1\pmod{p^2}$, never
	$1+p\pmod{p^2}$.

\end{proof}

\begin{theorem}\label{thm:7.8.4}\statementitems

	\begin{enumerate}[label={\alph*)}]
		\item (odd prime powers) For an odd prime $p$ and an integer $e\ge1$,
		      with $p\nmid a$, the congruence
		      $x^2\equiv a\pmod{p^e}$ has exactly $2$ solutions if $a$ is a
		      quadratic residue modulo $p$, and no solutions otherwise.
		\item (powers of $2$) For odd $a$, the congruence
		      $x^2\equiv a\pmod2$ has one solution;
		      the congruence $x^2\equiv a\pmod4$ has $2$ or $0$ solutions according as
		      $a\equiv1$ or $3\pmod4$; and, for $e\ge3$,
		      the congruence $x^2\equiv a\pmod{2^e}$ has $4$ or $0$ solutions according as
		      $a\equiv1\pmod8$ or $a\not\equiv1\pmod8$.
		\item (composite moduli) Let
		      \[
			      m=2^\alpha p_1^{\alpha_1}\cdots p_t^{\alpha_t},
		      \]
		      where $\alpha\ge0$, the $p_i$ are distinct odd primes, and
		      $\alpha_i\ge1$. If $\gcd(a,m)=1$, then
		      $x^2\equiv a\pmod m$ is solvable if and only if
		      $\leg{a}{p_i}=1$ for every $i$ and the power-of-$2$ condition is met:
		      no additional restriction when $\alpha\le1$,
		      the condition $a\equiv1\pmod4$ when $\alpha=2$, and
		      the condition $a\equiv1\pmod8$ when $\alpha\ge3$. When the congruence is solvable, the number
		      of roots is, respectively,
		      \[
			      2^t,\qquad 2^{t+1},\qquad 2^{t+2}
		      \]
		      according as $\alpha\le1$, $\alpha=2$, or $\alpha\ge3$.
		\item (unit quadratic residues) For an odd prime $p$ and
		      $e\ge1$, there are exactly $\varphi(p^e)/2$ unit quadratic
		      residues modulo $p^e$, and they are precisely the roots of
		      \[
			      a^{\varphi(p^e)/2}\equiv1\pmod{p^e}.
		      \]
	\end{enumerate}

\end{theorem}

\begin{proof} First let $p$ be odd. Modulo $p$, a nonzero quadratic
	residue has exactly two square roots by Theorem~\ref{thm:7.8.1},
	applied with exponent $2$. Since $p\nmid2$, Theorem~\ref{thm:7.8.3},
	applied with $n=2$, shows that each of these two roots lifts uniquely
	to a root modulo $p^e$. If $a$ is not a quadratic residue modulo $p$,
	reduction modulo $p$ rules out a root modulo $p^e$. This proves the
	odd-prime-power assertion.

	We next treat powers of $2$. Modulo $2$, there is only one odd class.
	Modulo $4$, every odd $x$ satisfies $x^2\equiv1\pmod4$, giving the
	two roots $x\equiv\pm1\pmod4$ when $a\equiv1\pmod4$ and no roots
	when $a\equiv3\pmod4$.

	Now let $e\ge3$. Necessity follows by reducing modulo $8$: if
	$x^2\equiv a\pmod{2^e}$, then $x$ is odd, say $x=2j+1$, and
	\[
		x^2=4j(j+1)+1\equiv1\pmod8
	\]
	because $j(j+1)$ is even. Thus $a\equiv1\pmod8$.

	Conversely, suppose that $a\equiv1\pmod8$ and take $x_3=1$.
	Suppose that $3\le j<e$ and write
	$x_j^2=a+2^j t$. Set $x_{j+1}=x_j+2^{j-1}s$. Since $2j-2\ge j+1$,
	\[
		x_{j+1}^2\equiv
		a+2^j(t+s x_j)\pmod{2^{j+1}}.
	\]
	The integer $x_j$ is odd, so choosing $s\equiv t\pmod2$ produces a
	root modulo $2^{j+1}$. Induction therefore gives a root modulo
	$2^e$.
	This correction preserves $x_j$ modulo $2^{j-1}$, but not necessarily
	modulo $2^j$: it proves existence at each level, not that every root
	lifts in its original residue class. For example, $3$ is a root of
	$x^2\equiv1\pmod8$, but both of its lifts $3,11\pmod{16}$ square to $9$.
	The correction instead selects $7\pmod{16}$, which squares to $1$.

	Fix one such root $x_0$. Since $x_0$ is a unit, $x$ is another root
	if and only if $x\equiv x_0u\pmod{2^e}$ for a unit $u$ satisfying
	$u^2\equiv1\pmod{2^e}$. For odd $u$,
	we have $\gcd(u-1,u+1)=2$; hence
	$2^e\mid(u-1)(u+1)$ if and only if one factor is divisible by
	$2^{e-1}$. There are exactly four possibilities:
	\[
		u\equiv\pm1,\ \pm(1+2^{e-1})\pmod{2^e}.
	\]
	They give the four distinct roots
	\[
		x\equiv\pm x_0,\ \pm(1+2^{e-1})x_0\pmod{2^e}.
	\]
	This proves the power-of-$2$ assertion.

	For the composite-modulus assertion, the localization principle of
	Theorem~\ref{thm:7.8.2} identifies the solution set modulo $m$ with
	the product of the solution sets modulo its prime-power factors.
	The preceding local criteria give the stated solvability conditions,
	and multiplying the local root counts gives $2^t$, $2^{t+1}$, or
	$2^{t+2}$ according as the power-of-$2$ component contributes
	$1$, $2$, or $4$ roots.

	Finally, let $p$ be odd, put $N=\varphi(p^e)$, and choose a generator
	$g$ of the cyclic group $\ZmodUnits{p^e}$ using
	Theorem~\ref{thm:6.6.3}. Write a unit as $a=g^j$ with
	$0\le j<N$. It is a square if and only if $2\mid j$. Moreover,
	\[
		a^{N/2}\equiv1\pmod{p^e}
		\iff N\mid\frac{jN}{2}
		\iff 2\mid j.
	\]
	Thus the roots of this congruence are exactly the $N/2$ even powers
	of $g$, which are precisely the unit quadratic residues.

\end{proof}
\begin{altproof}[Alternative proof of the power-of-$2$ part of Theorem~\ref{thm:7.8.4}]
	Modulo $2$, there is one unit. Modulo $4$, both units square to $1$,
	so the congruence $x^2\equiv a\pmod4$ has two roots when
	$a\equiv1\pmod4$ and none when $a\equiv3\pmod4$.

	Now let $e\ge3$. By Theorems~\ref{thm:6.6.3}
	and~\ref{thm:6.4.4}, every unit has a unique representation
	\[
		a\equiv(-1)^b5^c\pmod{2^e},
		\qquad
		b\in\{0,1\},\quad 0\le c<2^{e-2},
	\]
	and
	\[
		\ZmodUnits{2^e}\cong C_2\times C_{2^{e-2}}.
	\]
	The image of the squaring homomorphism consists exactly of the
	elements for which $b=0$ and $c$ is even. This is equivalent to
	$a\equiv1\pmod8$: indeed,
	\[
		5^c\equiv
		\begin{cases}
			1\pmod8,&c\text{ even},\\
			5\pmod8,&c\text{ odd},
		\end{cases}
	\]
	and multiplying by $(-1)^b$ leaves the residue $1\pmod8$ only when
	$b=0$ and $c$ is even.

	The kernel of squaring has two elements in each cyclic factor and
	therefore has four elements. Every nonempty fiber is a coset of this
	kernel (Theorem~\ref{thm:6.6.5}(a)), so $x^2\equiv a\pmod{2^e}$ has four roots when
	$a\equiv1\pmod8$ and none otherwise.

\end{altproof}
\subsection{Quadratic Algebras over \texorpdfstring{$\mathbb{F}_p$}{F-p} and Cipolla's Algorithm}

\begin{definition}[{quadratic algebra $(\mathbb{Z}/n\mathbb{Z})[\sqrt{d}]$}]\label{def:7.9.1}
	Fix integers $n\ge2$ and $d$, treat $\sqrt d$ as a formal symbol, and
	set
	\[
		R_n(d)=(\mathbb{Z}/n\mathbb{Z})[\sqrt d]
		=\{a+b\sqrt d\mid a,b\in\mathbb{Z}/n\mathbb{Z}\}.
	\]
	Addition is componentwise, and multiplication is defined by
	\[
		(a+b\sqrt d)(c+e\sqrt d)
		=(ac+dbe)+(ae+bc)\sqrt d,
	\]
	with all coefficients taken modulo $n$.

	For $z=a+b\sqrt d$, the conjugate and norm are
	\[
		\bar z=a-b\sqrt d,
		\qquad
		N(z)=z\bar z=a^2-db^2\in\mathbb{Z}/n\mathbb{Z}.
	\]
	The ring operations, conjugation, and norm are given by integer
	polynomial identities, so those identities persist modulo every $n$.

	When $n=p$ is an odd prime with $p\nmid d$, the algebra is the field
	$\mathbb{F}_p[\sqrt d]$ when $\leg{d}{p}=-1$, and it is isomorphic to
	$\mathbb{F}_p\times\mathbb{F}_p$ when $\leg{d}{p}=1$.

	For the rest of this subsection, $p$ is an odd prime,
	$d$ is an integer with $p\nmid d$, and
	$R=R_p(d)=\mathbb{F}_p[\sqrt d]$ is this quadratic algebra with
	$n=p$; elements are written $z=a+b\sqrt d$ with
	$a,b\in\mathbb{F}_p$.

\end{definition}

\begin{theorem}\label{thm:7.9.1} Assume $\leg{d}{p}=-1$. Then:
	\begin{enumerate}[label={\alph*)}]
		\item $R$ is a field with $p^{2}$ elements;
		\item $R^{\times}$ is cyclic of order $p^{2}-1$.
	\end{enumerate}

\end{theorem}

\begin{proof} (a) By definition, an element of $R$ is represented by a
	unique pair $(a,b)\in\mathbb{F}_p^2$ written as $a+b\sqrt{d}$; hence
	$|R|=p^2$. The given formulas make $R$ a commutative ring with
	identity $1 = 1 + 0\sqrt{d}$. Let $z = a + b\sqrt{d} \neq 0$ and put
	$N = z\bar z = a^2 - d b^2 \in \mathbb{F}_p$. If $N = 0$ and $b
		\neq 0$, then $d \equiv (a b^{-1})^2 \pmod{p}$ would make $d$ a residue,
	contradicting $\leg{d}{p} = -1$; when $b = 0$, one has $N =
		a^2 \neq 0$ since $a \neq 0$. Hence $N \in \mathbb{F}_p^{\times}$ and
	$z^{-1} = N^{-1}\bar z$, so every nonzero element is invertible and
	$R$ is a field of order $p^2$.

	(b) By Theorem~\ref{thm:6.10.1}, the multiplicative group of a finite field
	is cyclic; applying this to $F=R$ from part (a), we find that $R^{\times}$ is cyclic of
	order $|R|-1=p^{2}-1$.

\end{proof}
\begin{theorem}[Frobenius automorphism]\label{thm:7.9.2} Assume
	$\leg{d}{p}=-1$. Then the map $\phi:R\to R$, $\phi(z)=z^{p}$, is a ring
	automorphism of $R$ fixing $\mathbb{F}_p$ pointwise, and for every
	$z = a + b\sqrt{d} \in R$, we have $z^p = a - b\sqrt{d} = \bar z$; in particular,
	conjugation $z\mapsto\bar z$ is a ring automorphism of $R$.

\end{theorem}

\begin{proof} Since $R$ has characteristic $p$ and $p \mid
		\binom{p}{k}$ for $1 \le k \le p-1$ (Theorem~\ref{thm:3.5.2}), we have $(x + y)^p = x^p + y^p$, so $z^p = a^p + b^p(\sqrt{d})^p$. By Fermat's little theorem, $a^p = a$, $b^p =
		b$, and $(\sqrt{d})^p = \sqrt{d} \cdot d^{\frac{p-1}{2}} =
		\leg{d}{p}\sqrt{d} = -\sqrt{d}$ by Euler's criterion. Hence
	$z^p = a - b\sqrt{d} = \bar z$. The identity $(x+y)^p=x^p+y^p$ proves that
	$\phi$ is additive, and commutativity of $R$ gives multiplicativity, so $\phi$ is a ring
	homomorphism. It is injective because $R$ is a field by
	Theorem~\ref{thm:7.9.1}(a) and $\phi(z)=0$ implies $z=0$. Since $R$ is finite, $\phi$ is bijective.
	It fixes $\mathbb{F}_p$ pointwise by Fermat's little theorem
	(Theorem~\ref{thm:3.3.2}). Thus $\phi$, which coincides with conjugation, is a ring automorphism.

\end{proof}
\begin{theorem}\label{thm:7.9.3} Assume
	$\leg{d}{p}=-1$ and define, for $z = a + b\sqrt{d} \in R$,
	the norm $N(z) = z\, z^p = z\bar z = a^2 - d b^2$. Then:
	\begin{enumerate}[label={\alph*)}]
		\item $N : R^{\times} \to \mathbb{F}_p^{\times}$ is a surjective
		      group homomorphism;
		\item the norm-one elements $\{z \mid N(z) = 1\}$ form a cyclic
		      subgroup of order $p + 1$.
	\end{enumerate}

\end{theorem}

\begin{proof} (a) We first prove only the homomorphism
	assertion; surjectivity will follow after part (b). Conjugation
	$z\mapsto\bar z=z^p$ is a ring homomorphism (Theorem~\ref{thm:7.9.2}), so $N(zw) = (z\bar z)(w\bar w)
		= N(z)N(w)$, and $N(z) \in \mathbb{F}_p$ is nonzero for $z \neq 0$,
	so $N$ takes values in $\mathbb{F}_p^{\times}$. (b) By Theorem~\ref{thm:7.9.1}(b),
	$R^{\times}$ is cyclic of order $p^2-1$. Since $N(z)=z^{p+1}$, the norm-one elements
	are the solutions of $z^{p+1} = 1$, which by Theorem~\ref{thm:6.6.5}(b) form the unique subgroup of
	order $\gcd(p+1, p^2-1) = p+1$ in the cyclic group $R^{\times}$.
	Since this subgroup is the kernel of $N$, Theorem~\ref{thm:6.6.5}(a) shows that $\operatorname{im} N$ has order $\tfrac{p^2-1}{p+1} =
		p-1$, which proves surjectivity in (a).

\end{proof}
\begin{theorem}[Cipolla's algorithm]\label{thm:7.9.4} Let $p$ be an odd prime
	and let $a\in\mathbb{F}_p^{\times}$ be a quadratic residue. There exists
	$y\in\mathbb{F}_p$ with $\leg{y^2-a}{p}=-1$. For any such $y$, put
	$d=y^2-a$, let $R=\mathbb{F}_p[\sqrt{d}]$ be the ring of Definition~\ref{def:7.9.1}, and set
	\[
		z=y+\sqrt{d},\qquad x=z^{(p+1)/2}.
	\]
	Then $R$ is a field with $p^2$ elements and $x\in\mathbb{F}_p$ satisfies
	$x^2=a$.

\end{theorem}

\begin{proof} Put $\chi(t)=\leg{t}{p}$, with
	$\chi(0)=0$. For every $c\in\mathbb{F}_p$,
	\[
		\#\{w\in\mathbb{F}_p \mid w^2=c\}=1+\chi(c),
	\]
	as one checks in the three cases: if $c=0$, the only
	solution is $w=0$ and $\chi(0)=0$; if $c$ is a nonzero square, say
	$c=w_0^{2}$, then $w^{2}=c$ forces $(w-w_0)(w+w_0)=0$, hence
	$w=\pm w_0$, and these are distinct because $2w_0=0$ would give
	$w_0=0$, so there are $2=1+\chi(c)$ solutions; if $c$ is a
	non-residue, there is no solution and $1+\chi(c)=0$.
	Hence, if $M$ is the number of pairs $(y,w)$ satisfying
	$w^2=y^2-a$, then
	\[
		M=p+\sum_{y\in\mathbb{F}_p}\chi(y^2-a).
	\]
	On the other hand, $(y-w)(y+w)=a$, and every
	$u=y-w\in\mathbb{F}_p^{\times}$ determines the other factor $a/u$ and
	therefore determines $(y,w)$ uniquely, since $2$ is invertible;
	explicitly, $y=\tfrac12(u+a/u)$ and $w=\tfrac12(a/u-u)$, and $u\neq0$
	follows from $a\neq0$. Conversely, every
	$u\in\mathbb{F}_p^{\times}$ arises this way, so
	$(y,w)\leftrightarrow u$ is a bijection. Thus $M=p-1$ and
	\[
		\sum_{y\in\mathbb{F}_p}\chi(y^2-a)=-1.
	\]
	Because $a$ is a nonzero square, exactly two values of $y$ make
	$y^2-a=0$. If $N_+$ and $N_-$ count the remaining values for which
	the character equals $+1$ and $-1$, respectively, then
	$N_++N_-=p-2$ and $N_+-N_-=-1$. Thus
	$N_-=(p-1)/2>0$, proving that a suitable $y$ exists. In particular,
	a uniformly random $y\in\mathbb{F}_p$ is admissible with exact
	probability $(p-1)/(2p)$.

	For such a choice, $\leg{d}{p}=-1$, so Theorem~\ref{thm:7.9.1}(a) shows that $R$ is a
	field with $p^2$ elements. The norm formula of Theorem~\ref{thm:7.9.3} gives
	\[
		x^2=z^{p+1}=N(z)=y^2-d=a.
	\]
	In particular, $x\neq0$, and Euler's criterion gives
	$x^{p-1}=a^{(p-1)/2}=1$, so $x^p=x$. Write
	$x=u+v\sqrt{d}$. Theorem~\ref{thm:7.9.2} gives
	$x^p=u-v\sqrt{d}$; comparison with $x=u+v\sqrt{d}$ and the fact that
	$p$ is odd yield $v=0$. Hence $x\in\mathbb{F}_p$.

\end{proof}
\begin{theorem}\label{thm:7.9.5} Assume $\leg{d}{p}=1$ and choose
	$s\in\mathbb{F}_p^{\times}$ with $s^2=d$. Then the ring $R$ of Definition~\ref{def:7.9.1} is not a field: the map $\Phi : R \to \mathbb{F}_p \times \mathbb{F}_p$,
	$\Phi(a + b\sqrt{d}) = (a + bs,\ a - bs)$, is a ring isomorphism, so
	$R \cong \mathbb{F}_p \times \mathbb{F}_p$.

\end{theorem}

\begin{proof} The map $\Phi$ is additive with $\Phi(1) = (1, 1)$,
	and it is multiplicative because $(a + b\sqrt{d})(c + e\sqrt{d}) = (ac
		+ dbe) + (ae + bc)\sqrt{d}$ maps to $\big(ac + dbe + (ae + bc)s,\ ac +
		dbe - (ae + bc)s\big) = \big((a + bs)(c + es),\ (a - bs)(c -
		es)\big)$, using $s^{2} = d$. If $\Phi(a + b\sqrt{d}) = (0, 0)$, then
	$a + bs = a - bs = 0$, so $2bs = 0$; as $p$ is odd and $s \neq 0$, we get $b
		= 0$ and then $a = 0$, so $\Phi$ is injective, hence bijective since
	$|R| = p^{2} = |\mathbb{F}_p \times \mathbb{F}_p|$. In particular,
	$\sqrt{d} - s$ and $\sqrt{d} + s$ are zero divisors with product $d -
		s^{2} = 0$.

\end{proof}
\subsection{The Quadratic Gauss Sum}

\begin{definition}[quadratic Gauss sum]\label{def:7.10.1} For an odd
	prime $p$, put $\zeta_p=e^{2\pi i/p}$ and
	\[
		G_p=\sum_{t=1}^{p-1}\leg{t}{p}\zeta_p^{\,t}.
	\]
	The notation $G_p$ is reserved for this complex Gauss sum.

\end{definition}
\begin{theorem}\label{thm:7.10.1}
	Let $p$ be an odd prime, let $A$ be a commutative ring with identity,
	and let $\zeta\in A$ satisfy
	\[
		\zeta^p=1,\qquad
		1+\zeta+\cdots+\zeta^{p-1}=0.
	\]
	Put
	\[
		g=\sum_{t=1}^{p-1}\leg{t}{p}\zeta^t,
		\qquad p^*=(-1)^{(p-1)/2}p.
	\]
	Then $g^2=p^*\,1_A$, where the integer coefficients act through the
	canonical map $\mathbb{Z}\to A$.

\end{theorem}

\begin{proof} Expanding the square and writing the second
	nonzero residue as $s=tu$ gives
	\[
		g^2=\sum_{u=1}^{p-1}\leg{u}{p}
		\sum_{t=1}^{p-1}\zeta^{t(1+u)}.
	\]
	For $u\equiv-1\pmod{p}$, the inner sum is $(p-1)1_A$. Otherwise,
	multiplication by $1+u$ permutes the nonzero residue classes modulo
	$p$, so the inner sum is
	$\sum_{s=1}^{p-1}\zeta^s=-1_A$. Since
	$\sum_{u=1}^{p-1}\leg{u}{p}=0$ (Theorem~\ref{thm:7.2.2}),
	\[
		g^2=(p-1)\leg{-1}{p}\,1_A
		-\sum_{u\neq-1}\leg{u}{p}\,1_A
		=p\leg{-1}{p}\,1_A=p^*\,1_A.
	\]

\end{proof}
\begin{theorem}\label{thm:7.10.2} For every odd prime $p$,
	\[
		G_p^2=\leg{-1}{p}p=(-1)^{(p-1)/2}p,
		\qquad |G_p|=\sqrt{p}.
	\]

\end{theorem}

\begin{proof} Apply Theorem~\ref{thm:7.10.1} with
	$A=\mathbb{C}$ and $\zeta=\zeta_p$. The hypotheses hold for
	$\zeta_p=e^{2\pi i/p}$: $\zeta_p^{p}=1$, $\zeta_p\neq1$, and
	$1+\zeta_p+\cdots+\zeta_p^{p-1}=0$. This
	gives $G_p^2=p^*=p\leg{-1}{p}$. Complex conjugation gives
	$\overline{G_p}=\leg{-1}{p}G_p$ (replace $t$ by $-t$), and therefore
	$|G_p|^2=G_p\overline{G_p}=p$.

\end{proof}
\begin{theorem}\label{thm:7.10.3} For every odd prime $p$,
	\[
		G_p=\begin{cases}\sqrt{p},  & \text{if }p\equiv1\pmod{4}, \\[2pt]
		             i\sqrt{p}, & \text{if }p\equiv3\pmod{4}.\end{cases}
	\]

\end{theorem}

\begin{proof} We compute the trace and the determinant of
	the Fourier matrix $F=p^{-1/2}\big(\zeta_p^{\,jk}\big)_{0\le j,k\le
			p-1}$, following Schur; the few linear-algebraic facts used are proved
	as needed. Write $\zeta=\zeta_p$ and $M=(\zeta^{jk})_{0\le j,k\le p-1}$,
	so that $F=p^{-1/2}M$; the trace of $X=(x_{jk})_{0\le j,k\le p-1}$ is
	$\tr X=\sum_{j}x_{jj}$.

	First, $F^{2}$ is a permutation matrix, $F^{4}=I$, and $F$ is
	unitary. Indeed, $(M^{2})_{jk}=\sum_{l=0}^{p-1}\zeta^{l(j+k)}$ equals
	$p$ if $j+k\equiv0\pmod{p}$ and $0$ otherwise (the geometric sum), so
	$F^{2}=P$, where $P$ is the permutation matrix of the involution
	$j\mapsto-j$ on $\mathbb{Z}/p\mathbb{Z}$; since $P^{2}=I$, also
	$F^{4}=I$. The same computation with $j-k$ in place of $j+k$ gives
	$M\overline{M}^{\mathsf{T}}=pI$, that is,
	$F^{-1}=\overline{F}^{\mathsf{T}}$.

	Second, we determine the eigenvalues of $F$. For each $\omega$ with $\omega^{4}=1$ and
	each column vector $v\in\mathbb{C}^{p}$, put
	$\pi_\omega(v)=\tfrac14\big(v+\omega^{-1}Fv+\omega^{-2}F^{2}v+
		\omega^{-3}F^{3}v\big)$. Then $F\pi_\omega(v)=\omega\,\pi_\omega(v)$
	(using $F^{4}=I$ and $\omega^{4}=1$), and
	$\sum_{\omega}\pi_\omega(v)=v$, because $\sum_{\omega}\omega^{-j}$
	equals $4$ for $j=0$ and $0$ for $j=1,2,3$. Hence $\mathbb{C}^{p}$ has
	a basis of eigenvectors of $F$, and from $F^{4}=I$ every eigenvalue
	lies in $\{1,i,-1,-i\}$; write $m_1,m_i,m_{-1},m_{-i}$ for the
	multiplicities. In an eigenbasis, the matrix of $F$ is diagonal, so,
	using the facts that the trace is invariant under change of basis and that the
	determinant is multiplicative, we obtain
	\[
		\tr F=(m_1-m_{-1})+i(m_i-m_{-i}),\qquad
		\det F=(-1)^{m_{-1}}i^{\,m_i-m_{-i}}.
	\]
	Both facts are elementary: $\tr(XY)=\sum_{j}\sum_{l}
		x_{jl}y_{lj}=\tr(YX)$ gives
	$\tr(SXS^{-1})=\tr X$, and from the Leibniz
	formula $\det X=\sum_{\sigma}\sgn(\sigma)
		\prod_{j}x_{j,\sigma(j)}$ one obtains $\det(XY)=\det(X)\det(Y)$ by
	expanding $\prod_{j}\big(\sum_{k_j}x_{j,k_j}y_{k_j,\sigma(j)}\big)$:
	the inner sum over $\sigma$ is alternating in
	$(k_0,\ldots,k_{p-1})$, hence vanishes unless the $k_j$ are distinct,
	in which case $k_j=\tau(j)$ for a permutation $\tau$ and the sum
	equals $\sgn(\tau)\det Y$.

	Third, the trace is $G_p/\sqrt{p}$. Indeed,
	$\tr M=\sum_{j=0}^{p-1}\zeta^{j^{2}}$; here $j=0$
	contributes $1$ and each quadratic residue is $j^{2}$ for exactly two
	nonzero $j$, while each non-residue is $j^{2}$ for no such $j$, so
	\[
		\sum_{j=0}^{p-1}\zeta^{j^{2}}
		=1+2\!\!\sum_{t\text{ residue}}\!\!\zeta^{t}
		=1+\sum_{t=1}^{p-1}\Big(1+\leg{t}{p}\Big)\zeta^{t}=G_p,
	\]
	using $\sum_{t=1}^{p-1}\zeta^{t}=-1$ and the definition of $G_p$.

	Fourth, we use the Vandermonde determinant formula. With
	$V(x_0,\ldots,x_{n-1})=\det\big(x_k^{\,j}\big)_{0\le j,k\le n-1}$,
	subtracting $x_0$ times row $j-1$ from row $j$ for $j=n-1,\ldots,1$
	turns the first column into $(1,0,\ldots,0)^{\mathsf{T}}$ and the
	$(j,k)$-entry into $x_k^{\,j-1}(x_k-x_0)$; expanding along the first
	column and factoring $x_k-x_0$ out of column $k$ gives
	$V=\prod_{k=1}^{n-1}(x_k-x_0)\,V(x_1,\ldots,x_{n-1})$, whence
	$\det M=\prod_{0\le j<k\le p-1}(\zeta^{k}-\zeta^{j})$. Each factor is
	$\zeta^{k}-\zeta^{j}=e^{\pi i(j+k)/p}\cdot2i\sin\frac{\pi(k-j)}{p}$
	with $\sin\frac{\pi(k-j)}{p}>0$, and
	$\sum_{j<k}(j+k)=(p-1)\sum_{m=0}^{p-1}m=\tfrac{p(p-1)^{2}}{2}$,
	because each index $m$ occurs in exactly $p-1$ pairs; therefore
	\[
		\det M=(2i)^{p(p-1)/2}e^{\pi i(p-1)^{2}/2}
		\prod_{0\le j<k\le p-1}\sin\frac{\pi(k-j)}{p}
		=(2i)^{p(p-1)/2}\prod_{0\le j<k\le p-1}\sin\frac{\pi(k-j)}{p},
	\]
	since $(p-1)^{2}/2$ is even. As $F$ is unitary,
	$|\det F|^{2}=\det\big(F\overline{F}^{\mathsf{T}}\big)=1$;
	consequently, the positive factor of $p^{-p/2}\det M$ equals $1$ and
	\[
		\det F=i^{\,p(p-1)/2}=
		\begin{cases}(-1)^{(p-1)/4},     & p\equiv1\pmod{4}, \\
		             -i\,(-1)^{(p-3)/4}, & p\equiv3\pmod{4}.\end{cases}
	\]

	Finally, we determine the multiplicities. Besides the trace and determinant
	equations obtained above, one has $m_1+m_i+m_{-1}+m_{-i}=p$ and
	\[
		(m_1+m_{-1})-(m_i+m_{-i})=\tr F^{2}=\tr P=1,
	\]
	where the first equality uses the fact that the eigenvalues of $F^{2}$ are the
	squares of the eigenvalues of $F$, and the last equality uses the fact that $j=0$ is the only
	fixed point of $j\mapsto-j$. Suppose first that $p\equiv1\pmod{4}$. Then
	$G_p^{2}=p$ (Theorem~\ref{thm:7.10.2}), so $G_p=\pm\sqrt{p}$ is real;
	the trace equation gives $m_i=m_{-i}$ and $m_1-m_{-1}=\pm1$, and the
	two count equations give $m_1+m_{-1}=(p+1)/2$. Write $p=4t+1$: if
	$m_1-m_{-1}=-1$, then $m_{-1}=t+1\not\equiv t=(p-1)/4\pmod{2}$,
	contrary to $(-1)^{m_{-1}}=\det F=(-1)^{(p-1)/4}$. Hence $m_1-m_{-1}=1$
	and $G_p=\sqrt{p}$. Now suppose that $p\equiv3\pmod{4}$. Then
	$G_p^{2}=-p$, so $G_p=\pm i\sqrt{p}$; the trace equation gives
	$m_1=m_{-1}$ and $m_i-m_{-i}=\pm1$, and the count equations give
	$m_1=m_{-1}=(p+1)/4$. Write $p=4t+3$: the determinant equation
	$(-1)^{t+1}i^{\,m_i-m_{-i}}=-i\,(-1)^{t}$ reads
	$i^{\,m_i-m_{-i}}=i$, so $m_i-m_{-i}\equiv1\pmod{4}$; since $m_i-m_{-i}=\pm1$, it
	equals $1$, and $G_p=i\sqrt{p}$.

\end{proof}
\begin{altproof}[Alternative proof of Theorem~\ref{thm:7.5.1}]
	Let $p$ and $q$ be distinct odd primes. Put
	\[
		\Phi_p(T)=1+T+\cdots+T^{p-1}\in\mathbb{F}_q[T],\qquad
		R=\mathbb{F}_q[T]/(\Phi_p(T)),\qquad
		\zeta=T\bmod\Phi_p(T).
	\]
	Here the quotient notation means polynomial congruence modulo
	$\Phi_p$: polynomials $f$ and $h$ represent the same element if
	$\Phi_p\mid f-h$. Polynomial division over $\mathbb{F}_q$ gives a
	unique representative of degree less than $p-1$, including the zero
	polynomial. Add or multiply representatives and take the remainder
	modulo $\Phi_p$. These operations are well defined: if
	$f\equiv f'$ and $h\equiv h'$ modulo $\Phi_p$, then
	$f+h\equiv f'+h'$ and $fh\equiv f'h'$ modulo $\Phi_p$, by expansion.
	The ring laws therefore follow from those of polynomials.
	Distinct constants give distinct classes, so $\mathbb{F}_q$ embeds
	in $R$. In particular, $q\cdot1=0$, and no smaller positive integer
	multiple of $1$ is zero. This is what characteristic $q$ means for
	$R$. No irreducibility assumption on $\Phi_p$ is needed.

	Thus $R$ is a commutative ring of characteristic $q$,
	we have $1+\zeta+\cdots+\zeta^{p-1}=0$, and
	$\zeta^p=1$ because $T^p-1=(T-1)\Phi_p(T)$. Since $\deg\Phi_p=p-1\ge2$,
	we have $\zeta\neq1$, so $\zeta$ has order exactly $p$.

	Set
	\[
		g=\sum_{t=1}^{p-1}\leg{t}{p}\zeta^t,\qquad
		p^*=(-1)^{(p-1)/2}p,
	\]
	with the integer coefficients interpreted in $\mathbb{F}_q\subseteq R$.
	Theorem~\ref{thm:7.10.1} gives $g^2=p^*$. Since $q\nmid p$, the element $p^*$ is a unit of
	$\mathbb{F}_q$, and hence $g$ is a unit of $R$.

	The binomial theorem and $q\mid\binom{q}{k}$ for $0<k<q$ give
	$(x+y)^q=x^q+y^q$ in $R$. Also $(xy)^q=x^qy^q$, and Fermat's
	little theorem gives $a^q=a$ for $a\in\mathbb{F}_q$.
	Hence the Frobenius map $x\mapsto x^q$ is a ring homomorphism of $R$
	fixing $\mathbb{F}_q$. Substituting $t'=tq$ in the Gauss sum gives
	\[
		g^q=\sum_t\leg{t}{p}\zeta^{tq}
		=\sum_{t'}\leg{t'q^{-1}}{p}\zeta^{t'}
		=\leg{q}{p}g.
	\]
	On the other hand, Euler's criterion in $\mathbb{F}_q$ gives
	\[
		g^q=g(g^2)^{(q-1)/2}
		=g(p^*)^{(q-1)/2}
		=\leg{p^*}{q}g.
	\]
	Cancelling $g$ shows that $\leg{q}{p}=\leg{p^*}{q}$. These are signs, and
	$q$ is odd, so equality in $\mathbb{F}_q$ is equality in
	$\{\pm1\}$. Finally,
	\[
		\leg{p^*}{q}
		=\leg{-1}{q}^{(p-1)/2}\leg{p}{q}
		=(-1)^{\frac{p-1}{2}\frac{q-1}{2}}\leg{p}{q},
	\]
	which is the quadratic reciprocity law.

\end{altproof}
\subsection{Multiplicative Characters, Gauss Sums, and Jacobi Sums}\label{subsec:7.11}

\begin{definition}[multiplicative character]\label{def:7.11.1}
	Let $p$ be a prime. A multiplicative character modulo $p$ is a group
	homomorphism $\chi:\mathbb{F}_p^{\times}\to\mathbb{C}^{\times}$.
	The principal character $\varepsilon$ has value $1$ on every nonzero
	residue. We extend each character to $\mathbb{F}_p$ by
	\[
		\chi(0)=\begin{cases}1,&\chi=\varepsilon,\\0,&\chi\ne\varepsilon.\end{cases}
	\]
	Thus $\varepsilon$ is the constant-one function on $\mathbb{F}_p$.
	Character products and powers are defined on $\mathbb{F}_p^{\times}$
	and then extended by this rule. In particular, $\chi^0=\varepsilon$.
	If a product is principal, its extended value at $0$ is $1$, even
	when the pointwise product of the extended factors has value $0$.
	All identities involving products of character values at $0$ will
	therefore be checked separately.

\end{definition}

\begin{theorem}[the character group]\label{thm:7.11.1}
	For $a\in\mathbb{F}_p^{\times}$, every value $\chi(a)$ is a
	$(p-1)$-st root of unity. The characters form a cyclic group
	$\widehat{\mathbb{F}_p^{\times}}$ of order $p-1$ under pointwise
	multiplication on $\mathbb{F}_p^{\times}$.
	For odd $p$, its unique character of order $2$ is
	$\chi(a)=\leg{a}{p}$.

\end{theorem}

\begin{proof}
	For $p=2$, the unit group consists only of $1$, so its only character
	is principal and all the assertions follow. Now assume that $p$ is odd.
	Since $a^{p-1}=1$, we have $\chi(a)^{p-1}=1$.
	Fix a primitive root $g$ modulo $p$ (Theorem~\ref{thm:6.2.6}), and put
	$\omega=e^{2\pi i/(p-1)}$. Every nonzero residue is a unique
	$g^k$ with $0\le k<p-1$. A character is determined by its value
	at $g$, which can be any $(p-1)$-st root of unity. Thus the characters
	are exactly
	\[
		\chi_j(g^k)=\omega^{jk}\qquad(0\le j<p-1).
	\]
	These formulas are well defined because exponents are taken modulo
	$p-1$, and they give $\chi_j\chi_l=\chi_{j+l}$ on the unit group.
	Hence $\chi_1$ generates the character group. Its unique element of
	order $2$ sends $g^k$ to $(-1)^k$.
	A power $g^k$ is a square exactly when $k$ is even, so this character
	is the Legendre symbol. Its value at $0$ also agrees with that symbol.

\end{proof}

\begin{theorem}[orthogonality]\label{thm:7.11.2}
	For a character $\chi$ and a nonzero residue $a$, we have
	\[
		\sum_{x\in\mathbb{F}_p^{\times}}\chi(x)
		=\begin{cases}p-1,&\chi=\varepsilon,\\0,&\chi\ne\varepsilon,\end{cases}
		\qquad
		\sum_{\chi\in\widehat{\mathbb{F}_p^{\times}}}\chi(a)
		=\begin{cases}p-1,&a=1,\\0,&a\ne1.\end{cases}
	\]
	With the extensions in Definition~\ref{def:7.11.1},
	$\sum_{x\in\mathbb{F}_p}\chi(x)$ is $p$ for $\chi=\varepsilon$
	and $0$ otherwise.

\end{theorem}

\begin{proof}
	If $\chi\ne\varepsilon$, choose $b$ with $\chi(b)\ne1$.
	Multiplication by $b$ permutes the nonzero residues, so their
	character sum equals $\chi(b)$ times itself and must be $0$.
	The principal case is immediate. For the second sum, if $p=2$, then
	$a=1$ and $\varepsilon$ is the only character, so the sum is $1=p-1$.
	For odd $p$, let $g$, $\omega$, and the characters $\chi_j$ be as in the
	proof of Theorem~\ref{thm:7.11.1}, and write $a=g^k$ with
	$0\le k<p-1$. The second sum is then the geometric sum
	$\sum_{j=0}^{p-2}\omega^{jk}$.
	It equals $p-1$ when $k=0$ and $0$ otherwise.
	Adding the specified value at $0$ yields the final assertion.

\end{proof}

\begin{definition}[Gauss sum]\label{def:7.11.2}
	Put $\zeta_p=e^{2\pi i/p}$. For a character $\chi$ modulo $p$, define
	\[
		\mathcal{G}(\chi)=\sum_{t\in\mathbb{F}_p}\chi(t)\zeta_p^{\,t}.
	\]
	Here the exponent can be any integer representative of $t$.
	For odd $p$ and $\chi=\leg{\cdot}{p}$, this is the quadratic Gauss
	sum $G_p$ of Definition~\ref{def:7.10.1}.

\end{definition}

\begin{theorem}[Gauss sum identities]\label{thm:7.11.3}
	The principal Gauss sum is $\mathcal{G}(\varepsilon)=0$.
	For every nonprincipal character $\chi$, its complex conjugate
	$\bar\chi$ is its inverse on $\mathbb{F}_p^{\times}$, and
	\[
		|\mathcal{G}(\chi)|^2=p,
		\qquad
		\mathcal{G}(\chi)\mathcal{G}(\bar\chi)=\chi(-1)p.
	\]
	In particular, for odd $p$,
	$\mathcal{G}(\leg{\cdot}{p})^2=(-1)^{(p-1)/2}p$.

\end{theorem}

\begin{proof}
	The geometric sum $\sum_{t=0}^{p-1}\zeta_p^t$ is $0$, which gives
	the principal case. Character values on nonzero residues have
	absolute value $1$ by Theorem~\ref{thm:7.11.1}, so
	$\bar\chi=\chi^{-1}$ there. For nonprincipal $\chi$, the zero terms
	vanish. Set $s=tu$ in the square of the absolute value to obtain
	\[
		|\mathcal{G}(\chi)|^2
		=\sum_{u\in\mathbb{F}_p^{\times}}\chi(u)
		 \sum_{t\in\mathbb{F}_p^{\times}}\zeta_p^{\,t(u-1)}.
	\]
	The inner sum is $p-1$ for $u=1$ and $-1$ otherwise.
	Orthogonality therefore gives $|\mathcal{G}(\chi)|^2=p$.
	Replacing $t$ by $-t$ in the conjugate sum gives
	$\overline{\mathcal{G}(\chi)}=\bar\chi(-1)\mathcal{G}(\bar\chi)$.
	Since $\chi(-1)^2=1$, multiplication yields the stated product
	identity. The quadratic case follows from
	$\bar\chi=\chi$ and $\chi(-1)=(-1)^{(p-1)/2}$.

\end{proof}

\begin{definition}[Jacobi sum]\label{def:7.11.3}
	For characters $\chi,\lambda$ modulo $p$, put
	\[
		J(\chi,\lambda)
		=\sum_{x\in\mathbb{F}_p}\chi(x)\lambda(1-x).
	\]

\end{definition}

\begin{theorem}[Jacobi sum identities]\label{thm:7.11.4}
	With the zero-value convention in Definition~\ref{def:7.11.1},
	\[
		J(\varepsilon,\varepsilon)=p,\qquad
		J(\varepsilon,\psi)=J(\psi,\varepsilon)=0
		\quad(\psi\ne\varepsilon).
	\]
	For every nonprincipal $\chi$,
	\[
		J(\chi,\chi^{-1})=-\chi(-1).
	\]
	If $\chi$, $\lambda$, and $\chi\lambda$ are all nonprincipal, then
	\[
		J(\chi,\lambda)
		=\frac{\mathcal{G}(\chi)\mathcal{G}(\lambda)}
		       {\mathcal{G}(\chi\lambda)},
		\qquad |J(\chi,\lambda)|=\sqrt p.
	\]

\end{theorem}

\begin{proof}
	The first identities follow from the constant value of $\varepsilon$
	and the full-residue orthogonality relation.
	For $\chi\ne\varepsilon$, the terms at $x=0,1$ vanish, and
	\[
		J(\chi,\chi^{-1})
		=\sum_{x\ne0,1}\chi\!\left(\frac{x}{1-x}\right).
	\]
	The substitution $u=x/(1-x)$ is a bijection from
	$\mathbb{F}_p\setminus\{0,1\}$ to
	$\mathbb{F}_p\setminus\{0,-1\}$, with inverse $x=u/(1+u)$.
	Thus the last sum is $-\chi(-1)$ by orthogonality.

	Now suppose $\chi$, $\lambda$, and $\chi\lambda$ are nonprincipal.
	Group $\mathcal{G}(\chi)\mathcal{G}(\lambda)$ by $r=s+t$.
	For $r=0$, its coefficient is
	$\lambda(-1)\sum_{s\ne0}(\chi\lambda)(s)=0$.
	For $r\ne0$, set $s=rx$ and $t=r(1-x)$. Multiplicativity of the
	extended individual characters gives the coefficient
	$(\chi\lambda)(r)J(\chi,\lambda)$, including the zero terms.
	Consequently,
	\[
		\mathcal{G}(\chi)\mathcal{G}(\lambda)
		=J(\chi,\lambda)\mathcal{G}(\chi\lambda).
	\]
	The denominator has absolute value $\sqrt p$ by
	Theorem~\ref{thm:7.11.3}. Dividing proves the formula,
	and taking absolute values gives the last assertion.

\end{proof}

\begin{theorem}[a two-variable diagonal equation]\label{thm:7.11.5}
	Let $p$ be a prime and let $n\in\mathbb{Z}^+$ divide $p-1$.
	Choose a character $\chi$ of order $n$.
	The number $N$ of pairs $(x,y)\in\mathbb{F}_p^2$ with
	$x^n+y^n=1$ is
	\[
		N=p-\sum_{i=1}^{n-1}\chi^i(-1)
		  +\sum_{\substack{1\le i,j\le n-1\\n\nmid i+j}}
		      J(\chi^i,\chi^j).
	\]
	In particular,
	\[
		\left|N-p+\sum_{i=1}^{n-1}\chi^i(-1)\right|
		\le(n-1)(n-2)\sqrt p.
	\]
	When $n=1$, all the displayed sums are empty and $N=p$.

\end{theorem}

\begin{proof}
	Let $M_n(u)$ be the number of roots of $x^n=u$ in $\mathbb{F}_p$.
	We claim that
	\[
		M_n(u)=\sum_{j=0}^{n-1}\chi^j(u).
	\]
	For $u=0$, both sides equal $1$ by the chosen extensions.
	For $n=1$, both sides equal $1$ for every $u$. Now let $n>1$ and
	$u\ne0$. Then $p$ is odd; let $g$ be the primitive root fixed in the
	proof of Theorem~\ref{thm:7.11.1}, and write $u=g^k$ and
	$\chi(g)=e^{2\pi i c/n}$ with $c\in\mathbb{Z}$. Here
	$\gcd(c,n)=1$, because $\chi$ has order $n$ and is determined by
	its value at $g$.
	The geometric sum on the right equals
	$n$ if $n\mid k$ and $0$ otherwise. The equation
	$x^n=g^k$ is equivalent to $n\ell\equiv k\pmod{p-1}$ for
	$x=g^{\ell}$, and this congruence has exactly those numbers of
	solutions because $n\mid p-1$.
	It follows that
	\[
		N=\sum_{u\in\mathbb{F}_p}M_n(u)M_n(1-u)
		 =\sum_{i,j=0}^{n-1}J(\chi^i,\chi^j).
	\]
	The pair $(0,0)$ of indices contributes $p$. Every pair with exactly
	one zero index contributes $0$. For nonzero indices with
	$n\mid i+j$, we have $j=n-i$, and the contribution is
	$-\chi^i(-1)$. There are $n-1$ such pairs.
	The remaining $(n-1)(n-2)$ pairs have all three characters
	$\chi^i$, $\chi^j$, and $\chi^{i+j}$ nonprincipal.
	Each Jacobi sum then has absolute value $\sqrt p$ by
	Theorem~\ref{thm:7.11.4}. Summing these contributions
	proves the formula, and the triangle inequality proves the bound.

\end{proof}

\newpage

\section{Quadratic Representations of Integers}

\subsection{Sums of Two Squares and Uniqueness of the Representation}

\begin{theorem}[Thue's lemma]\label{thm:8.1.1} Let $p$ be a prime
	and suppose that $\gcd(a, p) = 1$. Then there exist integers $x, y$ with $0 < |x|,
		|y| < \sqrt{p}$ such that $x \equiv a y \pmod{p}$.

\end{theorem}

\begin{proof} Let $k = \lfloor \sqrt{p} \rfloor$. Since $p$ is
	prime, it is not a perfect square (if $p = c^{2}$ with $c \in \mathbb{Z}^+$,
	then $c \mid p$ and $1 < c < p$), so $k < \sqrt{p}$ and $(k+1)^2 > p$. Consider the $(k+1)^2$ pairs $(u, v)$
	with $0 \le u, v \le k$ and the map $(u, v) \mapsto u - a v \bmod p$.
	By the pigeonhole principle, two distinct pairs $(u_1, v_1) \neq (u_2, v_2)$ satisfy
	$u_1 - a v_1 \equiv u_2 - a v_2 \pmod{p}$. Put $x = u_1 - u_2$, $y =
		v_1 - v_2$; then $x \equiv a y \pmod{p}$, $|x|, |y| \le k < \sqrt{p}$,
	and $(x,y) \neq (0,0)$. If $y = 0$, then $x \equiv 0$ with $|x| < \sqrt{p}$
	forces $x = 0$, a contradiction; hence $y \neq 0$, and then $0 < |y| < \sqrt{p} \le p$ gives $p \nmid y$, so
	$p \nmid ay$ because $\gcd(a, p) = 1$; thus $x \equiv ay \not\equiv 0
		\pmod{p}$ and $x \neq 0$.

\end{proof}
\begin{theorem}[Fermat's two-square theorem]\label{thm:8.1.2} An
	odd prime $p$ is a sum of two squares if and only if $p \equiv 1 \pmod{4}$.

\end{theorem}

\begin{proof} ($\Rightarrow$) Modulo $4$, a square is $0$
	or $1$; if $p = x^2 + y^2$ is odd, exactly one of $x, y$ is odd, so
	$p \equiv 1 \pmod{4}$. ($\Leftarrow$) If $p \equiv 1 \pmod{4}$, then
	$-1$ is a quadratic residue modulo $p$ (Theorem~\ref{thm:7.2.4}), so we may choose $a$ with $a^2
		\equiv -1 \pmod{p}$. This congruence gives $p\nmid a$, so Thue's
	lemma applies: there are $x, y$ with $0 < |x|,
		|y| < \sqrt{p}$ and $x \equiv a y \pmod{p}$; then $x^2 \equiv -y^2
		\pmod{p}$, so $p \mid x^2 + y^2$, and $0 < x^2 + y^2 < 2p$ forces $x^2
		+ y^2 = p$.

\end{proof}
\begin{altproof}[Alternative proof of sufficiency in Theorem~\ref{thm:8.1.2}] ($\Leftarrow$) Assume $p \equiv 1
		\pmod{4}$ and let $S = \{(x, y, z) \in (\mathbb{Z}^+)^3 \mid x^2 + 4yz =
		p\}$, which is nonempty because $(1, 1, \tfrac{p-1}{4}) \in S$. This set is finite:
	from $x^2+4yz=p$ with $y,z\ge1$ we get $4y\le 4yz\le p-1$ and
	$4z\le 4yz\le p-1$, hence $1\le x<\sqrt{p}$ and $1\le y,z\le\tfrac{p-1}{4}$. Define
	$\iota$ on $S$
	by $\iota(x, y, z) = (x + 2z, z, y - x - z)$ if $x < y - z$, $(2y -
		x, y, x - y + z)$ if $y - z < x < 2y$, and $(x - 2y, x - y + z, y)$ if $x
		> 2y$. The boundary values $x = y - z$ and $x = 2y$ never occur (they
	would give $p = (y+z)^2$ or $p = 4y(y+z)$, contradicting primality).
	Each branch preserves the form and returns positive integers, so
	$\iota$ maps $S$ to $S$: the first gives
	$(x+2z)^2+4z(y-x-z)=x^2+4yz=p$ with $x+2z>0$, $z>0$, and $y-x-z>0$
	(as $x<y-z$); the middle gives $(2y-x)^2+4y(x-y+z)=x^2+4yz=p$ with
	$2y-x>0$ (as $x<2y$), $y>0$, and $x-y+z>0$ (as $x>y-z$); the last
	gives $(x-2y)^2+4(x-y+z)y=x^2+4yz=p$ with $x-2y>0$ (as $x>2y$),
	$x-y+z>0$ (as $x>2y\ge y+1>y-z$ since $z\ge1$), and $y>0$. The map is an involution; write $(X,Y,Z)$ for the image triple. The first branch ($x<y-z$) gives $(X,Y,Z)=(x+2z,z,y-x-z)$, where $X=x+2z>2z=2Y$, so the image lies in the region $x>2y$ and the third branch returns $(X-2Y,X-Y+Z,Y)=(x,y,z)$. Conversely, the third branch ($x>2y$) gives $(X,Y,Z)=(x-2y,x-y+z,y)$, where $Y-Z=x-2y+z=X+z>X$, so the image lies in the region $x<y-z$ and the first branch returns $(X+2Z,Z,Y-X-Z)=(x,y,z)$. The middle branch ($y-z<x<2y$) gives $(X,Y,Z)=(2y-x,y,x-y+z)$, where $X=2y-x<2y=2Y$ and $Y-Z=2y-x-z=X-z<X$, so the image again satisfies $Y-Z<X<2Y$ and a second application of the middle branch returns $(2Y-X,Y,X-Y+Z)=(x,y,z)$. A fixed point can occur only in the middle branch (the first and third interchange the disjoint regions $x<y-z$ and $x>2y$), where $2y-x=x$ forces $x=y$; then $p=x^2+4xz=x(x+4z)$ with $p$ prime and $x,z\ge1$ forces $x=y=1$ and $z=\tfrac{p-1}{4}$. Hence $\iota$ has the unique fixed point $(1,1,\tfrac{p-1}{4})$, and the remaining elements of $S$ split into orbits $\{s,\iota(s)\}$ of size exactly $2$, so the finite set $S$ has odd cardinality. The map $\tau(x, y, z) = (x, z, y)$ is also an involution of $S$; if it had no fixed point, its orbits would all have size $2$ and $|S|$ would be even. Hence $\tau$ has a
	fixed point, where $y = z$, giving $p = x^2 + (2y)^2$.

\end{altproof}
\begin{altproof}[Alternative proof of sufficiency in Theorem~\ref{thm:8.1.2}]
	($\Leftarrow$) Let $p\equiv1\pmod{4}$. Then $-1$ is a quadratic residue modulo $p$
	(Theorem~\ref{thm:7.2.4}): there is an integer $x_0$ with
	$x_0^{2}\equiv-1\pmod{p}$, and replacing $x_0$ by its representative in
	$\{-\tfrac{p-1}{2},\ldots,\tfrac{p-1}{2}\}$ (a complete residue system
	modulo $p$, as $p$ is odd) gives an integer $x$ with
	$|x|\le\tfrac{p-1}{2}$ and $x^{2}\equiv-1\pmod{p}$; writing $1+x^{2}=mp$ gives $0<m\le
		\tfrac{1}{p}\big(1+\tfrac{(p-1)^{2}}{4}\big)<p$. Thus the set of integers $m$ with
	$1\le m<p$ for which $mp$ is a sum of two squares is nonempty. Let $m$
	be its least element and write $mp=u^{2}+v^{2}$. If $m=1$, we are done,
	so assume $m>1$. Choose $u',v'$ with $u'\equiv u$, $v'\equiv v
		\pmod{m}$ and $|u'|,|v'|\le\tfrac{m}{2}$; then
	$u'^{2}+v'^{2}\equiv u^{2}+v^{2}=mp\equiv0\pmod{m}$, say
	$u'^{2}+v'^{2}=mr$ with $0\le r\le\tfrac{m}{2}<m$. If $r=0$, then $u'^{2}+v'^{2}=0$, so $u'=v'=0$ and therefore $m\mid u$ and
	$m\mid v$; then $m^{2}\mid u^{2}+v^{2}=mp$, so $m\mid p$, which is impossible for
	$1<m<p$. Hence
	$1\le r<m$. The Brahmagupta--Fibonacci identity, which follows by expanding both sides, gives
	\[
		(mp)(mr)=(u^{2}+v^{2})(u'^{2}+v'^{2})=(uu'+vv')^{2}+(uv'-vu')^{2},
	\]
	and modulo $m$ we have $uu'+vv'\equiv u^{2}+v^{2}\equiv0$ and
	$uv'-vu'\equiv uv-vu=0$, so both squares are squares of multiples of
	$m$. Dividing by $m^{2}$ exhibits $rp$ as a sum of two squares with
	$1\le r<m$, contradicting minimality. Hence $m=1$ and $p$ is a sum of
	two squares.

\end{altproof}
\begin{theorem}\label{thm:8.1.3} If a prime $p \equiv 3 \pmod{4}$ divides $x^2 + y^2$,
	then $p \mid x$ and $p \mid y$.

\end{theorem}

\begin{proof} If $p \nmid x$, then $x$ is invertible
	modulo $p$ and $(y x^{-1})^2 \equiv -1 \pmod{p}$, so $-1$ is a
	quadratic residue, which is impossible for $p \equiv 3 \pmod{4}$ (Theorem~\ref{thm:7.2.4}).
	Hence $p \mid x$. Substituting into $x^2+y^2\equiv0\pmod{p}$ gives
	$y^2\equiv0\pmod{p}$, so $p\mid y$.

\end{proof}
\begin{theorem}\label{thm:8.1.4} For a positive
	integer $n$, the number $2n$ is a sum of two squares if and only if $n$ is.

\end{theorem}

\begin{proof} If $2n = a^{2} + b^{2}$, then $a^{2} +
		b^{2}$ is even, so $a$ and $b$ have the same parity; hence
	$\tfrac{a+b}{2}, \tfrac{a-b}{2} \in \mathbb{Z}$ and $n =
		\big(\tfrac{a+b}{2}\big)^{2} + \big(\tfrac{a-b}{2}\big)^{2}$.
	Conversely, if $n = a^{2} + b^{2}$, then $2n = (a+b)^{2} + (a-b)^{2}$,
	so $n$ is a sum of two squares if and only if $2n$ is.

\end{proof}
\begin{theorem}\label{thm:8.1.5} A positive
	integer $n$ is a sum of two squares if and only if every prime $p
		\equiv 3 \pmod{4}$ occurs to an even power in $n$.

\end{theorem}

\begin{proof} The Brahmagupta--Fibonacci identity $(a^2 +
		b^2)(c^2 + d^2) = (ac - bd)^2 + (ad + bc)^2$, which follows by expanding both sides, shows that a product of sums
	of two squares is a sum of two squares. ($\Leftarrow$) Write $n=k^2m$, where $m$ is squarefree and every prime divisor
	of $m$ is either $2$ or congruent to $1\pmod{4}$. Since $2
		= 1^2 + 1^2$ and each such prime is a sum of two squares (Theorem~\ref{thm:8.1.2}), $m = u^2 + v^2$, so $n = (ku)^2 + (kv)^2$. ($\Rightarrow$) Suppose
	$n = x^2 + y^2$. Put $d = \gcd(x, y)$ and write $x = dx', y = dy'$ with $\gcd(x', y') = 1$. Let $p \equiv 3 \pmod{4}$ be prime. If $p \mid x'^2 + y'^2$, then Theorem~\ref{thm:8.1.3} forces $p \mid x', y'$, a contradiction. Thus $p \nmid x'^2 + y'^2$
	and $v_p(n) = 2v_p(d)$ is even.\qedhere

\end{proof}
\begin{theorem}\label{thm:8.1.6} A
	prime $p$ has the form $a^2 - ab + b^2$ with $a, b \in \mathbb{Z}$ if
	and only if $p = 3$ or $p \equiv 1 \pmod{3}$.

\end{theorem}

\begin{proof} ($\Rightarrow$) From $4(a^2 - ab + b^2) =
		(2a - b)^2 + 3b^2$, we obtain $p \equiv (2a - b)^2 \pmod{3}$, which is $0$ or
	$1$, so $p = 3$ or $p \equiv 1 \pmod{3}$. ($\Leftarrow$) For $p = 3$,
	we have $3 = 2^2 - 2 + 1$. If $p \equiv 1 \pmod{3}$, then
	$\ZmodUnits{p}$ has an element $g$ of order $3$ (Theorem~\ref{thm:7.7.2}), so $g^2 + g + 1 \equiv 0 \pmod{p}$. By Thue's lemma applied to
	$a = g$, there are $x, y$ with $0 < |x|, |y| < \sqrt{p}$ and $x \equiv
		gy \pmod{p}$; then $x^2 + xy + y^2 \equiv y^2(g^2 + g + 1) \equiv 0
		\pmod{p}$. Since $0 < x^2 + xy + y^2 < 3p$, the only multiples of $p$
	in this range are $p$ and $2p$. The value $2p$ is impossible: the
	form $x^2 + xy + y^2$ is even only when both $x$ and $y$ are even (in
	which case it is congruent to $0 \pmod{4}$) and is odd otherwise, so it is
	never congruent to $2 \pmod{4}$, whereas $2p \equiv 2 \pmod{4}$ because $p$
	is odd. (Equivalently, using $4(x^2 + xy + y^2) = (2x + y)^2 + 3y^2$:
	if $x^2 + xy + y^2 = 2p$, then reducing modulo $3$ gives $(2x + y)^2
		\equiv 2p \equiv 2 \pmod{3}$, which is impossible since a square is congruent to $0$
	or $1 \pmod{3}$ and $p \equiv 1 \pmod{3}$.) Hence $x^2 + xy + y^2 = p$,
	and $(a, b) = (x, -y)$ gives $p = a^2 - ab + b^2$.

\end{proof}
\begin{theorem}\label{thm:8.1.7} Let $p$ be a prime with $p = a^2 + b^2 = c^2 + d^2$
	for positive integers $a, b, c, d$. Then $\{a, b\} = \{c, d\}$.

\end{theorem}

\begin{proof} Since $p$ is prime, $\gcd(a, b) = \gcd(c, d) = 1$, and $0 < a, b, c, d < \sqrt{p}$. By the
	Brahmagupta--Fibonacci identity, $p^2 = (ac + bd)^2 + (ad - bc)^2 =
		(ac - bd)^2 + (ad + bc)^2$. Also $(ad - bc)(ad + bc) = a^2 d^2 - b^2
		c^2 = p(a^2 - c^2)$, so $p \mid (ad - bc)(ad + bc)$. Since $p$ is
	prime, it divides at least one of the two factors; and from $0 < a,
		b, c, d < \sqrt{p}$ we get $0 < ad, bc < p$, so $0 \le |ad - bc| < p$
	while $0 < ad + bc < 2p$. Hence $p \mid (ad - bc)$ forces $ad - bc =
		0$, that is, $ad = bc$, and $p \mid (ad + bc)$ forces $ad + bc = p$; so
	either $ad = bc$ or $ad + bc = p$. If $ad = bc$, then $\gcd(a, b) = 1$ and
	$a \mid bc$ give $a \mid c$, say $c = am$; then $d = bm$ and $m^2 p =
		p$, so $m = 1$ and $(c, d) = (a, b)$. If $ad + bc = p$, then $(ac -
		bd)^2 = p^2 - p^2 = 0$, so $ac = bd$; now $a \mid bd$ with $\gcd(a, b)=1$
	gives $a \mid d$, say $d = an$; then $c = bn$ and $n^2 p = p$, so $n
		= 1$ and $(c, d) = (b, a)$. In either case $\{c, d\} = \{a, b\}$.

\end{proof}
\begin{altproof}[Alternative proof of Theorem~\ref{thm:8.1.7}]
	If $p=2$, positivity forces $a=b=c=d=1$, so the result is immediate.
	Hence we may assume that $p$ is odd. From $p=a^{2}+b^{2}$ with $p$ prime, we
	get $p\nmid b$ (otherwise $p\mid
		a$ and $p^{2}\mid p$), so $b$ is invertible modulo $p$ and
	$(ab^{-1})^{2}\equiv-1\pmod{p}$; likewise $(cd^{-1})^{2}\equiv-1$. By
	Theorem~\ref{thm:3.6.2} the congruence $x^{2}\equiv-1\pmod{p}$ has at most two
	solutions, and $\pm ab^{-1}$ are two distinct ones (distinct because
	$p$ is odd and $p\nmid ab^{-1}$), so $cd^{-1}\equiv\pm ab^{-1}$, that
	is, $p\mid ad-bc$ or $p\mid ad+bc$. Since $0<a,b,c,d<\sqrt{p}$ we have
	$|ad-bc|<p$ and $0<ad+bc<2p$, so either $ad=bc$ or $ad+bc=p$. In the
	first case $\gcd(a,b)=1$ and $a\mid bc$ give $a\mid c$, say $c=am$,
	hence $d=bm$ and $p=m^{2}p$, so $m=1$ and $(c,d)=(a,b)$. In the second
	case the identity $p^{2}=(ac-bd)^{2}+(ad+bc)^{2}$ forces $ac=bd$, and
	the same divisibility argument gives $(c,d)=(b,a)$.

\end{altproof}
\begin{theorem}[Euler's four-square identity]\label{thm:8.1.8}
	For real numbers $x_1,\ldots,x_4,y_1,\ldots,y_4$,
	we have $(x_1^2 + x_2^2 + x_3^2 + x_4^2)(y_1^2 + y_2^2 + y_3^2 + y_4^2) =
		z_1^2 + z_2^2 + z_3^2 + z_4^2$, where
	\[
		\begin{aligned} z_1                  & = x_1 y_1 + x_2 y_2 + x_3 y_3 + x_4 y_4, & z_2
		                                     & = x_1 y_2 - x_2 y_1 + x_3 y_4 - x_4 y_3,                                            \\ z_3 & = x_1 y_3 - x_2 y_4
		                - x_3 y_1 + x_4 y_2, & z_4                                      & = x_1 y_4 + x_2 y_3 - x_3 y_2 - x_4 y_1.
		\end{aligned}
	\]

\end{theorem}

\begin{proof} Regard the two quadruples as quaternions
	$q = x_1 + x_2 i + x_3 j + x_4 k$ and $r = y_1 + y_2 i + y_3 j + y_4
		k$, where $i^2 = j^2 = k^2 = ijk = -1$. Every property of quaternions used
	below follows from a single embedding into complex matrices: writing $\mathrm{i}$
	for the complex unit and
	\[
		M(q) = \begin{pmatrix} x_1 + x_2\mathrm{i}  & x_3 + x_4\mathrm{i} \\
		                -x_3 + x_4\mathrm{i} & x_1 - x_2\mathrm{i}\end{pmatrix},
	\]
	one checks on basis elements that $M(1)$ is the identity matrix, that
	$M(i)^{2} = M(j)^{2} = M(k)^{2} = -M(1)$, and that
	$M(i)M(j) = M(k)$, $M(j)M(k) = M(i)$, $M(k)M(i) = M(j)$,
	$M(j)M(i) = -M(k)$, $M(k)M(j) = -M(i)$, $M(i)M(k) = -M(j)$. That is, $M$
	turns the full multiplication table of the basis $1,i,j,k$ into matrix
	identities, where each verification is a single $2 \times 2$ matrix product. Since $M$ is
	$\mathbb{R}$-linear and injective, it identifies the quaternions with a
	subring of the $2 \times 2$ complex matrices. Quaternion
	multiplication is therefore associative, as matrix multiplication is.
	Conjugation $\bar q = x_1 - x_2 i - x_3 j - x_4 k$ corresponds to the
	conjugate transpose, $M(\bar q) = M(q)^{*}$, so the rule
	$(AB)^{*} = B^{*}A^{*}$ for complex $2 \times 2$ matrices $A,B$
	gives $\overline{qr} = \bar r\,\bar q$. Finally $N(q) = q\bar q = \det
		M(q) = x_1^2 + x_2^2 + x_3^2 + x_4^2$ is a nonnegative real number, that is,
	a real multiple of the identity $1$, and such elements commute with every
	quaternion. Therefore
	\[
		N(q\bar r) = (q\bar r)\,\overline{(q\bar r)} = q\bar r\,r\bar q =
		q\,N(r)\,\bar q = N(r)\,q\bar q = N(q)\,N(r).
	\]
	On the other hand, multiplying out $q\bar r$ by the quaternion rules
	gives $q\bar r = z_1 - z_2 i - z_3 j - z_4 k$ with $z_1, \ldots, z_4$
	exactly as stated, so $N(q\bar r) = z_1^2 + z_2^2 + z_3^2 + z_4^2$.
	Comparing the two evaluations of $N(q\bar r)$ yields the identity; in
	particular, a product of sums of four squares is a sum of four squares.

\end{proof}
\begin{theorem}\label{thm:8.1.9} For every odd prime $p$ there
	exist integers $x, y$ and $m$ with $1 \le m < p$ such that $x^2 + y^2 + 1 = mp$.

\end{theorem}

\begin{proof} By Theorem~\ref{thm:6.10.3} there are integers $x, y$
	with $1 + x^{2} + y^{2} \equiv 0 \pmod{p}$. We may normalize any such
	solution so that $0 \le x, y \le \tfrac{p-1}{2}$: for every residue
	class $r$ modulo the odd prime $p$, either $r$ or $-r$ has a
	representative in this range, and changing the sign does not change the
	square modulo $p$. Thus
	\[
		0<x^2+y^2+1\le 2\Big(\frac{p-1}{2}\Big)^2+1=
		\frac{p^2-2p+3}{2}<p^2,
	\]
	so $m = (x^2 + y^2 + 1)/p$ satisfies $1 \le m < p$.

\end{proof}
\begin{theorem}[Lagrange's four-square theorem]\label{thm:8.1.10}
	Every nonnegative integer is a sum of four squares.

\end{theorem}

\begin{proof} By Theorem~\ref{thm:8.1.8} it suffices to represent
	an odd prime $p$ (as $0, 1, 2$ are sums of four squares). Let $m$ be
	the least positive integer with $mp = x_1^2 + x_2^2 + x_3^2 + x_4^2$;
	by Theorem~\ref{thm:8.1.9}, $1 \le m < p$. We claim that $m = 1$. Suppose $m > 1$.
	Choose $y_i \equiv x_i \pmod{m}$ with $-\tfrac{m}{2} < y_i \le \tfrac{m}{2}$; then $\sum y_i^2 = mr$ with $0 \le r \le m$. If $r = 0$, then all
	$y_i = 0$, so $m \mid x_i$ and $m^2 \mid mp$, that is, $m \mid p$,
	which is impossible. If $r = m$, then $\sum y_i^2=m^2$ while each
	$y_i^2\le m^2/4$; hence equality can hold only if every
	$y_i^2=m^2/4$, that is, all $|y_i|=\tfrac{m}{2}$. So $m$ is even, and since the
	representatives were chosen in the half-open interval
	$-\tfrac{m}{2}<y_i\le\tfrac{m}{2}$, this forces $y_i=\tfrac{m}{2}$ for every
	$i$. Thus $x_i\equiv\tfrac{m}{2}\pmod{m}$, say $x_i=\tfrac{m}{2}+mk_i$ with
	$k_i\in\mathbb{Z}$; squaring gives
	$x_i^2=\tfrac{m^2}{4}+m^2k_i+m^2k_i^2\equiv\tfrac{m^2}{4}\pmod{m^2}$, whence
	$mp=\sum x_i^2\equiv4\cdot\tfrac{m^2}{4}=m^2\equiv0\pmod{m^2}$,
	again giving $m \mid p$, which is impossible. Hence $1 \le r < m$. Applying Theorem~\ref{thm:8.1.8} to $(mp)(mr) = \sum z_i^2$, we verify that
	every $z_i \equiv 0 \pmod{m}$, using $y_i \equiv x_i \pmod{m}$: for the first
	component, $z_1=\sum_i x_iy_i\equiv\sum_i x_i^2=mp\equiv0\pmod{m}$;
	for the other three, replacing every $y_i$ by $x_i$ modulo $m$ makes
	the terms cancel in pairs, giving $z_2\equiv
		x_1x_2-x_2x_1+x_3x_4-x_4x_3=0$, $z_3\equiv
		x_1x_3-x_2x_4-x_3x_1+x_4x_2=0$, and $z_4\equiv
		x_1x_4+x_2x_3-x_3x_2-x_4x_1=0\pmod{m}$. Writing $z_i = m w_i$ gives $pr =
		\sum w_i^2$ with $r < m$, contradicting minimality. Therefore $m =
		1$.

\end{proof}
\subsection{Gaussian Integers}

\begin{definition}[Gaussian integers]\label{def:8.2.1}
	The ring of Gaussian integers is
	$\mathbb{Z}[i]=\{a+bi \mid a,b\in\mathbb{Z}\}$, where $i^2=-1$. The
	conjugate and norm of $\alpha=a+bi$ are
	\[
		\bar\alpha=a-bi,\qquad N(\alpha)=\alpha\bar\alpha=a^2+b^2.
	\]

	A unit of $\mathbb{Z}[i]$ is an element with a multiplicative inverse in
	$\mathbb{Z}[i]$; two elements are associates if they differ by a unit
	factor. A nonzero nonunit $\pi$ is irreducible if $\pi=\alpha\beta$
	forces $\alpha$ or $\beta$ to be a unit, and is prime (a Gaussian prime)
	if $\pi\mid\alpha\beta$ forces $\pi\mid\alpha$ or $\pi\mid\beta$.

\end{definition}
\begin{theorem}\label{thm:8.2.1}
	For all $\alpha,\beta\in\mathbb{Z}[i]$,
	we have $N(\alpha\beta)=N(\alpha)N(\beta)$.

\end{theorem}

\begin{proof} Conjugation is multiplicative, since direct expansion
	gives $\overline{\alpha\beta}=\bar\alpha\bar\beta$. Therefore
	\[
		N(\alpha\beta)=\alpha\beta\overline{\alpha\beta}
		=\alpha\beta\bar\alpha\bar\beta
		=(\alpha\bar\alpha)(\beta\bar\beta)=N(\alpha)N(\beta),
	\]
	since multiplication in $\mathbb{Z}[i]$ is commutative.

\end{proof}
\begin{theorem}\label{thm:8.2.2} Given $\alpha,\beta\in\mathbb{Z}[i]$ with
	$\beta\neq0$, there are $q,r\in\mathbb{Z}[i]$ such that
	$\alpha=\beta q+r$ and either $r=0$ or $N(r)<N(\beta)$. Consequently,
	the Euclidean algorithm, B{\'e}zout's identity, and
	unique factorization into Gaussian primes hold in $\mathbb{Z}[i]$, and every
	Gaussian irreducible is prime.
	Every Gaussian prime is also irreducible.
	The units are $\{\pm1,\pm i\}$.

\end{theorem}

\begin{proof} First, a Gaussian integer $u$ is a unit if and only if
	$N(u)=1$: an inverse in $\mathbb{Z}[i]$ forces $N(u)N(u^{-1})=1$, and conversely
	$N(u)=1$ gives $u\bar u=1$. The integer solutions of $N(u)=1$ are
	$u=\pm1,\pm i$.

	Write $\alpha/\beta=x+yi$ and choose
	$a,b\in\mathbb{Z}$ with $|x-a|,|y-b|\le\tfrac{1}{2}$. Set $q=a+bi$ and
	$r=\alpha-\beta q$. Then
	\[
		\frac{N(r)}{N(\beta)}=(x-a)^2+(y-b)^2\le\frac{1}{2}<1.
	\]
	Thus division with remainder is available. Repeating it on
	$(\alpha,\beta)$, $(\beta,r)$, and so on gives the Euclidean algorithm;
	since the norms of the successive remainders are nonnegative integers
	that strictly decrease, the algorithm terminates. Its last nonzero
	remainder $\delta$ divides both $\alpha$ and $\beta$, and back-substituting
	the division steps expresses $\delta=u\alpha+v\beta$ with
	$u,v\in\mathbb{Z}[i]$; hence every common divisor of $\alpha$ and $\beta$
	divides $\delta$, so $\delta$ is a greatest common divisor and
	B{\'e}zout's identity holds in $\mathbb{Z}[i]$.

	Every Gaussian prime is irreducible. Indeed, if a Gaussian prime
	$\pi$ satisfies $\pi=\alpha\beta$, then $\pi\mid\alpha$ or
	$\pi\mid\beta$. In the first case write $\alpha=\pi\gamma$ and
	cancel the nonzero complex number $\pi$ to get $1=\gamma\beta$,
	so $\beta$ is a unit. The other case is symmetric.

	To prove the converse, let $\pi$ be irreducible and suppose
	$\pi\mid\alpha\beta$. If $\pi\nmid\alpha$, then every common divisor of
	$\pi$ and $\alpha$ is a unit, so $\gcd(\pi, \alpha)=1$. By B{\'e}zout's identity there
	exist $u,v\in\mathbb{Z}[i]$ with $u\pi+v\alpha=1$; multiplying by
	$\beta$ gives $u\pi\beta+v\alpha\beta=\beta$, and both terms on the
	left are divisible by $\pi$, so $\pi\mid\beta$. Hence every
	irreducible is prime.

	Existence of factorization now follows by induction on the norm: if
	$\alpha$ is a nonzero nonunit and not irreducible, write
	$\alpha=\beta\gamma$ with both factors nonunits; then
	$1<N(\beta),N(\gamma)<N(\alpha)$ and the induction is complete. For
	uniqueness, if $\alpha=\pi_1\cdots\pi_r=\rho_1\cdots\rho_s$ are two
	factorizations into irreducibles, then $\pi_1\mid\rho_1\cdots\rho_s$,
	so, by the prime property of irreducibles just proved, $\pi_1$ divides some $\rho_j$; since both are
	irreducible, they are associates. Cancelling and iterating shows that the
	two factorizations agree up to order and multiplication by units.

\end{proof}
\begin{theorem}\label{thm:8.2.3}
	The rational primes factor in $\mathbb{Z}[i]$ as follows:
	\begin{enumerate}[label={\alph*)}]
		\item $2=-i(1+i)^2$, where $1+i$ is a Gaussian prime;
		\item if $p$ is a rational prime with $p\equiv1\pmod{4}$, then $p=\pi\bar\pi$ for nonassociate
		      Gaussian primes $\pi,\bar\pi$ of norm $p$;
		\item if $p$ is a rational prime with $p\equiv3\pmod{4}$, then $p$ is a Gaussian prime.
	\end{enumerate}

\end{theorem}

\begin{proof} The first identity is immediate, and $N(1+i)=2$ shows that
	$1+i$ cannot be a product of two nonunits; hence it is irreducible and
	therefore a Gaussian prime by Theorem~\ref{thm:8.2.2}. If
	$p\equiv1\pmod{4}$, Fermat's two-square theorem (Theorem~\ref{thm:8.1.2}) gives
	$p=a^2+b^2=(a+bi)(a-bi)$. Since a Gaussian integer of rational-prime
	norm cannot factor into two nonunits, $a\pm bi$ are Gaussian primes;
	they are nonassociate: since $p = a^2+b^2$ is an odd prime, we have
	$a,b\neq0$ and $a\neq\pm b$ (otherwise $p = a^2$ or $p = 2a^2$), while the
	associates of $a+bi$ are $\pm(a+bi)$ and $\pm i(a+bi) = \pm(-b+ai)$;
	comparison of real and imaginary parts shows that none of these four numbers equals
	$a-bi$.

	If $p\equiv3\pmod{4}$ and $p\mid N(a+bi)=a^2+b^2$, then either $p\mid b$
	or, if $p\nmid b$, the inverse $b^{-1}$ modulo $p$ exists and
	$(ab^{-1})^2\equiv-1\pmod{p}$, which is impossible by Theorem~\ref{thm:7.2.4}. Thus
	$p\mid b$, and then $p\mid a$; in other words, $p\mid N(\alpha)$
	implies $p\mid\alpha$ in $\mathbb{Z}[i]$. Now if $p\mid\alpha\beta$ in
	$\mathbb{Z}[i]$, write $\alpha\beta=p\gamma$; taking norms gives
	$N(\alpha)N(\beta)=p^2N(\gamma)$, so the rational prime $p$ divides
	$N(\alpha)$ or $N(\beta)$ and hence divides $\alpha$ or $\beta$ by the
	previous implication. Thus $p$ has the prime-divisor property in
	$\mathbb{Z}[i]$; and since $N(p)=p^2>1$, the element $p$ is neither $0$ nor a unit.
	Therefore $p$ is a prime element, that is, a Gaussian prime.

\end{proof}
\subsection{Jacobi's Two-Square Formula}

\begin{definition}[two-square counting function]\label{def:8.3.1}
	For $n\in\mathbb{Z}^+$, let
	\[
		r_2(n)=\#\{(x,y)\in\mathbb{Z}^2 \mid x^2+y^2=n\}.
	\]

\end{definition}
\begin{definition}[the nontrivial character modulo $4$]\label{def:8.3.2} For $d\in\mathbb{Z}$, put $\chi(d)=0$ for even $d$, $\chi(d)=1$ for
	$d\equiv1\pmod{4}$, and $\chi(d)=-1$ for $d\equiv3\pmod{4}$.

\end{definition}
\begin{theorem}[Jacobi's two-square formula]\label{thm:8.3.1} For
	$n\in\mathbb{Z}^+$,
	\[
		r_2(n)=4\sum_{d\mid n}\chi(d)=4\big(d_1(n)-d_3(n)\big),
	\]
	where $d_1(n)$ and $d_3(n)$ count the positive divisors of $n$
	congruent to $1$ and $3$ modulo $4$, respectively.

\end{theorem}

\begin{proof} Write
	\[
		n=2^a\prod_{i=1}^{r}p_i^{b_i}\prod_{j=1}^{s}q_j^{c_j},
		\qquad p_i\equiv1\pmod{4},
		\quad q_j\equiv3\pmod{4}.
	\]
	By Theorem~\ref{thm:8.1.5}, $r_2(n)=0$ if some $c_j$ is odd. To evaluate the
	right-hand side, write $f(n)=\sum_{d\mid n}\chi(d)$ and note first
	that $f$ is multiplicative: $\chi$ itself is multiplicative (a direct
	check from its definition modulo $4$), so Theorem~\ref{thm:5.1.3} applies. On
	prime powers, we have
	\[
		f(2^a)=1,\qquad f(p_i^{b_i})=b_i+1,\qquad
		f(q_j^{c_j})=\sum_{k=0}^{c_j}(-1)^k=\begin{cases}1, & c_j \text{ even}, \\
		             0, & c_j \text{ odd}.\end{cases}
	\]
	Hence if some $c_j$ is odd, the factor $f(q_j^{c_j})$ vanishes, so
	$f(n)=0$ and the right-hand side $4f(n)$ vanishes as well. Suppose
	therefore that all $c_j$ are even. In $\mathbb{Z}[i]$, Theorem~\ref{thm:8.2.3} gives
	$p_i=\pi_i\bar\pi_i$, while every $q_j$ remains prime and
	$2=-i(1+i)^2$. If $N(z)=n$, then $z\mid z\bar z=n$ in
	$\mathbb{Z}[i]$, so only the Gaussian primes $1+i$, $\pi_i$,
	$\bar\pi_i$, $q_j$ can occur in $z$; write
	$z=u(1+i)^{\alpha}\prod_i\pi_i^{k_i}\bar\pi_i^{\ell_i}\prod_j q_j^{h_j}$
	with a unit $u$. Conjugation swaps $\pi_i$ and $\bar\pi_i$ and sends
	$1+i$ to its associate $1-i=-i(1+i)$, so, up to a unit,
	$z\bar z=(1+i)^{2\alpha}\prod_i p_i^{k_i+\ell_i}\prod_j q_j^{2h_j}$.
	Comparing with
	$n=(-i)^{a}(1+i)^{2a}\prod_i p_i^{b_i}\prod_j q_j^{c_j}$ via unique
	factorization forces $\alpha=a$, $2h_j=c_j$ (which is possible because $c_j$ is
	even), and $k_i+\ell_i=b_i$; conversely every such choice gives
	$N(z)=n$, since $N(u)=1$, $N(1+i)=2$, $N(\pi_i)=N(\bar\pi_i)=p_i$,
	and $N(q_j^{c_j/2})=q_j^{c_j}$. Thus a Gaussian integer
	$z$ with $N(z)=n$ is determined by one of the four units and, for
	each $i$, a choice
	\[
		\pi_i^{k_i}\bar\pi_i^{b_i-k_i},
		\qquad 0\le k_i\le b_i;
	\]
	distinct choices give distinct $z$, since the $\pi_i,\bar\pi_i$ are
	nonassociate and factorization is unique. Hence
	\[
		r_2(n)=4\prod_{i=1}^{r}(b_i+1).
	\]

	By the multiplicativity shown above and the prime-power values,
	$f(n)=\prod_{i=1}^{r}(b_i+1)$ when all $c_j$ are even; note in
	particular that the power of
	$2$ contributes the factor $f(2^a)=1$ because $\chi(2)=0$, so the
	$2$-part does not change the divisor sum. Therefore $r_2(n)=4f(n)$,
	and the description of $\chi$ gives
	$f(n)=d_1(n)-d_3(n)$.

\end{proof}
\subsection{Pythagorean Triples}

\begin{theorem}\label{thm:8.4.1} Every solution of
	$x^2+y^2=z^2$ in positive integers with $\gcd(x,y,z)=1$ (a primitive
	solution) and with $x$ even has the
	form $x=2ab,\ y=a^2-b^2,\ z=a^2+b^2$ with $a>b\ge1$, $\gcd(a, b)=1$, and
	$a\not\equiv b\pmod{2}$.

\end{theorem}

\begin{proof} In a primitive triple, exactly one of $x,y$
	is even (if both were odd, then $x^2+y^2\equiv2\pmod{4}$, which is not a square; if both
	were even, the triple would not be primitive) and $z$ is odd. Here $x$ is even and $y$ is odd. Since $x^2=(z-y)(z+y)$ and $y,z$ are both odd, $z-y$ and $z+y$
	are even, so write $z-y=2s$, $z+y=2t$ with $s,t\in\mathbb{Z}^+$ (as
	$z>y$); then $x^2=4st$, and as $x$ is even, $st=(x/2)^2$. Any common
	divisor of $s,t$ divides $t-s=y$ and $t+s=z$, hence divides
	$\gcd(y, z)$; and $\gcd(y, z)=1$, since a prime $q$ dividing both $y$ and
	$z$ would divide $z^2-y^2=x^2$ and hence $x$, contradicting
	$\gcd(x,y,z)=1$. Therefore $\gcd(s, t)=1$. Two coprime positive
	integers whose product is a square are both squares (Theorem~\ref{thm:2.3.6}), so
	$s=b^2$ and $t=a^2$ for positive integers $a,b$; then $a>b$ (as
	$t>s$) and $\gcd(a, b)=1$ (as $\gcd(s, t)=1$). Hence $z=s+t=a^2+b^2$,
	$y=t-s=a^2-b^2$, and $x^2=4st=(2ab)^2$ gives $x=2ab$. Finally $a,b$
	have opposite parity, for otherwise $y=a^2-b^2$ and $z=a^2+b^2$
	would both be even, contradicting primitivity.

\end{proof}
\subsection{Fermat's Method of Infinite Descent}

\begin{theorem}[Fermat's Last Theorem for exponent $4$]\label{thm:8.5.1} The equation
	$x^{4}+y^{4}=z^{2}$ has no solution in positive integers; a fortiori
	$x^{4}+y^{4}=z^{4}$ is unsolvable (Fermat's Last Theorem for $n=4$).

\end{theorem}

\begin{proof} Among all positive solutions, choose one
	with $z$ minimal. If a prime $q$ divided both $x$ and $y$, then $q^{4}\mid
		z^{2}$, so $q^{2}\mid z$ and $(x/q,\,y/q,\,z/q^{2})$ would be a smaller
	solution; hence $\gcd(x, y)=1$ and $(x^{2},y^{2},z)$ is a primitive
	Pythagorean triple. In the primitive triple $(x^{2},y^{2},z)$, exactly one of $x^{2},y^{2}$ is even; relabel the two legs so that $y$ is the even one (then $x$ is odd). By Theorem~\ref{thm:8.4.1} there are coprime $a>b\ge1$ of opposite parity with $x^{2}=a^{2}-b^{2}$, $y^{2}=2ab$, $z=a^{2}+b^{2}$. From $x^{2}+b^{2}=a^{2}$ with $\gcd(a, b)=1$ we get a second
	primitive triple; here $x$ is odd, so $b$ is even and $a$ is odd, and
	Theorem~\ref{thm:8.4.1} again gives coprime $c>d\ge1$ with $x=c^{2}-d^{2}$,
	$b=2cd$, $a=c^{2}+d^{2}$. Now $y^{2}=2ab=4acd$, so $(y/2)^{2}=acd$.
	The factors $a,c,d$ are pairwise coprime: $\gcd(c, d)=1$, and any
	prime dividing both $a=c^{2}+d^{2}$ and $c$ would divide $d^{2}$,
	hence $d$, contradicting $\gcd(c, d)=1$; thus $\gcd(a, c)=1$, and
	likewise $\gcd(a, d)=1$. Their product $acd=(y/2)^{2}$ is a square,
	so, since the three are pairwise coprime, applying Theorem~\ref{thm:2.3.6} (first
	to $a$ and $cd$, then to $c$ and $d$) shows that each factor is a square:
	$a=e^{2}$, $c=f^{2}$, $d=g^{2}$. Then
	$e^{2}=a=c^{2}+d^{2}=f^{4}+g^{4}$, giving another positive solution with
	$e\le e^{2}=a\le a^{2}<a^{2}+b^{2}=z$, so $e<z$, contradicting
	minimality. So no solution exists; taking $z=w^{2}$ rules out
	$x^{4}+y^{4}=w^{4}$.

\end{proof}
\subsection{Binary Quadratic Forms and the Class Number}

\begin{definition}[binary quadratic forms and basic terminology]\label{def:8.6.1}
	A binary quadratic form is
	\[
		f(x,y)=ax^2+bxy+cy^2
	\]
	with $a,b,c\in\mathbb{Z}$. Its discriminant is
	$\disc(f):=b^2-4ac$, also denoted by $D$, and the form itself is
	denoted by $(a,b,c)$.

	The form is positive definite if $a>0$ and $D<0$. It is primitive if
	$\gcd(a,b,c)=1$.

\end{definition}
\begin{definition}[discriminants and the principal form]\label{def:8.6.2}
	Every discriminant $D=b^2-4ac$ satisfies
	$D\equiv0$ or $1\pmod{4}$. Conversely, for every integer $D$ satisfying
	one of these congruences, the principal form of discriminant $D$ is
	\[
		\begin{cases}
			x^2-\dfrac{D}{4}y^2,
				& D\equiv0\pmod{4},\\[4pt]
			x^2+xy+\dfrac{1-D}{4}y^2,
				& D\equiv1\pmod{4}.
		\end{cases}
	\]
	In either case, the displayed form has integer coefficients and
	discriminant $D$.

	A negative discriminant $D$ is fundamental if either
	$D\equiv1\pmod{4}$ and $D$ is squarefree, or $D=4m$ with
	$m\equiv2$ or $3\pmod{4}$ and $m$ squarefree.

\end{definition}
\begin{definition}[proper equivalence, reduction, and class number]\label{def:8.6.3}
	Two forms $f,g$ are properly equivalent, written $f\sim g$, if
	\[
		g(x,y)=f(\alpha x+\beta y,\,\gamma x+\delta y)
	\]
	for some integers $\alpha,\beta,\gamma,\delta$ with
	$\alpha\delta-\beta\gamma=1$.

	A positive-definite form $(a,b,c)$ is reduced if
	$|b|\le a\le c$ and $b\ge0$ whenever $|b|=a$ or $a=c$.

	For a discriminant $D<0$, the class number $h(D)$ is the number of
	proper-equivalence classes of primitive positive-definite forms of
	discriminant $D$.

\end{definition}
\begin{theorem}\label{thm:8.6.1} If $f\sim g$ via
	$(\alpha,\beta,\gamma,\delta)$ with
	$\alpha\delta-\beta\gamma=1$, then:
	\begin{enumerate}[label={\alph*)}]
		\item the map $T(x,y)=(\alpha x+\beta y,\gamma x+\delta y)$ is a
		      bijection of $\mathbb{Z}^2$ onto itself, which sends primitive
		      pairs (those with $\gcd(x, y)=1$) to primitive pairs, so $f$ and
		      $g$ represent exactly the same integers, with the same
		      multiplicities, and primitivity is likewise preserved;
		\item $\disc(g)=\disc(f)$;
		\item proper equivalence is an equivalence relation on forms of a
		      fixed discriminant.
	\end{enumerate}

\end{theorem}

\begin{proof} (a) Since $\alpha\delta-\beta\gamma=1$, the
	matrix $
		\begin{pmatrix}\alpha & \beta \\\gamma&\delta
		\end{pmatrix}$ has the integer inverse $
		\begin{pmatrix}\delta & -\beta \\-\gamma&\alpha
		\end{pmatrix}$, so $T$ is a bijection
	$\mathbb{Z}^2\to\mathbb{Z}^2$ with
	$T^{-1}(x',y')=(\delta x'-\beta y',-\gamma x'+\alpha y')$. If $\gcd(x,
		y)=1$ and a prime $\ell$ divides both $\alpha x+\beta y$ and $\gamma
		x+\delta y$, then $\ell$ divides
	$\delta(\alpha x+\beta y)-\beta(\gamma x+\delta
		y)=(\alpha\delta-\beta\gamma)x=x$ and $\alpha(\gamma x+\delta
		y)-\gamma(\alpha x+\beta y)=(\alpha\delta-\beta\gamma)y=y$, so
	$\ell\mid\gcd(x, y)=1$, a contradiction; hence $T$ preserves
	primitivity. Conversely, $T^{-1}$ is given by the integer matrix
	$\begin{pmatrix}\delta&-\beta\\-\gamma&\alpha\end{pmatrix}$, whose
	determinant is
	$\delta\alpha-(-\beta)(-\gamma)=\alpha\delta-\beta\gamma=1$; running
	the displayed computation with
	$(\alpha,\beta,\gamma,\delta)$ replaced by
	$(\delta,-\beta,-\gamma,\alpha)$ shows that $T^{-1}$ also carries
	primitive pairs to primitive pairs. Hence $T$ maps the set of
	primitive pairs bijectively onto itself, so primitivity is preserved
	in both directions. Since
	$g(x,y)=f(T(x,y))$ and $T$ is a primitivity-preserving bijection, $f$
	and $g$ represent the same integers with matching multiplicities.
	(b) Write $f(x,y)=(x,y)F(x,y)^{\mathsf{T}}$ with the symmetric matrix $F=
		\begin{pmatrix}a & b/2 \\b/2&c
		\end{pmatrix}$, so $\disc(f)=-4\det F$. If $M=
		\begin{pmatrix}\alpha & \beta \\\gamma&\delta
		\end{pmatrix}$, then $g(x,y)=f\big((x,y)M^{\mathsf{T}}\big)=(x,y)M^{\mathsf{T}} F
		M(x,y)^{\mathsf{T}}$, so $g$ has the symmetric matrix $G=M^{\mathsf{T}}FM$ and $\det
		G=(\det M)^2\det F=\det F$ (as $\det
		M=\alpha\delta-\beta\gamma=1$), giving
	$\disc(g)=-4\det G=-4\det F=\disc(f)$.
	(c) Reflexivity uses the identity matrix; symmetry uses the inverse
	matrix from (a), which also has determinant $1$; transitivity uses
	the fact that a product of two determinant-$1$ integer matrices again has
	determinant $1$ and integer entries.

\end{proof}
\begin{theorem}\label{thm:8.6.2} Every positive-definite form $f$ is properly equivalent to a
	reduced form of the same discriminant.

\end{theorem}

\begin{proof} Write $f=(a,b,c)$ and $D=b^2-4ac<0$. For
	$(x,y)\neq(0,0)$, completing the square gives
	$4a\,f(x,y)=(2ax+by)^2-Dy^2\ge|D|y^2$ (as $-D=|D|>0$); in particular
	$f(x,0)=ax^2\ge a$ for $x\neq0$, and for $y\neq0$,
	we have $f(x,y)\ge|D|y^2/(4a)\ge|D|/(4a)$. Thus for any bound $C$, all
	$(x,y)$ with $f(x,y)\le C$ have $|y|$ bounded (from the second
	inequality) and, for each such $y$, the inequality $|2ax+by|\le\sqrt{4aC}$ restricts
	$x$ to a finite set of values; so only finitely many lattice points
	satisfy $f(x,y)\le C$. The set of primitive values is nonempty (it contains $f(1,0)=a$). Since it has only finitely many elements below any bound, it has a minimum, say $m=f(x_0,y_0)$ with
	$\gcd(x_0, y_0)=1$. By B{\'e}zout's identity, choose $x_1,y_1$ with $x_0y_1-x_1y_0=1$;
	the substitution $(x,y)\mapsto(x_0x+x_1y,\,y_0x+y_1y)$ has
	determinant $x_0y_1-x_1y_0=1$, so it defines a proper equivalence to
	$g(x,y)=f(x_0x+x_1y,\,y_0x+y_1y)=(a',b',c')$ with $a'=f(x_0,y_0)=m$
	and $c'=f(x_1,y_1)$. Since $(x_1,y_1)$ is primitive (it is a column of a
	determinant-$1$ matrix, so any common divisor of $x_1$ and $y_1$ would
	divide the determinant $1$), we have $c'\ge m=a'$ by minimality of $m$. Now apply the
	proper equivalence $(x,y)\mapsto(x+ny,y)$ (determinant $1$ for any
	$n\in\mathbb{Z}$), which sends $(a',b',c')$ to $(a',\,b'+2a'n,\,c'')$
	for a suitable integer $c''=g(n,1)$; choose $n$ so that
	$-a'<b'+2a'n\le a'$ (which is possible since consecutive choices of $n$ shift
	$b'$ by $2a'$). Write the resulting form as $(a,b,c)$ with $|b|\le
		a$. Since $(n,1)$ is primitive, $c=g(n,1)\ge m=a$ by minimality of
	$m$ (note that $a=a'=m$ is unchanged by this substitution). Hence $|b|\le
		a\le c$; the normalization $-a'<b'+2a'n\le a'$ already gives $b\ge0$ whenever $|b|=a$ (then
	$b=a$). If the remaining boundary case $a=c$ with $b<0$ occurs, apply the proper equivalence
	$(x,y)\mapsto(-y,x)$ (determinant $1$), which sends $(a,b,c)=(a,b,a)$ to $(a,-b,a)$, replacing
	$b$ by $-b>0$ without changing $a$ or $c$; the result meets Definition~\ref{def:8.6.3}. This is a reduced
	form properly equivalent to $f$.

\end{proof}
\begin{theorem}\label{thm:8.6.3}
	For every $D<0$, each proper-equivalence class of positive-definite forms of discriminant $D$ contains exactly one reduced form. Consequently, only finitely many reduced forms have discriminant $D$, and $h(D)$ equals the number of primitive reduced forms of discriminant $D$.

\end{theorem}

\begin{proof} Existence follows from Theorem~\ref{thm:8.6.2}. Let $f=(a,b,c)$ be reduced. From
	\[
		4af(x,y)=(2ax+by)^2+|D|y^2
	\]
	and $|D|=4ac-b^2\ge3a^2$, we determine its least positive represented value. If $y=0$ and $x\neq0$, then $f(x,0)=ax^2\ge a$. If $|y|\ge2$, then
	\[
		f(x,y)\ge\frac{|D|y^2}{4a}\ge\frac{|D|}{a}\ge3a>a.
	\]
	If $|y|=1$, then
	\[
		f(x,\pm1)=ax^2\pm bx+c\ge ax^2-|b||x|+c\ge a|x|^2-a|x|+a\ge a.
	\]
	Thus the least positive value represented by $f$ is $a$, and it is attained at $\pm(1,0)$. We record exactly when it is attained elsewhere. By the case analysis above, a vector $(x,y)$ with $f(x,y)=a$ must have $|y|\le1$, since $|y|\ge2$ gives $f(x,y)\ge3a>a$. If $y=0$, then $ax^2=a$ forces $x=\pm1$. If $|y|=1$, the chain
	$f(x,\pm1)\ge ax^2-|b||x|+c\ge a|x|^2-a|x|+a$ shows that $f(x,\pm1)=a$ requires $a|x|^2-a|x|+a=a$, that is, $|x|\in\{0,1\}$, together with equality in $|b||x|=a|x|$ and in $c\ge a$ wherever these were used. Concretely:
	\begin{itemize}
		\item $x=0$, $|y|=1$: then $f(0,\pm1)=c$, so this is minimal exactly when $a=c$, and the extra minimal vectors are $\pm(0,1)$;
		\item $|x|=|y|=1$: then $f(x,y)=a\pm b+c$, which equals $a$ exactly when $c=\pm b$, and since $|b|\le a\le c$ this forces $|b|=a=c$, in which case the extra minimal vectors are $\pm(1,-1)$ for $b=a$ and $\pm(1,1)$ for $b=-a$.
	\end{itemize}
	So for $|b|\le a<c$ the only minimal vectors are $\pm(1,0)$; for $a=c$ (with $|b|<a$) they are $\pm(1,0),\pm(0,1)$; and for $|b|=a=c$ they are $\pm(1,0),\pm(0,1)$ and one further pair, so that all six vectors occur for the form $(a,a,a)$, that is, for $D=-3a^{2}$. Thus the reduced forms with extra minimal vectors are exactly those with $a=c$.
	For later use, note that $(x,y)\mapsto(-y,x)$ is a proper
	automorphism of $(a,0,a)$: it has determinant $1$, preserves
	$a(x^2+y^2)$, and interchanges the pairs
	$\{\pm(1,0)\}$ and $\{\pm(0,1)\}$ of minimal vectors. Likewise,
	$(x,y)\mapsto(-y,x+y)$ is a proper automorphism of $(a,a,a)$:
	direct substitution shows that it preserves $a(x^2+xy+y^2)$, and it
	cyclically permutes the three pairs
	$\{\pm(1,0)\}$, $\{\pm(0,1)\}$, and $\{\pm(1,-1)\}$.
	Thus, in these two exceptional cases, the automorphisms allow any pair
	of minimal vectors in their respective orbits to be moved to
	$\{\pm(1,0)\}$ in the uniqueness argument below.

	Now let $g=(a',b',c')$ be a reduced form properly equivalent to $f$. By Theorem~\ref{thm:8.6.1}(a), a proper equivalence is a bijection of $\mathbb{Z}^2$ carrying primitive pairs to primitive pairs and satisfying $g=f\circ T$, so $f$ and $g$ represent the same integers, with primitivity preserved; in particular, their least positive values agree and $a'=a$. If the equivalence sends $(1,0)$ to $\pm(1,0)$, composition with $\pm I$ makes its first column $(1,0)^{\mathsf{T}}$; its matrix is then $\begin{pmatrix}1&n\\0&1\end{pmatrix}$, and $b'=b+2an$. Since $|b|,|b'|\le a$, the only possibilities are $n=0$, or a boundary interchange between $b=-a$ and $b'=a$. The boundary rule excludes the negative choice, so $b'=b$ and then $c'=c$ from $b^2-4ac=D$.

	It remains to consider an equivalence whose first column is an
	additional minimal vector. Except for the last pair when
	$f=(a,a,a)$, composition with $-I$ makes this column $(0,1)^{\mathsf{T}}$, so its matrix
	has the form
	\[
		T=\begin{pmatrix}0&-1\\1&n\end{pmatrix}.
	\]
	Direct substitution gives
	\[
		g=(a,\,2an-b,\,an^2-bn+a).
	\]
	Since $g$ is reduced, $|2an-b|\le a$. If $0\le b<a$, this forces
	$n=0$, so $g=(a,-b,a)$; because $a'=c'=a$, the boundary rule for $g$
	forces $-b\ge0$, hence $b=0$ and $g=f$. If $b=a$, then
	$n\in\{0,1\}$; $n=0$ gives the forbidden form $(a,-a,a)$, while
	$n=1$ gives $g=f$. Finally, when $f=(a,a,a)$, the proper
	automorphism
	\[
		A=\begin{pmatrix}0&-1\\1&1\end{pmatrix}
	\]
	satisfies $f\circ A=f$ and permutes the three pairs of minimal
	vectors, reducing the remaining case to one already treated. Hence
	the reduced representative is unique.

	Moreover, $\gcd(a,b,c)=\gcd\bigl(f(1,0),f(0,1),f(1,1)\bigr)$, since
	$f(1,0)=a$, $f(0,1)=c$, and $f(1,1)=a+b+c$; hence the content of a form
	is the greatest common divisor of the values it takes at primitive
	pairs, and Theorem~\ref{thm:8.6.1}(a) shows that properly equivalent
	forms have the same content. In particular, primitivity is a class
	invariant.

	Finally, reducedness gives $|D|=4ac-b^2\ge3a^2$, so $1\le a\le\sqrt{|D|/3}$. For each of the finitely many possible values of $a$, the condition $|b|\le a$ leaves finitely many values of $b$, and $c=(b^2-D)/(4a)$ is then determined. Thus there are finitely many reduced forms. Since $h(D)$ counts only primitive classes (Definition~\ref{def:8.6.3}), uniqueness shows that $h(D)$ is exactly the number of primitive reduced forms of discriminant $D$.

\end{proof}
\begin{theorem}\label{thm:8.6.4} For an
	odd prime $p$, the equation $p = x^2 + 2y^2$ has an integer solution
	if and only if $\leg{-2}{p} = 1$, that is, $p \equiv 1, 3 \pmod{8}$.

\end{theorem}

\begin{proof} ($\Rightarrow$) If $p = x^2 + 2y^2$, then $p
		\nmid y$ (else $p \mid x$ and $p^2 \mid p$), so $(xy^{-1})^2 \equiv
		-2 \pmod{p}$ and $\leg{-2}{p} = 1$. ($\Leftarrow$)
	Choose $u$ with $u^2 \equiv -2 \pmod{p}$. Since $p$ is odd, this
	congruence gives $p\nmid u$, so Thue's lemma
	(Theorem~\ref{thm:8.1.1}) supplies nonzero $x,y$ with
	$0<|x|, |y| < \sqrt{p}$ and
	$x \equiv u y \pmod{p}$; then $x^2 \equiv -2y^2 \pmod{p}$, so $p \mid
		x^2 + 2y^2$ and $0 < x^2 + 2y^2 < 3p$, giving $x^2 + 2y^2 \in \{p,
		2p\}$. If it equals $2p$, then $x = 2x'$ is even and $p = y^2 +
		2x'^2$. In either case $p$ has the form $X^2 + 2Y^2$. Finally,
	parts (c), (e), and (f) of Theorem~\ref{thm:7.2.5} give
	$\leg{-2}{p}=\leg{-1}{p}\leg{2}{p}$; checking the four odd residue
	classes modulo $8$ shows that this product is $1$ exactly when
	$p\equiv1$ or $3\pmod{8}$.

\end{proof}
\begin{theorem}\label{thm:8.6.5} Let $D < 0$ have class
	number $h(D) = 1$. An odd
	prime $p \nmid D$ is represented by the principal form of
	discriminant $D$ (Definition~\ref{def:8.6.2}) if and only if $\leg{D}{p} = 1$. In
	particular, for $D=-7$ and every odd prime $p\neq7$, we have
	$p=x^2+xy+2y^2$ if and only if $\leg{-7}{p}=1$;
	for an odd prime $p$, we have $p = x^2 + y^2$ ($D = -4$) if and only if $p \equiv 1
		\pmod{4}$; and for an odd prime $p$, we have $p = x^2 + 2y^2$ ($D = -8$) if and
	only if $\leg{-2}{p} = 1$, which is Theorem~\ref{thm:8.6.4}.

\end{theorem}

\begin{proof} ($\Leftarrow$) Suppose $\leg{D}{p}=1$. Then $D$ is
	a quadratic residue modulo $p$, so by the Chinese remainder theorem
	(Theorem~\ref{thm:3.4.1}) we may choose $b$ with $b
		\equiv b_0 \pmod{p}$ (where $b_0^2 \equiv D \pmod{p}$) and $b \equiv D
		\pmod{2}$; then $b^2 \equiv D \pmod{p}$, and $b^2 \equiv D \pmod{4}$ (an
	odd square is $1 \pmod{4}$ and an even square is $0 \pmod{4}$, matching $D \equiv 1$ and $D \equiv 0 \pmod{4}$, respectively), so $b^2 \equiv D \pmod{4p}$. Set $c = \tfrac{b^2 - D}{4p}
		\in \mathbb{Z}$. The form $g(x, y) = p\,x^2 + b\,xy + c\,y^2$ has
	discriminant $b^2 - 4pc = D$, is primitive (since $p \nmid D$ gives
	$p \nmid b$) and positive definite, and represents $p$ via $g(1, 0) =
		p$. Because $h(D) = 1$, there
	is a single class of primitive positive-definite forms of
	discriminant $D$, so $g$ is properly equivalent to the principal
	form; since properly equivalent forms represent exactly the same
	integers, the principal form also represents $p$. Conversely, if the
	principal form (with middle coefficient $\beta \in \{0, 1\}$
	according as $D \equiv 0$ or $1 \pmod{4}$) represents $p = f(x,
		y)$, the identity $4\,f(x, y) = (2x + \beta y)^2 - D\,y^2$ gives $(2x
		+ \beta y)^2 \equiv D\,y^2 \pmod{p}$; here $p \nmid y$, for if
	$p\mid y$, then the representation $p=f(x,y)$ yields $p\mid x$ as well,
	so $p^2\mid(2x+\beta y)^2-Dy^2=4p$, which is impossible. Thus $y$ is invertible
	modulo $p$, so $D$ is a square modulo $p$ and $\leg{D}{p} = 1$.
	To verify the three stated class numbers, let
	$D\in\{-7,-4,-8\}$ and let $(a,b,c)$ be reduced. Then
	$3a^2\le|D|\le8$, so the positive integer $a$ must be $1$.
	Also $b\equiv D\pmod2$ and $|b|\le1$; when $|b|=1$,
	the boundary rule selects $b=1$. The value $c=(b^2-D)/4$ is then forced:
	\[
		\begin{array}{c|ccc}
		D & a & b & c \\
		\hline
		-7 & 1 & 1 & 2 \\
		-4 & 1 & 0 & 1 \\
		-8 & 1 & 0 & 2
		\end{array}
	\]
	Each listed form is primitive and reduced. By Theorems~\ref{thm:8.6.2}
	and~\ref{thm:8.6.3}, these are the unique reduced representatives,
	so all three class numbers are $1$. Their principal forms are
	$x^2+xy+2y^2$, $x^2+y^2$, and $x^2+2y^2$, respectively.
	The criterion just proved gives the stated equivalences: for $D=-4$,
	the symbol $\leg{-4}{p}=\leg{-1}{p}$ is $1$ exactly when $p\equiv1\pmod4$;
	for $D=-8$, we have $\leg{-8}{p}=\leg{-2}{p}$ because $4$ is a nonzero square
	modulo every odd prime.

\end{proof}
\begin{theorem}\label{thm:8.6.6}
	Let $D \equiv 0$ or $1 \pmod{4}$ and let $n \in \mathbb{Z}^+$. Then
	$n$ is properly represented by some integral binary quadratic form of
	discriminant $D$ (that is, $n = f(x, y)$ with $\gcd(x, y) = 1$ and
	$\disc(f) = D$) if and only if the congruence $b^2\equiv D\pmod{4n}$ has an integer solution $b$.

\end{theorem}

\begin{proof} ($\Rightarrow$) Suppose $n = f(x_0, y_0)$
	with $\gcd(x_0, y_0) = 1$. Choose $x_1, y_1 \in \mathbb{Z}$ with $x_0
		y_1 - x_1 y_0 = 1$; the substitution $(x, y) = (x_0 x' + x_1 y',\ y_0
		x' + y_1 y')$ lies in $\operatorname{SL}_2(\mathbb{Z})$ and carries
	$f$ into a form $g(x', y') = n x'^{2} + b x'y' + c y'^{2}$ with $g(1,
		0) = n$ and the same discriminant. Then $b^{2} - 4nc = D$, so $b^{2}
		\equiv D \pmod{4n}$. ($\Leftarrow$) If $b^{2} \equiv D \pmod{4n}$,
	then $4n \mid b^{2} - D$, so $c = \tfrac{b^{2} - D}{4n} \in
		\mathbb{Z}$, and the form $g(x, y) = n x^{2} + b xy + c y^{2}$ has
	discriminant $b^{2} - 4nc = D$ and represents $n = g(1, 0)$ properly.\qedhere

\end{proof}

\begin{theorem}\label{thm:8.6.7} A positive integer $n$ is represented by the
	principal form $x^{2} + xy + 2y^{2}$ of discriminant $D = -7$ if and
	only if every odd prime $p\neq 7$ with $\leg{-7}{p} = -1$
	(equivalently $p \equiv 3, 5, 6 \pmod{7}$) occurs to an even power in $n$.

\end{theorem}

\begin{proof} For an odd prime $p\neq7$, quadratic reciprocity and the
	first supplementary law give $\leg{-7}{p}=\leg{p}{7}$.
	The nonzero squares modulo $7$ are $1,2,4$, so the non-residue condition
	is equivalent to $p\equiv3,5,6\pmod7$.

	The form $f(x, y) = x^{2} + xy + 2y^{2}$ is
	the norm form of $\mathbb{Z}[\omega]$, where $\omega = \tfrac{1 +
			\sqrt{-7}}{2}$ (so $\omega^{2} = \omega - 2$), namely $f(x, y) = (x +
		y\omega)(x + y\bar\omega)$. From $(x + y\omega)(u + v\omega) = (xu -
		2yv) + (xv + yu + yv)\omega$ one gets the composition identity
	\[
		(x^{2} + xy + 2y^{2})(u^{2} + uv + 2v^{2}) = X^{2} + XY + 2Y^{2},
		\quad X = xu - 2yv,\ \ Y = xv + yu + yv,
	\]
	so the set of represented integers is closed under multiplication.
	Since $h(-7) = 1$, Theorem~\ref{thm:8.6.5} shows that an odd
	prime $p\neq 7$ has the form $p = f(x, y)$ if and only if $\leg{-7}{p} = 1$; moreover, the two
	remaining primes are represented directly, namely $2 = f(0, 1)$ and
	$7 = f(-1, 2)$. Conversely, suppose that an odd prime $p\neq 7$ with
	$\leg{-7}{p} = -1$ divides $f(x, y)$, and note that $4 f(x, y)
		= (2x + y)^{2} + 7 y^{2}$. If we had $p \nmid y$, we would get $(2x +
		y)^{2} \equiv -7 y^{2} \pmod{p}$, that is, $\leg{-7}{p} =
		1$, a contradiction. Hence $p \mid y$; it follows that $p \mid (2x + y)$, so $p
		\mid x$ and $p^{2} \mid f(x, y)$. ($\Rightarrow$) If $n$ is represented by $f$ and an odd prime $p\neq 7$ with
	$\leg{-7}{p} = -1$ divides $n$, the divisor property just shown
	gives $p \mid x$ and $p \mid y$, so $p^{2} \mid n$ and $n/p^{2} =
		f(x/p, y/p)$; iterating shows that $v_p(n)$ is even. ($\Leftarrow$)
	Conversely, if every prime $p$ with $\leg{-7}{p} = -1$ occurs to an
	even power in $n$, write $n = s^{2} m$ with $m$ squarefree; then
	every prime dividing $m$ is $2$, $7$, or an odd prime with
	$\leg{-7}{p} = 1$. Each is represented by $f$ ($2 = f(0, 1)$, $7
		= f(-1, 2)$, and every such odd $p = f(x, y)$ by Theorem~\ref{thm:8.6.5}),
	so by the composition identity $m = f(u, v)$ for some integers $u,
		v$, whence $n = s^{2} f(u, v) = f(su, sv)$.

\end{proof}
\begin{theorem}\label{thm:8.6.8} The discriminant $D = -20$ has class number $h(-20) =
		2$, and its reduced primitive positive-definite forms are the principal
	form $x^{2} + 5y^{2}$ and $2x^{2} + 2xy + 3y^{2}$, which lie in
	different classes. For a prime $p \neq 2, 5$:

	\begin{enumerate}[label={\alph*)}]
		\item $p$ is represented by one of the two forms if and only if
			      $\leg{-5}{p} = 1$, or equivalently $p \equiv 1, 3, 7, 9 \pmod{20}$;
		\item $p = x^{2} + 5y^{2}$ if and only if $p \equiv 1, 9 \pmod{20}$;
		\item $p = 2x^{2} + 2xy + 3y^{2}$ if and only if $p \equiv 3, 7 \pmod{20}$.
	\end{enumerate}

\end{theorem}

\begin{proof} Since $\leg{-20}{p} =
		\leg{-1}{p}\leg{4}{p}\leg{5}{p}
		= \leg{-5}{p}$, the hypothesis $\leg{-5}{p} = 1$ is equivalent to
	$\leg{-20}{p} = 1$. Theorem~\ref{thm:8.6.6} with $n = p$ requires a solution of
	$b^{2} \equiv -20 \pmod{4p}$: if $\leg{-20}{p} = 1$, pick $b_0$ with
	$b_0^{2} \equiv -20 \pmod{p}$ and, by the Chinese remainder theorem,
	an even $b$ with $b \equiv b_0 \pmod{p}$; then $b^{2} \equiv -20
		\pmod{p}$ and $b^{2} \equiv 0 \equiv -20 \pmod{4}$, so $b^{2} \equiv
		-20 \pmod{4p}$. Conversely, a solution of $b^{2} \equiv -20 \pmod{4p}$
	gives $\leg{-20}{p} = 1$ because $p \nmid 20$. The resulting form
	$g = (p, b, c)$ with $c = \tfrac{b^{2}+20}{4p}$ is primitive: a common
	prime divisor $e$ of $p, b, c$ would divide $b^{2} - 4pc = -20$ and
	also $p$, and hence $p = e \in \{2, 5\}$, which is excluded. Hence $p$ is
	represented by some primitive form of discriminant $-20$ if and only if
	$\leg{-5}{p} = 1$; and $\leg{-5}{p} =
		\leg{-1}{p}\leg{p}{5}$ depends only on
	$p \bmod 20$, and checking the eight residue classes coprime to $20$ shows that
	$\leg{-5}{p} = 1$ if and only if $p \equiv 1, 3, 7, 9 \pmod{20}$.
	The two listed forms are exactly the reduced primitive forms of
	discriminant $-20$: a reduced form $(a,b,c)$ satisfies $|D| = 4ac - b^2 \ge 4a^2 - a^2 = 3a^2$ (from $|b| \le a \le c$), so $a \le
		\sqrt{|D|/3} = \sqrt{20/3} < 2.6$, whence $a \in \{1,
		2\}$, while $b \equiv D \equiv 0 \pmod{2}$ and $4ac = b^{2} + 20$. For
	$a = 1$, only $b = 0,\ c = 5$ works, giving $(1, 0, 5) = x^{2} +
		5y^{2}$, and for $a = 2$, only $b = \pm 2,\ c = 3$ works, where the
	boundary rule $b \ge 0$ (as $|b| = a$) selects $(2, 2, 3) = 2x^{2} +
		2xy + 3y^{2}$. Both are primitive, so $h(-20) = 2$.
	To distinguish the two forms, note that $5 \equiv 1 \pmod{4}$ gives $x^{2} + 5y^{2}
		\equiv x^{2} + y^{2} \pmod{4}$, so an odd value of $x^{2} + 5y^{2}$ is
	congruent to $1 \pmod{4}$; hence $p = x^{2} + 5y^{2}$ forces $p \equiv 1
		\pmod{4}$, that is, $p \equiv 1, 9 \pmod{20}$. An odd prime value of
	$2x^{2} + 2xy + 3y^{2}$ forces $y$ to be odd, whence $2x(x + y) \equiv 0
		\pmod{4}$ and the value is congruent to $3 \pmod{4}$, that is, $p \equiv 3, 7
		\pmod{20}$. As every such $p$ is represented by one of the two forms
	and the sets $\{1, 9\}$ and $\{3, 7\}$ are disjoint, the three equivalences
	follow.

\end{proof}
\subsection{The Quadratic Rings \texorpdfstring{$\mathbb{Z}[\sqrt{d}]$}{Z[sqrt d]}, Residue Units, and Lifting the Exponent}\label{subsec:8.7}

\begin{definition}[{the quadratic ring $\mathbb{Z}[\sqrt{d}]$}]\label{def:8.7.1}
	Let $d$ be an integer that is not a square, and fix a complex square
	root $\sqrt d$. Put
	\[
		A_d=\mathbb{Z}[\sqrt d]
		=\{a+b\sqrt d\mid a,b\in\mathbb{Z}\}.
	\]
	Each element has a unique pair of integer coordinates, because $d$
	is not the square of a rational number
	(Theorem~\ref{thm:2.3.7}). Addition is
	componentwise, and multiplication follows the rule
	\[
		(a+b\sqrt d)(c+e\sqrt d)
		=(ac+dbe)+(ae+bc)\sqrt d.
	\]
	Thus $A_d$ is a subring of $\mathbb{C}$.
	For $\alpha=a+b\sqrt d$, define
	\[
		\bar\alpha=a-b\sqrt d,
		\qquad N(\alpha)=\alpha\bar\alpha=a^2-db^2.
	\]
	Divisibility $\gamma\mid\alpha$ means $\alpha=\gamma\delta$ for some
	$\delta\in A_d$. In particular, an integer $n\neq0$ divides
	$a+b\sqrt d$ exactly when it divides both $a$ and $b$.

	A unit has an inverse in $A_d$. Two elements are associates if one
	is a unit times the other. A nonzero nonunit $\pi$ is irreducible if
	$\pi=\alpha\beta$ with $\alpha,\beta\in A_d$ forces $\alpha$ or
	$\beta$ to be a unit. These definitions agree with
	those for Gaussian integers when $d=-1$.

\end{definition}

\begin{theorem}[conjugation, norm, and cancellation]\label{thm:8.7.1}
	For $\alpha,\beta\in A_d$, conjugation preserves sums and products,
	and
	\[
		N(\alpha\beta)=N(\alpha)N(\beta).
	\]
	A nonzero element has nonzero norm. If $\gamma\neq0$ and
	$\alpha\gamma=\beta\gamma$, then $\alpha=\beta$.
	An element $\alpha$ is a unit exactly when $N(\alpha)=\pm1$.

\end{theorem}

\begin{proof}
	The coordinate formulas of Definition~\ref{def:8.7.1} give
	$\overline{\alpha+\beta}=\bar\alpha+\bar\beta$ and
	$\overline{\alpha\beta}=\bar\alpha\bar\beta$. Thus
	\[
		N(\alpha\beta)=\alpha\beta\bar\alpha\bar\beta
		=(\alpha\bar\alpha)(\beta\bar\beta).
	\]
	If $a^2-db^2=0$ and $b\neq0$, then $d=(a/b)^2$, so $d$ is the square of an
	integer by Theorem~\ref{thm:2.3.7}. This contradicts the
	choice of $d$. If $b=0$, the equation gives $a=0$.
	Consequently $N(\alpha)=0$ only for $\alpha=0$.

	From $(\alpha-\beta)\gamma=0$ we obtain
	$N(\alpha-\beta)N(\gamma)=0$, so $\alpha=\beta$ when
	$\gamma\neq0$. Finally, an inverse forces
	$N(\alpha)N(\alpha^{-1})=1$, hence $N(\alpha)=\pm1$.
	Conversely, if $N(\alpha)=\varepsilon\in\{1,-1\}$, then
	$\alpha(\varepsilon\bar\alpha)=1$.

\end{proof}

\begin{definition}[ideals, domains, and division]\label{def:8.7.2}
	Let $R$ be a commutative ring with identity
	(Definition~\ref{def:6.6.4}). An ideal of $R$ is a subset containing
	$0$, closed under sums and negation, and closed under multiplication
	by every element of $R$. It is principal if it has the form
	$\gamma R=\{\gamma x\mid x\in R\}$ for one element $\gamma$.
	The ring $R$ is a domain if $1\neq0$ in $R$ and $R$ has no zero
	divisors. A domain is a principal ideal domain (PID) if every ideal
	in it is principal. Each $A_d$ is a domain, because it is a subring
	of $\mathbb{C}$.

	We call $A_d$ Euclidean if there is a function
	$\delta:A_d\setminus\{0\}\to\mathbb{Z}_{\ge0}$ such that, for
	every $\alpha,\beta\in A_d$ with $\beta\neq0$, one can write
	$\alpha=\beta q+r$ with $q,r\in A_d$ and either $r=0$ or
	$\delta(r)<\delta(\beta)$.
	It has unique factorization if every nonzero nonunit is a product
	of irreducibles, uniquely determined up to order and associates. Such a ring
	is also called a unique factorization domain (UFD).

\end{definition}

\begin{theorem}[failure of unique factorization]\label{thm:8.7.2}
	In $A_{-5}=\mathbb{Z}[\sqrt{-5}]$, the units are $\pm1$ and
	\[
		6=2\cdot3=(1+\sqrt{-5})(1-\sqrt{-5})
	\]
	gives two factorizations into irreducibles that do not agree up to
	order and associates. The set
	\[
		I=\{2x+(1+\sqrt{-5})y\mid x,y\in A_{-5}\}
	\]
	cannot equal $\gamma A_{-5}$ for any $\gamma\in A_{-5}$.
	Thus this concrete ideal is not principal, and $A_{-5}$ is neither
	a PID nor Euclidean.

\end{theorem}

\begin{proof}
	Here $N(a+b\sqrt{-5})=a^2+5b^2$ is positive for nonzero elements.
	The equation $N(\alpha)=1$ gives precisely $\alpha=\pm1$.
	There is no element of norm $2$ or $3$: if $b\neq0$, the norm is at
	least $5$, and if $b=0$, it is a square. The four displayed factors
	have norms $4,9,6,6$. A factorization into two nonunits would require
	positive norms greater than $1$ whose product is one of these
	numbers. It would therefore require a norm $2$ or $3$, which is
	impossible. All four factors are irreducible. Associates have the
	same norm here, so the factor of norm $4$ in the first factorization
	has no associate in the second.

	The set $I$ is closed under addition, negation, and multiplication
	by elements of $A_{-5}$. If $I=\gamma A_{-5}$, then $\gamma\neq0$
	and $\gamma$ divides both $2$ and $1+\sqrt{-5}$. Hence
	$N(\gamma)$ divides both $4$ and $6$, so it is $1$ or $2$.
	Norm $2$ is impossible, and norm $1$ would make $\gamma$ a unit,
	giving $I=A_{-5}$. However, the map
	\[
		a+b\sqrt{-5}\longmapsto a+b\pmod2
	\]
	preserves sums and products, since $-5\equiv1\pmod2$.
	It sends both generators of $I$ to $0$ and sends $1$ to $1$.
	Thus $1\notin I$, a contradiction.

	Suppose a size function $\delta:A_{-5}\setminus\{0\}\to
	\mathbb{Z}_{\ge0}$ allowed division with a zero remainder or a
	remainder of smaller size. Choose $\gamma\in I\setminus\{0\}$
	with $\delta(\gamma)$ minimal. For any $\alpha\in I$, division gives
	$\alpha=\gamma q+r$. Since $r\in I$, minimality forces $r=0$.
	This would give $I=\gamma A_{-5}$, which we have excluded.

\end{proof}

\begin{definition}[residue pairs and their units]\label{def:8.7.3}
	For $n\ge2$, write $\alpha\equiv\beta\pmod n$ in $A_d$ when both
	coordinates of $\alpha-\beta$ are divisible by $n$. Reducing the two
	coordinates gives the ring $R_n(d)$ of Definition~\ref{def:7.9.1}.
	Thus a residue is a pair $a+b\sqrt d$ with $a,b$ taken modulo $n$,
	and the multiplication rule remains
	\[
		(a+b\sqrt d)(c+e\sqrt d)
		=(ac+dbe)+(ae+bc)\sqrt d\pmod n.
	\]
	For $\gcd(n,2d)=1$, let $\psi_d(n)$ count the residues with a
	multiplicative inverse. Set $\psi_d(1)=1$.

\end{definition}

\begin{theorem}[reduction modulo a prime]\label{thm:8.7.3}
	Let $p$ be an odd prime with $p\nmid d$. Reduction of coordinates
	identifies the residues of $A_d$ modulo $p$ with
	\[
		\mathbb{F}_p[T]/(T^2-d),
	\]
	where two polynomials represent the same element when their
	difference is a multiple of $T^2-d$.
	If $\leg{d}{p}=-1$, this ring is a field with $p^2$ elements.
	If $\leg{d}{p}=1$ and $s^2=d$ in $\mathbb{F}_p$, the map
	\[
		a+b\sqrt d\longmapsto(a+bs,a-bs)
	\]
	identifies it with $\mathbb{F}_p\times\mathbb{F}_p$.
	These two cases are called inert and split, respectively.

\end{theorem}

\begin{proof}
	Division by the monic polynomial $T^2-d$ leaves a unique remainder
	$a+bT$. A nonzero multiple of $T^2-d$ has degree at least $2$, so
	two such remainders represent the same polynomial class exactly
	when their coordinates agree. Replacing $T^2$ by $d$ gives precisely
	the product rule for $R_p(d)$. This proves the first identification.

	By Definition~\ref{def:8.7.3}, the residues of $A_d$ modulo
	$p$ form the ring $R=R_p(d)$ of Definition~\ref{def:7.9.1}, and $p$
	and $d$ satisfy the standing assumptions made there.
	If $\leg{d}{p}=-1$, this ring is a field with $p^2$ elements by
	Theorem~\ref{thm:7.9.1}(a). If $\leg{d}{p}=1$ and $s^2=d$ in
	$\mathbb{F}_p$, then $s\neq0$ because $p\nmid d$, and
	Theorem~\ref{thm:7.9.5} shows that the stated map is a ring
	isomorphism onto $\mathbb{F}_p\times\mathbb{F}_p$.

\end{proof}

\begin{theorem}[unit counts]\label{thm:8.7.4}
	A residue $z=a+b\sqrt d$ modulo any $n\ge2$ is a unit exactly when
	\[
		\gcd(a^2-db^2,n)=1.
	\]
	For positive integers $m,n$ with $\gcd(m,n)=1$ and
	$\gcd(mn,2d)=1$,
	\[
		\psi_d(mn)=\psi_d(m)\psi_d(n).
	\]
	For an odd prime $p\nmid d$ and $k\ge1$,
	\[
		\psi_d(p^k)=p^{2k-2}(p-1)\bigl(p-\leg{d}{p}\bigr).
	\]
	Consequently, for every $n\ge1$ with $\gcd(n,2d)=1$,
	\[
		\psi_d(n)=n^2\prod_{p\mid n}
		\left(1-\frac1p\right)
		\left(1-\frac{\leg{d}{p}}p\right),
	\]
	where the empty product equals $1$.

\end{theorem}

\begin{proof}
	The polynomial norm identity remains valid modulo $n$. If $zw=1$,
	then $N(z)N(w)=1\pmod n$, so $N(z)$ is coprime to $n$.
	Conversely, if $cN(z)\equiv1\pmod n$, then $c\bar z$ is an inverse
	of $z$. This also shows that the criterion is independent of the
	chosen integer representatives of $a$ and $b$.

	For coprime $m,n\ge2$, the Chinese remainder theorem
	(Theorem~\ref{thm:3.4.1}), applied to each coordinate, gives a
	bijection
	\[
		R_{mn}(d)\longrightarrow R_m(d)\times R_n(d).
	\]
	The coordinate formulas show that it preserves sums, products, and the identity. A unit maps to units in both components. Conversely,
	the Chinese remainder theorem combines their two inverses into an
	inverse modulo $mn$.
	This proves multiplicativity. Cases with $m=1$ or $n=1$ follow
	from the definition of $\psi_d(1)$.

	For $k=1$, the inert case has $p^2-1$ units. In the split case a
	pair $(u,v)$ is invertible exactly when $u,v\neq0$, giving
	$(p-1)^2$ units. Both counts equal
	$(p-1)(p-\leg{d}{p})$.
	A pair modulo $p^k$ is a unit exactly when its norm is nonzero
	modulo $p$, or equivalently its reduction modulo $p$ is a unit.
	Each coordinate modulo $p$ has $p^{k-1}$ lifts, so each pair has
	$p^{2k-2}$ lifts. This proves the prime-power formula.
	Multiplicativity and the prime factorization of $n$ give the product
	formula.

\end{proof}

\begin{definition}[inert prime valuation]\label{def:8.7.4}
	Let $p$ be an odd prime with $\leg{d}{p}=-1$. For
	$\gamma=a+b\sqrt d\in A_d$, put
	\[
		w_p(\gamma)=\min\{v_p(a),v_p(b)\},
	\]
	with $v_p(0)=+\infty$. Thus $w_p(0)=+\infty$, and for nonzero
	$\gamma$ the value is the largest $k\ge0$ such that $p^k$ divides
	both coordinates. We use $(+\infty)+k=+\infty$, including $k=+\infty$.
	For a nonzero rational integer $m$, this gives
	$w_p(m)=v_p(m)$.

\end{definition}

\begin{theorem}[valuation laws]\label{thm:8.7.5}
	The value $w_p(\gamma)$ is a finite nonnegative integer for
	$\gamma\neq0$. For all $\gamma,\delta\in A_d$,
	\[
		w_p(\gamma\delta)=w_p(\gamma)+w_p(\delta),
		\qquad
		w_p(\gamma+\delta)\ge
		\min\{w_p(\gamma),w_p(\delta)\}.
	\]
	The latter inequality is an equality when the two values differ.
	For $\gamma\neq0$,
	\[
		v_p\bigl(N(\gamma)\bigr)=2w_p(\gamma).
	\]

\end{theorem}

\begin{proof}
	A nonzero element has at least one nonzero integer coordinate,
	whose ordinary valuation is finite. Suppose first that both
	$\gamma$ and $\delta$ are nonzero. Write
	$\gamma=p^a\gamma_0$ and $\delta=p^b\delta_0$, where
	$a=w_p(\gamma)$, $b=w_p(\delta)$, and both reduced pairs are
	nonzero modulo $p$. The inert residue ring is a field by
	Theorem~\ref{thm:8.7.3}. Hence their product is nonzero
	modulo $p$, giving $w_p(\gamma\delta)=a+b$.
	If either factor is zero, the same formula follows from the
	convention for $+\infty$.

	If either summand is zero, both assertions concerning sums hold directly.
	Otherwise both valuations are finite. The sum inequality follows
	by dividing both coordinates by the
	common power $p^{\min\{a,b\}}$. If $a<b<\infty$, reduction after
	division by $p^a$ leaves the nonzero pair $\gamma_0$.
	This gives equality. The case $b<a$ is symmetric.
	Finally, a nonzero pair modulo $p$ is a unit of the inert residue
	field, so its norm is nonzero modulo $p$ by
	Theorem~\ref{thm:8.7.4}. Thus $p\nmid N(\gamma_0)$ and
	$N(\gamma)=p^{2a}N(\gamma_0)$, proving the norm formula.

\end{proof}

\begin{definition}[admissible pair]\label{def:8.7.5}
	For the rest of this subsection, $p$ is an odd prime with
	$\leg{d}{p}=-1$, and
	$\alpha,\beta\in A_d$ satisfy
	\[
		\alpha\neq\beta,\qquad p\mid\alpha-\beta,
		\qquad p\nmid N(\alpha).
	\]
	Call $(\alpha,\beta)$ an admissible pair. Both reduce to the same
	nonzero unit modulo $p$, so $w_p(\alpha)=w_p(\beta)=0$.
	The value $w_p(\alpha-\beta)$ is finite and at least $1$ (Definition~\ref{def:8.7.1}).

\end{definition}

\begin{theorem}[exponents coprime to $p$]\label{thm:8.7.6}
	For an admissible pair and an integer $m\ge1$ with $p\nmid m$,
	\[
		w_p(\alpha^m-\beta^m)=w_p(\alpha-\beta).
	\]
	In particular, this holds for a rational prime exponent $q\neq p$.

\end{theorem}

\begin{proof}
	The finite geometric identity gives
	\[
		\alpha^m-\beta^m=(\alpha-\beta)S,
		\qquad S=\sum_{j=0}^{m-1}\alpha^j\beta^{m-1-j}.
	\]
	Modulo $p$, the sum is $m\beta^{m-1}$, which is nonzero because
	$p\nmid m$ and $\beta$ is a unit. Hence $w_p(S)=0$.
	The product law proves the formula and shows that
	$\alpha^m-\beta^m\neq0$.

\end{proof}

\begin{theorem}[exponent $p$]\label{thm:8.7.7}
	For an admissible pair,
	\[
		w_p(\alpha^p-\beta^p)=w_p(\alpha-\beta)+1.
	\]

\end{theorem}

\begin{proof}
	Put $h=\alpha-\beta=p^r\gamma$, where
	$r=w_p(\alpha-\beta)\ge1$ and $p\nmid\gamma$.
	The binomial expansion gives
	\[
		\alpha^p-\beta^p
		=ph\beta^{p-1}
		+\sum_{j=2}^{p-1}\binom pj h^j\beta^{p-j}+h^p.
	\]
	The first term has valuation $r+1$.
	For $2\le j\le p-1$, the coefficient $\binom pj$ is divisible
	by $p$, so the corresponding term is divisible by $p^{1+jr}$,
	and $1+jr\ge r+2$.
	The last term is divisible by $p^{pr}$, with $pr\ge r+2$
	because $p\ge3$ and $r\ge1$. Thus, for some $\eta\in A_d$,
	\[
		\alpha^p-\beta^p
		=p^{r+1}\bigl(\gamma\beta^{p-1}+p\eta\bigr).
	\]
	The element in parentheses is nonzero modulo $p$, since both
	$\gamma$ and $\beta$ are nonzero there and the residue ring is a
	field. Its valuation is $0$, proving the formula.

\end{proof}

\begin{theorem}[lifting the exponent in the inert case]\label{thm:8.7.8}
	For an admissible pair and every integer $n\ge1$,
	\[
		w_p(\alpha^n-\beta^n)
		=w_p(\alpha-\beta)+v_p(n).
	\]
	In particular, $\alpha^n\neq\beta^n$. When $\alpha,\beta$ are
	rational integers, this is the ordinary odd-prime LTE formula of
	Theorem~\ref{thm:4.8.3}.

\end{theorem}

\begin{proof}
	Write $n=p^a m$ with $p\nmid m$. Theorem~\ref{thm:8.7.6}
	gives a nonzero difference $\alpha^m-\beta^m$ with the original
	valuation. The elements $\alpha^m$ and $\beta^m$ are still congruent modulo
	$p$, and their reductions modulo $p$ remain units.
	Apply Theorem~\ref{thm:8.7.7} successively to
	\[
		(\alpha^m,\beta^m),\quad
		(\alpha^{pm},\beta^{pm}),\quad\ldots,\quad
		(\alpha^{p^{a-1}m},\beta^{p^{a-1}m}).
	\]
	Each step gives a finite valuation, so the next difference is
	nonzero and the next pair is admissible. The valuation increases
	by $1$ at each of the $a$ steps. Since $a=v_p(n)$, the result follows.
	For $a=0$, the coprime-exponent result already proves the formula.

\end{proof}

\subsection{Eisenstein Integers and Residue Counts}

\begin{definition}[Eisenstein integers]\label{def:8.8.1}
	Let $\omega=(-1+\sqrt{-3})/2$, so that $\omega^2+\omega+1=0$.
	The Eisenstein integers form the ring
	\[
		E=\mathbb{Z}[\omega]=\{a+b\omega\mid a,b\in\mathbb{Z}\}.
	\]
	For $\alpha=a+b\omega$, its conjugate and norm are
	\[
		\bar\alpha=a+b\omega^2=(a-b)-b\omega,
		\qquad N(\alpha)=\alpha\bar\alpha=a^2-ab+b^2.
	\]
	A unit has an inverse in $E$. Two elements are associates if they differ
	by a unit factor. A nonzero nonunit $\pi$ is irreducible if
	$\pi=\alpha\beta$ forces one factor to be a unit. It is prime if
	$\pi\mid\alpha\beta$ forces $\pi\mid\alpha$ or $\pi\mid\beta$.

\end{definition}

\begin{theorem}[norm and units]\label{thm:8.8.1}
	For $\alpha,\beta\in E$, conjugation is multiplicative and
	$N(\alpha\beta)=N(\alpha)N(\beta)$. The norm is a nonnegative integer,
	and it is zero only at zero. The units of $E$ are exactly
	$\{\pm1,\pm\omega,\pm\omega^2\}$.

\end{theorem}

\begin{proof}
	The relation $\omega^2=-1-\omega$ gives
	\[
		(a+b\omega)(c+d\omega)=(ac-bd)+(ad+bc-bd)\omega.
	\]
	Conjugation replaces $\omega$ by $\omega^2$, which satisfies the same
	relation, so direct substitution proves its multiplicativity. Thus
	$N(\alpha\beta)=\alpha\beta\bar\alpha\bar\beta=N(\alpha)N(\beta)$.
	Also
	\[
		N(a+b\omega)=(a-b/2)^2+3b^2/4.
	\]
	This proves positivity away from zero. If $u$ is a unit, multiplicativity
	forces $N(u)=1$. Conversely, $N(u)=1$ gives $u\bar u=1$.
	The identity $4N(a+b\omega)=(2a-b)^2+3b^2$ shows that $N=1$ forces
	$b\in\{-1,0,1\}$. Checking these three values gives
	\[
		(a,b)=(1,0),(-1,0),(0,1),(0,-1),(1,1),(-1,-1),
	\]
	which are the six listed units. Since $E$ is a subring of $\mathbb{C}$,
	it also has cancellation for nonzero factors.

\end{proof}

\begin{theorem}[division and factorization]\label{thm:8.8.2}
	Given $\alpha,\beta\in E$ with $\beta\neq0$, there are $q,r\in E$
	such that $\alpha=\beta q+r$ and either $r=0$ or $N(r)<N(\beta)$.
	The Euclidean algorithm and B{\'e}zout's identity hold in $E$.
	Every irreducible is prime, every prime is irreducible, and every nonzero
	nonunit has a factorization into primes. Such a factorization is unique
	up to the order of the factors and multiplication by units.

\end{theorem}

\begin{proof}
	Using the identity $\alpha/\beta=\alpha\bar\beta/N(\beta)$, write
	$\alpha/\beta=x+y\omega$ with $x,y\in\mathbb{Q}$. Choose integers $a,b$
	with $|x-a|,|y-b|\le1/2$, and put $q=a+b\omega$ and $r=\alpha-\beta q$.
	Then
	\[
		\frac{N(r)}{N(\beta)}=(x-a)^2-(x-a)(y-b)+(y-b)^2\le\frac34<1.
	\]
	Successive nonzero remainders have strictly decreasing positive integer
	norms, so the Euclidean algorithm terminates. Its last nonzero remainder
	$\delta$ divides both inputs, and back-substitution gives
	$\delta=u\alpha+v\beta$. Every common divisor therefore divides $\delta$.
	This establishes the greatest-common-divisor property and B{\'e}zout's identity.
	The same statements hold when one input is zero by taking the other input
	as the greatest common divisor.

	Let $\pi$ be irreducible and suppose $\pi\mid\alpha\beta$ but
	$\pi\nmid\alpha$. A greatest common divisor of $\pi$ and $\alpha$ must
	be a unit, since every divisor of $\pi$ is a unit or an associate of $\pi$.
	Multiplying its B{\'e}zout identity by its inverse gives
	$u\pi+v\alpha=1$. Multiplication by $\beta$ now proves $\pi\mid\beta$.
	Thus every irreducible is prime. Conversely, if a prime $\pi$ equals
	$\alpha\beta$, it divides one factor. If $\alpha=\pi\gamma$, cancellation
	in $\pi=\pi\gamma\beta$ gives $\gamma\beta=1$. Hence $\beta$ is a unit,
	and a prime is irreducible.

	Existence of factorization follows by induction on the norm. A reducible
	nonzero nonunit factors into two nonunits, each with norm strictly between
	$1$ and the original norm. For uniqueness, a prime in one factorization
	divides some prime in the other. The two primes are associates.
	Cancel them and repeat to match all factors, with the remaining unit
	absorbed into the unit factor.

\end{proof}

\begin{theorem}[principal ideals in the two concrete rings]\label{thm:8.8.3}
	Let $R$ be either $\mathbb{Z}[i]$ or $E$. Every ideal $I$ of $R$, in the
	sense of Definition~\ref{def:8.7.2}, has the form
	\[
		I=\delta R=\{\delta z\mid z\in R\}
	\]
	for some $\delta\in R$. Thus both rings are principal ideal domains.

\end{theorem}

\begin{proof}
	For $I=\{0\}$, take $\delta=0$. Otherwise, since the norms of nonzero elements are positive integers, we may choose a nonzero
	$\delta\in I$ of least norm.
	For any $\alpha\in I$, division with remainder in
	Theorem~\ref{thm:8.2.2} or Theorem~\ref{thm:8.8.2} gives
	$\alpha=\delta q+r$, with $r=0$ or $N(r)<N(\delta)$.
	The ideal properties give $r=\alpha-\delta q\in I$.
	Minimality of $N(\delta)$ forces $r=0$, so $I\subseteq\delta R$.
	The reverse inclusion follows from $\delta\in I$ and closure under
	multiplication. Both rings are subrings of $\mathbb{C}$ and have no zero
	divisors, so they are domains with the stated principal ideal property.

\end{proof}

\begin{theorem}[rational primes in the Eisenstein integers]\label{thm:8.8.4}
	The rational primes factor in $E$ as follows:
	\begin{enumerate}[label={\alph*)}]
		\item $3=-\omega^2(1-\omega)^2$, where $1-\omega$ is prime of norm $3$.
		\item If $p\equiv1\pmod3$, then $p=\pi\bar\pi$ for nonassociate primes
		      $\pi,\bar\pi$ with $N(\pi)=p$.
		\item If $p\equiv2\pmod3$, then $p$ is prime in $E$, with $N(p)=p^2$.
	\end{enumerate}
	Every prime of $E$ is an associate of one of the prime elements in this list.

\end{theorem}

\begin{proof}
	An element of rational-prime norm is irreducible by multiplicativity of
	the norm, and hence prime by Theorem~\ref{thm:8.8.2}.
	Now $N(1-\omega)=3$ and $(1-\omega)^2=-3\omega$, which prove (a).
	For (b), Theorem~\ref{thm:8.1.6} gives integers $a,b$ with
	$p=a^2-ab+b^2$. Put $\pi=a+b\omega$, so that $p=\pi\bar\pi$ and both
	factors have norm $p$.

	To show they are nonassociate, suppose $\pi\mid\bar\pi$. Then
	$\pi\mid\pi-\bar\pi=b(\omega-\omega^2)$. Since
	$N(\omega-\omega^2)=3$ and $N(\pi)=p\neq3$, the element $\pi$ cannot
	divide $\omega-\omega^2$. Its prime-divisor property gives $\pi\mid b$,
	and $\pi\mid a+b\omega$ then gives $\pi\mid a$. Taking norms yields
	$p\mid b^2$ and $p\mid a^2$, so the rational prime $p$ divides both
	$a$ and $b$. This would force $p^2\mid a^2-ab+b^2=p$, a contradiction.

	For (c), first suppose $p\mid a^2-ab+b^2$. If $p\nmid b$, put
	$t\equiv ab^{-1}\pmod p$. Then $t^2-t+1=0$, so $u=-t$ satisfies
	$u^2+u+1=0$ and $u^3=1$ modulo $p$. Since $p\neq3$, this $u$ is not $1$.
	It has order $3$, contradicting $3\nmid p-1$ by Fermat's little theorem.
	Thus $p\mid b$, and then $p\mid a$. We have proved
	\[
		p\mid N(\alpha)\quad\Longrightarrow\quad p\mid\alpha\text{ in }E.
	\]
	If $p\mid\alpha\beta$, taking norms shows that $p$ divides $N(\alpha)$
	or $N(\beta)$, and the implication just proved shows that $p$ divides
	$\alpha$ or $\beta$. Since its norm $p^2$ exceeds $1$, $p$ is a nonunit in $E$, and hence is prime in $E$.
	This argument also covers $p=2$.

	Finally, let $\rho$ be any prime of $E$. It divides the positive integer
	$N(\rho)=\rho\bar\rho>1$. Factoring this integer into rational primes
	shows that $\rho$ divides some rational prime $p$. Apply (a), (b), or
	(c), and use the prime-divisor property to see that $\rho$ divides one of
	the listed prime factors. Since that factor is irreducible, $\rho$ is its
	associate.

\end{proof}

\begin{definition}[congruence classes]\label{def:8.8.2}
	Let $R$ be either $\mathbb{Z}[i]$ or $E$, and let $0\neq\alpha\in R$.
	Write $z\equiv w\pmod\alpha$ if $\alpha\mid z-w$ in $R$.
	The class of $z$ is written $[z]$, and the set of these classes is
	denoted by $R/(\alpha)$, with addition and
	multiplication defined by $[z]+[w]=[z+w]$ and $[z][w]=[zw]$.
	These operations are well defined: divisibility of $z-z'$ and $w-w'$ by
	$\alpha$ implies divisibility of $(z+w)-(z'+w')$ and
	$zw-z'w'=(z-z')w+z'(w-w')$ by $\alpha$.

\end{definition}

\begin{theorem}[counting integer-pair classes]\label{thm:8.8.5}
	Let $v=(a,b)$ and $w=(c,d)$ be integer vectors with $ad-bc\neq0$.
	Modulo the subgroup $L=\{mv+nw\mid m,n\in\mathbb{Z}\}$ of
	$\mathbb{Z}^2$, there are exactly $|ad-bc|$ classes.

\end{theorem}

\begin{proof}
	Set $g=\gcd(a,b)>0$, and choose $s,t\in\mathbb{Z}$ with $sa+tb=g$.
	The integer matrix
	\[
		U=\begin{pmatrix}s&t\\-b/g&a/g\end{pmatrix}
	\]
	has determinant $1$ and an integer inverse. It therefore gives a bijection
	on integer vectors and on their classes. Its images of $v$ and $w$ are
	\[
		Uv=(g,0),\qquad Uw=(C,D),\qquad D=(ad-bc)/g\neq0.
	\]
	For any $(x,y)\in\mathbb{Z}^2$, subtract a multiple of $(C,D)$ so that
	$0\le y<|D|$. Then subtract a multiple of $(g,0)$ so that $0\le x<g$.
	Every class has a representative in this rectangle. If two such
	representatives differ by $m(g,0)+n(C,D)$, their second coordinates differ
	by $nD$. The range forces $n=0$, and the first-coordinate range then
	forces $m=0$. The representatives are unique, so there are
	$g|D|=|ad-bc|$ classes.

\end{proof}

\begin{theorem}[Gaussian and Eisenstein residue counts]\label{thm:8.8.6}
	For every nonzero $\alpha$ in either $R=\mathbb{Z}[i]$ or $R=E$,
	\[
		|R/(\alpha)|=N(\alpha).
	\]
	If $\pi$ is prime in $R$, then $R/(\pi)$ is a field of $N(\pi)$ elements.
	Its multiplicative group has order $N(\pi)-1$ and is cyclic.

\end{theorem}

\begin{proof}
	In $\mathbb{Z}[i]$, multiplication by $\alpha=a+bi$ sends the basis
	$1,i$ to the vectors $(a,b)$ and $(-b,a)$ in integer coordinates.
	The multiples of $\alpha$ are exactly their integer combinations.
	Theorem~\ref{thm:8.8.5} therefore gives $a^2+b^2=N(\alpha)$ classes.
	In $E$, multiplication by $\alpha=a+b\omega$ sends $1,\omega$ to
	$(a,b)$ and $(-b,a-b)$. The same theorem gives
	$a(a-b)+b^2=a^2-ab+b^2=N(\alpha)$ classes.

	Let $[z]\neq[0]$ modulo a prime $\pi$, so $\pi\nmid z$.
	A greatest common divisor of $\pi$ and $z$ is a unit, since primes are
	irreducible. B{\'e}zout's identity gives $u\pi+vz=1$, which shows that
	$[v][z]=[1]$. For the Gaussian ring these properties follow from
	Theorem~\ref{thm:8.2.2}.
	For $E$ they follow from Theorem~\ref{thm:8.8.2}.
	Thus every nonzero class is invertible. The class of $1$ is nonzero
	because $\pi$ is a nonunit, so this is a field. There are $N(\pi)-1$
	nonzero classes, and their group is cyclic by Theorem~\ref{thm:6.10.1}.

\end{proof}

\begin{theorem}[Fermat congruence and cube roots]\label{thm:8.8.7}
	For a prime $\pi\in E$ and $\pi\nmid\alpha$,
	\[
		\alpha^{N(\pi)-1}\equiv1\pmod\pi.
	\]
	If $\pi\nmid3$, then $3\mid N(\pi)-1$ and the three distinct cube roots
	of $1$ modulo $\pi$ are $1,\omega,\omega^2$. Hence
	$\alpha^{(N(\pi)-1)/3}$ is congruent to exactly one of them.

\end{theorem}

\begin{proof}
	Put $q=N(\pi)$ and work in the field $E/(\pi)$ from
	Theorem~\ref{thm:8.8.6}. Multiplication by a nonzero class $[\alpha]$
	permutes its $q-1$ nonzero classes. Multiplying all these classes and
	cancelling their nonzero product proves $[\alpha]^{q-1}=1$.

	The differences between $1,\omega,\omega^2$ are unit multiples of
	$1-\omega$: use $1+\omega=-\omega^2$ and
	$\omega^2-\omega=\omega(\omega-1)$. Since
	$3=-\omega^2(1-\omega)^2$, the condition $\pi\nmid3$ ensures that
	none of these differences vanishes modulo $\pi$. Thus $[\omega]$ has
	order $3$, and Lagrange's theorem (Theorem~\ref{thm:6.6.2}) gives
	$3\mid q-1$.
	The factorization
	\[
		X^3-1=(X-1)(X-\omega)(X-\omega^2)
	\]
	over this field shows that the three listed classes are all the cube roots
	of $1$. The Fermat congruence shows that
	$[\alpha]^{(q-1)/3}$ is one of them.

\end{proof}

\begin{definition}[cubic residue character]\label{def:8.8.3}
	Let $\pi$ be a prime of $E$ with $\pi\nmid3$, and put $q=N(\pi)$.
	For $\alpha\in E$ with $\pi\nmid\alpha$, define
	$\chi_\pi(\alpha)\in\{1,\omega,\omega^2\}$ by
	\[
		\alpha^{(q-1)/3}\equiv\chi_\pi(\alpha)\pmod\pi.
	\]
	The value is unique by Theorem~\ref{thm:8.8.7}.
	Set $\chi_\pi(\alpha)=0$ when $\pi\mid\alpha$.
	A nonzero class $[\alpha]\in E/(\pi)$ is a cubic residue if
	$[x]^3=[\alpha]$ for some nonzero class $[x]$.

\end{definition}

\begin{theorem}[cubic character criterion]\label{thm:8.8.8}
	The character $\chi_\pi$ depends only on the class modulo $\pi$ and satisfies
	\[
		\chi_\pi(\alpha\beta)=\chi_\pi(\alpha)\chi_\pi(\beta)
		\qquad(\alpha,\beta\in E).
	\]
	For $\pi\nmid\alpha$, the class of $\alpha$ is a cubic residue if and
	only if $\chi_\pi(\alpha)=1$. There are $(q-1)/3$ nonzero cubic residues,
	and each has exactly three cube roots modulo $\pi$.

\end{theorem}

\begin{proof}
	The defining power depends only on the class. When neither factor is
	zero modulo $\pi$, multiplication of the defining congruences gives the
	multiplicative rule. The three possible values multiply within
	$\{1,\omega,\omega^2\}$, and their distinct classes determine the product
	uniquely. If a factor is zero, the product is zero in the field, and both
	sides of the rule are zero.

	Choose a generator $g$ of the cyclic group $(E/(\pi))^\times$, which
	has order $q-1$. Write $[\alpha]=g^j$.
	The equation $g^{3k}=g^j$ is solvable exactly when $3\mid j$, with
	three solutions for $k$ modulo $q-1$. Also
	\[
		(g^j)^{(q-1)/3}=1\quad\Longleftrightarrow\quad3\mid j.
	\]
	This proves the criterion. Exactly $(q-1)/3$ exponents $j$ modulo $q-1$
	are divisible by $3$, which gives the count and the three-root statement.

\end{proof}

\newpage

\section{Prime Numbers: Primality Tests, Factorization, and Distribution}

\subsection{Quadratic-Residue and Order-Based Primality Criteria}

\begin{theorem}[Miller--Rabin criterion for $n\equiv3\pmod{4}$]\label{thm:9.1.1} Let $n>1$ be an integer with $n\equiv3\pmod{4}$. If $n$ is prime,
	then $a^{(n-1)/2}\equiv\pm1\pmod{n}$ for every
	$a\in\ZmodUnits{n}$. Conversely, if $n\equiv3\pmod{4}$ is
	composite, there exists $a\in\ZmodUnits{n}$ with
	$a^{(n-1)/2}\not\equiv\pm1\pmod{n}$.

\end{theorem}

\begin{proof} If $n$ is prime, then
	$\big(a^{(n-1)/2}\big)^{2}\equiv a^{n-1}\equiv1\pmod{n}$, and modulo a
	prime, $x^{2}\equiv1$ has only the roots $\pm1$. Now suppose that $n\equiv3\pmod{4}$ is
	composite. If $n=p^{k}$ with $p$ an odd prime and $k\ge2$, take $a=1+p$. Since
	$v_p(a-1)=v_p(p)=1$, the lifting-the-exponent lemma (Theorem~\ref{thm:4.8.3}) gives
	$v_p(a^{m}-1)=1+v_p(m)$ for every $m\in\mathbb{Z}^+$, so the order
	of $a$ modulo $p^{k}$ is $p^{k-1}$ (the least $m$ with $v_p(m)\ge
		k-1$). Since $p^{k}-1\equiv-1\pmod{p}$, we have $p\nmid(p^{k}-1)$, so
	$p^{k-1}\nmid(n-1)$ for $k\ge2$, meaning
	$a^{n-1}\not\equiv1\pmod{n}$; consequently $a^{(n-1)/2}$ cannot
	be $\pm1$ modulo $n$, since squaring either value would force
	$a^{n-1}\equiv1$. If $n=p^{k}n'$ with $\gcd(p, n')=1$ and $n'>1$, choose
	(by the Chinese remainder theorem, Theorem~\ref{thm:3.4.1})
	$a$ with $a\equiv-1\pmod{p^{k}}$ and $a\equiv1\pmod{n'}$; since
	$(n-1)/2$ is odd, $a^{(n-1)/2}\equiv-1\pmod{p^{k}}$ and
	$a^{(n-1)/2}\equiv1\pmod{n'}$; this is not congruent to $1\pmod{n}$ (as it is
	congruent to $-1\not\equiv1\pmod{p^{k}}$, since $p^k$ is odd) and not
	congruent to $-1\pmod{n}$ (as it is congruent to $1\not\equiv-1\pmod{n'}$, since
	$n'$ is odd and $n'>1$), so it is not $\pm1$ modulo $n$.

\end{proof}
\begin{theorem}[P{\'e}pin's test]\label{thm:9.1.2} Let
	$\mathcal{F}_k=2^{2^{k}}+1$ be the $k$-th Fermat number ($k\ge1$). Then $\mathcal{F}_k$
	is prime if and only if $3^{(\mathcal{F}_k-1)/2}\equiv-1\pmod{\mathcal{F}_k}$.

\end{theorem}

\begin{proof} If the congruence holds, then
	$3^{\mathcal{F}_k-1}\equiv1$ while $3^{(\mathcal{F}_k-1)/2}\not\equiv1\pmod{\mathcal{F}_k}$; since
	$\mathcal{F}_k-1=2^{2^{k}}$, Theorem~\ref{thm:6.1.2} (with the prime $2$ and $t=2^{k}$) shows that the order of $3$ modulo $\mathcal{F}_k$ is exactly $\mathcal{F}_k-1$. By Lagrange's theorem for finite groups (Theorem~\ref{thm:6.6.2}) this order divides $\varphi(\mathcal{F}_k)$, so $\mathcal{F}_k-1\le\varphi(\mathcal{F}_k)\le \mathcal{F}_k-1$, giving $\varphi(\mathcal{F}_k)=\mathcal{F}_k-1$. Since $\varphi(n)=n-1$ forces $n$ to be prime (a composite $n$ has a divisor $p$ with $1<p<n$, which is not coprime to $n$, so $\varphi(n)\le n-2$), the number $\mathcal{F}_k$ is prime. Conversely, if $\mathcal{F}_k$ is prime, then
	$\mathcal{F}_k\equiv1\pmod{4}$, so by reciprocity
	$\leg{3}{\mathcal{F}_k}=\leg{\mathcal{F}_k}{3}$; since
	$2^{2^{k}}\equiv1\pmod{3}$, we have $\mathcal{F}_k\equiv2\pmod{3}$, so
	$\leg{\mathcal{F}_k}{3}=\leg{2}{3}=-1$, and by
	Euler's criterion (Theorem~\ref{thm:7.1.2}) $3^{(\mathcal{F}_k-1)/2}\equiv-1\pmod{\mathcal{F}_k}$.

\end{proof}
\begin{definition}[Euler--Jacobi congruence, pseudoprimes, and witnesses]\label{def:9.1.1}
	Let $N>1$ be odd and suppose that $\gcd(b,N)=1$ (such a $b$ is called an
	admissible base for $N$). The Euler--Jacobi congruence for
	the base $b$ is
	\[
		b^{(N-1)/2}\equiv\leg{b}{N}\pmod{N},
	\]
	where $\leg{b}{N}$ is the Jacobi symbol.

	An odd composite integer $N$ is an Euler--Jacobi pseudoprime to an
	admissible base $b$ if the Euler--Jacobi congruence holds.

	For an odd composite $N>1$, an admissible base $b$ for which the
	Euler--Jacobi congruence fails is a Solovay--Strassen witness for the
	compositeness of $N$.

\end{definition}
\begin{theorem}\label{thm:9.1.3}
	If $N$ is an odd prime and $\gcd(b,N)=1$, then
	\[
		b^{(N-1)/2}\equiv\leg{b}{N}\pmod{N}.
	\]

\end{theorem}

\begin{proof} For prime $N$, the Jacobi symbol $\leg{b}{N}$ is the
	Legendre symbol. Euler's criterion (Theorem~\ref{thm:7.1.2}) gives the
	displayed congruence.

\end{proof}
\begin{theorem}\label{thm:9.1.4} Let $N>1$ be an odd composite number that
	is not a Carmichael number (Definition~\ref{def:6.5.1}). Then
	\[
		\#\{\,b\in\ZmodUnits{N} \mid b^{N-1}\equiv 1 \pmod{N}\,\}
		\le \tfrac12\,\varphi(N).
	\]
	Equivalently, at least half of the bases $b$ coprime to $N$ satisfy
	$b^{N-1}\not\equiv1\pmod{N}$, and so certify that $N$ is composite.

\end{theorem}

\begin{proof} Let $G=\ZmodUnits{N}$ and put
	$F=\{b\in G \mid b^{N-1}=1\}$. If $b,c\in F$, then
	$(bc^{-1})^{N-1}=b^{N-1}(c^{N-1})^{-1}=1$, so $F$ is a subgroup of $G$. Since
	$N$ is composite and not a Carmichael number, by the definition of a
	Carmichael number there exists $b_0\in G$ with $b_0^{N-1}\neq 1$, that is,
	$F\neq G$, so $F$ is a proper subgroup. By Lagrange's theorem for finite groups (Theorem~\ref{thm:6.6.2}), $[G:F]\ge 2$, hence
	$|F|\le\frac12|G|=\frac12\varphi(N)$.

\end{proof}
\begin{theorem}\label{thm:9.1.5} Every odd composite $N>1$ has a Solovay--Strassen witness
	(Definition~\ref{def:9.1.1}), that is, an admissible base $b$ for which
	$b^{(N-1)/2}\not\equiv\leg{b}{N}\pmod{N}$.

\end{theorem}

\begin{proof} If $N$ is squarefree, write $N=pm$ with $p$
	prime. Since $N$ is composite, $m>1$, and squarefreeness gives $\gcd(p,
		m)=1$; moreover $m$ is odd, hence $m\ge3$ and $1\not\equiv-1\pmod{m}$.
	Choose $u$ with $\leg{u}{p}=-1$ (half of $\ZmodUnits{p}$ consists
	of non-residues, Theorem~\ref{thm:7.2.2}) and, by the Chinese remainder theorem, an
	integer $b$ with $b\equiv u\pmod{p}$ and $b\equiv1\pmod{m}$. Then
	$\gcd(b, p)=\gcd(u, p)=1$ and $\gcd(b, m)=1$, so $\gcd(b, N)=1$ and the
	Jacobi symbol is defined:
	$\leg{b}{N}=\leg{b}{p}\leg{b}{m}=(-1)\cdot1=-1$. But
	$b^{(N-1)/2}\equiv1\pmod{m}$ while $1\not\equiv-1\pmod{m}$, which rules
	out $b^{(N-1)/2}\equiv-1\pmod{N}$. If $N$ is not squarefree, fix a prime $p$ with $p^{2}\mid N$ and write $N=p^{k}s$ with $k\ge2$ and $\gcd(p, s)=1$; by the Chinese remainder theorem choose $b$ with $b\equiv1+p\pmod{p^{2}}$ and $b\equiv1\pmod{s}$, so that $p\nmid b$ (as $b\equiv1\pmod{p}$) and $\gcd(b, s)=1$, whence $\gcd(b, N)=1$. Then $b^{N-1}\equiv(1+p)^{N-1}\equiv1+(N-1)p\pmod{p^{2}}$, and since $p\mid N$ forces $p\nmid N-1$, this is not congruent to $1\pmod{p^{2}}$; as $p^{2}\mid N$ it follows that $b^{N-1}\not\equiv1\pmod{N}$. Hence $b^{(N-1)/2}\not\equiv\pm1\pmod{N}$, so it cannot equal $\leg{b}{N}\in\{\pm1\}$ and the Euler--Jacobi congruence fails.

\end{proof}
\begin{theorem}[Proth's theorem]\label{thm:9.1.6} Let $n=2^{m}h+1$
	with $m,h\in\mathbb{Z}^+$, $h<2^{m}$, and $h$ odd. Then $n$ is prime if and only if
	there exists an integer $a$ with $a^{(n-1)/2}\equiv-1\pmod{n}$.

\end{theorem}

\begin{proof} For necessity, suppose that $n$ is prime. It then has
	a quadratic non-residue $a$ (exactly half the elements of $\ZmodUnits{n}$ are non-residues),
	and Euler's criterion gives $a^{(n-1)/2}\equiv\leg{a}{n}=-1\pmod{n}$,
	so a suitable $a$ exists. Conversely, suppose
	$a^{(n-1)/2}\equiv-1\pmod{n}$ for some $a$; then
	$\gcd(a, n)=1$. Let $P$ be a prime factor of $n$ and $d$ the order of
	$a$ modulo $P$. From $a^{(n-1)/2}\equiv-1$ and $a^{n-1}\equiv1\pmod{P}$ we get $d\nmid\tfrac{n-1}{2}$, $d\mid n-1$, and $d\mid P-1$; since
	$n-1=2^{m}h$ with $h$ odd, $d\mid n-1$ gives $v_2(d)\le m$ while
	$d\nmid 2^{m-1}h$ gives $v_2(d)\ge m$, so $v_2(d)=m$, that is, $2^{m}\mid d$, and $P\equiv1\pmod{2^{m}}$.
	Since $P$ was an arbitrary prime factor, every prime factor of $n$ is
	congruent to $1\pmod{2^m}$. Suppose $P<n$; multiplying the remaining prime factors
	with multiplicity gives $n/P\equiv1\pmod{2^m}$ and $n/P>1$, so we may write
	$n=(2^{m}x+1)(2^{m}y+1)$ with $x,y\ge1$; then
	$h=2^{m}xy+x+y\ge2^{m}$, contradicting $h<2^{m}$. Hence $P=n$, so $n$
	is prime.

\end{proof}
\begin{theorem}\label{thm:9.1.7} Let
	$n>1$ be odd and composite. The set of Euler liars
	$E(n)=\{a\in\ZmodUnits{n} \mid a^{(n-1)/2}\equiv\leg{a}{n}\pmod{n}\}$ (where $\leg{a}{n}$ is the Jacobi
	symbol) is a subgroup of $\ZmodUnits{n}$ of index at least
	$2$; hence at most half of the admissible bases are Euler liars, that is, $|E(n)|\le\varphi(n)/2$.

\end{theorem}

\begin{proof} For $a,b\in\ZmodUnits{n}$, the
	Jacobi symbol is multiplicative,
	$\leg{ab}{n}=\leg{a}{n}\leg{b}{n}$, with values in $\{\pm1\}$. If $a,b\in E(n)$, then
	$(ab)^{(n-1)/2}=a^{(n-1)/2}b^{(n-1)/2}\equiv\leg{a}{n}\leg{b}{n}=\leg{ab}{n}\pmod{n}$, so
	$ab\in E(n)$; and $1\in E(n)$. Being a nonempty subset of the finite group $\ZmodUnits{n}$ closed under multiplication, $E(n)$ is a subgroup: for $a\in E(n)$, two powers coincide, say $a^r=a^s$ with $r>s\ge0$; multiplying in the ambient unit group by $a^{-s}$ gives $a^{r-s}=1$. Thus $a^{-1}=a^{r-s-1}\in E(n)$. Hence $E(n)\le\ZmodUnits{n}$. By Theorem~\ref{thm:9.1.5} there is a base
	$b\in\ZmodUnits{n}$ with $b^{(n-1)/2}\not\equiv\leg{b}{n}\pmod{n}$, that is, $b\notin E(n)$; hence $E(n)$ is a proper
	subgroup. By Lagrange's theorem for finite groups (Theorem~\ref{thm:6.6.2}) its index is at least $2$, so
	$|E(n)|\le\varphi(n)/2$.

\end{proof}
\begin{theorem}[Lucas primality test]\label{thm:9.1.8} Let $n>1$
	and suppose there is an integer $a$ with $a^{n-1}\equiv1\pmod{n}$ and
	$a^{(n-1)/q}\not\equiv1\pmod{n}$ for every prime $q\mid n-1$. Then $n$ is prime.

\end{theorem}

\begin{proof} Since $a^{n-1}\equiv1\pmod{n}$,
	we have $\gcd(a, n)=1$, so $a\in\ZmodUnits{n}$; let
	$d=\delta_n(a)$. Then $d\mid n-1$. If $d<n-1$, choose a
	prime $q\mid\tfrac{n-1}{d}$; then $d\mid\tfrac{n-1}{q}$, whence
	$a^{(n-1)/q}\equiv1\pmod{n}$, contradicting the hypothesis. Hence
	$d=n-1$. As $d\mid\varphi(n)$, we have $\varphi(n)\ge n-1$; but $\varphi(n)\le n-1$ for every $n>1$, with equality if and only if $n$ is prime (a composite $n$ has a divisor $p$ with $1<p<n$ that is not coprime to $n$, so $\varphi(n)\le n-2$). Therefore $\varphi(n)=n-1$ and $n$ is prime.

\end{proof}
\begin{theorem}[Pocklington's criterion]\label{thm:9.1.9} Let
	$n>1$ be an integer, and write $n-1=q^{k}r$ with $q$ prime,
	$k,r\in\mathbb{Z}^+$, and $q\nmid r$. Suppose there is an
	integer $a$ with $a^{n-1}\equiv1\pmod{n}$ and
	$\gcd\left(a^{(n-1)/q}-1,\,n\right)=1$. Then every prime factor $p$
	of $n$ satisfies $p\equiv1\pmod{q^{k}}$.
	In particular, if $q^k+1>\sqrt n$, then $n$ is prime.

\end{theorem}

\begin{proof} Fix a prime $p\mid n$. From
	$a^{n-1}\equiv1\pmod{n}$ we get $\gcd(a, n)=1$, so $\gcd(a, p)=1$; let
	$d=\delta_p(a)$. Reducing $a^{n-1}\equiv1$ modulo $p$
	gives $d\mid n-1=q^{k}r$. If $d\mid\tfrac{n-1}{q}=q^{k-1}r$, then
	$a^{(n-1)/q}\equiv1\pmod{p}$, so $p\mid a^{(n-1)/q}-1$; with $p\mid n$
	this contradicts $\gcd(a^{(n-1)/q}-1, n)=1$. Hence $d\nmid q^{k-1}r$
	while $d\mid q^{k}r$, which forces $q^{k}\mid d$. By Fermat's little
	theorem $d\mid p-1$, so $q^{k}\mid p-1$, that is, $p\equiv1\pmod{q^{k}}$.

	If $q^k+1>\sqrt n$, every prime divisor $p$ of $n$ satisfies
	$p\ge q^k+1>\sqrt n$. If $n$ were composite, we could write $n=uv$
	with $1<u\le v$. Then $u\le\sqrt n$, and a prime divisor
	$p$ of $u$ would satisfy $p\mid n$ and $p\le u\le\sqrt n$,
	a contradiction.

\end{proof}
\begin{altproof}[Alternative proof of sufficiency in Theorem~\ref{thm:9.1.6}] Assume $a^{(n-1)/2}\equiv-1\pmod{n}$ and write $n-1=2^{m}h$
	with $h$ odd. In Theorem~\ref{thm:9.1.9}, take $q=2$, $k=m$, and
	$r=h$. Squaring
	gives $a^{n-1}\equiv1\pmod{n}$, and $a^{(n-1)/2}-1\equiv-2\pmod{n}$ with
	$n$ odd, so $\gcd\big(a^{(n-1)/2}-1,\,n\big)=\gcd(-2,n)=1$. Theorem~\ref{thm:9.1.9} therefore gives $P\equiv1\pmod{2^{m}}$ for every prime factor $P$
	of $n$. If $n$ were composite, then grouping its prime factors into two
	nontrivial parts, each congruent to $1\pmod{2^{m}}$, would give
	$n=(2^{m}x+1)(2^{m}y+1)$ with $x,y\ge1$, whence
	$h=2^{m}xy+x+y\ge2^{m}$, contradicting $h<2^{m}$. Hence $n$ is prime.

\end{altproof}
\subsection{The Quadratic Character of Factors of Mersenne Numbers}

\begin{theorem}\label{thm:9.2.1} Let $p$ be an odd prime and $q$ a
	prime divisor of the Mersenne number $M_p = 2^p - 1$. Then:
	\begin{enumerate}[label={\alph*)}]
		\item $q\equiv 1 \pmod{2p}$;
		\item $q \equiv \pm 1 \pmod{8}$.
	\end{enumerate}

\end{theorem}

\begin{proof} Since $M_p = 2^{p} - 1$ is odd, its prime
	divisor $q$ is odd. From $q \mid 2^p - 1$, we obtain $2^p \equiv 1 \pmod{q}$, so $\delta_q(2) \mid p$; since $p$ is prime, and since
	$\delta_q(2) = 1$ would give $q \mid 1$, we have
	$\delta_q(2) = p$. By Fermat's little theorem, $p \mid q - 1$, and since
	$q, p$ are odd, $2p \mid q - 1$, that is, $q \equiv 1 \pmod{2p}$. For
	(b), note that $p+1$ is even and $2 \equiv 2^{p+1} = \big(2^{(p+1)/2}\big)^2
		\pmod{q}$, so $\leg{2}{q} = 1$, giving $q \equiv \pm 1
		\pmod{8}$ by the second supplementary law (Theorem~\ref{thm:7.2.5}(f)).

\end{proof}
\begin{theorem}\label{thm:9.2.2} Let $p$ be a prime with $p \equiv 3 \pmod{4}$. If $q = 2p
		+ 1$ is also prime, then $q \mid M_p = 2^{p} - 1$.

\end{theorem}

\begin{proof} Since $p \equiv 3 \pmod{4}$, write $p = 4k +
		3$; then $q = 2p + 1 = 8k + 7 \equiv -1 \pmod{8}$, so by the second
	supplementary law (Theorem~\ref{thm:7.2.5}(f)) $\leg{2}{q} = 1$, that is, $2$ is a
	quadratic residue modulo $q$. By Euler's criterion (Theorem~\ref{thm:7.1.2}),
	$2^{(q-1)/2} \equiv \leg{2}{q} \equiv 1 \pmod{q}$; since
	$(q-1)/2 = p$, this gives $2^{p} \equiv 1 \pmod{q}$, that is, $q \mid
		2^{p} - 1 = M_p$.

\end{proof}
\begin{altproof}[Alternative proof of Theorem~\ref{thm:9.2.2}]
	Let $q=2p+1$ be prime; then $q\ge7$, and $d=\delta_q(2)$ divides
	$q-1=2p$, so $d\in\{1,2,p,2p\}$. The values $d=1$ and $d=2$ are
	impossible, since they would give $q\mid1$ or $q\mid3$ while $q>3$. If
	$d=2p=q-1$, then $2^{(q-1)/2}\not\equiv1\pmod{q}$, so
	$2^{(q-1)/2}\equiv-1$ and Euler's criterion (Theorem~\ref{thm:7.1.2}) shows that $2$ is
	a quadratic non-residue modulo $q$; but $p\equiv3\pmod{4}$ gives
	$q=2p+1\equiv7\pmod{8}$, whence $\leg{2}{q}=1$ by Theorem~\ref{thm:7.2.5}(f),
	a contradiction. Hence $d\mid p$, so $2^{p}\equiv1\pmod{q}$
	and $q\mid M_p$.

\end{altproof}
\subsection{The Lucas--Lehmer Test}

\begin{definition}[Lucas--Lehmer sequence]\label{def:9.3.1} Define
	$(S_k)_{k\ge0}$ by $S_0=4$ and $S_{k+1}=S_k^2-2$.

\end{definition}
\begin{theorem}\label{thm:9.3.1} In
	$\mathbb{Z}[\sqrt{3}]$, set $\omega=2+\sqrt{3}$ and $\bar\omega=2-\sqrt{3}$,
	so that $\omega\bar\omega=1$. Then
	$S_k=\omega^{2^k}+\bar\omega^{2^k}$ for every $k\ge0$.

\end{theorem}

\begin{proof} We argue by induction on $k$. For $k=0$,
	we have $\omega+\bar\omega=4=S_0$. Assuming the formula for $k$, we obtain
	$S_{k+1}=S_k^2-2=\big(\omega^{2^k}+\bar\omega^{2^k}\big)^2-2=\omega^{2^{k+1}}+\bar\omega^{2^{k+1}}+2(\omega\bar\omega)^{2^k}-2=\omega^{2^{k+1}}+\bar\omega^{2^{k+1}}$.

\end{proof}
\begin{theorem}[Lucas--Lehmer test]\label{thm:9.3.2} Let $p$ be an
	odd prime and $M_p=2^p-1$. Then $M_p$ is prime if and only if $M_p\mid S_{p-2}$.

\end{theorem}

\begin{proof} First suppose that $M\mid S_{p-2}$, where $M=M_p$.
	By Theorem~\ref{thm:9.3.1},
	$\omega^{2^{p-2}}+\bar\omega^{2^{p-2}}\equiv0\pmod{M}$ in
	$\mathbb{Z}[\sqrt{3}]$. Multiplying by $\omega^{2^{p-2}}$ and using
	$\omega\bar\omega=1$ gives $\omega^{2^{p-1}}\equiv-1\pmod{M}$, and
	squaring gives $\omega^{2^{p}}\equiv1\pmod{M}$. Suppose $M$ were
	composite and let $q$ be its least prime factor, so $q\le\sqrt{M}$.
	Reducing coefficients modulo $q$ identifies
	$R=\mathbb{Z}[\sqrt{3}]/(q)$ with the ring $R_q$ of pairs $(a,b)$ of
	residues modulo $q$ under the rule $(a,b)(c,e)=(ac+3be,ae+bc)$; these
	operations are given by integer polynomial identities and so persist modulo $q$
	(Definition~\ref{def:7.9.1} with $d=3$); thus $R_q$ is a finite commutative
	ring with identity $(1,0)$ and exactly $q^{2}$ elements.
	Moreover, $\omega=(2,1)$ and $\bar\omega=(2,-1)$ satisfy
	$\omega\bar\omega=(1,0)$, so the image of $\omega$ in $R$ is a unit with inverse $\bar\omega$. This
	unit satisfies $\omega^{2^{p}}=1$ but $\omega^{2^{p-1}}=-1\neq1$ (as
	$q$ is odd); hence, by Theorem~\ref{thm:6.1.2} applied in the group $R^{\times}$ with
	the prime $2$ and $t=p$, the element $\omega$ has order exactly $2^{p}$ in
	$R^{\times}$.
	Therefore $2^{p}\le|R^{\times}|\le q^2-1\le M-1=2^{p}-2$, a contradiction. So $M$ is prime.

	Conversely, suppose that $M$ is prime. Since
	$p\ge3$ is odd, we have $M\equiv7\pmod{8}$, $M\equiv1\pmod{3}$, and
	$M\equiv3\pmod{4}$; the supplementary law (Theorem~\ref{thm:7.2.5}(f)) and
	quadratic reciprocity (Theorem~\ref{thm:7.5.1})
	give $\leg{2}{M}=1$ and $\leg{3}{M}=-1$. Hence
	$3$ is a quadratic non-residue modulo $M$. By
	Theorems~\ref{thm:7.9.1} and~\ref{thm:7.9.2} with the prime $M$ in place of $p$ and with $d=3$,
	the ring $F=\mathbb{F}_M[\sqrt{3}]$ is a field with $M^2$ elements and the
	Frobenius automorphism is conjugation. In particular,
	$(1+\sqrt{3})^{M}=1-\sqrt{3}$ in $F$. Since
	$(1+\sqrt{3})^2=4+2\sqrt{3}=2\omega$,
	we have $\omega^{2^{p-1}}=\omega^{(M+1)/2}=\dfrac{(1+\sqrt{3})^{M+1}}{2^{(M+1)/2}}$;
	here $(1+\sqrt{3})^{M+1}\equiv(1-\sqrt{3})(1+\sqrt{3})=-2$ and
	$2^{(M+1)/2}=2\cdot2^{(M-1)/2}\equiv2\leg{2}{M}=2\pmod{M}$,
	so $\omega^{2^{p-1}}\equiv-1\pmod{M}$. Therefore
	$S_{p-2}=\omega^{2^{p-2}}+\omega^{-2^{p-2}}=\omega^{-2^{p-2}}\big(\omega^{2^{p-1}}+1\big)\equiv0\pmod{M}$, that is, $M\mid S_{p-2}$.

\end{proof}
\begin{altproof}[Alternative proof of sufficiency in Theorem~\ref{thm:9.3.2}]
	For any modulus $n\ge2$, the ring
	$R_n(3)=(\mathbb{Z}/n\mathbb{Z})[\sqrt{3}]$ of
	Definition~\ref{def:7.9.1} with $d=3$ is the set of pairs $(x,y)$ modulo
	$n$ with addition defined componentwise and multiplication defined by
	$(x+y\sqrt{3})(u+v\sqrt{3})=(xu+3yv)+(xv+yu)\sqrt{3}$; its operations
	are given by integer polynomial identities and so persist modulo any $n$, with
	identity $(1,0)$ and norm $N(x,y)=x^{2}-3y^{2}$ satisfying
	$N\big((x,y)\cdot(u,v)\big)\equiv N(x,y)N(u,v)\pmod{n}$. For an odd prime
	$q\neq3$, it is a field when $\leg{3}{q}=-1$
	(Theorem~\ref{thm:7.9.1}) and isomorphic to
	$\mathbb{F}_q\times\mathbb{F}_q$ when
	$\leg{3}{q}=1$ (Theorem~\ref{thm:7.9.5}).

	Put $\alpha=(2,1)$ and $\beta=(2,-1)$ in $R_n(3)$. Then
	$\alpha\beta=(4-3,\,-2+2)=(1,0)$, so $\alpha$ and $\beta$ are
	mutually inverse units, and $N(\alpha)=1$. Write
	$\alpha^{2^k}=(A_k,B_k)$. Making the substitutions $y\mapsto-y,\ v\mapsto-v$
	throughout the product formula only negates the second coordinate, so
	$\beta^{2^k}=(A_k,-B_k)$ and hence $\alpha^{2^k}+\beta^{2^k}=2A_k$.
	Squaring both sides gives $(A_{k+1},B_{k+1})=(A_k^2+3B_k^2,\ 2A_kB_k)$, and
	since $N(\alpha^{2^k})=N(\alpha)^{2^k}=1$ (that is, $A_k^2-3B_k^2=1$),
	we get $2A_{k+1}=2A_k^2+6B_k^2=2A_k^2+2(A_k^2-1)=(2A_k)^2-2$. As $A_0=2$
	gives $2A_0=4=S_0$, and $2A_{k+1}=(2A_k)^2-2$ is exactly the
	recursion defining $S_k$, induction gives
	$S_k\equiv2A_k\equiv\alpha^{2^k}+\beta^{2^k}\pmod{M_p}$ (working in
	$R_{M_p}(3)$) for every $k\ge0$.

	Assume $S_{p-2}\equiv0\pmod{M_p}$ and suppose, for a contradiction,
	that $M_p$ is composite; then $M_p$ has a prime factor
	$q\le\sqrt{M_p}$. Reduction modulo $q$ carries the ring $R_{M_p}(3)$ onto
	$R_q(3)$ (further reducing each coordinate modulo $q$), so
	$\alpha^{2^{p-2}}\equiv-\beta^{2^{p-2}}\pmod{q}$. Since
	$\beta=\alpha^{-1}$, multiplying both sides by $\alpha^{2^{p-2}}$
	gives $\alpha^{2^{p-1}}\equiv-1\pmod{q}$ (that is, $\alpha^{2^{p-1}}=(-1,0)$ in $R_q(3)$); squaring gives
	$\alpha^{2^{p}}\equiv1\pmod{q}$. As $q$ is odd, $(-1,0)\neq(1,0)$ in
	$R_q(3)$, so $\alpha^{2^{p-1}}\not\equiv1$; hence Theorem~\ref{thm:6.1.2}, applied in the group
	of units of $R_q(3)$ with the prime $2$ and $t=p$ (every divisor of $2^p$
	other than $2^p$ itself divides $2^{p-1}$), shows that the multiplicative
	order of $\alpha$ in $R_q(3)$ is exactly $2^p$.

	The units of $R_q(3)$ form a group under multiplication, and this group has at most
	$q^2-1$ elements because $R_q(3)$ has exactly $q^2$ elements in total
	and the zero element $(0,0)$ is never a unit. By Lagrange's theorem for finite groups (Theorem~\ref{thm:6.6.2})
	the order of $\alpha$ divides the order of this group, so $2^p\le
		q^2-1$. But $q\le\sqrt{M_p}=\sqrt{2^p-1}<\sqrt{2^p}$ gives $q^2<2^p$,
	so $q^2-1<2^p-1<2^p$, contradicting $2^p\le q^2-1$.

	Hence $M_p$ has no prime factor at most $\sqrt{M_p}$, so $M_p$ is prime.

\end{altproof}
\subsection{The Strong Probable-Prime Test (Miller--Rabin)}

\begin{definition}[strong probable prime]\label{def:9.4.1} Let
	$n>1$ be odd, and write $n-1=2^{s}d$ with $d$ odd and $s\ge1$. For $a$ with
	$\gcd(a, n)=1$, we say that $n$ passes the strong test to base $a$ if
	$a^{d}\equiv1\pmod{n}$, or $a^{2^{r}d}\equiv-1\pmod{n}$ for some $0\le r<s$.

\end{definition}
\begin{theorem}\label{thm:9.4.1} If $n$ is
	an odd prime, then $n$ passes the strong test to every base $a$ with
	$\gcd(a, n)=1$.

\end{theorem}

\begin{proof} By Fermat's little theorem
	$a^{n-1}=a^{2^{s}d}\equiv1\pmod{n}$. Let $k$ be minimal with $0\le
		k\le s$ such that $a^{2^{k}d}\equiv1\pmod{n}$ (it exists since $k=s$
	works). If $k=0$, then $a^{d}\equiv1\pmod{n}$ and $n$ passes. If
	$k\ge1$, then $\big(a^{2^{k-1}d}\big)^2\equiv1\pmod{n}$; since $n$ is
	prime, $n\mid(x-1)(x+1)$ implies $x\equiv\pm1\pmod{n}$; applying this
	with $x=a^{2^{k-1}d}$ gives
	$a^{2^{k-1}d}\equiv\pm1\pmod{n}$. It cannot be congruent to $1$, as that
	would contradict the minimality of $k$. Hence
	$a^{2^{k-1}d}\equiv-1\pmod{n}$, so $n$ passes with $r=k-1$.

\end{proof}
\begin{theorem}\label{thm:9.4.2} For every odd composite $n>1$, there exists $a$ with
	$\gcd(a, n)=1$ such that $n$ fails the strong test to base $a$.

\end{theorem}

\begin{proof} Write $n-1=2^{s}d$ with $d$ odd and $s\ge1$.

	First suppose that $n=p^{k}$ is a prime power with $k\ge2$; the prime
	$p$ is odd because $n$ is odd. Take $a=1+p$, so $\gcd(a, p)=1$ and
	hence $\gcd(a, n)=1$. Since
	$v_p(a-1)=v_p(p)=1$, the lifting-the-exponent lemma (Theorem~\ref{thm:4.8.3}) gives
	$v_p(a^{t}-1)=1+v_p(t)$ for every $t\ge1$; in particular
	$a^t\equiv1\pmod{p^k}$ if and only if $p^{k-1}\mid t$, so $a$ has order exactly
	$p^{k-1}$ modulo $p^k$. Since $p^k-1\equiv-1\pmod{p}$, we have $p\nmid(n-1)$,
	so $p\nmid d$ and (as $k\ge2$) $p^{k-1}\nmid d$; hence
	$a^{d}\not\equiv1\pmod{n}$. For the second condition, observe that every
	$a^{2^rd}$ lies in the cyclic group $\langle a\rangle$ of odd order
	$p^{k-1}$, so its order is odd; but $-1$ has order exactly $2$ in
	$\ZmodUnits{p^k}$ (as $p^k\ge p^2\ge9>2$). An
	element of odd order can never equal an element of order $2$, so
	$a^{2^{r}d}\not\equiv-1\pmod{n}$ for any $r$. Hence $n$ fails the
	strong test to base $a$.

	Now suppose that $n=p^{k}m$ with $p$ an odd prime, $k\ge1$,
	$\gcd(p, m)=1$, and $m>1$; since $n$ is odd, $m$ is odd, so $m\ge3$.
	By the
	Chinese remainder theorem (Theorem~\ref{thm:3.4.1}) choose $a$ with
	$a\equiv-1\pmod{p^k}$ and $a\equiv1\pmod{m}$; then $\gcd(a, n)=1$. As
	$d$ is odd, $a^{d}\equiv(-1)^d=-1\pmod{p^{k}}$ and $a^{d}\equiv1\pmod{m}$. Since $p^{k}\ge3$, we have $-1\not\equiv1\pmod{p^k}$, so
	$a^{d}\not\equiv1\pmod{n}$; since $m\ge3$, likewise $1\not\equiv-1\pmod{m}$, so
	$a^{d}\not\equiv-1\pmod{n}$ either. For $1\le r<s$, we have
	$a^{2^{r}d}=(a^{d})^{2^{r}}\equiv(-1)^{2^r}=1\pmod{p^{k}}$ and
	$a^{2^{r}d}\equiv1\pmod{m}$, so $a^{2^{r}d}\equiv1\pmod{n}$; as $n\ge9$,
	we have $1\not\equiv-1\pmod{n}$, so $a^{2^{r}d}\not\equiv-1\pmod{n}$. Thus
	$a^{2^{r}d}\not\equiv-1\pmod{n}$ for every $0\le r<s$, and also
	$a^{d}\not\equiv1\pmod{n}$, so $n$ fails the strong test to base $a$.

	Every odd composite $n$ falls into one of these two cases, which completes the
	proof. Theorem~\ref{thm:9.1.1} is the special case $n\equiv3\pmod{4}$, where
	$s=1$ and the strong test reduces to the condition
	$a^{(n-1)/2}\equiv\pm1\pmod{n}$.

\end{proof}
\subsection{Factorization Methods: Fermat and Pollard \texorpdfstring{$p-1$}{p-1}}

\begin{theorem}\label{thm:9.5.1} Let $n>1$ be odd. For an integer $t \ge \lceil\sqrt
		{n}\rceil$, we call $t$ good if $t^2-n$ is a perfect square, say
	$t^2 - n = s^2$ with an integer $s \ge 0$. Every good $t$ yields a factorization
	$n = (t-s)(t+s)$ into positive odd integers. Conversely, every
	factorization $n = ab$ with $1 \le a \le b$ arises from the good pair
	$t = \frac{a+b}{2}$, $s = \frac{b-a}{2}$. Now let
	$t^{*}$ be the least good $t \ge \lceil\sqrt{n}\rceil$. Then $n$ is prime
	if and only if $t^{*} = \frac{n+1}{2}$; otherwise $n =
		(t^{*}-s^{*})(t^{*}+s^{*})$ is a nontrivial factorization, where
	$s^{*} = \sqrt{(t^{*})^2-n}$.

\end{theorem}

\begin{proof} First suppose that $t$ is good. If $t^2-n=s^2$, then
	$n=t^2-s^2=(t-s)(t+s)$. Since $t\ge\lceil\sqrt
		{n}\rceil>0$ and $s^2=t^2-n\le t^2$ with equality only if $n=0$, we get
	$0\le s<t$, so $t-s\ge1>0$. As $t-s$ and $t+s$ have the same parity
	(their sum $2t$ is even) and their product $n$ is odd, both must be odd.

	Conversely, suppose that $n=ab$ with $1\le a\le b$. Then
	$a,b$ are odd (divisors of odd $n$), so $t=\frac{a+b}{2}$ and
	$s=\frac{b-a}{2}$ are integers, $s\ge0$, and $t^2-s^2=ab=n$. By AM--GM,
	$t=\frac{a+b}{2}\ge\sqrt{ab}=\sqrt{n}$, so $t\ge\lceil\sqrt{n}\rceil$ (as
	$t\in\mathbb{Z}$); hence $t$ is good with $t^2-n=s^2$.

	It remains to prove the primality criterion. The two constructions above are mutually
	inverse (from a good $t$ with $t^2-n=s^2$ we recover the
	factorization $(t-s,t+s)$, and applying the second construction to it
	returns $t$), so good values of $t$ correspond bijectively to
	factorizations $n=ab$ with $1\le a\le b$. If $n$ is prime, its only
	such factorization is $a=1,b=n$, giving the unique good
	$t=\frac{n+1}{2}$; hence $t^*=\frac{n+1}{2}$. If $n$ is composite,
	write $n=ab$ with $1<a\le\sqrt{n}\le b<n$. The function $f(x)=x+\frac{n}{x}$ satisfies $f(x)-f(y)=(x-y)\big(1-\frac{n}{xy}\big)$, which has
	the sign of $-(x-y)$ whenever $xy<n$; taking $x=a\le\sqrt{n}$ and
	$y=1$ gives $xy=a\le\sqrt{n}<n$ (as $n>1$), so $f(a)<f(1)=1+n$, that is, $a+b<n+1$. Hence the good value $t=\frac{a+b}{2}$ corresponding to $(a,b)$
	satisfies $t<\frac{n+1}{2}$. Therefore $t^*\le t<\frac{n+1}{2}$, so
	$t^*\neq\frac{n+1}{2}$. By the bijection, $t^*$ corresponds to a
	factorization different from $(1,n)$ (which corresponds only to
	$t=\frac{n+1}{2}$), so $t^*-s^*>1$ and $n=(t^*-s^*)(t^*+s^*)$ is
	nontrivial.

\end{proof}
\begin{definition}[Pollard $p-1$ setup]\label{def:9.5.1} Fix $B\in\mathbb{Z}^+$ and let
	$M=\lcm(1, 2, \ldots, B)$. For an integer $n>1$ and $a\in\mathbb{Z}$
	with $\gcd(a, n)=1$, define $d = \gcd(a^M-1, \, n)$.

\end{definition}
\begin{theorem}\label{thm:9.5.2} In the setting of
	Definition~\ref{def:9.5.1}, let $p$ be a prime factor of $n$ and write
	$p-1=\prod_{i=1}^{r}\ell_i^{e_i}$, where the $\ell_i$ are distinct
	primes. Suppose every maximal prime-power factor
	$\ell_i^{e_i}\parallel p-1$ is at most $B$, so that $(p-1)\mid M$.
	Then $p \mid \gcd(a^M-1, n)$. If, moreover, $n \nmid
		a^M-1$, then $d=\gcd(a^M-1, n)$ is a nontrivial divisor of $n$ ($1<d<n$).

\end{theorem}

\begin{proof} Since $\gcd(a, n)=1$ and $p\mid n$,
	we have $\gcd(a, p)=1$, so by Fermat's little theorem (Theorem~\ref{thm:3.3.2})
	$a^{p-1}\equiv1\pmod{p}$. As $(p-1)\mid M$, write $M=(p-1)k$; then
	$a^M=(a^{p-1})^k\equiv1\pmod{p}$, that is, $p\mid a^M-1$. Since also
	$p\mid n$, the prime $p$ divides $\gcd(a^M-1, n)=d$, so $d>1$. If $n\nmid
		a^M-1$, let $x$ be the least nonnegative residue of $a^M-1$ modulo
	$n$; then $0<x<n$ and $\gcd(x, n)=\gcd(a^M-1, n)=d$ (the gcd with $n$ is
	unaffected by adding multiples of $n$), so $d\le x<n$. Together with
	$d>1$, this shows that $d$ is a nontrivial divisor of $n$.

\end{proof}
\begin{theorem}\label{thm:9.5.3}
	Let $N\in\mathbb{Z}^+$ with $N>1$, and let $x,y\in\mathbb{Z}$ satisfy
	$x^{2}\equiv y^{2}\pmod{N}$ with $x\not\equiv\pm y\pmod{N}$.
	Then $\gcd(N,x-y)$ and $\gcd(N,x+y)$ are both nontrivial divisors of $N$.

\end{theorem}

\begin{proof} Write $d=\gcd(N,x-y)$. Starting from $N\mid(x-y)(x+y)$, we argue in
	two steps. If $d=1$, then $N$ is coprime to $x-y$, so $N\mid x+y$,
	that is, $x\equiv-y\pmod{N}$, contrary to hypothesis. If $d=N$, then
	$N\mid x-y$, that is, $x\equiv y\pmod{N}$, again contrary to hypothesis. Hence
	$1<d<N$, so $d$ is a nontrivial divisor. The same argument with the roles of
	$x-y$ and $x+y$ interchanged shows that $1<\gcd(N,x+y)<N$.

\end{proof}
\subsection{Elementary Estimates}

\begin{definition}[prime-counting function and enumeration of primes]\label{def:9.6.1} For
	$x \ge 0$, let $\pi(x)$ denote the number of primes not exceeding $x$. We write
	$p_1=2<p_2=3<p_3=5<\cdots$ for the primes in increasing order.

\end{definition}
\begin{altproof}[Alternative proof of Theorem~\ref{thm:2.1.2}]
	Suppose there were only finitely many primes, say $p_1,\ldots,p_k$. Every
	positive integer $n$ factors uniquely as $n=ms^{2}$ with $m$ squarefree:
	take $s^{2}$ to be the largest square dividing $n$. A squarefree $m$ is a
	product of distinct primes from the list, so there are at most $2^{k}$
	possible values of $m$; and $n\le x$ forces $s\le\sqrt{x}$. Hence
	\[
		\lfloor x\rfloor=\#\{n\in\mathbb{Z}^+ \mid n\le x\}\le 2^{k}\sqrt{x}
	\]
	for every real $x\ge1$; but for $x>4^{k+1}$ we have
	$\lfloor x\rfloor>x-1\ge\tfrac{x}{2}>2^{k}\sqrt{x}$ (since
	$\sqrt{x}>2^{k+1}$), a contradiction.

\end{altproof}
\begin{theorem}\label{thm:9.6.1}
	For every real $x\ge2$,
	\[
		\pi(x)\ge\frac{\ln x}{2\ln 2}.
	\]

\end{theorem}

\begin{proof} Apply the counting argument of the preceding
	alternative proof without the hypothesis that the set of primes is finite.
	Every positive integer $n\le x$ factors as $n=ms^{2}$ with $m$ squarefree and
	$s\le\sqrt{x}$, and a squarefree $m$ with $m \le x$ is a product of distinct
	primes not exceeding $x$, so there are at most $2^{\pi(x)}$ choices of $m$.
	Hence $\lfloor x\rfloor\le 2^{\pi(x)}\sqrt{x}$. For an integer $n\ge1$ this
	reads $n\le 2^{\pi(n)}\sqrt{n}$, that is, $\sqrt{n}\le 2^{\pi(n)}$. Equality $n=4^{\pi(n)}$ is impossible for $n\ge2$.
	Indeed, it would force $\pi(n)\ge2$: otherwise $\pi(n)=1$
	would give $n=4$, but $\pi(4)=2$.
	Consequently, $n\ge16$, so $3$ is a prime not exceeding $n$.
	Also $\sqrt n=2^{\pi(n)}$ is an integer. There are exactly
	$2^{\pi(n)}\sqrt n=n$ candidate pairs $(m,s)$, where $m$ is
	a product of a subset of the primes not exceeding $n$ and
	$s$ is an integer with $1\le s\le\sqrt n$.
	Every integer in $\{1,\ldots,n\}$ determines a distinct such pair.
	However, $(3,\sqrt n)$ is not in the image, since
	$3(\sqrt n)^2=3n>n$. An injection between two finite sets of
	$n$ elements must be surjective, a contradiction.
	Thus $n<4^{\pi(n)}$, and integrality gives $n+1\le4^{\pi(n)}$. Taking
	$n=\lfloor x\rfloor\ge2$ gives $x<\lfloor x\rfloor+1\le 4^{\pi(x)}$, so
	$\sqrt{x}<2^{\pi(x)}$, and taking logarithms gives
	$\tfrac12\ln x<\pi(x)\ln 2$, which is the assertion.

\end{proof}
\begin{theorem}[Euler's theorem on prime reciprocals]\label{thm:9.6.2} $\sum_{p} \tfrac{1}{p} = \infty$.

\end{theorem}

\begin{proof} Suppose the sum converges; then there is
	$k$ with $\sum_{i>k}\tfrac{1}{p_i} < \tfrac{1}{2}$. Call $p_1,\ldots,p_k$
	the ``small'' primes. For $N \in \mathbb{Z}^+$, the number of $n \le N$
	divisible by some large prime $p_i$ ($i>k$) is at most
	$\sum_{i>k}\big\lfloor \tfrac{N}{p_i}\big\rfloor \le
		N\sum_{i>k}\tfrac{1}{p_i} < \tfrac{N}{2}$. Every remaining $n \le N$
	has only small prime factors; write $n = s^2 m$ with $m$
	squarefree. Then $m$ is a product of a subset of $\{p_1,\ldots,p_k\}$ (at
	most $2^k$ choices) and $s \le \sqrt{N}$ (at most $\sqrt{N}$ choices),
	so there are at most $2^k\sqrt{N}$ such $n$. Hence $N \le \tfrac{N}{2}
		+ 2^k\sqrt{N}$, that is, $\sqrt{N} \le 2^{k+1}$, which is false for large $N$, a
	contradiction.

\end{proof}
\begin{theorem}\label{thm:9.6.3}
	$\lim_{n\to\infty}\tfrac{\pi(n)}{n} = 0$.

\end{theorem}

\begin{proof} Fix $r$ and let $m = p_1\cdots p_r$. Among
	$1,\ldots,n$, the number of integers coprime to $m$ is
	$n\prod_{i=1}^r(1-\tfrac{1}{p_i}) + O(2^r)$ (inclusion--exclusion,
	error at most $2^r$). Every prime greater than $p_r$ is coprime to $m$, so
	$\pi(n) \le r + n\prod_{i=1}^r(1-\tfrac{1}{p_i}) + 2^r$. Dividing by
	$n$ and letting $n\to\infty$, we obtain $\limsup_n \tfrac{\pi(n)}{n} \le
		\prod_{i=1}^r(1-\tfrac{1}{p_i})$. Since $\sum\tfrac{1}{p_i}=\infty$
	(Theorem~\ref{thm:9.6.2}), $\prod_{i=1}^r(1-\tfrac{1}{p_i})\to 0$ as
	$r\to\infty$: indeed, $1-x\le e^{-x}$ for all real $x$ (the difference
	$e^{-x}-(1-x)$ vanishes at $x=0$ and has derivative $1-e^{-x}$, which is
	nonnegative for $x\ge0$ and nonpositive for $x\le0$), so
	$\prod_{i=1}^{r}(1-\tfrac{1}{p_i})\le\exp\big(-\sum_{i=1}^{r}\tfrac{1}{p_i}\big)\to0$,
	forcing the $\limsup$ to be $0$.

\end{proof}
\begin{theorem}\label{thm:9.6.4} For $n\in\mathbb{Z}^+$:
	\begin{enumerate}[label={\alph*)}]
		\item $\tfrac{4^n}{2n+1}\le\binom{2n}{n}\le 4^n$;
		\item for every prime $p$,
		      we have $v_p\!\binom{2n}{n}=\sum_{k\ge1}\big(\lfloor\tfrac{2n}{p^k}\rfloor-2\lfloor\tfrac{n}{p^k}\rfloor\big)$,
		      with each summand in $\{0,1\}$ and
		      $p^{v_p(\binom{2n}{n})}\le 2n$;
		\item if $n\ge3$ and $p$ is a prime with $\tfrac{2n}{3}<p\le n$, then
		      $p\nmid\binom{2n}{n}$;
		\item $\prod_{p\le x}p<4^{x}$ for all real $x\ge1$.
	\end{enumerate}

\end{theorem}

\begin{proof} (a) For $0\le k<2n$,
	\[
		\frac{\binom{2n}{k+1}}{\binom{2n}{k}}=\frac{2n-k}{k+1},
	\]
	which is greater than $1$ for $k<n$ and less than $1$ for $k\ge n$.
	Thus $\binom{2n}{n}$ is the largest of the
	$2n+1$ binomial coefficients $\binom{2n}{k}$, whose sum is $4^n$, giving both
	bounds. (b) By Legendre's formula (Theorem~\ref{thm:4.7.1})
	$v_p(m!)=\sum_{k\ge1}\lfloor m/p^k\rfloor$, so the displayed
	expression holds; each summand satisfies $\lfloor 2t\rfloor-2\lfloor
		t\rfloor\in\{0,1\}$, and summands vanish once $p^k>2n$, so at most
	$\log_p(2n)$ are nonzero and $p^{v_p}\le 2n$. (c) For such $p$,
	the inequality $p\le n<2p$ gives $\lfloor n/p\rfloor=1$ and $2p\le 2n<3p$ gives
	$\lfloor 2n/p\rfloor=2$, so the $k=1$ summand is $0$; and $p^2>(2n/3)^2\ge 2n$ for $n\ge5$ (for $n=3,4$ the only such
	prime is $p=3$, and $p^2=9>2n$ directly rules out $k\ge2$), so
	$v_p=0$. (d) It suffices to treat integer values of $x$: for real $x\ge1$,
	we have $\prod_{p\le x}p=\prod_{p\le\lfloor x\rfloor}p$ while
	$4^{\lfloor x\rfloor}\le 4^{x}$, so the bound at $\lfloor x\rfloor$
	yields it at $x$. We argue by induction on the integer $n$. For the base
	cases, when $n=1$, the empty product is $1<4$; when $n=2$, it is $2<16$;
	and when $n=3$, it is $2\cdot3=6<64=4^3$. If $n\ge4$ is even, it is not
	prime, so $\prod_{p\le n}p=\prod_{p\le n-1}p<4^{n-1}<4^{n}$ by induction. If
	$n=2m+1$ is odd, every prime $p$ with $m+1<p\le 2m+1$ divides
	$\binom{2m+1}{m}=\frac{(2m+1)!}{m!\,(m+1)!}$, since such $p$ divides
	the numerator $(2m+1)!$ but neither $m!$ nor $(m+1)!$ (as $p>m+1$);
	hence $\prod_{m+1<p\le 2m+1}p\le\binom{2m+1}{m}$. Since
	$\binom{2m+1}{m}$ and $\binom{2m+1}{m+1}$ are two equal terms in
	$\sum_{k}\binom{2m+1}{k}=2^{2m+1}$, we have $2\binom{2m+1}{m}\le
		2^{2m+1}$, that is, $\binom{2m+1}{m}\le 4^m$, so $\prod_{m+1<p\le
			2m+1}p\le 4^m$. With $\prod_{p\le m+1}p<4^{m+1}$ (induction, as
	$m+1<n$), we get $\prod_{p\le 2m+1}p=\prod_{p\le
			m+1}p\cdot\prod_{m+1<p\le 2m+1}p<4^{m+1}4^m=4^{2m+1}$.

\end{proof}
\begin{theorem}[Bertrand's postulate]\label{thm:9.6.5} For every
	$n\in\mathbb{Z}^+$ there is a prime $p$ with $n < p \le 2n$.

\end{theorem}

\begin{proof} Suppose that, for some $n\ge5$, there is no prime in
	$(n,2n]$. Factor $\binom{2n}{n}=\prod_{p\in\Primes} p^{v_p}$: primes $p>2n$ do
	not occur, none in $(\tfrac{2n}{3},n]$ occur by Theorem~\ref{thm:9.6.4}(c),
	and none in $(n,2n]$ by assumption, so only $p\le\tfrac{2n}{3}$ occur.
	Primes $p\le\sqrt{2n}$ contribute $p^{v_p}\le 2n$ each (at most
	$\sqrt{2n}$ of them), a product not exceeding $(2n)^{\sqrt{2n}}$; primes
	$\sqrt{2n}<p\le\tfrac{2n}{3}$ have $v_p\le1$, contributing
	at most $\prod_{p\le 2n/3}p<4^{2n/3}$ by Theorem~\ref{thm:9.6.4}(d). Hence
	$\binom{2n}{n}<(2n)^{\sqrt{2n}}4^{2n/3}$. With
	$\binom{2n}{n}\ge\tfrac{4^n}{2n+1}$ from Theorem~\ref{thm:9.6.4}(a),
	$4^{n/3}<(2n+1)(2n)^{\sqrt{2n}}$, that is, $\tfrac{2n}{3}\ln2<\ln(2n+1)+\sqrt{2n}\,\ln(2n)$. Put
	$F(x)=\tfrac{2x}{3}\ln2-\ln(2x+1)-\sqrt{2x}\,\ln(2x)$.
	We use the threshold $512$. Since $\sqrt{1024}=32$ and $1025<2^{11}$,
	\[
		F(512)=\frac{64}{3}\ln2-\ln1025
		>\frac{31}{3}\ln2>0.
	\]
	The function $u(t)=(\ln t+2)/\sqrt t$ has
	$u'(t)=-\ln t/(2t^{3/2})<0$ for $t>1$. Therefore, for $x\ge512$,
	\[
		\begin{aligned}
		F'(x)&=\frac23\ln2-u(2x)-\frac{2}{2x+1}\\
		&\ge\frac{17}{48}\ln2-\frac1{16}-\frac2{1025}
		>\frac{17}{96}-\frac18=\frac5{96}>0,
		\end{aligned}
	\]
	where $\ln2=\int_1^2dt/t>1/2$ and $2/1025<1/16$.
	Hence $F(x)>0$ for every $x\ge512$, contradicting $F(n)<0$.
	For the remaining $n\le511$, use the primes
	$2,3,5,7,13,23,43,83,163,317,631$, each less than twice its
	predecessor. For $n=1$, the prime $2\in(1,2]$ works directly. For $2\le n<631$, let $p$ be the smallest listed prime with $p>n$ and let $q$ be its predecessor in the list (which exists, since then $p\ge3$); then $q\le n$ (as $p$ is the smallest listed prime exceeding $n$) and $p<2q\le 2n$, so $p\in(n,2n]$. Since $511<631$, every $n\le511$ is covered. Hence the
	postulate holds for all $n\ge1$.

\end{proof}
\begin{theorem}[Bonse's inequality]\label{thm:9.6.6} For every integer
	$n\ge4$, we have $p_{n+1}^{2} < p_1 p_2 \cdots p_n$.

\end{theorem}

\begin{proof} Write $P_n = \prod_{i=1}^n p_i$. For $n=4$,
	we have $p_5^2 = 121 < 210 = 2\cdot3\cdot5\cdot7$. Now assume
	$p_{n+1}^2 < P_n$ for some $n\ge4$. Then $P_n > p_{n+1}^2 >
		4p_{n+1}$ (since $p_{n+1}\ge 11 > 4$). By Bertrand's postulate
	(Theorem~\ref{thm:9.6.5}) $p_{n+2} < 2p_{n+1}$, so $p_{n+2}^2 < 4p_{n+1}^2 =
		p_{n+1}\cdot 4p_{n+1} < p_{n+1}\cdot P_n = P_{n+1}$.

\end{proof}
\begin{theorem}\label{thm:9.6.7} There
	exist constants $0<c_1<c_2$ with $c_1\tfrac{x}{\ln x} < \pi(x) <
		c_2\tfrac{x}{\ln x}$ for all $x\ge2$.

\end{theorem}

\begin{proof} For the upper bound, Theorem~\ref{thm:9.6.4}(d) gives
	$\theta(x):=\sum_{p\le x}\ln p<x\ln4$; since each prime $p\in(\sqrt
		{x},x]$ has $\ln p>\tfrac{1}{2}\ln x$, we get $\tfrac{1}{2}\ln x\,(\pi(x)-\pi(\sqrt
		{x}))\le\theta(x)<x\ln4$, so $\pi(x)<\sqrt{x}+\tfrac{2\ln4\,x}{\ln
			x}<c_2\tfrac{x}{\ln x}$ for a suitable $c_2$. For the lower bound,
	Theorem~\ref{thm:9.6.4}(a) and (b) give $\tfrac{4^n}{2n+1}\le\binom{2n}{n}=\prod_{p\le
			2n}p^{v_p}\le(2n)^{\pi(2n)}$, so
	$\pi(2n)\ge\tfrac{n\ln4-\ln(2n+1)}{\ln(2n)}
		\ge c\tfrac{2n}{\ln(2n)}$
	for some $c>0$ and all sufficiently large $n$. For sufficiently
	large real $x$, take $n=\lfloor x/2\rfloor$; then
	$2n\le x$ and $2n\ge x-2\ge x/2$, so $\ln(2n)\le\ln x$. By monotonicity,
	\[
		\pi(x)\ge\pi(2n)\ge c\frac{2n}{\ln(2n)}
		\ge c\frac{x/2}{\ln x}.
	\]
	Hence the required strict lower bound holds for all sufficiently large
	$x$ with any $0<c_1<c/2$; shrinking $c_1$ further, if necessary,
	handles the bounded range $2\le x\le x_0$.

\end{proof}
\begin{theorem}[Schur's theorem]\label{thm:9.6.8} For every nonconstant
	polynomial $f\in\mathbb{Z}[x]$, infinitely many primes divide at
	least one nonzero value among $f(1),f(2),f(3),\ldots$.

\end{theorem}

\begin{proof} If $f(0)=0$, then $x\mid f(x)$, so $p\mid f(p)$
	for every prime $p$. A nonzero polynomial has only finitely many roots
	(Theorem~\ref{thm:3.6.2}, over $\mathbb{Q}$), so $f(p)\neq0$ for all but
	finitely many primes, proving this case. Otherwise, let $c=f(0)\neq0$ and suppose
	only the primes $p_1,\ldots,p_r$ divide nonzero values of $f$; put
	$Q=p_1\cdots p_r$. Writing $f(x)=\sum_{j\ge0}a_jx^{j}$ with $a_0=c$, we obtain
	$f(cQt)=c+\sum_{j\ge1}a_j(cQt)^{j}=c\big(1+Qt\,h(t)\big)$ for some
	$h\in\mathbb{Z}[t]$. Thus $g(t):=1+Qt\,h(t)\equiv1\pmod{Q}$, so
	$\gcd(g(t), Q)=1$; and since $f$ is nonconstant, $g$ is nonconstant,
	so $|g(t)|>1$ for all large $|t|$. Choose such a $t$ with the same sign as $c$, so that the argument $cQt$ is a positive integer. Any prime $q\mid g(t)$ then divides
	$f(cQt)$ yet is coprime to $Q$, so $q\notin\{p_1,\ldots,p_r\}$, a
	contradiction. Since $f(cQt)=c\,g(t)\neq0$, this proves that infinitely many primes divide nonzero values of $f$.

\end{proof}
\begin{theorem}[Mertens's first theorem]\label{thm:9.6.9} As $x\to\infty$,
	we have $\displaystyle \sum_{p \le x} \frac{\ln p}{p} = \ln x + O(1)$.

\end{theorem}

\begin{proof} Write $n=\lfloor x\rfloor$. Then
	$n=x+O(1)$. For $n\ge2$, monotonicity of $\ln t$ gives
	\[
		\int_1^n\ln t\,dt\le\ln(n!)\le\int_1^n\ln t\,dt+\ln n.
	\]
	Hence $\ln(n!)=n\ln n-n+O(\ln n)$.
	Applying the mean value theorem to $t\ln t$ yields $n\ln n=x\ln x+O(\ln x)$, and hence $\ln(n!)=x\ln x-x+O(\ln x)$. On the other hand, by Legendre's formula (Theorem~\ref{thm:4.7.1}),
	\[
		\ln(n!) = \sum_{p \le x} \Big(\sum_{k \ge 1} \big\lfloor \tfrac{x}{p^{k}} \big\rfloor\Big) \ln p,
	\]
	using $\lfloor n/p^{k}\rfloor = \lfloor x/p^{k}\rfloor$. The terms with $k \ge 2$ contribute at most $\sum_{p} \ln p \sum_{k \ge 2} \frac{x}{p^{k}} \le x \sum_{p} \frac{\ln p}{p(p-1)} = O(x)$; the series converges because $\sum_{p\in\Primes} (\ln p)/(p(p-1))$ is bounded by a constant multiple of $\sum_{n\ge2}(\ln n)/n^2$, which converges by the integral test. For $k = 1$, we have $\lfloor x/p\rfloor = x/p + O(1)$,
	so $\sum_{p \le x} \lfloor x/p\rfloor \ln p = x \sum_{p \le x} \frac{\ln p}{p} + O\big(\sum_{p \le x} \ln p\big) = x \sum_{p \le x} \frac{\ln p}{p} + O(x)$, with the error controlled by $\sum_{p \le x} \ln p < x \ln 4$ (Theorem~\ref{thm:9.6.4}(d)). Combining these estimates, we obtain $x \ln x - x + O(\ln x) = x \sum_{p \le x} \frac{\ln p}{p} + O(x)$; dividing by $x$ gives the claim.

\end{proof}
\begin{theorem}[Mertens's second theorem]\label{thm:9.6.10} There is a
	constant $M$ (the Mertens constant) such that, as $x\to\infty$,
	we have $\displaystyle \sum_{p \le x} \frac{1}{p} = \ln \ln x + M +
	O\big(\tfrac{1}{\ln x}\big)$; in particular
	$\sum_{p \le x} \frac{1}{p} = \ln \ln x + O(1)$.

\end{theorem}

\begin{proof} It suffices to consider the case $x\ge2$. Let $A(t) = \sum_{p \le t} \frac{\ln p}{p} = \ln t + r(t)$, where $r(t) = O(1)$ by Theorem~\ref{thm:9.6.9}. Write $a_n=\tfrac{\ln n}{n}$ when $n$ is prime and $a_n=0$ otherwise, so $A(t)=\sum_{n\le t}a_n$.
	Let $N=\lfloor x\rfloor$ and define $w(t)=1/\ln t$ for $t\ge2$. Substituting
	$a_n=A(n)-A(n-1)$ and telescoping gives Abel's identity
	\[
		\sum_{2\le n\le N}a_nw(n)=A(N)w(N)-\sum_{2\le n<N}A(n)\bigl(w(n+1)-w(n)\bigr).
	\]
	Since $A$ is constant on each interval $[n,n+1)$ and
	$w(n+1)-w(n)=\int_{n}^{n+1}w'(t)\,dt$ with
	$w'(t)=-1/(t\ln^{2}t)$, the sum subtracted on the right equals
	$\int_{2}^{N}A(t)w'(t)\,dt
	=-\int_{2}^{N}\frac{A(t)}{t\ln^{2}t}\,dt$, so
	\[
		\sum_{p\le N}\frac1p=\frac{A(N)}{\ln N}
		+\int_{2}^{N}\frac{A(t)}{t\ln^{2}t}\,dt.
	\]
	Finally, $A$ is constant on $[N,x]$, whence
	$A(N)w(N)=A(x)w(x)+\int_{N}^{x}\frac{A(t)}{t\ln^{2}t}\,dt$; since also
	$\sum_{p\le x}1/p=\sum_{p\le N}1/p$, the two right-hand sides combine into
	\[
		\sum_{p \le x} \frac{1}{p} = \frac{A(x)}{\ln x} + \int_{2}^{x} \frac{A(t)}{t \ln^{2} t}\, dt.
	\]
	Substitute $A(t)=\ln t+r(t)$. Since $r(t)=O(1)$ and $r$ is bounded
	on every compact interval in $[2,\infty)$, we may choose $C>0$ with
	$|r(t)|\le C$ for every $t\ge2$. The integral
	$\int_2^\infty r(t)/(t\ln^2t)\,dt$ converges absolutely, because
	$\int_2^\infty dt/(t\ln^2t)=1/\ln2$. Define
	\[
		M=1-\ln\ln2+\int_2^\infty\frac{r(t)}{t\ln^2t}\,dt.
	\]
	The preceding identity now gives the exact remainder formula
	\[
		\sum_{p\le x}\frac1p-\ln\ln x-M
		=\frac{r(x)}{\ln x}-\int_x^\infty\frac{r(t)}{t\ln^2t}\,dt.
	\]
	Consequently, for every $x\ge2$,
	\[
		\left|\sum_{p\le x}\frac1p-\ln\ln x-M\right|
		\le\frac{C}{\ln x}+C\int_x^\infty\frac{dt}{t\ln^2t}
		=\frac{2C}{\ln x}.
	\]
	This establishes both assertions.

\end{proof}
\subsection{Sieve Methods}

\begin{definition}[sifting function]\label{def:9.7.1} For real numbers
	$x\ge1$ and $z\ge2$, put $P(z)=\prod_{p\le z}p$ and $\Phi(x,z)=\#\{1\le n\le
		x \mid \gcd(n, P(z))=1\}$ (the number of positive integers up to $x$ with no prime factor at most $z$).

\end{definition}
\begin{theorem}[Legendre's sieve formula]\label{thm:9.7.1} In the
	setting of Definition~\ref{def:9.7.1}, $\displaystyle\Phi(x,z)=\sum_{d\mid
			P(z)}\mu(d)\Big\lfloor\frac{x}{d}\Big\rfloor$.

\end{theorem}

\begin{proof} By the M{\"o}bius identity (Theorem~\ref{thm:5.3.2}(a)),
	$\sum_{d\mid\gcd(n, P(z))}\mu(d)=[\gcd(n, P(z))=1]$. Summing over $n\le
		x$ and interchanging the order of summation gives $\Phi(x,z)=\sum_{d\mid P(z)}\mu(d)\,\#\{n\le
		x \mid d\mid n\}=\sum_{d\mid P(z)}\mu(d)\lfloor x/d\rfloor$.

\end{proof}
\begin{theorem}\label{thm:9.7.2}
	In the setting of Definition~\ref{def:9.7.1},
	$\displaystyle\Phi(x,z)=x\prod_{p\le
			z}\Big(1-\frac{1}{p}\Big)+O\big(2^{\pi(z)}\big)$, and $\pi(x)\le
		z+\Phi(x,z)$. Taking $z=c\ln x$ for sufficiently small fixed $c>0$
	gives the quantitative density-zero bound $\pi(x)=O\big(x/\ln\ln
		x\big)$.

\end{theorem}

\begin{proof} Substitute $\lfloor x/d\rfloor=x/d+O(1)$ into the sieve formula of
	Theorem~\ref{thm:9.7.1}; since $P(z)$ has $2^{\pi(z)}$ squarefree divisors
	and $\sum_{d\mid P(z)}\mu(d)/d=\prod_{p\le z}(1-1/p)$
	(multiplicativity), $\Phi(x,z)=x\prod_{p\le z}(1-1/p)+O(2^{\pi(z)})$.
	Every prime in $(z,x]$ has no prime factor at most $z$, so
	$\pi(x)\le\pi(z)+\Phi(x,z)\le z+\Phi(x,z)$. For $m=\lfloor z\rfloor$,
	expanding the finite Euler product gives $\prod_{p\le
			z}(1-1/p)^{-1}=\sum_{n:\,p\mid n\Rightarrow p\le z}1/n\ge\sum_{n\le
			m}1/n=H_m$, because every prime factor of $n\le m$ is at most $z$.
	Therefore $\prod_{p\le z}(1-1/p)\le1/H_m=O(1/\ln z)$. Now take
	$z=c\ln x$ with $0<c<1/\ln2$. Then
	$2^{\pi(z)}\le2^z=x^{c\ln2}=o(x/\ln\ln x)$, while $x\prod_{p\le
			z}(1-1/p)=O(x/\ln z)=O(x/\ln\ln x)$ and $z=O(\ln x)$.
	Substitution in the preceding bounds yields $\pi(x)=O(x/\ln\ln x)$.

\end{proof}
\begin{theorem}[Bonferroni inequalities]\label{thm:9.7.3} For
	finite sets $E_1,\ldots,E_r$ in a finite universe, define
	$T_k=\sum_{1\le i_1<\cdots<i_k\le r}
	|E_{i_1}\cap\cdots\cap E_{i_k}|$ for $1\le k\le r$, and put $T_k=0$
	for $k>r$. Then, for every $m\in\mathbb{Z}^+$,
	\[
		T_1-T_2+\cdots-T_{2m}\le\left|\bigcup_{i=1}^rE_i\right|\le
		T_1-T_2+\cdots+T_{2m-1}.
	\]

\end{theorem}

\begin{proof} A point of the union belonging to exactly $j\ge1$ of the
	sets contributes $\binom{j}{k}$ to $T_k$. Its contribution to the
	truncation through $h$ is $\sum_{k=1}^{h}(-1)^{k-1}\binom{j}{k}=1-(-1)^h\binom{j-1}{h}$ for $h<j$ (equivalently
	$\sum_{k=0}^{h}(-1)^{k}\binom{j}{k}=(-1)^{h}\binom{j-1}{h}$, which
		follows by induction on $h$ from Pascal's rule
	$\binom{j}{k}=\binom{j-1}{k-1}+\binom{j-1}{k}$ for $1\le k\le j$, obtained by counting the $k$-element subsets of a $j$-element set according to whether they contain a fixed element), and is $1$ for $h\ge j$ (the terms with $k>j$ vanish and the full alternating sum is $(1-1)^{j}=0$). Thus
		every point of the union contributes at least $1$ to an odd truncation and at
		most $1$ to an even truncation, while points outside the union
		contribute $0$. Summing pointwise proves both inequalities.

\end{proof}
\newpage

\section{Diophantine Approximation and Lattice Points}

\subsection{Rational Approximation}

\begin{theorem}[Dirichlet's approximation theorem]\label{thm:10.1.1} For an irrational $\alpha$,
	there are infinitely many rationals $\tfrac{p}{q}$ in lowest terms, with
	arbitrarily large $q\ge1$, satisfying $\big|\alpha - \tfrac{p}{q}\big| <
		\tfrac{1}{q^2}$.

\end{theorem}

\begin{proof} For $N\ge1$, the $N+1$ fractional parts
	$\{k\alpha\}$ ($0\le k\le N$) lie in $[0,1)$; by the pigeonhole principle, two of
	them (say those with $k=i$ and $k=j$, where $i>j$) lie within $\tfrac{1}{N}$ of each other, so $\lVert q\alpha\rVert <
		\tfrac{1}{N}$ with $q = i-j\le N$. Hence $|q\alpha - p| <
		\tfrac{1}{N}\le\tfrac{1}{q}$ for some integer $p$, giving $\big|\alpha -
		\tfrac{p}{q}\big| < \tfrac{1}{qN}\le\tfrac{1}{q^2}$. To see that there are infinitely many such rationals, suppose the set $S$ of rationals $p/q$ (in lowest terms, $q\ge1$) with $\big|\alpha-\tfrac{p}{q}\big|<\tfrac{1}{q^2}$ were finite. Since $\alpha$ is irrational, $\big|\alpha-\tfrac{p}{q}\big|>0$ for each such rational, so $\delta:=\min_{p/q\in S}\big|\alpha-\tfrac{p}{q}\big|>0$. Choose $N$ with $\tfrac{1}{N}<\delta$; the construction above produces $p/q$ with $q\le N$ and $\big|\alpha-\tfrac{p}{q}\big|<\tfrac{1}{qN}\le\tfrac{1}{q^2}$; writing this fraction in lowest terms as $p'/q'$, we get $q'\le q$ and
	$\big|\alpha-\tfrac{p'}{q'}\big|=\big|\alpha-\tfrac{p}{q}\big|<\tfrac{1}{q^{2}}\le\tfrac{1}{q'^{2}}$,
	hence $p'/q'\in S$, yet
	$\big|\alpha-\tfrac{p'}{q'}\big|<\tfrac{1}{N}<\delta$, contradicting
	the minimality of $\delta$. Therefore $S$ is infinite. For each fixed
	$q$, only finitely many integers $p$ can satisfy
	$\big|\alpha-\tfrac{p}{q}\big|<\tfrac{1}{q^2}$, so the denominators of
	the fractions in $S$ are unbounded.

\end{proof}
\begin{theorem}\label{thm:10.1.2} The number
	$e=\exp(1)=\sum_{k\ge0}\tfrac{1}{k!}$ is irrational.

\end{theorem}

\begin{proof} Suppose that $e=p/q$ with $p,q\in\mathbb{Z}^+$. Then
	$q!\,e=(q-1)!\,p$ is an integer. Write
	\[
		q!\,e=\sum_{k=0}^{q}\frac{q!}{k!}+R,\qquad
		R=\sum_{j\ge1}\frac{q!}{(q+j)!}.
	\]
	Each $q!/k!$ with $0\le k\le q$ is an integer, so $R$ is an integer.
	On the other hand, $R>0$, and
	\[
		\frac{q!}{(q+j)!}=\frac{1}{(q+1)(q+2)\cdots(q+j)}\le\frac{1}{(q+1)^{j}}
	\]
	for every $j\ge1$, with strict inequality for $j\ge2$. Hence
	\[
		0<R<\sum_{j\ge1}\frac{1}{(q+1)^{j}}=\frac{1}{q}\le1,
	\]
	which is impossible for an integer $R$.

\end{proof}
\begin{definition}[algebraic, transcendental, and Liouville numbers]\label{def:10.1.1}
	A real number $\alpha$ is algebraic of degree $n$ if it is a root of some
	nonzero polynomial in $\mathbb{Z}[x]$ of degree $n$ and of no nonzero
	polynomial in $\mathbb{Z}[x]$ of smaller degree; $\alpha$ is transcendental
	if it is not algebraic of any degree. An algebraic real number of degree $2$
	is called a real quadratic irrational. The monic polynomial of least degree
	in $\mathbb{Q}[x]$ with $\alpha$ as a root is its minimal polynomial; when
	$\alpha$ is a real quadratic irrational, the other
	root of this quadratic is its algebraic conjugate, denoted by $\alpha'$.
	A real number $\alpha$ is a Liouville
	number if it is irrational and, for every $N\in\mathbb{Z}^+$, there are
	integers $p$ and $q\ge2$ with
	$0<\bigl|\alpha-\tfrac{p}{q}\bigr|<q^{-N}$.

\end{definition}
\begin{theorem}[Liouville's theorem]\label{thm:10.1.3} If $\alpha$ is a real
	algebraic number of degree $n\ge2$, there exists a constant $c>0$ such that $\big|\alpha
		- \tfrac{p}{q}\big| > \tfrac{c}{q^{n}}$ for all rationals $\tfrac{p}{q}$ (written with $q \ge 1$).
	Consequently, every Liouville number is transcendental.

\end{theorem}

\begin{proof} Let $f\in\mathbb{Z}[x]$ be a nonzero
	polynomial of least degree with $\alpha$ as a root; its degree is $n$. Let $M =
		\max_{|x-\alpha|\le1}|f'(x)|$. This maximum exists by continuity and is positive because $f'$ is a nonzero polynomial. For $\tfrac{p}{q}$ with $\big|\alpha -
		\tfrac{p}{q}\big|\le1$, we have $f(\tfrac{p}{q})\neq0$, for otherwise polynomial
	long division over $\mathbb{Q}$ would give
	$f(x)=(x-\tfrac{p}{q})g(x)$ with $g\in\mathbb{Q}[x]$ of degree $n-1$;
	since $\alpha\neq\tfrac{p}{q}$ ($\alpha$ has degree $n\ge2$, so it is
	not rational), we would have $g(\alpha)=0$, and clearing denominators would give a
	nonzero integer polynomial of degree $n-1$ vanishing at $\alpha$,
	contradicting the choice of $f$.
	Moreover, $q^n f(\tfrac{p}{q})\in\mathbb{Z}\setminus\{0\}$, so $|f(\tfrac{p}{q})|\ge\tfrac{1}{q^n}$. By the mean value theorem, $|f(\tfrac{p}{q})| =
		|f'(\xi)|\,\big|\alpha - \tfrac{p}{q}\big|\le M\big|\alpha - \tfrac{p}{q}\big|$, so $\big|\alpha - \tfrac{p}{q}\big|\ge\tfrac{1}{Mq^n}$; taking $c=\min\{\tfrac{1}{2}, \tfrac{1}{2M}\}$ gives the strict bound $\big|\alpha - \tfrac{p}{q}\big|\ge\tfrac{1}{Mq^n}\ge\tfrac{2c}{q^n}>\tfrac{c}{q^n}$ in this range (as $c\le\tfrac{1}{2M}$), while for $\big|\alpha - \tfrac{p}{q}\big|>1$ it is immediate since $\tfrac{c}{q^n}\le c\le\tfrac{1}{2}<1$.
	Finally, a Liouville number is irrational, so it is not algebraic of
	degree $1$. If it were algebraic of degree $n\ge2$, we could choose $N>n$ so
	large that $2^{n-N}<c$. Its defining approximation with denominator
	$q\ge2$ would then satisfy
	$0<|\alpha-p/q|<q^{-N}\le2^{n-N}q^{-n}<cq^{-n}$, contradicting the
	bound just proved.

\end{proof}
\begin{theorem}[a Liouville number]\label{thm:10.1.4} The number
	\[
		\lambda=\sum_{n\ge1}10^{-n!}
	\]
	is a Liouville number in the sense of Definition~\ref{def:10.1.1}.
	Consequently, $\lambda$ is transcendental by Theorem~\ref{thm:10.1.3};
	in particular, transcendental real numbers exist.

\end{theorem}

\begin{proof} The series converges, because
	$0<10^{-n!}\le10^{-n}$. For $m\in\mathbb{Z}^+$, put
	\[
		q_m=10^{m!},\qquad
		p_m=q_m\sum_{n=1}^{m}10^{-n!}=\sum_{n=1}^{m}10^{m!-n!}\in\mathbb{Z}.
	\]
	The integers $n!$ with $n>m$ are distinct and at least $(m+1)!$, so
	\[
		0<\lambda-\frac{p_m}{q_m}=\sum_{n>m}10^{-n!}
		\le\sum_{j\ge(m+1)!}10^{-j}=\frac{10}{9}\cdot10^{-(m+1)!}.\qquad(*)
	\]

	First, $\lambda$ is irrational. Suppose that $\lambda=p/q$ with
	$p\in\mathbb{Z}$ and $q\in\mathbb{Z}^+$. For every $m$, the number
	$\lambda-p_m/q_m=(pq_m-p_mq)/(qq_m)$ is positive, so its numerator is a
	positive integer and $\lambda-p_m/q_m\ge1/(qq_m)=10^{-m!}/q$. Combined
	with $(*)$, this gives $10^{(m+1)!-m!}\le\tfrac{10}{9}q$, that is,
	$10^{m\cdot m!}\le\tfrac{10}{9}q$, for every $m\in\mathbb{Z}^+$. This is
	impossible, since the left-hand side is unbounded.

	Next, let $N\in\mathbb{Z}^+$ and take $m=N$. Then $q_N\ge10\ge2$, and
	$(*)$ gives
	\[
		0<\lambda-\frac{p_N}{q_N}\le\frac{10}{9}\cdot10^{-(N+1)!}<10^{1-(N+1)!}.
	\]
	Since $(N+1)!-1=N\cdot N!+(N!-1)\ge N\cdot N!$, it follows that
	$0<\lambda-p_N/q_N<10^{-N\cdot N!}=q_N^{-N}$. Hence $\lambda$ is a
	Liouville number, and Theorem~\ref{thm:10.1.3} shows that it is
	transcendental.

\end{proof}
\begin{theorem}\label{thm:10.1.5}
	Let $k\in\mathbb{Z}^+$, let $\xi_1,\ldots,\xi_k$ be real numbers, and let
	$Q\ge2$ be an integer. Then
	there exist an integer $q$ with $1\le q\le Q^{k}$ and integers
	$p_1,\ldots,p_k$ such that
	\[
		\left|\xi_i-\frac{p_i}{q}\right|<\frac{1}{qQ}\qquad(1\le i\le k).
	\]

\end{theorem}

\begin{proof} Partition the cube $[0,1)^{k}$ into $Q^{k}$ boxes of
	the form $\prod_{i=1}^{k}[c_i/Q,(c_i+1)/Q)$ with $0\le c_i<Q$. Consider the
	$Q^{k}+1$ points
	\[
		P_q=\big(\{q\xi_1\},\ldots,\{q\xi_k\}\big)\in[0,1)^{k},\qquad q=0,1,\ldots,Q^{k}.
	\]
	By the pigeonhole principle, two of them, say $P_{q_1}$ and $P_{q_2}$ with
	$q_1<q_2$, lie in the same box, so
	$\big|\{q_2\xi_i\}-\{q_1\xi_i\}\big|<1/Q$ for every $i$. Put $q=q_2-q_1$, so
	$1\le q\le Q^{k}$, and let $p_i=\lfloor q_2\xi_i\rfloor-\lfloor q_1\xi_i\rfloor$.
	Then
	\[
		q\xi_i-p_i=\{q_2\xi_i\}-\{q_1\xi_i\},
	\]
	so $|q\xi_i-p_i|<1/Q$, and dividing by $q$ gives the claim.

\end{proof}
\subsection{Beatty Sequences}

\begin{theorem}[Beatty's theorem]\label{thm:10.2.1} If $\alpha,\beta>1$ are
	irrational with $\tfrac{1}{\alpha} + \tfrac{1}{\beta} = 1$, then the Beatty
	sequences $B_\alpha = \{\lfloor n\alpha\rfloor \mid n\in\mathbb{Z}^+\}$ and $B_\beta =
		\{\lfloor n\beta\rfloor \mid n\in\mathbb{Z}^+\}$ partition $\mathbb{Z}^+$.

\end{theorem}

\begin{proof} For $N\ge1$, we have $\lfloor n\alpha\rfloor\le
		N$ if and only if $n\alpha<N+1$, that is, if and only if $n\le\lfloor(N+1)/\alpha\rfloor$, so $B_\alpha$
	has $\lfloor(N+1)/\alpha\rfloor$ terms in $[1,N]$; likewise $B_\beta$
	has $\lfloor(N+1)/\beta\rfloor$. Since $\tfrac{N+1}{\alpha} +
		\tfrac{N+1}{\beta} = N+1$ with both terms irrational, their fractional
	parts sum to an integer in $(0,2)$, hence to $1$, so their floors sum
	to $(N+1)-1 = N$. Each sequence has no repetitions: since $\alpha>1$, we have $\lfloor(n+1)\alpha\rfloor\ge\lfloor n\alpha\rfloor+1$, and similarly for $\beta$. As $N$ increases to $N+1$, the combined count rises by exactly $1$, so precisely one of the two sequences contains the value $N+1$; hence every positive integer occurs exactly once.

\end{proof}
\begin{definition}[equidistribution in the unit interval]\label{def:10.2.1}
	A sequence $(x_n)$ in $[0,1)$ is equidistributed in $[0,1)$ if, for
	every half-open interval $I=[a,b)\subseteq[0,1)$,
	\[
		\frac{1}{N}\#\{1\le n\le N\mid x_n\in I\}\longrightarrow |I|
	\]
	as $N\to\infty$, where $|I|$ denotes the length of $I$.

\end{definition}
\begin{theorem}[Weyl's equidistribution theorem]\label{thm:10.2.2} For every
	irrational $\gamma$, the sequence $(\{M\gamma\})_{M\ge1}$ is
	equidistributed in $[0,1)$; in particular,
	$\tfrac{1}{T}\sum_{M=1}^{T}\{M\gamma\}\to\int_0^1 x\,dx=\tfrac{1}{2}$ as $T\to\infty$.

\end{theorem}

\begin{proof} Fix an integer $q\ge1$ for
	which some integer $p$ satisfies $\gcd(p, q)=1$ and $|\gamma-\tfrac{p}{q}|<\tfrac{1}{q^2}$; by Theorem~\ref{thm:10.1.1} infinitely many such $q$ exist,
	so $q$ may be taken arbitrarily large. Put $\delta=\gamma-\tfrac{p}{q}$,
	so $|\delta|<\tfrac{1}{q^2}$. Group the positive integers into
	consecutive blocks of length $q$. On a block
	$B=\{M_0+1,\ldots,M_0+q\}$, the numbers $\{\tfrac{Mp}{q}\}$ run through
	the $q$ equally spaced points $0,\tfrac{1}{q},\ldots,\tfrac{q-1}{q}$ exactly
	once, since $\gcd(p, q)=1$. Write $M\gamma=\tfrac{Mp}{q}+M\delta$ with
	$M\delta=M_0\delta+(M-M_0)\delta$. Across the block, the term
	$(M-M_0)\delta$ varies by less than $q|\delta|<\tfrac{1}{q}$. Hence, up
	to a common shift $t:=M_0\delta$ and a per-term perturbation of size
	less than $\tfrac{1}{q}$, the block values $\{M\gamma\}$ are the points $\{\tfrac{r}{q}+t\}$ for $r=0,\ldots,q-1$. For any real $t$, these $q$ points are again
	equally spaced with gap $\tfrac{1}{q}$: they equal $t_0+\tfrac{r}{q}$ with
	$t_0=\{qt\}/q\in[0,\tfrac{1}{q})$. We bound the two block
	statistics with explicit constants. For the count, among the $q$ equally
	spaced points $t_0+\tfrac{r}{q}$, the number lying in an interval
	$[a,b)\subseteq[0,1)$ equals $\lfloor(b-a)q\rfloor$ or
	$\lceil(b-a)q\rceil$, hence differs from $(b-a)q$ by at most $1$. Now
	restore the per-term perturbations, each of size less than $\tfrac{1}{q}$: a point
	can enter or leave $[a,b)$ only if it lies within
	$\tfrac{1}{q}$ of an endpoint $a$ or $b$, where distances are measured on the circle
	$\mathbb{R}/\mathbb{Z}$. Here we use the following elementary count: an
	arc of length $L$ on $\mathbb{R}/\mathbb{Z}$ contains at most $\lfloor
		Lq\rfloor + 1$ of the $q$ points $t_0 + \tfrac{r}{q}$ ($0 \le r \le
		q-1$), because any $j$ of these points that lie in a common arc span,
	along that arc, a distance of at least $\tfrac{j-1}{q}$, so $j$ points
	force $\tfrac{j-1}{q} \le L$, that is, $j \le Lq + 1$. Now the points at
	distance less than $\tfrac{1}{q}$ from a given endpoint lie in an arc of length
	$\tfrac{2}{q}$, which, by this count, contains at most
	$\lfloor\tfrac{2}{q}\cdot q\rfloor + 1 = 3$ of them; so at most $6$
	points are affected by the two endpoints. Therefore
	$\big|\#\{M\in B \mid \{M\gamma\}\in[a,b)\}-(b-a)q\big|\le 7$.
	For the sum, the unperturbed block sum is
	$\sum_{r=0}^{q-1}(t_0+\tfrac{r}{q})=qt_0+\tfrac{q-1}{2}$, which lies in
	$\big[\tfrac{q-1}{2},\tfrac{q+1}{2}\big)$ because $qt_0\in[0,1)$, hence
	is within $1$ of $\tfrac{q}{2}$. When the perturbations are restored, each of
	the $q$ terms moves by less than $\tfrac{1}{q}$, changing the sum by at most
	$q\cdot\tfrac{1}{q}=1$ in total, except for wrap-around at an endpoint
	of $[0,1)$. All perturbations have the sign of $\delta$: if $\delta>0$,
	only points within $\tfrac{1}{q}$ of $1$ can wrap, while if $\delta<0$,
	only points within $\tfrac{1}{q}$ of $0$ can wrap. In either case, these
	points lie in an arc of length $\tfrac{1}{q}$, which, by the same count, contains at most
	$\lfloor\tfrac{1}{q}\cdot q\rfloor + 1 = 2$ of the $q$ points, and each
	of them changes the sum by less than $1$; so the wrap-around changes the
	sum by less than $2$. Therefore
	$\big|\sum_{M\in B}\{M\gamma\}-\tfrac{q}{2}\big|\le 1+1+2=4$.
	Thus for every full block, with the absolute constants $C_1=7$ and
	$C_2=4$, we have $\big|\#\{M\in B \mid \{M\gamma\}\in[a,b)\}-(b-a)q\big|\le C_1$
	and $\big|\sum_{M\in B}\{M\gamma\}-\tfrac{q}{2}\big|\le C_2$. When we sum
	over the $\lfloor T/q\rfloor$ full blocks in $[1,T]$, the per-block
	errors accumulate to at most $C_i\lfloor T/q\rfloor\le C_i\,T/q$,
	while for the fewer than $q$ leftover terms, if we write
	$T=q\lfloor T/q\rfloor+r$ with $0\le r<q$, both the actual count and
	its expected value $(b-a)r$ lie in $[0,r]$, so they differ by at
	most $q$, and likewise for the sum; hence $\big|\#\{M\le
		T \mid \{M\gamma\}\in[a,b)\}-(b-a)T\big|\le C_1\tfrac{T}{q}+q$ and
	$\big|\sum_{M=1}^T\{M\gamma\}-\tfrac{T}{2}\big|\le C_2\tfrac{T}{q}+q$.
	Dividing by $T$ and letting $T\to\infty$ with $q$ fixed gives
	$\limsup_T|\tfrac{1}{T}\#\{M\le T \mid \{M\gamma\}\in[a,b)\}-(b-a)|\le C/q$
	and $\limsup_T|\tfrac{1}{T}\sum_{M=1}^T\{M\gamma\}-\tfrac{1}{2}|\le C/q$ for
	an absolute constant $C$. Since $q$ may be taken arbitrarily large,
	both left-hand sides are $0$: the sequence $(\{M\gamma\})$ is
	equidistributed in $[0,1)$ and
	$\tfrac{1}{T}\sum_{M=1}^T\{M\gamma\}\to\tfrac{1}{2}$.

\end{proof}
\begin{theorem}\label{thm:10.2.3} Suppose $\mathbb{Z}^+$
	is partitioned into finitely many Beatty sequences
	$B_{\alpha_1},\ldots,B_{\alpha_m}$ with each $\alpha_i>1$ irrational.
	Then $m=2$; in particular, no partition into three or more Beatty
	sequences exists.

\end{theorem}

\begin{proof} For irrational $\alpha>1$ and $N\ge0$,
	we have $\lfloor k\alpha\rfloor\le N$ if and only if $k<\tfrac{N+1}{\alpha}$, so
	$B_\alpha$ has exactly $\big\lfloor\tfrac{N+1}{\alpha}\big\rfloor$
	terms in $[1,N]$. If the $B_{\alpha_i}$ partition $\mathbb{Z}^+$,
	adding these counts gives
	$\sum_{i=1}^{m}\big\lfloor\tfrac{M}{\alpha_i}\big\rfloor=M-1$ for all
	integers $M\ge1$. Writing $\lfloor x\rfloor=x-\{x\}$ and letting
	$M\to\infty$ forces $\sum_i\tfrac{1}{\alpha_i}=1$, and then
	$\sum_{i=1}^{m}\{M/\alpha_i\}=1$ for every $M\ge1$. Each $1/\alpha_i$ is irrational, since a rational reciprocal would make $\alpha_i$ rational. Averaging over $M=1,\ldots,T$ and applying Theorem~\ref{thm:10.2.2} to each $1/\alpha_i$ gives $\tfrac{1}{T}\sum_{M=1}^{T}\{M/\alpha_i\}\to\tfrac{1}{2}$;
	hence the average of $\sum_{i=1}^{m}\{M/\alpha_i\}$ tends to $m/2$ while the right-hand side
	remains equal to $1$, so $m=2$.

\end{proof}
\subsection{Continued Fractions}

\begin{definition}[continued fraction expansion]\label{def:10.3.1}
	For a real number $\theta$, put $\theta_0=\theta$. Whenever $\theta_n$ is
	defined, set $a_n=\lfloor\theta_n\rfloor$. If $\theta_n=a_n$, stop;
	otherwise put $\theta_{n+1}=1/(\theta_n-a_n)$. If the algorithm stops
	at index $N$, then $a_N=\theta_N$ and
	$\theta=[a_0;a_1,\ldots,a_N]$ is a finite continued fraction.
	The resulting finite or infinite sequence $[a_0;a_1,\ldots]$ is the
	continued fraction expansion
	of $\theta$; the integers $a_n$ are its partial quotients, and the $\theta_n$ are its complete quotients. An infinite continued fraction is
	eventually periodic if there are integers $r\ge0$ and $k\ge1$ such that
	$a_{n+k}=a_n$ for every $n\ge r$; it is purely periodic if one can take
	$r=0$. We write
	$[a_0;\ldots,a_{r-1},\overline{a_r,\ldots,a_{r+k-1}}]$ when the displayed
	block of $k$ partial quotients repeats indefinitely; for $r=0$ this is
	written $[\overline{a_0,\ldots,a_{k-1}}]$.

\end{definition}
\begin{theorem}\label{thm:10.3.1}
	In the continued fraction expansion of $\theta$, we have $\theta_n>1$ and
	hence $a_n\ge1$ for every $n\ge1$ at which $a_n$ is defined. If the
	expansion terminates at an index $N\ge1$, then its last partial quotient
	satisfies $a_N\ge2$.

\end{theorem}

\begin{proof} For $n\ge1$, the recurrence gives
	$\theta_n=1/(\theta_{n-1}-a_{n-1})>1$, because
	$0<\theta_{n-1}-a_{n-1}<1$ by the definition of $a_{n-1}=\lfloor
		\theta_{n-1}\rfloor$ together with $\theta_{n-1}\notin\mathbb{Z}$; hence
	$a_n=\lfloor\theta_n\rfloor\ge1$. At termination, $a_N=\theta_N>1$ is an
	integer, so $a_N\ge2$.

\end{proof}
\begin{theorem}\label{thm:10.3.2} The continued fraction expansion of $\theta$
	terminates if and only if $\theta\in\mathbb{Q}$; for $\theta=p/q$
	with $q\ge1$ and $\gcd(p, q)=1$, it coincides with Euclid's algorithm applied to $(p,q)$.

\end{theorem}

\begin{proof} A terminating continued fraction expansion is a finite nested
	fraction of integers, hence rational. Conversely, apply the algorithm to
	$\theta=p/q$: set $s_0=p$, $s_1=q$, so $\theta_0=s_0/s_1$. Each step performs
	the Euclidean division $s_n=a_ns_{n+1}+s_{n+2}$ with
	$0\le s_{n+2}<s_{n+1}$; since $\theta_n=s_n/s_{n+1}$, this gives
	$\theta_{n+1}=1/(\theta_n-a_n)=s_{n+1}/s_{n+2}$, which is exactly one step of the Euclidean algorithm. The remainders $q=s_1>s_2>\cdots\ge0$ strictly decrease,
	so there is a least $N\ge0$ with $s_{N+2}=0$. At that stage
	$s_N=a_Ns_{N+1}$ with $s_{N+1}\neq0$, so
	$\theta_N=s_N/s_{N+1}=a_N$ is an integer and the algorithm stops after
	producing $a_0,a_1,\ldots,a_N$; that is, the continued fraction expansion
	terminates.

\end{proof}
\begin{definition}[convergents]\label{def:10.3.2} Given
	$a_0\in\mathbb{Z}$ and $a_n\in\mathbb{Z}^+$ ($n\ge1$), set
	$p_{-1}=1,\ q_{-1}=0,\ p_0=a_0,\ q_0=1$ and, for $n\ge1$,
	put $p_n=a_np_{n-1}+p_{n-2}$ and $q_n=a_nq_{n-1}+q_{n-2}$. The fraction
	$p_n/q_n$ is the $n$-th convergent.

\end{definition}
\begin{theorem}\label{thm:10.3.3} For every $n\ge0$:
	\begin{enumerate}[label={\alph*)}]
		\item $p_nq_{n-1}-p_{n-1}q_n=(-1)^{n+1}$, so
		      $\gcd(p_n, q_n)=1$;
		\item $[a_0;\ldots,a_n]=p_n/q_n$;
		\item if $\theta=[a_0;\ldots,a_n,\beta]$ for a real $\beta>0$, then
		      $\theta=\dfrac{p_n\beta+p_{n-1}}{q_n\beta+q_{n-1}}$ and, for
		      $n\ge1$, the number $\theta$ lies strictly between
		      $p_{n-1}/q_{n-1}$ and $p_n/q_n$.
	\end{enumerate}

\end{theorem}

\begin{proof} The recursion can be expressed by the matrix identity $
		\begin{pmatrix}p_n & p_{n-1} \\ q_n&q_{n-1}
		\end{pmatrix}=
		\begin{pmatrix}a_0 & 1 \\ 1&0
		\end{pmatrix}\cdots
		\begin{pmatrix}a_n & 1 \\ 1&0
		\end{pmatrix}$; taking determinants (each factor has determinant
	$-1$) gives (a). For (c), the case $n=0$ follows directly from
	\[
		[a_0;\beta]=a_0+\frac{1}{\beta}
		=\frac{p_0\beta+p_{-1}}{q_0\beta+q_{-1}}.
	\]
	For $n\ge1$, assume the formula at index $n-1$ and put
	$\gamma=a_n+1/\beta$. Then
	\[
		\theta=[a_0;\ldots,a_{n-1},\gamma]
		=\frac{p_{n-1}\gamma+p_{n-2}}{q_{n-1}\gamma+q_{n-2}}
		=\frac{p_n\beta+p_{n-1}}{q_n\beta+q_{n-1}},
	\]
	which completes the induction.
	By (a), for $n\ge1$,
	\[
		\theta-\frac{p_n}{q_n}
		=\frac{(-1)^n}{q_n(q_n\beta+q_{n-1})},\qquad
		\theta-\frac{p_{n-1}}{q_{n-1}}
		=\frac{(-1)^{n+1}\beta}{q_{n-1}(q_n\beta+q_{n-1})}.
	\]
	The two differences have opposite signs, so $\theta$ lies strictly
	between the two convergents. Letting $\beta\to\infty$ in the identity
	gives $[a_0;\ldots,a_n]=\lim_{\beta\to\infty}
		\dfrac{p_n\beta+p_{n-1}}{q_n\beta+q_{n-1}}=\dfrac{p_n}{q_n}$, which is
	(b).

\end{proof}
\begin{theorem}\label{thm:10.3.4} If
	$\theta$ is irrational with convergents $p_n/q_n$, then, for every $n\ge0$,
	we have $\big|\theta-\tfrac{p_n}{q_n}\big|<\dfrac{1}{q_nq_{n+1}}\le\dfrac{1}{q_n^{2}}$;
	in particular, $p_n/q_n\to\theta$.

\end{theorem}

\begin{proof} By Theorem~\ref{thm:10.3.3}(a),
	$\tfrac{p_n}{q_n}-\tfrac{p_{n+1}}{q_{n+1}}=\dfrac{(-1)^{n+1}}{q_nq_{n+1}}$;
	by Theorem~\ref{thm:10.3.3}(c) applied at index $n+1$ (with $\beta=\theta_{n+2}$), the number $\theta$ lies between $p_n/q_n$ and $p_{n+1}/q_{n+1}$, so
	$\big|\theta-\tfrac{p_n}{q_n}\big|<\tfrac{1}{q_nq_{n+1}}$. Since
	$q_{n+1}\ge q_n$, the second inequality follows. Moreover,
	$q_{n+2}=a_{n+2}q_{n+1}+q_n\ge q_{n+1}+q_n$, so the positive
	denominators $q_n$ tend to infinity. Convergence follows.

\end{proof}
\begin{theorem}\label{thm:10.3.5} Let $\theta$
	be irrational with convergents $p_n/q_n$. Then:
	\begin{enumerate}[label={\alph*)}]
		\item for $n\ge0$, $p\in\mathbb{Z}$, and $q\in\mathbb{Z}^+$, if
		      $q<q_{n+1}$, then $|q\theta-p|\ge|q_n\theta-p_n|$;
		\item for $n\ge1$, $p\in\mathbb{Z}$, and $q\in\mathbb{Z}^+$,
		      the inequality $|\theta-p/q|<|\theta-p_n/q_n|$ forces $q>q_n$.
	\end{enumerate}

\end{theorem}

\begin{proof} (a) Solve $p=p_nu+p_{n+1}v$,
	$q=q_nu+q_{n+1}v$ for integers $u,v$, which is possible with $(u,v)\neq(0,0)$
	since $p_nq_{n+1}-p_{n+1}q_n=\pm1$ (Theorem~\ref{thm:10.3.3}(a)). If $v=0$, then
	$q\theta-p=u(q_n\theta-p_n)$ with $|u|\ge1$. Otherwise, $0<q<q_{n+1}$
	forces $u,v$ to have opposite signs: if $u,v\ge0$ with $v\ge1$, then
	$q=q_nu+q_{n+1}v\ge q_{n+1}$; if $u,v\le0$, then $q\le0$. Also, the
	consecutive convergent errors $q_n\theta-p_n$ and
	$q_{n+1}\theta-p_{n+1}$ have opposite signs (Theorem~\ref{thm:10.3.3}(c),
	applied at index $n+1$). Hence
	in
	$q\theta-p=u(q_n\theta-p_n)+v(q_{n+1}\theta-p_{n+1})$,
	the two summands have the same sign, so
	$|q\theta-p|\ge|q_n\theta-p_n|$. For (b), suppose instead that
	$q\le q_n$. Since $n\ge1$, the recurrence in
	Definition~\ref{def:10.3.2} gives $q_n<q_{n+1}$, so (a) yields
	\[
		\left|\theta-\frac{p}{q}\right|
		=\frac{|q\theta-p|}{q}
		\ge\frac{|q_n\theta-p_n|}{q}
		\ge\frac{|q_n\theta-p_n|}{q_n}
		=\left|\theta-\frac{p_n}{q_n}\right|.
	\]
	This is the contrapositive of (b).

\end{proof}
\begin{theorem}[Legendre's theorem on continued fractions]\label{thm:10.3.6} Let $\theta$ be
	irrational. Then:
	\begin{enumerate}[label={\alph*)}]
		\item for every $n\ge0$, at least one of the two consecutive
		      convergents $p_n/q_n$ and $p_{n+1}/q_{n+1}$ satisfies
		      $\big|\theta-\tfrac{p_j}{q_j}\big|<\tfrac{1}{2q_j^{2}}$
		      for its index $j\in\{n,n+1\}$;
		\item conversely, if $p\in\mathbb{Z}$, $q\in\mathbb{Z}^+$,
		      $\gcd(p, q)=1$, and
		      $\big|\theta-\tfrac{p}{q}\big|<\tfrac{1}{2q^{2}}$, then $p/q$
		      is a convergent of $\theta$.
	\end{enumerate}

\end{theorem}

\begin{proof} (a) If both convergents failed to satisfy the inequality, then, since
	$\theta$ lies between them, we would have
	$\tfrac{1}{q_nq_{n+1}}=\big|\tfrac{p_n}{q_n}-\tfrac{p_{n+1}}{q_{n+1}}\big|\ge\tfrac{1}{2q_n^2}+\tfrac{1}{2q_{n+1}^2}$,
	contradicting AM--GM, which gives
	$\tfrac{1}{2q_n^2}+\tfrac{1}{2q_{n+1}^2}>\tfrac{1}{q_nq_{n+1}}$ whenever
	$q_n\neq q_{n+1}$. (The degenerate case $q_n=q_{n+1}$ arises only for
	$n=0$ with $a_1=1$, in which case the convergents are $a_0$ and $a_0+1$; then
	$|\theta-a_0|+|\theta-(a_0+1)|=1$. If both distances were
	$\tfrac{1}{2}$, then $\theta=a_0+\tfrac{1}{2}$ would be rational;
	therefore irrationality forces one of them to be
	less than $\tfrac{1}{2}$, which is the required bound because both denominators
	are $1$, so (a) still holds.) (b) The recurrence $q_{n+1}=a_{n+1}q_n+q_{n-1}$ makes $(q_n)$ strictly increasing from $q_1$ onward and unbounded.
	Since $q_0=1\le q$, let $n$ be the largest index such that $q_n\le q$.
	Then $q_n\le q<q_{n+1}$. If $p/q\neq p_n/q_n$, then
	$\tfrac{1}{qq_n}\le\big|\tfrac{p}{q}-\tfrac{p_n}{q_n}\big|\le\big|\theta-\tfrac{p}{q}\big|+\big|\theta-\tfrac{p_n}{q_n}\big|$. Now $\big|\theta-\tfrac{p}{q}\big|<\tfrac{1}{2q^2}\le\tfrac{1}{2qq_n}$, and by Theorem~\ref{thm:10.3.5}(a) we have
	$|q_n\theta-p_n|\le|q\theta-p|<\tfrac{1}{2q}$, which gives
	$\big|\theta-\tfrac{p_n}{q_n}\big|<\tfrac{1}{2qq_n}$; hence
	$\tfrac{1}{qq_n}<\tfrac{1}{qq_n}$, a contradiction. So $p/q=p_n/q_n$.

\end{proof}
\begin{theorem}[Lagrange's theorem on periodic continued fractions]\label{thm:10.3.7} The continued fraction expansion of an irrational
	$\theta$ is eventually periodic if and only if $\theta$ is a real
	quadratic irrational.

\end{theorem}

\begin{proof} If
	$\psi=[\overline{b_0,\ldots,b_{k-1}}]$ is purely periodic, then
	$\psi=\dfrac{P\psi+P'}{Q\psi+Q'}$ for the convergent coefficients
	over one period, with $Q>0$. Thus
	$Q\psi^2+(Q'-P)\psi-P'=0$. The number $\psi$, whose expansion is infinite, is
	irrational by Theorem~\ref{thm:10.3.2}, so $\psi$ has degree exactly $2$.
	Any eventually periodic $\theta=[a_0;\ldots,a_{m-1},\psi]$ is a
	rational M{\"o}bius image of $\psi$ with determinant $\pm1$.
	The inverse transformation is rational as well: substituting it into the quadratic
	equation for $\psi$ and clearing denominators gives a nonzero polynomial
	of degree at most $2$ for $\theta$. Since $\theta$ is irrational,
	its degree is exactly $2$. Conversely, choose $A,B,C\in\mathbb{Z}$
	with $A\neq0$ such that $Ax^2+Bx+C$ is a nonzero integer multiple of
	the minimal polynomial of $\theta$, and put
	$f(x,y)=Ax^2+Bxy+Cy^2$.
	Set
	\[
		f_n(x,y)=f(p_nx+p_{n-1}y,\ q_nx+q_{n-1}y)
		=A_nx^2+B_nxy+C_ny^2.
	\]
	All $f_n$ share the discriminant of $f$ because $p_nq_{n-1}-p_{n-1}q_n=\pm1$.
	Since $\theta=(p_n\theta_{n+1}+p_{n-1})/(q_n\theta_{n+1}+q_{n-1})$,
	the complete quotient $\theta_{n+1}$ is a root of $f_n(x,1)=0$.
	By Theorem~\ref{thm:10.3.4}, $q_n^2\big|\tfrac{p_n}{q_n}-\theta\big|=q_n|q_n\theta-p_n|<1$, and $\big|\tfrac{p_n}{q_n}-\theta'\big|\le\big|\tfrac{p_n}{q_n}-\theta\big|+|\theta-\theta'|<1+|\theta-\theta'|$ once $n$ is large; hence $|A_n|=|f(p_n,q_n)|=|A|\,q_n^2\big|\tfrac{p_n}{q_n}-\theta\big|\big|\tfrac{p_n}{q_n}-\theta'\big|<|A|\,(1+|\theta-\theta'|)$ stays bounded, so $A_n$ takes finitely many integer values. The same bound applies to $C_n=f(p_{n-1},q_{n-1})=A_{n-1}$, and then $B_n$ is pinned down by the fixed discriminant $B_n^2-4A_nC_n=B^2-4AC$; hence only finitely many triples $(A_n,B_n,C_n)$ occur. None of the polynomials $f_n(x,1)$ is zero (their discriminant is the fixed nonzero discriminant), and each has at most two roots; therefore the union of all their roots is finite. Thus the complete quotients $\theta_{n+1}$ lie in a finite set, so
	$\theta_i=\theta_j$ for some $1\le i<j$.
	Each complete quotient determines the next continued-fraction step.
	Hence $\theta_{i+k}=\theta_{j+k}$ and $a_{i+k}=a_{j+k}$ for every
	$k\ge0$, so the expansion is eventually periodic.

\end{proof}
\begin{theorem}\label{thm:10.3.8} Call a real quadratic irrational $\theta$ reduced if
	$\theta>1$ and its conjugate satisfies $-1<\theta'<0$. Then the continued fraction expansion
	of a reduced $\theta$ is purely periodic.

\end{theorem}

\begin{proof} If $\theta$ is reduced, then
	$a_0=\lfloor\theta\rfloor\ge1$, $\theta_1=1/(\theta-a_0)>1$, and
	conjugation gives $\theta_1'=1/(\theta'-a_0)\in(-1,0)$ because
	$\theta'-a_0<-1$; so every complete quotient $\theta_n$ is reduced.
	From $\theta_n'=a_n+1/\theta_{n+1}'$ and $-1<\theta_n'<0$,
	we get $-1/\theta_{n+1}'=a_n-\theta_n'\in(a_n,a_n+1)$.
	Therefore $a_n=\lfloor-1/\theta_{n+1}'\rfloor\ge1$, so
	$\theta_{n+1}$ determines $a_n$ and hence
	$\theta_n=a_n+1/\theta_{n+1}$; the shift
	$\theta_n\mapsto\theta_{n+1}$ is thus injective on the complete
	quotients. By Theorem~\ref{thm:10.3.7}, the partial quotients satisfy
	$a_n=a_{n+k}$ for all $n\ge j$ and some $k\ge1$, so $\theta_j$ and
	$\theta_{j+k}$ are the same infinite continued fraction and
	$\theta_j=\theta_{j+k}$;
	injectivity implies $\theta_0=\theta_k$, giving pure
	periodicity.

\end{proof}
\begin{theorem}\label{thm:10.3.9} For a
	non-square $d\in\mathbb{Z}^+$, we have $\sqrt
		{d}=[a_0;\overline{a_1,\ldots,a_r}]$ with $a_0=\lfloor\sqrt{d}\rfloor$;
	the expansion is periodic from $a_1$ on.

\end{theorem}

\begin{proof} Since $d$ is not a square, $\sqrt{d}$ is irrational
	(Theorem~\ref{thm:2.3.7}). Put $a_0=\lfloor\sqrt{d}\rfloor$ and
	$\theta_1=1/(\sqrt{d}-a_0)$. Then $\theta_1>1$ (since $0<\sqrt
		{d}-a_0<1$) and $\theta_1'=1/(-\sqrt{d}-a_0)\in(-1,0)$ (since $\sqrt
		{d}+a_0>1$), so $\theta_1$ is reduced and, by Theorem~\ref{thm:10.3.8}, its expansion is purely
	periodic.

\end{proof}
\begin{theorem}\label{thm:10.3.10}
	For a non-square $d\in\mathbb{Z}^+$, the equation $x^2-dy^2=1$ has an integer solution
	$(x,y)$ with $xy\neq0$; moreover, every positive integer solution $(x,y)$ has $x/y$ equal
	to a convergent of $\sqrt{d}$.

\end{theorem}

\begin{proof} By Theorem~\ref{thm:10.3.9}, write $\sqrt
		{d}=[a_0;\overline{a_1,\ldots,a_\ell}]$ and let $n$ be the least even
	multiple of the period length $\ell$ (so $n=\ell$ or $n=2\ell$). Then
	$\theta_{n+1}=\theta_1=1/(\sqrt{d}-a_0)$ and, by Theorem~\ref{thm:10.3.3}(c),
	\[
		\sqrt{d}=\dfrac{p_n\theta_1+p_{n-1}}{q_n\theta_1+q_{n-1}}.
	\]

	Clear the denominator and use $\theta_1(\sqrt{d}-a_0)=1$:
	multiplying $\sqrt{d}(q_n\theta_1+q_{n-1})=p_n\theta_1+p_{n-1}$ by
	$\sqrt{d}-a_0$ gives
	$\sqrt{d}\,(q_n-a_0q_{n-1})+dq_{n-1}=\sqrt{d}\,p_{n-1}+(p_n-a_0p_{n-1})$;
	since $1,\sqrt{d}$ are linearly independent over $\mathbb{Q}$ (a
	rational relation would make $\sqrt{d}$ rational, contrary to
	Theorem~\ref{thm:2.3.7}), comparing rational and irrational parts gives
	the coefficient identities $p_{n-1}=q_n-a_0 q_{n-1}$ and
	$q_{n-1}d=p_n-a_0 p_{n-1}$. Substituting these into the determinant
	identity of Theorem~\ref{thm:10.3.3}(a) yields
	\[
		p_{n-1}^2-d q_{n-1}^2=p_{n-1}q_n-q_{n-1}p_n=(-1)^{n}.
	\]

	Since $n$ is even, $p_{n-1}^2-d q_{n-1}^2=1$, so
	$(x,y)=(p_{n-1},q_{n-1})$ is a nontrivial solution. Conversely, if
	$x,y>0$ and $x^2-dy^2=1$, then $x-y\sqrt{d}=1/(x+y\sqrt{d})>0$ and
	\[
		\Bigl|\frac{x}{y}-\sqrt{d}\Bigr|=\frac{1}{y(x+y\sqrt{d})}<\frac{1}{2y^2},
	\]
	because $x>y\sqrt{d}$ forces $x+y\sqrt{d}>2y\sqrt{d}\ge 2\sqrt{2}\, y>2y$
	(using $d\ge2$). Moreover, any common divisor of $x$ and $y$ divides
	$x^{2}-dy^{2}=1$, so $\gcd(x,y)=1$. Theorem~\ref{thm:10.3.6}(b) therefore
	shows that $x/y$ is a
	convergent of $\sqrt{d}$. Thus the continued-fraction argument proves
	both existence and that $x/y$ is a convergent of $\sqrt{d}$ for every positive solution $(x,y)$.

\end{proof}
\subsection{Farey Fractions}

\begin{definition}[Farey sequence]\label{def:10.4.1} For $n\ge1$,
	the Farey sequence $F_n$ is the increasing list of all rationals
	$h/k$ in $[0,1]$ with $1\le k\le n$, $0\le h\le k$, and $\gcd(h, k)=1$.

\end{definition}
\begin{theorem}\label{thm:10.4.1} If $\tfrac{h}{k}<\tfrac{h'}{k'}$ are consecutive in $F_n$, then $h'k-hk'=1$.
	Moreover, the first fraction to appear between them, in $F_{k+k'}$, is
	their mediant $\tfrac{h+h'}{k+k'}$.

\end{theorem}

\begin{proof} Since $\gcd(h, k)=1$, the equation $kx-hy=1$
	has integer solutions, all of the form $(x_0+ht,\,y_0+kt)$; choose
	the one with $n-k<y\le n$ and set $\tfrac{a}{b}=\tfrac{x}{y}$. Then
	$\gcd(a, b)=1$ (as $ka-hb=1$), and $a/b=h/k+1/(kb)>0$. Also $h/k<1$ gives $k-h\ge1$, and $k(a-b)=1-(k-h)b\le0$, so $a\le b$. Together with $b\le n$, this proves $a/b\in F_n$ and $\tfrac{a}{b}>\tfrac{h}{k}$. If some $\tfrac{p}{q}\in F_n$ had
	$\tfrac{h}{k}<\tfrac{p}{q}<\tfrac{a}{b}$, then $\tfrac{p}{q}-\tfrac{h}{k}\ge\tfrac{1}{kq}$ and $\tfrac{a}{b}-\tfrac{p}{q}\ge\tfrac{1}{bq}$, whence
	$\tfrac{1}{kb}=\tfrac{a}{b}-\tfrac{h}{k}\ge\tfrac{b+k}{kbq}$, giving $q\ge
		b+k>n$, which is impossible. So $\tfrac{a}{b}$ is the successor
	$\tfrac{h'}{k'}$ and $h'k-hk'=1$. Finally, the mediant satisfies
	$\tfrac{h}{k}<\tfrac{h+h'}{k+k'}<\tfrac{h'}{k'}$, and the identity
	$(h+h')k-h(k+k')=h'k-hk'=1$ shows that it is in lowest terms; its
	denominator is $k+k'$. Moreover, if $\tfrac{p}{q}$ lies strictly between
	$\tfrac{h}{k}$ and $\tfrac{h'}{k'}$, then the same estimate gives
	$q\ge k+k'$, with equality only for the mediant: equality forces both
	difference estimates to hold with equality, so $pk-hq=1$ and $h'q-pk'=1$, and
	subtracting gives $p(k+k')=q(h+h')$, whence $q=k+k'$ implies $p=h+h'$.
	Hence the first
	fraction inserted between the neighbors, the one in $F_{k+k'}$, is
	exactly the mediant.

\end{proof}
\subsection{Hurwitz's Theorem}

\begin{theorem}[Hurwitz's theorem]\label{thm:10.5.1} For every irrational
	$\alpha$, there are infinitely many rationals $\tfrac{p}{q}$, written in
	lowest terms with $q\ge1$, such that
	$\big|\alpha-\tfrac{p}{q}\big|<\dfrac{1}{\sqrt{5}\,q^2}$; the constant
	$\sqrt{5}$ is best possible.

\end{theorem}

\begin{proof} Let $p_n/q_n$ be the convergents of
	$\alpha$, let $\alpha_{n+1}$ be the $(n+1)$-th complete quotient, and let
	$\beta_n=q_{n-1}/q_n$ (which lies in $(0,1]$ for $n\ge1$, with
	$\beta_n<1$ for $n\ge2$). From
	$\alpha=\dfrac{\alpha_{n+1}p_n+p_{n-1}}{\alpha_{n+1}q_n+q_{n-1}}$
	(Theorem~\ref{thm:10.3.3}) one gets
	$\big|\alpha-\tfrac{p_n}{q_n}\big|=\dfrac{1}{q_n^2(\alpha_{n+1}+\beta_n)}$,
	so the inequality at index $n$ holds if and only if
	$\alpha_{n+1}+\beta_n>\sqrt{5}$. Suppose it failed at three consecutive
	indices $n-1,n,n+1$. With $u=\alpha_{n+1}>1$, $v=\beta_n\in(0,1]$
	and the identities $\alpha_n=a_n+1/\alpha_{n+1}$ and $1/\beta_n=a_n+\beta_{n-1}$,
	failure at $n-1$ becomes $1/u+1/v\le\sqrt{5}$ and failure at $n$
	becomes $u+v\le\sqrt{5}$. These give $u/(\sqrt{5}u-1)\le v\le\sqrt{5}-u$,
	hence $u^2-\sqrt{5}u+1\le0$, that is, $u\le\tfrac{\sqrt{5}+1}{2}=\phi$. The
	same argument applied to indices $n,n+1$ gives $\alpha_{n+2}\le\phi$;
	then $u=\alpha_{n+1}=a_{n+1}+1/\alpha_{n+2}\ge1+1/\phi=\phi$. Thus
	$u=\phi$ with equality throughout, which forces
	$v=\beta_n=\sqrt{5}-\phi=\tfrac{\sqrt{5}-1}{2}$; this value is irrational,
	contradicting the fact that $\beta_n=q_{n-1}/q_n$ is rational. Hence at least
	one of every three consecutive convergents satisfies the inequality,
	giving infinitely many such rationals. For optimality, take
	$\alpha=\phi=[\,\overline1\,]$: every complete quotient equals $\phi$
	and the denominator recurrence gives
	$\beta_n=[0;\underbrace{1,\ldots,1}_{n\text{ terms}}]\to[0;\overline1]=1/\phi$
	by Theorem~\ref{thm:10.3.4}. Therefore
	$q_n^2\big|\phi-\tfrac{p_n}{q_n}\big|=\tfrac{1}{\phi+\beta_n}\to\tfrac{1}{\phi+1/\phi}=\tfrac{1}{\sqrt{5}}$.
	Hence, for any $c>\sqrt{5}$, only finitely many convergents satisfy
	$\big|\phi-\tfrac{p}{q}\big|<\tfrac{1}{cq^2}$, and since $c>2$, every
	solution is a convergent by Theorem~\ref{thm:10.3.6}(b); so only finitely many
	$\tfrac{p}{q}$ satisfy this inequality.
	Since $1/(\sqrt{5}\,q^2)<1/q^2$, the existence assertion strengthens
	Theorem~\ref{thm:10.1.1}.

\end{proof}
\subsection{Kronecker's Theorem (One Dimension)}

\begin{theorem}[Kronecker's theorem]\label{thm:10.6.1} If $\alpha$ is
	irrational, then for every real $\gamma$ and every $\varepsilon>0$
	there are integers $n\ge1$ and $m$ with
	$|n\alpha-m-\gamma|<\varepsilon$; equivalently, the fractional parts
	$(\{n\alpha\})_{n\ge1}$ are dense in $[0,1]$.

\end{theorem}

\begin{proof} By Dirichlet's theorem (Theorem~\ref{thm:10.1.1}),
	choose $q\ge1$ with $0<\delta:=\lVert q\alpha\rVert<\min\{\varepsilon,1/2\}$.
	Put $\theta=\{q\alpha\}$, so $\theta=\delta$ or $\theta=1-\delta$,
	and let $M=\lceil1/\delta\rceil-1$.
	The points $0,\delta,\ldots,M\delta$ on $\mathbb{R}/\mathbb{Z}$
	have successive gaps $\delta$ and a final gap $1-M\delta\le\delta$.
	Thus the points $\{k\theta\}$ with $0\le k\le M$ have all circular gaps
	at most $\delta$: for $\theta=1-\delta$, simply reflect the circle.

	Translating these points by $\theta$ preserves their gaps and gives
	$\{k\theta\}$ with $1\le k\le M+1$. Every point of the circle is
	therefore at distance at most $\delta/2$ from one of these positive
	multiples. Choose $1\le k\le M+1$ such that
	$\lVert kq\alpha-\gamma\rVert\le\delta/2<\varepsilon$.
	Set $n=kq\ge1$ and choose $m$ nearest to $n\alpha-\gamma$; then
	$|n\alpha-m-\gamma|<\varepsilon$, as required.

\end{proof}
\subsection{Pell's Equation}

\begin{definition}[Pell's equation]\label{def:10.7.1} For a non-square
	$d\in\mathbb{Z}^+$, the equation $x^2-dy^2=1$ is Pell's equation and $x^2-dy^2=-1$
	is the negative Pell equation.

\end{definition}
\begin{theorem}\label{thm:10.7.1} Let $d\in\mathbb{Z}^+$ be a non-square.
	Pell's equation $x^2-dy^2=1$ has a least positive integer solution $(x_1,y_1)$,
	and its positive integer solutions are precisely the pairs $(x_n,y_n)$
	defined by
	$x_n+y_n\sqrt{d}=(x_1+y_1\sqrt{d})^{n}$ for $n\in\mathbb{Z}^+$.
	When the negative equation is solvable,
	with least positive integer solution $(u_1,v_1)$, its positive integer
	solutions are given by the odd powers $u_n+v_n\sqrt
		{d}=(u_1+v_1\sqrt{d})^{2n-1}$ for $n\in\mathbb{Z}^+$.

\end{theorem}

\begin{proof} Existence of a nontrivial solution was
	already established constructively from the continued fraction expansion
	of $\sqrt{d}$ in Theorem~\ref{thm:10.3.10}; here we give a self-contained
	pigeonhole proof and then determine the full group structure. By
	Theorem~\ref{thm:10.1.1} there are infinitely many integer pairs $(x,y)$
	with $y\ge1$ and $\big|x-y\sqrt{d}\big|<\tfrac{1}{y}$. Since $d\ge2$ and
	$y\ge1$, these pairs also satisfy
	$x>y\sqrt d-1/y\ge\sqrt2-1>0$. For these pairs,
	$\big|x+y\sqrt{d}\big|<2y\sqrt{d}+\tfrac{1}{y}$, so
	$\big|x^2-dy^2\big|<2\sqrt{d}+1$. Hence some nonzero integer $k$ with
	$|k|<2\sqrt{d}+1$ satisfies $x^2-dy^2=k$ for infinitely many pairs,
	and two of them, say $(r_1,s_1)\neq(r_2,s_2)$, agree modulo $|k|$ in both
	coordinates. Put $u+v\sqrt{d}=\tfrac{(r_1+s_1\sqrt{d})(r_2-s_2\sqrt
			{d})}{k}$; then $u=\tfrac{r_1r_2-ds_1s_2}{k}$ and
	$v=\tfrac{r_2s_1-r_1s_2}{k}$ are integers (modulo $|k|$, the numerators are congruent to
		$r_1^2-ds_1^2\equiv0$ and to $r_1s_1-r_1s_1=0$, respectively), and
	$u^2-dv^2=\tfrac{k\cdot k}{k^2}=1$. Here $v\neq0$: if $v=0$, then
	$u^2=1$, so $u=\pm1$ and therefore
	$r_1+s_1\sqrt{d}=\pm(r_2+s_2\sqrt{d})$, which is impossible for distinct
	positive pairs. Since the norm equation depends only on the squares of
	$u$ and $v$, replacing them by $|u|$ and $|v|$ gives a positive
	nontrivial solution. For every real bound $C$, only finitely many
	positive integer pairs satisfy $x+y\sqrt d\le C$, so the nonempty set
	of positive solutions has a least member. Order the positive solutions
	by the increasing value $x+y\sqrt{d}>1$ and let
	$\varepsilon=x_1+y_1\sqrt{d}$ be the least. If
	$\xi=x+y\sqrt{d}\in\mathbb{Z}[\sqrt{d}]$ has norm $1$, then
	$\xi^{-1}=x-y\sqrt{d}\in\mathbb{Z}[\sqrt{d}]$. Hence products and
	quotients of integer solutions are again integer solutions, by the
	multiplicativity of the norm. If a positive solution $\eta$ were
	not a power of $\varepsilon$, we could pick $n$ with
	$\varepsilon^n<\eta<\varepsilon^{n+1}$; then
	$\eta\varepsilon^{-n}=u+v\sqrt{d}\in\mathbb{Z}[\sqrt{d}]$ has norm $1$ and lies in
	$(1,\varepsilon)$, and $u-v\sqrt{d}=(u+v\sqrt{d})^{-1}\in(0,1)$; adding and subtracting
	these two inequalities gives $u=\tfrac{(u+v\sqrt{d})+(u-v\sqrt
			{d})}{2}>\tfrac12$ and $v\sqrt{d}=\tfrac{(u+v\sqrt{d})-(u-v\sqrt
			{d})}{2}>0$, so $u,v>0$; this gives a positive solution below $\varepsilon$, contradicting
	minimality. Hence every positive solution is $\varepsilon^n$, that is,
	$x_n+y_n\sqrt{d}=(x_1+y_1\sqrt{d})^n$.

	For the negative Pell equation, suppose $x^2-dy^2=-1$ is solvable.
	Consider all units $\xi=x+y\sqrt{d}>1$ with $x,y>0$ and
	$N(\xi)=\pm1$, and let $\eta$ be the least one. We justify the same
	descent with both possible norm signs. If such a $\xi$ is not a power
	of $\eta$, choose $m$ with $\eta^m<\xi<\eta^{m+1}$ and put
	$\zeta=\xi\eta^{-m}=u+v\sqrt{d}$. Since
	$\eta^{-1}=N(\eta)\overline\eta\in\mathbb{Z}[\sqrt{d}]$, both $u$ and
	$v$ are integers. If $N(\zeta)=1$, then
	$\overline\zeta=\zeta^{-1}\in(0,1)$; if $N(\zeta)=-1$, then
	$\overline\zeta=-\zeta^{-1}\in(-1,0)$. In either case
	\[
		u=\frac{\zeta+\overline\zeta}{2}>0,\qquad
		v\sqrt{d}=\frac{\zeta-\overline\zeta}{2}>0.
	\]
	Thus $\zeta\in(1,\eta)$ is a smaller member of the same set, a
	contradiction. Hence every such $\xi$ is a power $\eta^m$. Since
	$N(\eta^m)=N(\eta)^m$, if $N(\eta)=1$, then no power would have norm
	$-1$, contradicting solvability; hence $N(\eta)=-1$. Thus the norm
	$+1$ solutions are exactly the even powers of $\eta$ and the norm
	$-1$ solutions exactly the odd powers. In particular, the least
	positive solution of the negative equation is $\delta=\eta$, the
	fundamental solution of the positive equation is
	$\varepsilon=\eta^2=\delta^2$, and the negative solutions are
	$(u_n,v_n)$ defined by $u_n+v_n\sqrt{d}=\delta^{2n-1}$ ($n\ge1$).

\end{proof}
\begin{theorem}\label{thm:10.7.2} Let $d\in\mathbb{Z}^+$ be a non-square and $a\neq0$ an
	integer. If $x^2-dy^2=a$ has one integer solution, it has infinitely
	many: if $(x_0,y_0)$ is a solution and $(x_1,y_1)$ is the least
	positive solution of $x^2-dy^2=1$, then $x_k+y_k\sqrt{d}=(x_0+y_0\sqrt
		{d})(x_1+y_1\sqrt{d})^{k}$ solves $x^2-dy^2=a$ for every
	$k\in\mathbb{Z}_{\ge0}$, with
	distinct $k$ giving distinct solutions.

\end{theorem}

\begin{proof} The norm $N(x+y\sqrt{d})=x^2-dy^2$ is
	multiplicative: from $(x+y\sqrt{d})(u+v\sqrt{d})=(xu+dyv)+(xv+yu)\sqrt
		{d}$ one computes $N\big((x+y\sqrt{d})(u+v\sqrt
		{d})\big)=(xu+dyv)^2-d(xv+yu)^2=(x^2-dy^2)(u^2-dv^2)$, since the cross terms
	$2dxyuv$ cancel; that is, $N((x+y\sqrt{d})(u+v\sqrt{d}))=N(x+y\sqrt
		{d})N(u+v\sqrt{d})$. Hence $N\big((x_0+y_0\sqrt
		{d})(x_1+y_1\sqrt{d})^{k}\big)=N(x_0+y_0\sqrt{d})\,N(x_1+y_1\sqrt
		{d})^{k}=a\cdot1^{k}=a$, so each $(x_k,y_k)$ is a solution. Negating
	$(x_0,y_0)$ negates every pair $(x_k,y_k)$ produced by the displayed
	formula, so pairwise distinctness is unaffected. Replacing
	$(x_0,y_0)$ by $(-x_0,-y_0)$ if necessary, we may assume
	$x_0+y_0\sqrt{d}>0$. Since $x_1+y_1\sqrt{d}>1$, the numbers
	$(x_0+y_0\sqrt{d})(x_1+y_1\sqrt{d})^k$ are then positive and strictly
	increasing in $k$; hence the resulting solutions are pairwise
	distinct.

\end{proof}
\subsection{Pick's Theorem}

\begin{definition}[planar lattices and lattice polygons]\label{def:10.8.1}
	A planar lattice is a subgroup
	$\Lambda=\mathbb{Z}e_1+\mathbb{Z}e_2\subseteq\mathbb{R}^2$ generated by
	linearly independent vectors $e_1,e_2$. The associated half-open
	fundamental parallelogram is
	\[
		\{t_1e_1+t_2e_2\mid 0\le t_1,t_2<1\},
	\]
	and its area $d(\Lambda)=|\det(e_1,e_2)|$ is called the covolume of
	$\Lambda$; it is independent of the chosen lattice basis by
	Theorem~\ref{thm:10.8.1}. When no lattice
	is specified, a lattice point means a point of $\mathbb{Z}^2$. A lattice
	path is a polygonal path whose vertices are lattice points, and a lattice
	polygon is a polygon whose vertices are lattice points; it is simple if
	its boundary has no self-intersection apart from common endpoints of
	adjacent sides. A lattice triangle is primitive if its only lattice
	points are its three vertices.

\end{definition}
\begin{theorem}[covolume]\label{thm:10.8.1}
	Let $\Lambda=\mathbb{Z}e_1+\mathbb{Z}e_2$ be a planar lattice. If $f_1,f_2$ is
	another pair of linearly independent vectors generating $\Lambda$, then
	$|\det(f_1,f_2)|=|\det(e_1,e_2)|$.

\end{theorem}

\begin{proof} If $f_1,f_2$ is
	another pair of linearly independent vectors generating $\Lambda$, then
	each of the pairs $(e_1,e_2)$ and $(f_1,f_2)$ is obtained from the other
	by an integer matrix, these two matrices are mutually inverse by linear
	independence, so their determinants are integers with product $1$, and
	hence $\det(f_1,f_2)=\pm\det(e_1,e_2)$.

\end{proof}
\begin{theorem}[Pick's theorem]\label{thm:10.8.2} A simple lattice polygon
	with $I$ interior and $B$ boundary lattice points has area $S = I +
		\tfrac{B}{2} - 1$.

\end{theorem}

\begin{proof} Both the area and $f(P) := I + \tfrac{B}{2} - 1$ are additive when
	a polygon is cut into $P_1, P_2$ along a lattice path whose two endpoints
	lie on the boundary and whose interior vertices lie in the interior (where a
	lattice chord is the case of a path with no interior vertex). Area is
	clearly additive. For $f$, suppose the path contains $m$ lattice points in its interior, in addition to its two endpoints. These $m$ points
	are interior to $P$ but become boundary points of both $P_1$ and
	$P_2$, while the two endpoints of the path lie on $\partial P$; hence,
	distributing the boundary lattice points of $P$ between the two pieces
	(the two endpoints counted in each), we obtain $B_1 + B_2 = (B + 2) + 2m$ and
	$I_1 + I_2 = I - m$. Therefore
	\[
		f(P_1) + f(P_2) = (I - m) + \tfrac{(B + 2) + 2m}{2} - 2 = I +
		\tfrac{B}{2} - 1 = f(P).
	\]
	To reduce to primitive triangles, we first cut the polygon into lattice
	triangles by repeatedly drawing interior diagonals between non-adjacent
	vertices. To justify this, we show that every simple
	polygon with at least four vertices has an interior diagonal. At a
	convex-hull vertex $v$, the interior angle is less than $180^{\circ}$. If the
	triangle $T$ spanned by $v$ and its neighbors $u,w$ contains no other
	vertex, then $uw$ is an interior diagonal. Otherwise a vertex
	$z\notin\{u,v,w\}$ of the polygon lying in the closed triangle $T$ and
	maximizing the distance to the line $uw$ makes $vz$ an interior diagonal.
	Here $z$ is allowed to lie on the segment $uw$, which covers the case in
	which all polygon vertices in $T$ other than $u,v,w$ are collinear with
	$u$ and $w$, and no polygon vertex lies in the open interior of $T$. Moreover, $z$ cannot lie in the
	relative interior of the sides $vu$ or $vw$, since a vertex there would
	make the boundary meet itself, contrary to simplicity. To justify the diagonal construction, note that a polygon
	edge meeting the interior of $vz$ cannot cross the sides $vu$ or $vw$ of
	the simple polygon, so it would end at a vertex in $T$ strictly on $v$'s
	side of the line through $z$ parallel to $uw$, contradicting the choice of
	$z$. Thus induction on the number of vertices reduces the polygon to
	triangles.

	We now subdivide any lattice triangle still containing a lattice point
	$z$ other than its vertices: if $z$ is on an edge, join it to the opposite
	vertex, and if $z$ is interior and the triangle has vertices
	$V_1,V_2,V_3$, cut first along the lattice path
	$V_1 \to z \to V_2$ and then cut the remaining quadrilateral along the lattice
	chord $z \to V_3$; together these cuts join $z$ to all three vertices. In
	either case, every resulting triangle has strictly fewer lattice
	points in total than the original triangle, so induction on the number of lattice
	points shows that the subdivision process terminates in primitive
	triangles. A primitive triangle has $I = 0$ and $B = 3$, so $f =
		\tfrac{1}{2}$; and its area is also $\tfrac{1}{2}$: after translating a vertex to
	the origin and denoting the other two by $u,v$, we find that the area is
	$\tfrac{1}{2}|\det(u,v)|$, so it suffices to show that $\{u,v\}$ is a
	$\mathbb{Z}$-basis of $\mathbb{Z}^2$. If not, let
	$L=\mathbb{Z}u+\mathbb{Z}v$, a proper sublattice of $\mathbb{Z}^2$, and
	choose $w\in\mathbb{Z}^2\setminus L$. Writing $w=su+tv$ with rational
	$s,t$ and subtracting suitable integer multiples of $u$ and $v$, we
	may assume $0\le s,t<1$ and $(s,t)\neq(0,0)$; then
	$p:=su+tv\in\mathbb{Z}^2$ lies in the half-open fundamental
	parallelogram of $L$. If $s+t\le1$ then $p$ is a non-vertex lattice
	point of the triangle $\{0,u,v\}$, while if $s+t>1$ then
	$u+v-p=(1-s)u+(1-t)v$ is one; either case contradicts
	primitivity. Hence $|\det(u,v)|=1$ and the area equals $f=\tfrac{1}{2}$.
	Summing over the triangulation gives $S = f(P) = I + \tfrac{B}{2} -
		1$.

\end{proof}
\subsection{The Gauss Circle Problem}

\begin{theorem}\label{thm:10.9.1} Let $A(r) =
		\#\{(x,y)\in\mathbb{Z}^2 \mid x^2 + y^2 \le r^2\}$ for $r>0$. Then $A(r)
		= 1 + 4\lfloor r\rfloor + 4\sum_{i=1}^{\lfloor r\rfloor}\big\lfloor
		\sqrt{r^2 - i^2}\,\big\rfloor$, and $A(r) \sim \pi r^2$ as $r\to\infty$.

\end{theorem}

\begin{proof} For the exact formula, count $(x,y)$ with
	$x^2+y^2\le r^2$ by columns: for $x=0$ there are $2\lfloor
		r\rfloor+1$ admissible values of $y$, and for each $x$ with $1\le x\le\lfloor
		r\rfloor$ there are $2\big\lfloor\sqrt{r^2-x^2}\big\rfloor+1$ values
	of $y$, and likewise for $-x$. Summing gives $A(r)=(2\lfloor
		r\rfloor+1)+2\sum_{x=1}^{\lfloor
			r\rfloor}\big(2\lfloor\sqrt{r^2-x^2}\rfloor+1\big)=1+4\lfloor
		r\rfloor+4\sum_{x=1}^{\lfloor
			r\rfloor}\big\lfloor\sqrt{r^2-x^2}\big\rfloor$. For the asymptotic relation,
	assign to each counted lattice point the unit square
	$[x,x+1)\times[y,y+1)$. These squares are disjoint with total area
	$A(r)$. Each square lies within distance $\sqrt2$ of its lower-left
	corner, so their union lies in the closed disk centered at the origin
	of radius $r+\sqrt2$. It also contains the closed disk centered at
	the origin of radius $\max\{r-\sqrt2,0\}$.
	Indeed, write $\lVert z\rVert$ for the Euclidean distance from a point
	$z\in\mathbb{R}^2$ to the origin (this is not the quantity of
	Definition~\ref{def:1.2.4}). If $r\ge\sqrt2$ and $\lVert z\rVert\le r-\sqrt2$, let $v$
	be the lower-left lattice point of the half-open unit square containing
	$z$. Then $\lVert v\rVert\le\lVert z\rVert+\sqrt2\le r$, so that
	square is one of those counted. If $0<r<\sqrt2$, the inner disk is
	$\{0\}$, which lies in the counted square $[0,1)\times[0,1)$.
	Hence
	$\pi\big(\max\{r-\sqrt{2}, 0\}\big)^2\le A(r)\le\pi(r+\sqrt{2})^2$, so $A(r)/r^2\to\pi$,
	that is, $A(r)\sim\pi r^2$.

\end{proof}
\subsection{Regular Lattice Polygons}

\begin{definition}[regular polygon]\label{def:10.10.1} Identify $\mathbb{R}^2$
	with $\mathbb{C}$ via $(x,y)\mapsto x+iy$, so that the lattice points of
	$\mathbb{Z}^2$ are exactly the Gaussian integers $\mathbb{Z}[i]$
	(Definition~\ref{def:8.2.1}). For an integer $n\ge3$, put
	$\zeta_n=e^{2\pi i/n}$. A regular $n$-gon is a list of $n$ points
	\[
		P_k=c+w\zeta_n^{\,k}\qquad(0\le k\le n-1)
	\]
	with $c\in\mathbb{C}$ and $w\in\mathbb{C}\setminus\{0\}$; the points $P_k$
	are its vertices, its centre is $c$,
	and the rotation $z\mapsto c+\zeta_n(z-c)$ about $c$ sends each $P_k$ to
	$P_{k+1}$, with indices read modulo $n$.

\end{definition}
\begin{theorem}\label{thm:10.10.1} If all vertices of a regular $n$-gon
	($n\ge3$) are lattice points, then $n=4$. Conversely, the unit square, with
	vertices $0,1,1+i,i$, is a regular $4$-gon with lattice vertices.

\end{theorem}

\begin{proof} Let $P_k=c+w\zeta_n^{\,k}$ be lattice points. Put
	$u=P_1-P_0=w(\zeta_n-1)$ and $v=P_2-P_1=w\zeta_n(\zeta_n-1)=\zeta_n u$. Both
	are Gaussian integers, and $u\neq0$ because $w\neq0$ and $\zeta_n\neq1$.
	Hence $\zeta_n=v/u$ is a quotient of Gaussian integers. By unique
	factorization in $\mathbb{Z}[i]$ (Theorem~\ref{thm:8.2.2}), we may cancel
	common prime factors and write $\zeta_n=\alpha/\beta$ with
	$\alpha,\beta\in\mathbb{Z}[i]$, $\beta\neq0$, and no Gaussian prime dividing
	both $\alpha$ and $\beta$. From $\zeta_n^n=1$ we get $\alpha^n=\beta^n$. If a
	Gaussian prime $\pi$ divided $\beta$, it would divide $\alpha^n$, hence
	$\alpha$, which is impossible; so $\beta$ is a unit and
	$\zeta_n\in\mathbb{Z}[i]$. Then $N(\zeta_n)$ is a positive integer with
	$N(\zeta_n)^n=N(\zeta_n^n)=1$ (Theorem~\ref{thm:8.2.1}), so $N(\zeta_n)=1$
	and $\zeta_n\in\{\pm1,\pm i\}$ by Theorem~\ref{thm:8.2.2}. In particular
	$\zeta_n^4=1$. Since $e^{2\pi ik/n}=1$ exactly when $n\mid k$, the element
	$\zeta_n$ of $\mathbb{C}^{\times}$ has order $n$, so $n\mid4$. With $n\ge3$,
	this forces $n=4$.

	Conversely, with $\zeta_4=i$, $c=\frac{1+i}{2}$, and $w=-\frac{1+i}{2}$, the
	points $c+wi^k$ for $k=0,1,2,3$ are $0,1,1+i,i$.

\end{proof}
\subsection{Minkowski's Convex-Body Theorem}

\begin{theorem}[Minkowski's theorem, planar case]\label{thm:10.11.1} Let
	$\Lambda\subseteq\mathbb{R}^2$ be a planar lattice with covolume
	$d(\Lambda)$ (Definition~\ref{def:10.8.1}), and let $K\subseteq\mathbb{R}^2$
	be convex and symmetric about the origin ($K=-K$), with area
	$\area(K)>4\,d(\Lambda)$. Then $K$ contains a nonzero
	point of $\Lambda$. If $K$ is in addition compact,
	$\area(K)\ge 4\,d(\Lambda)$ already suffices. In particular, taking
	$\Lambda=\mathbb{Z}^2$, for which $d(\Lambda)=1$, we see that a convex region in
	$\mathbb{R}^2$ that is symmetric about the origin and has area greater than $4$
	contains a nonzero integer point $(a,b)\in\mathbb{Z}^2$.

\end{theorem}

\begin{proof} Write $\Lambda=A\mathbb{Z}^2$ for an
	invertible linear map $A$ whose columns are a basis of $\Lambda$, so
	$|\det A|=d(\Lambda)$ by Theorem~\ref{thm:10.8.1}. Then $A^{-1}$ carries $\Lambda$ onto
	$\mathbb{Z}^2$ and $K$ onto the convex symmetric set $K'=A^{-1}K$ of
	area $\area(K)/d(\Lambda)>4$, and points of $\Lambda$
	in $K$ correspond to integer points in $K'$; so it suffices to treat
	$\Lambda=\mathbb{Z}^2$ with $\area(K)>4$. Let
	$S=\tfrac{1}{2}K$, which is convex and symmetric and has area
	$\area(K)/4>1$. Partition the plane into the unit cells
	$Q_{m,n}=[m,m+1)\times[n,n+1)$ for $(m,n)\in\mathbb{Z}^2$, and
	translate each piece $S\cap Q_{m,n}$ by $(-m,-n)$ into the unit
	square $Q_{0,0}$. The images all lie in $Q_{0,0}$, of area $1$; since the pieces $S\cap Q_{m,n}$ partition $S$, countable additivity of area shows that the sum of the areas of the images equals $\area(S)>1$, so if the images were pairwise disjoint, they would fit inside $Q_{0,0}$ with total area at most $1$, a contradiction. Hence two images overlap:
	there are points $x,y\in S$ with $x-y\in\mathbb{Z}^2$ and $x\neq y$,
	that is, $x-y\in\mathbb{Z}^2\setminus\{0\}$. Since $x,y\in\tfrac{1}{2}K$ we
	have $2x,2y\in K$; by symmetry $-2y\in K$, and by convexity the
	midpoint $\tfrac{1}{2}(2x)+\tfrac{1}{2}(-2y)=x-y$ lies in $K$. Thus $x-y$ is
	a nonzero point of $\Lambda$ in $K$. For compact $K$ with
	$\area(K)=4\,d(\Lambda)$ (so $\area(K)=4$ after the reduction above), apply the strict case to
	$(1+\tfrac{1}{m})K$ for each $m\ge1$. Since $K$ is convex and symmetric,
	it contains $0=\tfrac12x+\tfrac12(-x)$ for any $x\in K$, so $tK\subseteq sK$
	whenever $0<t\le s$. Hence the resulting nonzero integer points lie in the
	bounded set $2K$, where each coordinate takes only finitely many values,
	so they are finite in number, and one of them, $z$, occurs for infinitely
	many $m$. For these $m$ we have $z/(1+\tfrac{1}{m})\in K$; these points
	converge to $z$, and since $K$ is closed, $z\in K$.

\end{proof}
\newpage

\section{Polynomials over \texorpdfstring{$\mathbb{Z}$}{Z}: Roots, Irreducibility, and Lifting}

\subsection{Integer and Rational Roots}

\begin{theorem}[rational root theorem]\label{thm:11.1.1} Let $f(x)
		= \sum_{i=0}^n a_i x^i\in\mathbb{Z}[x]$ with $a_n\neq0$. If
	$p\in\mathbb{Z}$ and $q\in\mathbb{Z}^+$ satisfy $\gcd(p,q)=1$ and
	$\tfrac{p}{q}$ is a root, then $p\mid a_0$ and $q\mid a_n$. In
	particular, the rational roots of a monic integer polynomial are
	integers dividing $a_0$.

\end{theorem}

\begin{proof} After denominators are cleared in $f(p/q)=0$, every term
	except $a_0q^n$ is divisible by $p$. Thus $p\mid a_0 q^n$, and
	$\gcd(p, q)=1$ gives $p\mid a_0$. Likewise, every term except $a_np^n$
	is divisible by $q$, so $q\mid a_n$.

\end{proof}
\begin{theorem}\label{thm:11.1.2} Let $f(x) = \sum_{i=0}^{n}a_i
		x^i\in\mathbb{Z}[x]$ and let $\tfrac{p}{q}$ with $\gcd(p, q) = 1$ and
	$q\ge1$ be a rational root of $f$. Then $p - q\mid f(1)$ and $p +
		q\mid f(-1)$.

\end{theorem}

\begin{proof} Since $f(p/q) = 0$,
	\[
		q^{n}f(1)=q^{n}\big(f(1)-f(p/q)\big)=\sum_{i=0}^{n}a_i
		q^{n-i}(q^{i}-p^{i}),
	\]
	and each $q^{i}-p^{i}$ is divisible by $q-p$, so $p-q\mid q^{n}f(1)$.
	Any prime dividing both $q$ and $p-q$ divides $p$, hence
	$\gcd(p-q, q) = 1$ and $p - q\mid f(1)$. In the same way
	\[
		q^{n}f(-1)=q^{n}\big(f(-1)-f(p/q)\big)=\sum_{i=0}^{n}a_i
		q^{n-i}\big((-q)^{i}-p^{i}\big)
	\]
	is divisible by $-q-p$, since $a-b\mid a^{i}-b^{i}$ with $a=-q$ and
	$b=p$; and $\gcd(p+q, q) = \gcd(p, q) = 1$ gives $p + q\mid f(-1)$.

\end{proof}
\subsection{Gauss's Lemma}

\begin{definition}[content and primitivity]\label{def:11.2.1} The
	content $c(f)$ of a nonzero polynomial $f\in\mathbb{Z}[x]$ is the nonnegative gcd of its coefficients; set $c(0)=0$. A nonzero polynomial is primitive if $c(f)=1$.

\end{definition}
\begin{theorem}[Gauss's lemma on primitive polynomials]\label{thm:11.2.1} The product of
	primitive polynomials is primitive; hence $c(fg)=c(f)c(g)$, and a
	primitive $f\in\mathbb{Z}[x]$ that factors over $\mathbb{Q}$ already
	factors into lower-degree polynomials in $\mathbb{Z}[x]$.

\end{theorem}

\begin{proof} Let $f,g$ be primitive and suppose a prime
	$p$ divides every coefficient of $fg$. Reduction modulo $p$ gives $\bar f\,\bar g
		= 0$ in $\mathbb{F}_p[x]$, in which a product of nonzero polynomials is
	nonzero (its leading coefficient is the product of the leading
	coefficients, and $\mathbb{F}_p$ is a field), so $\bar f = 0$ or
	$\bar g = 0$, that is, $p$ divides all coefficients of $f$ or of $g$,
	contradicting primitivity. For the content formula, if $f=0$ or $g=0$,
	both sides of $c(fg)=c(f)c(g)$ are zero. Otherwise, write $f = c(f)f_0$
	and $g = c(g)g_0$ with $f_0, g_0$
	primitive; then $fg = c(f)c(g)\,f_0g_0$ with $f_0g_0$ primitive by the
	first part, so $c(fg) = c(f)c(g)\,c(f_0g_0) = c(f)c(g)$. For the
	factorization statement, suppose a
	primitive $f$ factors as $f=GH$ with $G,H\in\mathbb{Q}[x]$ of
	positive degree. Every nonzero rational polynomial can be normalized as follows: multiply by a common positive denominator to obtain an integer polynomial, then divide by the gcd of its coefficients. Thus write $G=\tfrac{a}{b}G_0$ and $H=\tfrac{c}{d}H_0$ with $G_0,H_0\in\mathbb{Z}[x]$ primitive; then $bd\,f=ac\,G_0H_0$.
	Comparing the contents of both sides gives
	$|bd|=|bd|\,c(f)=|ac|\,c(G_0H_0)=|ac|$ (using the fact that $f$ is primitive
	and, by the first part, so is $G_0H_0$), so $ac/bd=\pm1$ and $f=\pm
		G_0H_0$ is a factorization into integer polynomials of the same
	positive degrees as $G,H$.

\end{proof}
\subsection{Irreducibility Criteria}

\begin{definition}[irreducibility over a field]\label{def:11.3.1}
	Let $F$ be a field. A nonconstant polynomial $f\in F[x]$ is
	irreducible over $F$ if it cannot be written as $f=gh$ with
	$g,h\in F[x]$ both of positive degree.

\end{definition}
\begin{theorem}[Eisenstein's criterion]\label{thm:11.3.1} Let $f = \sum_{i=0}^n
		a_i x^i\in\mathbb{Z}[x]$ be of degree $n\ge1$. If a prime $p$ satisfies $p\nmid a_n$,
	$p\mid a_i$ for all $i<n$, and $p^2\nmid a_0$, then $f$ is
	irreducible over $\mathbb{Q}$.

\end{theorem}

\begin{proof} Put $c=c(f)$ and $f_0=f/c$, so $f_0\in\mathbb{Z}[x]$ is
	primitive. Since $c\mid a_n$ and $p\nmid a_n$, we have $p\nmid c$.
	Division by $c$ therefore preserves each of the three Eisenstein
	conditions at $p$: the leading coefficient is still not divisible by $p$,
	the lower coefficients remain divisible by $p$, and the constant term
	is still not divisible by $p^2$. Multiplication by the nonzero rational
	constant $c$ does not affect reducibility over $\mathbb{Q}$. We may
	therefore replace $f$ by $f_0$ (retaining the notation $f=\sum a_i x^i$)
	and assume that $f$ is primitive. By Gauss's lemma
	(Theorem~\ref{thm:11.2.1}) it now suffices to rule out a factorization $f=gh$ in
	$\mathbb{Z}[x]$ with $\deg g,\deg h\ge1$. Now reduce modulo $p$. Since $p\nmid a_n$ and $a_n$ is the product of
	the leading coefficients of $g$ and $h$, neither leading
	coefficient is divisible by $p$, so $\deg\bar g=\deg g\ge1$ and
	$\deg\bar h=\deg h\ge1$. Moreover, $\bar f = \bar a_n x^n = \bar g\bar h$
	in $\mathbb{F}_p[x]$. Since $\mathbb{F}_p$ is a field, a product of nonzero polynomials over $\mathbb{F}_p$ is nonzero, so the lowest-degree nonzero term of a product is the product of the lowest-degree nonzero terms of the factors. As $\bar g\bar h=\bar a_n x^n$ has its only nonzero term in degree $n=\deg\bar g+\deg\bar h$, the lowest terms of $\bar g$ and $\bar h$ must lie in degrees $\deg\bar g$ and $\deg\bar h$, respectively; hence both are monomials $\bar g = \bar b\,x^{r}$ and $\bar h = \bar c\,x^{s}$ with $\bar b\bar c=\bar a_n$
	and $r+s=n$, where $r=\deg\bar g\ge1$ and $s=\deg\bar h\ge1$.
	Hence the constant terms of $\bar g$ and $\bar h$ vanish, that is, $p\mid
		g(0)$ and $p\mid h(0)$, so $a_0 = g(0)h(0)$ is divisible by $p^2$,
	contradicting $p^2\nmid a_0$.

\end{proof}
\begin{theorem}\label{thm:11.3.2} Let $f\in\mathbb{Z}[x]$ and let
	$p$ be a prime not dividing the leading coefficient of $f$. If the
	reduction $\bar f\in\mathbb{F}_p[x]$ of $f$ modulo $p$ is irreducible
	over $\mathbb{F}_p$, then $f$ is irreducible over $\mathbb{Q}$.

\end{theorem}

\begin{proof} Suppose $f = GH$ with $G,
		H\in\mathbb{Q}[x]$ and $\deg G, \deg H\ge1$. Write $f = c(f)f_0$
	with $f_0$ primitive and normalize $G = \tfrac{a}{b}G_0$, $H =
		\tfrac{c}{d}H_0$ with $G_0, H_0\in\mathbb{Z}[x]$ primitive, as in
	Theorem~\ref{thm:11.2.1}. Then $bd\,c(f)f_0 = ac\,G_0H_0$; comparing
	the contents of both sides (using Theorem~\ref{thm:11.2.1} for the
	primitivity of $G_0H_0$) gives $|bd|c(f) = |ac|$. Thus
	$ac/(bd)=\varepsilon c(f)$ for some $\varepsilon\in\{1,-1\}$, so
	$f=\big(\varepsilon c(f)G_0\big)H_0$ is a factorization in
	$\mathbb{Z}[x]$ into factors of the same positive degrees as
	$G, H$. Since $p$ does not divide the leading coefficient of $f$, it
	divides neither the leading coefficient of $\varepsilon c(f)G_0$ nor that of $H_0$, so both
	factors keep their positive degrees after reduction modulo $p$, and
	$\bar f$ factors nontrivially in $\mathbb{F}_p[x]$, a
	contradiction.

\end{proof}
\begin{theorem}\label{thm:11.3.3}\statementitems
	\begin{enumerate}[label={\alph*)}]
		\item For $f\in\mathbb{Z}[x]$ and $c\in\mathbb{Z}$, the polynomial $f$ is
		      irreducible over $\mathbb{Q}$ if and only if $f(x+c)$ is
		      irreducible over $\mathbb{Q}$.
		\item For every prime $p$, the $p$-th cyclotomic polynomial
		      $\Phi_p(x)=x^{p-1}+x^{p-2}+\cdots+1$
		      (Definition~\ref{def:6.11.1}) is irreducible over $\mathbb{Q}$.
	\end{enumerate}

\end{theorem}

\begin{proof} (a) If $f$ is constant, including the zero polynomial, then
	$f(x+c)$ is also constant; neither is irreducible under
	Definition~\ref{def:11.3.1}. For nonconstant $f$, a factorization $f=GH$
	over $\mathbb{Q}$ with $\deg G,\deg H\ge1$ gives
	$f(x+c)=G(x+c)H(x+c)$ with the same positive degrees for the factors.
	Shifting back by $-c$ gives the converse. (b) From
	$x^{p} - 1 = (x - 1)\Phi_p(x)$ (Theorem~\ref{thm:6.11.1}),
	\[
		\Phi_p(x+1)=\frac{(x+1)^{p}-1}{x}=\sum_{k=1}^{p}\binom{p}{k}x^{k-1}.
	\]
	The leading coefficient is $\binom{p}{p} = 1$; for $1\le k\le p-1$
	the coefficient $\binom{p}{k}$ is divisible by $p$ (Theorem~\ref{thm:3.5.2}); and the constant term is $\binom{p}{1} = p$, which is not
	divisible by $p^{2}$. Hence $F(x):=\Phi_p(x+1)$ is Eisenstein at $p$
	(Theorem~\ref{thm:11.3.1}) and is irreducible over $\mathbb{Q}$.
	Applying part (a) to $F$ with $c=-1$ shows that
	$F(x-1)=\Phi_p(x)$ is irreducible over $\mathbb{Q}$.

\end{proof}
\begin{theorem}\label{thm:11.3.4} For every $n\in\mathbb{Z}^+$, the
	cyclotomic polynomial $\Phi_n(x)$ (Definition~\ref{def:6.11.1}) is
	irreducible over $\mathbb{Q}$. Consequently, $x^{n} - 1 = \prod_{d\mid
			n}\Phi_d(x)$ (Theorem~\ref{thm:6.11.1}) is the factorization of
	$x^{n} - 1$ into monic irreducibles in $\mathbb{Q}[x]$.

\end{theorem}

\begin{proof} The case of prime $n$ is Theorem~\ref{thm:11.3.3}(b); the general case follows by reduction modulo
	primes not dividing $n$. Recall that $\Phi_n\in\mathbb{Z}[x]$ is monic of
	degree $\varphi(n)$ (Theorem~\ref{thm:6.11.2}). Put $\zeta = e^{2\pi
				i/n}$ and let $f\in\mathbb{Q}[x]$ be the minimal polynomial of
	$\zeta$: the monic polynomial of least degree with $\zeta$ as a root.
	It exists ($\Phi_n(\zeta) = 0$), it is unique (the difference of two
	such polynomials has smaller degree and still vanishes at $\zeta$),
	and it divides every polynomial in $\mathbb{Q}[x]$ vanishing at
	$\zeta$ (divide and evaluate the remainder at $\zeta$). It is also
	irreducible: if $f=f_1f_2$ with both factors of smaller degree, then
	$f_1(\zeta)f_2(\zeta)=0$ forces one of them to vanish at $\zeta$,
	contrary to minimality. Thus $f$ is the minimal polynomial of each of
	its roots. In particular, $\Phi_n=fh$ in $\mathbb{Q}[x]$, and $h$ is monic.
	Write $f=F/a$ and $h=H/b$, where $F,H\in\mathbb{Z}[x]$ are
	primitive with positive leading coefficients $a,b\in\mathbb{Z}^+$,
	respectively. Then $ab\Phi_n=FH$. Since $\Phi_n$ is monic, it is
	primitive. Comparing contents by Theorem~\ref{thm:11.2.1} gives
	$ab=c(ab\Phi_n)=c(FH)=1$. Thus $f,h\in\mathbb{Z}[x]$, both monic.

	We claim that for every prime $p\nmid n$ and every root $\pi$ of $f$,
	the power $\pi^{p}$ is again a root of $f$. Otherwise, since $\pi^{p}$ is a
	root of $\Phi_n = fh$ (it is a primitive $n$-th root of unity:
	$\gcd(p, n) = 1$ preserves the order $n$ of $\pi$), it would be a
	root of $h$, so $\pi$ would be a root of $h(x^{p})$; minimality
	gives $h(x^{p})=f(x)q(x)$ in $\mathbb{Q}[x]$. Since $f$ is monic,
	division of $h(x^p)$ by $f$ can be performed in $\mathbb{Z}[x]$.
	Uniqueness of the quotient and remainder over $\mathbb{Q}[x]$ therefore
	gives $q\in\mathbb{Z}[x]$. Reduce modulo $p$:
	in $\mathbb{F}_p[x]$ the binomial theorem and
	Theorem~\ref{thm:3.5.2} give $(u+v)^p=u^p+v^p$; applying this repeatedly
	and using Fermat's little theorem (Theorem~\ref{thm:3.3.2}) on the
	coefficients gives $\bar h(x^{p}) = \bar h(x)^{p}$. Hence
	$\bar f\mid\bar h^{p}$. We use two facts about $\mathbb{F}_p[x]$.
	First, every nonconstant polynomial $F$ has an irreducible factor: a
	divisor of $F$ of least positive degree is irreducible, since a
	factorization of it into two factors of positive degree would give a
	divisor of smaller positive degree. Second, if an irreducible $g$
	divides $AB$ but not $A$, then $g\mid B$: indeed $A\neq0$, so among the
	nonzero polynomials $uA+vg$ with $u,v\in\mathbb{F}_p[x]$ we may choose
	one, $r$, of least degree. Division of $g$ and of $A$ with remainder by $r$, possible
	because the leading coefficient of $r$ is invertible in $\mathbb{F}_p$,
	leaves remainders of the same form and of smaller degree, so by the
	choice of $r$ these remainders are $0$ and $r$ divides both $g$ and $A$. As $g$ is
	irreducible, either $r$ is a nonzero constant or $g$ is a constant
	multiple of $r$; the latter would give $g\mid A$. Hence $r$ is a nonzero
	constant, and multiplying $r=uA+vg$ by $r^{-1}B$ shows that
	$g\mid B$. By the first fact, $\bar f$, which is monic of positive
	degree, has an irreducible factor $g$; then $g\mid\bar h^{p}$, and
	repeated use of the second fact gives $g\mid\bar h$. Then $g^2$
	divides $\bar f\,\bar h = \bar\Phi_n$, which divides $x^{n} - 1$ in
	$\mathbb{F}_p[x]$; but a repeated irreducible factor of $x^{n} - 1$
	also divides its derivative $nx^{n-1}$. Since $p\nmid n$, the constant
	$n$ is invertible in $\mathbb{F}_p$, so $g\mid x^{n-1}$, and the second
	fact gives $g\mid x$. Thus $g$ is a nonzero constant multiple of $x$,
	and $x$ would divide $x^n-1$, which is false because the constant
	term of $x^n-1$ is $-1$; this is a contradiction.

	Now let $m\ge1$ with $\gcd(m, n) = 1$ and factor $m = p_1\cdots p_r$
	into primes; each $p_j\nmid n$. Applying the claim successively to
	the roots $\zeta, \zeta^{p_1}, \zeta^{p_1p_2}, \ldots$ of $f$ shows
	that $\zeta^{m}$ is a root of $f$. Hence every primitive $n$-th root
	of unity is a root of $f$, so $\deg f\ge\varphi(n) = \deg\Phi_n$.
	Since $f\mid\Phi_n$ and both are monic, $f = \Phi_n$: the cyclotomic
	polynomial is irreducible over $\mathbb{Q}$. The concluding
	statement follows at once: each factor $\Phi_d$ with $d\mid n$ is
	irreducible and monic.

\end{proof}
\subsection{Hensel's Lemma}

\begin{theorem}[Hensel's lemma]\label{thm:11.4.1} Let
	$f\in\mathbb{Z}[x]$, let $p$ be prime, and let $a\in\mathbb{Z}$. If
	$f(a)\equiv0\pmod{p}$ and
	$f'(a)\not\equiv0\pmod{p}$, then for every $k\ge1$ there is a unique
	$a_k\pmod{p^k}$ with $a_k\equiv a\pmod{p}$ and $f(a_k)\equiv0\pmod{p^k}$.

\end{theorem}

\begin{proof} For $k=1$, take $a_1\equiv a\pmod p$; this is the unique
	class satisfying the two required congruences. Assume that $a_k$ has
	been constructed; we seek $a_{k+1}=a_k+tp^k$.
	For $x,h\in\mathbb{Z}$, expanding each power $(x+h)^i$ by the
	binomial theorem shows that
	$f(x+h)-f(x)-hf'(x)$ is divisible by $h^2$. Applied with
	$h=tp^k$ (and $2k\ge k+1$), this gives
	$f(a_k + tp^k)\equiv f(a_k) +
		t p^k f'(a_k)\pmod{p^{k+1}}$. With $f(a_k) = p^k s$, the
	requirement $f(a_{k+1})\equiv0\pmod{p^{k+1}}$ becomes $s + t
		f'(a_k)\equiv0\pmod{p}$, which has a unique solution $t\pmod{p}$ since
	$f'(a_k)\equiv f'(a)\not\equiv0\pmod{p}$ is invertible. This proves existence. For uniqueness, if $b_{k+1}$ is any other root congruent to $a$ modulo $p$, then its reduction $b_k$ modulo $p^k$ is a root with the same initial class; by the induction hypothesis, $b_k\equiv a_k\pmod{p^k}$. Hence $b_{k+1}=a_k+tp^k$ for some $t\pmod{p}$, and the uniqueness of the solution modulo $p$ to the linear congruence above forces the same $t$ modulo $p$. Thus the lift is globally unique.

	Theorem~\ref{thm:7.8.3} is the case $f=x^n-c$ of this lemma, and the
	odd-prime lifting step in Theorem~\ref{thm:7.8.4}(a) is the case $f=x^2-c$ of the same lemma:
	for $f=x^n-c$ with $n\ge2$, the hypotheses $p\nmid n$ and
	$p\nmid c$ say exactly that $f'(x) = n x^{n-1}$ is nonzero at the root
	modulo $p$. The $p = 2$ case of Theorem~\ref{thm:7.8.4} is not covered by this
	lemma: there $f'(x) = 2x \equiv 0 \pmod{2}$ for every $x$, so the derivative
	condition fails identically, and the criterion $c\equiv1\pmod8$ is what
	remains in that singular case. The derivative condition cannot be
	dropped: for $f(x)=x^p-(1+p)$ and $a=1$, one has
	$f(1)=-p\equiv0\pmod p$ and $f'(1)=p\equiv0\pmod p$, but every possible
	lift has the form $x=1+pt$ and satisfies $x^p\equiv1\not\equiv1+p\pmod{p^2}$ by
	the binomial theorem.

\end{proof}

\begin{theorem}[Hensel's lemma, strong form]\label{thm:11.4.2} Let $f\in\mathbb{Z}[x]$, let $p$
	be prime, let $a\in\mathbb{Z}$ with $f'(a)\neq0$, and write
	$s = v_p(f'(a))<\infty$ (Definition~\ref{def:4.5.1}). If $v_p(f(a))\ge2s + 1$, then for every $m\ge2s + 1$
	there is a unique $a_m\pmod{p^{m-s}}$ with $a_m\equiv a\pmod{p^{s+1}}$
	and $f(a_m)\equiv0\pmod{p^{m}}$.

\end{theorem}

\begin{proof} For $x,h\in\mathbb{Z}$, expanding each power $(x+h)^i$ by
	the binomial theorem shows directly that
	\[
		h^2\mid f(x+h)-f(x)-hf'(x).
	\]
	Note also that $b
		\equiv a\pmod{p^{s+1}}$ implies $v_p(f'(b)) = s$: indeed $f'(b)\equiv
		f'(a)\pmod{p^{s+1}}$ term by term, and $v_p(f'(a)) = s < s + 1$.
	Consequently, for $m\ge2s+1$ and $b\equiv a\pmod{p^{s+1}}$,
	replacing $b$ by $b+up^{m-s}$ does not change $f(b)$ modulo $p^m$:
	the linear term has valuation at least $m$, and
	$2(m-s)\ge m$. Thus the root condition is well defined on the
	stated residue classes modulo $p^{m-s}$.

	For $m = 2s + 1$, take $a_m = a$; then $v_p(f(a_m))\ge m$ by
	hypothesis, and the uniqueness statement is modulo $p^{m-s} =
		p^{s+1}$, which is exactly the modulus of the prescribed congruence class. For the
	inductive step, let $m\ge2s + 1$ and suppose $a_m$ satisfies the required conditions; we seek
	$a_{m+1} = a_m + tp^{m-s}$ with $t\in\{0, 1, \ldots, p-1\}$. Since
	$2(m - s)\ge m + 1$ exactly when $m\ge2s + 1$, the divisibility statement
	gives
	\[
		f(a_m+tp^{m-s})\equiv f(a_m)+tp^{m-s}f'(a_m)\pmod{p^{m+1}}.
	\]
	Write $f(a_m) = p^{m}u$. Because $a_m\equiv a\pmod{p^{s+1}}$, we have
	$v_p(f'(a_m)) = s$, so $f'(a_m) = p^{s}w$ with $p\nmid w$, and the
	requirement $f(a_{m+1})\equiv0\pmod{p^{m+1}}$ becomes $u + tw\equiv0
		\pmod{p}$, which has a unique solution $t\pmod{p}$. Moreover,
	$a_{m+1}\equiv a_m\pmod{p^{m-s}}$ and $m - s\ge s + 1$, so
	$a_{m+1}\equiv a\pmod{p^{s+1}}$. For uniqueness, any $b\equiv a
		\pmod{p^{s+1}}$ with $f(b)\equiv0\pmod{p^{m+1}}$ also satisfies
	$f(b)\equiv0\pmod{p^{m}}$, so the inductive hypothesis gives $b\equiv
		a_m\pmod{p^{m-s}}$; write $b = a_m + tp^{m-s}$. The same linear
	congruence uniquely determines $t\pmod{p}$, whence $b\equiv a_{m+1}
		\pmod{p^{m+1-s}}$.

	For $f=x^2-c$ at $a=1$ with $p=2$, one has $s=v_2(2)=1$, and the
	hypothesis $v_2(f(1))\ge3$ is exactly $c\equiv1\pmod8$. Thus the strong
	form recovers the lifting and existence direction of the $p=2$ criterion
	in Theorem~\ref{thm:7.8.4}, which
	Theorem~\ref{thm:11.4.1} does not cover.

\end{proof}
\newpage

\section{Additive Number Theory and Partitions}

\subsection{Polynomial Tools for Additive Problems}

\begin{definition}[partial degree]\label{def:12.1.1} Let $F$ be a field and let
	$f\in F[x_1,\ldots,x_n]$ with $n\in\mathbb{Z}^+$. For nonzero $f$, the
	partial degree $\deg_{x_k}f$ ($1\le k\le n$) is the largest exponent of
	$x_k$ occurring in a monomial of $f$ with nonzero coefficient; set
	$\deg_{x_k}0=-\infty$.

\end{definition}
\begin{theorem}\label{thm:12.1.1} Let $n\in\mathbb{Z}^+$, let $F$ be a
	field, and let $S_1,\ldots,S_n\subseteq F$ be finite. If
	$g\in F[x_1,\ldots,x_n]$ satisfies $\deg_{x_k}g<|S_k|$ for every
	$1\le k\le n$ and vanishes on $S_1\times\cdots\times S_n$, then $g=0$.

\end{theorem}

\begin{proof} We induct on $n$. For $n=1$, a nonzero
	polynomial of degree less than $|S_1|$ has fewer than $|S_1|$ roots, so it cannot
	vanish on $S_1$. For $n>1$, the case $g=0$ is immediate, so assume
	$g\neq0$ and write
	$g=\sum_{r=0}^{d}g_r(x_1,\ldots,x_{n-1})x_n^{r}$ with $d<|S_n|$. For
	each fixed $(s_1,\ldots,s_{n-1})\in S_1\times\cdots\times S_{n-1}$,
	the one-variable polynomial $g(s_1,\ldots,s_{n-1},x_n)$ has degree less than $|S_n|$ and
	vanishes on $S_n$, hence is zero, so every
	$g_r(s_1,\ldots,s_{n-1})=0$. For $1\le j<n$, every coefficient
	polynomial satisfies $\deg_{x_j}g_r\le\deg_{x_j}g<|S_j|$. Thus each
	$g_r$ vanishes on the smaller grid and, by induction, $g_r=0$. Hence
	$g=0$.

\end{proof}
\begin{theorem}[Combinatorial Nullstellensatz]\label{thm:12.1.2} Let
	$n\in\mathbb{Z}^+$, let $F$ be a field, let
	$t_i\in\mathbb{Z}_{\ge0}$ for $1\le i\le n$, and let
	$f\in F[x_1,\ldots,x_n]$ have total degree
	$\sum_{i=1}^n t_i$. Suppose that the coefficient of $\prod_{i=1}^{n}x_i^{t_i}$
	is nonzero. If $S_1,\ldots,S_n\subseteq F$ with $|S_i| > t_i$, then
	$f(s_1,\ldots,s_n)\neq0$ for some
	$(s_1,\ldots,s_n)\in\prod_{i=1}^{n}S_i$.

\end{theorem}

\begin{proof} For each $i$, choose a subset
	$S_i'\subseteq S_i$ with $|S_i'|=t_i+1$. It suffices to find a point of $\prod_iS_i'$ at which the value is nonzero, so we replace $S_i$ by $S_i'$; in
	particular, all the sets are now finite. For each $i$, put
	\[
		g_i(x_i)=\prod_{u\in S_i}(x_i-u).
	\]
	This polynomial is monic of degree $|S_i|$. Perform successive polynomial divisions, one variable at a time: first divide by $g_1(x_1)$ while treating the other variables as coefficients, then by $g_2(x_2)$, and so on. Each division is possible because $g_i$ is monic: while some monomial has $x_i$-degree $r\ge|S_i|$, subtract $c\,x_i^{r-|S_i|}g_i(x_i)$, where $c$ is the coefficient of $x_i^{r}$ (a polynomial in the other variables), and no coefficient is ever divided. Since the coefficients of $g_i$ lie in $F$, these steps do not raise the degree in any other variable, so the bounds $\deg_{x_j}<|S_j|$ obtained for $j<i$ are preserved. The result is a polynomial $\bar f$ such that
	\[
		f-\bar f=\sum_{i=1}^{n}q_i(x_1,\ldots,x_n)g_i(x_i),\qquad \deg_{x_i}\bar f<|S_i|\quad(1\le i\le n).
	\]
	At every point $(s_1,\ldots,s_n)\in\prod_iS_i$, each $g_i(s_i)=0$, so $f(s_1,\ldots,s_n)=\bar f(s_1,\ldots,s_n)$.

	Assume, for contradiction, that $f$ vanishes on the whole grid. Then $\bar f$ also vanishes there, so Theorem~\ref{thm:12.1.1} gives $\bar f=0$. On the other hand, the coefficient of $x_1^{t_1}\cdots x_n^{t_n}$ is unchanged by the reductions. Indeed, $t_i<|S_i|$, so this monomial itself is never reduced. Every actual reduction replaces a power $x_i^r$ with $r\ge|S_i|$ by terms of strictly smaller $x_i$-degree and therefore strictly smaller total degree. Since every monomial of $f$ has total degree at most $\sum_i t_i$, no reduced monomial can produce the monomial $x_1^{t_1}\cdots x_n^{t_n}$ of total degree exactly $\sum_i t_i$. Its coefficient in $\bar f$ is therefore the same nonzero coefficient as in $f$, contradicting $\bar f=0$. Hence $f$ is nonzero at some point of the grid.

\end{proof}
\subsection{Sumsets and the Cauchy--Davenport Theorem}

\begin{definition}[sumsets and iterated sumsets]\label{def:12.2.1}
	An additive commutative monoid is a set $M$ with an associative and
	commutative operation $+$ and an identity element $0$. For subsets
	$A,B\subseteq M$, define
	\[
		A+B=\{a+b \mid a\in A,\ b\in B\}.
	\]
	For $h\in\mathbb{Z}^+$, the iterated sumset $hA$ is the sum of
	$h$ copies of $A$; thus $1A=A$ and $(h+1)A=hA+A$.

\end{definition}
\begin{theorem}[Cauchy--Davenport theorem]\label{thm:12.2.1} For a prime $p$
	and nonempty $A,B\subseteq\mathbb{Z}_p$, we have $|A+B|\ge\min\{p, \, |A|+|B|-1\}$.

\end{theorem}

\begin{proof} If $|A|+|B|>p$, then for every
	$c\in\mathbb{Z}_p$, the sets $A$ and $c-B$ have sizes summing to more
	than $p$, so they intersect and $c\in A+B$; hence $A+B=\mathbb{Z}_p$
	and $|A+B|=p$. Otherwise, $|A|+|B|\le p$; suppose, for contradiction,
	that $|A+B|\le|A|+|B|-2$. Pick $C\supseteq A+B$ with
	$|C|=|A|+|B|-2<p$ and set $f(x,y)=\prod_{c\in C}(x+y-c)$, of total
	degree $|A|+|B|-2$. The coefficient of $x^{|A|-1}y^{|B|-1}$ is
	$\binom{|A|+|B|-2}{|A|-1}$, nonzero in $\mathbb{Z}_p$ since
	$|A|+|B|-2<p$. With $S_1=A$ and $S_2=B$, Theorem~\ref{thm:12.1.2} gives $(a,b)\in
		A\times B$ with $f(a,b)\neq0$; but $a+b\in A+B\subseteq C$ forces a
	factor to vanish, a contradiction.

\end{proof}
\begin{altproof}[Alternative proof of Theorem~\ref{thm:6.10.3}] Let
	$A = \{x^{2} \bmod p \mid 0 \le x \le \tfrac{p-1}{2}\} \subseteq
		\mathbb{Z}_p$. If $x^2\equiv y^2\pmod{p}$ with
	$0\le x,y\le\tfrac{p-1}{2}$, then $p\mid(x-y)(x+y)$ and the
	given bounds force $x=y$; hence $A$ is nonempty and
	$|A|=\tfrac{p+1}{2}$. The Cauchy--Davenport theorem (Theorem~\ref{thm:12.2.1}) gives
	\[
		|A + A| \ge \min\{p,\ |A| + |A| - 1\} = \min\{p,\ p\} = p,
	\]
	so $A + A = \mathbb{Z}_p$. In particular, $-1 \in A + A$, that is,
	$-1 \equiv x^{2} + y^{2} \pmod{p}$ for some $x, y$, which is the
	assertion.

\end{altproof}
\begin{theorem}[Erd{\H{o}}s--Heilbronn theorem]\label{thm:12.2.2} For a
	prime $p$ and $A\subseteq\mathbb{Z}_p$ with $|A|\ge2$,
	\[
		\#\{a+b \mid a, b\in A,\ a\neq b\}\ge\min\{p,\ 2|A|-3\}.
	\]

\end{theorem}

\begin{proof} Write $k = |A|$ and let $R = \{a +
		b \mid a, b\in A,\ a\neq b\}$ be the restricted sumset.
	Suppose first that $2k - 3\le p$; we claim $|R|\ge2k - 3$.
	If $k=2$, say $A=\{a,b\}$ with $a\neq b$, then $R=\{a+b\}$ has size
	$1=2k-3$ and the claim holds; assume from now on that $k\ge3$.
	Suppose, for contradiction, that $|R|\le2k - 4$; pick $C\supseteq R$
	with $|C| = 2k - 4\le p -
		1$ and set
	\[
		f(x, y)=(x-y)\prod_{c\in C}(x+y-c),
	\]
	a polynomial over $\mathbb{F}_p$ of total degree $1 + 2k - 4 = 2k -
		3$. Its top homogeneous component is $(x - y)(x + y)^{2k-4}$, so the
	coefficient of $x^{k-1}y^{k-2}$ in $f$ equals the coefficient of
	$x^{k-1}y^{k-2}$ in $(x - y)(x + y)^{2k-4}$, namely
	\[
		\binom{2k-4}{k-2}-\binom{2k-4}{k-3}
		=\binom{2k-4}{k-2}\left(1-\frac{k-2}{k-1}\right)
		=\frac{1}{k-1}\binom{2k-4}{k-2}.
	\]
	The difference is a positive integer $M$ (it is a difference of
	integers, and it is positive by the display), and $M\not\equiv0\pmod{p}$: the display
	gives $\binom{2k-4}{k-2} = (k - 1)M$, so $p\mid M$ would give
	$p\mid\binom{2k-4}{k-2}$, hence also $p\mid(2k - 4)! =
		\binom{2k-4}{k-2}\big((k - 2)!\,(k - 2)!\big)$, which is impossible
	because $2k - 4\le p - 1$. Take $S_1 = S_2 = A$, so $|S_1| = k
		> k - 1$ and $|S_2| = k > k - 2$; the total degree $2k - 3$ equals
	$(k - 1) + (k - 2)$, and Theorem~\ref{thm:12.1.2} produces $(a,
		b)\in A\times A$ with $f(a, b)\neq0$. But then $a\neq b$ and $a +
		b\notin C$, contradicting $a + b\in R\subseteq C$.

	If instead $2k - 3 > p$, then $p$ is odd (since $p = 2$ would force $k\le2$, whence $2k - 3\le1 < p$) and $k\ge(p + 5)/2$. Choose $A'\subseteq A$
	with $|A'| = (p + 3)/2$; the case already proved, applied to $A'$, gives
	\[
		|R|\ge\#\{a+b \mid a, b\in A',\ a\neq b\}\ge\min\{p,\
		2\cdot\tfrac{p+3}{2}-3\}=p,
	\]
	so $R = \mathbb{Z}_p$ and the bound holds with equality.

\end{proof}
\subsection{The Erd{\H{o}}s--Ginzburg--Ziv Theorem}

\begin{theorem}[Erd{\H{o}}s--Ginzburg--Ziv theorem, prime case]\label{thm:12.3.1}
	For a prime $p$, among any $2p - 1$ integers, some $p$ of them have a sum
	divisible by $p$.

\end{theorem}

\begin{proof} Let the integers be $a_1, \ldots,
		a_{2p-1}$, and over $\mathbb{F}_p$ set $f_1(x) = \sum_{i=1}^{2p-1}
		a_i x_i^{p-1}$, $f_2(x) = \sum_{i=1}^{2p-1} x_i^{p-1}$. The sum of
	the degrees of the nonzero polynomials among $f_1,f_2$ is at most
	$2(p-1)<2p-1$, and $0$ is a common zero, so by
	Theorem~\ref{thm:6.10.2} there is a nonzero solution $(b_i)$. Let $I = \{i \mid b_i \neq
		0\} \neq \varnothing$; by Fermat's little theorem, $b_i^{p-1} \equiv 1$ for $i \in I$,
	so $f_2$ gives $|I| \equiv 0 \pmod{p}$; hence $|I| = p$, while $f_1$
	gives $\sum_{i \in I} a_i \equiv 0 \pmod{p}$.

\end{proof}
\begin{altproof}[Alternative proof of Theorem~\ref{thm:12.3.1}]
	Sort the residues as $r_1\le\cdots\le r_{2p-1}$ in
	$\{0,\ldots,p-1\}$. If $r_i=r_{i+p-1}$ for some $1\le i\le p-1$,
	then $r_i,\ldots,r_{i+p-1}$ are $p$ equal residues summing to $0$.
	Otherwise, each $B_i=\{r_i,r_{i+p-1}\}$ ($1\le i\le p-1$) has size
	$2$. Iterating the Cauchy--Davenport theorem
	(Theorem~\ref{thm:12.2.1}) gives
	\[
		|B_1+\cdots+B_j|\ge\min\{p,j+1\}\qquad(1\le j\le p-1):
	\]
	the case $j=1$ is $|B_1|=2$, and applying the theorem to
	$B_1+\cdots+B_j$ and $B_{j+1}$ increases the lower bound by $1$ until
	it reaches $p$. At $j=p-1$, the sumset is therefore all of
	$\mathbb{Z}/p\mathbb{Z}$. Pick elements, one from each $B_i$, whose sum is congruent to
	$-r_{2p-1}\pmod p$; since the pairs $\{i,i+p-1\}$ are disjoint, these
	$p-1$ elements together with $r_{2p-1}$ correspond to $p$ of the original
	integers whose sum is divisible by $p$.

\end{altproof}
\begin{theorem}[Erd{\H{o}}s--Ginzburg--Ziv theorem, general case]\label{thm:12.3.2}
	For every $n\in\mathbb{Z}^+$, among any $2n - 1$ integers, some $n$ of
	them have a sum divisible by $n$.

\end{theorem}

\begin{proof} We proceed by induction on the number of prime factors of $n$ (counted with multiplicity). The base case $n = 1$ is trivial, and the
	prime case is Theorem~\ref{thm:12.3.1}. It suffices to prove multiplicativity:
	if the statement holds for $n = a$ and $n = b$, it holds for $n =
		ab$. Given $2ab - 1$ integers, repeatedly extract disjoint blocks of
	$a$ elements whose sum is divisible by $a$: as long as at least $2a -
		1$ elements remain, the case $n = a$ provides such a block. After
	extracting $2b - 2$ blocks, $2ab - 1 - (2b-2)a = 2a - 1$ elements
	remain, so a $(2b-1)$-th block can still be extracted; thus we obtain
	$2b - 1$ disjoint blocks $G_1, \ldots, G_{2b-1}$ with $\sum_{j \in
			G_i} a_j = a\, s_i$. Applying the case $n = b$ to $s_1, \ldots,
		s_{2b-1}$, we obtain $b$ indices $i \in J$ satisfying $\sum_{i \in J} s_i \equiv
		0 \pmod{b}$. Then $\bigcup_{i \in J} G_i$ has $ab$ elements with sum
	$a \sum_{i \in J} s_i \equiv 0 \pmod{ab}$.

\end{proof}
\begin{theorem}\label{thm:12.3.3} For every $n\in\mathbb{Z}^+$ there
	exist $2n - 2$ integers such that no $n$ of them have a sum divisible by $n$:
	the bound of Theorem~\ref{thm:12.3.2} cannot be lowered.

\end{theorem}

\begin{proof} For $n=1$, take the empty list; it has no one-element
	subcollection. Now assume $n\ge2$ and take $n - 1$ copies of $0$ and $n - 1$
	copies of $1$. Any $n$ of these $2n - 2$ integers contain $k$ copies
	of $1$ with $1\le k\le n - 1$: choosing $n$ of the only $n - 1$
	available zeros is impossible, and so is choosing $n$ of the only
	$n - 1$ available ones. The sum is then $k\not\equiv0\pmod{n}$.

\end{proof}
\begin{definition}[Davenport constant]\label{def:12.3.1} For a finite abelian group
	$G$, written additively, the Davenport constant $D(G)$ is the least
	positive integer $\ell$ such that every sequence $g_1, \ldots,
		g_\ell$ of $\ell$ elements of $G$ has a nonempty subsequence whose
	sum is $0$.

\end{definition}
\begin{theorem}\label{thm:12.3.4}\statementitems
	\begin{enumerate}[label={\alph*)}]
		\item For every finite abelian group $G$, the Davenport constant $D(G)$ exists and
		      $D(G)\le|G|$.
		\item For every $n\in\mathbb{Z}^+$, we have $D(\mathbb{Z}_n)=n$.
	\end{enumerate}

\end{theorem}

\begin{proof} (a) Given $|G|$ elements $g_1, \ldots,
		g_{|G|}$, consider the $|G|$ partial sums $s_k = g_1 + \cdots + g_k$.
	If some $s_k = 0$, the first $k$ terms form the required subsequence.
	Otherwise, the $s_k$ take at most $|G| - 1$ distinct nonzero values,
	so $s_i = s_j$ for some $i < j$, and then $g_{i+1} + \cdots + g_j =
		0$. (b) Part (a) gives $D(\mathbb{Z}_n)\le n$. The sequence
	consisting of $n - 1$ copies of $1\in\mathbb{Z}_n$ has no nonempty
	zero-sum subsequence: every nonempty subsequence has sum $k$ with
	$1\le k\le n - 1$, which is nonzero modulo $n$. Hence
	$D(\mathbb{Z}_n)\ge n$.

\end{proof}
\subsection{Schnirelmann Density and Finite Additive Bases}

\begin{definition}[Schnirelmann density]\label{def:12.4.1} For
	$A\subseteq\mathbb{Z}_{\ge0}$ with $0\in A$, set $A(n)=\#\{a\in
		A \mid 1\le a\le n\}$ for $n\in\mathbb{Z}^+$ and
	$\sigma(A)=\inf_{n\in\mathbb{Z}^+}A(n)/n$. Here $\sigma$
	takes a set as its argument; it is unrelated to the divisor-sum function
	$\sigma(n)$ of Definition~\ref{def:5.1.2}.

\end{definition}
\begin{theorem}[Schnirelmann's inequality]\label{thm:12.4.1} If
	$A,B\subseteq\mathbb{Z}_{\ge0}$ with $0\in A\cap B$, then $\sigma(A+B)\ge\sigma(A)+\sigma(B)-\sigma(A)\sigma(B)$.

\end{theorem}

\begin{proof} Put $\alpha=\sigma(A)$ and
	$\beta=\sigma(B)$, and fix $n\ge1$. If $A(n)=0$, then
	$\alpha\le A(n)/n=0$, so $\alpha=0$; moreover $B\subseteq A+B$
	because $0\in A$, and therefore
	$(A+B)(n)\ge B(n)\ge\beta n=(\alpha+\beta-\alpha\beta)n$.
	Hence we may assume $k=A(n)\ge1$. Set $a_0=0$ and list
	$A\cap[1,n]$ as $a_1<\cdots<a_k$. Each
	$a_i\in A+B$ for $1\le i\le k$ (as $0\in B$); in the gap after $a_i$ of length
	$\ell_i$ (with $\ell_i=a_{i+1}-a_i-1$ for $0\le i\le k-1$, including
	the initial gap $\ell_0=a_1-1$, and $\ell_k=n-a_k$), the numbers $a_i+b$ with $b\in
		B$ and $1\le b\le\ell_i$ lie in $(A+B)\cap[1,n]$ and have not yet been counted; their number is
	at least $\beta\ell_i$. These gaps satisfy
	$\sum_{i=0}^{k}\ell_i=n-A(n)$, so summing over $i=0,\ldots,k$ gives $(A+B)(n)\ge A(n)+\beta(n-A(n))\ge\alpha
		n(1-\beta)+\beta n=(\alpha+\beta-\alpha\beta)n$; divide by $n$ and
	take the infimum.\qedhere

\end{proof}
\begin{theorem}\label{thm:12.4.2}
	If $0\in A\subseteq\mathbb{Z}_{\ge0}$ and $\sigma(A)>0$, then some
	finite iterated sumset $hA$ equals $\mathbb{Z}_{\ge0}$.

\end{theorem}

\begin{proof} Put $\alpha=\sigma(A)>0$. Repeated
	application of Theorem~\ref{thm:12.4.1} gives $\sigma(kA)\ge1-(1-\alpha)^k$ by
	induction. Choose $k$ so that this quantity exceeds $1/2$, and put
	$B=kA$. For every $n\ge1$, we have $B(n)>n/2$, hence
	$|B\cap[0,n]|=B(n)+1>(n+1)/2$. The reflected set $n-(B\cap[0,n])$
	has the same size, so these two subsets of $[0,n]$ intersect. Thus $n=b_1+b_2$ with
	$b_1,b_2\in B$, so $2B=2kA=\mathbb{Z}_{\ge0}$. Taking $h=2k$ proves
	the claim.

\end{proof}
\subsection{The Partition Function}

\begin{definition}[partition function]\label{def:12.5.1} For
	$n\ge1$, let $p(n)$ be the number of ways to write $n$ as an unordered
	sum of positive integers. The empty partition is the unique partition
	of $0$, so set $p(0)=1$; set $p(m)=0$ for $m<0$.

\end{definition}
\begin{theorem}[Euler's generating function]\label{thm:12.5.1} For $N\in\mathbb{Z}^+$,
	put $P_N(x)=\prod_{k=1}^{N}(1-x^k)^{-1}$. As formal power series,
	\[
		\sum_{n\ge0}p(n)x^n=\prod_{k\ge1}\frac{1}{1-x^k},
	\]
	where the infinite product denotes the coefficientwise limit of the
	partial products $P_N(x)$. Moreover, for every real $x$ with $0\le x<1$,
	the series $P(x):=\sum_{n\ge0}p(n)x^n$ converges, $P_N(x)$ increases to
	$P(x)$ as $N\to\infty$, and
	\[
		\ln P(x)=\sum_{k\ge1}\bigl(-\ln(1-x^k)\bigr).
	\]

\end{theorem}

\begin{proof} Let $p_N(n)$ count the partitions of $n$ with all parts at
	most $N$, with $p_N(0)=1$. Expanding the finitely many geometric series gives
	\[
		P_N(x)=\sum_{n\ge0}p_N(n)x^n,
	\]
	both formally and for real $0\le x<1$, where the left-hand side is then a
	product of finitely many absolutely convergent series: the coefficient of
	$x^n$ counts the tuples $(j_1,\ldots,j_N)$ of nonnegative integers with
	$\sum_{k=1}^Nkj_k=n$, where $j_k$ is the multiplicity of the part $k$.
	We have $0\le p_N(n)\le p(n)$, and $p_N(n)=p(n)$ whenever $n\le N$.
	Thus each coefficient stabilizes at $p(n)$, proving the formal identity.

	Now fix a real $x$ with $0\le x<1$. Every factor $(1-x^k)^{-1}$ is at least
	$1$, so $P_N(x)$ increases with $N$. Since $1+a\le e^{a}$ for real $a$ and
	$(1-x^k)^{-1}=1+\frac{x^k}{1-x^k}\le1+\frac{x^k}{1-x}$,
	\[
		P_N(x)\le\exp\Bigl(\sum_{k=1}^{N}\frac{x^k}{1-x}\Bigr)
		\le\exp\Bigl(\frac{x}{(1-x)^2}\Bigr).
	\]
	All terms below are nonnegative, so
	\[
		\sum_{n=0}^{N}p(n)x^n\le\sum_{n\ge0}p_N(n)x^n=P_N(x)
		\le\exp\Bigl(\frac{x}{(1-x)^2}\Bigr).
	\]
	Hence the partial sums of $\sum_{n\ge0}p(n)x^n$ are bounded, and the series
	converges to a limit $P(x)$. The inequality $p_N(n)\le p(n)$ now gives
	$P_N(x)\le P(x)$, while the left-hand inequality above gives
	$\sum_{n=0}^{N}p(n)x^n\le P_N(x)$; letting $N\to\infty$ shows that
	$P_N(x)\to P(x)$. Finally,
	$\ln P_N(x)=\sum_{k=1}^{N}\bigl(-\ln(1-x^k)\bigr)$, and continuity of $\ln$
	on $(0,\infty)$ gives the last identity.

\end{proof}
\begin{theorem}[Euler's partition theorem]\label{thm:12.5.2} For every
	$n\in\mathbb{Z}_{\ge0}$, the number of partitions of $n$ into distinct
	parts equals the number of partitions of $n$ into odd parts.

\end{theorem}

\begin{proof} We work coefficientwise with formal power series;
	for any fixed degree, only finitely many factors contribute.
	The generating functions for the two types of partitions are
	$\prod_{k\ge1}(1+x^k)$ and $\prod_{k\ge1}\tfrac{1}{1-x^{2k-1}}$. But
	$\prod_{k\ge1}(1+x^k)=\prod_{k\ge1}\dfrac{1-x^{2k}}{1-x^k}=\dfrac{\prod_{k\ge1}(1-x^{2k})}{\prod_{k\ge1}(1-x^k)}=\prod_{k\ge1}\dfrac{1}{1-x^{2k-1}}$,
	since the even factors cancel. Equal generating functions have equal
	coefficients.\qedhere

\end{proof}
\begin{altproof}[Alternative proof of Theorem~\ref{thm:12.5.2}] Define a map from partitions into odd parts to partitions into distinct parts as follows: for each odd part $m$ occurring with
	multiplicity $c_m\ge0$, write $c_m = \sum_{j}2^{e_j}$ in binary and
	replace the $c_m$ copies of $m$ by the parts $2^{e_j}m$. The
	resulting parts are distinct: if $2^{e}m = 2^{e'}m'$ with $m, m'$
	odd, comparing the exponents of $2$ on both sides gives $e = e'$ and
	then $m = m'$. Conversely, map a partition into distinct parts to a
	partition into odd parts by writing each part uniquely as $2^{e}m$
	with $m$ odd and replacing it by $2^{e}$ copies of $m$. The two maps
	are mutually inverse: splitting $2^{e}m$ into $2^{e}$ copies of $m$
	and then merging them again according to the binary expansion of the multiplicity recovers $2^{e}m$, and merging the $c_m$ copies of $m$
	into the parts $2^{e_j}m$ and then re-splitting recovers the $c_m$
	copies, since binary expansions are unique. Hence the two families of
	partitions are in bijection and have the same cardinality.

\end{altproof}
\begin{theorem}[Euler's pentagonal number theorem]\label{thm:12.5.3} As
	formal power series,
	\[
		\prod_{k\ge1}(1-x^k)=\sum_{j=-\infty}^{\infty}(-1)^j
		x^{j(3j-1)/2}=1+\sum_{j\ge1}(-1)^j\big(x^{j(3j-1)/2}+x^{j(3j+1)/2}\big).
	\]

\end{theorem}

\begin{proof} The constant term on each side is $1$, corresponding to
	the empty partition. Fix $n\ge1$. Expanding the product, we find that the coefficient
	of $x^n$ is $\sum(-1)^r$ over partitions of $n$ into $r$ distinct parts;
	this is a finite sum. Given such a partition with $r$ parts, let $s$ be its smallest
	part and $t$ the length of the top staircase (the
	longest run of consecutive integers $a,a-1,\ldots$ descending from the
	largest part $a$). Define an involution as follows. If $s\le t$,
	delete the smallest part $s$ and add $1$ to each of the $s$ largest
	parts. If $s>t$, subtract $1$ from each of the $t$ largest
	parts and adjoin a new smallest part $t$.

	Whenever the first move is defined, increasing the $s$ largest
	remaining parts preserves distinctness. Differences within the
	modified block are unchanged, and any gap to the remaining parts
	increases. Since $s\le t$, the first $s$ parts remain consecutive.
	If another part follows them, the gap to it is at least $2$.
	Thus the new top staircase has length exactly $s$, while the new
	smallest part is greater than $s$. The second move therefore
	reverses the first.

	For the second move, if $t<r$, maximality of the top staircase
	means that the gap between the $t$-th and $(t+1)$-th largest parts
	is at least $2$. After lowering the first $t$ parts, this gap
	remains at least $1$. The new part $t$ is smaller than every
	remaining part. If $t=r$, the lowered smallest part is $s-1$,
	so the only possible collision is $s-1=t$, the exceptional
	configuration described below. Outside that configuration, the
	new smallest part is $t$ and the new top staircase has length
	at least $t$, so the first move reverses the second.

	Both moves preserve the sum and change the number of parts by exactly $1$; hence they reverse the sign $(-1)^r$.

	A move fails in exactly two boundary configurations. First, when
	$s\le t$, the first move fails if and only if the top staircase
	contains all parts and $s=t=k$. Then the parts are exactly
	$k,k+1,\ldots,2k-1$ for some $k\ge1$; after deleting the smallest part,
	only $s-1$ parts remain, so there are not enough parts to increase $s$ of them.
	This family has $r=k$ parts and sum $\tfrac{k(3k-1)}{2}$. Second, when
	$s>t$, the second move fails if and only if the top staircase
	contains all parts and $s=t+1$. With $k=t$, the parts are
	exactly $k+1,\ldots,2k$; subtracting $1$ from the top $t$ parts
	and adjoining a new smallest part $t=s-1$ would duplicate the
	lowered smallest part. This family has $r=k$ parts and sum
	$\tfrac{k(3k+1)}{2}$. Every other partition is
	paired with a distinct partner and cancels, so only these two
	exceptional families survive, each contributing $(-1)^k$. This gives
	the stated signed sum over generalized pentagonal numbers.

\end{proof}
\begin{theorem}[Euler's recurrence for $p(n)$]\label{thm:12.5.4} If $n\in\mathbb{Z}^+$, then
	$p(n)=\sum_{j\ge1}(-1)^{j-1}\Big(p\big(n-\tfrac{j(3j-1)}{2}\big)+p\big(n-\tfrac{j(3j+1)}{2}\big)\Big)$.

\end{theorem}

\begin{proof} By Theorem~\ref{thm:12.5.1},
	$\big(\sum_{m\ge0}p(m)x^m\big)\prod_{k\ge1}(1-x^k)=1$. Replace the
	product by its pentagonal expansion (Theorem~\ref{thm:12.5.3}) and compare the
	coefficient of $x^n$ for $n\ge1$ (which is $0$ on the right):
	$p(n)+\sum_{j\ge1}(-1)^j\big(p(n-\tfrac{j(3j-1)}{2})+p(n-\tfrac{j(3j+1)}{2})\big)=0$,
	which rearranges to the claimed recurrence.

\end{proof}
\begin{definition}[Ferrers diagrams and conjugate partitions]
	\label{def:12.5.2} A partition
	$\lambda$ with parts $\lambda_1\ge\lambda_2\ge\cdots\ge\lambda_r\ge1$
	is represented by its Ferrers diagram, the left-justified array with
	$\lambda_i$ dots in row $i$.
	The conjugate of a partition $\lambda$ is the sequence $\lambda'$ given by
	$\lambda'_j=\#\{i \mid \lambda_i\ge j\}$ for
	$1\le j\le\lambda_1$; thus trailing zeros are omitted. Its Ferrers
	diagram is the transpose of the diagram of $\lambda$.

\end{definition}
\begin{theorem}\label{thm:12.5.5} For every $n\ge1$ and every partition
	$\lambda$ of $n$, the conjugate $\lambda'$ (Definition~\ref{def:12.5.2}) is again a partition of $n$; one has
	$(\lambda')' = \lambda$, and the number of parts of $\lambda'$ equals
	the largest part of $\lambda$. Consequently, for $1\le k\le n$ the
	number of partitions of $n$ with largest part $k$ equals the number
	of partitions of $n$ into exactly $k$ parts.

\end{theorem}

\begin{proof} The positive integers $\lambda'_j$ ($1\le j\le\lambda_1$)
	are weakly decreasing in $j$, and
	\[
		\sum_{j=1}^{\lambda_1}\lambda'_j
		=\sum_{j=1}^{\lambda_1}\#\{i \mid \lambda_i\ge j\}
		=\sum_{i=1}^{r}\sum_{j=1}^{\lambda_i}1
		=\sum_{i=1}^{r}\lambda_i=n,
	\]
	so $\lambda'$ is a partition of $n$. Its number of parts is
	$\max\{1\le j\le\lambda_1 \mid \lambda'_j\ge1\}
	=\#\{1\le j\le\lambda_1 \mid \lambda_1\ge j\} =
		\lambda_1$, the largest part of $\lambda$. Applying the definition
	twice gives
	\[
		((\lambda')')_i
		=\#\{j \mid 1\le j\le\lambda_1,\ \lambda'_j\ge i\}
		=\#\{j \mid 1\le j\le\lambda_i\}=\lambda_i,
	\]
	because $\lambda'_j=\#\{l \mid \lambda_l\ge j\}\ge i$ holds if and only if
	$\lambda_i\ge j$. Hence conjugation is an involution exchanging the
	largest part with the number of parts, which gives the stated
	equality of counts.

\end{proof}
\begin{theorem}\label{thm:12.5.6} For every $n\in\mathbb{Z}^+$,
	\[
		p(n)<e^{\pi\sqrt{2n/3}}.
	\]

\end{theorem}

\begin{proof} Fix a real $x$ with $0<x<1$, and let $P(x)=\sum_{m\ge0}p(m)x^m$ as
	in Theorem~\ref{thm:12.5.1}. All terms of this series are nonnegative, so
	$p(n)x^n\le P(x)$. For $0\le y<1$, we have
	$-\ln(1-y)=\sum_{m\ge1}y^m/m$. Theorem~\ref{thm:12.5.1} therefore gives
	\[
		\ln P(x)=\sum_{k\ge1}\sum_{m\ge1}\frac{x^{km}}{m}
		=\sum_{m\ge1}\frac1m\sum_{k\ge1}x^{km}
		=\sum_{m\ge1}\frac1m\cdot\frac{x^m}{1-x^m},
	\]
	where the order of summation may be interchanged because all terms are
	nonnegative. For $m\ge1$, since $x^j\ge x^{m-1}$ for $0\le j\le m-1$,
	\[
		1-x^m=(1-x)(1+x+\cdots+x^{m-1})\ge(1-x)\,m\,x^{m-1},
	\]
	so $\frac{x^m}{1-x^m}\le\frac{x}{m(1-x)}$. By Theorem~\ref{thm:5.7.3},
	\[
		\ln P(x)\le\frac{x}{1-x}\sum_{m\ge1}\frac1{m^2}=\frac{\pi^2}{6}\cdot\frac{x}{1-x}.
	\]
	Now write $x=e^{-t}$ with $t>0$. Since $e^{t}>1+t$, we have
	$\frac{x}{1-x}=\frac{1}{e^{t}-1}<\frac1t$, and therefore
	\[
		\ln p(n)\le\ln P(x)+nt<\frac{\pi^2}{6t}+nt.
	\]
	Choosing $t=\pi/\sqrt{6n}$ makes both terms on the right equal to
	$\pi\sqrt{n/6}$, so $\ln p(n)<2\pi\sqrt{n/6}=\pi\sqrt{2n/3}$.

\end{proof}
\subsection{Waring's Problem: The Case \texorpdfstring{$k=2$}{k=2}}

\begin{definition}[the functions $g(k)$ and $G(k)$]
	\label{def:12.6.1} For an integer $k\ge1$,
	let $g(k)$ be the least positive integer $s$, if one exists, such that
	every nonnegative integer is a sum of $s$ $k$-th powers of
	nonnegative integers. Put $g(k)=\infty$ if no such $s$ exists.

	Let $G(k)$ be the least positive integer $s$, if one exists, such that
	all sufficiently large nonnegative integers are sums of $s$ $k$-th
	powers of nonnegative integers; that is, there is an $N$ such that every
	$n\ge N$ is representable. Put $G(k)=\infty$ if no such $s$ exists.

\end{definition}
\begin{theorem}\label{thm:12.6.1} Every nonnegative integer is a
	sum of four squares, while there exist arbitrarily large integers
	that are not sums of three squares; in the notation of Definition~\ref{def:12.6.1}, $g(2)
		= G(2) = 4$.

\end{theorem}

\begin{proof} Lagrange's four-square theorem (Theorem~\ref{thm:8.1.10}) gives $g(2)\le4$, hence also $G(2)\le4$. A square
	is congruent to $0, 1$, or $4\pmod{8}$, and the sums of three
	elements of $\{0, 1, 4\}$ cover only the residues $0, 1, 2, 3, 4, 5,
		6\pmod{8}$; hence no integer $n\equiv7\pmod{8}$ is a sum of three
	squares. In particular, $7$ itself is not, so $g(2)\ge4$; and the
	class $n\equiv7\pmod{8}$ contains arbitrarily large integers, so
	$G(2)\ge4$.

\end{proof}
\end{document}
